\part*{Part II — Strategy}
\label{navigation:II}
\addcontentsline{toc}{part}{Part II — Strategy}
The model separates Business Capabilities, Business Enabling Capabilities and Enterprise Capabilities within Strategy. Follow those existing branches to move from health-service responsibilities to their enabling concerns and platform capabilities. The folder called Application Layer here remains inside Strategy: it contains Enterprise Capabilities, and its navigation label does not reclassify those capabilities as application components.

{\small Reading sources: \hyperref[element:ed342cb8656b4f80b55ad519ecbd2d3a]{[CT1] Care Access \& Coordination}; \hyperref[element:b5a92874f8ac4509a30af5253590edc1]{[CT1] Entity Information Enablement}; \hyperref[element:15e57021490c4e738aee9009ff84a7c3]{[CT1] Governed Information Operations}.}\par
Health Service Governance illustrates a boundary that matters when reading capability descriptions: it establishes obligations, authority and compliance expectations, while operational management progresses means within that envelope and assurance evaluates evidence. Follow the capability's stated boundary and its qualified links rather than reading all neighbouring responsibilities as one owner. Proposed placement, relevance or enablement remains independently qualified even where the responsibility is canonical.

{\small Reading sources: \hyperref[element:209aada781e844e28202a2535a264c21]{[CT1] Health Service Governance}.}\par
\begin{center}\begin{tikzpicture}[>=Stealth,every node/.style={draw=Ink!40,rounded corners,align=left,font=\small,text width=5.2cm,inner sep=6pt}]
\node (source) at (0,0) {\hyperref[element:b13ee2c5a4ef4fd6be8a44795696dfc2]{[CT1] Health Service Operations Enablement}};
\node (target0) at (8,4.25) {\hyperref[element:f117ccec0bfa4f738e6ddb646b3534a2]{[CT2] Bed \& Care-Place Management}};
\draw[->] (source.east) -- (3,0) |- (target0.west);
\node[draw=none,text width=1.8cm,font=\tiny,align=center,inner sep=0pt] at (4.15,4.63) {Composition\newline{}Proposed / current};
\node (target1) at (8,2.55) {\hyperref[element:3dffe63e03124efab255204ee4e9a5d9]{[CT2] Clinic \& Practice Operations}};
\draw[->] (source.east) -- (3,0) |- (target1.west);
\node[draw=none,text width=1.8cm,font=\tiny,align=center,inner sep=0pt] at (4.15,2.93) {Composition\newline{}Proposed / current};
\node (target2) at (8,0.85) {\hyperref[element:8a99f4a6e5054d6c997bacec04805c01]{[CT2] Clinical Competency Administration}};
\draw[->] (source.east) -- (3,0) |- (target2.west);
\node[draw=none,text width=1.8cm,font=\tiny,align=center,inner sep=0pt] at (4.15,1.23) {Composition\newline{}Proposed / current};
\node (target3) at (8,-0.85) {\hyperref[element:0c27be11b4a44cb19e992309b8146a17]{[CT2] Clinical Logistics Coordination}};
\draw[->] (source.east) -- (3,0) |- (target3.west);
\node[draw=none,text width=1.8cm,font=\tiny,align=center,inner sep=0pt] at (4.15,-0.47) {Composition\newline{}Proposed / current};
\node (target4) at (8,-2.55) {\hyperref[element:0511a705d74a4f48a0ce92590ae764fa]{[CT2] Clinical Qualification Administration}};
\draw[->] (source.east) -- (3,0) |- (target4.west);
\node[draw=none,text width=1.8cm,font=\tiny,align=center,inner sep=0pt] at (4.15,-2.17) {Composition\newline{}Proposed / current};
\node (target5) at (8,-4.25) {\hyperref[element:192ca88a289544dcbf6adaf643f39969]{[CT2] Clinical Rostering}};
\draw[->] (source.east) -- (3,0) |- (target5.west);
\node[draw=none,text width=1.8cm,font=\tiny,align=center,inner sep=0pt] at (4.15,-3.87) {Composition\newline{}Proposed / current};
\end{tikzpicture}\end{center}
{\small Selected outgoing Composition assertions; arrows retain source-to-target direction. Labels retain assertion standing and record state; element standing accompanies each description. \href{structure.json}{Complete relationship qualifications and provenance} remain source-defined. No omitted edge or semantic effect is inferred.}\par
\chapter*{II.1 Capabilities}
\phantomsection
\label{navigation:II.1}
\addcontentsline{toc}{chapter}{II.1 Capabilities}

\section*{II.1.1 Enterprise Capabilities}
\phantomsection
\label{navigation:II.1.1}
\addcontentsline{toc}{section}{II.1.1 Enterprise Capabilities}

\subsection*{II.1.1.1 [CT1] Coordinated Activity \& Interaction}
\phantomsection
\label{navigation:II.1.1.1}
\addcontentsline{toc}{subsection}{II.1.1.1 [CT1] Coordinated Activity \& Interaction}

\label{element:5187c3db390e444db3eb809c02a2f1c1}
\index{[CT1] Coordinated Activity \& Interaction (II.1.1.1)}
{\small Capability · Proposed · current}\par

Enable governed interaction, event participation, boundary exchange and accountable activity progression.



\textbf{Overview}



Requests, events, work and human interaction are related ways to cause or observe change, with distinct contracts and progression responsibilities.



\textbf{Intent}



Make requests, event participation, work and human interaction progress according to their different contracts.



\textbf{Boundary}



Exchange is not execution, notification is not outcome, and interaction is not clinical authority.


{\footnotesize \href{../../../../continuum/strategy/capability/5187c3db390e444db3eb809c02a2f1c1.md}{Canonical record: 5187c3db390e444db3eb809c02a2f1c1}; full qualifications and provenance.}\par
\subsubsection*{II.1.1.1.1 [CT2] Activity \& Execution}
\phantomsection
\label{navigation:II.1.1.1.1}
\addcontentsline{toc}{subsubsection}{II.1.1.1.1 [CT2] Activity \& Execution}

\label{element:e5ff80e655084524a8dcee8e4ded788a}
\index{[CT2] Activity \& Execution (II.1.1.1.1)}
{\small Capability · Canonical · current}\par

\begin{itemize}


\item \textbf{Architectural Scope}: Define, instantiate, match, assign, dispatch, supervise, and progress coordinated operational and clinical activity units.

\item \textbf{Functional Responsibilities}:

\item Instantiating operational work items (Work Orders, To Dos, synthetic platform Tasks) in response to triggers.

\item Matching and assigning work units to eligible human or automated performers and roles based on policy rules.

\item Dispatching work units to their assigned performers with acknowledgement according to the applicable interaction contract.

\item Supervising execution progress, tracking timeouts, managing escalations, and handling activity interruptions.

\item \textbf{Decoupling Note}: Free of worker threads, thread pools, execution daemons, or specific workflow runtime engines.

\end{itemize}


Activity \& Execution provides reusable instantiation, matching, dispatch and supervision for governed activity units. Human Work Orders, human decision To Dos and automated Synthetic Tasks retain different performer and authority semantics while sharing a need for accountable progression. This reuse supports operational coordination without making a runtime engine the definition of the work.



\textbf{Boundary}



Coordination does not confer clinical work allocation or judgement. The capability specifies neither worker threads nor a particular execution engine.


{\footnotesize \href{../../../../continuum/strategy/capability/e5ff80e655084524a8dcee8e4ded788a.md}{Canonical record: e5ff80e655084524a8dcee8e4ded788a}; full qualifications and provenance.}\par
\paragraph*{II.1.1.1.1.1 [CT3] Entity-Relative Operational Progression}
\phantomsection
\label{navigation:II.1.1.1.1.1}
\addcontentsline{toc}{paragraph}{II.1.1.1.1.1 [CT3] Entity-Relative Operational Progression}

\label{element:69c82474fb534152b151c92f706197d0}
\index{[CT3] Entity-Relative Operational Progression (II.1.1.1.1.1)}
{\small Capability · Proposed · current}\par

Coordinate serial, entity-relative operational inputs and their interested behaviours.



\textbf{Overview}



Patient, practitioner and bed concerns recur across different business contexts; each can need a coherent progression frame without one central entity service.



\textbf{Intent}



Maintain a coherent operational reference frame as one entity reacts to successive inputs.



\textbf{Boundary}



Interest fan-out is distinct from intra-entity concurrency; entity progression does not acquire originating facts.


{\footnotesize \href{../../../../continuum/strategy/capability/69c82474fb534152b151c92f706197d0.md}{Canonical record: 69c82474fb534152b151c92f706197d0}; full qualifications and provenance.}\par
\paragraph*{II.1.1.1.1.2 [CT3] Governed Activity Execution Control}
\phantomsection
\label{navigation:II.1.1.1.1.2}
\addcontentsline{toc}{paragraph}{II.1.1.1.1.2 [CT3] Governed Activity Execution Control}

\label{element:1e351cdfdc1045e384c3bd768eed8e80}
\index{[CT3] Governed Activity Execution Control (II.1.1.1.1.2)}
{\small Capability · Proposed · current}\par

Instantiate, dispatch, supervise and progress authorised operational undertakings with explicit outcome standing.



\textbf{Overview}



Reusable execution control respects work-unit archetype, participation and uncertain effects rather than equating dispatch with completion.



\textbf{Intent}



Progress accountable operational work without equating dispatch, execution and successful outcome.



\textbf{Boundary}



Does not own underlying clinical-work meaning, clinical decisions or originator business authority.


{\footnotesize \href{../../../../continuum/strategy/capability/1e351cdfdc1045e384c3bd768eed8e80.md}{Canonical record: 1e351cdfdc1045e384c3bd768eed8e80}; full qualifications and provenance.}\par
\paragraph*{II.1.1.1.1.3 [CT3] Operational Fulfilment Definition}
\phantomsection
\label{navigation:II.1.1.1.1.3}
\addcontentsline{toc}{paragraph}{II.1.1.1.1.3 [CT3] Operational Fulfilment Definition}

\label{element:2b52c0de17b146da962f303e99f90a4a}
\index{[CT3] Operational Fulfilment Definition (II.1.1.1.1.3)}
{\small Capability · Proposed · current}\par

Define the context requirements, expected outcomes and housekeeping of an operational undertaking.



\textbf{Overview}



The semantic execution envelope applies to one-step and multi-step activity and stays distinct from dispatch and supervision.



\textbf{Intent}



Keep the semantic promise of an undertaking separate from the machinery that executes it.



\textbf{Boundary}



A fulfilment definition does not execute, orchestrate or choreograph its activity.


{\footnotesize \href{../../../../continuum/strategy/capability/2b52c0de17b146da962f303e99f90a4a.md}{Canonical record: 2b52c0de17b146da962f303e99f90a4a}; full qualifications and provenance.}\par
\subsubsection*{II.1.1.1.2 [CT2] Event \& Subscription}
\phantomsection
\label{navigation:II.1.1.1.2}
\addcontentsline{toc}{subsubsection}{II.1.1.1.2 [CT2] Event \& Subscription}

\label{element:b768b129971e4bf5bccc03d33d2fcb4e}
\index{[CT2] Event \& Subscription (II.1.1.1.2)}
{\small Capability · Canonical · current}\par

\begin{itemize}


\item \textbf{Architectural Scope}: Publish, ingest, correlate, filter, and syndicate events and subscription-driven changes across the enterprise landscape.

\item \textbf{Functional Responsibilities}:

\item Ingesting operational and clinical event triggers from external systems and internal state changes.

\item Managing topic- and criteria-based event subscriptions for internal and external consumers.

\item Evaluating dynamic subscription filters to match event payloads with interested subscriber criteria.

\item Delivering asynchronous event notifications reliably to registered subscribers.

\item \textbf{Decoupling Note}: Defined strictly in terms of event publishing, criteria matching, and notification delivery; entirely independent of broker topics, partitions, or queues.

\end{itemize}


Event \& Subscription connects significant changes with participants whose registered interests match them. It makes event ingestion, correlation, interest criteria and notification reusable across clinical and operational contexts. The architectural meaning is the relationship between a change and an interested consumer, rather than the broker constructs used to carry it.



\textbf{Boundary}



Receiving a notification does not determine the consumer's business response; topics and queues do not define the capability.


{\footnotesize \href{../../../../continuum/strategy/capability/b768b129971e4bf5bccc03d33d2fcb4e.md}{Canonical record: b768b129971e4bf5bccc03d33d2fcb4e}; full qualifications and provenance.}\par
\subsubsection*{II.1.1.1.3 [CT2] Interaction \& Experience}
\phantomsection
\label{navigation:II.1.1.1.3}
\addcontentsline{toc}{subsubsection}{II.1.1.1.3 [CT2] Interaction \& Experience}

\label{element:4ec9757b1c6947b4b70b0f1159b760d3}
\index{[CT2] Interaction \& Experience (II.1.1.1.3)}
{\small Capability · Canonical · current}\par

\begin{itemize}


\item \textbf{Architectural Scope}: Provide contextual interaction between human actors, Harmonia, and enterprise presentation services, including visualization, alerting, and collaboration integration.

\item \textbf{Functional Responsibilities}:

\item Mediating presentation interactions via dedicated backend-for-frontend gateways decoupled from internal databases.

\item Rendering clinical timelines, provider directory management interfaces, and operational console dashboards.

\item Providing contextual notifications and alerts to human users across web and mobile channels.

\item Projecting longitudinal record context into secure multidisciplinary collaboration spaces.

\item \textbf{Boundary Note}: Does not absorb the healthcare semantics of Clinical Collaboration or medical documentation. Consumes established identity and security context; does not own authentication, consent, or clinical truth.

\end{itemize}


Interaction \& Experience makes contextual information and governed actions available to human participants through presentation and collaboration integration. Clinical timelines, directory administration and operational dashboards have different purposes, but share a need to retain established identity, security and information context. The recurring responsibility is usable interaction within those boundaries.



\textbf{Boundary}



Presentation does not own consent, authentication or clinical truth, and collaboration projection does not absorb clinical documentation semantics.


{\footnotesize \href{../../../../continuum/strategy/capability/4ec9757b1c6947b4b70b0f1159b760d3.md}{Canonical record: 4ec9757b1c6947b4b70b0f1159b760d3}; full qualifications and provenance.}\par
\subsubsection*{II.1.1.1.4 [CT2] Interoperability \& Exchange}
\phantomsection
\label{navigation:II.1.1.1.4}
\addcontentsline{toc}{subsubsection}{II.1.1.1.4 [CT2] Interoperability \& Exchange}

\label{element:0f2e5222874949d8b0cce234941ecb20}
\index{[CT2] Interoperability \& Exchange (II.1.1.1.4)}
{\small Capability · Canonical · current}\par

\begin{itemize}


\item \textbf{Architectural Scope}: Resolve destinations, transform schemas, route, transport, deliver, acknowledge, and coordinate standards-based information exchange across enterprise boundaries.

\item \textbf{Functional Responsibilities}:

\item Providing multi-protocol boundary adaptation for healthcare standards (HL7 v2, FHIR REST, DICOM, secure messaging).

\item Executing bidirectional structural and semantic transformations between legacy formats and canonical representations.

\item Ensuring reliable addressed delivery and fan-out distribution to external clinical endpoints.

\item Processing and distinguishing technical delivery acknowledgements from business acknowledgements.

\item \textbf{Boundary Rule}: Egress terminates Harmonia management of the emitted representation (Axiom AX-13).

\end{itemize}


Interoperability \& Exchange moves information across participating boundaries through destination resolution, transformation, delivery and acknowledgements. It preserves the difference between technical transmission and the recipient's business response so that an exchange can be understood under its actual contract. Standards govern the external representation while internal management semantics remain private.



\textbf{Boundary}



Egress ends Harmonia management of the emitted representation. Retained delivery evidence does not extend control over an external recipient or transfer payload authority.


{\footnotesize \href{../../../../continuum/strategy/capability/0f2e5222874949d8b0cce234941ecb20.md}{Canonical record: 0f2e5222874949d8b0cce234941ecb20}; full qualifications and provenance.}\par
\subsection*{II.1.1.2 [CT1] Governed Information Operations}
\phantomsection
\label{navigation:II.1.1.2}
\addcontentsline{toc}{subsection}{II.1.1.2 [CT1] Governed Information Operations}

\label{element:15e57021490c4e738aee9009ff84a7c3}
\index{[CT1] Governed Information Operations (II.1.1.2)}
{\small Capability · Proposed · current}\par

Maintain identifiable, meaningful, context-qualified and retrievable information through valid state progression.



\textbf{Overview}



Entity integrity, context, lifecycle, information maintenance, discovery and semantic conformance are constituent reusable abilities; none is an implementation subsystem.



\textbf{Intent}



Keep information identity, meaning, context, state and retrieval compatible across diverse healthcare uses.



\textbf{Boundary}



Information management and semantic governance do not grant originating domain authority or erase active/durable state separation.


{\footnotesize \href{../../../../continuum/strategy/capability/15e57021490c4e738aee9009ff84a7c3.md}{Canonical record: 15e57021490c4e738aee9009ff84a7c3}; full qualifications and provenance.}\par
\subsubsection*{II.1.1.2.1 [CT2] Context Management}
\phantomsection
\label{navigation:II.1.1.2.1}
\addcontentsline{toc}{subsubsection}{II.1.1.2.1 [CT2] Context Management}

\label{element:3cfd3e56c23a4943bd1b476ec248ac49}
\index{[CT2] Context Management (II.1.1.2.1)}
{\small Capability · Canonical · current}\par

\begin{itemize}


\item \textbf{Architectural Scope}: Establish, validate, propagate, and preserve the execution, clinical, security, and administrative context within which information and activities have meaning.

\item \textbf{Functional Responsibilities}:

\item Binding the subject of care, encounter, ordering provider, and healthcare service context to in-flight messages and transactions.

\item Propagating tamper-evident, attributable security context across asynchronous, multi-hop operational boundaries.

\item Ensuring context preservation during cross-protocol transformations across disparate healthcare exchange formats and protocols.

\item Detecting and rejecting requests where mandatory operational or clinical context is missing or inconsistent.

\item \textbf{Collaborative Nature}: Bounded reusable context responsibility may be realised through participation across gateways, activity execution and information-access boundaries. Participation does not transfer semantic ownership.

\end{itemize}


Context Management preserves the circumstances in which information and activity have meaning as they cross processing boundaries. Subject, encounter, provider, service and security context must remain attributable and consistent rather than be reconstructed from convenient local assumptions. Its reuse is especially significant in multi-hop activity where a syntactically valid message can otherwise lose the context required for legitimate interpretation.



\textbf{Boundary}



Context participation does not transfer semantic ownership. Generic context management and assurance-specific context semantics contribute together without allocating a new central Context Service.


{\footnotesize \href{../../../../continuum/strategy/capability/3cfd3e56c23a4943bd1b476ec248ac49.md}{Canonical record: 3cfd3e56c23a4943bd1b476ec248ac49}; full qualifications and provenance.}\par
\subsubsection*{II.1.1.2.2 [CT2] Information Management}
\phantomsection
\label{navigation:II.1.1.2.2}
\addcontentsline{toc}{subsubsection}{II.1.1.2.2 [CT2] Information Management}

\label{element:43051efd406b4040bff8e794b7a631e0}
\index{[CT2] Information Management (II.1.1.2.2)}
{\small Capability · Canonical · current}\par

\begin{itemize}


\item \textbf{Architectural Scope}: Validate, represent, persist, retrieve, version, and manage governed clinical, administrative, and operational information assets.

\item \textbf{Functional Responsibilities}:

\item Validating inbound information structures against canonical information constraints.

\item Providing authoritative durable persistence for historical health records and administrative registries.

\item Managing active distributed working memory representations for low-latency retrieval.

\item Governing document versioning, supersession, amendments, and retractions without data loss.

\item \textbf{Architectural Principle}: Enforces the foundational separation between \emph{active distributed state access} and \emph{authoritative durable state preservation} (Axiom AX-05).

\end{itemize}


Information Management provides the recurring ability to represent and maintain governed information through validation, preservation, retrieval and versioned change. It accommodates both responsive active use and authoritative durable preservation while keeping those responsibilities separate. This distinction allows availability and working state to support care without turning a cache into durable truth.



\textbf{Boundary}



Engine capabilities provide machinery under Harmonia semantics. Persistence success does not establish information authority or make the platform the originator of the represented facts.


{\footnotesize \href{../../../../continuum/strategy/capability/43051efd406b4040bff8e794b7a631e0.md}{Canonical record: 43051efd406b4040bff8e794b7a631e0}; full qualifications and provenance.}\par
\paragraph*{II.1.1.2.2.1 [CT3] Active Information Availability}
\phantomsection
\label{navigation:II.1.1.2.2.1}
\addcontentsline{toc}{paragraph}{II.1.1.2.2.1 [CT3] Active Information Availability}

\label{element:bb711f4d8fb74af2917869b9b05f7f84}
\index{[CT3] Active Information Availability (II.1.1.2.2.1)}
{\small Capability · Proposed · current}\par

Maintain responsive, reconstructable active information and coordinate its safe use.



\textbf{Overview}



Active representation, observation and convergence support availability while preserving an explicit trust basis for use.



\textbf{Intent}



Provide responsive managed use without promoting active representation into durable truth.



\textbf{Boundary}



An active generation is not an authoritative version; failure to converge does not establish valid active state.


{\footnotesize \href{../../../../continuum/strategy/capability/bb711f4d8fb74af2917869b9b05f7f84.md}{Canonical record: bb711f4d8fb74af2917869b9b05f7f84}; full qualifications and provenance.}\par
\paragraph*{II.1.1.2.2.2 [CT3] Authoritative Durable State Establishment}
\phantomsection
\label{navigation:II.1.1.2.2.2}
\addcontentsline{toc}{paragraph}{II.1.1.2.2.2 [CT3] Authoritative Durable State Establishment}

\label{element:0d5bd68ff50f410aaceed05377e1a4df}
\index{[CT3] Authoritative Durable State Establishment (II.1.1.2.2.2)}
{\small Capability · Proposed · current}\par

Atomically establish durable managed state and its authoritative version progression.



\textbf{Overview}



Durable acceptance gives managed information an explicit committed state from which applicable active representations can be reconstructed.



\textbf{Intent}



Give managed change an explicit committed state and authoritative progression basis.



\textbf{Boundary}



Commit is distinct from semantic source authority and from successful active convergence; applications do not gain a raw persistence path.


{\footnotesize \href{../../../../continuum/strategy/capability/0d5bd68ff50f410aaceed05377e1a4df.md}{Canonical record: 0d5bd68ff50f410aaceed05377e1a4df}; full qualifications and provenance.}\par
\paragraph*{II.1.1.2.2.3 [CT3] Binary Content Association \& Resolution}
\phantomsection
\label{navigation:II.1.1.2.2.3}
\addcontentsline{toc}{paragraph}{II.1.1.2.2.3 [CT3] Binary Content Association \& Resolution}

\label{element:5314959509dc4b558711ec169ddee5a0}
\index{[CT3] Binary Content Association \& Resolution (II.1.1.2.2.3)}
{\small Capability · Proposed · current}\par

Maintain content associations and custody knowledge sufficient for governed payload resolution, retrieval and assembly.



\textbf{Overview}



Document attachments, diagnostic media and operational artefacts share a need to locate usable content while distinguishing logical meaning from physical possession.



\textbf{Intent}



Make logical information and its usable payloads consistently associable while retaining custody boundaries.



\textbf{Boundary}



No universal binary custody, storage topology or persistence authority is implied.


{\footnotesize \href{../../../../continuum/strategy/capability/5314959509dc4b558711ec169ddee5a0.md}{Canonical record: 5314959509dc4b558711ec169ddee5a0}; full qualifications and provenance.}\par
\subsubsection*{II.1.1.2.3 [CT2] Managed Entity \& Relationship}
\phantomsection
\label{navigation:II.1.1.2.3}
\addcontentsline{toc}{subsubsection}{II.1.1.2.3 [CT2] Managed Entity \& Relationship}

\label{element:910634fb1621416c818a022ab0b5d3a1}
\index{[CT2] Managed Entity \& Relationship (II.1.1.2.3)}
{\small Capability · Canonical · current}\par

\begin{itemize}


\item \textbf{Architectural Scope}: Identify, classify, relate, and maintain the structural integrity of managed healthcare and administrative entities throughout the enterprise landscape.

\item \textbf{Functional Responsibilities}:

\item Uniquely identifying core domain entities (Persons, Practitioners, Roles, Organisations, Locations, Services, Devices).

\item Resolving and maintaining multi-authority identifier systems and aliases without destructive overwriting.

\item Representing and traversing multi-dimensional relationships (affiliations, organizational hierarchies, physical location trees, care team memberships).

\item Maintaining referential integrity across interconnected entity graphs.

\item \textbf{Anti-Responsibilities}: Does not own physical database relational tables or foreign key constraints (Domain 04/07).

\end{itemize}


This reusable capability keeps managed healthcare entities identifiable and their relationships structurally intelligible. It supplies common identity, alias and relationship behaviour recurring in persons, practitioners, organisations, locations and devices. Reuse concerns those architectural responsibilities, allowing the healthcare-specific meaning of each entity to remain with its governing context.



\textbf{Boundary}



Relationship integrity is not a physical database design. Correlating entities does not establish originating clinical identity authority or a new standalone Correlation capability.


{\footnotesize \href{../../../../continuum/strategy/capability/910634fb1621416c818a022ab0b5d3a1.md}{Canonical record: 910634fb1621416c818a022ab0b5d3a1}; full qualifications and provenance.}\par
\subsubsection*{II.1.1.2.4 [CT2] Managed State \& Lifecycle}
\phantomsection
\label{navigation:II.1.1.2.4}
\addcontentsline{toc}{subsubsection}{II.1.1.2.4 [CT2] Managed State \& Lifecycle}

\label{element:6296abeee5074b9e80f21aea3a297e42}
\index{[CT2] Managed State \& Lifecycle (II.1.1.2.4)}
{\small Capability · Canonical · current}\par

\begin{itemize}


\item \textbf{Architectural Scope}: Govern the current and historical state, effective temporal periods, state transition validations, and lifecycle progression of managed integration artifacts and entities.

\item \textbf{Functional Responsibilities}:

\item Defining and enforcing valid state transition state machines (e.g., \texttt{\seqsplit{REQUESTED}} \$\textbackslash{}to\$ \texttt{\seqsplit{VALIDATED}} \$\textbackslash{}to\$ \texttt{\seqsplit{IN\_PROGRESS}} \$\textbackslash{}to\$ \texttt{\seqsplit{COMPLETED}}).

\item Maintaining temporal validity intervals (effective start and end dates) for relationships, credentials, and states.

\item Supporting non-destructive state progression, preserving historical state versions for longitudinal review.

\item Rejecting illegal or out-of-order state transitions while accommodating idempotent replay.

\item \textbf{Anti-Responsibilities}: Does not define thread-level mutexes or database transaction isolation levels.

\end{itemize}


Managed State \& Lifecycle provides reusable responsibility for valid progression, temporal applicability and non-destructive change. A directory relationship, clinical document and operational work item have different business states, yet all need an intelligible distinction between current standing, historical state and invalid transitions. The capability captures that recurrence without replacing each subject's lifecycle meaning.



\textbf{Boundary}



Business lifecycle integrity does not prescribe locks, database isolation or worker implementation, and illustrative state names do not define every subject's transitions.


{\footnotesize \href{../../../../continuum/strategy/capability/6296abeee5074b9e80f21aea3a297e42.md}{Canonical record: 6296abeee5074b9e80f21aea3a297e42}; full qualifications and provenance.}\par
\subsubsection*{II.1.1.2.5 [CT2] Search \& Discovery}
\phantomsection
\label{navigation:II.1.1.2.5}
\addcontentsline{toc}{subsubsection}{II.1.1.2.5 [CT2] Search \& Discovery}

\label{element:6a45701634f04b43919d9152a61e5f2e}
\index{[CT2] Search \& Discovery (II.1.1.2.5)}
{\small Capability · Canonical · current}\par

\begin{itemize}


\item \textbf{Architectural Scope}: Index, search, resolve, filter, and discover information assets, entities, directory endpoints, and healthcare services across the enterprise.

\item \textbf{Functional Responsibilities}:

\item Providing multi-criteria indexed search across complex entity graphs (e.g., find active cardiologists at Location X accepting new referrals).

\item Resolving service and practitioner communication endpoints for addressed delivery.

\item Supporting bounded, paginated clinical queries across longitudinal patient records.

\item Applying dynamic security and consent filters to search results prior to egress.

\item \textbf{Anti-Responsibilities}: Does not dictate specific search engine indexes, inverted index implementations, or SQL query syntaxes.

\end{itemize}


Search \& Discovery makes information, entities and service endpoints findable through governed criteria and bounded queries. It addresses the recurring need to locate relevant material or a destination before it can be used. Security-qualified results preserve the distinction between knowing that something exists and being entitled to obtain or disclose it.



\textbf{Boundary}



Discovery does not decide policy or prescribe a search engine, index format or query language.


{\footnotesize \href{../../../../continuum/strategy/capability/6a45701634f04b43919d9152a61e5f2e.md}{Canonical record: 6a45701634f04b43919d9152a61e5f2e}; full qualifications and provenance.}\par
\subsubsection*{II.1.1.2.6 [CT2] Semantic Governance \& Conformance}
\phantomsection
\label{navigation:II.1.1.2.6}
\addcontentsline{toc}{subsubsection}{II.1.1.2.6 [CT2] Semantic Governance \& Conformance}

\label{element:9ee41437bd4c4bd18389f05c12108f97}
\index{[CT2] Semantic Governance \& Conformance (II.1.1.2.6)}
{\small Capability · Canonical · current}\par

\begin{itemize}


\item \textbf{Architectural Scope}: Govern canonical models, schemas, terminologies, value sets, mapping tables, constraint rules, and semantic conformance validation across the enterprise.

\item \textbf{Functional Responsibilities}:

\item Defining and publishing canonical healthcare data models and profile structures.

\item Governing clinical terminology mappings, value sets, and concept lookups (SNOMED CT, LOINC).

\item Validating inbound and outbound message representations against published semantic rules.

\item Governing semantic versioning, deprecation, and backward compatibility across schema revisions.

\item \textbf{Component Relationship}: Informs the Calliope strategic logical component, but Calliope is the responsibility centre, not the capability definition.

\end{itemize}


This capability makes canonical information models, terminology bindings, mappings and conformance rules governable across enterprise use. It addresses both what representations mean and whether they respect the published rules, keeping version and compatibility concerns explicit. A conformant representation can therefore be distinguished from one that merely uses a supported transport or parser.



\textbf{Boundary}



Conformance does not establish clinical truth, information authority or disclosure permission. Calliope is an established responsibility centre, not the definition of the capability.


{\footnotesize \href{../../../../continuum/strategy/capability/9ee41437bd4c4bd18389f05c12108f97.md}{Canonical record: 9ee41437bd4c4bd18389f05c12108f97}; full qualifications and provenance.}\par
\subsection*{II.1.1.3 [CT1] Trust, Accountability \& Assurance}
\phantomsection
\label{navigation:II.1.1.3}
\addcontentsline{toc}{subsection}{II.1.1.3 [CT1] Trust, Accountability \& Assurance}

\label{element:321de197612f49089d8b4a4cc1a2e003}
\index{[CT1] Trust, Accountability \& Assurance (II.1.1.3)}
{\small Capability · Proposed · current}\par

Make governed activity permissible, attributable, operationally dependable and independently assessable.



\textbf{Overview}



Policy, evidence, resilience and independent assurance contribute to confidence through different responsibilities. Keeping them as peers prevents technical success or recorded evidence from becoming a business assurance conclusion.



\textbf{Intent}



Keep permission, attribution, operational integrity and independent confidence distinguishable and usable together.



\textbf{Boundary}



Cross-cutting participation does not distribute semantic ownership, mandate centralisation or authorise the assured subject to determine its own conclusion.


{\footnotesize \href{../../../../continuum/strategy/capability/321de197612f49089d8b4a4cc1a2e003.md}{Canonical record: 321de197612f49089d8b4a4cc1a2e003}; full qualifications and provenance.}\par
\subsubsection*{II.1.1.3.1 [CT2] Information Authority Qualification}
\phantomsection
\label{navigation:II.1.1.3.1}
\addcontentsline{toc}{subsubsection}{II.1.1.3.1 [CT2] Information Authority Qualification}

\label{element:cbba9984dd404deabed41fdcc96088b1}
\index{[CT2] Information Authority Qualification (II.1.1.3.1)}
{\small Capability · Proposed · current}\par

Represent and qualify the source-authority and credibility standing of information independently of transport, persistence and concurrency.



\textbf{Overview}



Identity correlation, clinical-history synthesis, reference-content use and assurance evidence all need to distinguish the standing of an assertion from technical acceptance. This recurring ability is insufficiently explicit in the existing fourteen EC definitions and is proposed as a new peer responsibility.



\textbf{Intent}



Keep source credibility explicit even when the technical handling of competing assertions is equally valid.



\textbf{Boundary}



This is not independent service assurance, a universal conflict arbiter or clinical truth determination. Existing EC-13 governs authority nomenclature; this capability applies and preserves the qualified standing.


{\footnotesize \href{../../../../continuum/strategy/capability/cbba9984dd404deabed41fdcc96088b1.md}{Canonical record: cbba9984dd404deabed41fdcc96088b1}; full qualifications and provenance.}\par
\subsubsection*{II.1.1.3.2 [CT2] Operational Assurance}
\phantomsection
\label{navigation:II.1.1.3.2}
\addcontentsline{toc}{subsubsection}{II.1.1.3.2 [CT2] Operational Assurance}

\label{element:938cd44380ab431fa8dae13c702d83c1}
\index{[CT2] Operational Assurance (II.1.1.3.2)}
{\small Capability · Canonical · current}\par

\begin{itemize}


\item \textbf{Architectural Scope}: Ensure concurrency integrity, system resilience, duplicate detection, exception handling, automated recovery, and operational telemetry across platform operations.

\item \textbf{Functional Responsibilities}:

\item Enforcing sliding-window duplicate message detection and idempotent suppression.

\item Guaranteeing ingress dual-write safety (durable acceptance before emitting positive application ACK).

\item Managing structured failure recovery, exception escalation, and dead-letter isolation.

\item Emitting platform telemetry, operational health metrics, and queue processing metrics.

\item \textbf{Collaborative Nature}: Realized collaboratively across boundary gateways, messaging infrastructure, execution workers, and storage layers.

\end{itemize}


\par\medskip\textbf{EC-12 / EC-14 Boundary}\par


EC-12 retains the established scope and responsibilities above. It may contribute directly to generic operational/processing assurance, limited to that contribution in the \href{\detokenize{../../../../../docs/markdown/02-strategy/capability-maps/health-service-assurance-derivation.md#collaborative-enterprise-capability-contribution-matrix}}{approved assurance matrix}. It does not own the general semantics of assurance disposition, assurance-specific criteria applicability, evidentiary association, assessment against assurance criteria, adjudication, or independent assurance findings and conclusions. Those distinguishing responsibilities are introduced through EC-14 Service Guardian. Operational health and business-level processing assurance remain distinct; EC-12 is not expanded to absorb EC-14.



Operational Assurance preserves processing integrity under concurrency, duplicate delivery and failure. It brings together resilience, durable ingress acceptance, recovery and operational telemetry so that technical degradation does not silently weaken the managed operation's semantics. Its contribution is collaborative across operational boundaries while remaining a bounded architectural responsibility.



\textbf{Boundary}



Operational health is distinct from independent assessment of a governed subject. Assurance-specific criteria applicability, adjudication and conclusions belong to the established EC-14 contribution.


{\footnotesize \href{../../../../continuum/strategy/capability/938cd44380ab431fa8dae13c702d83c1.md}{Canonical record: 938cd44380ab431fa8dae13c702d83c1}; full qualifications and provenance.}\par
\subsubsection*{II.1.1.3.3 [CT2] Policy \& Control}
\phantomsection
\label{navigation:II.1.1.3.3}
\addcontentsline{toc}{subsubsection}{II.1.1.3.3 [CT2] Policy \& Control}

\label{element:0c5d986cc948421a8bebcab3b74d108f}
\index{[CT2] Policy \& Control (II.1.1.3.3)}
{\small Capability · Canonical · current}\par

\begin{itemize}


\item \textbf{Architectural Scope}: Evaluate and enforce security, consent, privacy, and regulatory policies including authority, access control, and information handling rules.

\item \textbf{Functional Responsibilities}:

\item Executing impartial, default-deny policy evaluations for all integration transactions.

\item Enforcing Role-Based and Attribute-Based Access Control (RBAC/ABAC) policies.

\item Evaluating patient consent directives and privacy restrictions at boundary and query checkpoints.

\item Enforcing statutory information disclosure boundaries (e.g., mental health, VIP records).

\item \textbf{Collaborative Nature}: Evaluated across boundary gateways, internal bus dispatch, storage interfaces, and presentation layers.

\end{itemize}


Policy \& Control supplies the reusable evaluation and enforcement of authority, consent, privacy and handling constraints. Its purpose is consistent governed action across ingress, information access, execution and presentation boundaries. Default-deny evaluation preserves the need for established authority rather than treating a caller's intent or a technically valid request as permission.



\textbf{Boundary}



Participation at multiple enforcement points does not distribute policy ownership or turn operational security context into persistent clinical content.


{\footnotesize \href{../../../../continuum/strategy/capability/0c5d986cc948421a8bebcab3b74d108f.md}{Canonical record: 0c5d986cc948421a8bebcab3b74d108f}; full qualifications and provenance.}\par
\paragraph*{II.1.1.3.3.1 [CT3] Operation Policy Resolution}
\phantomsection
\label{navigation:II.1.1.3.3.1}
\addcontentsline{toc}{paragraph}{II.1.1.3.3.1 [CT3] Operation Policy Resolution}

\label{element:7eca9fc4d2244b0facc83376e57ebbc4}
\index{[CT3] Operation Policy Resolution (II.1.1.3.3.1)}
{\small Capability · Proposed · current}\par

Resolve effective, temporally applicable permission criteria for an established identity and governed operation.



\textbf{Overview}



Permission depends on controlled criteria, role assignments, action and information scope; revocation is distinct from natural expiry.



\textbf{Intent}



Determine the effective permitted action under the applicable identity, information and temporal criteria.



\textbf{Boundary}



Does not authenticate identity or grant permission from caller-controlled tags or mere role names.


{\footnotesize \href{../../../../continuum/strategy/capability/7eca9fc4d2244b0facc83376e57ebbc4.md}{Canonical record: 7eca9fc4d2244b0facc83376e57ebbc4}; full qualifications and provenance.}\par
\paragraph*{II.1.1.3.3.2 [CT3] Operation-Bound Policy Enforcement}
\phantomsection
\label{navigation:II.1.1.3.3.2}
\addcontentsline{toc}{paragraph}{II.1.1.3.3.2 [CT3] Operation-Bound Policy Enforcement}

\label{element:cb32b32c657445b3b9b27bebc069ace9}
\index{[CT3] Operation-Bound Policy Enforcement (II.1.1.3.3.2)}
{\small Capability · Proposed · current}\par

Enforce resolved policy constraints at the boundary where a governed action is attempted.



\textbf{Overview}



Distributed participants apply the applicable effective authorisation and fail closed when authority is absent or untrustworthy.



\textbf{Intent}



Ensure an operation cannot proceed merely because its caller omitted or bypassed a voluntary policy check.



\textbf{Boundary}



Enforcement participation does not distribute policy ownership or make security context persistent content.


{\footnotesize \href{../../../../continuum/strategy/capability/cb32b32c657445b3b9b27bebc069ace9.md}{Canonical record: cb32b32c657445b3b9b27bebc069ace9}; full qualifications and provenance.}\par
\subsubsection*{II.1.1.3.4 [CT2] Provenance \& Traceability}
\phantomsection
\label{navigation:II.1.1.3.4}
\addcontentsline{toc}{subsubsection}{II.1.1.3.4 [CT2] Provenance \& Traceability}

\label{element:68a0755aa4e14a509ac5e6c0bbcff3e7}
\index{[CT2] Provenance \& Traceability (II.1.1.3.4)}
{\small Capability · Canonical · current}\par

\begin{itemize}


\item \textbf{Architectural Scope}: Preserve required origin attribution, meaningful transformation history, significant execution checkpoints, tamper-evident audit evidence and end-to-end traceability. Harmonia SHALL preserve provenance, audit and operational evidence where required by governing obligations, architectural significance, accountability, reconstruction, security or assurance needs. Routine technical activity SHALL NOT automatically require durable evidentiary recording.

\item \textbf{Functional Responsibilities}:

\item Recording attributable, verifiable origin metadata for information sources, authors, and timestamps.

\item Tracking information-significant transformation history and required mapping checkpoints as information crosses system seams.

\item Capturing non-PHI operational and security audit trails for regulatory compliance (HIPAA, GDPR).

\item Supporting end-to-end transaction tracing across asynchronous processing pipelines.

\item \textbf{Collaborative Nature}: Participating elements preserve required provenance and audit context according to policy. Telemetry, operational facts, provenance, audit and assurance evidence remain distinct. Being evidence is contextual, not an intrinsic property of every recorded fact; preservation does not establish an assurance conclusion.

\end{itemize}


Provenance \& Traceability preserves the meaningful attribution and history needed to explain information and significant activity. It connects source, transformation and required checkpoints without making every internal mechanism permanent evidence. Its recurring contribution supports accountability, reconstruction and assurance use while retaining the context that makes a recorded fact evidentially relevant.



\textbf{Boundary}



Telemetry, provenance, audit and assurance evidence remain distinct. Recording evidence neither establishes an assurance conclusion nor automatically requires retention of routine technical activity.


{\footnotesize \href{../../../../continuum/strategy/capability/68a0755aa4e14a509ac5e6c0bbcff3e7.md}{Canonical record: 68a0755aa4e14a509ac5e6c0bbcff3e7}; full qualifications and provenance.}\par
\paragraph*{II.1.1.3.4.1 [CT3] Governed Activity Audit Evidence}
\phantomsection
\label{navigation:II.1.1.3.4.1}
\addcontentsline{toc}{paragraph}{II.1.1.3.4.1 [CT3] Governed Activity Audit Evidence}

\label{element:c8d93ddc5aa84960964b78345fff1b29}
\index{[CT3] Governed Activity Audit Evidence (II.1.1.3.4.1)}
{\small Capability · Proposed · current}\par

Capture required significant activity and decision evidence under its applicable governed capture authority.



\textbf{Overview}



Denied actions and successful mutations may both require trustworthy evidence; the audited identity and capture authority remain distinct.



\textbf{Intent}



Preserve required significant action evidence under its proper capture authority.



\textbf{Boundary}



Mandatory capture authority is specific to the established audit responsibility, not a platform-wide permission shortcut.


{\footnotesize \href{../../../../continuum/strategy/capability/c8d93ddc5aa84960964b78345fff1b29.md}{Canonical record: c8d93ddc5aa84960964b78345fff1b29}; full qualifications and provenance.}\par
\paragraph*{II.1.1.3.4.2 [CT3] Information Attribution \& Lineage}
\phantomsection
\label{navigation:II.1.1.3.4.2}
\addcontentsline{toc}{paragraph}{II.1.1.3.4.2 [CT3] Information Attribution \& Lineage}

\label{element:8bc66c9bc1024a68a1bfd010b4dd1b79}
\index{[CT3] Information Attribution \& Lineage (II.1.1.3.4.2)}
{\small Capability · Proposed · current}\par

Preserve source attribution and significant transformation relationships for managed information.



\textbf{Overview}



Lineage explains how an assertion reached its present form without turning every infrastructure event into provenance.



\textbf{Intent}



Explain the information-significant origin and transformation of an assertion.



\textbf{Boundary}



Does not infer clinical credibility from a successful transport or establish independent assurance findings.


{\footnotesize \href{../../../../continuum/strategy/capability/8bc66c9bc1024a68a1bfd010b4dd1b79.md}{Canonical record: 8bc66c9bc1024a68a1bfd010b4dd1b79}; full qualifications and provenance.}\par
\subsubsection*{II.1.1.3.5 [CT2] Service Guardian}
\phantomsection
\label{navigation:II.1.1.3.5}
\addcontentsline{toc}{subsubsection}{II.1.1.3.5 [CT2] Service Guardian}

\label{element:d767f12f9b1c495883216ace464ee1c2}
\index{[CT2] Service Guardian (II.1.1.3.5)}
{\small Capability · Canonical · current}\par

\begin{itemize}


\item \textbf{Architectural Scope}: Service Guardian provides the reusable capability to independently establish whether governed service behaviour, state or outcomes remain within defined assurance expectations by associating applicable criteria and trustworthy evidence with a governed subject, assessing that evidence, adjudicating the assessment, and establishing reportable assurance findings and conclusions.

\item \textbf{Functional Responsibilities}:

\item Supplying assurance-specific semantics for assurance disposition and assurance criteria applicability.

\item Associating a governed subject/context with applicable criteria and sufficient trustworthy evidence for a particular assurance purpose.

\item Assessing evidence against applicable criteria and adjudicating assessed detail to establish reportable findings, exceptions and conclusions.

\item Preserving assurance independence, temporal/version-specific evidence and criteria applicability, and explicit inability to establish a conclusion where evidence is insufficient.

\item \textbf{Collaborative Nature}: Supplies the assurance-specific semantics required to bind existing reusable capabilities into governed assurance. It contributes directly to Assurance Design, Assurance Criteria Management and Governed Assurance; direct contribution is not sole realisation. The \href{\detokenize{../../../../../docs/markdown/02-strategy/capability-maps/health-service-assurance-derivation.md#collaborative-enterprise-capability-contribution-matrix}}{approved matrix} preserves the contributions of EC-01 through EC-13.

\item \textbf{Anti-Responsibilities}: Does not own the underlying governed subject, authoritative governance requirements, or source information used as evidence. Does not perform, manage, remediate, assign, delegate or operationally escalate the subject activity. Does not replace System Operations or EC-12 Operational Assurance, and does not provide Clinical Services Delivery Assurance.

\item \textbf{Allocation Boundary}: No strategic logical component, application component, runtime service or implementation construct is allocated EC-14. The catalogue defines a capability, not a software component or deployment boundary.

\end{itemize}


\par\medskip\textbf{Conceptual Behavioural Sequence}\par


\begin{quote}\small\ttfamily
Define Good Behaviour\\
    ↓\\
Gather Evidence\\
    ↓\\
Assess Evidence against Good Behaviour\\
    ↓\\
Adjudicate Assessed Detail\\
    ↓\\
Report\\

\end{quote}


This sequence explains the capability; it is not an implementation workflow specification or a derived Business Process. “Define Good Behaviour” identifies assurance expectations and criteria; it does not transfer ownership of governing requirements to Service Guardian. Source information ownership remains with its responsible capability; assurance owns its contextual evidentiary association.



The \href{\detokenize{../../../../../docs/markdown/02-strategy/capabilities/business-enabling-capabilities.md#health-service-assurance-approved-business-enabling-capabilities}}{Business Enabling assurance principles} govern disposition, generic processing assurance, evidence/criteria temporality, assessment/adjudication, independence, failure/recovery and explicit evidence insufficiency. These semantics are not universally reducible to PASS / FAIL / UNKNOWN and do not imply an assurance information model or persistence mechanism.



Service Guardian supplies the reusable assurance semantics that bind a governed subject to applicable criteria and trustworthy evidence, support assessment and adjudication, and establish reportable findings. Its contribution distinguishes explicit evidence insufficiency from a conclusion about the subject. Existing Enterprise Capabilities contribute collaboratively; EC-14 does not duplicate all their information, policy or execution responsibilities.



Independence concerns control of evaluation and conclusion. A subject may supply evidence but cannot solely determine, suppress or retrospectively change its own assurance outcome.



\textbf{Boundary}



The capability does not manage or remediate the subject, replace operational assurance, or assure clinical adequacy. No application component, runtime service or deployment boundary is allocated, and the explanatory sequence is not a derived Business Process.


{\footnotesize \href{../../../../continuum/strategy/capability/d767f12f9b1c495883216ace464ee1c2.md}{Canonical record: d767f12f9b1c495883216ace464ee1c2}; full qualifications and provenance.}\par
\section*{II.1.2 Business Capabilities}
\phantomsection
\label{navigation:II.1.2}
\addcontentsline{toc}{section}{II.1.2 Business Capabilities}

\subsection*{II.1.2.1 [CT1] Care Access \& Coordination}
\phantomsection
\label{navigation:II.1.2.1}
\addcontentsline{toc}{subsection}{II.1.2.1 [CT1] Care Access \& Coordination}

{\small Historical predecessor: R.2.1 \hyperref[element:0dec07b6ddbe48bda93a2295dd3e2d7d]{[CT1] Care Access \& Coordination}; Canonical ID: \texttt{\seqsplit{0dec07b6ddbe48bda93a2295dd3e2d7d}}. Superseded historical identity — direct correspondence · superseded.}\par
\label{element:ed342cb8656b4f80b55ad519ecbd2d3a}
\index{[CT1] Care Access \& Coordination (II.1.2.1)}
{\small Capability · Canonical · current}\par

Managing patient intake, referral triage, appointment scheduling, multidisciplinary care team coordination, cross-facility patient transfers, and care continuity across organisational boundaries.



This capability makes care reachable and operationally coherent across participating services. It encompasses entering care, obtaining scheduled opportunities and usable capacity, coordinating participation and responsibility transitions, and arranging bounded support logistics. Its decomposition separates access decisions, resource use, continuity and supporting work so that movement or message delivery cannot stand in for acceptance of care responsibility.



\textbf{Intent}



Help a person's care requirements progress across organisational and operational boundaries with the relevant access and continuity context intact.



\textbf{Boundary}



Coordination does not determine clinical work or confer clinical authority. Corporate resource governance and general hospital logistics are outside the bounded care-support responsibility.



\textbf{Harmonia Architectural Relevance}



Referrals, encounters, appointments, operational state and communication provide material information for continuity. Harmonia's integration and coordination support does not make it the owner of clinical handover or the originating workforce and resource facts.


{\footnotesize \href{../../../../continuum/strategy/capability/ed342cb8656b4f80b55ad519ecbd2d3a.md}{Canonical record: ed342cb8656b4f80b55ad519ecbd2d3a}; full qualifications and provenance.}\par
\subsubsection*{II.1.2.1.1 [CT2] Care Access Management}
\phantomsection
\label{navigation:II.1.2.1.1}
\addcontentsline{toc}{subsubsection}{II.1.2.1.1 [CT2] Care Access Management}

\label{element:acbcad8330e64b59be5e3ec4a100a544}
\index{[CT2] Care Access Management (II.1.2.1.1)}
{\small Capability · Canonical · current}\par

Establish and manage access to applicable care.



Care Access Management establishes the circumstances in which an individual seeks or enters care. Intake creates the access context, referral requests another participant's involvement, and eligibility determines whether applicable access conditions are met. These are separate contributions to access: a person can be eligible without having a booking, and a received referral can remain undecided.



\textbf{Intent}



Make entry requirements and referral disposition intelligible without implying that care has already been accepted or delivered.



\textbf{Boundary}



Eligibility does not allocate care, and referral receipt does not establish an assumption of clinical responsibility.


{\small Recorded unresolved matters: source-qualified matters; exact scope remains in the canonical record.}\par
{\footnotesize \href{../../../../continuum/strategy/capability/acbcad8330e64b59be5e3ec4a100a544.md}{Canonical record: acbcad8330e64b59be5e3ec4a100a544}; full qualifications and provenance.}\par
\paragraph*{II.1.2.1.1.1 [CT3] Care Eligibility Determination}
\phantomsection
\label{navigation:II.1.2.1.1.1}
\addcontentsline{toc}{paragraph}{II.1.2.1.1.1 [CT3] Care Eligibility Determination}

\label{element:0922a1141f82416ca9de12054c28a69a}
\index{[CT3] Care Eligibility Determination (II.1.2.1.1.1)}
{\small Capability · Canonical · current}\par

Determine whether applicable eligibility conditions for care access are satisfied.



Care Eligibility Determination establishes whether the conditions governing access to applicable care are met. It contributes a qualified access decision alongside intake and referral, so that an identifiable request can be distinguished from one that satisfies the relevant criteria. Its purpose is eligibility in the applicable context rather than clinical diagnosis or service allocation.



\textbf{Boundary}



Eligibility alone does not reserve capacity, allocate a care place or deliver care.


{\small Recorded unresolved matters: source-qualified matters; exact scope remains in the canonical record.}\par
{\footnotesize \href{../../../../continuum/strategy/capability/0922a1141f82416ca9de12054c28a69a.md}{Canonical record: 0922a1141f82416ca9de12054c28a69a}; full qualifications and provenance.}\par
\paragraph*{II.1.2.1.1.2 [CT3] Care Intake Management}
\phantomsection
\label{navigation:II.1.2.1.1.2}
\addcontentsline{toc}{paragraph}{II.1.2.1.1.2 [CT3] Care Intake Management}

\label{element:1582467fff3b4dafa348aa100e9ca2de}
\index{[CT3] Care Intake Management (II.1.2.1.1.2)}
{\small Capability · Canonical · current}\par

Establish and manage the intake context through which an individual seeks or enters care.



Care Intake Management establishes the access context through which a person seeks or enters care. It makes the entry circumstance manageable without assuming that every intake follows a referral or already satisfies eligibility conditions. Within Care Access Management, intake is the responsibility for entry context rather than a decision that all subsequent care arrangements are complete.



\textbf{Boundary}



An intake record does not establish acceptance of a referral, care responsibility or delivery of care.


{\small Recorded unresolved matters: source-qualified matters; exact scope remains in the canonical record.}\par
{\footnotesize \href{../../../../continuum/strategy/capability/1582467fff3b4dafa348aa100e9ca2de.md}{Canonical record: 1582467fff3b4dafa348aa100e9ca2de}; full qualifications and provenance.}\par
\paragraph*{II.1.2.1.1.3 [CT3] Care Referral Management}
\phantomsection
\label{navigation:II.1.2.1.1.3}
\addcontentsline{toc}{paragraph}{II.1.2.1.1.3 [CT3] Care Referral Management}

\label{element:ebdaa781bba542f8b744ef34a836a269}
\index{[CT3] Care Referral Management (II.1.2.1.1.3)}
{\small Capability · Canonical · current}\par

Manage a request for another participant/service to assess, manage or assume applicable care responsibility.



Care Referral Management concerns a request for another participant or service to assess, manage or assume applicable care responsibility. It preserves the distinction between the requester's intent and the recipient's disposition, allowing care coordination to recognise an unresolved referral rather than assume a successful handoff. Referral is therefore broader than sending a document.



\textbf{Boundary}



Submission or receipt does not establish acceptance, an appointment, delivered care or responsibility actually assumed.


{\small Recorded unresolved matters: source-qualified matters; exact scope remains in the canonical record.}\par
{\footnotesize \href{../../../../continuum/strategy/capability/ebdaa781bba542f8b744ef34a836a269.md}{Canonical record: ebdaa781bba542f8b744ef34a836a269}; full qualifications and provenance.}\par
\subsubsection*{II.1.2.1.2 [CT2] Care Coordination \& Transition}
\phantomsection
\label{navigation:II.1.2.1.2}
\addcontentsline{toc}{subsubsection}{II.1.2.1.2 [CT2] Care Coordination \& Transition}

\label{element:174a9b93d19c44828f20360e9d543925}
\index{[CT2] Care Coordination \& Transition (II.1.2.1.2)}
{\small Capability · Canonical · current}\par

Coordinate care participation, actual responsibility transition and discharge.



Care Coordination \& Transition supports continuity by keeping participation, actual responsibility assumption and discharge progression distinct. Care teams may collaborate while responsibility is shared or changes, but physical transit and information exchange do not themselves establish who has accepted care. The capability therefore addresses continuity of accountable participation as well as the practical circumstances of departure.



\textbf{Intent}



Avoid gaps in care continuity caused by treating a movement or communication as an accepted responsibility transition.



\textbf{Boundary}



Harmonia may assemble handover information without managing clinical handover. Arrival, discharge readiness and physical departure are not equivalent to responsibility acceptance.


{\small Recorded unresolved matters: source-qualified matters; exact scope remains in the canonical record.}\par
{\footnotesize \href{../../../../continuum/strategy/capability/174a9b93d19c44828f20360e9d543925.md}{Canonical record: 174a9b93d19c44828f20360e9d543925}; full qualifications and provenance.}\par
\paragraph*{II.1.2.1.2.1 [CT3] Care Discharge Coordination}
\phantomsection
\label{navigation:II.1.2.1.2.1}
\addcontentsline{toc}{paragraph}{II.1.2.1.2.1 [CT3] Care Discharge Coordination}

\label{element:62e436db8ce44a98ba855de68ac6b876}
\index{[CT3] Care Discharge Coordination (II.1.2.1.2.1)}
{\small Capability · Canonical · current}\par

Coordinate applicable discharge requirements and progression.



Care Discharge Coordination brings applicable discharge requirements and progression into a coherent operational context. It supports the transition out of a care circumstance while keeping readiness, physical exit and continuing responsibility distinguishable. The capability concerns coordinating those requirements rather than deciding that every clinical condition for discharge has been met.



\textbf{Boundary}



Discharge readiness and departure are not themselves evidence that another participant has accepted care responsibility.


{\small Recorded unresolved matters: source-qualified matters; exact scope remains in the canonical record.}\par
{\footnotesize \href{../../../../continuum/strategy/capability/62e436db8ce44a98ba855de68ac6b876.md}{Canonical record: 62e436db8ce44a98ba855de68ac6b876}; full qualifications and provenance.}\par
\paragraph*{II.1.2.1.2.2 [CT3] Care Responsibility Transition}
\phantomsection
\label{navigation:II.1.2.1.2.2}
\addcontentsline{toc}{paragraph}{II.1.2.1.2.2 [CT3] Care Responsibility Transition}

\label{element:f8c9f7e2dcf14ceb9f228daaf9fed45e}
\index{[CT3] Care Responsibility Transition (II.1.2.1.2.2)}
{\small Capability · Canonical · current}\par

Establish and coordinate who has actually assumed applicable care responsibility.



Care Responsibility Transition establishes who has actually assumed applicable care responsibility and coordinates that transition. Its purpose is to prevent physical movement or exchanged information from being mistaken for accountable acceptance. It gives care continuity a responsibility meaning that a transport checkpoint or delivered referral cannot establish alone.



\textbf{Boundary}



Patient arrival, message receipt and departure do not by themselves demonstrate accepted responsibility.


{\small Recorded unresolved matters: source-qualified matters; exact scope remains in the canonical record.}\par
{\footnotesize \href{../../../../continuum/strategy/capability/f8c9f7e2dcf14ceb9f228daaf9fed45e.md}{Canonical record: f8c9f7e2dcf14ceb9f228daaf9fed45e}; full qualifications and provenance.}\par
\paragraph*{II.1.2.1.2.3 [CT3] Care-Team Coordination}
\phantomsection
\label{navigation:II.1.2.1.2.3}
\addcontentsline{toc}{paragraph}{II.1.2.1.2.3 [CT3] Care-Team Coordination}

\label{element:bbd138aab70d4b06b43662fb3f5373db}
\index{[CT3] Care-Team Coordination (II.1.2.1.2.3)}
{\small Capability · Canonical · current}\par

Coordinate participating care-team context and collaboration around care.



Care-Team Coordination maintains the participating team context and collaboration around care. It supports continuity by making participation intelligible across people and organisational boundaries, without assuming that membership itself assigns every clinical responsibility. Within Care Coordination \& Transition it concerns collaboration around care rather than proof of a responsibility handoff.



\textbf{Boundary}



Harmonia's support for team information and collaboration does not confer management of clinical handover or clinical worklists.


{\small Recorded unresolved matters: source-qualified matters; exact scope remains in the canonical record.}\par
{\footnotesize \href{../../../../continuum/strategy/capability/bbd138aab70d4b06b43662fb3f5373db.md}{Canonical record: bbd138aab70d4b06b43662fb3f5373db}; full qualifications and provenance.}\par
\subsubsection*{II.1.2.1.3 [CT2] Care Scheduling \& Capacity Management}
\phantomsection
\label{navigation:II.1.2.1.3}
\addcontentsline{toc}{subsubsection}{II.1.2.1.3 [CT2] Care Scheduling \& Capacity Management}

\label{element:61dd4b247df14546bf012587164bfe91}
\index{[CT2] Care Scheduling \& Capacity Management (II.1.2.1.3)}
{\small Capability · Canonical · current}\par

Manage scheduled opportunities, care places, usable operational capacity and care-demand resourcing.



This capability connects care demand with scheduled opportunities, suitable care places and legitimately usable operational resources. Its children distinguish a booking, a place allocation, the readiness of that place, resource deployment, flow coordination and service-level usable capacity. An empty bed or an available practitioner is relevant evidence, but neither alone establishes that the service can accommodate care.



\textbf{Intent}



Support care requirements with an operationally credible combination of time, place, resources and capacity.



\textbf{Boundary}



Deployment consumes roster, qualification and privilege context without changing their authoritative records. Operational capacity differs from strategic capacity planning and from the unresolved collective service-readiness concern.


{\small Recorded unresolved matters: source-qualified matters; exact scope remains in the canonical record.}\par
{\footnotesize \href{../../../../continuum/strategy/capability/61dd4b247df14546bf012587164bfe91.md}{Canonical record: 61dd4b247df14546bf012587164bfe91}; full qualifications and provenance.}\par
\paragraph*{II.1.2.1.3.1 [CT3] Care Appointment Management}
\phantomsection
\label{navigation:II.1.2.1.3.1}
\addcontentsline{toc}{paragraph}{II.1.2.1.3.1 [CT3] Care Appointment Management}

\label{element:14dc21b28938425498ff7b04e49efd30}
\index{[CT3] Care Appointment Management (II.1.2.1.3.1)}
{\small Capability · Canonical · current}\par

Establish and manage a scheduled opportunity for care.



Care Appointment Management establishes a scheduled opportunity for care and maintains that opportunity as arrangements change. Within scheduling and capacity it supplies the temporal commitment around which participants can organise, while leaving the actual availability of places and resources distinguishable. A scheduled opportunity is useful coordination information even before attendance or care occurs.



\textbf{Boundary}



A booking does not establish attendance, intervention performance or satisfaction of diagnostic-service-specific scheduling constraints.


{\small Recorded unresolved matters: source-qualified matters; exact scope remains in the canonical record.}\par
{\footnotesize \href{../../../../continuum/strategy/capability/14dc21b28938425498ff7b04e49efd30.md}{Canonical record: 14dc21b28938425498ff7b04e49efd30}; full qualifications and provenance.}\par
\paragraph*{II.1.2.1.3.2 [CT3] Care Resource Deployment}
\phantomsection
\label{navigation:II.1.2.1.3.2}
\addcontentsline{toc}{paragraph}{II.1.2.1.3.2 [CT3] Care Resource Deployment}

\label{element:7854224c3b064116b0d48f01ce5537d8}
\index{[CT3] Care Resource Deployment (II.1.2.1.3.2)}
{\small Capability · Canonical · current}\par

Apply legitimately available and appropriately authorised resources to a current or planned operational care requirement, particularly using roster/on-call and service-coverage context.



Care Resource Deployment applies legitimately available resources to a current or planned care requirement. It considers coverage, role, qualification, engagement and privilege context to distinguish availability in principle from deployment to a particular need. Its contribution is the operational application of resources, without changing the authoritative facts that make their use legitimate.



\textbf{Boundary}



Deployment is neither roster authoring nor clinical work assignment and does not confer clinical authority or ownership of enterprise assets.


{\small Recorded unresolved matters: source-qualified matters; exact scope remains in the canonical record.}\par
{\footnotesize \href{../../../../continuum/strategy/capability/7854224c3b064116b0d48f01ce5537d8.md}{Canonical record: 7854224c3b064116b0d48f01ce5537d8}; full qualifications and provenance.}\par
\paragraph*{II.1.2.1.3.3 [CT3] Care-Place Allocation}
\phantomsection
\label{navigation:II.1.2.1.3.3}
\addcontentsline{toc}{paragraph}{II.1.2.1.3.3 [CT3] Care-Place Allocation}

\label{element:f8a4df47bea742a0869802472136d55f}
\index{[CT3] Care-Place Allocation (II.1.2.1.3.3)}
{\small Capability · Canonical · current}\par

Allocate an applicable care place to a care requirement.



Care-Place Allocation associates a care requirement with an applicable care place. Its outcome is the allocation itself, distinct from defining the place or establishing its readiness. This separation supports reasoning about reserved, allocated and usable places without assuming that an apparently empty location is fit for the requirement.



\textbf{Boundary}



Allocation does not establish clinical suitability beyond the applicable context or make an unready place ready.


{\small Recorded unresolved matters: source-qualified matters; exact scope remains in the canonical record.}\par
{\footnotesize \href{../../../../continuum/strategy/capability/f8a4df47bea742a0869802472136d55f.md}{Canonical record: f8a4df47bea742a0869802472136d55f}; full qualifications and provenance.}\par
\paragraph*{II.1.2.1.3.4 [CT3] Care-Place Readiness Management}
\phantomsection
\label{navigation:II.1.2.1.3.4}
\addcontentsline{toc}{paragraph}{II.1.2.1.3.4 [CT3] Care-Place Readiness Management}

\label{element:4dd0f1b9f9334c5c80066a8974fdfb7a}
\index{[CT3] Care-Place Readiness Management (II.1.2.1.3.4)}
{\small Capability · Canonical · current}\par

Establish and manage whether a particular care place is ready for allocation or use.



Care-Place Readiness Management establishes whether a particular place can be allocated or used in its operational context. It connects the place's condition with its intended use so that vacancy can be distinguished from readiness. This is a place-specific responsibility within scheduling and capacity, supporting credible allocation rather than merely counting physical spaces.



\textbf{Boundary}



Readiness of one place does not establish collective readiness or usable capacity of the whole service.


{\small Recorded unresolved matters: source-qualified matters; exact scope remains in the canonical record.}\par
{\footnotesize \href{../../../../continuum/strategy/capability/4dd0f1b9f9334c5c80066a8974fdfb7a.md}{Canonical record: 4dd0f1b9f9334c5c80066a8974fdfb7a}; full qualifications and provenance.}\par
\paragraph*{II.1.2.1.3.5 [CT3] Operational Care Capacity Management}
\phantomsection
\label{navigation:II.1.2.1.3.5}
\addcontentsline{toc}{paragraph}{II.1.2.1.3.5 [CT3] Operational Care Capacity Management}

\label{element:3040519e1fec43699c14a3cafb2d945d}
\index{[CT3] Operational Care Capacity Management (II.1.2.1.3.5)}
{\small Capability · Canonical · current}\par

Understand and manage the service’s usable capacity to accommodate care in its current operational context.



Operational Care Capacity Management establishes and manages the service's usable ability to accommodate care under current conditions. It considers capacity as a service-level concern rather than a simple count of staff, devices or vacant beds. This supports decisions grounded in what can actually be accommodated, while keeping current capacity distinct from plans for future provision.



\textbf{Boundary}



Resource availability, strategic capacity requirements and collective service readiness are distinct; the latter's placement remains unresolved.


{\small Recorded unresolved matters: source-qualified matters; exact scope remains in the canonical record.}\par
{\footnotesize \href{../../../../continuum/strategy/capability/3040519e1fec43699c14a3cafb2d945d.md}{Canonical record: 3040519e1fec43699c14a3cafb2d945d}; full qualifications and provenance.}\par
\paragraph*{II.1.2.1.3.6 [CT3] Operational Care Flow Coordination}
\phantomsection
\label{navigation:II.1.2.1.3.6}
\addcontentsline{toc}{paragraph}{II.1.2.1.3.6 [CT3] Operational Care Flow Coordination}

\label{element:d1dedb10ea694af5b88dbe1b6192f3ac}
\index{[CT3] Operational Care Flow Coordination (II.1.2.1.3.6)}
{\small Capability · Canonical · current}\par

Coordinate operational care flow and its progression within the care-delivery environment.



Operational Care Flow Coordination addresses progression of care flow within the delivery environment. It provides the coordinating responsibility around operational movement and impediments, rather than treating each scheduled opportunity or place allocation as an isolated fact. Its outcome concerns coherent flow under current conditions, not the clinical determination of what care should occur.



\textbf{Boundary}



Coordinating flow does not transfer clinical work, workforce governance or general hospital-logistics authority.


{\small Recorded unresolved matters: source-qualified matters; exact scope remains in the canonical record.}\par
{\footnotesize \href{../../../../continuum/strategy/capability/d1dedb10ea694af5b88dbe1b6192f3ac.md}{Canonical record: d1dedb10ea694af5b88dbe1b6192f3ac}; full qualifications and provenance.}\par
\subsubsection*{II.1.2.1.4 [CT2] Care Support Logistics Management}
\phantomsection
\label{navigation:II.1.2.1.4}
\addcontentsline{toc}{subsubsection}{II.1.2.1.4 [CT2] Care Support Logistics Management}

\label{element:92be036674254f719cf4bb613120128a}
\index{[CT2] Care Support Logistics Management (II.1.2.1.4)}
{\small Capability · Canonical · current}\par

Coordinate bounded operational support needed for care access and fulfilment.



Care Support Logistics Management coordinates the bounded physical and operational support needed to fulfil care. Patient transit, clinical-material movement and authorised support work have different subjects and completion evidence, even when the same operational environment supports them. This capability provides their common care-support frame without absorbing every logistical activity of a healthcare enterprise.



\textbf{Intent}



Ensure necessary supporting activity can progress with the relevant patient, material and operational constraints visible.



\textbf{Boundary}



Assignment does not establish a worker's qualification or clinical authority. Enterprise supply chains, facilities and workforce governance retain their own responsibilities.


{\small Recorded unresolved matters: source-qualified matters; exact scope remains in the canonical record.}\par
{\footnotesize \href{../../../../continuum/strategy/capability/92be036674254f719cf4bb613120128a.md}{Canonical record: 92be036674254f719cf4bb613120128a}; full qualifications and provenance.}\par
\paragraph*{II.1.2.1.4.1 [CT3] Care Support Work Coordination}
\phantomsection
\label{navigation:II.1.2.1.4.1}
\addcontentsline{toc}{paragraph}{II.1.2.1.4.1 [CT3] Care Support Work Coordination}

\label{element:f9f5c69b0f0b472287cd4ac56ccff6f3}
\index{[CT3] Care Support Work Coordination (II.1.2.1.4.1)}
{\small Capability · Canonical · current}\par

Coordinate authorised operational support work: what needs doing, who performs it and how progression, completion or escalation is coordinated.



Care Support Work Coordination organises authorised operational support work around care: what needs doing, who will perform it and how progress, completion or escalation is coordinated. It gives supporting activity an explicit progression responsibility, rather than assuming that issuing a request means the work occurred. The work remains bounded by applicable authority and resource constraints.



\textbf{Boundary}



Dispatch does not establish worker availability, qualification or privilege, and support coordination does not determine clinical work.


{\small Recorded unresolved matters: source-qualified matters; exact scope remains in the canonical record.}\par
{\footnotesize \href{../../../../continuum/strategy/capability/f9f5c69b0f0b472287cd4ac56ccff6f3.md}{Canonical record: f9f5c69b0f0b472287cd4ac56ccff6f3}; full qualifications and provenance.}\par
\paragraph*{II.1.2.1.4.2 [CT3] Clinical Material Logistics Coordination}
\phantomsection
\label{navigation:II.1.2.1.4.2}
\addcontentsline{toc}{paragraph}{II.1.2.1.4.2 [CT3] Clinical Material Logistics Coordination}

\label{element:4b66444de58a48a9a9433568bf3a60d3}
\index{[CT3] Clinical Material Logistics Coordination (II.1.2.1.4.2)}
{\small Capability · Canonical · current}\par

Coordinate clinically relevant material movement, custody, availability or disposition needed for care or diagnostic fulfilment.



This capability coordinates clinically relevant material movement and the custody, availability or disposition needed for care or diagnostic fulfilment. It makes the material's operational journey intelligible while retaining why the material matters to the clinical activity. The responsibility differs from moving a patient and from managing the diagnostic meaning of information derived from a specimen.



\textbf{Boundary}



Bounded clinical-material logistics does not establish a general procurement, supply-chain or enterprise inventory architecture.


{\small Recorded unresolved matters: source-qualified matters; exact scope remains in the canonical record.}\par
{\footnotesize \href{../../../../continuum/strategy/capability/4b66444de58a48a9a9433568bf3a60d3.md}{Canonical record: 4b66444de58a48a9a9433568bf3a60d3}; full qualifications and provenance.}\par
\paragraph*{II.1.2.1.4.3 [CT3] Patient Movement Coordination}
\phantomsection
\label{navigation:II.1.2.1.4.3}
\addcontentsline{toc}{paragraph}{II.1.2.1.4.3 [CT3] Patient Movement Coordination}

\label{element:f404b6647fe04a3390e57022bdef0843}
\index{[CT3] Patient Movement Coordination (II.1.2.1.4.3)}
{\small Capability · Canonical · current}\par

Coordinate physical patient transit and applicable mobility constraints.



Patient Movement Coordination concerns the person's physical transit and the mobility constraints relevant to it. Within care-support logistics it supports movement appropriate to the patient's circumstances without treating transport completion as a clinical transition. Its responsibility is useful wherever care requires a change of physical location.



\textbf{Boundary}



Transport progression does not establish clinical handover, destination acceptance of care or authority over the clinical plan.


{\small Recorded unresolved matters: source-qualified matters; exact scope remains in the canonical record.}\par
{\footnotesize \href{../../../../continuum/strategy/capability/f404b6647fe04a3390e57022bdef0843.md}{Canonical record: f404b6647fe04a3390e57022bdef0843}; full qualifications and provenance.}\par
\subsubsection*{II.1.2.1.5 [CT2] Service Delivery Preparedness Management}
\phantomsection
\label{navigation:II.1.2.1.5}
\addcontentsline{toc}{subsubsection}{II.1.2.1.5 [CT2] Service Delivery Preparedness Management}

\label{element:9665d894c4c9438999d8cc10afae695d}
\index{[CT2] Service Delivery Preparedness Management (II.1.2.1.5)}
{\small Capability · Proposed · current}\par

Establish whether service offerings can be delivered in their intended operating contexts through the combined standing of their dependencies.



\textbf{Overview}



Preparedness asks whether the service as a whole can operate, distinct from scheduling capacity or readiness of a single place. This is a proposed missing business frame for the already identified readiness concern.



\textbf{Intent}



Recognise whether the whole service can actually operate in its intended context.



\textbf{Boundary}



Does not own the dependencies, their management or clinical execution.


{\footnotesize \href{../../../../continuum/strategy/capability/9665d894c4c9438999d8cc10afae695d.md}{Canonical record: 9665d894c4c9438999d8cc10afae695d}; full qualifications and provenance.}\par
\paragraph*{II.1.2.1.5.1 [CT3] Health Service Operational Readiness}
\phantomsection
\label{navigation:II.1.2.1.5.1}
\addcontentsline{toc}{paragraph}{II.1.2.1.5.1 [CT3] Health Service Operational Readiness}

\label{element:cc91eb4017034939ab82b152bddf7b1c}
\index{[CT3] Health Service Operational Readiness (II.1.2.1.5.1)}
{\small Capability · Candidate · current}\par

Definition unresolved in the baseline. No definition has been synthesised from the name.



\textbf{Proposed overview}



Collective readiness is an access prerequisite and has a coherent proposed preparedness frame; it is not reduced to capacity, resource state or service alignment.



\textbf{Proposed definition}



Establish collective operational readiness of a service offering in a site or session context from its applicable dependencies.



\textbf{Proposed boundary}



The element remains Candidate. The definition/placement elaboration is Proposed and does not adjudicate the unresolved RAD question or transfer source-domain authority.



\medskip\hrule\medskip

{\small Recorded unresolved matters: source-qualified matters; exact scope remains in the canonical record.}\par
{\footnotesize \href{../../../../continuum/strategy/capability/cc91eb4017034939ab82b152bddf7b1c.md}{Canonical record: cc91eb4017034939ab82b152bddf7b1c}; full qualifications and provenance.}\par
\subsection*{II.1.2.2 [CT1] Clinical Quality, Safety \& Improvement}
\phantomsection
\label{navigation:II.1.2.2}
\addcontentsline{toc}{subsection}{II.1.2.2 [CT1] Clinical Quality, Safety \& Improvement}

{\small Historical predecessor: R.2.11 \hyperref[element:b967c383dbcd4c24abe6a83ea5e79e00]{[CT1] Clinical Quality, Safety \& Improvement}; Canonical ID: \texttt{\seqsplit{b967c383dbcd4c24abe6a83ea5e79e00}}. Superseded historical identity — direct correspondence · superseded.}\par
\label{element:63837bf848774572a06b935af95581a2}
\index{[CT1] Clinical Quality, Safety \& Improvement (II.1.2.2)}
{\small Capability · Canonical · current}\par

Operational clinical quality, safety and improvement: monitoring clinical outcomes; identifying quality and safety concerns; adverse event and incident management, reporting and response; infection control; operational quality and safety monitoring; continuous quality improvement; and coordination of operational responses to identified concerns.



This capability addresses operational clinical quality and safety, including concerns, incidents, response and improvement. Its established safety branch distinguishes risk identification, prospective mitigation and active response. Its quality branch distinguishes criteria, evidence, assessment and reporting. Together these make it possible to reason about clinical quality without equating an observation, a judgement and the action taken in response.



\textbf{Intent}



Support recognition and management of clinical safety concerns and informed improvement of clinical quality.



\textbf{Boundary}



This is neither a generic enterprise incident container nor a replacement for macro-level Health Service Assurance. Its bounded decomposition does not exhaust clinical quality and safety practice.



\textbf{Harmonia Architectural Relevance}



Information about outcomes, risks, incidents and operational responses materially supports these responsibilities. Integrating that evidence does not assign clinical quality judgement or originating clinical authority to the platform.


{\footnotesize \href{../../../../continuum/strategy/capability/63837bf848774572a06b935af95581a2.md}{Canonical record: 63837bf848774572a06b935af95581a2}; full qualifications and provenance.}\par
\subsubsection*{II.1.2.2.1 [CT2] Clinical Quality Management}
\phantomsection
\label{navigation:II.1.2.2.1}
\addcontentsline{toc}{subsubsection}{II.1.2.2.1 [CT2] Clinical Quality Management}

{\small Historical predecessor: R.2.6 \hyperref[element:dc1bece79eb7452e951d03a2ee6dfaa6]{[CT2] Clinical Quality Management}; Canonical ID: \texttt{\seqsplit{dc1bece79eb7452e951d03a2ee6dfaa6}}. Superseded historical identity — direct correspondence · superseded.}\par
\label{element:89761f2d38f846c39df1df0975964288}
\index{[CT2] Clinical Quality Management (II.1.2.2.1)}
{\small Capability · Canonical · current}\par

Manage the clinical-quality distinctions materially required for R1/R2.



Clinical Quality Management supplies the criteria and evidence basis for evaluating clinical quality and communicating established findings. Its four children preserve the difference between the expected standard, the material considered, the assessment reached and the report of that assessment. This supports scrutiny of why a quality conclusion is justified rather than treating the presence of measurements as a judgement.



\textbf{Intent}



Make clinical-quality conclusions understandable in relation to their criteria, evidence and assessment authority.



\textbf{Boundary}



Reporting does not originate the conclusion or perform subsequent improvement. Neither a separate monitoring provider nor generic Continuous Improvement placement is established by this decomposition.


{\small Recorded unresolved matters: source-qualified matters; exact scope remains in the canonical record.}\par
{\footnotesize \href{../../../../continuum/strategy/capability/89761f2d38f846c39df1df0975964288.md}{Canonical record: 89761f2d38f846c39df1df0975964288}; full qualifications and provenance.}\par
\paragraph*{II.1.2.2.1.1 [CT3] Clinical Quality Criteria Management}
\phantomsection
\label{navigation:II.1.2.2.1.1}
\addcontentsline{toc}{paragraph}{II.1.2.2.1.1 [CT3] Clinical Quality Criteria Management}

{\small Historical predecessor: R.2.2 \hyperref[element:7a51321b00464016920976b269981edd]{[CT3] Clinical Quality Criteria Management}; Canonical ID: \texttt{\seqsplit{7a51321b00464016920976b269981edd}}. Superseded historical identity — direct correspondence · superseded.}\par
\label{element:1f48527a9a7c4382a7f8a660ccb9f29f}
\index{[CT3] Clinical Quality Criteria Management (II.1.2.2.1.1)}
{\small Capability · Canonical · current}\par

Establish and manage measures, indicators, expectations and criteria against which clinical quality is considered.



Clinical Quality Criteria Management establishes the measures, indicators and expectations used to consider clinical quality. It makes the basis of a later judgement explicit so that evidence can be assessed against an identifiable expectation rather than an unstated standard. Criteria are distinct from the measurements or findings produced when they are applied.



\textbf{Boundary}



Managing criteria does not originate the underlying evidence or establish a quality conclusion.


{\small Recorded unresolved matters: source-qualified matters; exact scope remains in the canonical record.}\par
{\footnotesize \href{../../../../continuum/strategy/capability/1f48527a9a7c4382a7f8a660ccb9f29f.md}{Canonical record: 1f48527a9a7c4382a7f8a660ccb9f29f}; full qualifications and provenance.}\par
\paragraph*{II.1.2.2.1.2 [CT3] Clinical Quality Evidence Management}
\phantomsection
\label{navigation:II.1.2.2.1.2}
\addcontentsline{toc}{paragraph}{II.1.2.2.1.2 [CT3] Clinical Quality Evidence Management}

{\small Historical predecessor: R.2.3 \hyperref[element:f31902d474894670aeef3e0abbedccf0]{[CT3] Clinical Quality Evidence Management}; Canonical ID: \texttt{\seqsplit{f31902d474894670aeef3e0abbedccf0}}. Superseded historical identity — direct correspondence · superseded.}\par
\label{element:1b0a1c6945ed4fadba045ff4753c104c}
\index{[CT3] Clinical Quality Evidence Management (II.1.2.2.1.2)}
{\small Capability · Canonical · current}\par

Manage the body of applicable evidence used for clinical-quality consideration without owning its originating activity.



Clinical Quality Evidence Management maintains the body of material applicable to quality consideration. Its purpose is to retain evidence in a form and context that can support review without treating availability as proof of quality. The originating activity and its authority remain distinguishable from the later quality purpose for which its information is used.



\textbf{Boundary}



Evidentiary use transfers neither source ownership nor authority to determine the activity being reviewed.


{\small Recorded unresolved matters: source-qualified matters; exact scope remains in the canonical record.}\par
{\footnotesize \href{../../../../continuum/strategy/capability/1b0a1c6945ed4fadba045ff4753c104c.md}{Canonical record: 1b0a1c6945ed4fadba045ff4753c104c}; full qualifications and provenance.}\par
\paragraph*{II.1.2.2.1.3 [CT3] Clinical Quality Reporting}
\phantomsection
\label{navigation:II.1.2.2.1.3}
\addcontentsline{toc}{paragraph}{II.1.2.2.1.3 [CT3] Clinical Quality Reporting}

{\small Historical predecessor: R.2.4 \hyperref[element:e5812fcdf3f5465b9b1d8ba536075c5f]{[CT3] Clinical Quality Reporting}; Canonical ID: \texttt{\seqsplit{e5812fcdf3f5465b9b1d8ba536075c5f}}. Superseded historical identity — direct correspondence · superseded.}\par
\label{element:70f34839c23042d59469e6ef96944057}
\index{[CT3] Clinical Quality Reporting (II.1.2.2.1.3)}
{\small Capability · Canonical · current}\par

Communicate established clinical-quality measures, assessments, findings and conclusions to authorised consumers.



Clinical Quality Reporting communicates established quality measures, assessments and findings to authorised consumers. It makes the outcome of quality consideration usable beyond the assessment context while retaining its basis and standing. Reporting therefore represents a communication responsibility, not another opportunity to originate or silently change the reported judgement.



\textbf{Boundary}



Communicating a quality concern does not perform the operational or clinical improvement response.


{\small Recorded unresolved matters: source-qualified matters; exact scope remains in the canonical record.}\par
{\footnotesize \href{../../../../continuum/strategy/capability/70f34839c23042d59469e6ef96944057.md}{Canonical record: 70f34839c23042d59469e6ef96944057}; full qualifications and provenance.}\par
\paragraph*{II.1.2.2.1.4 [CT3] Clinical Quality Review \& Assessment Management}
\phantomsection
\label{navigation:II.1.2.2.1.4}
\addcontentsline{toc}{paragraph}{II.1.2.2.1.4 [CT3] Clinical Quality Review \& Assessment Management}

{\small Historical predecessor: R.2.5 \hyperref[element:4b17ec92f0af4c4b9556912dbdd0f29f]{[CT3] Clinical Quality Review \& Assessment Management}; Canonical ID: \texttt{\seqsplit{4b17ec92f0af4c4b9556912dbdd0f29f}}. Superseded historical identity — direct correspondence · superseded.}\par
\label{element:8318c35b302b4dcaa0e795fbfdc93e3b}
\index{[CT3] Clinical Quality Review \& Assessment Management (II.1.2.2.1.4)}
{\small Capability · Canonical · current}\par

Review evidence against applicable criteria and establish a clinical-quality assessment, finding or conclusion.



This capability evaluates quality evidence against applicable criteria and establishes a finding, assessment or conclusion. It is the judgement-bearing responsibility within Clinical Quality Management, distinct from maintaining the material and from communicating the result. Making that distinction supports scrutiny of the criterion and evidentiary basis for a quality assessment.



\textbf{Boundary}



An assessment is not a report of itself or the subsequent improvement action it may inform.


{\small Recorded unresolved matters: source-qualified matters; exact scope remains in the canonical record.}\par
{\footnotesize \href{../../../../continuum/strategy/capability/8318c35b302b4dcaa0e795fbfdc93e3b.md}{Canonical record: 8318c35b302b4dcaa0e795fbfdc93e3b}; full qualifications and provenance.}\par
\subsubsection*{II.1.2.2.2 [CT2] Clinical Safety Management}
\phantomsection
\label{navigation:II.1.2.2.2}
\addcontentsline{toc}{subsubsection}{II.1.2.2.2 [CT2] Clinical Safety Management}

{\small Historical predecessor: R.2.10 \hyperref[element:650ccec030fa46cc8c9d6b508a841778]{[CT2] Clinical Safety Management}; Canonical ID: \texttt{\seqsplit{650ccec030fa46cc8c9d6b508a841778}}. Superseded historical identity — direct correspondence · superseded.}\par
\label{element:28d1bf3e514e4c29837ebedaf9e4917b}
\index{[CT2] Clinical Safety Management (II.1.2.2.2)}
{\small Capability · Canonical · current}\par

Manage the clinical-risk distinctions materially required for R1/R2.



Clinical Safety Management distinguishes recognising a safety risk from reducing future exposure and actively dealing with a developing or materialised concern. These responsibilities may use overlapping evidence, but their purposes differ: characterisation makes the risk intelligible, mitigation changes prospective risk, and response addresses the situation requiring action.



\textbf{Intent}



Support proportionate safety action grounded in the nature of the identified clinical risk.



\textbf{Boundary}



Monitoring evidence is not automatically risk identification. Other capabilities can perform required actions without transferring their responsibilities; the three children are not a mandatory sequence or exhaustive safety taxonomy.


{\small Recorded unresolved matters: source-qualified matters; exact scope remains in the canonical record.}\par
{\footnotesize \href{../../../../continuum/strategy/capability/28d1bf3e514e4c29837ebedaf9e4917b.md}{Canonical record: 28d1bf3e514e4c29837ebedaf9e4917b}; full qualifications and provenance.}\par
\paragraph*{II.1.2.2.2.1 [CT3] Clinical Risk Identification}
\phantomsection
\label{navigation:II.1.2.2.2.1}
\addcontentsline{toc}{paragraph}{II.1.2.2.2.1 [CT3] Clinical Risk Identification}

{\small Historical predecessor: R.2.7 \hyperref[element:690da55f1f5b4c44ba58cf67a717a564]{[CT3] Clinical Risk Identification}; Canonical ID: \texttt{\seqsplit{690da55f1f5b4c44ba58cf67a717a564}}. Superseded historical identity — direct correspondence · superseded.}\par
\label{element:d54c19c6efaf42af82ad71a4142149d4}
\index{[CT3] Clinical Risk Identification (II.1.2.2.2.1)}
{\small Capability · Canonical · current}\par

Recognise and characterise a clinical safety risk from applicable evidence.



Clinical Risk Identification recognises and characterises a safety risk from relevant evidence. It establishes the concern to which mitigation or active response may be directed, rather than assuming that every monitored deviation is already a clinical risk. Its value is the meaningful characterisation that makes proportionate safety consideration possible.



\textbf{Boundary}



Evidence observation alone does not identify risk or establish a mitigation or response decision.


{\small Recorded unresolved matters: source-qualified matters; exact scope remains in the canonical record.}\par
{\footnotesize \href{../../../../continuum/strategy/capability/d54c19c6efaf42af82ad71a4142149d4.md}{Canonical record: d54c19c6efaf42af82ad71a4142149d4}; full qualifications and provenance.}\par
\paragraph*{II.1.2.2.2.2 [CT3] Clinical Risk Mitigation Management}
\phantomsection
\label{navigation:II.1.2.2.2.2}
\addcontentsline{toc}{paragraph}{II.1.2.2.2.2 [CT3] Clinical Risk Mitigation Management}

{\small Historical predecessor: R.2.8 \hyperref[element:0bd3311d89af45bb935fde20e719369c]{[CT3] Clinical Risk Mitigation Management}; Canonical ID: \texttt{\seqsplit{0bd3311d89af45bb935fde20e719369c}}. Superseded historical identity — direct correspondence · superseded.}\par
\label{element:2c232ade21ec45619c3ee4e928bd97ee}
\index{[CT3] Clinical Risk Mitigation Management (II.1.2.2.2.2)}
{\small Capability · Canonical · current}\par

Manage prospective/preventative measures reducing identified risk likelihood, exposure or consequence.



Clinical Risk Mitigation Management addresses prospective measures intended to reduce an identified risk's likelihood, exposure or consequences. It gives prevention a distinct responsibility alongside risk characterisation and active response. Required measures can involve other capabilities without transferring the ownership of their actions or information to safety management.



\textbf{Boundary}



Mitigation differs from containing or remediating a concern that has materialised or requires active intervention.


{\small Recorded unresolved matters: source-qualified matters; exact scope remains in the canonical record.}\par
{\footnotesize \href{../../../../continuum/strategy/capability/2c232ade21ec45619c3ee4e928bd97ee.md}{Canonical record: 2c232ade21ec45619c3ee4e928bd97ee}; full qualifications and provenance.}\par
\paragraph*{II.1.2.2.2.3 [CT3] Clinical Risk Response Management}
\phantomsection
\label{navigation:II.1.2.2.2.3}
\addcontentsline{toc}{paragraph}{II.1.2.2.2.3 [CT3] Clinical Risk Response Management}

{\small Historical predecessor: R.2.9 \hyperref[element:20f719ab4b6741b393ef570b4d1b05df]{[CT3] Clinical Risk Response Management}; Canonical ID: \texttt{\seqsplit{20f719ab4b6741b393ef570b4d1b05df}}. Superseded historical identity — direct correspondence · superseded.}\par
\label{element:5e36cb137e4e4782b2ddbf78cdf0da7b}
\index{[CT3] Clinical Risk Response Management (II.1.2.2.2.3)}
{\small Capability · Canonical · current}\par

Manage active response, containment and remediation where risk has materialised, is emerging or requires active intervention.



Clinical Risk Response Management manages active intervention, containment or remediation where a safety risk is emerging, materialised or otherwise demands action. Its purpose is to address the developing situation rather than only reduce future exposure. The responsibility remains clinical-risk-specific while the actual actions may belong to other established abilities.



\textbf{Boundary}



It is not a general enterprise incident-management container and does not acquire the clinical authority of responding participants.


{\small Recorded unresolved matters: source-qualified matters; exact scope remains in the canonical record.}\par
{\footnotesize \href{../../../../continuum/strategy/capability/5e36cb137e4e4782b2ddbf78cdf0da7b.md}{Canonical record: 5e36cb137e4e4782b2ddbf78cdf0da7b}; full qualifications and provenance.}\par
\subsection*{II.1.2.3 [CT1] Community Health \& Wellbeing}
\phantomsection
\label{navigation:II.1.2.3}
\addcontentsline{toc}{subsection}{II.1.2.3 [CT1] Community Health \& Wellbeing}

{\small Historical predecessor: R.2.19 \hyperref[element:175174c5cce549d2b154a4f6704721dc]{[CT1] Community Health \& Wellbeing}; Canonical ID: \texttt{\seqsplit{175174c5cce549d2b154a4f6704721dc}}. Superseded historical identity — direct correspondence · superseded.}\par
\label{element:b326b4aaf02a41a986ed22d184cbc152}
\index{[CT1] Community Health \& Wellbeing (II.1.2.3)}
{\small Capability · Canonical · current}\par

Support community health and wellbeing, including the bounded preventative-health abilities required for R1/R2. Wider environmental health, outreach and social determinants remain enterprise context without inferred additional decomposition.



Community Health \& Wellbeing frames organised activity supporting health beyond a single person's episode of care. The governed decomposition focuses on preventative health: identifying needs, defining programmes and cohorts, educating, observing programme operation and reporting. Wider outreach, environmental and social-determinant concerns remain enterprise context rather than additional modelled children.



\textbf{Intent}



Support organised prevention with a clear relationship between the health problem, intended population and programme activity.



\textbf{Boundary}



Programme management does not acquire the clinical or appointment responsibilities needed to deliver a programme.



\textbf{Harmonia Architectural Relevance}



Programme and cohort distinctions materially shape the information, notifications and care enablement required. This supports the Relevant classification without assigning wider public-health operating authority to Harmonia.


{\footnotesize \href{../../../../continuum/strategy/capability/b326b4aaf02a41a986ed22d184cbc152.md}{Canonical record: b326b4aaf02a41a986ed22d184cbc152}; full qualifications and provenance.}\par
\subsubsection*{II.1.2.3.1 [CT2] Preventative Health}
\phantomsection
\label{navigation:II.1.2.3.1}
\addcontentsline{toc}{subsubsection}{II.1.2.3.1 [CT2] Preventative Health}

{\small Historical predecessor: R.2.18 \hyperref[element:80affa7bcd284ecbbcbc5092f8497ddf]{[CT2] Preventative Health}; Canonical ID: \texttt{\seqsplit{80affa7bcd284ecbbcbc5092f8497ddf}}. Superseded historical identity — direct correspondence · superseded.}\par
\label{element:81be04a2d58548a4aea94f9aae90feb6}
\index{[CT2] Preventative Health (II.1.2.3.1)}
{\small Capability · Canonical · current}\par

Address preventative-health needs through organised activity with sufficient R1/R2 distinctions.



Preventative Health organises a response to a health problem before adverse outcomes arise or worsen. The established children distinguish why attention is warranted, how an organised programme is managed, who it targets, the education provided, what programme operation shows and what is reported. This allows a screening or immunisation programme to be explained without creating a new capability for every programme type.



\textbf{Intent}



Connect a preventative purpose with an identifiable target population and observable organised activity.



\textbf{Boundary}



Programme performance differs from population-health assessment and individual clinical monitoring. Unmodelled specialist responsibilities are not silently assigned to the existing children.


{\small Recorded unresolved matters: source-qualified matters; exact scope remains in the canonical record.}\par
{\footnotesize \href{../../../../continuum/strategy/capability/81be04a2d58548a4aea94f9aae90feb6.md}{Canonical record: 81be04a2d58548a4aea94f9aae90feb6}; full qualifications and provenance.}\par
\paragraph*{II.1.2.3.1.1 [CT3] Preventative Health Cohort Management}
\phantomsection
\label{navigation:II.1.2.3.1.1}
\addcontentsline{toc}{paragraph}{II.1.2.3.1.1 [CT3] Preventative Health Cohort Management}

{\small Historical predecessor: R.2.12 \hyperref[element:fef479f0fd5048b19b2aec770f0233fe]{[CT3] Preventative Health Cohort Management}; Canonical ID: \texttt{\seqsplit{fef479f0fd5048b19b2aec770f0233fe}}. Superseded historical identity — direct correspondence · superseded.}\par
\label{element:97625bb28e5745d5994e1bffd4037188}
\index{[CT3] Preventative Health Cohort Management (II.1.2.3.1.1)}
{\small Capability · Canonical · current}\par

Establish and manage the target population, inclusion/exclusion criteria, membership, stratification and applicable entry/exit context.



Preventative Health Cohort Management establishes who belongs to the population targeted for a preventative purpose. Inclusion, exclusion, stratification and entry or exit context make membership explainable rather than a list detached from its rationale. It provides a governed population context through which programme activity and reporting can be understood.



\textbf{Boundary}



A population at risk is not automatically a defined cohort, and cohort membership is not an individual clinical determination.


{\small Recorded unresolved matters: source-qualified matters; exact scope remains in the canonical record.}\par
{\footnotesize \href{../../../../continuum/strategy/capability/97625bb28e5745d5994e1bffd4037188.md}{Canonical record: 97625bb28e5745d5994e1bffd4037188}; full qualifications and provenance.}\par
\paragraph*{II.1.2.3.1.2 [CT3] Preventative Health Education Services}
\phantomsection
\label{navigation:II.1.2.3.1.2}
\addcontentsline{toc}{paragraph}{II.1.2.3.1.2 [CT3] Preventative Health Education Services}

{\small Historical predecessor: R.2.13 \hyperref[element:dbe10959237f405caeadef2a5748303a]{[CT3] Preventative Health Education Services}; Canonical ID: \texttt{\seqsplit{dbe10959237f405caeadef2a5748303a}}. Superseded historical identity — direct correspondence · superseded.}\par
\label{element:7541d9b3b65d47d286dabdc093cb0eed}
\index{[CT3] Preventative Health Education Services (II.1.2.3.1.2)}
{\small Capability · Canonical · current}\par

Provide education intended to prevent adverse health outcomes and influence preventative knowledge or behaviour.



Preventative Health Education Services provides education intended to influence knowledge or behaviour so that adverse health outcomes can be prevented. Its purpose distinguishes it from general organisational communication: the educational contribution is directed to a preventative-health need. It can support organised programme activity without being equivalent to programme management or clinical intervention.



\textbf{Boundary}



Communicating educational material does not establish individual clinical assessment, treatment or programme effectiveness.


{\small Recorded unresolved matters: source-qualified matters; exact scope remains in the canonical record.}\par
{\footnotesize \href{../../../../continuum/strategy/capability/7541d9b3b65d47d286dabdc093cb0eed.md}{Canonical record: 7541d9b3b65d47d286dabdc093cb0eed}; full qualifications and provenance.}\par
\paragraph*{II.1.2.3.1.3 [CT3] Preventative Health Problem Identification}
\phantomsection
\label{navigation:II.1.2.3.1.3}
\addcontentsline{toc}{paragraph}{II.1.2.3.1.3 [CT3] Preventative Health Problem Identification}

{\small Historical predecessor: R.2.14 \hyperref[element:b35d3521141c419fb676a79b2ede40bd]{[CT3] Preventative Health Problem Identification}; Canonical ID: \texttt{\seqsplit{b35d3521141c419fb676a79b2ede40bd}}. Superseded historical identity — direct correspondence · superseded.}\par
\label{element:c1567022671f433fa1eeb37a9d156d92}
\index{[CT3] Preventative Health Problem Identification (II.1.2.3.1.3)}
{\small Capability · Canonical · current}\par

Identify and characterise a preventative-health problem warranting organised attention.



This capability establishes the preventative-health problem warranting organised attention. It gives an intended programme a reason grounded in the concern being addressed rather than assuming that a population statistic already defines the response. The identified problem can frame subsequent cohort and programme considerations without becoming an individual's clinical diagnosis.



\textbf{Boundary}



Population evidence informs problem identification but does not by itself establish a programme, cohort or individual care need.


{\small Recorded unresolved matters: source-qualified matters; exact scope remains in the canonical record.}\par
{\footnotesize \href{../../../../continuum/strategy/capability/c1567022671f433fa1eeb37a9d156d92.md}{Canonical record: c1567022671f433fa1eeb37a9d156d92}; full qualifications and provenance.}\par
\paragraph*{II.1.2.3.1.4 [CT3] Preventative Health Programme Management}
\phantomsection
\label{navigation:II.1.2.3.1.4}
\addcontentsline{toc}{paragraph}{II.1.2.3.1.4 [CT3] Preventative Health Programme Management}

{\small Historical predecessor: R.2.15 \hyperref[element:d71def9e76a544ff9761c672bb24f3be]{[CT3] Preventative Health Programme Management}; Canonical ID: \texttt{\seqsplit{d71def9e76a544ff9761c672bb24f3be}}. Superseded historical identity — direct correspondence · superseded.}\par
\label{element:563d36b179d34cb6913253d06e8d77f6}
\index{[CT3] Preventative Health Programme Management (II.1.2.3.1.4)}
{\small Capability · Canonical · current}\par

Establish and manage an organised response to a preventative-health need/objective without absorbing its delivery capabilities.



Preventative Health Programme Management establishes and manages an organised response to a preventative need or objective. It keeps the programme purpose and coordinated activity intelligible while allowing appointments, interventions and other delivery abilities to retain their own responsibilities. Screening and immunisation illustrate programmes rather than creating additional capability tiers.



\textbf{Boundary}



Managing the programme does not acquire the clinical or operational authority of the activities delivering it.


{\small Recorded unresolved matters: source-qualified matters; exact scope remains in the canonical record.}\par
{\footnotesize \href{../../../../continuum/strategy/capability/563d36b179d34cb6913253d06e8d77f6.md}{Canonical record: 563d36b179d34cb6913253d06e8d77f6}; full qualifications and provenance.}\par
\paragraph*{II.1.2.3.1.5 [CT3] Preventative Health Programme Monitoring}
\phantomsection
\label{navigation:II.1.2.3.1.5}
\addcontentsline{toc}{paragraph}{II.1.2.3.1.5 [CT3] Preventative Health Programme Monitoring}

{\small Historical predecessor: R.2.16 \hyperref[element:e05f6310f6d544efaca5568604f27dd8]{[CT3] Preventative Health Programme Monitoring}; Canonical ID: \texttt{\seqsplit{e05f6310f6d544efaca5568604f27dd8}}. Superseded historical identity — direct correspondence · superseded.}\par
\label{element:945b40ced8004e1cbf7c98c731d17daf}
\index{[CT3] Preventative Health Programme Monitoring (II.1.2.3.1.5)}
{\small Capability · Canonical · current}\par

Observe programme operation, reach, participation, coverage, progress and applicable performance.



This capability observes how a preventative programme is operating, including reach, participation, coverage and progress. It supplies an account of programme performance in the programme's own scope, rather than substituting population-health state for evidence of programme activity. That account can inform reporting or consideration of the programme without itself being an assurance conclusion.



\textbf{Boundary}



Programme monitoring is distinct from individual health monitoring, population assessment and independent assurance.


{\small Recorded unresolved matters: source-qualified matters; exact scope remains in the canonical record.}\par
{\footnotesize \href{../../../../continuum/strategy/capability/945b40ced8004e1cbf7c98c731d17daf.md}{Canonical record: 945b40ced8004e1cbf7c98c731d17daf}; full qualifications and provenance.}\par
\paragraph*{II.1.2.3.1.6 [CT3] Preventative Health Reporting}
\phantomsection
\label{navigation:II.1.2.3.1.6}
\addcontentsline{toc}{paragraph}{II.1.2.3.1.6 [CT3] Preventative Health Reporting}

{\small Historical predecessor: R.2.17 \hyperref[element:4faef45b4ebe4fa69e57c3d81bfea05f]{[CT3] Preventative Health Reporting}; Canonical ID: \texttt{\seqsplit{4faef45b4ebe4fa69e57c3d81bfea05f}}. Superseded historical identity — direct correspondence · superseded.}\par
\label{element:d2ef01c4cab348df822e802ab084faae}
\index{[CT3] Preventative Health Reporting (II.1.2.3.1.6)}
{\small Capability · Canonical · current}\par

Report established preventative-health problems, programmes, cohorts, activity, performance and outcomes.



Preventative Health Reporting communicates established information about preventative problems, programmes, cohorts, activity and outcomes. It brings these different subjects together for a reporting purpose while preserving their distinct meanings. A programme's reach and an observed population outcome can both be relevant without implying that one proves the other.



\textbf{Boundary}



Reporting does not originate the underlying problem, cohort or evaluation conclusion.


{\small Recorded unresolved matters: source-qualified matters; exact scope remains in the canonical record.}\par
{\footnotesize \href{../../../../continuum/strategy/capability/d2ef01c4cab348df822e802ab084faae.md}{Canonical record: d2ef01c4cab348df822e802ab084faae}; full qualifications and provenance.}\par
\subsection*{II.1.2.4 [CT1] Connected Health Services}
\phantomsection
\label{navigation:II.1.2.4}
\addcontentsline{toc}{subsection}{II.1.2.4 [CT1] Connected Health Services}

{\small Historical predecessor: R.2.20 \hyperref[element:48f01ef8b0e6431c85cf4f40cf8f94c3]{[CT1] Connected Health Services}; Canonical ID: \texttt{\seqsplit{48f01ef8b0e6431c85cf4f40cf8f94c3}}. Superseded historical identity — direct correspondence · superseded.}\par
\label{element:3ef397463e584737b367e93f0bf1d85d}
\index{[CT1] Connected Health Services (II.1.2.4)}
{\small Capability · Canonical · current}\par

Exchange health information appropriately and provide digitally mediated communication, notifications and collaborative interaction between health-service participants.



Connected Health Services represents the organisation's ability to exchange health information and communicate digitally between health-service participants. Its two children distinguish exchange of health-information artefacts from participant communication, notification and collaboration. Both support connected care, but a user viewing information through a screen is not thereby participating in a communication channel.



\textbf{Intent}



Enable appropriate information movement and interaction between otherwise separate participants and services.



\textbf{Boundary}



This organisational capability is not defined by an interface engine or messaging product. General presentation remains a separate enabling concern.



\textbf{Harmonia Architectural Relevance}



Information exchange and digitally mediated interaction are direct platform responsibilities within their assigned contracts, supporting the Core classification without transferring payload or clinical authority.


{\footnotesize \href{../../../../continuum/strategy/capability/3ef397463e584737b367e93f0bf1d85d.md}{Canonical record: 3ef397463e584737b367e93f0bf1d85d}; full qualifications and provenance.}\par
\subsubsection*{II.1.2.4.1 [CT2] Digital Communication Channel}
\phantomsection
\label{navigation:II.1.2.4.1}
\addcontentsline{toc}{subsubsection}{II.1.2.4.1 [CT2] Digital Communication Channel}

\label{element:f5defe141228490598115fe41db89bd7}
\index{[CT2] Digital Communication Channel (II.1.2.4.1)}
{\small Capability · Canonical · current}\par

Provide digitally mediated communication, notification and collaborative interaction between patients, practitioners and other health-service participants.



Digital Communication Channel supports communication between patients, practitioners and other participants through digitally mediated messages, notifications and collaborative interaction. Its organisational purpose is the exchange between participants: for example, a notification can connect someone to a care event without itself delivering care or changing the event's originating facts.



\textbf{Boundary}



A presentation interface used to exercise another capability is not automatically this communication ability. No CT3 set or general presentation BEC placement is established here.


{\small Recorded unresolved matters: source-qualified matters; exact scope remains in the canonical record.}\par
{\footnotesize \href{../../../../continuum/strategy/capability/f5defe141228490598115fe41db89bd7.md}{Canonical record: f5defe141228490598115fe41db89bd7}; full qualifications and provenance.}\par
\subsubsection*{II.1.2.4.2 [CT2] Health Information Exchange}
\phantomsection
\label{navigation:II.1.2.4.2}
\addcontentsline{toc}{subsubsection}{II.1.2.4.2 [CT2] Health Information Exchange}

\label{element:baad6b6117084311acbf5fc4044dfabf}
\index{[CT2] Health Information Exchange (II.1.2.4.2)}
{\small Capability · Canonical · current}\par

Exchange health information appropriately between participating parties, systems and services.



As a Business Capability, Health Information Exchange is the organisational ability to move health information appropriately between participating parties, systems and services. It explains the enterprise purpose served by exchange without defining that purpose as a protocol, server or software product. Its placement under Connected Health Services distinguishes information exchange from the broader interaction and notification purposes of Digital Communication Channel.



\textbf{Boundary}



No CT3 decomposition is established. The identically named BEC is a distinct enabling responsibility; the naming conflict remains unresolved, and this entry does not merge the two or assign the sharing-programme concern a parent.


{\small Recorded unresolved matters: source-qualified matters; exact scope remains in the canonical record.}\par
{\footnotesize \href{../../../../continuum/strategy/capability/baad6b6117084311acbf5fc4044dfabf.md}{Canonical record: baad6b6117084311acbf5fc4044dfabf}; full qualifications and provenance.}\par
\paragraph*{II.1.2.4.2.1 [CT3] Health Information Sharing Programme Management}
\phantomsection
\label{navigation:II.1.2.4.2.1}
\addcontentsline{toc}{paragraph}{II.1.2.4.2.1 [CT3] Health Information Sharing Programme Management}

\label{element:6c5c2cb625414b149ac7f2ccb4b8a897}
\index{[CT3] Health Information Sharing Programme Management (II.1.2.4.2.1)}
{\small Capability · Candidate · current}\par

Definition unresolved in the baseline. No definition has been synthesised from the name.



\textbf{Proposed overview}



Programme-level coordination supplies a plausible frame for the preserved name without replacing the established exchange abilities.



\textbf{Proposed definition}



Coordinate the governed scope, participants and change context of a health-information-sharing programme.



\textbf{Proposed boundary}



The element remains Candidate. The definition/placement elaboration is Proposed and does not adjudicate the unresolved RAD question or transfer source-domain authority.



\medskip\hrule\medskip

{\small Recorded unresolved matters: source-qualified matters; exact scope remains in the canonical record.}\par
{\footnotesize \href{../../../../continuum/strategy/capability/6c5c2cb625414b149ac7f2ccb4b8a897.md}{Canonical record: 6c5c2cb625414b149ac7f2ccb4b8a897}; full qualifications and provenance.}\par
\subsection*{II.1.2.5 [CT1] Corporate Resources \& Enterprise Services}
\phantomsection
\label{navigation:II.1.2.5}
\addcontentsline{toc}{subsection}{II.1.2.5 [CT1] Corporate Resources \& Enterprise Services}

{\small Historical predecessor: R.2.21 \hyperref[element:a2dbd74847a44f4c879983229f5c9d08]{[CT1] Corporate Resources \& Enterprise Services}; Canonical ID: \texttt{\seqsplit{a2dbd74847a44f4c879983229f5c9d08}}. Superseded historical identity — direct correspondence · superseded.}\par
\label{element:b79628f156e14a8cbb82a0dea2fc57ff}
\index{[CT1] Corporate Resources \& Enterprise Services (II.1.2.5)}
{\small Capability · Canonical · current}\par

Provide the bounded clinical-activity charging information and enterprise support context required for R1/R2, without assigning enterprise finance, HR, facilities or ICT operating authority to Harmonia.



This capability retains the financial and corporate-support context material to R1/R2. The financial branch concentrates on information about clinical activity relevant to charging or cost; corporate support covers administrative, physical, technical, people and communication conditions affecting service activity. The limited decomposition makes these dependencies visible without assigning a complete corporate enterprise architecture to Harmonia.



\textbf{Intent}



Support health-service activity with the required corporate context and legitimate activity information.



\textbf{Boundary}



Evidence about cost-relevant activity does not perform billing or accounting. Consuming corporate resource state does not acquire the authority to operate those resources.



\textbf{Harmonia Architectural Relevance}



Clinical activity evidence provides material financial support. Corporate constraints may also affect enablement, but contextual consumption does not classify every corporate descendant as Relevant or Core.


{\footnotesize \href{../../../../continuum/strategy/capability/b79628f156e14a8cbb82a0dea2fc57ff.md}{Canonical record: b79628f156e14a8cbb82a0dea2fc57ff}; full qualifications and provenance.}\par
\subsubsection*{II.1.2.5.1 [CT2] Corporate Support Services}
\phantomsection
\label{navigation:II.1.2.5.1}
\addcontentsline{toc}{subsubsection}{II.1.2.5.1 [CT2] Corporate Support Services}

\label{element:c8777bccbe304e36b7d9fbd955d67579}
\index{[CT2] Corporate Support Services (II.1.2.5.1)}
{\small Capability · Canonical · current}\par

Provide enterprise support whose resources, availability, state or constraints may materially affect service activity.



Corporate Support Services frames enterprise support whose resources, availability or constraints affect health-service activity. Its children retain recognisable administrative, facilities, infrastructure, communications, people and operational-oversight responsibilities. This explains why their context matters to service delivery while keeping the people who operate those concerns accountable for them.



\textbf{Intent}



Make material support conditions visible to health-service activity without confusing consumption with ownership.



\textbf{Boundary}



No further R1/R2 decomposition is established. Context from these abilities does not settle the collective service-readiness placement question.


{\small Recorded unresolved matters: source-qualified matters; exact scope remains in the canonical record.}\par
{\footnotesize \href{../../../../continuum/strategy/capability/c8777bccbe304e36b7d9fbd955d67579.md}{Canonical record: c8777bccbe304e36b7d9fbd955d67579}; full qualifications and provenance.}\par
\paragraph*{II.1.2.5.1.1 [CT3] Back-Office Services}
\phantomsection
\label{navigation:II.1.2.5.1.1}
\addcontentsline{toc}{paragraph}{II.1.2.5.1.1 [CT3] Back-Office Services}

\label{element:9f4d3705403647f4a1e25bc3dece035f}
\index{[CT3] Back-Office Services (II.1.2.5.1.1)}
{\small Capability · Canonical · current}\par

Provide administrative/support services on which health-service activity depends.



Back-Office Services provides the administrative and support activity on which health-service operations depend. Within Corporate Support Services it retains an enterprise context for dependencies that may affect whether service activity can proceed. The model stops at this broad responsibility rather than allocating each administrative function to Harmonia.



\textbf{Boundary}



Consumption of support status does not transfer administrative operating authority or establish a more detailed decomposition.


{\small Recorded unresolved matters: source-qualified matters; exact scope remains in the canonical record.}\par
{\footnotesize \href{../../../../continuum/strategy/capability/9f4d3705403647f4a1e25bc3dece035f.md}{Canonical record: 9f4d3705403647f4a1e25bc3dece035f}; full qualifications and provenance.}\par
\paragraph*{II.1.2.5.1.2 [CT3] Communications \& Publicity}
\phantomsection
\label{navigation:II.1.2.5.1.2}
\addcontentsline{toc}{paragraph}{II.1.2.5.1.2 [CT3] Communications \& Publicity}

\label{element:1fec53fb12444abcac41094dace1d3c2}
\index{[CT3] Communications \& Publicity (II.1.2.5.1.2)}
{\small Capability · Canonical · current}\par

Manage organisational communication/publication arrangements relevant to information availability/distribution.



Communications \& Publicity manages organisational publication and communication arrangements relevant to information availability and distribution. It provides enterprise communication context without being identical to clinical record exchange or digitally mediated participant interaction. The governed decomposition retains that broad support role without adding specific publicity or communications functions.



\textbf{Boundary}



Organisational communication does not originate the clinical facts or governed conclusions it may convey.


{\small Recorded unresolved matters: source-qualified matters; exact scope remains in the canonical record.}\par
{\footnotesize \href{../../../../continuum/strategy/capability/1fec53fb12444abcac41094dace1d3c2.md}{Canonical record: 1fec53fb12444abcac41094dace1d3c2}; full qualifications and provenance.}\par
\paragraph*{II.1.2.5.1.3 [CT3] Facilities Management}
\phantomsection
\label{navigation:II.1.2.5.1.3}
\addcontentsline{toc}{paragraph}{II.1.2.5.1.3 [CT3] Facilities Management}

\label{element:9c82e623bd174409b90d24af6b58ebf0}
\index{[CT3] Facilities Management (II.1.2.5.1.3)}
{\small Capability · Canonical · current}\par

Manage physical facilities/services and applicable resource, location, availability and facility-state context.



Facilities Management manages physical facilities and services, including the state and availability context relevant to health activity. It contributes the real-world conditions that location definitions alone cannot establish. This distinguishes a place recognised in a directory from a facility being operated and maintained for use.



\textbf{Boundary}



Harmonia's consumption of facility constraints does not make it the facilities operator or establish collective service readiness.


{\small Recorded unresolved matters: source-qualified matters; exact scope remains in the canonical record.}\par
{\footnotesize \href{../../../../continuum/strategy/capability/9c82e623bd174409b90d24af6b58ebf0.md}{Canonical record: 9c82e623bd174409b90d24af6b58ebf0}; full qualifications and provenance.}\par
\paragraph*{II.1.2.5.1.4 [CT3] IT Infrastructure Management}
\phantomsection
\label{navigation:II.1.2.5.1.4}
\addcontentsline{toc}{paragraph}{II.1.2.5.1.4 [CT3] IT Infrastructure Management}

\label{element:9a37413582314e0dbac006f82900858b}
\index{[CT3] IT Infrastructure Management (II.1.2.5.1.4)}
{\small Capability · Canonical · current}\par

Operate/manage underlying technology infrastructure and its availability.



IT Infrastructure Management operates the underlying technology infrastructure and its availability. It supplies an enterprise support context for health-service activity that relies on functioning technology. Its business outcome is operational infrastructure, distinct from governance of health-information meaning or the application responsibilities realised on that infrastructure.



\textbf{Boundary}



Infrastructure machinery does not define Harmonia's information semantics or acquire clinical authority.


{\small Recorded unresolved matters: source-qualified matters; exact scope remains in the canonical record.}\par
{\footnotesize \href{../../../../continuum/strategy/capability/9a37413582314e0dbac006f82900858b.md}{Canonical record: 9a37413582314e0dbac006f82900858b}; full qualifications and provenance.}\par
\paragraph*{II.1.2.5.1.5 [CT3] Operational Oversight}
\phantomsection
\label{navigation:II.1.2.5.1.5}
\addcontentsline{toc}{paragraph}{II.1.2.5.1.5 [CT3] Operational Oversight}

\label{element:d6f39d94cd2e4bf2ab91ff58ba604ab0}
\index{[CT3] Operational Oversight (II.1.2.5.1.5)}
{\small Capability · Canonical · current}\par

Observe broader operational conditions affecting the organisation’s ability to work or deliver services.



Operational Oversight observes broader operating conditions affecting the organisation's ability to work and deliver services. It provides context beyond an individual work item or care place, helping make material constraints visible. Within corporate support it remains a broad enterprise responsibility rather than an inferred owner for every operational-readiness question.



\textbf{Boundary}



Oversight does not resolve collective service-readiness placement or constitute independent assurance merely through observation.


{\small Recorded unresolved matters: source-qualified matters; exact scope remains in the canonical record.}\par
{\footnotesize \href{../../../../continuum/strategy/capability/d6f39d94cd2e4bf2ab91ff58ba604ab0.md}{Canonical record: d6f39d94cd2e4bf2ab91ff58ba604ab0}; full qualifications and provenance.}\par
\paragraph*{II.1.2.5.1.6 [CT3] People Management}
\phantomsection
\label{navigation:II.1.2.5.1.6}
\addcontentsline{toc}{paragraph}{II.1.2.5.1.6 [CT3] People Management}

\label{element:df8b9d309b354efcb15991b3d75ed852}
\index{[CT3] People Management (II.1.2.5.1.6)}
{\small Capability · Canonical · current}\par

Manage people across the organisation, distinct from practitioner-specific Health Workforce Management.



People Management addresses people across the organisation as an enterprise support responsibility. It retains a scope broader than the practitioner-specific standing, privilege and coverage concerns of Health Workforce Management. This distinction allows health-service activity to consume relevant people context without redefining all workforce administration as clinical governance.



\textbf{Boundary}



General people administration does not itself establish practitioner clinical qualification or privilege.


{\small Recorded unresolved matters: source-qualified matters; exact scope remains in the canonical record.}\par
{\footnotesize \href{../../../../continuum/strategy/capability/df8b9d309b354efcb15991b3d75ed852.md}{Canonical record: df8b9d309b354efcb15991b3d75ed852}; full qualifications and provenance.}\par
\subsubsection*{II.1.2.5.2 [CT2] Financial \& Commercial Services}
\phantomsection
\label{navigation:II.1.2.5.2}
\addcontentsline{toc}{subsubsection}{II.1.2.5.2 [CT2] Financial \& Commercial Services}

\label{element:290b574a02634f559b35a28daad0dea8}
\index{[CT2] Financial \& Commercial Services (II.1.2.5.2)}
{\small Capability · Canonical · current}\par

Provide the financial/commercial business context materially required for R1/R2.



Financial \& Commercial Services provides the financial business frame needed for the bounded R1/R2 responsibility. Its established child identifies and supplies information about chargeable or cost-relevant clinical activity. The purpose is to make applicable activity evidence available to the responsible financial context without expanding the model into pricing, claims or revenue management.



\textbf{Intent}



Support legitimate funding, financial and charging purposes with relevant clinical-activity information.



\textbf{Boundary}



Not every exchanged request or activity is chargeable. Provision of evidence does not establish a charging decision or an enterprise financial system.


{\small Recorded unresolved matters: source-qualified matters; exact scope remains in the canonical record.}\par
{\footnotesize \href{../../../../continuum/strategy/capability/290b574a02634f559b35a28daad0dea8.md}{Canonical record: 290b574a02634f559b35a28daad0dea8}; full qualifications and provenance.}\par
\paragraph*{II.1.2.5.2.1 [CT3] Clinical Activity Charging Information Provision}
\phantomsection
\label{navigation:II.1.2.5.2.1}
\addcontentsline{toc}{paragraph}{II.1.2.5.2.1 [CT3] Clinical Activity Charging Information Provision}

\label{element:ce9655c00f044ca996c9d4a1222785f5}
\index{[CT3] Clinical Activity Charging Information Provision (II.1.2.5.2.1)}
{\small Capability · Canonical · current}\par

Identify and provide information about chargeable or cost-relevant clinical/health-service activity where required for applicable funding, financial or charging purposes.



This capability identifies and supplies information about clinical or health-service activity relevant to legitimate charging, funding or cost purposes. Requests, service types, quantities and utilisation context may contribute evidence, but chargeability cannot be inferred from the fact that Harmonia mediated an interaction. Its outcome is usable activity information for the responsible financial context.



\textbf{Boundary}



Evidence provision does not perform pricing, billing, claims, accounting or revenue management.


{\small Recorded unresolved matters: source-qualified matters; exact scope remains in the canonical record.}\par
{\footnotesize \href{../../../../continuum/strategy/capability/ce9655c00f044ca996c9d4a1222785f5.md}{Canonical record: ce9655c00f044ca996c9d4a1222785f5}; full qualifications and provenance.}\par
\subsection*{II.1.2.6 [CT1] Diagnostic, Therapeutic \& Clinical Support Services}
\phantomsection
\label{navigation:II.1.2.6}
\addcontentsline{toc}{subsection}{II.1.2.6 [CT1] Diagnostic, Therapeutic \& Clinical Support Services}

{\small Historical predecessor: R.2.22 \hyperref[element:e003b9cd0730459a84775d2837d1d173]{[CT1] Diagnostic, Therapeutic \& Clinical Support Services}; Canonical ID: \texttt{\seqsplit{e003b9cd0730459a84775d2837d1d173}}. Superseded historical identity — direct correspondence · superseded.}\par
\label{element:4d98c3503a9f4890a677f09b0c3b6f82}
\index{[CT1] Diagnostic, Therapeutic \& Clinical Support Services (II.1.2.6)}
{\small Capability · Canonical · current}\par

The specialised clinical services that support direct patient care, including pathology and laboratory testing, diagnostic imaging (radiology), pharmacy and medication dispensing, and allied health therapies.



Specialised services contribute evidence, professional interpretation, medication provision and therapeutic support to individual care. The established decomposition distinguishes diagnostic provision from medication and pharmacy services, while retaining wider therapeutic and allied-health scope as enterprise context. The catalogue deliberately explains the modelled branches without presenting them as an exhaustive inventory of clinical support.



\textbf{Intent}



Provide the specialised contributions that treating participants need to assess and care for individuals.



\textbf{Boundary}



Producing diagnostic evidence, authorising a report, dispensing medication and administering treatment establish different outcomes and authorities.



\textbf{Harmonia Architectural Relevance}



Requests, specimen associations, reports and medication events require reliable correlation and exchange. Supporting those information flows does not make Harmonia the diagnostic practitioner, pharmacy or therapy provider.


{\footnotesize \href{../../../../continuum/strategy/capability/4d98c3503a9f4890a677f09b0c3b6f82.md}{Canonical record: 4d98c3503a9f4890a677f09b0c3b6f82}; full qualifications and provenance.}\par
\subsubsection*{II.1.2.6.1 [CT2] Diagnostic Support Provision}
\phantomsection
\label{navigation:II.1.2.6.1}
\addcontentsline{toc}{subsubsection}{II.1.2.6.1 [CT2] Diagnostic Support Provision}

\label{element:827910cf74ae4061afbda52dde3da106}
\index{[CT2] Diagnostic Support Provision (II.1.2.6.1)}
{\small Capability · Canonical · current}\par

Provide diagnostic services and manage their evidence, interpretation and governed reporting.



Diagnostic Support Provision covers the distinct responsibilities through which a diagnostic service obtains evidence and makes its findings available. Diagnostic scheduling, service performance, specimen stewardship, data management, interpretation, report governance and report distribution each establish a different outcome. Their separation allows preliminary, amended or final information to retain its professional and evidentiary standing without treating delivery as authorisation.



\textbf{Intent}



Make diagnostic evidence and professionally governed findings usable by the appropriate care participants.



\textbf{Boundary}



Diagnostic interpretation is limited to the service's professional scope; patient-level Clinical Determination considers the wider clinical context. The children do not prescribe a universal diagnostic workflow or establish a Diagnostic Order Management capability.


{\small Recorded unresolved matters: source-qualified matters; exact scope remains in the canonical record.}\par
{\footnotesize \href{../../../../continuum/strategy/capability/827910cf74ae4061afbda52dde3da106.md}{Canonical record: 827910cf74ae4061afbda52dde3da106}; full qualifications and provenance.}\par
\paragraph*{II.1.2.6.1.1 [CT3] Diagnostic Data Management}
\phantomsection
\label{navigation:II.1.2.6.1.1}
\addcontentsline{toc}{paragraph}{II.1.2.6.1.1 [CT3] Diagnostic Data Management}

\label{element:aeb8e6d2225a4df38e61ae116850b32c}
\index{[CT3] Diagnostic Data Management (II.1.2.6.1.1)}
{\small Capability · Canonical · current}\par

Manage acquired or produced diagnostic evidence as information, including identity, association, integrity, organisation, availability and lifecycle.



Diagnostic Data Management maintains acquired or produced diagnostic evidence as usable information. Its responsibility includes identity, association, integrity, organisation and lifecycle so that evidence can be found and understood in the right diagnostic context. It addresses the evidence record independently of the professional conclusion subsequently drawn from it.



\textbf{Boundary}



Managing diagnostic data neither interprets that data nor governs the authoritative expression of findings in a report.


{\small Recorded unresolved matters: source-qualified matters; exact scope remains in the canonical record.}\par
{\footnotesize \href{../../../../continuum/strategy/capability/aeb8e6d2225a4df38e61ae116850b32c.md}{Canonical record: aeb8e6d2225a4df38e61ae116850b32c}; full qualifications and provenance.}\par
\paragraph*{II.1.2.6.1.2 [CT3] Diagnostic Interpretation}
\phantomsection
\label{navigation:II.1.2.6.1.2}
\addcontentsline{toc}{paragraph}{II.1.2.6.1.2 [CT3] Diagnostic Interpretation}

\label{element:82ad867109f1438eb1b0916ed9cd820f}
\index{[CT3] Diagnostic Interpretation (II.1.2.6.1.2)}
{\small Capability · Canonical · current}\par

Establish what diagnostic evidence demonstrates within the diagnostic service’s professional scope.



Diagnostic Interpretation establishes what diagnostic evidence demonstrates within the service's professional scope. It turns evidence into diagnostic findings without claiming the wider patient-level judgement that considers all clinical circumstances. Separating interpretation from acquisition and reporting makes the professional contribution clear even when one service performs all three.



\textbf{Boundary}



Patient-level Clinical Determination remains distinct, as do report authorisation and distribution.


{\small Recorded unresolved matters: source-qualified matters; exact scope remains in the canonical record.}\par
{\footnotesize \href{../../../../continuum/strategy/capability/82ad867109f1438eb1b0916ed9cd820f.md}{Canonical record: 82ad867109f1438eb1b0916ed9cd820f}; full qualifications and provenance.}\par
\paragraph*{II.1.2.6.1.3 [CT3] Diagnostic Report Distribution}
\phantomsection
\label{navigation:II.1.2.6.1.3}
\addcontentsline{toc}{paragraph}{II.1.2.6.1.3 [CT3] Diagnostic Report Distribution}

\label{element:c8fdba3d18e143acb2bd77ae828afd4b}
\index{[CT3] Diagnostic Report Distribution (II.1.2.6.1.3)}
{\small Capability · Canonical · current}\par

Provide an established report to applicable authorised recipients or consuming contexts under applicable access, privacy, security and communication controls.



Diagnostic Report Distribution provides an already established report to applicable authorised recipients or consuming contexts. Its purpose is to make the governed findings available where they are needed, while respecting access, privacy, security and communication controls. Distribution retains the report's standing rather than originating or silently reinterpreting its findings.



\textbf{Boundary}



Delivery does not establish clinical review, report authorisation or a patient-level clinical decision.


{\small Recorded unresolved matters: source-qualified matters; exact scope remains in the canonical record.}\par
{\footnotesize \href{../../../../continuum/strategy/capability/c8fdba3d18e143acb2bd77ae828afd4b.md}{Canonical record: c8fdba3d18e143acb2bd77ae828afd4b}; full qualifications and provenance.}\par
\paragraph*{II.1.2.6.1.4 [CT3] Diagnostic Report Management}
\phantomsection
\label{navigation:II.1.2.6.1.4}
\addcontentsline{toc}{paragraph}{II.1.2.6.1.4 [CT3] Diagnostic Report Management}

\label{element:977ed170c34445ddb5095945f54f6907}
\index{[CT3] Diagnostic Report Management (II.1.2.6.1.4)}
{\small Capability · Canonical · current}\par

Establish and govern the authoritative expression of diagnostic findings, including applicable creation, review, amendment, verification, finalisation and authorisation/signing.



Diagnostic Report Management governs the authoritative expression of diagnostic findings. Creation, review, verification, amendment, finalisation and signing matter because they establish the standing of the report rather than merely its existence as data. The responsibility preserves a meaningful distinction between evidence, interpretation and the report that a service is prepared to authorise.



\textbf{Boundary}



Signing is behaviour within this capability, not another CT3. Sending a report does not establish its authorisation.


{\small Recorded unresolved matters: source-qualified matters; exact scope remains in the canonical record.}\par
{\footnotesize \href{../../../../continuum/strategy/capability/977ed170c34445ddb5095945f54f6907.md}{Canonical record: 977ed170c34445ddb5095945f54f6907}; full qualifications and provenance.}\par
\paragraph*{II.1.2.6.1.5 [CT3] Diagnostic Service Appointment Management}
\phantomsection
\label{navigation:II.1.2.6.1.5}
\addcontentsline{toc}{paragraph}{II.1.2.6.1.5 [CT3] Diagnostic Service Appointment Management}

\label{element:5410a0275dfa4460ad49100fa5a2f2ec}
\index{[CT3] Diagnostic Service Appointment Management (II.1.2.6.1.5)}
{\small Capability · Canonical · current}\par

Manage the diagnostic-service-specific scheduled opportunity, including preparation, modality, urgency, duration and other service-specific constraints.



Diagnostic Service Appointment Management makes a scheduled diagnostic opportunity credible in the service's own context. Preparation, modality, urgency and duration can materially affect whether an appointment is usable, so the responsibility is richer than recording a generic time slot. It supports service access while remaining separate from performance of the diagnostic activity.



\textbf{Boundary}



A diagnostic booking establishes neither evidence acquisition nor a report and is distinct from generic Care Appointment Management.


{\small Recorded unresolved matters: source-qualified matters; exact scope remains in the canonical record.}\par
{\footnotesize \href{../../../../continuum/strategy/capability/5410a0275dfa4460ad49100fa5a2f2ec.md}{Canonical record: 5410a0275dfa4460ad49100fa5a2f2ec}; full qualifications and provenance.}\par
\paragraph*{II.1.2.6.1.6 [CT3] Diagnostic Service Delivery}
\phantomsection
\label{navigation:II.1.2.6.1.6}
\addcontentsline{toc}{paragraph}{II.1.2.6.1.6 [CT3] Diagnostic Service Delivery}

\label{element:7a1ae953164d456c93637f8e35b4940e}
\index{[CT3] Diagnostic Service Delivery (II.1.2.6.1.6)}
{\small Capability · Canonical · current}\par

Perform the diagnostic activity that obtains or produces diagnostic evidence.



Diagnostic Service Delivery performs the activity that obtains or produces diagnostic evidence. Within Diagnostic Support Provision it is the point at which the diagnostic procedure actually occurs, rather than being scheduled or requested. The evidence obtained can then be managed and interpreted under the distinct responsibilities that give it professional meaning.



\textbf{Boundary}



Performance does not itself establish interpretation, report authorisation or report delivery.


{\small Recorded unresolved matters: source-qualified matters; exact scope remains in the canonical record.}\par
{\footnotesize \href{../../../../continuum/strategy/capability/7a1ae953164d456c93637f8e35b4940e.md}{Canonical record: 7a1ae953164d456c93637f8e35b4940e}; full qualifications and provenance.}\par
\paragraph*{II.1.2.6.1.7 [CT3] Diagnostic Specimen Management}
\phantomsection
\label{navigation:II.1.2.6.1.7}
\addcontentsline{toc}{paragraph}{II.1.2.6.1.7 [CT3] Diagnostic Specimen Management}

\label{element:93ab99a055314129984a18ea85db5cc2}
\index{[CT3] Diagnostic Specimen Management (II.1.2.6.1.7)}
{\small Capability · Canonical · current}\par

Maintain a specimen’s identity, integrity, suitability, custody and applicable lifecycle context.



Diagnostic Specimen Management preserves the specimen context needed to trust its use: which specimen it is, whether it remains suitable, and how integrity and custody are maintained. It concerns the physical diagnostic subject across its applicable lifecycle rather than only the results generated from it. This distinction protects the relationship between evidence and the material from which it arose.



\textbf{Boundary}



Specimen stewardship is separate from diagnostic-data management and interpretation of results.


{\small Recorded unresolved matters: source-qualified matters; exact scope remains in the canonical record.}\par
{\footnotesize \href{../../../../continuum/strategy/capability/93ab99a055314129984a18ea85db5cc2.md}{Canonical record: 93ab99a055314129984a18ea85db5cc2}; full qualifications and provenance.}\par
\subsubsection*{II.1.2.6.2 [CT2] Medication \& Pharmacy Services}
\phantomsection
\label{navigation:II.1.2.6.2}
\addcontentsline{toc}{subsubsection}{II.1.2.6.2 [CT2] Medication \& Pharmacy Services}

\label{element:ab2a4e6a7acc4868bc4e7b6c0d292cb4}
\index{[CT2] Medication \& Pharmacy Services (II.1.2.6.2)}
{\small Capability · Canonical · current}\par

Provide professional medication/pharmacy services under applicable controls.



This capability frames professional medication provision and medication-specific advice. Prescription dispensing, non-prescription provision, advice and dose administration aids answer different needs, even where exercised by one pharmacy. The decomposition keeps the authorised regimen, provision of medicine and actual administration separate so that a supply event is not mistaken for treatment performed.



\textbf{Intent}



Support safe, professionally controlled access to and use of medication within the established pharmacy responsibility.



\textbf{Boundary}



Advice does not confer prescribing authority, and dose aids do not establish a regimen. Whether specialised compounding/preparation warrants a separate capability remains unresolved under OI-BC-001.


{\small Recorded unresolved matters: source-qualified matters; exact scope remains in the canonical record.}\par
{\footnotesize \href{../../../../continuum/strategy/capability/ab2a4e6a7acc4868bc4e7b6c0d292cb4.md}{Canonical record: ab2a4e6a7acc4868bc4e7b6c0d292cb4}; full qualifications and provenance.}\par
\paragraph*{II.1.2.6.2.1 [CT3] Medication Advice}
\phantomsection
\label{navigation:II.1.2.6.2.1}
\addcontentsline{toc}{paragraph}{II.1.2.6.2.1 [CT3] Medication Advice}

\label{element:b867a2fde009437196a5fed682abf073}
\index{[CT3] Medication Advice (II.1.2.6.2.1)}
{\small Capability · Canonical · current}\par

Provide professional medication-specific advice about appropriate use, precautions, interactions, administration, adherence or safe handling.



Medication Advice provides professional guidance about medication use and its precautions, interactions, administration, adherence or handling. It supports understanding of how medication can be used safely without turning every advisory interaction into a prescribing decision. Within pharmacy services, the advice outcome is distinct from provision of the medicine itself.



\textbf{Boundary}



Advice does not confer authority to diagnose, prescribe or alter treatment; detailed scope remains subject to refinement.


{\small Recorded unresolved matters: source-qualified matters; exact scope remains in the canonical record.}\par
{\footnotesize \href{../../../../continuum/strategy/capability/b867a2fde009437196a5fed682abf073.md}{Canonical record: b867a2fde009437196a5fed682abf073}; full qualifications and provenance.}\par
\paragraph*{II.1.2.6.2.2 [CT3] Medication Compounding \& Preparation}
\phantomsection
\label{navigation:II.1.2.6.2.2}
\addcontentsline{toc}{paragraph}{II.1.2.6.2.2 [CT3] Medication Compounding \& Preparation}

\label{element:07380933da66427c865ca8e64d414e59}
\index{[CT3] Medication Compounding \& Preparation (II.1.2.6.2.2)}
{\small Capability · Proposed · current}\par

Prepare or compound a clinically required medication formulation under the applicable professional and product-governance mandate.



\textbf{Overview}



Prepare or compound a clinically required medication formulation under the applicable professional and product-governance mandate. The proposed decomposition separates this outcome from the other responsibilities within Medication \& Pharmacy Services.



\textbf{Intent}



Provide the required medication formulation when ordinary supply alone does not satisfy the authorised preparation need.



\textbf{Boundary}



Preparation differs from dispensing and administration; this proposal assigns no compounding activity to Harmonia.


{\footnotesize \href{../../../../continuum/strategy/capability/07380933da66427c865ca8e64d414e59.md}{Canonical record: 07380933da66427c865ca8e64d414e59}; full qualifications and provenance.}\par
\paragraph*{II.1.2.6.2.3 [CT3] Medication Dose Administration Aid Provision}
\phantomsection
\label{navigation:II.1.2.6.2.3}
\addcontentsline{toc}{paragraph}{II.1.2.6.2.3 [CT3] Medication Dose Administration Aid Provision}

\label{element:bd597608766b4d709703977c02cc5f07}
\index{[CT3] Medication Dose Administration Aid Provision (II.1.2.6.2.3)}
{\small Capability · Canonical · current}\par

Prepare and provide organised medication doses supporting administration according to an already established regimen.



This capability prepares and supplies organised doses to support administration under an already established medication regimen. Its purpose is to make that regimen easier to follow in practice while preserving the difference between organising doses and choosing therapy. The business responsibility is not defined by any particular commercial packaging arrangement.



\textbf{Boundary}



Preparing a dose aid does not establish the regimen or demonstrate that a dose was administered.


{\small Recorded unresolved matters: source-qualified matters; exact scope remains in the canonical record.}\par
{\footnotesize \href{../../../../continuum/strategy/capability/bd597608766b4d709703977c02cc5f07.md}{Canonical record: bd597608766b4d709703977c02cc5f07}; full qualifications and provenance.}\par
\paragraph*{II.1.2.6.2.4 [CT3] Non-Prescription Medication Provision}
\phantomsection
\label{navigation:II.1.2.6.2.4}
\addcontentsline{toc}{paragraph}{II.1.2.6.2.4 [CT3] Non-Prescription Medication Provision}

\label{element:7bbe4810fb974a8ab6d47b30d6475df3}
\index{[CT3] Non-Prescription Medication Provision (II.1.2.6.2.4)}
{\small Capability · Canonical · current}\par

Provide medication not requiring an applicable prescription, subject to professional, regulatory and safety controls.



Non-Prescription Medication Provision makes medication available where the applicable conditions do not require a prescription. It remains a professional, regulated and safety-constrained responsibility rather than an unrestricted supply transaction. Its distinction from prescription dispensing preserves the different authority basis under which provision occurs.



\textbf{Boundary}



Provision neither establishes diagnosis nor proves use or administration of the medication.


{\small Recorded unresolved matters: source-qualified matters; exact scope remains in the canonical record.}\par
{\footnotesize \href{../../../../continuum/strategy/capability/7bbe4810fb974a8ab6d47b30d6475df3.md}{Canonical record: 7bbe4810fb974a8ab6d47b30d6475df3}; full qualifications and provenance.}\par
\paragraph*{II.1.2.6.2.5 [CT3] Prescription Medication Dispensing}
\phantomsection
\label{navigation:II.1.2.6.2.5}
\addcontentsline{toc}{paragraph}{II.1.2.6.2.5 [CT3] Prescription Medication Dispensing}

\label{element:fef5b79eea77438d9d453a7daef744fc}
\index{[CT3] Prescription Medication Dispensing (II.1.2.6.2.5)}
{\small Capability · Canonical · current}\par

Professionally prepare and provide medication according to an applicable authorised prescription/order and dispensing requirements.



Prescription Medication Dispensing is the professional preparation and provision of medication against an applicable authorised prescription or order. It fulfils a pharmacy responsibility grounded in the prescribed intent and dispensing requirements. The resulting dispensing event provides different evidence from the prescribing decision and from later administration to the individual.



\textbf{Boundary}



Dispensing does not establish that medication was administered; unresolved specialised compounding must not be silently assigned here.


{\small Recorded unresolved matters: source-qualified matters; exact scope remains in the canonical record.}\par
{\footnotesize \href{../../../../continuum/strategy/capability/fef5b79eea77438d9d453a7daef744fc.md}{Canonical record: fef5b79eea77438d9d453a7daef744fc}; full qualifications and provenance.}\par
\subsection*{II.1.2.7 [CT1] Health Information \& Knowledge Management}
\phantomsection
\label{navigation:II.1.2.7}
\addcontentsline{toc}{subsection}{II.1.2.7 [CT1] Health Information \& Knowledge Management}

{\small Historical predecessor: R.2.23 \hyperref[element:a1c557bd263047d0b3ff566cf50b98c1]{[CT1] Health Information \& Knowledge Management}; Canonical ID: \texttt{\seqsplit{a1c557bd263047d0b3ff566cf50b98c1}}. Superseded historical identity — direct correspondence · superseded.}\par
\label{element:cca5d5abec204b339eb5b9bd924ed64e}
\index{[CT1] Health Information \& Knowledge Management (II.1.2.7)}
{\small Capability · Canonical · current}\par

Manage governed health-information content, structures, lifecycle, catalogues, compartments, correlation and retrieval, alongside reusable clinical knowledge and its governed standing.



This capability governs both health information and reusable clinical knowledge while preserving their different meanings. Health Information Management concerns content, structures, lifecycle, discovery context, grouping, correlation and retrieval. Clinical Knowledge Management concerns identifying and formulating reusable knowledge and establishing its governed standing. A source clinical fact can contribute to knowledge without becoming a general clinical conclusion merely through reuse.



\textbf{Intent}



Keep information interpretable and usable across time and context while making the standing and limits of reusable knowledge explicit.



\textbf{Boundary}



Managing or preserving information does not originate its clinical facts. Governing structures and binding concepts does not acquire the authority of the terminology they reference.



\textbf{Harmonia Architectural Relevance}



These are substantive information-management responsibilities within Harmonia's assigned boundary. The Core classification remains bounded by information authority, active/durable-state separation and clinical decision authority.


{\footnotesize \href{../../../../continuum/strategy/capability/cca5d5abec204b339eb5b9bd924ed64e.md}{Canonical record: cca5d5abec204b339eb5b9bd924ed64e}; full qualifications and provenance.}\par
\subsubsection*{II.1.2.7.1 [CT2] Clinical Knowledge Management}
\phantomsection
\label{navigation:II.1.2.7.1}
\addcontentsline{toc}{subsubsection}{II.1.2.7.1 [CT2] Clinical Knowledge Management}

\label{element:76b375248aee49aeb5a5044176a688d8}
\index{[CT2] Clinical Knowledge Management (II.1.2.7.1)}
{\small Capability · Canonical · current}\par

Identify, extract, formulate, associate, govern and share reusable clinical knowledge.



Clinical Knowledge Management turns potentially useful source material into explicit reusable knowledge with an intelligible basis and governed standing. Context identification and extraction preserve where material came from; formulation expresses applicability and rationale; concept association supplies semantic context; review and publication govern its intended use. These responsibilities distinguish source information, formulated knowledge and approved knowledge without prescribing a compulsory workflow.



\textbf{Intent}



Make reusable clinical knowledge available with the qualifications and approval standing needed to judge its use.



\textbf{Boundary}



Publication does not establish applicability to an individual or a clinical decision. Concept association does not redefine terminology, and formulation does not constitute approval.


{\small Recorded unresolved matters: source-qualified matters; exact scope remains in the canonical record.}\par
{\footnotesize \href{../../../../continuum/strategy/capability/76b375248aee49aeb5a5044176a688d8.md}{Canonical record: 76b375248aee49aeb5a5044176a688d8}; full qualifications and provenance.}\par
\paragraph*{II.1.2.7.1.1 [CT3] Clinical Knowledge Concept Association}
\phantomsection
\label{navigation:II.1.2.7.1.1}
\addcontentsline{toc}{paragraph}{II.1.2.7.1.1 [CT3] Clinical Knowledge Concept Association}

\label{element:3666bf40ec894f5aac3714dc76df392a}
\index{[CT3] Clinical Knowledge Concept Association (II.1.2.7.1.1)}
{\small Capability · Canonical · current}\par

Associate knowledge with applicable governed clinical concepts without acquiring concept authority.



Clinical Knowledge Concept Association relates reusable knowledge to the governed clinical concepts relevant to its meaning. It makes the semantic context of the knowledge intelligible without allowing local association to redefine the referenced concepts. The association assists interpretation while preserving the distinction between knowledge authority and concept authority.



\textbf{Boundary}



Associating knowledge with a terminology concept does not acquire terminology governance or establish approval of the knowledge.


{\small Recorded unresolved matters: source-qualified matters; exact scope remains in the canonical record.}\par
{\footnotesize \href{../../../../continuum/strategy/capability/3666bf40ec894f5aac3714dc76df392a.md}{Canonical record: 3666bf40ec894f5aac3714dc76df392a}; full qualifications and provenance.}\par
\paragraph*{II.1.2.7.1.2 [CT3] Clinical Knowledge Context Extraction}
\phantomsection
\label{navigation:II.1.2.7.1.2}
\addcontentsline{toc}{paragraph}{II.1.2.7.1.2 [CT3] Clinical Knowledge Context Extraction}

\label{element:0911997eb4584172afd10b1d37946061}
\index{[CT3] Clinical Knowledge Context Extraction (II.1.2.7.1.2)}
{\small Capability · Canonical · current}\par

Extract relevant material while preserving its source/context relationship.



Clinical Knowledge Context Extraction selects relevant material while retaining its relationship to the source and context. Its purpose is to make potential knowledge material available without stripping away the circumstances needed to understand it. Extracted material remains distinguishable from the reusable knowledge subsequently formulated from it.



\textbf{Boundary}



Extraction does not establish approval or turn source information into an individual clinical conclusion.


{\small Recorded unresolved matters: source-qualified matters; exact scope remains in the canonical record.}\par
{\footnotesize \href{../../../../continuum/strategy/capability/0911997eb4584172afd10b1d37946061.md}{Canonical record: 0911997eb4584172afd10b1d37946061}; full qualifications and provenance.}\par
\paragraph*{II.1.2.7.1.3 [CT3] Clinical Knowledge Context Identification}
\phantomsection
\label{navigation:II.1.2.7.1.3}
\addcontentsline{toc}{paragraph}{II.1.2.7.1.3 [CT3] Clinical Knowledge Context Identification}

\label{element:a0facf9015514223a928e39fcbbb6bee}
\index{[CT3] Clinical Knowledge Context Identification (II.1.2.7.1.3)}
{\small Capability · Canonical · current}\par

Identify sources/contexts potentially containing reusable clinical knowledge.



This capability identifies sources or contexts that may contain reusable clinical knowledge. It establishes where potentially useful material can be sought without assuming that all source information is already formulated knowledge. Within knowledge management it provides a contextual starting point rather than an approval or publication outcome.



\textbf{Boundary}



Identifying a source does not extract its material, establish reusable knowledge or determine its clinical applicability.


{\small Recorded unresolved matters: source-qualified matters; exact scope remains in the canonical record.}\par
{\footnotesize \href{../../../../continuum/strategy/capability/a0facf9015514223a928e39fcbbb6bee.md}{Canonical record: a0facf9015514223a928e39fcbbb6bee}; full qualifications and provenance.}\par
\paragraph*{II.1.2.7.1.4 [CT3] Clinical Knowledge Information Formulation}
\phantomsection
\label{navigation:II.1.2.7.1.4}
\addcontentsline{toc}{paragraph}{II.1.2.7.1.4 [CT3] Clinical Knowledge Information Formulation}

\label{element:05421be130ff4020b119c314aca46573}
\index{[CT3] Clinical Knowledge Information Formulation (II.1.2.7.1.4)}
{\small Capability · Canonical · current}\par

Formulate explicit reusable knowledge, including applicability, qualifications, rationale and evidence relationships.



This capability expresses reusable clinical knowledge with its applicability, qualifications, rationale and evidence relationships. It moves beyond extracted material by making the knowledge claim and its limits explicit. That formulation supports later review and use while retaining the distinction between a proposed knowledge expression and one approved for a governed purpose.



\textbf{Boundary}



Formulation does not approve knowledge, govern terminology or define health-information structures.


{\small Recorded unresolved matters: source-qualified matters; exact scope remains in the canonical record.}\par
{\footnotesize \href{../../../../continuum/strategy/capability/05421be130ff4020b119c314aca46573.md}{Canonical record: 05421be130ff4020b119c314aca46573}; full qualifications and provenance.}\par
\paragraph*{II.1.2.7.1.5 [CT3] Clinical Knowledge Publication \& Sharing}
\phantomsection
\label{navigation:II.1.2.7.1.5}
\addcontentsline{toc}{paragraph}{II.1.2.7.1.5 [CT3] Clinical Knowledge Publication \& Sharing}

\label{element:962a5e3048424ff7b2413e21a8e51f53}
\index{[CT3] Clinical Knowledge Publication \& Sharing (II.1.2.7.1.5)}
{\small Capability · Canonical · current}\par

Make approved knowledge available to applicable authorised consumers and use contexts.



Clinical Knowledge Publication \& Sharing makes approved knowledge available to applicable authorised consumers and contexts. It preserves the standing and limits under which knowledge was approved rather than reducing publication to unrestricted dissemination. Its contribution is availability for appropriate use, distinct from both approval and the individual clinical application that may follow.



\textbf{Boundary}



Published knowledge is not a patient-specific conclusion, decision or instruction merely because a clinician can access it.


{\small Recorded unresolved matters: source-qualified matters; exact scope remains in the canonical record.}\par
{\footnotesize \href{../../../../continuum/strategy/capability/962a5e3048424ff7b2413e21a8e51f53.md}{Canonical record: 962a5e3048424ff7b2413e21a8e51f53}; full qualifications and provenance.}\par
\paragraph*{II.1.2.7.1.6 [CT3] Clinical Knowledge Review \& Approval Management}
\phantomsection
\label{navigation:II.1.2.7.1.6}
\addcontentsline{toc}{paragraph}{II.1.2.7.1.6 [CT3] Clinical Knowledge Review \& Approval Management}

\label{element:ee04d00077d548a7a0a2dffd424c3d78}
\index{[CT3] Clinical Knowledge Review \& Approval Management (II.1.2.7.1.6)}
{\small Capability · Canonical · current}\par

Manage review, endorsement, approval/rejection, conditions and applicable history for intended governed use.



This capability manages the review and governed standing of reusable knowledge, including endorsement, rejection, conditions and history. It establishes how the knowledge may legitimately be treated for its intended use rather than assuming that formulation makes it trustworthy or approved. Temporal standing matters where a later review changes the conditions for use.



\textbf{Boundary}



Approval does not establish applicability to a particular individual or replace the responsible clinical judgement.


{\small Recorded unresolved matters: source-qualified matters; exact scope remains in the canonical record.}\par
{\footnotesize \href{../../../../continuum/strategy/capability/ee04d00077d548a7a0a2dffd424c3d78.md}{Canonical record: ee04d00077d548a7a0a2dffd424c3d78}; full qualifications and provenance.}\par
\subsubsection*{II.1.2.7.2 [CT2] Health Information Management}
\phantomsection
\label{navigation:II.1.2.7.2}
\addcontentsline{toc}{subsubsection}{II.1.2.7.2 [CT2] Health Information Management}

\label{element:eb2984f0ca79420e83012f9e06a005fb}
\index{[CT2] Health Information Management (II.1.2.7.2)}
{\small Capability · Canonical · current}\par

Manage health information as governed, semantically meaningful content.



Health Information Management maintains health information as meaningful governed content rather than isolated persisted values. Its children distinguish the content itself from the structures that explain it, lifecycle and history, discovery aids, contextual groupings, relationships and retrieval. These distinctions let an architect ask whether information exists, what it means, which state applies and whether it may be obtained without conflating those questions.



\textbf{Intent}



Preserve usable health information whose meaning, relationships and temporal standing can be understood.



\textbf{Boundary}



Discoverability is not permission, compartment membership is not an access decision, and correlation is not identity or clinical interpretation. Current state and immutable historical operational truth remain distinct.


{\small Recorded unresolved matters: source-qualified matters; exact scope remains in the canonical record.}\par
{\footnotesize \href{../../../../continuum/strategy/capability/eb2984f0ca79420e83012f9e06a005fb.md}{Canonical record: eb2984f0ca79420e83012f9e06a005fb}; full qualifications and provenance.}\par
\paragraph*{II.1.2.7.2.1 [CT3] Health Information Catalogue and Index Management}
\phantomsection
\label{navigation:II.1.2.7.2.1}
\addcontentsline{toc}{paragraph}{II.1.2.7.2.1 [CT3] Health Information Catalogue and Index Management}

\label{element:4c834101307546d39a0ac36b07358faf}
\index{[CT3] Health Information Catalogue and Index Management (II.1.2.7.2.1)}
{\small Capability · Canonical · current}\par

Manage descriptive catalogues/indexes through which information existence and location can be established.



This capability maintains the descriptive aids through which information existence and location can be established. It makes information discoverable without equating the catalogue entry with the content it describes. Within information management it answers where relevant information may be found before the separate question of obtaining it under an authorised context.



\textbf{Boundary}



Discoverability does not authorise access, retrieve the content or communicate it to a recipient.


{\small Recorded unresolved matters: source-qualified matters; exact scope remains in the canonical record.}\par
{\footnotesize \href{../../../../continuum/strategy/capability/4c834101307546d39a0ac36b07358faf.md}{Canonical record: 4c834101307546d39a0ac36b07358faf}; full qualifications and provenance.}\par
\paragraph*{II.1.2.7.2.2 [CT3] Health Information Compartmentalisation}
\phantomsection
\label{navigation:II.1.2.7.2.2}
\addcontentsline{toc}{paragraph}{II.1.2.7.2.2 [CT3] Health Information Compartmentalisation}

\label{element:660d78d6a34041939d1c290266d47f03}
\index{[CT3] Health Information Compartmentalisation (II.1.2.7.2.2)}
{\small Capability · Canonical · current}\par

Manage contextual association, partitioning and bounded information groupings.



Health Information Compartmentalisation establishes contextual associations and bounded groupings of information. It makes the relationship of information to a business context explicit, allowing groupings to carry meaning beyond physical storage location. That context can inform governed use while remaining distinct from sensitivity classification and a particular access decision.



\textbf{Boundary}



Compartment membership neither establishes permission nor replaces classification or terminology governance.


{\small Recorded unresolved matters: source-qualified matters; exact scope remains in the canonical record.}\par
{\footnotesize \href{../../../../continuum/strategy/capability/660d78d6a34041939d1c290266d47f03.md}{Canonical record: 660d78d6a34041939d1c290266d47f03}; full qualifications and provenance.}\par
\paragraph*{II.1.2.7.2.3 [CT3] Health Information Data Correlation}
\phantomsection
\label{navigation:II.1.2.7.2.3}
\addcontentsline{toc}{paragraph}{II.1.2.7.2.3 [CT3] Health Information Data Correlation}

\label{element:6c41e90a8cf24ed689d9d6a363c8d3fe}
\index{[CT3] Health Information Data Correlation (II.1.2.7.2.3)}
{\small Capability · Canonical · current}\par

Establish meaningful relationships between independently managed information items/sets.



Health Information Data Correlation establishes meaningful relationships between independently managed information items or sets. It helps information be understood together without requiring a merger of their identities, histories or originating authority. The responsibility is the association itself, rather than a clinical conclusion inferred from the associated content.



\textbf{Boundary}



Correlation does not establish person identity, clinical interpretation or ownership transfer between the related items.


{\small Recorded unresolved matters: source-qualified matters; exact scope remains in the canonical record.}\par
{\footnotesize \href{../../../../continuum/strategy/capability/6c41e90a8cf24ed689d9d6a363c8d3fe.md}{Canonical record: 6c41e90a8cf24ed689d9d6a363c8d3fe}; full qualifications and provenance.}\par
\paragraph*{II.1.2.7.2.4 [CT3] Health Information Data Management}
\phantomsection
\label{navigation:II.1.2.7.2.4}
\addcontentsline{toc}{paragraph}{II.1.2.7.2.4 [CT3] Health Information Data Management}

\label{element:9565887a92164432933ea97e95af2368}
\index{[CT3] Health Information Data Management (II.1.2.7.2.4)}
{\small Capability · Canonical · current}\par

Manage actual health-information content within applicable governed structures.



Health Information Data Management concerns the actual content maintained within applicable governed structures. It keeps the represented information distinct from definitions describing how that information should be organised and interpreted. Within Health Information Management, its responsibility is meaningful content rather than mere successful persistence of values.



\textbf{Boundary}



Managing content does not originate its clinical facts or acquire the authority of its semantic references.


{\small Recorded unresolved matters: source-qualified matters; exact scope remains in the canonical record.}\par
{\footnotesize \href{../../../../continuum/strategy/capability/9565887a92164432933ea97e95af2368.md}{Canonical record: 9565887a92164432933ea97e95af2368}; full qualifications and provenance.}\par
\paragraph*{II.1.2.7.2.5 [CT3] Health Information Data Retrieval}
\phantomsection
\label{navigation:II.1.2.7.2.5}
\addcontentsline{toc}{paragraph}{II.1.2.7.2.5 [CT3] Health Information Data Retrieval}

\label{element:082760fc4804489a81eaed3e4fd9daa9}
\index{[CT3] Health Information Data Retrieval (II.1.2.7.2.5)}
{\small Capability · Canonical · current}\par

Obtain identified information under an applicable authorised request/context.



Health Information Data Retrieval obtains identified information under an applicable authorised request and context. It supplies the content whose existence may already be known through a catalogue or index. The distinction matters because locating an item, obtaining it and presenting it to a person each establish different outcomes.



\textbf{Boundary}



Retrieval does not itself grant permission, provide participant communication or confer authority over the retrieved facts.


{\small Recorded unresolved matters: source-qualified matters; exact scope remains in the canonical record.}\par
{\footnotesize \href{../../../../continuum/strategy/capability/082760fc4804489a81eaed3e4fd9daa9.md}{Canonical record: 082760fc4804489a81eaed3e4fd9daa9}; full qualifications and provenance.}\par
\paragraph*{II.1.2.7.2.6 [CT3] Health Information Lifecycle Management}
\phantomsection
\label{navigation:II.1.2.7.2.6}
\addcontentsline{toc}{paragraph}{II.1.2.7.2.6 [CT3] Health Information Lifecycle Management}

\label{element:874cd3b294184f928dbb7f656abd97f0}
\index{[CT3] Health Information Lifecycle Management (II.1.2.7.2.6)}
{\small Capability · Canonical · current}\par

Govern state, amendment, correction, supersession, currency, retention and disposition while preserving provenance and historical meaning.



Health Information Lifecycle Management governs how information's state and temporal standing change through amendment, correction, supersession, retention and disposition. It allows current knowledge to change without retrospectively rewriting what was previously received, known or managed. That distinction supports interpretation of both the present representation and its historical operational meaning.



\textbf{Boundary}



Current authoritative state differs from historical truth; preserving either does not confer originating authority for the represented facts.


{\small Recorded unresolved matters: source-qualified matters; exact scope remains in the canonical record.}\par
{\footnotesize \href{../../../../continuum/strategy/capability/874cd3b294184f928dbb7f656abd97f0.md}{Canonical record: 874cd3b294184f928dbb7f656abd97f0}; full qualifications and provenance.}\par
\paragraph*{II.1.2.7.2.7 [CT3] Health Information Structure Management}
\phantomsection
\label{navigation:II.1.2.7.2.7}
\addcontentsline{toc}{paragraph}{II.1.2.7.2.7 [CT3] Health Information Structure Management}

\label{element:c3b3e39dd141448e96b7dd7b4c4516f2}
\index{[CT3] Health Information Structure Management (II.1.2.7.2.7)}
{\small Capability · Canonical · current}\par

Manage structures, element/reference definitions, constraints and semantic bindings giving content unambiguous meaning.



This capability manages the structures, element definitions, references and bindings that make content interpretable. It addresses how a represented item can have an unambiguous meaning without treating the structure as the clinical fact itself. Semantic bindings identify applicable governed references while preserving their external authority.



\textbf{Boundary}



Structure management is not terminology governance and does not prescribe FHIR resources, physical schemas or storage technology.


{\small Recorded unresolved matters: source-qualified matters; exact scope remains in the canonical record.}\par
{\footnotesize \href{../../../../continuum/strategy/capability/c3b3e39dd141448e96b7dd7b4c4516f2.md}{Canonical record: c3b3e39dd141448e96b7dd7b4c4516f2}; full qualifications and provenance.}\par
\subsection*{II.1.2.8 [CT1] Health Products \& Clinical Technology}
\phantomsection
\label{navigation:II.1.2.8}
\addcontentsline{toc}{subsection}{II.1.2.8 [CT1] Health Products \& Clinical Technology}

{\small Historical predecessor: R.2.24 \hyperref[element:a631a8600b8b49df9ef452c48c91349a]{[CT1] Health Products \& Clinical Technology}; Canonical ID: \texttt{\seqsplit{a631a8600b8b49df9ef452c48c91349a}}. Superseded historical identity — direct correspondence · superseded.}\par
\label{element:54388e3202c2403089b5c476c78c4567}
\index{[CT1] Health Products \& Clinical Technology (II.1.2.8)}
{\small Capability · Canonical · current}\par

Manage health products and clinical technology in the care environment, with distinct device inventory, functional/release, deployment, utilisation, integration and lifecycle responsibilities. Product procurement and supply context is retained without inferring further Health Product Management decomposition.



This capability manages health products and clinical technology as business concerns in the care environment. Its established clinical-technology branch distinguishes knowledge of the device estate, functional state, deployment, user enablement, integration and lifecycle disposition. Product procurement and supply remain relevant context without a newly inferred Health Product Management decomposition.



\textbf{Intent}



Keep the technology context needed for safe service activity and information exchange coherent and governed.



\textbf{Boundary}



Device-specific enablement, clinical privilege and use in care are different concerns. The seven established device children do not resolve the separate maintenance-membership discrepancy.



\textbf{Harmonia Architectural Relevance}



Integration provision and device-data integration materially affect exchanged information and its care context. Supply-chain, maintenance and clinical-use authority are not transferred to Harmonia.


{\footnotesize \href{../../../../continuum/strategy/capability/54388e3202c2403089b5c476c78c4567.md}{Canonical record: 54388e3202c2403089b5c476c78c4567}; full qualifications and provenance.}\par
\subsubsection*{II.1.2.8.1 [CT2] Clinical Technology Management}
\phantomsection
\label{navigation:II.1.2.8.1}
\addcontentsline{toc}{subsubsection}{II.1.2.8.1 [CT2] Clinical Technology Management}

\label{element:6c91476a45eb4432b1f90b47a4c2ad35}
\index{[CT2] Clinical Technology Management (II.1.2.8.1)}
{\small Capability · Canonical · current}\par

Manage clinical technology as a coherent business concern rather than only the device information Harmonia happens to use.



Clinical Technology Management treats a device as more than a data endpoint. Its responsibilities distinguish what device exists, what functionality and release it has, where and to whom it is deployed, how authorised users are supported, whether exchange is provisioned, what data is integrated and whether it should remain in service. These distinctions preserve the business meaning of the device estate beyond the subset of information used by Harmonia.



\textbf{Intent}



Maintain intelligible device capability, deployment and lifecycle context for clinical service use.



\textbf{Boundary}



Provisioning integration does not demonstrate that data has flowed. Monitoring, restoring device condition and deciding lifecycle disposition remain distinct; maintenance's tier membership requires adjudication.


{\small Recorded unresolved matters: source-qualified matters; exact scope remains in the canonical record.}\par
{\footnotesize \href{../../../../continuum/strategy/capability/6c91476a45eb4432b1f90b47a4c2ad35.md}{Canonical record: 6c91476a45eb4432b1f90b47a4c2ad35}; full qualifications and provenance.}\par
\paragraph*{II.1.2.8.1.1 [CT3] Clinical Device Data Integration}
\phantomsection
\label{navigation:II.1.2.8.1.1}
\addcontentsline{toc}{paragraph}{II.1.2.8.1.1 [CT3] Clinical Device Data Integration}

\label{element:382a6fe3546f491f940cba1c1e5df5ef}
\index{[CT3] Clinical Device Data Integration (II.1.2.8.1.1)}
{\small Capability · Canonical · current}\par

Acquire, exchange or convey device data through established integration capability.



Clinical Device Data Integration acquires, exchanges or conveys device data using an established integration capability. It represents actual information participation rather than the prior configuration enabling it. The distinction makes it possible to preserve the device and care context of received data while retaining separate responsibility for its subsequent interpretation.



\textbf{Boundary}



Data movement does not confer clinical meaning, originating authority or device-maintenance responsibility.


{\small Recorded unresolved matters: source-qualified matters; exact scope remains in the canonical record.}\par
{\footnotesize \href{../../../../continuum/strategy/capability/382a6fe3546f491f940cba1c1e5df5ef.md}{Canonical record: 382a6fe3546f491f940cba1c1e5df5ef}; full qualifications and provenance.}\par
\paragraph*{II.1.2.8.1.2 [CT3] Clinical Device Feature \& Release Management}
\phantomsection
\label{navigation:II.1.2.8.1.2}
\addcontentsline{toc}{paragraph}{II.1.2.8.1.2 [CT3] Clinical Device Feature \& Release Management}

\label{element:23d87ead4cdc40b785b006ca80f6d7b1}
\index{[CT3] Clinical Device Feature \& Release Management (II.1.2.8.1.2)}
{\small Capability · Canonical · current}\par

Manage functional/release state, including applicable software, firmware, configuration and enabled device characteristics.



This capability governs the functional state of a device, including applicable software, firmware, configuration and enabled characteristics. Two instances of the same model may therefore have materially different capabilities in service. The responsibility makes those differences intelligible without reducing device governance to physical inventory or location.



\textbf{Boundary}



Release state is distinct from maintenance condition, deployment and professional authority to use the device.


{\small Recorded unresolved matters: source-qualified matters; exact scope remains in the canonical record.}\par
{\footnotesize \href{../../../../continuum/strategy/capability/23d87ead4cdc40b785b006ca80f6d7b1.md}{Canonical record: 23d87ead4cdc40b785b006ca80f6d7b1}; full qualifications and provenance.}\par
\paragraph*{II.1.2.8.1.3 [CT3] Clinical Device Integration Provision}
\phantomsection
\label{navigation:II.1.2.8.1.3}
\addcontentsline{toc}{paragraph}{II.1.2.8.1.3 [CT3] Clinical Device Integration Provision}

\label{element:7abab1d7d68f48cea565df752fad34c4}
\index{[CT3] Clinical Device Integration Provision (II.1.2.8.1.3)}
{\small Capability · Canonical · current}\par

Establish and maintain the device’s capability and configuration to participate in applicable information exchanges.



Clinical Device Integration Provision establishes the device capability and configuration required to participate in applicable information exchanges. Its outcome is an integration-ready arrangement, which must remain distinguishable from data actually acquired or exchanged. This allows an architect to discuss whether connectivity has been enabled without assuming that a particular clinical feed has operated successfully.



\textbf{Boundary}



Provisioning does not demonstrate data integration, clinical interpretation or downstream authority for the data.


{\small Recorded unresolved matters: source-qualified matters; exact scope remains in the canonical record.}\par
{\footnotesize \href{../../../../continuum/strategy/capability/7abab1d7d68f48cea565df752fad34c4.md}{Canonical record: 7abab1d7d68f48cea565df752fad34c4}; full qualifications and provenance.}\par
\paragraph*{II.1.2.8.1.4 [CT3] Clinical Device Inventory Management}
\phantomsection
\label{navigation:II.1.2.8.1.4}
\addcontentsline{toc}{paragraph}{II.1.2.8.1.4 [CT3] Clinical Device Inventory Management}

\label{element:f8205d4cf8144f34b2da1139ff45c868}
\index{[CT3] Clinical Device Inventory Management (II.1.2.8.1.4)}
{\small Capability · Canonical · current}\par

Establish what devices exist and their managed asset identity, specification, ownership/custody and inventory context.



Clinical Device Inventory Management establishes the managed identity and inventory context of the device estate. Specification, ownership and custody make it possible to distinguish which devices exist from their current release, deployment or use. That stable asset context supports subsequent technology-management responsibilities without claiming their different outcomes.



\textbf{Boundary}



Inventory presence does not establish readiness, functional release state or current assignment to a care requirement.


{\small Recorded unresolved matters: source-qualified matters; exact scope remains in the canonical record.}\par
{\footnotesize \href{../../../../continuum/strategy/capability/f8205d4cf8144f34b2da1139ff45c868.md}{Canonical record: f8205d4cf8144f34b2da1139ff45c868}; full qualifications and provenance.}\par
\paragraph*{II.1.2.8.1.5 [CT3] Clinical Device Lifecycle Management}
\phantomsection
\label{navigation:II.1.2.8.1.5}
\addcontentsline{toc}{paragraph}{II.1.2.8.1.5 [CT3] Clinical Device Lifecycle Management}

\label{element:186667f5a3844d10a236af73bc91a0de}
\index{[CT3] Clinical Device Lifecycle Management (II.1.2.8.1.5)}
{\small Capability · Canonical · current}\par

Govern introduction, continued service, review, replacement, withdrawal, decommissioning and retirement within the overall device estate.



Clinical Device Lifecycle Management governs whether devices enter, remain in or leave the managed estate. Introduction, review, replacement, withdrawal and retirement concern the device's overall service disposition. This gives technology management a longer-lived responsibility than the device's current loan, release or data-exchange session.



\textbf{Boundary}



Lifecycle disposition differs from monitoring and restoring device condition; the maintenance-membership discrepancy remains unresolved.


{\small Recorded unresolved matters: source-qualified matters; exact scope remains in the canonical record.}\par
{\footnotesize \href{../../../../continuum/strategy/capability/186667f5a3844d10a236af73bc91a0de.md}{Canonical record: 186667f5a3844d10a236af73bc91a0de}; full qualifications and provenance.}\par
\paragraph*{II.1.2.8.1.6 [CT3] Clinical Device Loan \& Deployment Management}
\phantomsection
\label{navigation:II.1.2.8.1.6}
\addcontentsline{toc}{paragraph}{II.1.2.8.1.6 [CT3] Clinical Device Loan \& Deployment Management}

\label{element:5ac802cf322d408f82c4a28cd24cb43b}
\index{[CT3] Clinical Device Loan \& Deployment Management (II.1.2.8.1.6)}
{\small Capability · Canonical · current}\par

Manage device loan, issue, assignment, placement, return and deployment.



Clinical Device Loan \& Deployment Management manages where devices are placed and the issue, assignment, return or loan context under which they are available. It connects an identifiable device with its operational deployment rather than assuming ownership or inventory presence explains its current use. This supports an intelligible account of the device estate as placements change.



\textbf{Boundary}



Deployment is not device-specific user enablement, clinical use or evidence of clinical privilege.


{\small Recorded unresolved matters: source-qualified matters; exact scope remains in the canonical record.}\par
{\footnotesize \href{../../../../continuum/strategy/capability/5ac802cf322d408f82c4a28cd24cb43b.md}{Canonical record: 5ac802cf322d408f82c4a28cd24cb43b}; full qualifications and provenance.}\par
\paragraph*{II.1.2.8.1.7 [CT3] Clinical Device Maintenance Management}
\phantomsection
\label{navigation:II.1.2.8.1.7}
\addcontentsline{toc}{paragraph}{II.1.2.8.1.7 [CT3] Clinical Device Maintenance Management}

\label{element:731ea36d36d8466d84196ccda2e61cf9}
\index{[CT3] Clinical Device Maintenance Management (II.1.2.8.1.7)}
{\small Capability · Candidate · current}\par

Definition unresolved in the baseline. No definition has been synthesised from the name.



\textbf{Proposed overview}



Maintenance/restoration is materially distinct from overall lifecycle disposition and monitoring; include the Candidate concern as a proposed eighth CT3.



\textbf{Proposed definition}



Maintain or restore acceptable device condition through testing, calibration, inspection, maintenance and repair.



\textbf{Proposed boundary}



The element remains Candidate. The definition/placement elaboration is Proposed and does not adjudicate the unresolved RAD question or transfer source-domain authority.



\medskip\hrule\medskip

{\small Recorded unresolved matters: source-qualified matters; exact scope remains in the canonical record.}\par
{\footnotesize \href{../../../../continuum/strategy/capability/731ea36d36d8466d84196ccda2e61cf9.md}{Canonical record: 731ea36d36d8466d84196ccda2e61cf9}; full qualifications and provenance.}\par
\paragraph*{II.1.2.8.1.8 [CT3] Clinical Device Utilisation Support}
\phantomsection
\label{navigation:II.1.2.8.1.8}
\addcontentsline{toc}{paragraph}{II.1.2.8.1.8 [CT3] Clinical Device Utilisation Support}

\label{element:0100cdb80a0c41c5930b1ca66248adf9}
\index{[CT3] Clinical Device Utilisation Support (II.1.2.8.1.8)}
{\small Capability · Canonical · current}\par

Enable and support otherwise appropriately authorised users to use the particular device safely and effectively.



Clinical Device Utilisation Support helps otherwise appropriately authorised users use a particular device safely and effectively. It represents the device-specific enablement needed beyond general professional qualifications. Within technology management it addresses the practical support for use, without making that support the clinical action performed with the device.



\textbf{Boundary}



Professional qualification or privilege alone does not establish device-specific enablement; enablement does not confer clinical authority.


{\small Recorded unresolved matters: source-qualified matters; exact scope remains in the canonical record.}\par
{\footnotesize \href{../../../../continuum/strategy/capability/0100cdb80a0c41c5930b1ca66248adf9.md}{Canonical record: 0100cdb80a0c41c5930b1ca66248adf9}; full qualifications and provenance.}\par
\subsection*{II.1.2.9 [CT1] Health Research \& Innovation}
\phantomsection
\label{navigation:II.1.2.9}
\addcontentsline{toc}{subsection}{II.1.2.9 [CT1] Health Research \& Innovation}

{\small Historical predecessor: R.2.25 \hyperref[element:2cc31384c8864c73989025cbe96dcfa7]{[CT1] Health Research \& Innovation}; Canonical ID: \texttt{\seqsplit{2cc31384c8864c73989025cbe96dcfa7}}. Superseded historical identity — direct correspondence · superseded.}\par
\label{element:582bf62145fd42118e5d57bb714841e8}
\index{[CT1] Health Research \& Innovation (II.1.2.9)}
{\small Capability · Canonical · current}\par

Provide existing health information for authorised research under applicable constraints and evaluate health innovation; deliberately under-decomposed for R1/R2.



This capability frames the deliberately bounded research and innovation responsibilities material to Harmonia's R1/R2 architecture. Research Information Support concerns provision of existing information for authorised research; Health Innovation concerns evaluation of proposed or emerging innovation. These establish useful enterprise context without implying a complete research organisation or innovation lifecycle.



\textbf{Intent}



Support legitimate research information needs and informed consideration of changes to health-service or clinical practice.



\textbf{Boundary}



Information provision does not own trials, ethics approval or specialist research analysis. Evaluating innovation does not authorise adoption or implementation.



\textbf{Harmonia Architectural Relevance}



Governed provision of existing information materially enables authorised research. Innovation evaluation has no separate platform assignment; the CT1 classification must not be inherited by each descendant.


{\footnotesize \href{../../../../continuum/strategy/capability/582bf62145fd42118e5d57bb714841e8.md}{Canonical record: 582bf62145fd42118e5d57bb714841e8}; full qualifications and provenance.}\par
\subsubsection*{II.1.2.9.1 [CT2] Health Innovation}
\phantomsection
\label{navigation:II.1.2.9.1}
\addcontentsline{toc}{subsubsection}{II.1.2.9.1 [CT2] Health Innovation}

\label{element:f263989e25c9469cac1bbad6491432f0}
\index{[CT2] Health Innovation (II.1.2.9.1)}
{\small Capability · Canonical · current}\par

Consider health-service and clinical innovation within the R1/R2 business context.



Health Innovation provides the business frame for considering changes in health-service and clinical practice. Its established evaluation child examines suitability, value and implications so that an innovation can be understood in the context in which it might be used. The model stops at that evaluative responsibility rather than presenting a complete pipeline from idea to operational adoption.



\textbf{Intent}



Make consideration of innovation informed by its prospective benefits, consequences and suitability.



\textbf{Boundary}



Evaluation does not establish an adoption decision, professional authority or an implementation programme.


{\small Recorded unresolved matters: source-qualified matters; exact scope remains in the canonical record.}\par
{\footnotesize \href{../../../../continuum/strategy/capability/f263989e25c9469cac1bbad6491432f0.md}{Canonical record: f263989e25c9469cac1bbad6491432f0}; full qualifications and provenance.}\par
\paragraph*{II.1.2.9.1.1 [CT3] Health Innovation Evaluation}
\phantomsection
\label{navigation:II.1.2.9.1.1}
\addcontentsline{toc}{paragraph}{II.1.2.9.1.1 [CT3] Health Innovation Evaluation}

\label{element:adc555bcb54e4de290869c7deebd128c}
\index{[CT3] Health Innovation Evaluation (II.1.2.9.1.1)}
{\small Capability · Canonical · current}\par

Evaluate proposed/emerging innovation for suitability, value, implications and potential adoption.



Health Innovation Evaluation is the organisation's ability to form a considered view of a proposed or emerging innovation before treating it as suitable for use. Within Health Innovation it brings suitability, value and implications into the same evaluative context, so that potential benefit can be considered alongside consequences for the health-service or clinical setting. The outcome supports informed consideration of adoption rather than making adoption inevitable.



\textbf{Intent}



Give decision makers an intelligible basis for considering whether an innovation merits further adoption consideration.



\textbf{Boundary}



Evaluation does not establish adoption, authorise a clinical change or create an implementation responsibility. The single established child does not imply a complete innovation lifecycle.


{\small Recorded unresolved matters: source-qualified matters; exact scope remains in the canonical record.}\par
{\footnotesize \href{../../../../continuum/strategy/capability/adc555bcb54e4de290869c7deebd128c.md}{Canonical record: adc555bcb54e4de290869c7deebd128c}; full qualifications and provenance.}\par
\subsubsection*{II.1.2.9.2 [CT2] Research Information Support}
\phantomsection
\label{navigation:II.1.2.9.2}
\addcontentsline{toc}{subsubsection}{II.1.2.9.2 [CT2] Research Information Support}

\label{element:ef18f999eef3488ba584d263a551732b}
\index{[CT2] Research Information Support (II.1.2.9.2)}
{\small Capability · Canonical · current}\par

Provide the existing health information needed for authorised research within the deliberately bounded R1/R2 responsibility.



Research Information Support represents the ability to make existing health information available for a legitimate research purpose. Its single established child concentrates on governed provision, preserving the constraints under which information can be used rather than treating technical availability as research permission. The narrow decomposition reflects R1/R2 sufficiency rather than a claim to encompass every form of research support.



\textbf{Intent}



Support authorised research with information supplied under the applicable ethical, privacy and governance conditions.



\textbf{Boundary}



Research governance, trial oversight and specialist analysis remain unmodelled and are not absorbed into provision.


{\small Recorded unresolved matters: source-qualified matters; exact scope remains in the canonical record.}\par
{\footnotesize \href{../../../../continuum/strategy/capability/ef18f999eef3488ba584d263a551732b.md}{Canonical record: ef18f999eef3488ba584d263a551732b}; full qualifications and provenance.}\par
\paragraph*{II.1.2.9.2.1 [CT3] Governed Research Information Provision}
\phantomsection
\label{navigation:II.1.2.9.2.1}
\addcontentsline{toc}{paragraph}{II.1.2.9.2.1 [CT3] Governed Research Information Provision}

\label{element:6907401e6786470083c1d34a65aa225b}
\index{[CT3] Governed Research Information Provision (II.1.2.9.2.1)}
{\small Capability · Canonical · current}\par

Provide existing health information for an authorised research purpose under applicable ethics, governance, privacy, consent, security and information-use constraints.



Governed Research Information Provision makes existing health information available for an authorised research purpose under its applicable constraints. It represents a substantive information contribution while keeping ethics, governance, consent, privacy, security and permitted use material to provision. The responsibility concerns supplying existing information legitimately, rather than running the research that uses it.



\textbf{Boundary}



Availability is not research authorisation. Trials, genomic analysis and research-governance responsibilities are not subsumed by this deliberately narrow capability.


{\small Recorded unresolved matters: source-qualified matters; exact scope remains in the canonical record.}\par
{\footnotesize \href{../../../../continuum/strategy/capability/6907401e6786470083c1d34a65aa225b.md}{Canonical record: 6907401e6786470083c1d34a65aa225b}; full qualifications and provenance.}\par
\subsection*{II.1.2.10 [CT1] Health Rights, Advocacy \& Participation}
\phantomsection
\label{navigation:II.1.2.10}
\addcontentsline{toc}{subsection}{II.1.2.10 [CT1] Health Rights, Advocacy \& Participation}

{\small Historical predecessor: R.2.26 \hyperref[element:58376e5c53bf4af0a54f99f21ac52666]{[CT1] Health Rights, Advocacy \& Participation}; Canonical ID: \texttt{\seqsplit{58376e5c53bf4af0a54f99f21ac52666}}. Superseded historical identity — direct correspondence · superseded.}\par
\label{element:7ab5c82a500b45f39267883a6e78c482}
\index{[CT1] Health Rights, Advocacy \& Participation (II.1.2.10)}
{\small Capability · Canonical · current}\par

Protecting patient rights, facilitating informed consent, managing advance care directives, supporting patient advocacy, and enabling consumer participation in health decisions.



This capability represents the organisation's responsibility to protect patients' rights and enable their voice in health decisions. Consent, advance care directives, advocacy and consumer participation concern the legitimacy and personal consequences of care, rather than merely the availability of clinical information. The governed model retains this broad enterprise responsibility without an established subordinate capability set.



\textbf{Intent}



Support informed participation and respect for the individual's rights in decisions affecting their health.



\textbf{Boundary}



Recording a directive or restricting access can support rights without performing advocacy or managing consumer participation.



\textbf{Harmonia Architectural Relevance}



The established Adjacent classification reflects the use and supply of consent and record-access evidence. It does not assign the underlying rights, advocacy or participation activity to Harmonia.


{\footnotesize \href{../../../../continuum/strategy/capability/7ab5c82a500b45f39267883a6e78c482.md}{Canonical record: 7ab5c82a500b45f39267883a6e78c482}; full qualifications and provenance.}\par
\subsection*{II.1.2.11 [CT1] Health Service Assurance}
\phantomsection
\label{navigation:II.1.2.11}
\addcontentsline{toc}{subsection}{II.1.2.11 [CT1] Health Service Assurance}

{\small Historical predecessor: R.2.27 \hyperref[element:6a82dc1423374e809dbf542eeaa69482]{[CT1] Health Service Assurance}; Canonical ID: \texttt{\seqsplit{6a82dc1423374e809dbf542eeaa69482}}. Superseded historical identity — direct correspondence · superseded.}\par
\label{element:14860102964c441681d41ec3da423fe0}
\index{[CT1] Health Service Assurance (II.1.2.11)}
{\small Capability · Canonical · current}\par

Establish and maintain justified confidence that health services can be, and are being, delivered in accordance with applicable expectations, obligations, constraints and required levels of performance.



Health Service Assurance establishes justified confidence in the ability to deliver health services appropriately. Specification defines how confidence will be established; evaluation determines what the available evidence supports. Its subject is macro-level service delivery, with workforce, clinical, information, security and other concerns contributing evidence or criteria without becoming separate assurance trees.



Where independent governed assurance is required, the subject cannot control its own evaluation or conclusion. Independence concerns control of the assurance judgement rather than requiring a separate organisation or technology installation.



\textbf{Intent}



Distinguish justified confidence, qualified conclusions and insufficient evidence so that assurance communicates what is actually supported.



\textbf{Boundary}



Assurance manages its own evaluation and does not assign, remediate or progress the subject activity. Harmonia does not thereby assume responsibility for clinical adequacy or professional practice.



\textbf{Harmonia Architectural Relevance}



Approved independent governed assurance materially supports service confidence. The three established assurance BECs provide bounded support; they are not remapped to the newer BC descendants and are not assigned implementation components.


{\footnotesize \href{../../../../continuum/strategy/capability/14860102964c441681d41ec3da423fe0.md}{Canonical record: 14860102964c441681d41ec3da423fe0}; full qualifications and provenance.}\par
\subsubsection*{II.1.2.11.1 [CT2] Assurance Evaluation}
\phantomsection
\label{navigation:II.1.2.11.1}
\addcontentsline{toc}{subsubsection}{II.1.2.11.1 [CT2] Assurance Evaluation}

\label{element:d1b20e0fb91e4bd1950ad57b84e2cafd}
\index{[CT2] Assurance Evaluation (II.1.2.11.1)}
{\small Capability · Canonical · current}\par

Determine what justified confidence the applicable evidence supports.



Assurance Evaluation determines what confidence the evidence legitimately supports under the established approach and criteria. Evidence Assessment considers the material against that basis; Conclusion Determination establishes the resulting position, including qualifications and unresolved concerns. The distinction prevents accumulated evidence or execution completion from being mistaken for a justified assurance conclusion.



\textbf{Intent}



State the supported level and limits of confidence without concealing insufficient evidence.



\textbf{Boundary}



Insufficient evidence establishes neither assurance nor non-assurance. Operational response remains with the activity's responsible capability, and the existing assurance BEC ownership is unchanged.


{\small Recorded unresolved matters: source-qualified matters; exact scope remains in the canonical record.}\par
{\footnotesize \href{../../../../continuum/strategy/capability/d1b20e0fb91e4bd1950ad57b84e2cafd.md}{Canonical record: d1b20e0fb91e4bd1950ad57b84e2cafd}; full qualifications and provenance.}\par
\paragraph*{II.1.2.11.1.1 [CT3] Assurance Conclusion Determination}
\phantomsection
\label{navigation:II.1.2.11.1.1}
\addcontentsline{toc}{paragraph}{II.1.2.11.1.1 [CT3] Assurance Conclusion Determination}

\label{element:10f538feae4042829fafeae472882128}
\index{[CT3] Assurance Conclusion Determination (II.1.2.11.1.1)}
{\small Capability · Canonical · current}\par

Establish the legitimately supported conclusion, including insufficiency, qualification and unresolved concerns where applicable.



Assurance Conclusion Determination establishes what confidence is legitimately supported by the assessment. It preserves qualifications and unresolved concerns rather than forcing incomplete evidence into a positive or negative result. Within Assurance Evaluation it is the responsibility that gives the conclusion its standing, distinct from gathering or assessing the supporting material.



\textbf{Boundary}



Insufficiency proves neither assurance nor non-assurance, and the conclusion does not manage the subject's operational response.


{\small Recorded unresolved matters: source-qualified matters; exact scope remains in the canonical record.}\par
{\footnotesize \href{../../../../continuum/strategy/capability/10f538feae4042829fafeae472882128.md}{Canonical record: 10f538feae4042829fafeae472882128}; full qualifications and provenance.}\par
\paragraph*{II.1.2.11.1.2 [CT3] Assurance Evidence Assessment}
\phantomsection
\label{navigation:II.1.2.11.1.2}
\addcontentsline{toc}{paragraph}{II.1.2.11.1.2 [CT3] Assurance Evidence Assessment}

\label{element:1b295763b37f4622bd1bb80c6ea3741c}
\index{[CT3] Assurance Evidence Assessment (II.1.2.11.1.2)}
{\small Capability · Canonical · current}\par

Evaluate available evidence against the established assurance approach and criteria.



Assurance Evidence Assessment evaluates available evidence against the established approach and criteria. It produces the assessed basis from which a legitimate conclusion may be determined, retaining qualifications and evidentiary limits. The responsibility is evaluative rather than the operational production of the source information or remediation of what it reveals.



\textbf{Boundary}



Assessment is distinct from conclusion determination; unavailable or insufficient evidence must remain explicit.


{\small Recorded unresolved matters: source-qualified matters; exact scope remains in the canonical record.}\par
{\footnotesize \href{../../../../continuum/strategy/capability/1b295763b37f4622bd1bb80c6ea3741c.md}{Canonical record: 1b295763b37f4622bd1bb80c6ea3741c}; full qualifications and provenance.}\par
\subsubsection*{II.1.2.11.2 [CT2] Health Service Assurance Specification}
\phantomsection
\label{navigation:II.1.2.11.2}
\addcontentsline{toc}{subsubsection}{II.1.2.11.2 [CT2] Health Service Assurance Specification}

\label{element:ec5c4d1219464433b2cee5c5a988fef7}
\index{[CT2] Health Service Assurance Specification (II.1.2.11.2)}
{\small Capability · Canonical · current}\par

Specify how macro-level service-delivery assurance will be undertaken and what sufficiently supports confidence.



This capability defines the approach and evaluation criteria through which service-delivery assurance can be justified. Approach Definition addresses scope, method, responsibility, evidence and applicable independence; Criteria Governance establishes the basis for judging sufficiency. Their separation distinguishes the design of an assurance activity from the normative or evaluative material it applies.



\textbf{Intent}



Make the intended basis of confidence explicit before a conclusion is treated as assurance.



\textbf{Boundary}



An approach cannot waive a higher-level independence requirement. Contributing specialist concerns do not create separate Business Assurance hierarchies.


{\small Recorded unresolved matters: source-qualified matters; exact scope remains in the canonical record.}\par
{\footnotesize \href{../../../../continuum/strategy/capability/ec5c4d1219464433b2cee5c5a988fef7.md}{Canonical record: ec5c4d1219464433b2cee5c5a988fef7}; full qualifications and provenance.}\par
\paragraph*{II.1.2.11.2.1 [CT3] Assurance Approach Definition}
\phantomsection
\label{navigation:II.1.2.11.2.1}
\addcontentsline{toc}{paragraph}{II.1.2.11.2.1 [CT3] Assurance Approach Definition}

\label{element:d672f8831df54f2099927f820d0025f5}
\index{[CT3] Assurance Approach Definition (II.1.2.11.2.1)}
{\small Capability · Canonical · current}\par

Define scope, evidence expectations, responsibilities, methods, applicable degree of independence and timing/frequency for service-delivery assurance.



Assurance Approach Definition establishes how service-delivery assurance will be undertaken, including scope, evidence, methods, responsibilities, timing and required independence. It gives assurance activity a defensible approach rather than assuming any available observation can support confidence. The approach must respect applicable higher-level requirements when deciding what constitutes sufficient evaluation.



\textbf{Boundary}



Self-attestation cannot replace required independent governed assurance, and an approach does not itself establish a conclusion.


{\small Recorded unresolved matters: source-qualified matters; exact scope remains in the canonical record.}\par
{\footnotesize \href{../../../../continuum/strategy/capability/d672f8831df54f2099927f820d0025f5.md}{Canonical record: d672f8831df54f2099927f820d0025f5}; full qualifications and provenance.}\par
\paragraph*{II.1.2.11.2.2 [CT3] Assurance Evaluation Criteria Governance}
\phantomsection
\label{navigation:II.1.2.11.2.2}
\addcontentsline{toc}{paragraph}{II.1.2.11.2.2 [CT3] Assurance Evaluation Criteria Governance}

\label{element:f3e32f2a2c9d47a89eaf124647161522}
\index{[CT3] Assurance Evaluation Criteria Governance (II.1.2.11.2.2)}
{\small Capability · Canonical · current}\par

Establish and maintain criteria against which sufficient assurance can be determined.



This capability establishes and maintains the criteria against which assurance sufficiency is judged. It makes the evaluative basis of confidence explicit across the different concerns that can affect service delivery. Criteria can draw on clinical, workforce, technology, information or other contexts without turning those subjects into independent assurance capability trees.



\textbf{Boundary}



Criteria governance is distinct from applying criteria to evidence and does not remap the existing assurance BEC responsibilities.


{\small Recorded unresolved matters: source-qualified matters; exact scope remains in the canonical record.}\par
{\footnotesize \href{../../../../continuum/strategy/capability/f3e32f2a2c9d47a89eaf124647161522.md}{Canonical record: f3e32f2a2c9d47a89eaf124647161522}; full qualifications and provenance.}\par
\subsection*{II.1.2.12 [CT1] Health Service Direction \& Stewardship}
\phantomsection
\label{navigation:II.1.2.12}
\addcontentsline{toc}{subsection}{II.1.2.12 [CT1] Health Service Direction \& Stewardship}

{\small Historical predecessor: R.2.28 \hyperref[element:360bf50f92484448a01cc57f2dbce23c]{[CT1] Health Service Direction \& Stewardship}; Canonical ID: \texttt{\seqsplit{360bf50f92484448a01cc57f2dbce23c}}. Superseded historical identity — direct correspondence · superseded.}\par
\label{element:d51e7125bf9c4d2c9c59b60346b4d86f}
\index{[CT1] Health Service Direction \& Stewardship (II.1.2.12)}
{\small Capability · Canonical · current}\par

Model needed health services and their viability requirements, support establishment/improvement and align services with the wider ecosystem of dependencies, resources, financial arrangements and organisational change.



This capability considers what health services should exist, what makes them viable and how they fit within the wider service ecosystem. Modelling frames the need and required abilities; alignment addresses dependencies, resources, financial arrangements and surrounding change. The distinction keeps the rationale for a service separate from its identifiable offering and its current operational delivery state.



\textbf{Intent}



Support viable, coherent service development and improvement within organisational and ecosystem constraints.



\textbf{Boundary}



Identifying dependencies and resources does not acquire their operating authority. Service modelling and alignment do not establish the authoritative service catalogue.



\textbf{Harmonia Architectural Relevance}



Service need, viability and alignment facts materially shape service information and derivation requirements. That relevance does not assign executive authority or establish a future BEC mapping.


{\footnotesize \href{../../../../continuum/strategy/capability/d51e7125bf9c4d2c9c59b60346b4d86f.md}{Canonical record: d51e7125bf9c4d2c9c59b60346b4d86f}; full qualifications and provenance.}\par
\subsubsection*{II.1.2.12.1 [CT2] Health Service Alignment}
\phantomsection
\label{navigation:II.1.2.12.1}
\addcontentsline{toc}{subsubsection}{II.1.2.12.1 [CT2] Health Service Alignment}

\label{element:d4fc5367352741a99bd334f7e8bf5856}
\index{[CT2] Health Service Alignment (II.1.2.12.1)}
{\small Capability · Canonical · current}\par

Position the service coherently in its wider health-service ecosystem.



Health Service Alignment positions a service among the other services, resources and arrangements on which it depends. It considers reciprocal dependencies, required resources, financial fit and the organisational changes needed for introduction or material change. Where modelling asks what service is needed, alignment asks how it can fit coherently into its environment.



\textbf{Intent}



Expose the dependencies and surrounding commitments that affect service viability.



\textbf{Boundary}



Alignment does not establish the service offering, own the required resources or settle the unresolved ancestry of collective operational readiness.


{\small Recorded unresolved matters: source-qualified matters; exact scope remains in the canonical record.}\par
{\footnotesize \href{../../../../continuum/strategy/capability/d4fc5367352741a99bd334f7e8bf5856.md}{Canonical record: d4fc5367352741a99bd334f7e8bf5856}; full qualifications and provenance.}\par
\paragraph*{II.1.2.12.1.1 [CT3] Health Service Co-Dependency Management}
\phantomsection
\label{navigation:II.1.2.12.1.1}
\addcontentsline{toc}{paragraph}{II.1.2.12.1.1 [CT3] Health Service Co-Dependency Management}

\label{element:f031ee3ddc5e47c48d4308d4077b3967}
\index{[CT3] Health Service Co-Dependency Management (II.1.2.12.1.1)}
{\small Capability · Canonical · current}\par

Identify and manage what the service needs from others and what others need from it.



Health Service Co-Dependency Management identifies and manages reciprocal needs between a service and its wider ecosystem. It makes explicit both what the service requires from others and what others rely on it to provide. This supports alignment by exposing dependencies that a service description considered in isolation may conceal.



\textbf{Boundary}



Recognising a dependency does not acquire the provider's responsibility or establish a new technical Service dependency.


{\small Recorded unresolved matters: source-qualified matters; exact scope remains in the canonical record.}\par
{\footnotesize \href{../../../../continuum/strategy/capability/f031ee3ddc5e47c48d4308d4077b3967.md}{Canonical record: f031ee3ddc5e47c48d4308d4077b3967}; full qualifications and provenance.}\par
\paragraph*{II.1.2.12.1.2 [CT3] Health Service Financial Management}
\phantomsection
\label{navigation:II.1.2.12.1.2}
\addcontentsline{toc}{paragraph}{II.1.2.12.1.2 [CT3] Health Service Financial Management}

\label{element:2fbf31e974574f288f087179e4023dc2}
\index{[CT3] Health Service Financial Management (II.1.2.12.1.2)}
{\small Capability · Canonical · current}\par

Establish the service’s fit within funding, financial and charging arrangements without absorbing enterprise finance/accounting.



Health Service Financial Management establishes how the service fits applicable funding, financial and charging arrangements. It connects service viability with its surrounding financial context without attempting to become the enterprise finance function. The responsibility is service-specific alignment rather than the production of all financial transactions associated with care.



\textbf{Boundary}



Financial fit does not confer pricing, accounting, billing or enterprise revenue-management authority.


{\small Recorded unresolved matters: source-qualified matters; exact scope remains in the canonical record.}\par
{\footnotesize \href{../../../../continuum/strategy/capability/2fbf31e974574f288f087179e4023dc2.md}{Canonical record: 2fbf31e974574f288f087179e4023dc2}; full qualifications and provenance.}\par
\paragraph*{II.1.2.12.1.3 [CT3] Health Service Implementation Planning}
\phantomsection
\label{navigation:II.1.2.12.1.3}
\addcontentsline{toc}{paragraph}{II.1.2.12.1.3 [CT3] Health Service Implementation Planning}

\label{element:67bd795d91ce4dd49fdd481a3c451366}
\index{[CT3] Health Service Implementation Planning (II.1.2.12.1.3)}
{\small Capability · Canonical · current}\par

Establish and coordinate surrounding organisational/service changes needed to introduce or materially change the service.



Health Service Implementation Planning establishes and coordinates the organisational and surrounding service changes needed for introduction or material change. It considers how the service fits its environment, complementing modelling's establishment support with an alignment view of interdependent change. Its purpose is coherence of the planned transition rather than ownership of every change activity.



\textbf{Boundary}



Planning does not perform all supporting changes or define the authoritative service offering.


{\small Recorded unresolved matters: source-qualified matters; exact scope remains in the canonical record.}\par
{\footnotesize \href{../../../../continuum/strategy/capability/67bd795d91ce4dd49fdd481a3c451366.md}{Canonical record: 67bd795d91ce4dd49fdd481a3c451366}; full qualifications and provenance.}\par
\paragraph*{II.1.2.12.1.4 [CT3] Health Service Resource Identification}
\phantomsection
\label{navigation:II.1.2.12.1.4}
\addcontentsline{toc}{paragraph}{II.1.2.12.1.4 [CT3] Health Service Resource Identification}

\label{element:7d46575486c847cd8fdbb665b3ed1090}
\index{[CT3] Health Service Resource Identification (II.1.2.12.1.4)}
{\small Capability · Canonical · current}\par

Identify required workforce, facilities, equipment, clinical technology, ICT, administration and support resources without owning them.



This capability establishes the resources a service requires across workforce, facilities, equipment, clinical technology, ICT and support. It makes resource needs visible in the service's alignment context without treating identification as procurement, ownership or allocation. Required resources can therefore be discussed distinctly from resources presently available or deployed.



\textbf{Boundary}



Identifying resources does not operate them, assign them to care or establish collective service readiness.


{\small Recorded unresolved matters: source-qualified matters; exact scope remains in the canonical record.}\par
{\footnotesize \href{../../../../continuum/strategy/capability/7d46575486c847cd8fdbb665b3ed1090.md}{Canonical record: 7d46575486c847cd8fdbb665b3ed1090}; full qualifications and provenance.}\par
\subsubsection*{II.1.2.12.2 [CT2] Health Service Modelling}
\phantomsection
\label{navigation:II.1.2.12.2}
\addcontentsline{toc}{subsubsection}{II.1.2.12.2 [CT2] Health Service Modelling}

\label{element:9a3e90db97b94d1888a2a5c5effb8999}
\index{[CT2] Health Service Modelling (II.1.2.12.2)}
{\small Capability · Canonical · current}\par

Establish what service should exist and what it needs to become viable.



Health Service Modelling establishes the service needed and the ordinary business abilities required to make it viable. Its children distinguish recognising unmet need, identifying required abilities, supporting establishment and evaluating an existing service for improvement. The model can therefore revisit service assumptions without treating the original definition as permanently sufficient.



\textbf{Intent}



Connect a service's purpose with the abilities and preparation required for it to exist usefully.



\textbf{Boundary}



Capability Identification uses ordinary business meaning; it does not allocate Harmonia capability tiers. Implementation support does not take over recruitment, facilities, procurement or ICT responsibilities.


{\small Recorded unresolved matters: source-qualified matters; exact scope remains in the canonical record.}\par
{\footnotesize \href{../../../../continuum/strategy/capability/9a3e90db97b94d1888a2a5c5effb8999.md}{Canonical record: 9a3e90db97b94d1888a2a5c5effb8999}; full qualifications and provenance.}\par
\paragraph*{II.1.2.12.2.1 [CT3] Health Service Capability Identification}
\phantomsection
\label{navigation:II.1.2.12.2.1}
\addcontentsline{toc}{paragraph}{II.1.2.12.2.1 [CT3] Health Service Capability Identification}

\label{element:53f64a402905438aa87f9d16eff876da}
\index{[CT3] Health Service Capability Identification (II.1.2.12.2.1)}
{\small Capability · Canonical · current}\par

Identify the ordinary business abilities needed to satisfy the service need across clinical, workforce, facility, technology, information and support concerns.



This capability identifies the ordinary business abilities a proposed service needs across clinical, workforce, facility, information, technology and support concerns. It broadens service modelling beyond a clinical description to the conditions needed for viability. Its contribution is an understanding of required abilities, leaving their fulfilment and governing authority with the appropriate responsibilities.



\textbf{Boundary}



The word capability here does not allocate Harmonia CT tiers or establish new catalogue elements.


{\small Recorded unresolved matters: source-qualified matters; exact scope remains in the canonical record.}\par
{\footnotesize \href{../../../../continuum/strategy/capability/53f64a402905438aa87f9d16eff876da.md}{Canonical record: 53f64a402905438aa87f9d16eff876da}; full qualifications and provenance.}\par
\paragraph*{II.1.2.12.2.2 [CT3] Health Service Continuous Improvement}
\phantomsection
\label{navigation:II.1.2.12.2.2}
\addcontentsline{toc}{paragraph}{II.1.2.12.2.2 [CT3] Health Service Continuous Improvement}

\label{element:b2b2b91b1a094ef1b13ce6870ff9716f}
\index{[CT3] Health Service Continuous Improvement (II.1.2.12.2.2)}
{\small Capability · Canonical · current}\par

Evaluate the established service and identify correction/improvement needs that may feed back into modelling and implementation support.



This capability evaluates an established service and identifies correction or improvement needs across its service-wide purpose. It allows experience of actual operation to inform reconsideration of the service model and implementation support. The scope is broader than reporting a clinical-quality finding, while remaining anchored to the service being improved.



\textbf{Boundary}



It does not establish generic Continuous Improvement placement or transfer the authority of specialist activities affected by improvement.


{\small Recorded unresolved matters: source-qualified matters; exact scope remains in the canonical record.}\par
{\footnotesize \href{../../../../continuum/strategy/capability/b2b2b91b1a094ef1b13ce6870ff9716f.md}{Canonical record: b2b2b91b1a094ef1b13ce6870ff9716f}; full qualifications and provenance.}\par
\paragraph*{II.1.2.12.2.3 [CT3] Health Service Implementation Support}
\phantomsection
\label{navigation:II.1.2.12.2.3}
\addcontentsline{toc}{paragraph}{II.1.2.12.2.3 [CT3] Health Service Implementation Support}

\label{element:f74f5944573b408d94fe8c88e33e128a}
\index{[CT3] Health Service Implementation Support (II.1.2.12.2.3)}
{\small Capability · Canonical · current}\par

Support establishment of the framed service through applicable organisational preparation and programmes of work.



Health Service Implementation Support assists establishment of the framed service through organisational preparation and programmes of work. It connects the service model to practical readiness for introduction while recognising that recruitment, training, equipment and technology changes retain their own owners. Support therefore contributes to establishment without becoming an all-purpose implementation authority.



\textbf{Boundary}



It does not absorb HR, facilities, procurement or ICT responsibilities or establish operational-readiness ancestry.


{\small Recorded unresolved matters: source-qualified matters; exact scope remains in the canonical record.}\par
{\footnotesize \href{../../../../continuum/strategy/capability/f74f5944573b408d94fe8c88e33e128a.md}{Canonical record: f74f5944573b408d94fe8c88e33e128a}; full qualifications and provenance.}\par
\paragraph*{II.1.2.12.2.4 [CT3] Health Service Need Identification}
\phantomsection
\label{navigation:II.1.2.12.2.4}
\addcontentsline{toc}{paragraph}{II.1.2.12.2.4 [CT3] Health Service Need Identification}

\label{element:976c90e05d0e44b1b24c89bf61eac005}
\index{[CT3] Health Service Need Identification (II.1.2.12.2.4)}
{\small Capability · Canonical · current}\par

Establish an unmet health-service need requiring response.



Health Service Need Identification establishes the unmet service need for which an organisational response is being considered. It gives service modelling a reason grounded in what is missing rather than beginning with a predetermined solution. That identified need can then frame the abilities and resources required without assuming that a new offering has already been established.



\textbf{Boundary}



Identifying need does not define an offering, approve funding or determine the surrounding implementation arrangements.


{\small Recorded unresolved matters: source-qualified matters; exact scope remains in the canonical record.}\par
{\footnotesize \href{../../../../continuum/strategy/capability/976c90e05d0e44b1b24c89bf61eac005.md}{Canonical record: 976c90e05d0e44b1b24c89bf61eac005}; full qualifications and provenance.}\par
\subsection*{II.1.2.13 [CT1] Health Service Governance}
\phantomsection
\label{navigation:II.1.2.13}
\addcontentsline{toc}{subsection}{II.1.2.13 [CT1] Health Service Governance}

{\small Historical predecessor: R.2.29 \hyperref[element:24f3dd30223142efa35a5826552d53ab]{[CT1] Health Service Governance}; Canonical ID: \texttt{\seqsplit{24f3dd30223142efa35a5826552d53ab}}. Superseded historical identity — direct correspondence · superseded.}\par
\label{element:209aada781e844e28202a2535a264c21}
\index{[CT1] Health Service Governance (II.1.2.13)}
{\small Capability · Canonical · current}\par

Establish and maintain the obligations, authorities, constraints, policies and governance framework under which health-service activity must operate, and define how adherence will be assessed, checked and reported.



Health Service Governance establishes the normative envelope within which health-service activity operates: what obligations apply, who may decide, what constraints bind decisions and how adherence will be assessed and reported. Its framework and compliance branches distinguish the authoritative requirements from the arrangements used to judge satisfaction. Governance therefore shapes legitimate activity without becoming the manager of every governed operation.



\textbf{Intent}



Make obligations, decision authority and compliance expectations explicit and traceable.



\textbf{Boundary}



Operational management chooses and progresses means within the governing envelope; assurance evaluates evidence. Neither responsibility is absorbed merely because governance defines its constraints.



\textbf{Harmonia Architectural Relevance}



Formal relevance remains Unresolved. The requirements plainly constrain managed activity, but the RAD has not assigned enough platform participation or enabling derivation to adjudicate its relevance classification.


{\small Recorded unresolved matters: source-qualified matters; exact scope remains in the canonical record.}\par
{\footnotesize \href{../../../../continuum/strategy/capability/209aada781e844e28202a2535a264c21.md}{Canonical record: 209aada781e844e28202a2535a264c21}; full qualifications and provenance.}\par
\subsubsection*{II.1.2.13.1 [CT2] Compliance Governance}
\phantomsection
\label{navigation:II.1.2.13.1}
\addcontentsline{toc}{subsubsection}{II.1.2.13.1 [CT2] Compliance Governance}

\label{element:02a52157967f4b498807cac70e3410e9}
\index{[CT2] Compliance Governance (II.1.2.13.1)}
{\small Capability · Canonical · current}\par

Define how adherence to the authoritative WHAT will be determined, checked and reported.



Compliance Governance defines how adherence to established obligations, authorities and constraints will be determined. Assessment Definition specifies the evidence and basis for evaluation; Reporting \& Checkpoint Definition specifies where adherence is checked and what must be communicated. This makes the compliance test explicit without treating the design of a test as its successful execution.



\textbf{Intent}



Provide a defensible basis for checking and reporting whether the governing envelope has been respected.



\textbf{Boundary}



Defining assessment arrangements does not perform assurance or manage the governed activity. The obligations being assessed retain their authoritative source.


{\small Recorded unresolved matters: source-qualified matters; exact scope remains in the canonical record.}\par
{\footnotesize \href{../../../../continuum/strategy/capability/02a52157967f4b498807cac70e3410e9.md}{Canonical record: 02a52157967f4b498807cac70e3410e9}; full qualifications and provenance.}\par
\paragraph*{II.1.2.13.1.1 [CT3] Compliance Assessment Definition}
\phantomsection
\label{navigation:II.1.2.13.1.1}
\addcontentsline{toc}{paragraph}{II.1.2.13.1.1 [CT3] Compliance Assessment Definition}

\label{element:8a6db4bd0ed845dda125896c4e3ad3d0}
\index{[CT3] Compliance Assessment Definition (II.1.2.13.1.1)}
{\small Capability · Canonical · current}\par

Define the evidence, measures, criteria, controls and assessment arrangements sufficient to evaluate compliance with established obligations, authorities and constraints.



Compliance Assessment Definition specifies the evidence, measures, criteria and arrangements sufficient to evaluate adherence. It makes the assessment basis explicit after obligations, authorities and constraints have been established. This preserves the distinction between defining what compliance means and determining whether the available evidence demonstrates it in a particular case.



\textbf{Boundary}



Assessment definition does not replace the authoritative obligation, perform assurance or establish a compliance finding.


{\small Recorded unresolved matters: source-qualified matters; exact scope remains in the canonical record.}\par
{\footnotesize \href{../../../../continuum/strategy/capability/8a6db4bd0ed845dda125896c4e3ad3d0.md}{Canonical record: 8a6db4bd0ed845dda125896c4e3ad3d0}; full qualifications and provenance.}\par
\paragraph*{II.1.2.13.1.2 [CT3] Governance Reporting \& Checkpoint Definition}
\phantomsection
\label{navigation:II.1.2.13.1.2}
\addcontentsline{toc}{paragraph}{II.1.2.13.1.2 [CT3] Governance Reporting \& Checkpoint Definition}

\label{element:0ff9f332acdc435e98f4f527c74f60c9}
\index{[CT3] Governance Reporting \& Checkpoint Definition (II.1.2.13.1.2)}
{\small Capability · Canonical · current}\par

Establish where/how adherence is checked and what governance/compliance information must be reported.



This capability specifies where adherence must be checked and what governance or compliance information must be reported. It gives the governance framework observable expectations without treating reporting requirements as completed reporting. Within Compliance Governance it complements the assessment basis by making checking and communication obligations explicit.



\textbf{Boundary}



Defining a checkpoint does not perform the check or take over progression of the governed operation.


{\small Recorded unresolved matters: source-qualified matters; exact scope remains in the canonical record.}\par
{\footnotesize \href{../../../../continuum/strategy/capability/0ff9f332acdc435e98f4f527c74f60c9.md}{Canonical record: 0ff9f332acdc435e98f4f527c74f60c9}; full qualifications and provenance.}\par
\subsubsection*{II.1.2.13.2 [CT2] Governance Framework Management}
\phantomsection
\label{navigation:II.1.2.13.2}
\addcontentsline{toc}{subsubsection}{II.1.2.13.2 [CT2] Governance Framework Management}

\label{element:b39f53b3ffb441f388a89b45cc517deb}
\index{[CT2] Governance Framework Management (II.1.2.13.2)}
{\small Capability · Canonical · current}\par

Establish the authoritative WHAT and legitimate working parameters.



Governance Framework Management establishes what must be satisfied and the legitimate authority and bounds for organisational decisions. Its two children separate obligations attributable to source authority from the allocation and limitation of decision authority. Together they supply the authoritative reference against which operational arrangements and compliance assessments can be considered.



\textbf{Intent}



Keep required outcomes and legitimate decision parameters distinguishable and traceable to their authority.



\textbf{Boundary}



Interpreting an external obligation or engaging with a regulator does not transfer that body's authority. Operational means remain with management.


{\small Recorded unresolved matters: source-qualified matters; exact scope remains in the canonical record.}\par
{\footnotesize \href{../../../../continuum/strategy/capability/b39f53b3ffb441f388a89b45cc517deb.md}{Canonical record: b39f53b3ffb441f388a89b45cc517deb}; full qualifications and provenance.}\par
\paragraph*{II.1.2.13.2.1 [CT3] Authority \& Constraint Governance}
\phantomsection
\label{navigation:II.1.2.13.2.1}
\addcontentsline{toc}{paragraph}{II.1.2.13.2.1 [CT3] Authority \& Constraint Governance}

\label{element:ad7f41dc48a046b185f044f0a24517ca}
\index{[CT3] Authority \& Constraint Governance (II.1.2.13.2.1)}
{\small Capability · Canonical · current}\par

Establish who may make, approve or govern organisational decisions/artefacts and the bounds within which authority may be exercised.



This capability establishes who may make or approve organisational decisions and the bounds within which that authority applies. It complements obligation governance by distinguishing required outcomes from legitimate decision powers. A participant can therefore be responsible for acting within constraints without being authorised to redefine the governing requirement.



\textbf{Boundary}



Authority allocation does not itself perform an operational decision or a runtime access-control evaluation.


{\small Recorded unresolved matters: source-qualified matters; exact scope remains in the canonical record.}\par
{\footnotesize \href{../../../../continuum/strategy/capability/ad7f41dc48a046b185f044f0a24517ca.md}{Canonical record: ad7f41dc48a046b185f044f0a24517ca}; full qualifications and provenance.}\par
\paragraph*{II.1.2.13.2.2 [CT3] Policy / Obligation Governance}
\phantomsection
\label{navigation:II.1.2.13.2.2}
\addcontentsline{toc}{paragraph}{II.1.2.13.2.2 [CT3] Policy / Obligation Governance}

\label{element:73c6e5b36bdf4bf0ad2b8b50c9b9a00b}
\index{[CT3] Policy / Obligation Governance (II.1.2.13.2.2)}
{\small Capability · Canonical · current}\par

Identify, interpret, establish and maintain authoritative legislative, regulatory, policy, standards, guideline, contractual and other obligations, attributable and traceable to source authority.



Policy / Obligation Governance establishes and maintains the authoritative requirements under which the organisation operates. It interprets applicable legislative, regulatory, contractual, policy and other sources while preserving traceability to their authority. The purpose is a legitimate account of what must be satisfied, rather than an operational choice about how to satisfy it.



\textbf{Boundary}



Interpreting obligations does not acquire the regulator's authority or determine operational means.


{\small Recorded unresolved matters: source-qualified matters; exact scope remains in the canonical record.}\par
{\footnotesize \href{../../../../continuum/strategy/capability/73c6e5b36bdf4bf0ad2b8b50c9b9a00b.md}{Canonical record: 73c6e5b36bdf4bf0ad2b8b50c9b9a00b}; full qualifications and provenance.}\par
\subsection*{II.1.2.14 [CT1] Individual Care Delivery}
\phantomsection
\label{navigation:II.1.2.14}
\addcontentsline{toc}{subsection}{II.1.2.14 [CT1] Individual Care Delivery}

{\small Historical predecessor: R.2.48 \hyperref[element:162deb7e292042f78f677a14476a3a1d]{[CT1] Individual Care Delivery}; Canonical ID: \texttt{\seqsplit{162deb7e292042f78f677a14476a3a1d}}. Superseded historical identity — direct correspondence · superseded.}\par
\label{element:b964a11cde594318aa14ff935414158a}
\index{[CT1] Individual Care Delivery (II.1.2.14)}
{\small Capability · Canonical · current}\par

Assess, clinically determine, plan, intervene, monitor and review individual care independently of setting, mode, time pattern or organisational arrangement.



Individual Care Delivery is the organisation's ability to understand an individual's health and act on that understanding through purposeful care. Its six constituent abilities distinguish establishing findings, interpreting their clinical meaning, choosing a response, performing care, observing change and evaluating care against its objectives. These distinctions describe responsibilities and outcomes rather than a compulsory sequence: a consultation or continuing episode may revisit several of them.



The same ability can be exercised in a hospital, a general practice, a home or a virtual consultation. Setting and time pattern affect the circumstances of care without changing the responsibility represented by this capability.



\textbf{Intent}



Support clinically meaningful care whose intent, performance, observed effects and continuing appropriateness can be distinguished.



\textbf{Boundary}



Care access, appointments and operational coordination enable care but do not establish that it has been performed. Clinical judgement and the authority to deliver care remain with the responsible participants; possession or coordination of their information does not transfer that authority to Harmonia.



\textbf{Harmonia Architectural Relevance}



Harmonia supplies and connects information used in individual care. The approved Health Information Access support trace is bounded: records can inform care, but direct assessment can establish findings independently, and access to a proposed intervention does not authorise its performance.


{\footnotesize \href{../../../../continuum/strategy/capability/b964a11cde594318aa14ff935414158a.md}{Canonical record: b964a11cde594318aa14ff935414158a}; full qualifications and provenance.}\par
\subsubsection*{II.1.2.14.1 [CT2] Clinical Determination}
\phantomsection
\label{navigation:II.1.2.14.1}
\addcontentsline{toc}{subsubsection}{II.1.2.14.1 [CT2] Clinical Determination}

{\small Historical predecessor: R.2.33 \hyperref[element:018ccd818ec142fba1ffb086d900acc6]{[CT2] Clinical Determination}; Canonical ID: \texttt{\seqsplit{018ccd818ec142fba1ffb086d900acc6}}. Superseded historical identity — direct correspondence · superseded.}\par
\label{element:39219464b6304b59bd7c3de4e237066e}
\index{[CT2] Clinical Determination (II.1.2.14.1)}
{\small Capability · Canonical · current}\par

The capability to interpret assessed health information and establish clinically meaningful conclusions about an individual's health state, condition, risk or prognosis.



Clinical Determination turns assessed findings into clinically interpreted understanding of a person's condition, risk or prognosis. It represents the professional responsibility for deciding what the findings mean in the individual's clinical context, including the significance of changes identified through monitoring. This is the interpretive responsibility within Individual Care Delivery, rather than another method of collecting evidence.



\textbf{Intent}



Establish clinical conclusions on which an intended response can be based.



\textbf{Boundary}



Determining clinical meaning does not itself select the care response. No CT3 taxonomy has been adjudicated; diagnosis and prognostic judgement illustrate this responsibility without becoming additional catalogue elements.


{\small Recorded unresolved matters: source-qualified matters; exact scope remains in the canonical record.}\par
{\footnotesize \href{../../../../continuum/strategy/capability/39219464b6304b59bd7c3de4e237066e.md}{Canonical record: 39219464b6304b59bd7c3de4e237066e}; full qualifications and provenance.}\par
\paragraph*{II.1.2.14.1.1 [CT3] Clinical Finding Interpretation}
\phantomsection
\label{navigation:II.1.2.14.1.1}
\addcontentsline{toc}{paragraph}{II.1.2.14.1.1 [CT3] Clinical Finding Interpretation}

{\small Historical predecessor: R.2.30 \hyperref[element:cf8ddfca7a8e4ed893ec39b795200d77]{[CT3] Clinical Finding Interpretation}; Canonical ID: \texttt{\seqsplit{cf8ddfca7a8e4ed893ec39b795200d77}}. Superseded historical identity — direct correspondence · superseded.}\par
\label{element:deb78ba3b1514074a8ca5bc5da0673f5}
\index{[CT3] Clinical Finding Interpretation (II.1.2.14.1.1)}
{\small Capability · Proposed · current}\par

Interpret observed clinical findings and their significance for the individual care concern.



\textbf{Overview}



Interpret observed clinical findings and their significance for the individual care concern. The proposed decomposition separates this outcome from the other responsibilities within Clinical Determination.



\textbf{Intent}



Make the clinical significance of observations intelligible before diagnostic or response decisions are taken.



\textbf{Boundary}



Does not itself establish a diagnosis, prescribe intervention or confer this clinical responsibility on Harmonia.


{\footnotesize \href{../../../../continuum/strategy/capability/deb78ba3b1514074a8ca5bc5da0673f5.md}{Canonical record: deb78ba3b1514074a8ca5bc5da0673f5}; full qualifications and provenance.}\par
\paragraph*{II.1.2.14.1.2 [CT3] Diagnosis \& Problem Determination}
\phantomsection
\label{navigation:II.1.2.14.1.2}
\addcontentsline{toc}{paragraph}{II.1.2.14.1.2 [CT3] Diagnosis \& Problem Determination}

{\small Historical predecessor: R.2.31 \hyperref[element:e575c583783546ef85cf898523d18797]{[CT3] Diagnosis \& Problem Determination}; Canonical ID: \texttt{\seqsplit{e575c583783546ef85cf898523d18797}}. Superseded historical identity — direct correspondence · superseded.}\par
\label{element:d965032637e64275b7a61c0413834590}
\index{[CT3] Diagnosis \& Problem Determination (II.1.2.14.1.2)}
{\small Capability · Proposed · current}\par

Establish and qualify the clinical problems or diagnoses applicable to an individual.



\textbf{Overview}



Establish and qualify the clinical problems or diagnoses applicable to an individual. The proposed decomposition separates this outcome from the other responsibilities within Clinical Determination.



\textbf{Intent}



Establish the clinically accountable problem statement that subsequent care decisions address.



\textbf{Boundary}



Clinical determination remains with the mandated professional; a coded record is not evidence that Harmonia made the diagnosis.


{\footnotesize \href{../../../../continuum/strategy/capability/d965032637e64275b7a61c0413834590.md}{Canonical record: d965032637e64275b7a61c0413834590}; full qualifications and provenance.}\par
\paragraph*{II.1.2.14.1.3 [CT3] Health Risk \& Prognostic Determination}
\phantomsection
\label{navigation:II.1.2.14.1.3}
\addcontentsline{toc}{paragraph}{II.1.2.14.1.3 [CT3] Health Risk \& Prognostic Determination}

{\small Historical predecessor: R.2.32 \hyperref[element:3215f0b0e1ce45e8ab16dae272ea6fe7]{[CT3] Health Risk \& Prognostic Determination}; Canonical ID: \texttt{\seqsplit{3215f0b0e1ce45e8ab16dae272ea6fe7}}. Superseded historical identity — direct correspondence · superseded.}\par
\label{element:5bba0ea969814b5fa7ad11b67ab40a0c}
\index{[CT3] Health Risk \& Prognostic Determination (II.1.2.14.1.3)}
{\small Capability · Proposed · current}\par

Establish clinically meaningful risk and expected-course judgements for an individual.



\textbf{Overview}



Establish clinically meaningful risk and expected-course judgements for an individual. The proposed decomposition separates this outcome from the other responsibilities within Clinical Determination.



\textbf{Intent}



Make anticipated clinical risk and course explicit for the individual care concern.



\textbf{Boundary}



Risk/prognostic judgement is distinct from operational urgency, source-information authority and platform assurance.


{\footnotesize \href{../../../../continuum/strategy/capability/5bba0ea969814b5fa7ad11b67ab40a0c.md}{Canonical record: 5bba0ea969814b5fa7ad11b67ab40a0c}; full qualifications and provenance.}\par
\subsubsection*{II.1.2.14.2 [CT2] Clinical Intervention}
\phantomsection
\label{navigation:II.1.2.14.2}
\addcontentsline{toc}{subsubsection}{II.1.2.14.2 [CT2] Clinical Intervention}

{\small Historical predecessor: R.2.34 \hyperref[element:958345ee65f446d3b1259b8895ab100f]{[CT2] Clinical Intervention}; Canonical ID: \texttt{\seqsplit{958345ee65f446d3b1259b8895ab100f}}. Superseded historical identity — direct correspondence · superseded.}\par
\label{element:5359d1b0ff814a0e871384469f5ebcd7}
\index{[CT2] Clinical Intervention (II.1.2.14.2)}
{\small Capability · Canonical · current}\par

The capability to perform clinical actions intended to alter, maintain or otherwise manage an individual's health state in accordance with an established care intent.



Clinical Intervention is the ability to perform a clinical action under an established care intent. It marks the transition from intended care to care actually performed, whether the action seeks to alter, maintain or otherwise manage health. Its exercise can establish that treatment occurred even when the desired clinical result was not achieved.



\textbf{Intent}



Put the intended clinical response into practice through appropriately authorised care.



\textbf{Boundary}



Orders and appointments express or arrange future activity; they are not evidence that the intervention occurred. Therapy, medication and procedural examples do not constitute an approved three-child taxonomy: that proposed decomposition remains provisional.


{\small Recorded unresolved matters: source-qualified matters; exact scope remains in the canonical record.}\par
{\footnotesize \href{../../../../continuum/strategy/capability/5359d1b0ff814a0e871384469f5ebcd7.md}{Canonical record: 5359d1b0ff814a0e871384469f5ebcd7}; full qualifications and provenance.}\par
\subsubsection*{II.1.2.14.3 [CT2] Individual Care Planning}
\phantomsection
\label{navigation:II.1.2.14.3}
\addcontentsline{toc}{subsubsection}{II.1.2.14.3 [CT2] Individual Care Planning}

{\small Historical predecessor: R.2.37 \hyperref[element:574e61e230e643f78c58649adb4db313]{[CT2] Individual Care Planning}; Canonical ID: \texttt{\seqsplit{574e61e230e643f78c58649adb4db313}}. Superseded historical identity — direct correspondence · superseded.}\par
\label{element:f036a9866459429a8fb5f5e10727b5e8}
\index{[CT2] Individual Care Planning (II.1.2.14.3)}
{\small Capability · Canonical · current}\par

The capability to establish the intended clinical response to an individual's assessed and clinically interpreted health state, including the objectives, interventions and monitoring required to progress their care.



Individual Care Planning establishes the intended course of care in response to assessed and clinically interpreted health information. It connects clinical understanding to objectives, intended interventions and monitoring needs. Its outcome is an expression of what should happen, allowing intended care to remain distinguishable from requests, operational arrangements and actions subsequently performed.



\textbf{Intent}



Make the clinical response and its objectives explicit enough to guide care and later evaluate its effectiveness.



\textbf{Boundary}



Prescribing can express clinical intent without establishing dispensing or administration. Scheduling or allocating resources does not perform the planned intervention. No CT3 decomposition is established.


{\small Recorded unresolved matters: source-qualified matters; exact scope remains in the canonical record.}\par
{\footnotesize \href{../../../../continuum/strategy/capability/f036a9866459429a8fb5f5e10727b5e8.md}{Canonical record: f036a9866459429a8fb5f5e10727b5e8}; full qualifications and provenance.}\par
\paragraph*{II.1.2.14.3.1 [CT3] Care Goal Definition}
\phantomsection
\label{navigation:II.1.2.14.3.1}
\addcontentsline{toc}{paragraph}{II.1.2.14.3.1 [CT3] Care Goal Definition}

{\small Historical predecessor: R.2.35 \hyperref[element:8bfbcc0a455d423c90ae2f1cfb732ca8]{[CT3] Care Goal Definition}; Canonical ID: \texttt{\seqsplit{8bfbcc0a455d423c90ae2f1cfb732ca8}}. Superseded historical identity — direct correspondence · superseded.}\par
\label{element:609e1bef516b44bb976b5dbebe7f684b}
\index{[CT3] Care Goal Definition (II.1.2.14.3.1)}
{\small Capability · Proposed · current}\par

Establish the intended outcomes and priorities against which an individual care response is planned.



\textbf{Overview}



Establish the intended outcomes and priorities against which an individual care response is planned. The proposed decomposition separates this outcome from the other responsibilities within Individual Care Planning.



\textbf{Intent}



Establish the intended clinical benefit and priorities against which care responses can be considered.



\textbf{Boundary}



Goal definition remains distinct from selecting an intervention and from recording an externally agreed care plan.


{\footnotesize \href{../../../../continuum/strategy/capability/609e1bef516b44bb976b5dbebe7f684b.md}{Canonical record: 609e1bef516b44bb976b5dbebe7f684b}; full qualifications and provenance.}\par
\paragraph*{II.1.2.14.3.2 [CT3] Care Response Planning}
\phantomsection
\label{navigation:II.1.2.14.3.2}
\addcontentsline{toc}{paragraph}{II.1.2.14.3.2 [CT3] Care Response Planning}

{\small Historical predecessor: R.2.36 \hyperref[element:d64569acaf2c44f3b2a5332f74fff7c7]{[CT3] Care Response Planning}; Canonical ID: \texttt{\seqsplit{d64569acaf2c44f3b2a5332f74fff7c7}}. Superseded historical identity — direct correspondence · superseded.}\par
\label{element:2636e43903e44a97ad5b8db32692ce8b}
\index{[CT3] Care Response Planning (II.1.2.14.3.2)}
{\small Capability · Proposed · current}\par

Determine the intended clinical response and its sequencing in relation to agreed care goals.



\textbf{Overview}



Determine the intended clinical response and its sequencing in relation to agreed care goals. The proposed decomposition separates this outcome from the other responsibilities within Individual Care Planning.



\textbf{Intent}



Make the intended clinical response coherent in relation to the agreed care goals.



\textbf{Boundary}



Planning is distinct from intervention delivery and from the operational definition of a synthetic activity.


{\footnotesize \href{../../../../continuum/strategy/capability/2636e43903e44a97ad5b8db32692ce8b.md}{Canonical record: 2636e43903e44a97ad5b8db32692ce8b}; full qualifications and provenance.}\par
\subsubsection*{II.1.2.14.4 [CT2] Individual Care Review}
\phantomsection
\label{navigation:II.1.2.14.4}
\addcontentsline{toc}{subsubsection}{II.1.2.14.4 [CT2] Individual Care Review}

{\small Historical predecessor: R.2.40 \hyperref[element:e1048f6c635d4e16b5384cc2cb45a71b]{[CT2] Individual Care Review}; Canonical ID: \texttt{\seqsplit{e1048f6c635d4e16b5384cc2cb45a71b}}. Superseded historical identity — direct correspondence · superseded.}\par
\label{element:625a93f132f242f5a8e2ae1dc19979fb}
\index{[CT2] Individual Care Review (II.1.2.14.4)}
{\small Capability · Canonical · current}\par

The capability to evaluate the progress and effectiveness of an individual's care against its intended objectives and establish the continuing appropriateness of the current care intent.



Individual Care Review asks whether care is achieving its intended objectives and whether the existing care intent remains suitable. Its children separate effectiveness from appropriateness: an intervention may have an effect yet no longer be suitable to continue, or remain a reasonable choice despite an inadequate response. Review therefore evaluates the relationship between care, outcomes and intent rather than simply repeating assessment.



\textbf{Intent}



Establish whether continuing the current approach is justified and whether reconsideration is needed.



\textbf{Boundary}



Recognising a need to change care does not establish the replacement plan. Monitoring supplies relevant observations; planning establishes revised intent.


{\small Recorded unresolved matters: source-qualified matters; exact scope remains in the canonical record.}\par
{\footnotesize \href{../../../../continuum/strategy/capability/625a93f132f242f5a8e2ae1dc19979fb.md}{Canonical record: 625a93f132f242f5a8e2ae1dc19979fb}; full qualifications and provenance.}\par
\paragraph*{II.1.2.14.4.1 [CT3] Care Appropriateness Review}
\phantomsection
\label{navigation:II.1.2.14.4.1}
\addcontentsline{toc}{paragraph}{II.1.2.14.4.1 [CT3] Care Appropriateness Review}

{\small Historical predecessor: R.2.39 \hyperref[element:905eb2de1ac341ac96331e3d7f76079b]{[CT3] Care Intent Appropriateness Review}; Canonical ID: \texttt{\seqsplit{905eb2de1ac341ac96331e3d7f76079b}}. Superseded historical identity — evolved correspondence · superseded.}\par
\label{element:5eac5fb73e31416288d8c992805b33cb}
\index{[CT3] Care Appropriateness Review (II.1.2.14.4.1)}
{\small Capability · Canonical · current}\par

Evaluate whether the current intended care remains suitable in the individual’s current context.



Care Appropriateness Review evaluates whether the current care intent remains suitable for the individual in their present context. Its question concerns continuing suitability rather than whether treatment has produced a measurable effect. That distinction permits a review to recognise that effective care may no longer be appropriate to continue.



\textbf{Boundary}



An appropriateness judgement can identify a need to reconsider care; establishing revised intent remains Individual Care Planning.


{\small Recorded unresolved matters: source-qualified matters; exact scope remains in the canonical record.}\par
{\footnotesize \href{../../../../continuum/strategy/capability/5eac5fb73e31416288d8c992805b33cb.md}{Canonical record: 5eac5fb73e31416288d8c992805b33cb}; full qualifications and provenance.}\par
\paragraph*{II.1.2.14.4.2 [CT3] Care Effectiveness Review}
\phantomsection
\label{navigation:II.1.2.14.4.2}
\addcontentsline{toc}{paragraph}{II.1.2.14.4.2 [CT3] Care Effectiveness Review}

{\small Historical predecessor: R.2.38 \hyperref[element:31f49cdea9f44f188459503580c2b86e]{[CT3] Care Effectiveness Review}; Canonical ID: \texttt{\seqsplit{31f49cdea9f44f188459503580c2b86e}}. Superseded historical identity — direct correspondence · superseded.}\par
\label{element:79cfeaefeb174a54ba43343d37f035c1}
\index{[CT3] Care Effectiveness Review (II.1.2.14.4.2)}
{\small Capability · Canonical · current}\par

Evaluate whether delivered care is achieving or has achieved its intended goals and outcomes.



Care Effectiveness Review evaluates delivered care against its intended goals and outcomes. It connects observed response to the reason the care was undertaken, making a judgement that monitoring alone does not supply. An effectiveness finding can inform review of appropriateness while preserving the possibility that a reasonable care choice has not achieved its objective.



\textbf{Boundary}



Observed improvement is not automatically proof that the intended objectives have been achieved or that continuing care is appropriate.


{\small Recorded unresolved matters: source-qualified matters; exact scope remains in the canonical record.}\par
{\footnotesize \href{../../../../continuum/strategy/capability/79cfeaefeb174a54ba43343d37f035c1.md}{Canonical record: 79cfeaefeb174a54ba43343d37f035c1}; full qualifications and provenance.}\par
\subsubsection*{II.1.2.14.5 [CT2] Individual Health Assessment}
\phantomsection
\label{navigation:II.1.2.14.5}
\addcontentsline{toc}{subsubsection}{II.1.2.14.5 [CT2] Individual Health Assessment}

{\small Historical predecessor: R.2.44 \hyperref[element:9111446aa2d045bd9acd98542b0eef1f]{[CT2] Individual Health Assessment}; Canonical ID: \texttt{\seqsplit{9111446aa2d045bd9acd98542b0eef1f}}. Superseded historical identity — direct correspondence · superseded.}\par
\label{element:045f819bb3694d3c9db7f20886e2993f}
\index{[CT2] Individual Health Assessment (II.1.2.14.5)}
{\small Capability · Canonical · current}\par

The capability to establish an understanding of an individual's current health state through the collection, consideration and evaluation of relevant clinical, health and contextual information.



Assessment establishes what is currently known about an individual's health through history, examination and measurement. Within Individual Care Delivery it supplies assessed findings, including findings obtained directly from the person rather than retrieved from an existing record. Its three children distinguish ways of establishing knowledge, rather than professions, clinical specialties or locations.



\textbf{Intent}



Give subsequent clinical interpretation a defensible account of the person's current state and relevant history.



\textbf{Boundary}



Finding an abnormal measurement is distinguishable from deciding what it means clinically. Diagnosis, prognosis, selection of care and performance of care belong to adjacent abilities.


{\small Recorded unresolved matters: source-qualified matters; exact scope remains in the canonical record.}\par
{\footnotesize \href{../../../../continuum/strategy/capability/045f819bb3694d3c9db7f20886e2993f.md}{Canonical record: 045f819bb3694d3c9db7f20886e2993f}; full qualifications and provenance.}\par
\paragraph*{II.1.2.14.5.1 [CT3] Clinical Examination}
\phantomsection
\label{navigation:II.1.2.14.5.1}
\addcontentsline{toc}{paragraph}{II.1.2.14.5.1 [CT3] Clinical Examination}

{\small Historical predecessor: R.2.41 \hyperref[element:8bc15e2b8f114deebcbbe3f2013a8147]{[CT3] Clinical Examination}; Canonical ID: \texttt{\seqsplit{8bc15e2b8f114deebcbbe3f2013a8147}}. Superseded historical identity — direct correspondence · superseded.}\par
\label{element:8855d4d9e2954e0ab66ab65a556fc61a}
\index{[CT3] Clinical Examination (II.1.2.14.5.1)}
{\small Capability · Canonical · current}\par

Establish health findings through clinical examination.



Clinical Examination establishes findings through examination of the individual. It distinguishes directly acquired clinical knowledge from the historical account and the measurement-focused contribution of its assessment siblings. Its outcome is a set of examination findings that can inform clinical interpretation, irrespective of the setting or profession exercising the ability.



\textbf{Boundary}



An examination finding is not itself a diagnosis, care plan or performed treatment.


{\small Recorded unresolved matters: source-qualified matters; exact scope remains in the canonical record.}\par
{\footnotesize \href{../../../../continuum/strategy/capability/8855d4d9e2954e0ab66ab65a556fc61a.md}{Canonical record: 8855d4d9e2954e0ab66ab65a556fc61a}; full qualifications and provenance.}\par
\paragraph*{II.1.2.14.5.2 [CT3] Health History Assessment}
\phantomsection
\label{navigation:II.1.2.14.5.2}
\addcontentsline{toc}{paragraph}{II.1.2.14.5.2 [CT3] Health History Assessment}

{\small Historical predecessor: R.2.42 \hyperref[element:bcaacf446f45443eafd8e21a1fd74ee1]{[CT3] Health History Assessment}; Canonical ID: \texttt{\seqsplit{bcaacf446f45443eafd8e21a1fd74ee1}}. Superseded historical identity — direct correspondence · superseded.}\par
\label{element:7ea7e3f74e2940a7934b588a1a038834}
\index{[CT3] Health History Assessment (II.1.2.14.5.2)}
{\small Capability · Canonical · current}\par

Establish relevant health knowledge from the individual’s history and applicable historical/contextual information.



Health History Assessment establishes clinically relevant understanding from a person's history and its surrounding context. Within assessment it contributes knowledge that a current examination or measurement alone may not reveal, while requiring consideration of the relevance of historical information to the present situation. Access to previous records can assist, but the responsibility is the assessment of history rather than retrieval itself.



\textbf{Boundary}



Historical information does not become a current clinical conclusion merely because it is available; interpretation of the assessed findings remains Clinical Determination.


{\small Recorded unresolved matters: source-qualified matters; exact scope remains in the canonical record.}\par
{\footnotesize \href{../../../../continuum/strategy/capability/7ea7e3f74e2940a7934b588a1a038834.md}{Canonical record: 7ea7e3f74e2940a7934b588a1a038834}; full qualifications and provenance.}\par
\paragraph*{II.1.2.14.5.3 [CT3] Health Status Measurement}
\phantomsection
\label{navigation:II.1.2.14.5.3}
\addcontentsline{toc}{paragraph}{II.1.2.14.5.3 [CT3] Health Status Measurement}

{\small Historical predecessor: R.2.43 \hyperref[element:3717906f7f8840d3a8a3e3877543dd18]{[CT3] Health Status Measurement}; Canonical ID: \texttt{\seqsplit{3717906f7f8840d3a8a3e3877543dd18}}. Superseded historical identity — direct correspondence · superseded.}\par
\label{element:a796bfb3335d44f6aeb9970628ed8554}
\index{[CT3] Health Status Measurement (II.1.2.14.5.3)}
{\small Capability · Canonical · current}\par

Establish individual health observations or measurements.



Health Status Measurement establishes observations or measurements of an individual's health. Its contribution to assessment is evidence about the measured state, with a purpose distinct from developing a view over time or deciding the clinical meaning of a result. Measurements can inform both present assessment and later monitoring without collapsing those abilities.



\textbf{Boundary}



Obtaining a value does not establish a temporal trend or the clinical significance of that value.


{\small Recorded unresolved matters: source-qualified matters; exact scope remains in the canonical record.}\par
{\footnotesize \href{../../../../continuum/strategy/capability/a796bfb3335d44f6aeb9970628ed8554.md}{Canonical record: a796bfb3335d44f6aeb9970628ed8554}; full qualifications and provenance.}\par
\subsubsection*{II.1.2.14.6 [CT2] Individual Health Monitoring}
\phantomsection
\label{navigation:II.1.2.14.6}
\addcontentsline{toc}{subsubsection}{II.1.2.14.6 [CT2] Individual Health Monitoring}

{\small Historical predecessor: R.2.47 \hyperref[element:31e7c27556314cafb74d040441d2dac6]{[CT2] Individual Health Monitoring}; Canonical ID: \texttt{\seqsplit{31e7c27556314cafb74d040441d2dac6}}. Superseded historical identity — direct correspondence · superseded.}\par
\label{element:c7d0e498e1ad44138df940af2f91cd2d}
\index{[CT2] Individual Health Monitoring (II.1.2.14.6)}
{\small Capability · Canonical · current}\par

The capability to observe and evaluate an individual's health state over time in order to identify clinically meaningful change, stability or deviation from an expected state or trajectory.



Monitoring develops a view of an individual's health over time so that change, stability or deviation can be recognised. Health State Monitoring establishes the aggregate temporally situated view; Care Response Monitoring considers that view in relation to a particular condition and treatment. Together they distinguish observation over time from a single measurement and from the clinical judgement prompted by a trend.



\textbf{Intent}



Make relevant trajectories visible while preserving what was observed and what still requires interpretation.



\textbf{Boundary}



Monitoring does not establish the clinical significance of a change or judge whether care achieved its objectives. Those responsibilities remain Clinical Determination and Individual Care Review respectively.


{\small Recorded unresolved matters: source-qualified matters; exact scope remains in the canonical record.}\par
{\footnotesize \href{../../../../continuum/strategy/capability/c7d0e498e1ad44138df940af2f91cd2d.md}{Canonical record: c7d0e498e1ad44138df940af2f91cd2d}; full qualifications and provenance.}\par
\paragraph*{II.1.2.14.6.1 [CT3] Care Response Monitoring}
\phantomsection
\label{navigation:II.1.2.14.6.1}
\addcontentsline{toc}{paragraph}{II.1.2.14.6.1 [CT3] Care Response Monitoring}

{\small Historical predecessor: R.2.45 \hyperref[element:615f93d05fbb4a8481f07778bd4f4803]{[CT3] Care Response Monitoring}; Canonical ID: \texttt{\seqsplit{615f93d05fbb4a8481f07778bd4f4803}}. Superseded historical identity — direct correspondence · superseded.}\par
\label{element:d6951c9357fa40a4943cde581d95d5b1}
\index{[CT3] Care Response Monitoring (II.1.2.14.6.1)}
{\small Capability · Canonical · current}\par

Scope that health-state view to a particular problem/condition and the treatment applied to it.



Care Response Monitoring considers the monitored health-state view in relation to a particular problem or condition and the treatment applied. It makes a response trajectory visible in its treatment context, helping distinguish a general change in health from a change relevant to the care being followed. It consumes the aggregate view established through Health State Monitoring.



\textbf{Boundary}



Identifying the response does not judge treatment effectiveness against objectives or determine a replacement treatment.


{\small Recorded unresolved matters: source-qualified matters; exact scope remains in the canonical record.}\par
{\footnotesize \href{../../../../continuum/strategy/capability/d6951c9357fa40a4943cde581d95d5b1.md}{Canonical record: d6951c9357fa40a4943cde581d95d5b1}; full qualifications and provenance.}\par
\paragraph*{II.1.2.14.6.2 [CT3] Health State Monitoring}
\phantomsection
\label{navigation:II.1.2.14.6.2}
\addcontentsline{toc}{paragraph}{II.1.2.14.6.2 [CT3] Health State Monitoring}

{\small Historical predecessor: R.2.46 \hyperref[element:346f024f74b447a8b6056b9d521837ff]{[CT3] Health State Monitoring}; Canonical ID: \texttt{\seqsplit{346f024f74b447a8b6056b9d521837ff}}. Superseded historical identity — direct correspondence · superseded.}\par
\label{element:d77f22909f8b47dba4cbd95982e576df}
\index{[CT3] Health State Monitoring (II.1.2.14.6.2)}
{\small Capability · Canonical · current}\par

Establish an aggregate, temporally situated view of the individual’s health from applicable monitored data.



Health State Monitoring brings applicable monitored data into an aggregate view situated in time. Its purpose is to make the individual's changing or stable health state intelligible across observations, rather than leaving each observation isolated. Within monitoring it provides the broader view that can also be considered in relation to a particular treatment or condition.



\textbf{Boundary}



Aggregation is not clinical interpretation. Care Response Monitoring uses the established view rather than independently recreating it.


{\small Recorded unresolved matters: source-qualified matters; exact scope remains in the canonical record.}\par
{\footnotesize \href{../../../../continuum/strategy/capability/d77f22909f8b47dba4cbd95982e576df.md}{Canonical record: d77f22909f8b47dba4cbd95982e576df}; full qualifications and provenance.}\par
\subsection*{II.1.2.15 [CT1] Population Health \& Health-System Planning}
\phantomsection
\label{navigation:II.1.2.15}
\addcontentsline{toc}{subsection}{II.1.2.15 [CT1] Population Health \& Health-System Planning}

{\small Historical predecessor: R.2.49 \hyperref[element:a5003436a7f348b1bff60ba593adfa27]{[CT1] Population Health \& Health-System Planning}; Canonical ID: \texttt{\seqsplit{a5003436a7f348b1bff60ba593adfa27}}. Superseded historical identity — direct correspondence · superseded.}\par
\label{element:c63b38b37e9b4e9b97dc87ed54604def}
\index{[CT1] Population Health \& Health-System Planning (II.1.2.15)}
{\small Capability · Canonical · current}\par

Understand population health through assessment/surveillance and applicable event/notification management; assess health-service demand and plan required capacity and its distribution.



This capability connects an understanding of population health with consideration of future health-service needs. Surveillance addresses population state, significant events and notification; planning assesses service demand and required capacity and distribution. Population evidence can inform planning without predetermining its conclusions or becoming the operational allocation of resources.



\textbf{Intent}



Support informed public-health awareness and a reasoned view of the services a population will require.



\textbf{Boundary}



Population conclusions differ from individual clinical findings, and planned capacity differs from capacity currently usable by a service.



\textbf{Harmonia Architectural Relevance}



Governed population information and formal notifications materially support health activity. Consuming planning context does not make the platform the health-system planner or owner of its resources.


{\footnotesize \href{../../../../continuum/strategy/capability/c63b38b37e9b4e9b97dc87ed54604def.md}{Canonical record: c63b38b37e9b4e9b97dc87ed54604def}; full qualifications and provenance.}\par
\subsubsection*{II.1.2.15.1 [CT2] Health-System Planning}
\phantomsection
\label{navigation:II.1.2.15.1}
\addcontentsline{toc}{subsubsection}{II.1.2.15.1 [CT2] Health-System Planning}

\label{element:9ca4df07eceb4150a0190cc488270ce2}
\index{[CT2] Health-System Planning (II.1.2.15.1)}
{\small Capability · Canonical · current}\par

Assess health-service demand and plan the capacity and footprint required to meet it.



Health-System Planning considers the capacity and service footprint required to meet assessed demand. Demand Assessment establishes the basis and scale of need, Capacity Planning determines what is required, and Distribution Planning addresses where that capacity should sit across the service network. The decomposition preserves the difference between understanding demand and deciding a response to it.



\textbf{Intent}



Support a defensible relationship between expected demand and the planned scale and distribution of health services.



\textbf{Boundary}



Plans do not establish available resources or operational allocations. Detailed commissioning, capital, funding and workforce planning remain deliberately unmodelled.


{\small Recorded unresolved matters: source-qualified matters; exact scope remains in the canonical record.}\par
{\footnotesize \href{../../../../continuum/strategy/capability/9ca4df07eceb4150a0190cc488270ce2.md}{Canonical record: 9ca4df07eceb4150a0190cc488270ce2}; full qualifications and provenance.}\par
\paragraph*{II.1.2.15.1.1 [CT3] Health Service Capacity Planning}
\phantomsection
\label{navigation:II.1.2.15.1.1}
\addcontentsline{toc}{paragraph}{II.1.2.15.1.1 [CT3] Health Service Capacity Planning}

\label{element:c1d1db0e2c854077a41721615e86119d}
\index{[CT3] Health Service Capacity Planning (II.1.2.15.1.1)}
{\small Capability · Canonical · current}\par

Determine required strategic/tactical capacity to address assessed demand.



Health Service Capacity Planning determines the strategic or tactical capacity required to address assessed demand. It considers what should be available, preserving a distinction between requirement and the resources currently in service. Within Health-System Planning it gives demand a capacity response without performing operational allocation.



\textbf{Boundary}



Planned capacity is neither established resource availability nor the service's currently usable operational capacity.


{\small Recorded unresolved matters: source-qualified matters; exact scope remains in the canonical record.}\par
{\footnotesize \href{../../../../continuum/strategy/capability/c1d1db0e2c854077a41721615e86119d.md}{Canonical record: c1d1db0e2c854077a41721615e86119d}; full qualifications and provenance.}\par
\paragraph*{II.1.2.15.1.2 [CT3] Health Service Demand Assessment}
\phantomsection
\label{navigation:II.1.2.15.1.2}
\addcontentsline{toc}{paragraph}{II.1.2.15.1.2 [CT3] Health Service Demand Assessment}

\label{element:e56d354eafd94ff7a495fc2ca2b8a593}
\index{[CT3] Health Service Demand Assessment (II.1.2.15.1.2)}
{\small Capability · Canonical · current}\par

Assess current/expected health-service demand across a population, geography, service domain or health-system context.



Health Service Demand Assessment establishes a view of current or expected demand within a population, geography or service context. It translates relevant evidence into a demand judgement rather than assuming population-health state directly determines service requirements. Its assumptions and scope matter because capacity planning depends on what demand has actually been assessed.



\textbf{Boundary}



System-level demand is distinct from a particular person's care requirement and does not itself determine planned capacity.


{\small Recorded unresolved matters: source-qualified matters; exact scope remains in the canonical record.}\par
{\footnotesize \href{../../../../continuum/strategy/capability/e56d354eafd94ff7a495fc2ca2b8a593.md}{Canonical record: e56d354eafd94ff7a495fc2ca2b8a593}; full qualifications and provenance.}\par
\paragraph*{II.1.2.15.1.3 [CT3] Health Service Distribution Planning}
\phantomsection
\label{navigation:II.1.2.15.1.3}
\addcontentsline{toc}{paragraph}{II.1.2.15.1.3 [CT3] Health Service Distribution Planning}

\label{element:4fc525fccbe142a58ee43dcc58b11c92}
\index{[CT3] Health Service Distribution Planning (II.1.2.15.1.3)}
{\small Capability · Canonical · current}\par

Determine the distribution of planned capacity across a geographical, organisational, site or service-network footprint.



Health Service Distribution Planning determines how planned capacity should be distributed across a geographical, organisational or service-network footprint. It addresses the placement of the planned response rather than only its aggregate scale. This allows the planning context to distinguish enough total capacity from capacity positioned where it is needed.



\textbf{Boundary}



A planned footprint does not establish recognised location records, identifiable offerings or actual resource deployment.


{\small Recorded unresolved matters: source-qualified matters; exact scope remains in the canonical record.}\par
{\footnotesize \href{../../../../continuum/strategy/capability/4fc525fccbe142a58ee43dcc58b11c92.md}{Canonical record: 4fc525fccbe142a58ee43dcc58b11c92}; full qualifications and provenance.}\par
\subsubsection*{II.1.2.15.2 [CT2] Population Health Surveillance}
\phantomsection
\label{navigation:II.1.2.15.2}
\addcontentsline{toc}{subsubsection}{II.1.2.15.2 [CT2] Population Health Surveillance}

\label{element:83e568fa8af14d7791a7253765ea390f}
\index{[CT2] Population Health Surveillance (II.1.2.15.2)}
{\small Capability · Canonical · current}\par

Understand population-health state and manage organised surveillance, significant events and applicable notification.



Population Health Surveillance establishes patterns and change across a defined population and manages the surveillance arrangements through which they are considered. Its children separate longitudinal assessment, surveillance management, event identification, formal notification and reporting. These distinctions matter because an observed pattern can be significant without being notifiable, and identifying a notifiable event does not demonstrate that the notification obligation has been discharged.



\textbf{Intent}



Make population-health concerns and applicable communication obligations visible in their population and temporal context.



\textbf{Boundary}



Notification management does not define the obligation or own the operational response. Surveillance evidence can inform other abilities without establishing their findings.


{\small Recorded unresolved matters: source-qualified matters; exact scope remains in the canonical record.}\par
{\footnotesize \href{../../../../continuum/strategy/capability/83e568fa8af14d7791a7253765ea390f.md}{Canonical record: 83e568fa8af14d7791a7253765ea390f}; full qualifications and provenance.}\par
\paragraph*{II.1.2.15.2.1 [CT3] Notifiable Health Event Management}
\phantomsection
\label{navigation:II.1.2.15.2.1}
\addcontentsline{toc}{paragraph}{II.1.2.15.2.1 [CT3] Notifiable Health Event Management}

\label{element:e605436254ed4dad819f452448e054a1}
\index{[CT3] Notifiable Health Event Management (II.1.2.15.2.1)}
{\small Capability · Canonical · current}\par

Manage formal notification applicability, required information, recipient, status, submission, correction/supplementation and applicable follow-up.



Notifiable Health Event Management addresses the formal communication responsibilities applicable to a health event. It manages applicability, required information, recipient and submission status, including correction or supplementation where relevant. Its purpose is to make the notification obligation and its discharge distinguishable from both event identification and ordinary reporting.



\textbf{Boundary}



Notification management does not define the governing obligation or own the operational response to the event.


{\small Recorded unresolved matters: source-qualified matters; exact scope remains in the canonical record.}\par
{\footnotesize \href{../../../../continuum/strategy/capability/e605436254ed4dad819f452448e054a1.md}{Canonical record: e605436254ed4dad819f452448e054a1}; full qualifications and provenance.}\par
\paragraph*{II.1.2.15.2.2 [CT3] Population Health Event Identification}
\phantomsection
\label{navigation:II.1.2.15.2.2}
\addcontentsline{toc}{paragraph}{II.1.2.15.2.2 [CT3] Population Health Event Identification}

\label{element:e084ba370bb545d0a4c28c37e2e8e850}
\index{[CT3] Population Health Event Identification (II.1.2.15.2.2)}
{\small Capability · Canonical · current}\par

Identify and characterise significant population-health events/patterns from applicable evidence.



Population Health Event Identification recognises and characterises a significant event or pattern from relevant evidence. It establishes the population-level concern rather than merely noting that observations changed. The resulting event context can inform notification or response decisions while retaining a separate responsibility for those decisions.



\textbf{Boundary}



An identified event is not automatically notifiable, and notifiability does not demonstrate that an obligation was discharged.


{\small Recorded unresolved matters: source-qualified matters; exact scope remains in the canonical record.}\par
{\footnotesize \href{../../../../continuum/strategy/capability/e084ba370bb545d0a4c28c37e2e8e850.md}{Canonical record: e084ba370bb545d0a4c28c37e2e8e850}; full qualifications and provenance.}\par
\paragraph*{II.1.2.15.2.3 [CT3] Population Health Longitudinal Assessment}
\phantomsection
\label{navigation:II.1.2.15.2.3}
\addcontentsline{toc}{paragraph}{II.1.2.15.2.3 [CT3] Population Health Longitudinal Assessment}

\label{element:4d52b8acda504329a595be27a736e63c}
\index{[CT3] Population Health Longitudinal Assessment (II.1.2.15.2.3)}
{\small Capability · Canonical · current}\par

Establish population-health state, patterns and change over time from applicable evidence.



This capability establishes population-health state, patterns and change over time from relevant evidence. It interprets information at a population scope rather than treating a collection of individual observations as a population conclusion by itself. Its longitudinal purpose makes the temporal and population basis important to the meaning of the assessment.



\textbf{Boundary}



Population assessment is not preventative-programme performance, an individual health conclusion or a service-demand decision.


{\small Recorded unresolved matters: source-qualified matters; exact scope remains in the canonical record.}\par
{\footnotesize \href{../../../../continuum/strategy/capability/4d52b8acda504329a595be27a736e63c.md}{Canonical record: 4d52b8acda504329a595be27a736e63c}; full qualifications and provenance.}\par
\paragraph*{II.1.2.15.2.4 [CT3] Population Health Reporting}
\phantomsection
\label{navigation:II.1.2.15.2.4}
\addcontentsline{toc}{paragraph}{II.1.2.15.2.4 [CT3] Population Health Reporting}

\label{element:50b8b0884f8444c4b11d9e0f7a7754a3}
\index{[CT3] Population Health Reporting (II.1.2.15.2.4)}
{\small Capability · Canonical · current}\par

Communicate established population-health assessments, surveillance, events, trends and findings.



Population Health Reporting communicates established assessments, surveillance findings, events and trends. Its role is to make population understanding available with its scope intact, rather than replace the assessment or event-identification work underlying it. This supports use of the findings by other abilities without assuming what those abilities must conclude.



\textbf{Boundary}



General reporting is distinct from formal notification management and does not establish a response or planning decision.


{\small Recorded unresolved matters: source-qualified matters; exact scope remains in the canonical record.}\par
{\footnotesize \href{../../../../continuum/strategy/capability/50b8b0884f8444c4b11d9e0f7a7754a3.md}{Canonical record: 50b8b0884f8444c4b11d9e0f7a7754a3}; full qualifications and provenance.}\par
\paragraph*{II.1.2.15.2.5 [CT3] Population Health Surveillance Management}
\phantomsection
\label{navigation:II.1.2.15.2.5}
\addcontentsline{toc}{paragraph}{II.1.2.15.2.5 [CT3] Population Health Surveillance Management}

\label{element:f0791768ef494e619006ee0d89d0d9d8}
\index{[CT3] Population Health Surveillance Management (II.1.2.15.2.5)}
{\small Capability · Canonical · current}\par

Establish and manage surveillance purposes, populations, concerns and regimes.



Population Health Surveillance Management establishes the purposes, populations and regimes under which surveillance operates. It concerns the organised basis for watching health concerns, rather than any particular event detected through that activity. Keeping the regime distinct from its findings makes it possible to understand why evidence is being collected and considered.



\textbf{Boundary}



Managing surveillance neither identifies every significant event nor determines the operational response to one.


{\small Recorded unresolved matters: source-qualified matters; exact scope remains in the canonical record.}\par
{\footnotesize \href{../../../../continuum/strategy/capability/f0791768ef494e619006ee0d89d0d9d8.md}{Canonical record: f0791768ef494e619006ee0d89d0d9d8}; full qualifications and provenance.}\par
\subsection*{II.1.2.16 [CT1] Security, Privacy \& Digital Trust}
\phantomsection
\label{navigation:II.1.2.16}
\addcontentsline{toc}{subsection}{II.1.2.16 [CT1] Security, Privacy \& Digital Trust}

{\small Historical predecessor: R.2.50 \hyperref[element:b01441b2f5584fc2b9c3988643a7c93b]{[CT1] Security, Privacy \& Digital Trust}; Canonical ID: \texttt{\seqsplit{b01441b2f5584fc2b9c3988643a7c93b}}. Superseded historical identity — direct correspondence · superseded.}\par
\label{element:71f0431fdfcd4f779e16d9e5ea609833}
\index{[CT1] Security, Privacy \& Digital Trust (II.1.2.16)}
{\small Capability · Canonical · current}\par

Establish adequate identity confidence, coherent access controls, privacy/protection safeguards, governed activity attribution and accountability for security-relevant activity and information use.



This capability establishes confidence in identity, coherent control of access, protection of information and accountability for its use. Its three branches distinguish identity/access governance, protection and privacy, and attribution/accountability. The scope includes physical and procedural arrangements as well as digital controls, preventing a software authorisation check from standing in for the organisation's whole trust responsibility.



\textbf{Intent}



Support information use and governed activity within legitimate identity, authority, privacy and protection constraints.



\textbf{Boundary}



Identifying an actor, permitting an action, recording participation and judging appropriateness are different outcomes. Evidence production is not itself independent assurance.



\textbf{Harmonia Architectural Relevance}



Security and governance are intrinsic to Harmonia-managed operations. The Core classification concerns assigned platform controls and evidence; it does not acquire every organisational security or privacy responsibility.


{\footnotesize \href{../../../../continuum/strategy/capability/71f0431fdfcd4f779e16d9e5ea609833.md}{Canonical record: 71f0431fdfcd4f779e16d9e5ea609833}; full qualifications and provenance.}\par
\subsubsection*{II.1.2.16.1 [CT2] Digital Accountability \& Trust}
\phantomsection
\label{navigation:II.1.2.16.1}
\addcontentsline{toc}{subsubsection}{II.1.2.16.1 [CT2] Digital Accountability \& Trust}

\label{element:9a600233ae1f4f3289c09a0b08ce029d}
\index{[CT2] Digital Accountability \& Trust (II.1.2.16.1)}
{\small Capability · Canonical · current}\par

Establish governed attribution and accountability for activity and information use.



Digital Accountability \& Trust makes activity attributable and evaluates whether security-relevant activity and information use were appropriate. Attribution records participation together with acting capacity and relevant permission or constraint context; accountability uses that context to consider conduct against policy and obligations. A record of who acted is therefore necessary context, rather than proof that the action was permitted.



\textbf{Intent}



Support meaningful responsibility for information use and security-relevant conduct.



\textbf{Boundary}



Authentication, attribution and authorisation must remain distinguishable. Expected downstream evidentiary use does not establish new BC-to-BEC relationships.


{\small Recorded unresolved matters: source-qualified matters; exact scope remains in the canonical record.}\par
{\footnotesize \href{../../../../continuum/strategy/capability/9a600233ae1f4f3289c09a0b08ce029d.md}{Canonical record: 9a600233ae1f4f3289c09a0b08ce029d}; full qualifications and provenance.}\par
\paragraph*{II.1.2.16.1.1 [CT3] Governed Activity Attribution}
\phantomsection
\label{navigation:II.1.2.16.1.1}
\addcontentsline{toc}{paragraph}{II.1.2.16.1.1 [CT3] Governed Activity Attribution}

\label{element:646b3fd80a8b4f269691db7942e0f809}
\index{[CT3] Governed Activity Attribution (II.1.2.16.1.1)}
{\small Capability · Canonical · current}\par

Preserve who participated, what was done, the attributable person/role/unit/organisation capacity and applicable permission/authority/constraint context.



Governed Activity Attribution preserves who participated and in what acting capacity, together with the relevant permission, authority and constraint context. It makes conduct explainable beyond a bare actor identifier and timestamp. Within Digital Accountability \& Trust, this supplies the contextual record needed to distinguish participation from entitlement.



\textbf{Boundary}



Attributing an action to a person or role does not establish that the action was permitted or appropriate.


{\small Recorded unresolved matters: source-qualified matters; exact scope remains in the canonical record.}\par
{\footnotesize \href{../../../../continuum/strategy/capability/646b3fd80a8b4f269691db7942e0f809.md}{Canonical record: 646b3fd80a8b4f269691db7942e0f809}; full qualifications and provenance.}\par
\paragraph*{II.1.2.16.1.2 [CT3] Security \& Information-Use Accountability}
\phantomsection
\label{navigation:II.1.2.16.1.2}
\addcontentsline{toc}{paragraph}{II.1.2.16.1.2 [CT3] Security \& Information-Use Accountability}

\label{element:a33d0cc14c384bd48012def4ded54a67}
\index{[CT3] Security \& Information-Use Accountability (II.1.2.16.1.2)}
{\small Capability · Canonical · current}\par

Evaluate whether attributed security-relevant activity/information use was appropriate within policy, obligation and authority constraints, and establish accountability.



This capability evaluates attributed activity and information use against applicable policy, obligations and authority constraints. It establishes accountability rather than simply recording participation. Its purpose depends on preserving the circumstances of action so that appropriateness is considered within the relevant mandate instead of inferred from identity or technical success.



\textbf{Boundary}



An authenticated actor and a successful operation are not proof of legitimate information use.


{\small Recorded unresolved matters: source-qualified matters; exact scope remains in the canonical record.}\par
{\footnotesize \href{../../../../continuum/strategy/capability/a33d0cc14c384bd48012def4ded54a67.md}{Canonical record: a33d0cc14c384bd48012def4ded54a67}; full qualifications and provenance.}\par
\subsubsection*{II.1.2.16.2 [CT2] Digital Identity \& Access Governance}
\phantomsection
\label{navigation:II.1.2.16.2}
\addcontentsline{toc}{subsubsection}{II.1.2.16.2 [CT2] Digital Identity \& Access Governance}

\label{element:b9cb3666f9134ba29eb896e67eba7250}
\index{[CT2] Digital Identity \& Access Governance (II.1.2.16.2)}
{\small Capability · Canonical · current}\par

Govern sufficient identity confidence and coherent access controls across the organisation.



This capability governs who may act and how controlled concerns are identified and protected across the organisation. Its children separate confidence in a claimed identity, identity models, contextual identity quality, authentication, access-control context, assurance, collaboration and enforcement. That separation prevents successful authentication from being treated as sufficient authority for every use.



\textbf{Intent}



Make identity confidence and control requirements appropriate to the circumstances in which access or activity occurs.



\textbf{Boundary}



Federation or Single Sign-On does not guarantee identity quality. Enforcement spans physical, procedural and digital controls; it is broader than a runtime policy engine.


{\small Recorded unresolved matters: source-qualified matters; exact scope remains in the canonical record.}\par
{\footnotesize \href{../../../../continuum/strategy/capability/b9cb3666f9134ba29eb896e67eba7250.md}{Canonical record: b9cb3666f9134ba29eb896e67eba7250}; full qualifications and provenance.}\par
\paragraph*{II.1.2.16.2.1 [CT3] Access Control Assurance \& Reporting}
\phantomsection
\label{navigation:II.1.2.16.2.1}
\addcontentsline{toc}{paragraph}{II.1.2.16.2.1 [CT3] Access Control Assurance \& Reporting}

\label{element:ba0bf23c2b5b4de1bdfbc5df90f57a2f}
\index{[CT3] Access Control Assurance \& Reporting (II.1.2.16.2.1)}
{\small Capability · Canonical · current}\par

Establish organisational assurance/evidence of whether access-control arrangements are appropriately established and operating.



This capability establishes organisational evidence and assurance about whether access-control arrangements are appropriately established and operating. It makes control effectiveness a distinct concern from defining requirements or performing an individual enforcement action. Reporting communicates that established position without making an isolated successful check proof of end-to-end control.



\textbf{Boundary}



Subject-specific access-control assurance does not create another macro-level Health Service Assurance hierarchy.


{\small Recorded unresolved matters: source-qualified matters; exact scope remains in the canonical record.}\par
{\footnotesize \href{../../../../continuum/strategy/capability/ba0bf23c2b5b4de1bdfbc5df90f57a2f.md}{Canonical record: ba0bf23c2b5b4de1bdfbc5df90f57a2f}; full qualifications and provenance.}\par
\paragraph*{II.1.2.16.2.2 [CT3] Access Control Collaboration}
\phantomsection
\label{navigation:II.1.2.16.2.2}
\addcontentsline{toc}{paragraph}{II.1.2.16.2.2 [CT3] Access Control Collaboration}

\label{element:c1a39b83d27a4a148f997c0a709e82ac}
\index{[CT3] Access Control Collaboration (II.1.2.16.2.2)}
{\small Capability · Canonical · current}\par

Coordinate arrangements across independently controlled domains/parties while retaining their control ownership.



Access Control Collaboration coordinates arrangements across parties or domains that retain independent control ownership. Its purpose is coherent cross-boundary control where a strong local arrangement alone cannot establish the whole result. Collaboration therefore concerns agreed control participation rather than silently centralising every participant's authority.



\textbf{Boundary}



Coordination does not transfer the parties' control ownership or guarantee identity quality across domains.


{\small Recorded unresolved matters: source-qualified matters; exact scope remains in the canonical record.}\par
{\footnotesize \href{../../../../continuum/strategy/capability/c1a39b83d27a4a148f997c0a709e82ac.md}{Canonical record: c1a39b83d27a4a148f997c0a709e82ac}; full qualifications and provenance.}\par
\paragraph*{II.1.2.16.2.3 [CT3] Access Control Context Management}
\phantomsection
\label{navigation:II.1.2.16.2.3}
\addcontentsline{toc}{paragraph}{II.1.2.16.2.3 [CT3] Access Control Context Management}

\label{element:a51f7a8045d24d0f982320025ba3a6ae}
\index{[CT3] Access Control Context Management (II.1.2.16.2.3)}
{\small Capability · Canonical · current}\par

Determine what information, functions, activities or other concerns require control, under what circumstances and how requirements are modelled.



Access Control Context Management establishes what requires control, the circumstances in which control applies and how those requirements are modelled. It gives access governance an explicit subject and context rather than reducing it to assembling runtime attributes. The responsibility can encompass information, functions and activities under physical, procedural or digital arrangements.



\textbf{Boundary}



Defining control context does not itself enforce the controls or establish that they operate effectively.


{\small Recorded unresolved matters: source-qualified matters; exact scope remains in the canonical record.}\par
{\footnotesize \href{../../../../continuum/strategy/capability/a51f7a8045d24d0f982320025ba3a6ae.md}{Canonical record: a51f7a8045d24d0f982320025ba3a6ae}; full qualifications and provenance.}\par
\paragraph*{II.1.2.16.2.4 [CT3] Access Control Enforcement}
\phantomsection
\label{navigation:II.1.2.16.2.4}
\addcontentsline{toc}{paragraph}{II.1.2.16.2.4 [CT3] Access Control Enforcement}

\label{element:b72a7b3fb7384870a9790b04d3c9e94a}
\index{[CT3] Access Control Enforcement (II.1.2.16.2.4)}
{\small Capability · Canonical · current}\par

Ensure physical, procedural and digital control mechanisms coherently impose intended controls.



Access Control Enforcement ensures that intended controls are coherently imposed through the applicable physical, procedural and digital mechanisms. It represents the organisational ability to make controls effective, rather than equating enforcement with one software decision. Its outcome depends on the controlled circumstances and the arrangements intended to apply there.



\textbf{Boundary}



A local software authorisation result alone does not establish coherent end-to-end organisational enforcement.


{\small Recorded unresolved matters: source-qualified matters; exact scope remains in the canonical record.}\par
{\footnotesize \href{../../../../continuum/strategy/capability/b72a7b3fb7384870a9790b04d3c9e94a.md}{Canonical record: b72a7b3fb7384870a9790b04d3c9e94a}; full qualifications and provenance.}\par
\paragraph*{II.1.2.16.2.5 [CT3] Identity Assurance}
\phantomsection
\label{navigation:II.1.2.16.2.5}
\addcontentsline{toc}{paragraph}{II.1.2.16.2.5 [CT3] Identity Assurance}

\label{element:56ccd6d3bf5b4ec980a333dfadf8d9fa}
\index{[CT3] Identity Assurance (II.1.2.16.2.5)}
{\small Capability · Canonical · current}\par

Establish sufficient confidence that a person is who they claim to be through appropriate evidence and processes.



Identity Assurance establishes confidence that a person is who they claim to be using appropriate evidence and processes. It concerns the basis and strength of the identity claim, rather than only successful use of an authentication mechanism. Within identity governance it makes confidence a question that can be considered separately from contextual sufficiency and access authority.



\textbf{Boundary}



An assured identity is not automatically sufficient for every context or permitted to perform every action.


{\small Recorded unresolved matters: source-qualified matters; exact scope remains in the canonical record.}\par
{\footnotesize \href{../../../../continuum/strategy/capability/56ccd6d3bf5b4ec980a333dfadf8d9fa.md}{Canonical record: 56ccd6d3bf5b4ec980a333dfadf8d9fa}; full qualifications and provenance.}\par
\paragraph*{II.1.2.16.2.6 [CT3] Identity Identification \& Authentication}
\phantomsection
\label{navigation:II.1.2.16.2.6}
\addcontentsline{toc}{paragraph}{II.1.2.16.2.6 [CT3] Identity Identification \& Authentication}

\label{element:ba4ad62548984d63ad75da3561566431}
\index{[CT3] Identity Identification \& Authentication (II.1.2.16.2.6)}
{\small Capability · Canonical · current}\par

Identify an actor and authenticate the asserted identity.



This capability identifies the acting participant and authenticates the asserted identity. It establishes who is acting under the applicable identity arrangements so that later control and accountability can be grounded in an identified actor. The outcome remains distinct from confidence established through identity evidence and from authority to access a controlled concern.



\textbf{Boundary}



Authentication neither establishes every required identity-quality condition nor authorises the requested activity.


{\small Recorded unresolved matters: source-qualified matters; exact scope remains in the canonical record.}\par
{\footnotesize \href{../../../../continuum/strategy/capability/ba4ad62548984d63ad75da3561566431.md}{Canonical record: ba4ad62548984d63ad75da3561566431}; full qualifications and provenance.}\par
\paragraph*{II.1.2.16.2.7 [CT3] Identity Modelling \& Interoperability}
\phantomsection
\label{navigation:II.1.2.16.2.7}
\addcontentsline{toc}{paragraph}{II.1.2.16.2.7 [CT3] Identity Modelling \& Interoperability}

\label{element:c40a00b0fe56405d99b6bb6ede72e18e}
\index{[CT3] Identity Modelling \& Interoperability (II.1.2.16.2.7)}
{\small Capability · Canonical · current}\par

Establish identity models and interoperability arrangements appropriate to organisational identity use.



This capability establishes identity models and arrangements through which identities can be understood across organisational use. It addresses the coherence needed when parties exchange or rely on identity information without assuming that their identity assertions have identical quality. Its purpose is meaningful interoperability, rather than a particular federation or Single Sign-On technology.



\textbf{Boundary}



Interoperability does not guarantee assurance quality or transfer each party's control ownership.


{\small Recorded unresolved matters: source-qualified matters; exact scope remains in the canonical record.}\par
{\footnotesize \href{../../../../continuum/strategy/capability/c40a00b0fe56405d99b6bb6ede72e18e.md}{Canonical record: c40a00b0fe56405d99b6bb6ede72e18e}; full qualifications and provenance.}\par
\paragraph*{II.1.2.16.2.8 [CT3] Identity Qualification Control}
\phantomsection
\label{navigation:II.1.2.16.2.8}
\addcontentsline{toc}{paragraph}{II.1.2.16.2.8 [CT3] Identity Qualification Control}

\label{element:3d53064f2fb7475b8f49fb58ec8ad4a5}
\index{[CT3] Identity Qualification Control (II.1.2.16.2.8)}
{\small Capability · Canonical · current}\par

Determine whether presented identity characteristics, provenance and assurance quality are sufficient for its intended context.



Identity Qualification Control evaluates whether presented identity characteristics, provenance and assurance are sufficient for the intended context. It adds a contextual judgement beyond whether the actor can be authenticated. This permits an identity adequate for one purpose to remain insufficient for another without claiming that the identity is necessarily false.



\textbf{Boundary}



Identity qualification does not establish practitioner clinical qualification or a general permission to act.


{\small Recorded unresolved matters: source-qualified matters; exact scope remains in the canonical record.}\par
{\footnotesize \href{../../../../continuum/strategy/capability/3d53064f2fb7475b8f49fb58ec8ad4a5.md}{Canonical record: 3d53064f2fb7475b8f49fb58ec8ad4a5}; full qualifications and provenance.}\par
\subsubsection*{II.1.2.16.3 [CT2] Health Information Protection}
\phantomsection
\label{navigation:II.1.2.16.3}
\addcontentsline{toc}{subsubsection}{II.1.2.16.3 [CT2] Health Information Protection}

\label{element:4ae0de8fa16840d09ab2ef3f5545a5e7}
\index{[CT2] Health Information Protection (II.1.2.16.3)}
{\small Capability · Canonical · current}\par

Govern privacy consequences and adequate protection of the information estate and its environment.



Health Information Protection governs the consequences of information exposure and the safeguards protecting information and its operating environment. Privacy safeguards, estate protection, access-practice monitoring, compartment assurance and classification assurance address different questions. Together they distinguish appropriate arrangements from evidence of how information is actually handled, including physical and printed handling.



\textbf{Intent}



Protect information and the people affected by its use, with an intelligible basis for evaluating protective arrangements.



\textbf{Boundary}



Compartment assurance does not decide access. Subject-specific protection assurance does not create additional macro-level Health Service Assurance trees.


{\small Recorded unresolved matters: source-qualified matters; exact scope remains in the canonical record.}\par
{\footnotesize \href{../../../../continuum/strategy/capability/4ae0de8fa16840d09ab2ef3f5545a5e7.md}{Canonical record: 4ae0de8fa16840d09ab2ef3f5545a5e7}; full qualifications and provenance.}\par
\paragraph*{II.1.2.16.3.1 [CT3] Health Information Access Monitoring}
\phantomsection
\label{navigation:II.1.2.16.3.1}
\addcontentsline{toc}{paragraph}{II.1.2.16.3.1 [CT3] Health Information Access Monitoring}

\label{element:01854c7311404bbf9f698fce4656b54d}
\index{[CT3] Health Information Access Monitoring (II.1.2.16.3.1)}
{\small Capability · Canonical · current}\par

Evaluate actual organisational access/handling practice, including controls, accessibility, compartments, classification and physical/printed handling.



Health Information Access Monitoring evaluates actual access and handling practice across organisational circumstances. Digital audit information can contribute, but so can physical access, printed handling, compartments and classification practice. Its purpose is to understand what is happening in use rather than merely observe software events or assume configured controls describe reality.



\textbf{Boundary}



Evidence supply does not transfer business accountability, and observation is not automatically independent assurance.


{\small Recorded unresolved matters: source-qualified matters; exact scope remains in the canonical record.}\par
{\footnotesize \href{../../../../continuum/strategy/capability/01854c7311404bbf9f698fce4656b54d.md}{Canonical record: 01854c7311404bbf9f698fce4656b54d}; full qualifications and provenance.}\par
\paragraph*{II.1.2.16.3.2 [CT3] Health Information Classification Assurance}
\phantomsection
\label{navigation:II.1.2.16.3.2}
\addcontentsline{toc}{paragraph}{II.1.2.16.3.2 [CT3] Health Information Classification Assurance}

\label{element:4a4a871081bf45829eab2a93c9d1c21f}
\index{[CT3] Health Information Classification Assurance (II.1.2.16.3.2)}
{\small Capability · Canonical · current}\par

Assure appropriate classification according to sensitivity and potential consequences of exposure/mishandling.



This capability evaluates whether classification is appropriate to information sensitivity and the consequences of exposure or mishandling. It examines the application of classification rather than merely the presence of a label. The responsibility supports protection decisions while retaining the distinction between the classification basis and the control applied to a particular request.



\textbf{Boundary}



Classification assurance does not define compartment membership or grant permission to access classified information.


{\small Recorded unresolved matters: source-qualified matters; exact scope remains in the canonical record.}\par
{\footnotesize \href{../../../../continuum/strategy/capability/4a4a871081bf45829eab2a93c9d1c21f.md}{Canonical record: 4a4a871081bf45829eab2a93c9d1c21f}; full qualifications and provenance.}\par
\paragraph*{II.1.2.16.3.3 [CT3] Health Information Compartment Assurance}
\phantomsection
\label{navigation:II.1.2.16.3.3}
\addcontentsline{toc}{paragraph}{II.1.2.16.3.3 [CT3] Health Information Compartment Assurance}

\label{element:be47d201de1a45a88f48234d69deb2eb}
\index{[CT3] Health Information Compartment Assurance (II.1.2.16.3.3)}
{\small Capability · Canonical · current}\par

Assure appropriate establishment and application of compartmentalisation without establishing access decisions.



This capability evaluates whether contextual information compartments are appropriately established and applied. It treats the correctness of the grouping as a concern distinct from permission to disclose a particular item. Such assurance can make compartment use more trustworthy without changing the business relationships that define membership.



\textbf{Boundary}



Assuring compartmentalisation does not establish access decisions or sensitivity classification.


{\small Recorded unresolved matters: source-qualified matters; exact scope remains in the canonical record.}\par
{\footnotesize \href{../../../../continuum/strategy/capability/be47d201de1a45a88f48234d69deb2eb.md}{Canonical record: be47d201de1a45a88f48234d69deb2eb}; full qualifications and provenance.}\par
\paragraph*{II.1.2.16.3.4 [CT3] Health Information Safeguarding}
\phantomsection
\label{navigation:II.1.2.16.3.4}
\addcontentsline{toc}{paragraph}{II.1.2.16.3.4 [CT3] Health Information Safeguarding}

\label{element:6e89db75de5b4c4c9274f78c8d9e665a}
\index{[CT3] Health Information Safeguarding (II.1.2.16.3.4)}
{\small Capability · Canonical · current}\par

Ensure information and its operational environment are protected and resilient against loss, compromise, inappropriate access or operational failure.



Health Information Safeguarding protects information and its operational environment against loss, compromise, inappropriate access and operational failure. It frames resilience and protection as organisational outcomes, allowing physical controls, recovery arrangements and testing to contribute without defining the capability as a product. The concern is broader than any individual access transaction.



\textbf{Boundary}



Estate protection remains distinct from evaluating privacy consequences and from assurance conclusions about service delivery.


{\small Recorded unresolved matters: source-qualified matters; exact scope remains in the canonical record.}\par
{\footnotesize \href{../../../../continuum/strategy/capability/6e89db75de5b4c4c9274f78c8d9e665a.md}{Canonical record: 6e89db75de5b4c4c9274f78c8d9e665a}; full qualifications and provenance.}\par
\paragraph*{II.1.2.16.3.5 [CT3] Health Privacy Safeguard Management}
\phantomsection
\label{navigation:II.1.2.16.3.5}
\addcontentsline{toc}{paragraph}{II.1.2.16.3.5 [CT3] Health Privacy Safeguard Management}

\label{element:089cd8744ed6410388daeff463507fc1}
\index{[CT3] Health Privacy Safeguard Management (II.1.2.16.3.5)}
{\small Capability · Canonical · current}\par

Establish organisational safeguards and assurance processes concerning privacy consequences of information supply, exposure or use.



This capability establishes safeguards and assurance arrangements concerning the privacy consequences of supplying, exposing or using information. It focuses on the effect on people and legitimate information use, distinguishing that concern from protection of the information estate itself. Technical access restrictions may contribute without exhausting the organisational privacy responsibility.



\textbf{Boundary}



Privacy safeguards do not replace broader estate resilience, access decisions or originating consent authority.


{\small Recorded unresolved matters: source-qualified matters; exact scope remains in the canonical record.}\par
{\footnotesize \href{../../../../continuum/strategy/capability/089cd8744ed6410388daeff463507fc1.md}{Canonical record: 089cd8744ed6410388daeff463507fc1}; full qualifications and provenance.}\par
\subsection*{II.1.2.17 [CT1] Standards, Semantics \& Reference Governance}
\phantomsection
\label{navigation:II.1.2.17}
\addcontentsline{toc}{subsection}{II.1.2.17 [CT1] Standards, Semantics \& Reference Governance}

{\small Historical predecessor: R.2.51 \hyperref[element:3b65a3d007f54aa9b875e655475e1786]{[CT1] Standards, Semantics \& Reference Governance}; Canonical ID: \texttt{\seqsplit{3b65a3d007f54aa9b875e655475e1786}}. Superseded historical identity — direct correspondence · superseded.}\par
\label{element:5ffeb7401d314f06b05d79c3aff98242}
\index{[CT1] Standards, Semantics \& Reference Governance (II.1.2.17)}
{\small Capability · Canonical · current}\par

Govern organisational adoption, oversight and appropriate use of health-information models, standards, terminology, classifications and semantic mappings.



This capability establishes the organisational standing and appropriate use of models, standards and semantic artefacts. Standards governance addresses discovery, adoption, education, compliance and impacts; clinical-semantics governance addresses terminology, classifications and mappings. It provides the governed meaning within which information can be represented and exchanged, rather than merely maintaining technical schemas.



\textbf{Intent}



Support consistent, appropriate interpretation and use of health-information standards and semantic references.



\textbf{Boundary}



Using a terminology binding, mapping or classification does not transfer its governing authority. The established peer set must not be expanded simply to create a symmetrical taxonomy.



\textbf{Harmonia Architectural Relevance}



Semantic and conformance governance is a substantive platform responsibility. Its organisational meaning constrains technical representations; implementation products do not define that meaning.


{\footnotesize \href{../../../../continuum/strategy/capability/5ffeb7401d314f06b05d79c3aff98242.md}{Canonical record: 5ffeb7401d314f06b05d79c3aff98242}; full qualifications and provenance.}\par
\subsubsection*{II.1.2.17.1 [CT2] Clinical Semantics Governance}
\phantomsection
\label{navigation:II.1.2.17.1}
\addcontentsline{toc}{subsubsection}{II.1.2.17.1 [CT2] Clinical Semantics Governance}

\label{element:5cc8319e1dde4c818cc248a47233b67b}
\index{[CT2] Clinical Semantics Governance (II.1.2.17.1)}
{\small Capability · Canonical · current}\par

Govern clinical meaning and appropriate organisational use of semantic artefacts.



Clinical Semantics Governance maintains the standing and appropriate use of terminology, classifications and semantic mappings. Terminology provides consistent meaning; classification concerns sensitivity and consequences of exposure or mishandling. The established children distinguish governing artefacts from overseeing actual use and permit controlled terminology localisation where adopted material does not meet legitimate needs.



\textbf{Intent}



Keep meaning and classification intelligible across organisational use without concealing semantic drift or local variation.



\textbf{Boundary}



Compartmentalisation concerns contextual relationships rather than classification or terminology. Classification Localisation is not an established peer; mappings must preserve the separate meanings and authorities they connect.


{\small Recorded unresolved matters: source-qualified matters; exact scope remains in the canonical record.}\par
{\footnotesize \href{../../../../continuum/strategy/capability/5cc8319e1dde4c818cc248a47233b67b.md}{Canonical record: 5cc8319e1dde4c818cc248a47233b67b}; full qualifications and provenance.}\par
\paragraph*{II.1.2.17.1.1 [CT3] Classification Governance}
\phantomsection
\label{navigation:II.1.2.17.1.1}
\addcontentsline{toc}{paragraph}{II.1.2.17.1.1 [CT3] Classification Governance}

\label{element:4520ce37454e4de68580d75fdb08efe6}
\index{[CT3] Classification Governance (II.1.2.17.1.1)}
{\small Capability · Canonical · current}\par

Govern applicable classification definitions and their organisational standing.



Classification Governance establishes the definitions and organisational standing of applicable classifications. In this governed context, classification concerns sensitivity and the consequences of exposure or mishandling, rather than the clinical meaning supplied by terminology. Its outcome supplies a basis for consistent application and protection considerations.



\textbf{Boundary}



Classification is orthogonal to contextual compartment membership and does not by itself determine an access decision.


{\small Recorded unresolved matters: source-qualified matters; exact scope remains in the canonical record.}\par
{\footnotesize \href{../../../../continuum/strategy/capability/4520ce37454e4de68580d75fdb08efe6.md}{Canonical record: 4520ce37454e4de68580d75fdb08efe6}; full qualifications and provenance.}\par
\paragraph*{II.1.2.17.1.2 [CT3] Classification Utilisation Oversight}
\phantomsection
\label{navigation:II.1.2.17.1.2}
\addcontentsline{toc}{paragraph}{II.1.2.17.1.2 [CT3] Classification Utilisation Oversight}

\label{element:1d7c50aa80464613864dd6efb3fb56cc}
\index{[CT3] Classification Utilisation Oversight (II.1.2.17.1.2)}
{\small Capability · Canonical · current}\par

Oversee appropriate organisational application of classifications.



Classification Utilisation Oversight evaluates appropriate organisational application of classifications. It addresses whether information is treated consistently with the governed sensitivity meaning, rather than assuming that an established classification scheme guarantees correct use. This retains a distinct responsibility for practice alongside definition governance.



\textbf{Boundary}



Oversight neither creates classification definitions nor establishes a new localisation peer.


{\small Recorded unresolved matters: source-qualified matters; exact scope remains in the canonical record.}\par
{\footnotesize \href{../../../../continuum/strategy/capability/1d7c50aa80464613864dd6efb3fb56cc.md}{Canonical record: 1d7c50aa80464613864dd6efb3fb56cc}; full qualifications and provenance.}\par
\paragraph*{II.1.2.17.1.3 [CT3] Semantic Mapping Governance}
\phantomsection
\label{navigation:II.1.2.17.1.3}
\addcontentsline{toc}{paragraph}{II.1.2.17.1.3 [CT3] Semantic Mapping Governance}

\label{element:17dc6ab9f4394049a5327415d19628f2}
\index{[CT3] Semantic Mapping Governance (II.1.2.17.1.3)}
{\small Capability · Canonical · current}\par

Govern mappings between applicable semantic artefacts without collapsing their source meanings or authorities.



Semantic Mapping Governance maintains mappings between applicable semantic artefacts while preserving the meaning and authority of each side. Its purpose is intelligible correspondence rather than assumed equivalence. A governed mapping can support exchange or interpretation without suggesting that the mapped artefacts have become the same concept or share one authority.



\textbf{Boundary}



A mapping does not erase source distinctions or give a consuming system authority to redefine either source.


{\small Recorded unresolved matters: source-qualified matters; exact scope remains in the canonical record.}\par
{\footnotesize \href{../../../../continuum/strategy/capability/17dc6ab9f4394049a5327415d19628f2.md}{Canonical record: 17dc6ab9f4394049a5327415d19628f2}; full qualifications and provenance.}\par
\paragraph*{II.1.2.17.1.4 [CT3] Terminology Governance}
\phantomsection
\label{navigation:II.1.2.17.1.4}
\addcontentsline{toc}{paragraph}{II.1.2.17.1.4 [CT3] Terminology Governance}

\label{element:e926a47881fa482088859f59fa6fe7e4}
\index{[CT3] Terminology Governance (II.1.2.17.1.4)}
{\small Capability · Canonical · current}\par

Manage governed terminology/reference-set lifecycle and standing, including acceptance, release, update, supersession/retirement and authority/provider coordination.



Terminology Governance maintains the lifecycle and standing of terminology and reference sets used by the organisation. Acceptance, release, updates and retirement require coordination with the appropriate authority or provider so that a code's meaning is not treated as timeless or locally self-defined. This supplies a governed semantic basis for information bindings and knowledge associations.



\textbf{Boundary}



Governing terminology differs from binding a structure to it or overseeing its actual use.


{\small Recorded unresolved matters: source-qualified matters; exact scope remains in the canonical record.}\par
{\footnotesize \href{../../../../continuum/strategy/capability/e926a47881fa482088859f59fa6fe7e4.md}{Canonical record: e926a47881fa482088859f59fa6fe7e4}; full qualifications and provenance.}\par
\paragraph*{II.1.2.17.1.5 [CT3] Terminology Localisation}
\phantomsection
\label{navigation:II.1.2.17.1.5}
\addcontentsline{toc}{paragraph}{II.1.2.17.1.5 [CT3] Terminology Localisation}

\label{element:3aa80e1d84ea4e009b149c4f9a23f0f8}
\index{[CT3] Terminology Localisation (II.1.2.17.1.5)}
{\small Capability · Canonical · current}\par

Manage controlled local responses where adopted terminology does not meet legitimate organisational semantic requirements.



Terminology Localisation manages controlled local responses where adopted terminology does not satisfy legitimate organisational semantic needs. It makes the local variation explicit rather than allowing ad hoc usage to conceal a gap or change in meaning. Within clinical-semantics governance, this preserves a distinction between governed local adaptation and uncontrolled drift.



\textbf{Boundary}



The established localisation responsibility concerns terminology; no corresponding Classification Localisation capability is inferred.


{\small Recorded unresolved matters: source-qualified matters; exact scope remains in the canonical record.}\par
{\footnotesize \href{../../../../continuum/strategy/capability/3aa80e1d84ea4e009b149c4f9a23f0f8.md}{Canonical record: 3aa80e1d84ea4e009b149c4f9a23f0f8}; full qualifications and provenance.}\par
\paragraph*{II.1.2.17.1.6 [CT3] Terminology Utilisation Oversight}
\phantomsection
\label{navigation:II.1.2.17.1.6}
\addcontentsline{toc}{paragraph}{II.1.2.17.1.6 [CT3] Terminology Utilisation Oversight}

\label{element:4ca436cd90b745daa290b067c57a0450}
\index{[CT3] Terminology Utilisation Oversight (II.1.2.17.1.6)}
{\small Capability · Canonical · current}\par

Oversee actual terminology/reference-set use and identify inappropriate, inconsistent or drifting usage.



Terminology Utilisation Oversight considers how terminology and reference sets are actually used. It identifies inappropriate application, inconsistency or semantic drift that the existence of an approved terminology does not prevent. Its purpose is to preserve coherent use without making each consuming capability the owner of terminology meaning.



\textbf{Boundary}



Use oversight does not redefine terminology or automatically authorise a local semantic variation.


{\small Recorded unresolved matters: source-qualified matters; exact scope remains in the canonical record.}\par
{\footnotesize \href{../../../../continuum/strategy/capability/4ca436cd90b745daa290b067c57a0450.md}{Canonical record: 4ca436cd90b745daa290b067c57a0450}; full qualifications and provenance.}\par
\subsubsection*{II.1.2.17.2 [CT2] Health Information Standards Governance}
\phantomsection
\label{navigation:II.1.2.17.2}
\addcontentsline{toc}{subsubsection}{II.1.2.17.2 [CT2] Health Information Standards Governance}

\label{element:b228039d8e764f9b98cb810b68af329d}
\index{[CT2] Health Information Standards Governance (II.1.2.17.2)}
{\small Capability · Canonical · current}\par

Govern organisational adoption and use of health-information models and applicable standards.



This capability governs which information models and standards the organisation adopts and how their use is supported and evaluated. The model distinguishes discovering candidates, adopting models, educating users, addressing conformity and overseeing governance. Standards impact identification has a separate established place because its implications may extend into clinical practice, policy and obligations beyond information-system implementation.



\textbf{Intent}



Make organisational standards use deliberate and its consequences understandable.



\textbf{Boundary}



Oversight evaluates governance without performing every peer activity. Adoption includes its own model-impact consideration; the listed abilities do not establish a sequential adoption Process.


{\small Recorded unresolved matters: source-qualified matters; exact scope remains in the canonical record.}\par
{\footnotesize \href{../../../../continuum/strategy/capability/b228039d8e764f9b98cb810b68af329d.md}{Canonical record: b228039d8e764f9b98cb810b68af329d}; full qualifications and provenance.}\par
\paragraph*{II.1.2.17.2.1 [CT3] Information Model Adoption}
\phantomsection
\label{navigation:II.1.2.17.2.1}
\addcontentsline{toc}{paragraph}{II.1.2.17.2.1 [CT3] Information Model Adoption}

\label{element:0c3ce0778d6342b6b100e1e30e777fbe}
\index{[CT3] Information Model Adoption (II.1.2.17.2.1)}
{\small Capability · Canonical · current}\par

Adopt applicable information models, including intrinsic assessment of adoption impacts.



Information Model Adoption establishes the organisation's use of an applicable information model, including consideration of the consequences of adopting it. It gives a discovered model an organisational standing distinct from mere availability. The intrinsic impact consideration helps adoption remain a business governance responsibility rather than a mechanical import of definitions.



\textbf{Boundary}



Adoption does not perform every education, compliance or governance-oversight responsibility associated with subsequent use.


{\small Recorded unresolved matters: source-qualified matters; exact scope remains in the canonical record.}\par
{\footnotesize \href{../../../../continuum/strategy/capability/0c3ce0778d6342b6b100e1e30e777fbe.md}{Canonical record: 0c3ce0778d6342b6b100e1e30e777fbe}; full qualifications and provenance.}\par
\paragraph*{II.1.2.17.2.2 [CT3] Information Model Compliance}
\phantomsection
\label{navigation:II.1.2.17.2.2}
\addcontentsline{toc}{paragraph}{II.1.2.17.2.2 [CT3] Information Model Compliance}

\label{element:160db14dbe5149db91223a93f6ee9102}
\index{[CT3] Information Model Compliance (II.1.2.17.2.2)}
{\small Capability · Canonical · current}\par

Address organisational conformity to applicable information-model requirements.



Information Model Compliance addresses conformity to applicable model requirements. It makes actual organisational use accountable to the adopted semantic frame rather than treating model availability or user education as proof of conformity. Within standards governance it concerns the requirements for use, not the separate task of discovering or adopting models.



\textbf{Boundary}



Conformity does not redefine the model or absorb governance oversight of the wider adoption arrangements.


{\small Recorded unresolved matters: source-qualified matters; exact scope remains in the canonical record.}\par
{\footnotesize \href{../../../../continuum/strategy/capability/160db14dbe5149db91223a93f6ee9102.md}{Canonical record: 160db14dbe5149db91223a93f6ee9102}; full qualifications and provenance.}\par
\paragraph*{II.1.2.17.2.3 [CT3] Information Model Discovery}
\phantomsection
\label{navigation:II.1.2.17.2.3}
\addcontentsline{toc}{paragraph}{II.1.2.17.2.3 [CT3] Information Model Discovery}

\label{element:f33dc5b8ba9643329bb7b9431e1f9a6f}
\index{[CT3] Information Model Discovery (II.1.2.17.2.3)}
{\small Capability · Canonical · current}\par

Identify information models relevant to the organisation’s health-information needs.



Information Model Discovery identifies candidate models relevant to the organisation's health-information needs. It supports informed standards governance by making applicable possibilities visible before adoption. Its business purpose is finding suitable semantic frames rather than selecting a software representation or assuming every discovered model should be used.



\textbf{Boundary}



Discovery does not adopt a model, establish compliance or grant it organisational authority.


{\small Recorded unresolved matters: source-qualified matters; exact scope remains in the canonical record.}\par
{\footnotesize \href{../../../../continuum/strategy/capability/f33dc5b8ba9643329bb7b9431e1f9a6f.md}{Canonical record: f33dc5b8ba9643329bb7b9431e1f9a6f}; full qualifications and provenance.}\par
\paragraph*{II.1.2.17.2.4 [CT3] Information Model Education}
\phantomsection
\label{navigation:II.1.2.17.2.4}
\addcontentsline{toc}{paragraph}{II.1.2.17.2.4 [CT3] Information Model Education}

\label{element:7dcb6427f9134fb7a4a63a5ec1d5c85b}
\index{[CT3] Information Model Education (II.1.2.17.2.4)}
{\small Capability · Canonical · current}\par

Support organisational understanding and appropriate use of adopted information models.



Information Model Education supports understanding and appropriate use of adopted models. It addresses the people and practice dimension of semantic consistency: a published model is not necessarily a model understood in the contexts where it is applied. Its responsibility therefore complements adoption and compliance without becoming either.



\textbf{Boundary}



Educational support neither changes the authoritative model nor establishes that actual use complies with it.


{\small Recorded unresolved matters: source-qualified matters; exact scope remains in the canonical record.}\par
{\footnotesize \href{../../../../continuum/strategy/capability/7dcb6427f9134fb7a4a63a5ec1d5c85b.md}{Canonical record: 7dcb6427f9134fb7a4a63a5ec1d5c85b}; full qualifications and provenance.}\par
\paragraph*{II.1.2.17.2.5 [CT3] Information Model Governance Oversight}
\phantomsection
\label{navigation:II.1.2.17.2.5}
\addcontentsline{toc}{paragraph}{II.1.2.17.2.5 [CT3] Information Model Governance Oversight}

\label{element:1e71b707beab450f9cbfaf0bf88b6ec4}
\index{[CT3] Information Model Governance Oversight (II.1.2.17.2.5)}
{\small Capability · Canonical · current}\par

Oversee and evaluate effectiveness, consistency and appropriate operation of information-model governance.



This capability evaluates the effectiveness and consistency of information-model governance itself. It asks whether discovery, adoption and use are being governed appropriately, rather than repeating each operational peer. Its purpose is coherent oversight of the arrangements that give models standing and guide their application.



\textbf{Boundary}



Oversight does not itself perform discovery, adoption, education or compliance activities.


{\small Recorded unresolved matters: source-qualified matters; exact scope remains in the canonical record.}\par
{\footnotesize \href{../../../../continuum/strategy/capability/1e71b707beab450f9cbfaf0bf88b6ec4.md}{Canonical record: 1e71b707beab450f9cbfaf0bf88b6ec4}; full qualifications and provenance.}\par
\paragraph*{II.1.2.17.2.6 [CT3] Standards Adoption}
\phantomsection
\label{navigation:II.1.2.17.2.6}
\addcontentsline{toc}{paragraph}{II.1.2.17.2.6 [CT3] Standards Adoption}

\label{element:0d59f956f5714d3985398f09aeb7c541}
\index{[CT3] Standards Adoption (II.1.2.17.2.6)}
{\small Capability · Canonical · current}\par

Adopt applicable standards within the organisation.



Standards Adoption establishes which applicable standards the organisation adopts. It gives standards an intentional place in organisational practice rather than allowing familiar products or conventions to define requirements by default. The business responsibility concerns adoption in context, including standards whose effects extend beyond information representation.



\textbf{Boundary}



Adoption does not establish a complete standards lifecycle taxonomy or perform the distinct identification of broader standards impacts.


{\small Recorded unresolved matters: source-qualified matters; exact scope remains in the canonical record.}\par
{\footnotesize \href{../../../../continuum/strategy/capability/0d59f956f5714d3985398f09aeb7c541.md}{Canonical record: 0d59f956f5714d3985398f09aeb7c541}; full qualifications and provenance.}\par
\paragraph*{II.1.2.17.2.7 [CT3] Standards Impact Identification}
\phantomsection
\label{navigation:II.1.2.17.2.7}
\addcontentsline{toc}{paragraph}{II.1.2.17.2.7 [CT3] Standards Impact Identification}

\label{element:6b9c96c7d9174e7fbfd11267673d42f3}
\index{[CT3] Standards Impact Identification (II.1.2.17.2.7)}
{\small Capability · Canonical · current}\par

Identify standards impacts across clinical/work practice, permitted medications, procedures, policy, safety, obligations and reporting as applicable.



Standards Impact Identification recognises the consequences of standards for relevant organisational activity. Impacts may concern clinical practice, medications, procedures, policy, safety, obligations or reporting, so the responsibility cannot be reduced to technical compatibility. It makes those consequences visible for the responsible capabilities to consider within their own authority.



\textbf{Boundary}



Identifying an impact neither adopts the standard nor acquires the operational authority needed to respond to it.


{\small Recorded unresolved matters: source-qualified matters; exact scope remains in the canonical record.}\par
{\footnotesize \href{../../../../continuum/strategy/capability/6b9c96c7d9174e7fbfd11267673d42f3.md}{Canonical record: 6b9c96c7d9174e7fbfd11267673d42f3}; full qualifications and provenance.}\par
\subsection*{II.1.2.18 [CT1] Workforce \& Organisational Capability}
\phantomsection
\label{navigation:II.1.2.18}
\addcontentsline{toc}{subsection}{II.1.2.18 [CT1] Workforce \& Organisational Capability}

{\small Historical predecessor: R.2.52 \hyperref[element:b1161fe1a315420fbc7d048a1ea58126]{[CT1] Workforce \& Organisational Capability}; Canonical ID: \texttt{\seqsplit{b1161fe1a315420fbc7d048a1ea58126}}. Superseded historical identity — direct correspondence · superseded.}\par
\label{element:631b43664bfb4c50ac8cf4958b2c347d}
\index{[CT1] Workforce \& Organisational Capability (II.1.2.18)}
{\small Capability · Canonical · current}\par

Govern practitioner qualification, engagement/privilege and enterprise service coverage, together with organisation identity/structure, recognised locations and identifiable healthcare service offerings.



This capability governs the workforce and organisational conditions under which health services can operate. Its workforce branch addresses professional standing, participation, privilege and enterprise coverage; its structural branch addresses organisations, recognised locations and identifiable service offerings. Directories can represent these facts, but their representation is separate from business authority to establish or change them.



\textbf{Intent}



Make the organisational and workforce context of care explicit enough to constrain legitimate participation and service availability.



\textbf{Boundary}



Qualification, engagement, privilege, service coverage and deployment are distinct. Organisation or location association does not substitute for material Healthcare Service context.



\textbf{Harmonia Architectural Relevance}



These facts materially constrain care enablement, routing and operational context. Directory support preserves the originating workforce and organisational governance rather than making it platform-owned.


{\footnotesize \href{../../../../continuum/strategy/capability/631b43664bfb4c50ac8cf4958b2c347d.md}{Canonical record: 631b43664bfb4c50ac8cf4958b2c347d}; full qualifications and provenance.}\par
\subsubsection*{II.1.2.18.1 [CT2] Health Workforce Management}
\phantomsection
\label{navigation:II.1.2.18.1}
\addcontentsline{toc}{subsubsection}{II.1.2.18.1 [CT2] Health Workforce Management}

\label{element:7ec329cdc070472eb68966345d82031a}
\index{[CT2] Health Workforce Management (II.1.2.18.1)}
{\small Capability · Canonical · current}\par

Govern practitioners’ professional standing, engagement/privilege and whole-enterprise service coverage.



Health Workforce Management governs practitioners' professional standing, authorised participation and the ability to cover required enterprise work. Its three children separate qualification from engagement and privilege, and both from whole-enterprise coverage. This prevents an available or rostered person from being treated as automatically qualified and privileged for a particular activity.



\textbf{Intent}



Support legitimate practitioner participation and an informed understanding of service-coverage gaps and risks.



\textbf{Boundary}



Coverage is broader than rostering. Applying an available resource to a care requirement and progressing a particular support work item remain separate operational responsibilities.


{\small Recorded unresolved matters: source-qualified matters; exact scope remains in the canonical record.}\par
{\footnotesize \href{../../../../continuum/strategy/capability/7ec329cdc070472eb68966345d82031a.md}{Canonical record: 7ec329cdc070472eb68966345d82031a}; full qualifications and provenance.}\par
\paragraph*{II.1.2.18.1.1 [CT3] Practitioner Engagement \& Privilege Management}
\phantomsection
\label{navigation:II.1.2.18.1.1}
\addcontentsline{toc}{paragraph}{II.1.2.18.1.1 [CT3] Practitioner Engagement \& Privilege Management}

\label{element:36d5bf4c19d34172b72f4b77099ff235}
\index{[CT3] Practitioner Engagement \& Privilege Management (II.1.2.18.1.1)}
{\small Capability · Canonical · current}\par

Manage practitioners as service participants, their roles/capacities, applicable privileges and qualification context above ordinary HR employment administration.



This capability manages how practitioners participate in services and the privileges applicable to that participation. It connects professional standing with roles and permitted capacities in the service context, beyond ordinary employment administration. Separating engagement and privilege from qualification makes it possible to recognise a qualified practitioner who is not authorised for a particular service activity.



\textbf{Boundary}



Representation or consumption of privilege information does not transfer the underlying workforce decision authority to Harmonia.


{\small Recorded unresolved matters: source-qualified matters; exact scope remains in the canonical record.}\par
{\footnotesize \href{../../../../continuum/strategy/capability/36d5bf4c19d34172b72f4b77099ff235.md}{Canonical record: 36d5bf4c19d34172b72f4b77099ff235}; full qualifications and provenance.}\par
\paragraph*{II.1.2.18.1.2 [CT3] Practitioner Qualification Management}
\phantomsection
\label{navigation:II.1.2.18.1.2}
\addcontentsline{toc}{paragraph}{II.1.2.18.1.2 [CT3] Practitioner Qualification Management}

\label{element:279bab5de1b74c17a5ee2db5cae8d385}
\index{[CT3] Practitioner Qualification Management (II.1.2.18.1.2)}
{\small Capability · Canonical · current}\par

Govern professional/clinical standing and qualifications, including applicable review, escalation, complaints, penalties/issues, suspension and revocation.



Practitioner Qualification Management governs professional and clinical standing, including review and changes such as suspension or revocation. It makes qualification a continuing governed concern rather than a static directory attribute. That standing can constrain engagement and privilege while remaining distinguishable from either an employment relationship or permission for a particular activity.



\textbf{Boundary}



Qualification alone does not establish engagement, clinical privilege, device-specific enablement or operational assignment.


{\small Recorded unresolved matters: source-qualified matters; exact scope remains in the canonical record.}\par
{\footnotesize \href{../../../../continuum/strategy/capability/279bab5de1b74c17a5ee2db5cae8d385.md}{Canonical record: 279bab5de1b74c17a5ee2db5cae8d385}; full qualifications and provenance.}\par
\paragraph*{II.1.2.18.1.3 [CT3] Workforce Service Coverage Management}
\phantomsection
\label{navigation:II.1.2.18.1.3}
\addcontentsline{toc}{paragraph}{II.1.2.18.1.3 [CT3] Workforce Service Coverage Management}

\label{element:cfb04e13895745619171a3e8f15d349c}
\index{[CT3] Workforce Service Coverage Management (II.1.2.18.1.3)}
{\small Capability · Canonical · current}\par

Observe whether required enterprise work/service obligations can be fulfilled, identify coverage gaps/risks and adjust coverage in response.



Workforce Service Coverage Management considers whether required enterprise work and service obligations can be fulfilled and how coverage should change when gaps or risks arise. Its concern is the adequacy of coverage, broader than the existence of a roster. This supplies a service-wide workforce view distinct from deploying a person to a particular care requirement.



\textbf{Boundary}



Coverage management is not individual support-work progression or automatic authority to assign clinical work.


{\small Recorded unresolved matters: source-qualified matters; exact scope remains in the canonical record.}\par
{\footnotesize \href{../../../../continuum/strategy/capability/cfb04e13895745619171a3e8f15d349c.md}{Canonical record: cfb04e13895745619171a3e8f15d349c}; full qualifications and provenance.}\par
\subsubsection*{II.1.2.18.2 [CT2] Healthcare Organisation \& Service Structure Management}
\phantomsection
\label{navigation:II.1.2.18.2}
\addcontentsline{toc}{subsubsection}{II.1.2.18.2 [CT2] Healthcare Organisation \& Service Structure Management}

\label{element:de8c1b408c774e59899662495f4d6ced}
\index{[CT2] Healthcare Organisation \& Service Structure Management (II.1.2.18.2)}
{\small Capability · Canonical · current}\par

Manage organisational identity/structure, recognised locations and identifiable organisational service offerings.



This capability establishes the organisation, location and service-offering context through which health services are recognised. Its children distinguish the identity and structure of participants, the definition of places, and what identifiable services are offered under temporal and conditional availability constraints. The distinction allows a service to exist without assuming it is offered everywhere or currently available.



\textbf{Intent}



Provide an intelligible account of who offers what service, where and under which conditions.



\textbf{Boundary}



Location definition is not facilities operation. Offering availability does not establish that all resources and dependencies are collectively ready to deliver the service.


{\small Recorded unresolved matters: source-qualified matters; exact scope remains in the canonical record.}\par
{\footnotesize \href{../../../../continuum/strategy/capability/de8c1b408c774e59899662495f4d6ced.md}{Canonical record: de8c1b408c774e59899662495f4d6ced}; full qualifications and provenance.}\par
\paragraph*{II.1.2.18.2.1 [CT3] Healthcare Location Definition Management}
\phantomsection
\label{navigation:II.1.2.18.2.1}
\addcontentsline{toc}{paragraph}{II.1.2.18.2.1 [CT3] Healthcare Location Definition Management}

\label{element:d77a5c8f1818465c98125acc4bc2ac81}
\index{[CT3] Healthcare Location Definition Management (II.1.2.18.2.1)}
{\small Capability · Canonical · current}\par

Establish and manage recognised locations and their applicable business characteristics and relationships.



Healthcare Location Definition Management establishes recognised places and their business characteristics and relationships. It gives health-service activity an intelligible location context while remaining separate from maintaining the physical facilities. A defined location can be referenced in offerings and operational activity without implying that it is currently usable or ready for care.



\textbf{Boundary}



Location definition is not facilities operation, care-place readiness or service availability.


{\small Recorded unresolved matters: source-qualified matters; exact scope remains in the canonical record.}\par
{\footnotesize \href{../../../../continuum/strategy/capability/d77a5c8f1818465c98125acc4bc2ac81.md}{Canonical record: d77a5c8f1818465c98125acc4bc2ac81}; full qualifications and provenance.}\par
\paragraph*{II.1.2.18.2.2 [CT3] Healthcare Organisation Identity \& Structure Management}
\phantomsection
\label{navigation:II.1.2.18.2.2}
\addcontentsline{toc}{paragraph}{II.1.2.18.2.2 [CT3] Healthcare Organisation Identity \& Structure Management}

\label{element:4684348b4d144e61aa7dd203facfebdd}
\index{[CT3] Healthcare Organisation Identity \& Structure Management (II.1.2.18.2.2)}
{\small Capability · Canonical · current}\par

Establish and manage identity and structural relationships of participating organisations and organisational units.



This capability establishes the identity and structural relationships of participating organisations and organisational units. It provides the business account of who the participants are and how they relate, which directory representations can then convey. The distinction matters because a represented hierarchy should not create organisational authority merely through document or interface nesting.



\textbf{Boundary}



Organisation identity does not substitute for location, practitioner or material Healthcare Service context.


{\small Recorded unresolved matters: source-qualified matters; exact scope remains in the canonical record.}\par
{\footnotesize \href{../../../../continuum/strategy/capability/4684348b4d144e61aa7dd203facfebdd.md}{Canonical record: 4684348b4d144e61aa7dd203facfebdd}; full qualifications and provenance.}\par
\paragraph*{II.1.2.18.2.3 [CT3] Healthcare Service Offering Management}
\phantomsection
\label{navigation:II.1.2.18.2.3}
\addcontentsline{toc}{paragraph}{II.1.2.18.2.3 [CT3] Healthcare Service Offering Management}

\label{element:9fe177e4422b4853acb1dd0d0b74c4ba}
\index{[CT3] Healthcare Service Offering Management (II.1.2.18.2.3)}
{\small Capability · Canonical · current}\par

Establish/manage what identifiable service is offered, by whom, where, when and under which availability conditions/constraints.



This capability establishes the identifiable service offered, including provider, place, timing and availability conditions. It makes a service discoverable as an offering rather than merely an abstract service idea. Temporal and conditional availability are material: an offering may exist but be unavailable in the present context, and another pathway may cover different hours.



\textbf{Boundary}



An available offering does not establish collective readiness to deliver it; organisation and location alone do not supply the service context.


{\small Recorded unresolved matters: source-qualified matters; exact scope remains in the canonical record.}\par
{\footnotesize \href{../../../../continuum/strategy/capability/9fe177e4422b4853acb1dd0d0b74c4ba.md}{Canonical record: 9fe177e4422b4853acb1dd0d0b74c4ba}; full qualifications and provenance.}\par
\section*{II.1.3 Business Enabling Capabilities}
\phantomsection
\label{navigation:II.1.3}
\addcontentsline{toc}{section}{II.1.3 Business Enabling Capabilities}

\subsection*{II.1.3.1 [CT1] Care Administration Enablement}
\phantomsection
\label{navigation:II.1.3.1}
\addcontentsline{toc}{subsection}{II.1.3.1 [CT1] Care Administration Enablement}

{\small Historical predecessor: R.2.71 \hyperref[element:b611a1fcfadb4b30af9ba74a4c993d8a]{[CT1] Care Administration Enablement}; Canonical ID: \texttt{\seqsplit{b611a1fcfadb4b30af9ba74a4c993d8a}}. Superseded historical identity — direct correspondence · superseded.}\par
\label{element:28df27d55c514a37995145ab46945f78}
\index{[CT1] Care Administration Enablement (II.1.3.1)}
{\small Capability · Proposed · current}\par

Administer the information and progression of care access, requests, encounters, records and plans across responsible care systems.



\textbf{Overview}



Administer the information and progression of care access, requests, encounters, records and plans across responsible care systems.



\textbf{Intent}



Administer the information and progression of care access, requests, encounters, records and plans across responsible care systems.



\textbf{Boundary}



Information administration does not perform the clinical act, own host booking engines or infer Transfer of Care from referral acceptance.


{\footnotesize \href{../../../../continuum/strategy/capability/28df27d55c514a37995145ab46945f78.md}{Canonical record: 28df27d55c514a37995145ab46945f78}; full qualifications and provenance.}\par
\subsubsection*{II.1.3.1.1 [CT2] Care Access Transaction Enablement [superseded]}
\phantomsection
\label{navigation:II.1.3.1.1}
\addcontentsline{toc}{subsubsection}{II.1.3.1.1 [CT2] Care Access Transaction Enablement [superseded]}

\label{element:b65973b576154d3a96cf458bf33df53a}
\index{[CT2] Care Access Transaction Enablement (II.1.3.1.1)}
{\small Capability · Proposed · superseded}\par
\href{../../../../continuum/strategy/capability/b65973b576154d3a96cf458bf33df53a.md}{Historical canonical definition and provenance}.\par
\subsubsection*{II.1.3.1.2 [CT2] Care Documentation \& Communication Enablement [superseded]}
\phantomsection
\label{navigation:II.1.3.1.2}
\addcontentsline{toc}{subsubsection}{II.1.3.1.2 [CT2] Care Documentation \& Communication Enablement [superseded]}

\label{element:5ad2e381ec484d1d9062a65eee842fa7}
\index{[CT2] Care Documentation \& Communication Enablement (II.1.3.1.2)}
{\small Capability · Proposed · superseded}\par
\href{../../../../continuum/strategy/capability/5ad2e381ec484d1d9062a65eee842fa7.md}{Historical canonical definition and provenance}.\par
\subsubsection*{II.1.3.1.3 [CT2] Care Financial Information Enablement [superseded]}
\phantomsection
\label{navigation:II.1.3.1.3}
\addcontentsline{toc}{subsubsection}{II.1.3.1.3 [CT2] Care Financial Information Enablement [superseded]}

\label{element:59a756b89c604e3e97ee5dede2d6ed89}
\index{[CT2] Care Financial Information Enablement (II.1.3.1.3)}
{\small Capability · Proposed · superseded}\par
\href{../../../../continuum/strategy/capability/59a756b89c604e3e97ee5dede2d6ed89.md}{Historical canonical definition and provenance}.\par
\subsubsection*{II.1.3.1.4 [CT2] Care Fulfilment Administration Enablement [superseded]}
\phantomsection
\label{navigation:II.1.3.1.4}
\addcontentsline{toc}{subsubsection}{II.1.3.1.4 [CT2] Care Fulfilment Administration Enablement [superseded]}

\label{element:9f0f4379c1954524ba5e691d0c7d95ab}
\index{[CT2] Care Fulfilment Administration Enablement (II.1.3.1.4)}
{\small Capability · Proposed · superseded}\par
\href{../../../../continuum/strategy/capability/9f0f4379c1954524ba5e691d0c7d95ab.md}{Historical canonical definition and provenance}.\par
\subsubsection*{II.1.3.1.5 [CT2] Care Plan Administration}
\phantomsection
\label{navigation:II.1.3.1.5}
\addcontentsline{toc}{subsubsection}{II.1.3.1.5 [CT2] Care Plan Administration}

{\small Historical predecessor: R.2.54 \hyperref[element:e31407cea7624632a8ed2dc53b09d726]{[CT2] Care Plan Administration}; Canonical ID: \texttt{\seqsplit{e31407cea7624632a8ed2dc53b09d726}}. Superseded historical identity — direct correspondence · superseded.}\par
\label{element:cf06eec5699f406fa30a00d2f76e3098}
\index{[CT2] Care Plan Administration (II.1.3.1.5)}
{\small Capability · Canonical · current}\par

Care Plan Administration receives care plans and coordinates the distribution of updates to authorised team members. It supports shared understanding of established clinical intent across participating organisations. The responsibility concerns the plan as managed information, distinct from deciding the intended course of care or performing the actions expressed by it.



\textbf{Overview}



Care Plan Administration receives care plans and coordinates the distribution of updates to authorised team members. It supports shared understanding of established clinical intent across participating organisations. The responsibility concerns the plan as managed information, distinct from deciding the intended course of care or performing the actions expressed by it.



\textbf{Intent}



Care Plan Administration receives care plans and coordinates the distribution of updates to authorised team members. It supports shared understanding of established clinical intent across participating organisations. The responsibility concerns the plan as managed information, distinct from deciding the intended course of care or performing the actions expressed by it.



\textbf{Boundary}



Information administration does not confer authority over the represented person, organisation, service, place, device or clinical activity.



\medskip\hrule\medskip

{\footnotesize \href{../../../../continuum/strategy/capability/cf06eec5699f406fa30a00d2f76e3098.md}{Canonical record: cf06eec5699f406fa30a00d2f76e3098}; full qualifications and provenance.}\par
\paragraph*{II.1.3.1.5.1 [CT3] Care Plan Lifecycle \& Coordination}
\phantomsection
\label{navigation:II.1.3.1.5.1}
\addcontentsline{toc}{paragraph}{II.1.3.1.5.1 [CT3] Care Plan Lifecycle \& Coordination}

{\small Historical predecessor: R.2.53 \hyperref[element:689bb3f8a24249b79f787e1f48eb268d]{[CT3] Care Plan Lifecycle \& Coordination}; Canonical ID: \texttt{\seqsplit{689bb3f8a24249b79f787e1f48eb268d}}. Superseded historical identity — direct correspondence · superseded.}\par
\label{element:4cb7db65392f4fab801d7794af4d4b29}
\index{[CT3] Care Plan Lifecycle \& Coordination (II.1.3.1.5.1)}
{\small Capability · Proposed · current}\par

Administer received care plans, goals, revisions and cross-participant coordination while retaining clinical planning authority at its source.



\textbf{Overview}



Administer received care plans, goals, revisions and cross-participant coordination while retaining clinical planning authority at its source. The reviewed hierarchy assigns this enduring ability to its stated business-enabling context.



\textbf{Intent}



Administer received care plans, goals, revisions and cross-participant coordination while retaining clinical planning authority at its source.



\textbf{Boundary}



Information administration and coordination preserve the authority of the responsible source, clinical actor and governed domain; they confer no additional real-world management authority.


{\footnotesize \href{../../../../continuum/strategy/capability/4cb7db65392f4fab801d7794af4d4b29.md}{Canonical record: 4cb7db65392f4fab801d7794af4d4b29}; full qualifications and provenance.}\par
\subparagraph*{II.1.3.1.5.1.1 [Feature] Care Goal \& Plan Progression Representation}
\phantomsection
\label{navigation:II.1.3.1.5.1.1}
\addcontentsline{toc}{subparagraph}{II.1.3.1.5.1.1 [Feature] Care Goal \& Plan Progression Representation}

\label{element:7332de92e6c14bc8b071cf18f78f9525}
\index{[Feature] Care Goal \& Plan Progression Representation (II.1.3.1.5.1.1)}
{\small Capability / Feature · Proposed · current}\par

Represent externally agreed care goals, participants and plan changes with their effective standing.



\textbf{Overview}



Plan ingestion/distribution is complemented by the ability to distinguish the applicable goal and plan context as a shared plan progresses.



\textbf{Boundary}



Representing a goal does not choose treatment, determine clinical completion or grant plan approval authority.


{\footnotesize \href{../../../../continuum/strategy/capability/feature/7332de92e6c14bc8b071cf18f78f9525.md}{Canonical record: 7332de92e6c14bc8b071cf18f78f9525}; full qualifications and provenance.}\par
\subparagraph*{II.1.3.1.5.1.2 [Feature] Care Plan Ingestion \& Distribution}
\phantomsection
\label{navigation:II.1.3.1.5.1.2}
\addcontentsline{toc}{subparagraph}{II.1.3.1.5.1.2 [Feature] Care Plan Ingestion \& Distribution}

\label{element:6d4b374a6bca4b58a86c912a94bbd6b1}
\index{[Feature] Care Plan Ingestion \& Distribution (II.1.3.1.5.1.2)}
{\small Capability / Feature · Canonical · current}\par

Ingest multidisciplinary care plans and syndicate updates to authorized care team members.



Care-plan exchange keeps authorised team members aware of established multidisciplinary intent and its updates. It provides information continuity while preserving clinical planning authority and the distinction between a distributed plan and care actually performed.


{\footnotesize \href{../../../../continuum/strategy/capability/feature/6d4b374a6bca4b58a86c912a94bbd6b1.md}{Canonical record: 6d4b374a6bca4b58a86c912a94bbd6b1}; full qualifications and provenance.}\par
\paragraph*{II.1.3.1.5.2 [CT3] Community Care Coordination}
\phantomsection
\label{navigation:II.1.3.1.5.2}
\addcontentsline{toc}{paragraph}{II.1.3.1.5.2 [CT3] Community Care Coordination}

\label{element:8c8c6266d9ed43e09dbe44fcdd669132}
\index{[CT3] Community Care Coordination (II.1.3.1.5.2)}
{\small Capability · Canonical · current}\par

Definition unresolved in the baseline. No definition has been synthesised from the name.



\medskip\hrule\medskip

{\footnotesize \href{../../../../continuum/strategy/capability/8c8c6266d9ed43e09dbe44fcdd669132.md}{Canonical record: 8c8c6266d9ed43e09dbe44fcdd669132}; full qualifications and provenance.}\par
\subparagraph*{II.1.3.1.5.2.1 [Feature] Community Care Plan Synchronization}
\phantomsection
\label{navigation:II.1.3.1.5.2.1}
\addcontentsline{toc}{subparagraph}{II.1.3.1.5.2.1 [Feature] Community Care Plan Synchronization}

\label{element:30a0a2a6f1044d9ca9337b655f303e24}
\index{[Feature] Community Care Plan Synchronization (II.1.3.1.5.2.1)}
{\small Capability / Feature · Canonical · current}\par

Synchronize community care plan updates between hospital outreach teams and community health organizations.



Community plan synchronisation keeps hospital outreach and community participants informed of established plan changes. It supports shared clinical intent without allowing synchronisation to become authority to choose that intent or allocate clinical work.


{\footnotesize \href{../../../../continuum/strategy/capability/feature/30a0a2a6f1044d9ca9337b655f303e24.md}{Canonical record: 30a0a2a6f1044d9ca9337b655f303e24}; full qualifications and provenance.}\par
\subparagraph*{II.1.3.1.5.2.2 [Feature] Community Encounter Ingestion}
\phantomsection
\label{navigation:II.1.3.1.5.2.2}
\addcontentsline{toc}{subparagraph}{II.1.3.1.5.2.2 [Feature] Community Encounter Ingestion}

\label{element:7afe8ae69ed7491192d884576ca0dc15}
\index{[Feature] Community Encounter Ingestion (II.1.3.1.5.2.2)}
{\small Capability / Feature · Canonical · current}\par

Ingest home visit notes, community nursing observations, and social support assessments.



Community encounter ingestion brings home-visit and community assessments into the longitudinal information context. It supports continuity across organisational boundaries while preserving source attribution and the difference between the documentation received and care performed by the platform.


{\footnotesize \href{../../../../continuum/strategy/capability/feature/7afe8ae69ed7491192d884576ca0dc15.md}{Canonical record: 7afe8ae69ed7491192d884576ca0dc15}; full qualifications and provenance.}\par
\paragraph*{II.1.3.1.5.3 [CT3] Outreach Care Continuity}
\phantomsection
\label{navigation:II.1.3.1.5.3}
\addcontentsline{toc}{paragraph}{II.1.3.1.5.3 [CT3] Outreach Care Continuity}

\label{element:8526d971b07f40de84ef05b0d192089c}
\index{[CT3] Outreach Care Continuity (II.1.3.1.5.3)}
{\small Capability · Canonical · current}\par

Definition unresolved in the baseline. No definition has been synthesised from the name.



\medskip\hrule\medskip

{\footnotesize \href{../../../../continuum/strategy/capability/8526d971b07f40de84ef05b0d192089c.md}{Canonical record: 8526d971b07f40de84ef05b0d192089c}; full qualifications and provenance.}\par
\subparagraph*{II.1.3.1.5.3.1 [Feature] Outreach Encounter Outcome Ingestion}
\phantomsection
\label{navigation:II.1.3.1.5.3.1}
\addcontentsline{toc}{subparagraph}{II.1.3.1.5.3.1 [Feature] Outreach Encounter Outcome Ingestion}

\label{element:464ee3ad8d3648d992c6e51e3632d593}
\index{[Feature] Outreach Encounter Outcome Ingestion (II.1.3.1.5.3.1)}
{\small Capability / Feature · Canonical · current}\par

Ingest and reconcile clinical documentation gathered during rural and remote outreach visits.



Outreach outcome ingestion brings documentation collected during remote visits back into the governed record. It reconnects previously separated activity with the regional context without erasing when, where and by whom the source information was established.


{\footnotesize \href{../../../../continuum/strategy/capability/feature/464ee3ad8d3648d992c6e51e3632d593.md}{Canonical record: 464ee3ad8d3648d992c6e51e3632d593}; full qualifications and provenance.}\par
\subparagraph*{II.1.3.1.5.3.2 [Feature] Outreach Visit Context Provision}
\phantomsection
\label{navigation:II.1.3.1.5.3.2}
\addcontentsline{toc}{subparagraph}{II.1.3.1.5.3.2 [Feature] Outreach Visit Context Provision}

\label{element:44049d029a41437685272dfe94db4a2a}
\index{[Feature] Outreach Visit Context Provision (II.1.3.1.5.3.2)}
{\small Capability / Feature · Canonical · current}\par

Package and deliver portable clinical summaries and offline-capable care summaries for remote outreach clinics.



Outreach context provision prepares portable summaries for care beyond the usual connected environment. Their usefulness depends on retaining clinical and temporal context; an offline-capable copy is not automatically current authoritative state or a new origin of clinical facts.


{\footnotesize \href{../../../../continuum/strategy/capability/feature/44049d029a41437685272dfe94db4a2a.md}{Canonical record: 44049d029a41437685272dfe94db4a2a}; full qualifications and provenance.}\par
\subsubsection*{II.1.3.1.6 [CT2] Client-Centred Administration}
\phantomsection
\label{navigation:II.1.3.1.6}
\addcontentsline{toc}{subsubsection}{II.1.3.1.6 [CT2] Client-Centred Administration}

{\small Historical predecessor: R.2.55 \hyperref[element:232b3fb2ff204c9b9ff7a60ff6a7a32d]{[CT2] Client-Centred Administration}; Canonical ID: \texttt{\seqsplit{232b3fb2ff204c9b9ff7a60ff6a7a32d}}. Superseded historical identity — direct correspondence · superseded.}\par
\label{element:8752b7a937754b2faa296300c6a48911}
\index{[CT2] Client-Centred Administration (II.1.3.1.6)}
{\small Capability · Canonical · current}\par

Client-Centred Administration receives demographic and registration changes and makes verified updates available to participating clinical systems. It connects administrative events from external PAS environments with the client context used across care. The enabling responsibility concerns integration of changes, rather than independently deciding the authoritative demographic truth whenever a source sends an update.



\textbf{Overview}



Client-Centred Administration receives demographic and registration changes and makes verified updates available to participating clinical systems. It connects administrative events from external PAS environments with the client context used across care. The enabling responsibility concerns integration of changes, rather than independently deciding the authoritative demographic truth whenever a source sends an update.



\textbf{Intent}



Client-Centred Administration receives demographic and registration changes and makes verified updates available to participating clinical systems. It connects administrative events from external PAS environments with the client context used across care. The enabling responsibility concerns integration of changes, rather than independently deciding the authoritative demographic truth whenever a source sends an update.



\textbf{Boundary}



Information administration does not confer authority over the represented person, organisation, service, place, device or clinical activity.



\medskip\hrule\medskip

{\footnotesize \href{../../../../continuum/strategy/capability/8752b7a937754b2faa296300c6a48911.md}{Canonical record: 8752b7a937754b2faa296300c6a48911}; full qualifications and provenance.}\par
\paragraph*{II.1.3.1.6.1 [CT3] Client Intake \& Demographic Update}
\phantomsection
\label{navigation:II.1.3.1.6.1}
\addcontentsline{toc}{paragraph}{II.1.3.1.6.1 [CT3] Client Intake \& Demographic Update}

\label{element:f51f8d6b7556424987c76d9a7465c67c}
\index{[CT3] Client Intake \& Demographic Update (II.1.3.1.6.1)}
{\small Capability · Proposed · current}\par

Administer the established client intake \& demographic update information and coordination responsibility while preserving its source authority.



\textbf{Overview}



Administer the established client intake \& demographic update information and coordination responsibility while preserving its source authority. The reviewed hierarchy assigns this enduring ability to its stated business-enabling context.



\textbf{Intent}



Administer the established client intake \& demographic update information and coordination responsibility while preserving its source authority.



\textbf{Boundary}



Information administration and coordination preserve the authority of the responsible source, clinical actor and governed domain; they confer no additional real-world management authority.


{\footnotesize \href{../../../../continuum/strategy/capability/f51f8d6b7556424987c76d9a7465c67c.md}{Canonical record: f51f8d6b7556424987c76d9a7465c67c}; full qualifications and provenance.}\par
\subparagraph*{II.1.3.1.6.1.1 [Feature] Demographic Change Distribution}
\phantomsection
\label{navigation:II.1.3.1.6.1.1}
\addcontentsline{toc}{subparagraph}{II.1.3.1.6.1.1 [Feature] Demographic Change Distribution}

\label{element:bfd2ce654a79404dbef2694d30fb39fc}
\index{[Feature] Demographic Change Distribution (II.1.3.1.6.1.1)}
{\small Capability / Feature · Canonical · current}\par

Fan out verified client demographic updates to downstream clinical subsystems.



Change distribution propagates verified demographic updates to participating clinical systems. It supports consistent subject context across the network while keeping delivery to each destination distinguishable from that destination's eventual incorporation or business response.


{\footnotesize \href{../../../../continuum/strategy/capability/feature/bfd2ce654a79404dbef2694d30fb39fc.md}{Canonical record: bfd2ce654a79404dbef2694d30fb39fc}; full qualifications and provenance.}\par
\subparagraph*{II.1.3.1.6.1.2 [Feature] Demographic Update Ingestion}
\phantomsection
\label{navigation:II.1.3.1.6.1.2}
\addcontentsline{toc}{subparagraph}{II.1.3.1.6.1.2 [Feature] Demographic Update Ingestion}

\label{element:594ccbd95f5841e4b874fd1bade34a5d}
\index{[Feature] Demographic Update Ingestion (II.1.3.1.6.1.2)}
{\small Capability / Feature · Canonical · current}\par

Ingest and validate patient admission and registration updates from external PAS systems.



Demographic ingestion receives administrative changes from upstream admission and registration systems. It makes externally established updates available for governed processing without treating each incoming assertion as an independent master-person adjudication by Harmonia.


{\footnotesize \href{../../../../continuum/strategy/capability/feature/594ccbd95f5841e4b874fd1bade34a5d.md}{Canonical record: 594ccbd95f5841e4b874fd1bade34a5d}; full qualifications and provenance.}\par
\paragraph*{II.1.3.1.6.2 [CT3] Screening and Immunisation Enablement}
\phantomsection
\label{navigation:II.1.3.1.6.2}
\addcontentsline{toc}{paragraph}{II.1.3.1.6.2 [CT3] Screening and Immunisation Enablement}

\label{element:3e9016eaebae4ed0b67a5cacba2652f5}
\index{[CT3] Screening and Immunisation Enablement (II.1.3.1.6.2)}
{\small Capability · Canonical · current}\par

Definition unresolved in the baseline. No definition has been synthesised from the name.



\medskip\hrule\medskip

{\footnotesize \href{../../../../continuum/strategy/capability/3e9016eaebae4ed0b67a5cacba2652f5.md}{Canonical record: 3e9016eaebae4ed0b67a5cacba2652f5}; full qualifications and provenance.}\par
\subparagraph*{II.1.3.1.6.2.1 [Feature] Immunisation Event Submission}
\phantomsection
\label{navigation:II.1.3.1.6.2.1}
\addcontentsline{toc}{subparagraph}{II.1.3.1.6.2.1 [Feature] Immunisation Event Submission}

\label{element:47bab2681bf943899e83fab8c9941676}
\index{[Feature] Immunisation Event Submission (II.1.3.1.6.2.1)}
{\small Capability / Feature · Canonical · current}\par

Ingest point-of-care immunisation administration records and syndicate to national immunization registers.



Immunisation submission connects a source-recorded administration event with the national register context. It supports continuity and reporting without performing the immunisation or making successful transmission equivalent to every downstream registry outcome.


{\footnotesize \href{../../../../continuum/strategy/capability/feature/47bab2681bf943899e83fab8c9941676.md}{Canonical record: 47bab2681bf943899e83fab8c9941676}; full qualifications and provenance.}\par
\subparagraph*{II.1.3.1.6.2.2 [Feature] Screening Recall Notification Distribution}
\phantomsection
\label{navigation:II.1.3.1.6.2.2}
\addcontentsline{toc}{subparagraph}{II.1.3.1.6.2.2 [Feature] Screening Recall Notification Distribution}

\label{element:7ec85487625349a7b3292585fadb20c7}
\index{[Feature] Screening Recall Notification Distribution (II.1.3.1.6.2.2)}
{\small Capability / Feature · Canonical · current}\par

Distribute clinical screening reminders and recall notices to designated healthcare providers.



Screening recall distribution conveys established reminders to the designated care providers. Its notification responsibility is expressible through approved event and workflow composition; the catalogue does not invent a specialised Function or transfer the clinical decision to recall the patient.


{\footnotesize \href{../../../../continuum/strategy/capability/feature/7ec85487625349a7b3292585fadb20c7.md}{Canonical record: 7ec85487625349a7b3292585fadb20c7}; full qualifications and provenance.}\par
\subsubsection*{II.1.3.1.7 [CT2] Clinical Charging \& Claims Administration}
\phantomsection
\label{navigation:II.1.3.1.7}
\addcontentsline{toc}{subsubsection}{II.1.3.1.7 [CT2] Clinical Charging \& Claims Administration}

{\small Historical predecessor: R.2.56 \hyperref[element:44433108cc8f4f5f9d7581f9e5b29eb5]{[CT2] Clinical Charging \& Claims Administration}; Canonical ID: \texttt{\seqsplit{44433108cc8f4f5f9d7581f9e5b29eb5}}. Superseded historical identity — direct correspondence · superseded.}\par
\label{element:af78ab7c89454a608cd4d34ec7703331}
\index{[CT2] Clinical Charging \& Claims Administration (II.1.3.1.7)}
{\small Capability · Canonical · current}\par

This capability supplies clinical activity summaries and charging records to billing contexts. It preserves a route from activity information to the systems responsible for financial processing without treating the export as a completed bill or claim. No Feature decomposition or comprehensive revenue-cycle responsibility is established.



\textbf{Overview}



This capability supplies clinical activity summaries and charging records to billing contexts. It preserves a route from activity information to the systems responsible for financial processing without treating the export as a completed bill or claim. No Feature decomposition or comprehensive revenue-cycle responsibility is established.



\textbf{Intent}



This capability supplies clinical activity summaries and charging records to billing contexts. It preserves a route from activity information to the systems responsible for financial processing without treating the export as a completed bill or claim. No Feature decomposition or comprehensive revenue-cycle responsibility is established.



\textbf{Boundary}



Information administration does not confer authority over the represented person, organisation, service, place, device or clinical activity.



\medskip\hrule\medskip

{\footnotesize \href{../../../../continuum/strategy/capability/af78ab7c89454a608cd4d34ec7703331.md}{Canonical record: af78ab7c89454a608cd4d34ec7703331}; full qualifications and provenance.}\par
\subsubsection*{II.1.3.1.8 [CT2] Clinical Communication Administration}
\phantomsection
\label{navigation:II.1.3.1.8}
\addcontentsline{toc}{subsubsection}{II.1.3.1.8 [CT2] Clinical Communication Administration}

{\small Historical predecessor: R.2.57 \hyperref[element:bd4d6f4a7f03402ba7857383c7be7a4c]{[CT2] Clinical Communication Administration}; Canonical ID: \texttt{\seqsplit{bd4d6f4a7f03402ba7857383c7be7a4c}}. Superseded historical identity — direct correspondence · superseded.}\par
\label{element:b38f2c2811724734b5887b6dbaaabbfc}
\index{[CT2] Clinical Communication Administration (II.1.3.1.8)}
{\small Capability · Canonical · current}\par

This capability manages secure clinical communication transactions, delivery and receipt evidence, and the required enduring audit trail. Verified provider and team endpoints support appropriate addressing while transaction correlation makes the communication history intelligible. The responsibility preserves separate meanings for transport delivery, recipient acknowledgement and clinical or business disposition.



\textbf{Overview}



This capability manages secure clinical communication transactions, delivery and receipt evidence, and the required enduring audit trail. Verified provider and team endpoints support appropriate addressing while transaction correlation makes the communication history intelligible. The responsibility preserves separate meanings for transport delivery, recipient acknowledgement and clinical or business disposition.



\textbf{Intent}



This capability manages secure clinical communication transactions, delivery and receipt evidence, and the required enduring audit trail. Verified provider and team endpoints support appropriate addressing while transaction correlation makes the communication history intelligible. The responsibility preserves separate meanings for transport delivery, recipient acknowledgement and clinical or business disposition.



\textbf{Boundary}



Information administration does not confer authority over the represented person, organisation, service, place, device or clinical activity.



\medskip\hrule\medskip

{\footnotesize \href{../../../../continuum/strategy/capability/b38f2c2811724734b5887b6dbaaabbfc.md}{Canonical record: b38f2c2811724734b5887b6dbaaabbfc}; full qualifications and provenance.}\par
\paragraph*{II.1.3.1.8.1 [CT3] Secure Clinical Communication}
\phantomsection
\label{navigation:II.1.3.1.8.1}
\addcontentsline{toc}{paragraph}{II.1.3.1.8.1 [CT3] Secure Clinical Communication}

\label{element:eddb86c7922145439b9dcf0259f4aed5}
\index{[CT3] Secure Clinical Communication (II.1.3.1.8.1)}
{\small Capability · Proposed · current}\par

Administer the established secure clinical communication information and coordination responsibility while preserving its source authority.



\textbf{Overview}



Administer the established secure clinical communication information and coordination responsibility while preserving its source authority. The reviewed hierarchy assigns this enduring ability to its stated business-enabling context.



\textbf{Intent}



Administer the established secure clinical communication information and coordination responsibility while preserving its source authority.



\textbf{Boundary}



Information administration and coordination preserve the authority of the responsible source, clinical actor and governed domain; they confer no additional real-world management authority.


{\footnotesize \href{../../../../continuum/strategy/capability/eddb86c7922145439b9dcf0259f4aed5.md}{Canonical record: eddb86c7922145439b9dcf0259f4aed5}; full qualifications and provenance.}\par
\subparagraph*{II.1.3.1.8.1.1 [Feature] Clinical Communication Distribution}
\phantomsection
\label{navigation:II.1.3.1.8.1.1}
\addcontentsline{toc}{subparagraph}{II.1.3.1.8.1.1 [Feature] Clinical Communication Distribution}

\label{element:8527f36d5d5a460b92403f9d84f36490}
\index{[Feature] Clinical Communication Distribution (II.1.3.1.8.1.1)}
{\small Capability / Feature · Canonical · current}\par

Distribute secure clinical communications between verified healthcare providers and designated care team endpoints.



Clinical communication distribution connects verified providers and designated team endpoints through governed addressed delivery. The responsibility concerns reaching the intended participant, rather than assuming that delivered information has been read, incorporated or acted upon clinically.


{\footnotesize \href{../../../../continuum/strategy/capability/feature/8527f36d5d5a460b92403f9d84f36490.md}{Canonical record: 8527f36d5d5a460b92403f9d84f36490}; full qualifications and provenance.}\par
\subparagraph*{II.1.3.1.8.1.2 [Feature] Communication Audit Logging}
\phantomsection
\label{navigation:II.1.3.1.8.1.2}
\addcontentsline{toc}{subparagraph}{II.1.3.1.8.1.2 [Feature] Communication Audit Logging}

\label{element:b9df49b2f221486392f5cc393369d888}
\index{[Feature] Communication Audit Logging (II.1.3.1.8.1.2)}
{\small Capability / Feature · Canonical · current}\par

Maintain an enduring, non-PHI, tamper-evident log of all clinical communication transactions.



Communication audit evidence preserves the attributable transaction history required for scrutiny of clinical communication. It records meaningful communication facts under the zero-PHI constraints rather than turning the audit trail into a copy of message content or a permanent diagnostic trace.


{\footnotesize \href{../../../../continuum/strategy/capability/feature/b9df49b2f221486392f5cc393369d888.md}{Canonical record: b9df49b2f221486392f5cc393369d888}; full qualifications and provenance.}\par
\subparagraph*{II.1.3.1.8.1.3 [Feature] Delivery Acknowledgement Tracking}
\phantomsection
\label{navigation:II.1.3.1.8.1.3}
\addcontentsline{toc}{subparagraph}{II.1.3.1.8.1.3 [Feature] Delivery Acknowledgement Tracking}

\label{element:6884eb06bb794c5b8513b9b63d69af7e}
\index{[Feature] Delivery Acknowledgement Tracking (II.1.3.1.8.1.3)}
{\small Capability / Feature · Canonical · current}\par

Correlate recipient delivery and read receipts with outbound clinical communications, distinguishing transport delivery from recipient acknowledgement. Receipt/read evidence does not by itself establish business acceptance, clinical review, incorporation, approval, finality or originating authority.



Acknowledgement tracking associates delivery and read evidence with the communication transaction. Each receipt retains its contractual meaning; neither delivery nor reading alone proves business acceptance, clinical review, approval, finality or originating authority.


{\footnotesize \href{../../../../continuum/strategy/capability/feature/6884eb06bb794c5b8513b9b63d69af7e.md}{Canonical record: 6884eb06bb794c5b8513b9b63d69af7e}; full qualifications and provenance.}\par
\subsubsection*{II.1.3.1.9 [CT2] Clinical Record Administration}
\phantomsection
\label{navigation:II.1.3.1.9}
\addcontentsline{toc}{subsubsection}{II.1.3.1.9 [CT2] Clinical Record Administration}

{\small Historical predecessor: R.2.59 \hyperref[element:be3f2e695eda403fbce8118daf83d386]{[CT2] Clinical Record Administration}; Canonical ID: \texttt{\seqsplit{be3f2e695eda403fbce8118daf83d386}}. Superseded historical identity — direct correspondence · superseded.}\par
\label{element:ce9e7e7b22054cefb707a650d5be7219}
\index{[CT2] Clinical Record Administration (II.1.3.1.9)}
{\small Capability · Canonical · current}\par

Clinical Record Administration manages clinical documents as governed information with authorship, classification, versions and supersession history. It distinguishes a new document from an amendment or withdrawal and maintains metadata for discovery. Its purpose is to preserve an intelligible documentary record across change, rather than merely accumulate files.



\textbf{Overview}



Clinical Record Administration manages clinical documents as governed information with authorship, classification, versions and supersession history. It distinguishes a new document from an amendment or withdrawal and maintains metadata for discovery. Its purpose is to preserve an intelligible documentary record across change, rather than merely accumulate files.



\textbf{Intent}



Clinical Record Administration manages clinical documents as governed information with authorship, classification, versions and supersession history. It distinguishes a new document from an amendment or withdrawal and maintains metadata for discovery. Its purpose is to preserve an intelligible documentary record across change, rather than merely accumulate files.



\textbf{Boundary}



Information administration does not confer authority over the represented person, organisation, service, place, device or clinical activity.



\medskip\hrule\medskip

{\footnotesize \href{../../../../continuum/strategy/capability/ce9e7e7b22054cefb707a650d5be7219.md}{Canonical record: ce9e7e7b22054cefb707a650d5be7219}; full qualifications and provenance.}\par
\paragraph*{II.1.3.1.9.1 [CT3] Clinical Documentation Lifecycle Administration}
\phantomsection
\label{navigation:II.1.3.1.9.1}
\addcontentsline{toc}{paragraph}{II.1.3.1.9.1 [CT3] Clinical Documentation Lifecycle Administration}

{\small Historical predecessor: R.2.58 \hyperref[element:5b1636373a8d4e788b3ddbc1659a47ef]{[CT3] Clinical Documentation Lifecycle Administration}; Canonical ID: \texttt{\seqsplit{5b1636373a8d4e788b3ddbc1659a47ef}}. Superseded historical identity — direct correspondence · superseded.}\par
\label{element:df683100067b47e785037c27cfd17a3c}
\index{[CT3] Clinical Documentation Lifecycle Administration (II.1.3.1.9.1)}
{\small Capability · Proposed · current}\par

Administer ingestion, indexing, versioning, supersession and governed availability of clinical documentation.



\textbf{Overview}



Administer ingestion, indexing, versioning, supersession and governed availability of clinical documentation. The reviewed hierarchy assigns this enduring ability to its stated business-enabling context.



\textbf{Intent}



Administer ingestion, indexing, versioning, supersession and governed availability of clinical documentation.



\textbf{Boundary}



Information administration and coordination preserve the authority of the responsible source, clinical actor and governed domain; they confer no additional real-world management authority.


{\footnotesize \href{../../../../continuum/strategy/capability/df683100067b47e785037c27cfd17a3c.md}{Canonical record: df683100067b47e785037c27cfd17a3c}; full qualifications and provenance.}\par
\subparagraph*{II.1.3.1.9.1.1 [Feature] Clinical Document Ingestion}
\phantomsection
\label{navigation:II.1.3.1.9.1.1}
\addcontentsline{toc}{subparagraph}{II.1.3.1.9.1.1 [Feature] Clinical Document Ingestion}

\label{element:54c79edadc7a478fb594ff4eb47f8368}
\index{[Feature] Clinical Document Ingestion (II.1.3.1.9.1.1)}
{\small Capability / Feature · Canonical · current}\par

Ingest structured and unstructured clinical documents (discharge summaries, consult notes, letters).



Clinical document ingestion receives authored documents into governed information management. It supports both structured and unstructured contributions while preserving that receipt establishes a managed contribution, not independent clinical verification of every assertion in the document.


{\footnotesize \href{../../../../continuum/strategy/capability/feature/54c79edadc7a478fb594ff4eb47f8368.md}{Canonical record: 54c79edadc7a478fb594ff4eb47f8368}; full qualifications and provenance.}\par
\subparagraph*{II.1.3.1.9.1.2 [Feature] Document Metadata Indexing}
\phantomsection
\label{navigation:II.1.3.1.9.1.2}
\addcontentsline{toc}{subparagraph}{II.1.3.1.9.1.2 [Feature] Document Metadata Indexing}

\label{element:0c85a269e90a4329b83d934f9061cd49}
\index{[Feature] Document Metadata Indexing (II.1.3.1.9.1.2)}
{\small Capability / Feature · Canonical · current}\par

Index clinical document metadata (author, specialty, encounter, date, confidentiality) for rapid discovery.



Metadata indexing makes documents discoverable by clinically and operationally relevant attributes. It describes content for search without making the index the document itself or granting access merely because a matching record can be located.


{\footnotesize \href{../../../../continuum/strategy/capability/feature/0c85a269e90a4329b83d934f9061cd49.md}{Canonical record: 0c85a269e90a4329b83d934f9061cd49}; full qualifications and provenance.}\par
\subparagraph*{II.1.3.1.9.1.3 [Feature] Document Versioning \& Supersession}
\phantomsection
\label{navigation:II.1.3.1.9.1.3}
\addcontentsline{toc}{subparagraph}{II.1.3.1.9.1.3 [Feature] Document Versioning \& Supersession}

\label{element:c8e2e623dd4e4c3b8de6c3bfb8c2a3ee}
\index{[Feature] Document Versioning \& Supersession (II.1.3.1.9.1.3)}
{\small Capability / Feature · Canonical · current}\par

Manage document classification, author attribution, amendment attribution, versioning, and supersession or withdrawal without destructive data loss.



Versioning and supersession make document change understandable without destructive loss of the previous account. Author and amendment attribution distinguish who established each contribution; withdrawal or correction changes current standing without retrospectively erasing what Harmonia previously received or managed.


{\footnotesize \href{../../../../continuum/strategy/capability/feature/c8e2e623dd4e4c3b8de6c3bfb8c2a3ee.md}{Canonical record: c8e2e623dd4e4c3b8de6c3bfb8c2a3ee}; full qualifications and provenance.}\par
\subsubsection*{II.1.3.1.10 [CT2] Diagnostic Administration}
\phantomsection
\label{navigation:II.1.3.1.10}
\addcontentsline{toc}{subsubsection}{II.1.3.1.10 [CT2] Diagnostic Administration}

{\small Historical predecessor: R.2.60 \hyperref[element:2b1a27d5d9944e05b4cf0fcaad4ca433]{[CT2] Diagnostic Administration}; Canonical ID: \texttt{\seqsplit{2b1a27d5d9944e05b4cf0fcaad4ca433}}. Superseded historical identity — direct correspondence · superseded.}\par
\label{element:3cc720ba2e8b4d7a94bb349391d61895}
\index{[CT2] Diagnostic Administration (II.1.3.1.10)}
{\small Capability · Canonical · current}\par

Diagnostic Administration correlates requests with specimens and encounters and integrates diagnostic reports into their request context. It connects the laboratory or imaging contribution with the participants expecting an outcome, preserving report identity and standing as updates occur. Its enabling role differs from performing the diagnostic procedure or professionally interpreting evidence.



\textbf{Overview}



Diagnostic Administration correlates requests with specimens and encounters and integrates diagnostic reports into their request context. It connects the laboratory or imaging contribution with the participants expecting an outcome, preserving report identity and standing as updates occur. Its enabling role differs from performing the diagnostic procedure or professionally interpreting evidence.



\textbf{Intent}



Diagnostic Administration correlates requests with specimens and encounters and integrates diagnostic reports into their request context. It connects the laboratory or imaging contribution with the participants expecting an outcome, preserving report identity and standing as updates occur. Its enabling role differs from performing the diagnostic procedure or professionally interpreting evidence.



\textbf{Boundary}



Information administration does not confer authority over the represented person, organisation, service, place, device or clinical activity.



\medskip\hrule\medskip

{\footnotesize \href{../../../../continuum/strategy/capability/3cc720ba2e8b4d7a94bb349391d61895.md}{Canonical record: 3cc720ba2e8b4d7a94bb349391d61895}; full qualifications and provenance.}\par
\paragraph*{II.1.3.1.10.1 [CT3] Diagnostic Service Integration}
\phantomsection
\label{navigation:II.1.3.1.10.1}
\addcontentsline{toc}{paragraph}{II.1.3.1.10.1 [CT3] Diagnostic Service Integration}

\label{element:9ded2f3a28d24017bb921a0cdcf65a2b}
\index{[CT3] Diagnostic Service Integration (II.1.3.1.10.1)}
{\small Capability · Canonical · current}\par

Definition unresolved in the baseline. No definition has been synthesised from the name.



\medskip\hrule\medskip

{\footnotesize \href{../../../../continuum/strategy/capability/9ded2f3a28d24017bb921a0cdcf65a2b.md}{Canonical record: 9ded2f3a28d24017bb921a0cdcf65a2b}; full qualifications and provenance.}\par
\subparagraph*{II.1.3.1.10.1.1 [Feature] Diagnostic History Consolidation}
\phantomsection
\label{navigation:II.1.3.1.10.1.1}
\addcontentsline{toc}{subparagraph}{II.1.3.1.10.1.1 [Feature] Diagnostic History Consolidation}

\label{element:0b3fa0b7374842ff8ce199191a3479bd}
\index{[Feature] Diagnostic History Consolidation (II.1.3.1.10.1.1)}
{\small Capability / Feature · Canonical · current}\par

Consolidate historical diagnostic trends (e.g., cumulative laboratory results) across multiple regional providers into an integrated view.



Diagnostic history consolidation makes prior evidence visible across regional providers. A combined historical view can support comparison over time without itself establishing a diagnostic interpretation, reconciling conflicting clinical assertions or replacing each report's source standing.


{\footnotesize \href{../../../../continuum/strategy/capability/feature/0b3fa0b7374842ff8ce199191a3479bd.md}{Canonical record: 0b3fa0b7374842ff8ce199191a3479bd}; full qualifications and provenance.}\par
\subparagraph*{II.1.3.1.10.1.2 [Feature] Laboratory \& Imaging Result Distribution}
\phantomsection
\label{navigation:II.1.3.1.10.1.2}
\addcontentsline{toc}{subparagraph}{II.1.3.1.10.1.2 [Feature] Laboratory \& Imaging Result Distribution}

\label{element:93d42becda8f4850b73ff2c93bb7e8f3}
\index{[Feature] Laboratory \& Imaging Result Distribution (II.1.3.1.10.1.2)}
{\small Capability / Feature · Canonical · current}\par

Reliably distribute structured diagnostic reports (pathology, radiology) to ordering clinicians and designated secondary recipients.



Result distribution connects established diagnostic reports with ordering and designated secondary recipients. It supports continuity of diagnostic information across care settings while retaining the distinction between delivery, recipient review and the clinical decision based on the result.


{\footnotesize \href{../../../../continuum/strategy/capability/feature/93d42becda8f4850b73ff2c93bb7e8f3.md}{Canonical record: 93d42becda8f4850b73ff2c93bb7e8f3}; full qualifications and provenance.}\par
\paragraph*{II.1.3.1.10.2 [CT3] Diagnostic Workflow Management}
\phantomsection
\label{navigation:II.1.3.1.10.2}
\addcontentsline{toc}{paragraph}{II.1.3.1.10.2 [CT3] Diagnostic Workflow Management}

\label{element:270ed0e32b294458aa52cf57df7e55db}
\index{[CT3] Diagnostic Workflow Management (II.1.3.1.10.2)}
{\small Capability · Proposed · current}\par

Administer the established diagnostic workflow management information and coordination responsibility while preserving its source authority.



\textbf{Overview}



Administer the established diagnostic workflow management information and coordination responsibility while preserving its source authority. The reviewed hierarchy assigns this enduring ability to its stated business-enabling context.



\textbf{Intent}



Administer the established diagnostic workflow management information and coordination responsibility while preserving its source authority.



\textbf{Boundary}



Information administration and coordination preserve the authority of the responsible source, clinical actor and governed domain; they confer no additional real-world management authority.


{\footnotesize \href{../../../../continuum/strategy/capability/270ed0e32b294458aa52cf57df7e55db.md}{Canonical record: 270ed0e32b294458aa52cf57df7e55db}; full qualifications and provenance.}\par
\subparagraph*{II.1.3.1.10.2.1 [Feature] Diagnostic Preliminary \& Corrected Report Handling}
\phantomsection
\label{navigation:II.1.3.1.10.2.1}
\addcontentsline{toc}{subparagraph}{II.1.3.1.10.2.1 [Feature] Diagnostic Preliminary \& Corrected Report Handling}

\label{element:7449ee4dada7424ea6d39785258b4981}
\index{[Feature] Diagnostic Preliminary \& Corrected Report Handling (II.1.3.1.10.2.1)}
{\small Capability / Feature · Proposed · current}\par

Receive, qualify and distribute preliminary, corrected and addendum diagnostic information with its relationship to prior reports.



\textbf{Overview}



The finalised-report Feature remains unchanged. The broader documented diagnostic behaviour needs a separate frame for information whose preliminary or correction standing is material.



\textbf{Boundary}



Does not originate diagnostic interpretation or silently overwrite a previously received report.


{\footnotesize \href{../../../../continuum/strategy/capability/feature/7449ee4dada7424ea6d39785258b4981.md}{Canonical record: 7449ee4dada7424ea6d39785258b4981}; full qualifications and provenance.}\par
\subparagraph*{II.1.3.1.10.2.2 [Feature] Diagnostic Report Distribution}
\phantomsection
\label{navigation:II.1.3.1.10.2.2}
\addcontentsline{toc}{subparagraph}{II.1.3.1.10.2.2 [Feature] Diagnostic Report Distribution}

\label{element:3813b3b5175e461ba74805d8fa2aa648}
\index{[Feature] Diagnostic Report Distribution (II.1.3.1.10.2.2)}
{\small Capability / Feature · Canonical · current}\par

Distribute finalized laboratory, pathology, and imaging reports to ordering clinicians and designated care teams.



This Feature distributes finalised diagnostic reports to the ordering and designated care participants. Its narrower finalised-report scope is preserved even though broader Diagnostic Administration behaviour also handles preliminary, corrected and addendum information; those behaviours do not silently redefine the Feature.


{\footnotesize \href{../../../../continuum/strategy/capability/feature/3813b3b5175e461ba74805d8fa2aa648.md}{Canonical record: 3813b3b5175e461ba74805d8fa2aa648}; full qualifications and provenance.}\par
\subparagraph*{II.1.3.1.10.2.3 [Feature] Diagnostic Report Ingestion \& Binding}
\phantomsection
\label{navigation:II.1.3.1.10.2.3}
\addcontentsline{toc}{subparagraph}{II.1.3.1.10.2.3 [Feature] Diagnostic Report Ingestion \& Binding}

\label{element:a00823445a8a44c1a00f3ec93a7e0e6f}
\index{[Feature] Diagnostic Report Ingestion \& Binding (II.1.3.1.10.2.3)}
{\small Capability / Feature · Canonical · current}\par

Ingest finalized laboratory and radiology reports, verifying content integrity and binding them to the parent order.



Report ingestion binds a finalised diagnostic contribution to the order it answers and verifies content integrity. The report's professional standing comes from its diagnostic source; successful ingestion does not independently authorise its findings or establish a patient-level clinical conclusion.


{\footnotesize \href{../../../../continuum/strategy/capability/feature/a00823445a8a44c1a00f3ec93a7e0e6f.md}{Canonical record: a00823445a8a44c1a00f3ec93a7e0e6f}; full qualifications and provenance.}\par
\subparagraph*{II.1.3.1.10.2.4 [Feature] Diagnostic Request Correlation}
\phantomsection
\label{navigation:II.1.3.1.10.2.4}
\addcontentsline{toc}{subparagraph}{II.1.3.1.10.2.4 [Feature] Diagnostic Request Correlation}

\label{element:30cd73b8d7d941648affbb379196f532}
\index{[Feature] Diagnostic Request Correlation (II.1.3.1.10.2.4)}
{\small Capability / Feature · Canonical · current}\par

Correlate diagnostic test orders with collected clinical specimens and patient encounters.



Diagnostic request correlation connects the order with its specimen and encounter context. It helps distinguish evidence belonging to the intended fulfilment from superficially similar results, without treating association as proof of specimen suitability or diagnostic interpretation.


{\footnotesize \href{../../../../continuum/strategy/capability/feature/30cd73b8d7d941648affbb379196f532.md}{Canonical record: 30cd73b8d7d941648affbb379196f532}; full qualifications and provenance.}\par
\subsubsection*{II.1.3.1.11 [CT2] Episode \& Encounter Administration}
\phantomsection
\label{navigation:II.1.3.1.11}
\addcontentsline{toc}{subsubsection}{II.1.3.1.11 [CT2] Episode \& Encounter Administration}

{\small Historical predecessor: R.2.61 \hyperref[element:2778f988e8d44e3091749542c39d38a4]{[CT2] Episode \& Encounter Administration}; Canonical ID: \texttt{\seqsplit{2778f988e8d44e3091749542c39d38a4}}. Superseded historical identity — direct correspondence · superseded.}\par
\label{element:61d7067aee0540259b3daca4287d66d2}
\index{[CT2] Episode \& Encounter Administration (II.1.3.1.11)}
{\small Capability · Canonical · current}\par

This capability maintains identifiable encounter contexts and their governed progression, including association with care teams and locations. It enables information from an admission, visit or community contact to be understood within the relevant episode of activity. Encounter progression is an administrative and operational account, separate from the clinical acts that occur during it.



\textbf{Overview}



This capability maintains identifiable encounter contexts and their governed progression, including association with care teams and locations. It enables information from an admission, visit or community contact to be understood within the relevant episode of activity. Encounter progression is an administrative and operational account, separate from the clinical acts that occur during it.



\textbf{Intent}



This capability maintains identifiable encounter contexts and their governed progression, including association with care teams and locations. It enables information from an admission, visit or community contact to be understood within the relevant episode of activity. Encounter progression is an administrative and operational account, separate from the clinical acts that occur during it.



\textbf{Boundary}



Information administration does not confer authority over the represented person, organisation, service, place, device or clinical activity.



\medskip\hrule\medskip

{\footnotesize \href{../../../../continuum/strategy/capability/61d7067aee0540259b3daca4287d66d2.md}{Canonical record: 61d7067aee0540259b3daca4287d66d2}; full qualifications and provenance.}\par
\paragraph*{II.1.3.1.11.1 [CT3] Acute Care Monitoring \& Event Coordination}
\phantomsection
\label{navigation:II.1.3.1.11.1}
\addcontentsline{toc}{paragraph}{II.1.3.1.11.1 [CT3] Acute Care Monitoring \& Event Coordination}

\label{element:ee17b26b8c874cc3af50d07a787ce488}
\index{[CT3] Acute Care Monitoring \& Event Coordination (II.1.3.1.11.1)}
{\small Capability · Canonical · current}\par

Definition unresolved in the baseline. No definition has been synthesised from the name.



\medskip\hrule\medskip

{\footnotesize \href{../../../../continuum/strategy/capability/ee17b26b8c874cc3af50d07a787ce488.md}{Canonical record: ee17b26b8c874cc3af50d07a787ce488}; full qualifications and provenance.}\par
\subparagraph*{II.1.3.1.11.1.1 [Feature] Acute Care Clinical Context Provision}
\phantomsection
\label{navigation:II.1.3.1.11.1.1}
\addcontentsline{toc}{subparagraph}{II.1.3.1.11.1.1 [Feature] Acute Care Clinical Context Provision}

\label{element:90de621f7ac1410eb6fc770cb07a9e03}
\index{[Feature] Acute Care Clinical Context Provision (II.1.3.1.11.1.1)}
{\small Capability / Feature · Canonical · current}\par

Provide consolidated longitudinal history and active medication lists to acute care clinicians at point of triage.



Acute context provision brings relevant history and active medications together for clinicians facing a current presentation. It supports rapid understanding without substituting the consolidated view for clinical assessment, triage or a decision about treatment.


{\footnotesize \href{../../../../continuum/strategy/capability/feature/90de621f7ac1410eb6fc770cb07a9e03.md}{Canonical record: 90de621f7ac1410eb6fc770cb07a9e03}; full qualifications and provenance.}\par
\subparagraph*{II.1.3.1.11.1.2 [Feature] Acute Clinical State Event Capture}
\phantomsection
\label{navigation:II.1.3.1.11.1.2}
\addcontentsline{toc}{subparagraph}{II.1.3.1.11.1.2 [Feature] Acute Clinical State Event Capture}

\label{element:566f6934c9b04627b0876f91bedd8f30}
\index{[Feature] Acute Clinical State Event Capture (II.1.3.1.11.1.2)}
{\small Capability / Feature · Canonical · current}\par

Ingest clinical deterioration notifications and acute clinical status updates.



Acute-state capture records deterioration notifications and current clinical status as externally established information. Approved reusable clinical-information composition is sufficient for this responsibility; routing a critical-result alert does not become a specialised capture Function merely because both concern urgent care.


{\footnotesize \href{../../../../continuum/strategy/capability/feature/566f6934c9b04627b0876f91bedd8f30.md}{Canonical record: 566f6934c9b04627b0876f91bedd8f30}; full qualifications and provenance.}\par
\paragraph*{II.1.3.1.11.2 [CT3] Encounter Lifecycle Management}
\phantomsection
\label{navigation:II.1.3.1.11.2}
\addcontentsline{toc}{paragraph}{II.1.3.1.11.2 [CT3] Encounter Lifecycle Management}

\label{element:2ebf1bb20ab34755b41567c156cc5928}
\index{[CT3] Encounter Lifecycle Management (II.1.3.1.11.2)}
{\small Capability · Proposed · current}\par

Administer the established encounter lifecycle management information and coordination responsibility while preserving its source authority.



\textbf{Overview}



Administer the established encounter lifecycle management information and coordination responsibility while preserving its source authority. The reviewed hierarchy assigns this enduring ability to its stated business-enabling context.



\textbf{Intent}



Administer the established encounter lifecycle management information and coordination responsibility while preserving its source authority.



\textbf{Boundary}



Information administration and coordination preserve the authority of the responsible source, clinical actor and governed domain; they confer no additional real-world management authority.


{\footnotesize \href{../../../../continuum/strategy/capability/2ebf1bb20ab34755b41567c156cc5928.md}{Canonical record: 2ebf1bb20ab34755b41567c156cc5928}; full qualifications and provenance.}\par
\subparagraph*{II.1.3.1.11.2.1 [Feature] Encounter Bed/Location Association}
\phantomsection
\label{navigation:II.1.3.1.11.2.1}
\addcontentsline{toc}{subparagraph}{II.1.3.1.11.2.1 [Feature] Encounter Bed/Location Association}

\label{element:47f0fca36aa4454eae95f4b151c7076f}
\index{[Feature] Encounter Bed/Location Association (II.1.3.1.11.2.1)}
{\small Capability / Feature · Canonical · current}\par

Dynamically maintain the association between an encounter, care team, and physical bed or ward location.



Encounter association keeps the relevant place and care-team context aligned with the episode as the patient moves. Location changes are meaningful operational facts, but they do not themselves establish a transfer of clinical responsibility.


{\footnotesize \href{../../../../continuum/strategy/capability/feature/47f0fca36aa4454eae95f4b151c7076f.md}{Canonical record: 47f0fca36aa4454eae95f4b151c7076f}; full qualifications and provenance.}\par
\subparagraph*{II.1.3.1.11.2.2 [Feature] Encounter Context Creation}
\phantomsection
\label{navigation:II.1.3.1.11.2.2}
\addcontentsline{toc}{subparagraph}{II.1.3.1.11.2.2 [Feature] Encounter Context Creation}

\label{element:817f1046d7e84699a2e2050b84285f95}
\index{[Feature] Encounter Context Creation (II.1.3.1.11.2.2)}
{\small Capability / Feature · Canonical · current}\par

Create and maintain authoritative encounter tracking contexts for hospital admissions, outpatient visits, and community contacts.



Encounter context creation supplies an identifiable frame for information concerning a particular admission, visit or contact. Its governed tracking authority concerns the encounter representation within Harmonia; it does not make Harmonia the originator of every clinical fact associated with that encounter.


{\footnotesize \href{../../../../continuum/strategy/capability/feature/817f1046d7e84699a2e2050b84285f95.md}{Canonical record: 817f1046d7e84699a2e2050b84285f95}; full qualifications and provenance.}\par
\subparagraph*{II.1.3.1.11.2.3 [Feature] Encounter History Retrieval}
\phantomsection
\label{navigation:II.1.3.1.11.2.3}
\addcontentsline{toc}{subparagraph}{II.1.3.1.11.2.3 [Feature] Encounter History Retrieval}

\label{element:4088e1408f554143a0f1469efd1b43a7}
\index{[Feature] Encounter History Retrieval (II.1.3.1.11.2.3)}
{\small Capability / Feature · Proposed · current}\par

Expose qualified episode, movement and encounter-history context.



\textbf{Overview}



Consumers need to interpret prior care-place and episode relationships, not just the latest encounter state; historical retrieval remains bounded by source attribution and access policy.



\textbf{Boundary}



A historical view does not retrospectively alter clinical events or acquire longitudinal-record ownership.


{\footnotesize \href{../../../../continuum/strategy/capability/feature/4088e1408f554143a0f1469efd1b43a7.md}{Canonical record: 4088e1408f554143a0f1469efd1b43a7}; full qualifications and provenance.}\par
\subparagraph*{II.1.3.1.11.2.4 [Feature] Encounter State Progression}
\phantomsection
\label{navigation:II.1.3.1.11.2.4}
\addcontentsline{toc}{subparagraph}{II.1.3.1.11.2.4 [Feature] Encounter State Progression}

\label{element:c049039bf57e4e4aa97aab5967196383}
\index{[Feature] Encounter State Progression (II.1.3.1.11.2.4)}
{\small Capability / Feature · Canonical · current}\par

Track encounter status transitions (PLANNED → ARRIVED → IN\_PROGRESS → ON\_LEAVE → DISCHARGED).



Encounter progression preserves the difference between planned, arrived, active, leave and discharged circumstances. These states make administrative activity interpretable without implying a clinical outcome. The source's ON\_LEAVE transition detail remains incomplete in the corresponding Process and is not repaired here.


{\footnotesize \href{../../../../continuum/strategy/capability/feature/c049039bf57e4e4aa97aab5967196383.md}{Canonical record: c049039bf57e4e4aa97aab5967196383}; full qualifications and provenance.}\par
\subsubsection*{II.1.3.1.12 [CT2] Funding \& Coverage Administration}
\phantomsection
\label{navigation:II.1.3.1.12}
\addcontentsline{toc}{subsubsection}{II.1.3.1.12 [CT2] Funding \& Coverage Administration}

{\small Historical predecessor: R.2.62 \hyperref[element:58beafcbcd084331be710eb00d27f6cf]{[CT2] Funding \& Coverage Administration}; Canonical ID: \texttt{\seqsplit{58beafcbcd084331be710eb00d27f6cf}}. Superseded historical identity — direct correspondence · superseded.}\par
\label{element:640432d413cf47928ea2b3d221f92313}
\index{[CT2] Funding \& Coverage Administration (II.1.3.1.12)}
{\small Capability · Canonical · current}\par

Funding \& Coverage Administration represents externally supplied insurance coverage and eligibility flags needed as care context. It allows those administrative conditions to be available without implying that Harmonia determines entitlements. The established Adjacent standing and absence of Features bound the current responsibility to this information contribution.



\textbf{Overview}



Funding \& Coverage Administration represents externally supplied insurance coverage and eligibility flags needed as care context. It allows those administrative conditions to be available without implying that Harmonia determines entitlements. The established Adjacent standing and absence of Features bound the current responsibility to this information contribution.



\textbf{Intent}



Funding \& Coverage Administration represents externally supplied insurance coverage and eligibility flags needed as care context. It allows those administrative conditions to be available without implying that Harmonia determines entitlements. The established Adjacent standing and absence of Features bound the current responsibility to this information contribution.



\textbf{Boundary}



Information administration does not confer authority over the represented person, organisation, service, place, device or clinical activity.



\medskip\hrule\medskip

{\footnotesize \href{../../../../continuum/strategy/capability/640432d413cf47928ea2b3d221f92313.md}{Canonical record: 640432d413cf47928ea2b3d221f92313}; full qualifications and provenance.}\par
\subsubsection*{II.1.3.1.13 [CT2] Medication Administration}
\phantomsection
\label{navigation:II.1.3.1.13}
\addcontentsline{toc}{subsubsection}{II.1.3.1.13 [CT2] Medication Administration}

{\small Historical predecessor: R.2.63 \hyperref[element:5a38bf2527a8404bac714121f2dc76a9]{[CT2] Medication Administration}; Canonical ID: \texttt{\seqsplit{5a38bf2527a8404bac714121f2dc76a9}}. Superseded historical identity — direct correspondence · superseded.}\par
\label{element:c907e2a95629435b9951941ad4e2ebe9}
\index{[CT2] Medication Administration (II.1.3.1.13)}
{\small Capability · Canonical · current}\par

In this enabling catalogue, Medication Administration integrates medication orders, dispensing evidence and point-of-care administration records. It preserves the distinct events through which intended therapy becomes supplied and, where reported, administered. This administrative integration supports a coherent medication account across systems without itself dispensing or administering a medicine.



\textbf{Overview}



In this enabling catalogue, Medication Administration integrates medication orders, dispensing evidence and point-of-care administration records. It preserves the distinct events through which intended therapy becomes supplied and, where reported, administered. This administrative integration supports a coherent medication account across systems without itself dispensing or administering a medicine.



\textbf{Intent}



In this enabling catalogue, Medication Administration integrates medication orders, dispensing evidence and point-of-care administration records. It preserves the distinct events through which intended therapy becomes supplied and, where reported, administered. This administrative integration supports a coherent medication account across systems without itself dispensing or administering a medicine.



\textbf{Boundary}



Information administration does not confer authority over the represented person, organisation, service, place, device or clinical activity.



\medskip\hrule\medskip

{\footnotesize \href{../../../../continuum/strategy/capability/c907e2a95629435b9951941ad4e2ebe9.md}{Canonical record: c907e2a95629435b9951941ad4e2ebe9}; full qualifications and provenance.}\par
\paragraph*{II.1.3.1.13.1 [CT3] Medication Order \& Dispense Management}
\phantomsection
\label{navigation:II.1.3.1.13.1}
\addcontentsline{toc}{paragraph}{II.1.3.1.13.1 [CT3] Medication Order \& Dispense Management}

\label{element:9874098228764b62b1f640ea749b742d}
\index{[CT3] Medication Order \& Dispense Management (II.1.3.1.13.1)}
{\small Capability · Proposed · current}\par

Administer the established medication order \& dispense management information and coordination responsibility while preserving its source authority.



\textbf{Overview}



Administer the established medication order \& dispense management information and coordination responsibility while preserving its source authority. The reviewed hierarchy assigns this enduring ability to its stated business-enabling context.



\textbf{Intent}



Administer the established medication order \& dispense management information and coordination responsibility while preserving its source authority.



\textbf{Boundary}



Information administration and coordination preserve the authority of the responsible source, clinical actor and governed domain; they confer no additional real-world management authority.


{\footnotesize \href{../../../../continuum/strategy/capability/9874098228764b62b1f640ea749b742d.md}{Canonical record: 9874098228764b62b1f640ea749b742d}; full qualifications and provenance.}\par
\subparagraph*{II.1.3.1.13.1.1 [Feature] Dispense Event Tracking}
\phantomsection
\label{navigation:II.1.3.1.13.1.1}
\addcontentsline{toc}{subparagraph}{II.1.3.1.13.1.1 [Feature] Dispense Event Tracking}

\label{element:430bf990d4194255b7ae09c2fe27367e}
\index{[Feature] Dispense Event Tracking (II.1.3.1.13.1.1)}
{\small Capability / Feature · Canonical · current}\par

Ingest medication dispense verification records from pharmacy management systems.



Dispense tracking receives the pharmacy's evidence that medication was supplied under the applicable context. It connects supply with the medication account without inferring that the patient received a dose or that Harmonia performed the professional dispensing activity.


{\footnotesize \href{../../../../continuum/strategy/capability/feature/430bf990d4194255b7ae09c2fe27367e.md}{Canonical record: 430bf990d4194255b7ae09c2fe27367e}; full qualifications and provenance.}\par
\subparagraph*{II.1.3.1.13.1.2 [Feature] Medication Administration Record (MAR) Ingestion}
\phantomsection
\label{navigation:II.1.3.1.13.1.2}
\addcontentsline{toc}{subparagraph}{II.1.3.1.13.1.2 [Feature] Medication Administration Record (MAR) Ingestion}

\label{element:779d9cb66bc84585a8dd1f456b0ee37c}
\index{[Feature] Medication Administration Record (MAR) Ingestion (II.1.3.1.13.1.2)}
{\small Capability / Feature · Canonical · current}\par

Ingest point-of-care medication administration events.



MAR ingestion captures source-reported point-of-care administration events. It preserves the crucial distinction between intended or supplied medication and an administration recorded as performed, without originating that bedside event or its clinical authority.


{\footnotesize \href{../../../../continuum/strategy/capability/feature/779d9cb66bc84585a8dd1f456b0ee37c.md}{Canonical record: 779d9cb66bc84585a8dd1f456b0ee37c}; full qualifications and provenance.}\par
\subparagraph*{II.1.3.1.13.1.3 [Feature] Medication Order Ingestion}
\phantomsection
\label{navigation:II.1.3.1.13.1.3}
\addcontentsline{toc}{subparagraph}{II.1.3.1.13.1.3 [Feature] Medication Order Ingestion}

\label{element:5247b47f181243bb975ae2cd70db65ba}
\index{[Feature] Medication Order Ingestion (II.1.3.1.13.1.3)}
{\small Capability / Feature · Canonical · current}\par

Ingest inpatient and discharge prescription orders.



Medication order ingestion preserves the intended therapy expressed by inpatient or discharge prescriptions. The received order is evidence of clinical intent, distinct from pharmacy dispensing and the patient's subsequent administration record.


{\footnotesize \href{../../../../continuum/strategy/capability/feature/5247b47f181243bb975ae2cd70db65ba.md}{Canonical record: 5247b47f181243bb975ae2cd70db65ba}; full qualifications and provenance.}\par
\paragraph*{II.1.3.1.13.2 [CT3] Medication Therapy Provision and Monitoring}
\phantomsection
\label{navigation:II.1.3.1.13.2}
\addcontentsline{toc}{paragraph}{II.1.3.1.13.2 [CT3] Medication Therapy Provision and Monitoring}

\label{element:d1821b4f6d5746f9a697dfe3b792a186}
\index{[CT3] Medication Therapy Provision and Monitoring (II.1.3.1.13.2)}
{\small Capability · Canonical · current}\par

Definition unresolved in the baseline. No definition has been synthesised from the name.



\medskip\hrule\medskip

{\footnotesize \href{../../../../continuum/strategy/capability/d1821b4f6d5746f9a697dfe3b792a186.md}{Canonical record: d1821b4f6d5746f9a697dfe3b792a186}; full qualifications and provenance.}\par
\subparagraph*{II.1.3.1.13.2.1 [Feature] Medication Administration Tracking}
\phantomsection
\label{navigation:II.1.3.1.13.2.1}
\addcontentsline{toc}{subparagraph}{II.1.3.1.13.2.1 [Feature] Medication Administration Tracking}

\label{element:6e01c72ef9c149e2ac231463810c2265}
\index{[Feature] Medication Administration Tracking (II.1.3.1.13.2.1)}
{\small Capability / Feature · Canonical · current}\par

Distribute medication therapy orders, dispensing updates, and administration records across clinical care boundaries.



Medication tracking conveys orders, dispensing changes and administration information across care boundaries while keeping their meanings separate. This supports a coherent therapy account without treating a supplied medicine as an administered dose or a distributed order as treatment performed.


{\footnotesize \href{../../../../continuum/strategy/capability/feature/6e01c72ef9c149e2ac231463810c2265.md}{Canonical record: 6e01c72ef9c149e2ac231463810c2265}; full qualifications and provenance.}\par
\subparagraph*{II.1.3.1.13.2.2 [Feature] Medication History Reconciliation Ingestion}
\phantomsection
\label{navigation:II.1.3.1.13.2.2}
\addcontentsline{toc}{subparagraph}{II.1.3.1.13.2.2 [Feature] Medication History Reconciliation Ingestion}

\label{element:561e1ef130e04f38b854665bb8fb9c77}
\index{[Feature] Medication History Reconciliation Ingestion (II.1.3.1.13.2.2)}
{\small Capability / Feature · Canonical · current}\par

Ingest reconciled medication histories compiled during clinical transitions of care.



This Feature receives medication histories already reconciled by clinical participants during care transitions. The platform preserves and connects that account; ingestion is distinct from performing the clinical reconciliation or independently deciding which regimen is correct.


{\footnotesize \href{../../../../continuum/strategy/capability/feature/561e1ef130e04f38b854665bb8fb9c77.md}{Canonical record: 561e1ef130e04f38b854665bb8fb9c77}; full qualifications and provenance.}\par
\subsubsection*{II.1.3.1.14 [CT2] Order Administration}
\phantomsection
\label{navigation:II.1.3.1.14}
\addcontentsline{toc}{subsubsection}{II.1.3.1.14 [CT2] Order Administration}

{\small Historical predecessor: R.2.65 \hyperref[element:ebb735a2915744819d4001d9b8425cc6]{[CT2] Order Administration}; Canonical ID: \texttt{\seqsplit{ebb735a2915744819d4001d9b8425cc6}}. Superseded historical identity — direct correspondence · superseded.}\par
\label{element:c3c5abd7987f4433a46cbfed8167914c}
\index{[CT2] Order Administration (II.1.3.1.14)}
{\small Capability · Canonical · current}\par

Order Administration manages the closed loop around a request for defined fulfilment. It receives the request, resolves and routes to a destination, distinguishes delivery from business disposition, tracks progress and associates the eventual outcome. This makes the difference between an order sent and an order fulfilled visible throughout the requester's interaction.



\textbf{Overview}



Order Administration manages the closed loop around a request for defined fulfilment. It receives the request, resolves and routes to a destination, distinguishes delivery from business disposition, tracks progress and associates the eventual outcome. This makes the difference between an order sent and an order fulfilled visible throughout the requester's interaction.



\textbf{Intent}



Order Administration manages the closed loop around a request for defined fulfilment. It receives the request, resolves and routes to a destination, distinguishes delivery from business disposition, tracks progress and associates the eventual outcome. This makes the difference between an order sent and an order fulfilled visible throughout the requester's interaction.



\textbf{Boundary}



Information administration does not confer authority over the represented person, organisation, service, place, device or clinical activity.



\medskip\hrule\medskip

{\footnotesize \href{../../../../continuum/strategy/capability/c3c5abd7987f4433a46cbfed8167914c.md}{Canonical record: c3c5abd7987f4433a46cbfed8167914c}; full qualifications and provenance.}\par
\paragraph*{II.1.3.1.14.1 [CT3] Closed-Loop Order Progression}
\phantomsection
\label{navigation:II.1.3.1.14.1}
\addcontentsline{toc}{paragraph}{II.1.3.1.14.1 [CT3] Closed-Loop Order Progression}

{\small Historical predecessor: R.2.64 \hyperref[element:2ad29cebc1e94053a300d869e7797232]{[CT3] Closed-Loop Order Progression}; Canonical ID: \texttt{\seqsplit{2ad29cebc1e94053a300d869e7797232}}. Superseded historical identity — direct correspondence · superseded.}\par
\label{element:2f79ebd7392c46ef92def6f25bc9c7ce}
\index{[CT3] Closed-Loop Order Progression (II.1.3.1.14.1)}
{\small Capability · Proposed · current}\par

Administer the traceable progression of an order from request through destination, delivery, acknowledgement, fulfilment and outcome association.



\textbf{Overview}



Administer the traceable progression of an order from request through destination, delivery, acknowledgement, fulfilment and outcome association. The reviewed hierarchy assigns this enduring ability to its stated business-enabling context.



\textbf{Intent}



Administer the traceable progression of an order from request through destination, delivery, acknowledgement, fulfilment and outcome association.



\textbf{Boundary}



Information administration and coordination preserve the authority of the responsible source, clinical actor and governed domain; they confer no additional real-world management authority.


{\footnotesize \href{../../../../continuum/strategy/capability/2f79ebd7392c46ef92def6f25bc9c7ce.md}{Canonical record: 2f79ebd7392c46ef92def6f25bc9c7ce}; full qualifications and provenance.}\par
\subparagraph*{II.1.3.1.14.1.1 [Feature] Order Cancellation \& Modification Coordination}
\phantomsection
\label{navigation:II.1.3.1.14.1.1}
\addcontentsline{toc}{subparagraph}{II.1.3.1.14.1.1 [Feature] Order Cancellation \& Modification Coordination}

\label{element:6d12d9ddf4f8425ca60b289f3356f683}
\index{[Feature] Order Cancellation \& Modification Coordination (II.1.3.1.14.1.1)}
{\small Capability / Feature · Canonical · current}\par

Coordinate order cancellation and amendment requests, ensuring synchronization with performing systems.



Cancellation and amendment coordination carries a requested change into the performing-system context. It preserves the difference between asking to change an order and evidence that the change took effect, rather than silently rewriting the original request as though no prior intent existed.


{\footnotesize \href{../../../../continuum/strategy/capability/feature/6d12d9ddf4f8425ca60b289f3356f683.md}{Canonical record: 6d12d9ddf4f8425ca60b289f3356f683}; full qualifications and provenance.}\par
\subparagraph*{II.1.3.1.14.1.2 [Feature] Order Closed-Loop Progression Tracking}
\phantomsection
\label{navigation:II.1.3.1.14.1.2}
\addcontentsline{toc}{subparagraph}{II.1.3.1.14.1.2 [Feature] Order Closed-Loop Progression Tracking}

\label{element:9e63463942b94a6ba3bc045a04d3580b}
\index{[Feature] Order Closed-Loop Progression Tracking (II.1.3.1.14.1.2)}
{\small Capability / Feature · Canonical · current}\par

Track order progression across destination determination, routing, delivery, receipt acknowledgement (distinguishing technical delivery acknowledgement from business acknowledgement), execution progress, completion, and outcome association.



Closed-loop tracking keeps delivery, technical acknowledgement, business disposition, execution and returned outcome distinct. Its purpose is to expose what has actually been established at each stage, preventing a successful handoff from being reported as eventual fulfilment.


{\footnotesize \href{../../../../continuum/strategy/capability/feature/9e63463942b94a6ba3bc045a04d3580b.md}{Canonical record: 9e63463942b94a6ba3bc045a04d3580b}; full qualifications and provenance.}\par
\subparagraph*{II.1.3.1.14.1.3 [Feature] Order Destination Resolution \& Routing}
\phantomsection
\label{navigation:II.1.3.1.14.1.3}
\addcontentsline{toc}{subparagraph}{II.1.3.1.14.1.3 [Feature] Order Destination Resolution \& Routing}

\label{element:9f8f1542efc64e89ad226689291f2635}
\index{[Feature] Order Destination Resolution \& Routing (II.1.3.1.14.1.3)}
{\small Capability / Feature · Canonical · current}\par

Resolve fulfillment endpoints and reliably route orders to performing service systems.



Destination resolution connects an order with the service endpoint intended to fulfil it. Routing advances the request towards that destination without proving execution or changing the original clinical intent simply because a transport route was chosen.


{\footnotesize \href{../../../../continuum/strategy/capability/feature/9f8f1542efc64e89ad226689291f2635.md}{Canonical record: 9f8f1542efc64e89ad226689291f2635}; full qualifications and provenance.}\par
\subparagraph*{II.1.3.1.14.1.4 [Feature] Order Outcome Association}
\phantomsection
\label{navigation:II.1.3.1.14.1.4}
\addcontentsline{toc}{subparagraph}{II.1.3.1.14.1.4 [Feature] Order Outcome Association}

\label{element:4ad2e11162f34775adfea01b949d2803}
\index{[Feature] Order Outcome Association (II.1.3.1.14.1.4)}
{\small Capability / Feature · Proposed · current}\par

Associate an externally established outcome with the applicable bounded order and its progression.



\textbf{Overview}



Execution completion and returned outcome are different checkpoints. This Feature makes the result link explicit without equating an order acknowledgement with an outcome.



\textbf{Boundary}



Does not interpret diagnostic findings or declare fulfilment from transport delivery alone.


{\footnotesize \href{../../../../continuum/strategy/capability/feature/4ad2e11162f34775adfea01b949d2803.md}{Canonical record: 4ad2e11162f34775adfea01b949d2803}; full qualifications and provenance.}\par
\subparagraph*{II.1.3.1.14.1.5 [Feature] Order Request Ingestion}
\phantomsection
\label{navigation:II.1.3.1.14.1.5}
\addcontentsline{toc}{subparagraph}{II.1.3.1.14.1.5 [Feature] Order Request Ingestion}

\label{element:fc1f8168fd6949beb810e4ff57a2368e}
\index{[Feature] Order Request Ingestion (II.1.3.1.14.1.5)}
{\small Capability / Feature · Canonical · current}\par

Ingest clinical diagnostic, medication, and procedure order requests with validated ordering clinician context.



Order ingestion establishes the received bounded fulfilment request with its ordering-clinician context. It begins the information account needed for a closed loop while leaving the performing service's business acceptance and execution distinguishable from technical ingress acceptance.


{\footnotesize \href{../../../../continuum/strategy/capability/feature/fc1f8168fd6949beb810e4ff57a2368e.md}{Canonical record: fc1f8168fd6949beb810e4ff57a2368e}; full qualifications and provenance.}\par
\subsubsection*{II.1.3.1.15 [CT2] Procedure Administration}
\phantomsection
\label{navigation:II.1.3.1.15}
\addcontentsline{toc}{subsubsection}{II.1.3.1.15 [CT2] Procedure Administration}

{\small Historical predecessor: R.2.66 \hyperref[element:466a4057001846bb8ae0481453d3e410]{[CT2] Procedure Administration}; Canonical ID: \texttt{\seqsplit{466a4057001846bb8ae0481453d3e410}}. Superseded historical identity — direct correspondence · superseded.}\par
\label{element:649dd7cc26094e539386a050086dfb91}
\index{[CT2] Procedure Administration (II.1.3.1.15)}
{\small Capability · Canonical · current}\par

Procedure Administration integrates procedural bookings and the documentation of procedures performed. It keeps the scheduled case distinct from operative, anaesthetic and post-procedure information received later. That separation supports continuity across scheduling and clinical documentation without inferring that a booking establishes performance.



\textbf{Overview}



Procedure Administration integrates procedural bookings and the documentation of procedures performed. It keeps the scheduled case distinct from operative, anaesthetic and post-procedure information received later. That separation supports continuity across scheduling and clinical documentation without inferring that a booking establishes performance.



\textbf{Intent}



Procedure Administration integrates procedural bookings and the documentation of procedures performed. It keeps the scheduled case distinct from operative, anaesthetic and post-procedure information received later. That separation supports continuity across scheduling and clinical documentation without inferring that a booking establishes performance.



\textbf{Boundary}



Information administration does not confer authority over the represented person, organisation, service, place, device or clinical activity.



\medskip\hrule\medskip

{\footnotesize \href{../../../../continuum/strategy/capability/649dd7cc26094e539386a050086dfb91.md}{Canonical record: 649dd7cc26094e539386a050086dfb91}; full qualifications and provenance.}\par
\paragraph*{II.1.3.1.15.1 [CT3] Procedure \& Surgical Administration}
\phantomsection
\label{navigation:II.1.3.1.15.1}
\addcontentsline{toc}{paragraph}{II.1.3.1.15.1 [CT3] Procedure \& Surgical Administration}

\label{element:3c4a53dd8c8f468f83ca5d92397dd7e8}
\index{[CT3] Procedure \& Surgical Administration (II.1.3.1.15.1)}
{\small Capability · Proposed · current}\par

Administer the established procedure \& surgical administration information and coordination responsibility while preserving its source authority.



\textbf{Overview}



Administer the established procedure \& surgical administration information and coordination responsibility while preserving its source authority. The reviewed hierarchy assigns this enduring ability to its stated business-enabling context.



\textbf{Intent}



Administer the established procedure \& surgical administration information and coordination responsibility while preserving its source authority.



\textbf{Boundary}



Information administration and coordination preserve the authority of the responsible source, clinical actor and governed domain; they confer no additional real-world management authority.


{\footnotesize \href{../../../../continuum/strategy/capability/3c4a53dd8c8f468f83ca5d92397dd7e8.md}{Canonical record: 3c4a53dd8c8f468f83ca5d92397dd7e8}; full qualifications and provenance.}\par
\subparagraph*{II.1.3.1.15.1.1 [Feature] Procedural Documentation Ingestion}
\phantomsection
\label{navigation:II.1.3.1.15.1.1}
\addcontentsline{toc}{subparagraph}{II.1.3.1.15.1.1 [Feature] Procedural Documentation Ingestion}

\label{element:324aefc4578f4cae934101ae5ad74ff8}
\index{[Feature] Procedural Documentation Ingestion (II.1.3.1.15.1.1)}
{\small Capability / Feature · Canonical · current}\par

Ingest operative reports, anesthesia records, and post-procedure summaries.



Procedural documentation ingestion preserves the operative, anaesthetic and post-procedure account supplied by the clinical service. It connects performed-care documentation with the wider record without giving the integration process authority to originate or judge the clinical findings.


{\footnotesize \href{../../../../continuum/strategy/capability/feature/324aefc4578f4cae934101ae5ad74ff8.md}{Canonical record: 324aefc4578f4cae934101ae5ad74ff8}; full qualifications and provenance.}\par
\subparagraph*{II.1.3.1.15.1.2 [Feature] Procedure Booking Ingestion}
\phantomsection
\label{navigation:II.1.3.1.15.1.2}
\addcontentsline{toc}{subparagraph}{II.1.3.1.15.1.2 [Feature] Procedure Booking Ingestion}

\label{element:12e0f64cff2043339be17fb38708247a}
\index{[Feature] Procedure Booking Ingestion (II.1.3.1.15.1.2)}
{\small Capability / Feature · Canonical · current}\par

Ingest scheduled surgical case bookings and procedural requests.



Procedure booking ingestion makes scheduled cases and procedural requests available in the wider care context. It conveys intended or arranged activity while retaining the distinction between a booked procedure and one actually performed.


{\footnotesize \href{../../../../continuum/strategy/capability/feature/12e0f64cff2043339be17fb38708247a.md}{Canonical record: 12e0f64cff2043339be17fb38708247a}; full qualifications and provenance.}\par
\subsubsection*{II.1.3.1.16 [CT2] Referral Administration}
\phantomsection
\label{navigation:II.1.3.1.16}
\addcontentsline{toc}{subsubsection}{II.1.3.1.16 [CT2] Referral Administration}

{\small Historical predecessor: R.2.68 \hyperref[element:5b0ccf4b30054145952c298b348ed266]{[CT2] Referral Administration}; Canonical ID: \texttt{\seqsplit{5b0ccf4b30054145952c298b348ed266}}. Superseded historical identity — direct correspondence · superseded.}\par
\label{element:ae0de25600ee416aa24d2a167b64f230}
\index{[CT2] Referral Administration (II.1.3.1.16)}
{\small Capability · Canonical · current}\par

Referral Administration receives referral material, associates the supporting clinical context and tracks the recipient's explicit disposition and outcome. Its purpose is to keep the request and its progression intelligible across care boundaries. It distinguishes a request for another participant's involvement from an order for bounded fulfilment and from actual acceptance of care responsibility.



\textbf{Overview}



Referral Administration receives referral material, associates the supporting clinical context and tracks the recipient's explicit disposition and outcome. Its purpose is to keep the request and its progression intelligible across care boundaries. It distinguishes a request for another participant's involvement from an order for bounded fulfilment and from actual acceptance of care responsibility.



\textbf{Intent}



Referral Administration receives referral material, associates the supporting clinical context and tracks the recipient's explicit disposition and outcome. Its purpose is to keep the request and its progression intelligible across care boundaries. It distinguishes a request for another participant's involvement from an order for bounded fulfilment and from actual acceptance of care responsibility.



\textbf{Boundary}



Information administration does not confer authority over the represented person, organisation, service, place, device or clinical activity.



\medskip\hrule\medskip

{\footnotesize \href{../../../../continuum/strategy/capability/ae0de25600ee416aa24d2a167b64f230.md}{Canonical record: ae0de25600ee416aa24d2a167b64f230}; full qualifications and provenance.}\par
\paragraph*{II.1.3.1.16.1 [CT3] Referral Intake \& Progression}
\phantomsection
\label{navigation:II.1.3.1.16.1}
\addcontentsline{toc}{paragraph}{II.1.3.1.16.1 [CT3] Referral Intake \& Progression}

{\small Historical predecessor: R.2.67 \hyperref[element:7180f6265d5c4687a001220c6126e301]{[CT3] Referral Intake \& Progression}; Canonical ID: \texttt{\seqsplit{7180f6265d5c4687a001220c6126e301}}. Superseded historical identity — direct correspondence · superseded.}\par
\label{element:d11bdb9464b2476894aa372ab9b929a7}
\index{[CT3] Referral Intake \& Progression (II.1.3.1.16.1)}
{\small Capability · Proposed · current}\par

Receive and correlate referral requests, supporting information, destinations, disposition and progression through responsible care systems.



\textbf{Overview}



Receive and correlate referral requests, supporting information, destinations, disposition and progression through responsible care systems. The reviewed hierarchy assigns this enduring ability to its stated business-enabling context.



\textbf{Intent}



Receive and correlate referral requests, supporting information, destinations, disposition and progression through responsible care systems.



\textbf{Boundary}



Information administration and coordination preserve the authority of the responsible source, clinical actor and governed domain; they confer no additional real-world management authority.


{\footnotesize \href{../../../../continuum/strategy/capability/d11bdb9464b2476894aa372ab9b929a7.md}{Canonical record: d11bdb9464b2476894aa372ab9b929a7}; full qualifications and provenance.}\par
\subparagraph*{II.1.3.1.16.1.1 [Feature] Referral Destination Coordination}
\phantomsection
\label{navigation:II.1.3.1.16.1.1}
\addcontentsline{toc}{subparagraph}{II.1.3.1.16.1.1 [Feature] Referral Destination Coordination}

\label{element:7d7ddffcda9d4854a02be5c3265ad775}
\index{[Feature] Referral Destination Coordination (II.1.3.1.16.1.1)}
{\small Capability / Feature · Proposed · current}\par

Coordinate qualified referral destination resolution and convey the request under its applicable disposition contract.



\textbf{Overview}



Referral intake and status capture leave a distinct routing responsibility: finding the intended service and preserving redirection without turning technical delivery into assumed care responsibility.



\textbf{Boundary}



Does not make clinical triage decisions or infer Transfer of Care from acceptance.


{\footnotesize \href{../../../../continuum/strategy/capability/feature/7d7ddffcda9d4854a02be5c3265ad775.md}{Canonical record: 7d7ddffcda9d4854a02be5c3265ad775}; full qualifications and provenance.}\par
\subparagraph*{II.1.3.1.16.1.2 [Feature] Referral Document Ingestion}
\phantomsection
\label{navigation:II.1.3.1.16.1.2}
\addcontentsline{toc}{subparagraph}{II.1.3.1.16.1.2 [Feature] Referral Document Ingestion}

\label{element:1d1c5f8e8ca1498fb68dba601a7079dc}
\index{[Feature] Referral Document Ingestion (II.1.3.1.16.1.2)}
{\small Capability / Feature · Canonical · current}\par

Ingest structured clinical referral messages across integration boundaries.



Referral ingestion establishes the received referral material within the integration boundary. It preserves the request for another participant's involvement without treating successful receipt as clinical acceptance, responsibility assumption or a promise that care will be delivered.


{\footnotesize \href{../../../../continuum/strategy/capability/feature/1d1c5f8e8ca1498fb68dba601a7079dc.md}{Canonical record: 1d1c5f8e8ca1498fb68dba601a7079dc}; full qualifications and provenance.}\par
\subparagraph*{II.1.3.1.16.1.3 [Feature] Referral Status \& Outcome Tracking}
\phantomsection
\label{navigation:II.1.3.1.16.1.3}
\addcontentsline{toc}{subparagraph}{II.1.3.1.16.1.3 [Feature] Referral Status \& Outcome Tracking}

\label{element:836fa704f7b24bfc8f7eeb72a36bb897}
\index{[Feature] Referral Status \& Outcome Tracking (II.1.3.1.16.1.3)}
{\small Capability / Feature · Canonical · current}\par

Track referral operational progression and communicate referral outcomes (acceptance, decline, or redirection) to referring providers.



Status and outcome tracking makes explicit referral dispositions visible to the referring participant. Acceptance, decline and redirection remain meaningful alternatives; a transport acknowledgement does not fill in a missing business decision or establish a universal referral progression.


{\footnotesize \href{../../../../continuum/strategy/capability/feature/836fa704f7b24bfc8f7eeb72a36bb897.md}{Canonical record: 836fa704f7b24bfc8f7eeb72a36bb897}; full qualifications and provenance.}\par
\subparagraph*{II.1.3.1.16.1.4 [Feature] Referral Supporting Information}
\phantomsection
\label{navigation:II.1.3.1.16.1.4}
\addcontentsline{toc}{subparagraph}{II.1.3.1.16.1.4 [Feature] Referral Supporting Information}

\label{element:c8a50c0e932847d7b7d0289101aacd98}
\index{[Feature] Referral Supporting Information (II.1.3.1.16.1.4)}
{\small Capability / Feature · Canonical · current}\par

Collate and associate supporting clinical documentation, diagnostic history, and patient context with referral requests.



Supporting information gives the referral its relevant history and clinical context. Association makes the material available for the recipient's consideration while preserving source provenance and the distinction between evidence supplied and the recipient's own clinical decision.


{\footnotesize \href{../../../../continuum/strategy/capability/feature/c8a50c0e932847d7b7d0289101aacd98.md}{Canonical record: c8a50c0e932847d7b7d0289101aacd98}; full qualifications and provenance.}\par
\subsubsection*{II.1.3.1.17 [CT2] Scheduling Administration}
\phantomsection
\label{navigation:II.1.3.1.17}
\addcontentsline{toc}{subsubsection}{II.1.3.1.17 [CT2] Scheduling Administration}

{\small Historical predecessor: R.2.70 \hyperref[element:1a6630940ffb44fcad4543b93a367b67]{[CT2] Scheduling Administration}; Canonical ID: \texttt{\seqsplit{1a6630940ffb44fcad4543b93a367b67}}. Superseded historical identity — direct correspondence · superseded.}\par
\label{element:35fb3744c900470f9895056246a13fb1}
\index{[CT2] Scheduling Administration (II.1.3.1.17)}
{\small Capability · Canonical · current}\par

Scheduling Administration integrates appointment information established by departmental scheduling systems. It receives scheduled events and propagates cancellation, rescheduling and attendance changes so that care-coordination contexts do not diverge. Its concern is the coherence of scheduling information across participants rather than originating every booking decision.



\textbf{Overview}



Scheduling Administration integrates appointment information established by departmental scheduling systems. It receives scheduled events and propagates cancellation, rescheduling and attendance changes so that care-coordination contexts do not diverge. Its concern is the coherence of scheduling information across participants rather than originating every booking decision.



\textbf{Intent}



Scheduling Administration integrates appointment information established by departmental scheduling systems. It receives scheduled events and propagates cancellation, rescheduling and attendance changes so that care-coordination contexts do not diverge. Its concern is the coherence of scheduling information across participants rather than originating every booking decision.



\textbf{Boundary}



Information administration does not confer authority over the represented person, organisation, service, place, device or clinical activity.



\medskip\hrule\medskip

{\footnotesize \href{../../../../continuum/strategy/capability/35fb3744c900470f9895056246a13fb1.md}{Canonical record: 35fb3744c900470f9895056246a13fb1}; full qualifications and provenance.}\par
\paragraph*{II.1.3.1.17.1 [CT3] Appointment \& Booking Administration}
\phantomsection
\label{navigation:II.1.3.1.17.1}
\addcontentsline{toc}{paragraph}{II.1.3.1.17.1 [CT3] Appointment \& Booking Administration}

{\small Historical predecessor: R.2.69 \hyperref[element:cdc16dd0f84049f8af3f1ec76f1e1428]{[CT3] Appointment \& Booking Administration}; Canonical ID: \texttt{\seqsplit{cdc16dd0f84049f8af3f1ec76f1e1428}}. Superseded historical identity — direct correspondence · superseded.}\par
\label{element:7b3ba9bcac564dd4a7a08631581db2ca}
\index{[CT3] Appointment \& Booking Administration (II.1.3.1.17.1)}
{\small Capability · Proposed · current}\par

Administer externally supplied appointment and booking information and communicate changes in their scheduling status.



\textbf{Overview}



Administer externally supplied appointment and booking information and communicate changes in their scheduling status. The reviewed hierarchy assigns this enduring ability to its stated business-enabling context.



\textbf{Intent}



Administer externally supplied appointment and booking information and communicate changes in their scheduling status.



\textbf{Boundary}



Information administration and coordination preserve the authority of the responsible source, clinical actor and governed domain; they confer no additional real-world management authority.


{\footnotesize \href{../../../../continuum/strategy/capability/7b3ba9bcac564dd4a7a08631581db2ca.md}{Canonical record: 7b3ba9bcac564dd4a7a08631581db2ca}; full qualifications and provenance.}\par
\subparagraph*{II.1.3.1.17.1.1 [Feature] Appointment Notification Ingestion}
\phantomsection
\label{navigation:II.1.3.1.17.1.1}
\addcontentsline{toc}{subparagraph}{II.1.3.1.17.1.1 [Feature] Appointment Notification Ingestion}

\label{element:1d5bc2b1742544728357fbea42a55ec4}
\index{[Feature] Appointment Notification Ingestion (II.1.3.1.17.1.1)}
{\small Capability / Feature · Canonical · current}\par

Ingest scheduled appointment events from departmental scheduling systems.



Appointment ingestion receives scheduling facts established by departmental systems. It supports shared awareness of planned care opportunities without claiming authority over the underlying booking decision or proving that attendance and clinical fulfilment will occur.


{\footnotesize \href{../../../../continuum/strategy/capability/feature/1d5bc2b1742544728357fbea42a55ec4.md}{Canonical record: 1d5bc2b1742544728357fbea42a55ec4}; full qualifications and provenance.}\par
\subparagraph*{II.1.3.1.17.1.2 [Feature] Appointment Status Synchronization}
\phantomsection
\label{navigation:II.1.3.1.17.1.2}
\addcontentsline{toc}{subparagraph}{II.1.3.1.17.1.2 [Feature] Appointment Status Synchronization}

\label{element:a7a9932600d84434a1a37005f726747d}
\index{[Feature] Appointment Status Synchronization (II.1.3.1.17.1.2)}
{\small Capability / Feature · Canonical · current}\par

Propagate booking cancellations, reschedules, and attendances across care coordination systems.



Status synchronisation keeps cancellation, rescheduling and attendance information coherent across coordination contexts. The purpose is propagation of the externally established booking change rather than independent scheduling or inference of a clinical outcome from appointment state.


{\footnotesize \href{../../../../continuum/strategy/capability/feature/a7a9932600d84434a1a37005f726747d.md}{Canonical record: a7a9932600d84434a1a37005f726747d}; full qualifications and provenance.}\par
\subsection*{II.1.3.2 [CT1] Care Context Integration Enablement [superseded]}
\phantomsection
\label{navigation:II.1.3.2}
\addcontentsline{toc}{subsection}{II.1.3.2 [CT1] Care Context Integration Enablement [superseded]}

\label{element:a84a7e96e7874a4b87d18c53b1b5cf29}
\index{[CT1] Care Context Integration Enablement (II.1.3.2)}
{\small Capability · Proposed · superseded}\par
\href{../../../../continuum/strategy/capability/a84a7e96e7874a4b87d18c53b1b5cf29.md}{Historical canonical definition and provenance}.\par
\subsubsection*{II.1.3.2.1 [CT2] Facility \& Episodic Care Integration [superseded]}
\phantomsection
\label{navigation:II.1.3.2.1}
\addcontentsline{toc}{subsubsection}{II.1.3.2.1 [CT2] Facility \& Episodic Care Integration [superseded]}

\label{element:c1eb24ed7e36417cb426288ff57512c0}
\index{[CT2] Facility \& Episodic Care Integration (II.1.3.2.1)}
{\small Capability · Proposed · superseded}\par
\href{../../../../continuum/strategy/capability/c1eb24ed7e36417cb426288ff57512c0.md}{Historical canonical definition and provenance}.\par
\subsubsection*{II.1.3.2.2 [CT2] Longitudinal \& Community Care Integration [superseded]}
\phantomsection
\label{navigation:II.1.3.2.2}
\addcontentsline{toc}{subsubsection}{II.1.3.2.2 [CT2] Longitudinal \& Community Care Integration [superseded]}

\label{element:8430d50adc504ecc959f15b7cbbcd095}
\index{[CT2] Longitudinal \& Community Care Integration (II.1.3.2.2)}
{\small Capability · Proposed · superseded}\par
\href{../../../../continuum/strategy/capability/8430d50adc504ecc959f15b7cbbcd095.md}{Historical canonical definition and provenance}.\par
\subsubsection*{II.1.3.2.3 [CT2] Specialist Clinical Context Integration [superseded]}
\phantomsection
\label{navigation:II.1.3.2.3}
\addcontentsline{toc}{subsubsection}{II.1.3.2.3 [CT2] Specialist Clinical Context Integration [superseded]}

\label{element:062802a425d14530b0ce9434c5980d85}
\index{[CT2] Specialist Clinical Context Integration (II.1.3.2.3)}
{\small Capability · Proposed · superseded}\par
\href{../../../../continuum/strategy/capability/062802a425d14530b0ce9434c5980d85.md}{Historical canonical definition and provenance}.\par
\subsection*{II.1.3.3 [CT1] Clinical Service Dependency Management}
\phantomsection
\label{navigation:II.1.3.3}
\addcontentsline{toc}{subsection}{II.1.3.3 [CT1] Clinical Service Dependency Management}

{\small Historical predecessor: R.2.88 \hyperref[element:984ee17123eb4f92b3e55e55d24459b8]{[CT1] Clinical Service Dependency Management}; Canonical ID: \texttt{\seqsplit{984ee17123eb4f92b3e55e55d24459b8}}. Superseded historical identity — direct correspondence · superseded.}\par
\label{element:232db9b8b47c450cb0aa9acaf539c09d}
\index{[CT1] Clinical Service Dependency Management (II.1.3.3)}
{\small Capability · Proposed · current}\par

Identify and track operational and infrastructure dependencies necessary for clinical services to function, including care-place and environmental readiness, devices, materiel, forms, workstations, connectivity, infrastructure and basic working materials.



\textbf{Overview}



Identify and track operational and infrastructure dependencies necessary for clinical services to function, including care-place and environmental readiness, devices, materiel, forms, workstations, connectivity, infrastructure and basic working materials. The reviewed hierarchy assigns this enduring ability to its stated business-enabling context.



\textbf{Intent}



Identify and track operational and infrastructure dependencies necessary for clinical services to function, including care-place and environmental readiness, devices, materiel, forms, workstations, connectivity, infrastructure and basic working materials.



\textbf{Boundary}



Information administration and coordination preserve the authority of the responsible source, clinical actor and governed domain; they confer no additional real-world management authority.


{\footnotesize \href{../../../../continuum/strategy/capability/232db9b8b47c450cb0aa9acaf539c09d.md}{Canonical record: 232db9b8b47c450cb0aa9acaf539c09d}; full qualifications and provenance.}\par
\subsubsection*{II.1.3.3.1 [CT2] Acute Care}
\phantomsection
\label{navigation:II.1.3.3.1}
\addcontentsline{toc}{subsubsection}{II.1.3.3.1 [CT2] Acute Care}

{\small Historical predecessor: R.2.72 \hyperref[element:7b591957f7424424abc80f699e9f1960]{[CT2] Acute Care}; Canonical ID: \texttt{\seqsplit{7b591957f7424424abc80f699e9f1960}}. Superseded historical identity — direct correspondence · superseded.}\par
\label{element:9bb988c526404522b83aeed5069e0eb5}
\index{[CT2] Acute Care (II.1.3.3.1)}
{\small Capability · Proposed · current}\par

Identify and track the operational and infrastructure dependencies necessary for acute care services to function, retaining each dependency provider's responsibility.



\textbf{Overview}



Identify and track the operational and infrastructure dependencies necessary for acute care services to function, retaining each dependency provider's responsibility. The reviewed hierarchy assigns this enduring ability to its stated business-enabling context.



\textbf{Intent}



Identify and track the operational and infrastructure dependencies necessary for acute care services to function, retaining each dependency provider's responsibility.



\textbf{Boundary}



Information administration and coordination preserve the authority of the responsible source, clinical actor and governed domain; they confer no additional real-world management authority.


{\footnotesize \href{../../../../continuum/strategy/capability/9bb988c526404522b83aeed5069e0eb5.md}{Canonical record: 9bb988c526404522b83aeed5069e0eb5}; full qualifications and provenance.}\par
\subsubsection*{II.1.3.3.2 [CT2] Allied Health}
\phantomsection
\label{navigation:II.1.3.3.2}
\addcontentsline{toc}{subsubsection}{II.1.3.3.2 [CT2] Allied Health}

{\small Historical predecessor: R.2.73 \hyperref[element:6cd45c27d0c344b9807174dbd812d572]{[CT2] Allied Health}; Canonical ID: \texttt{\seqsplit{6cd45c27d0c344b9807174dbd812d572}}. Superseded historical identity — direct correspondence · superseded.}\par
\label{element:63cd714a3e2a4354aa5383b189fea566}
\index{[CT2] Allied Health (II.1.3.3.2)}
{\small Capability · Canonical · current}\par

Identify and track the operational and infrastructure dependencies necessary for allied health services to function, retaining each dependency provider's responsibility.



\textbf{Overview}



Identify and track the operational and infrastructure dependencies necessary for allied health services to function, retaining each dependency provider's responsibility.



\textbf{Intent}



Identify and track the operational and infrastructure dependencies necessary for allied health services to function, retaining each dependency provider's responsibility.



\textbf{Boundary}



Information administration does not confer authority over the represented person, organisation, service, place, device or clinical activity.



\medskip\hrule\medskip

{\footnotesize \href{../../../../continuum/strategy/capability/63cd714a3e2a4354aa5383b189fea566.md}{Canonical record: 63cd714a3e2a4354aa5383b189fea566}; full qualifications and provenance.}\par
\subsubsection*{II.1.3.3.3 [CT2] Ambulatory Care}
\phantomsection
\label{navigation:II.1.3.3.3}
\addcontentsline{toc}{subsubsection}{II.1.3.3.3 [CT2] Ambulatory Care}

{\small Historical predecessor: R.2.74 \hyperref[element:e996b61b6e7b4619913f427aed2b55b1]{[CT2] Ambulatory Care}; Canonical ID: \texttt{\seqsplit{e996b61b6e7b4619913f427aed2b55b1}}. Superseded historical identity — direct correspondence · superseded.}\par
\label{element:f74d7c3a39c540c4b462576e8f09fc58}
\index{[CT2] Ambulatory Care (II.1.3.3.3)}
{\small Capability · Canonical · current}\par

Identify and track the operational and infrastructure dependencies necessary for ambulatory care services to function, retaining each dependency provider's responsibility.



\textbf{Overview}



Identify and track the operational and infrastructure dependencies necessary for ambulatory care services to function, retaining each dependency provider's responsibility.



\textbf{Intent}



Identify and track the operational and infrastructure dependencies necessary for ambulatory care services to function, retaining each dependency provider's responsibility.



\textbf{Boundary}



Information administration does not confer authority over the represented person, organisation, service, place, device or clinical activity.



\medskip\hrule\medskip

{\footnotesize \href{../../../../continuum/strategy/capability/f74d7c3a39c540c4b462576e8f09fc58.md}{Canonical record: f74d7c3a39c540c4b462576e8f09fc58}; full qualifications and provenance.}\par
\subsubsection*{II.1.3.3.4 [CT2] Clinical Assessment}
\phantomsection
\label{navigation:II.1.3.3.4}
\addcontentsline{toc}{subsubsection}{II.1.3.3.4 [CT2] Clinical Assessment}

{\small Historical predecessor: R.2.75 \hyperref[element:1c3b553c87704e27ba211fbbd02d95cb]{[CT2] Clinical Assessment}; Canonical ID: \texttt{\seqsplit{1c3b553c87704e27ba211fbbd02d95cb}}. Superseded historical identity — direct correspondence · superseded.}\par
\label{element:a7d84c95e48c453797419813ca80a66c}
\index{[CT2] Clinical Assessment (II.1.3.3.4)}
{\small Capability · Canonical · current}\par

Identify and track the operational and infrastructure dependencies necessary for clinical assessment services to function, retaining each dependency provider's responsibility.



\textbf{Overview}



Identify and track the operational and infrastructure dependencies necessary for clinical assessment services to function, retaining each dependency provider's responsibility.



\textbf{Intent}



Identify and track the operational and infrastructure dependencies necessary for clinical assessment services to function, retaining each dependency provider's responsibility.



\textbf{Boundary}



Information administration does not confer authority over the represented person, organisation, service, place, device or clinical activity.



\medskip\hrule\medskip

{\footnotesize \href{../../../../continuum/strategy/capability/a7d84c95e48c453797419813ca80a66c.md}{Canonical record: a7d84c95e48c453797419813ca80a66c}; full qualifications and provenance.}\par
\subsubsection*{II.1.3.3.5 [CT2] Community Care}
\phantomsection
\label{navigation:II.1.3.3.5}
\addcontentsline{toc}{subsubsection}{II.1.3.3.5 [CT2] Community Care}

{\small Historical predecessor: R.2.76 \hyperref[element:4b7177f3f40449518fba381acae6ba1e]{[CT2] Community Care}; Canonical ID: \texttt{\seqsplit{4b7177f3f40449518fba381acae6ba1e}}. Superseded historical identity — direct correspondence · superseded.}\par
\label{element:942aa78cad004ef99de6b3aae8ddd014}
\index{[CT2] Community Care (II.1.3.3.5)}
{\small Capability · Proposed · current}\par

Identify and track the operational and infrastructure dependencies necessary for community care services to function, retaining each dependency provider's responsibility.



\textbf{Overview}



Identify and track the operational and infrastructure dependencies necessary for community care services to function, retaining each dependency provider's responsibility. The reviewed hierarchy assigns this enduring ability to its stated business-enabling context.



\textbf{Intent}



Identify and track the operational and infrastructure dependencies necessary for community care services to function, retaining each dependency provider's responsibility.



\textbf{Boundary}



Information administration and coordination preserve the authority of the responsible source, clinical actor and governed domain; they confer no additional real-world management authority.


{\footnotesize \href{../../../../continuum/strategy/capability/942aa78cad004ef99de6b3aae8ddd014.md}{Canonical record: 942aa78cad004ef99de6b3aae8ddd014}; full qualifications and provenance.}\par
\subsubsection*{II.1.3.3.6 [CT2] Diagnostic Services}
\phantomsection
\label{navigation:II.1.3.3.6}
\addcontentsline{toc}{subsubsection}{II.1.3.3.6 [CT2] Diagnostic Services}

{\small Historical predecessor: R.2.77 \hyperref[element:03a57e0c9ea34f7fb7b78e5f5c85379f]{[CT2] Diagnostic Services}; Canonical ID: \texttt{\seqsplit{03a57e0c9ea34f7fb7b78e5f5c85379f}}. Superseded historical identity — direct correspondence · superseded.}\par
\label{element:b825ebab42f444d8b738b4c1f489f785}
\index{[CT2] Diagnostic Services (II.1.3.3.6)}
{\small Capability · Proposed · current}\par

Identify and track the operational and infrastructure dependencies necessary for diagnostic services services to function, retaining each dependency provider's responsibility.



\textbf{Overview}



Identify and track the operational and infrastructure dependencies necessary for diagnostic services services to function, retaining each dependency provider's responsibility. The reviewed hierarchy assigns this enduring ability to its stated business-enabling context.



\textbf{Intent}



Identify and track the operational and infrastructure dependencies necessary for diagnostic services services to function, retaining each dependency provider's responsibility.



\textbf{Boundary}



Information administration and coordination preserve the authority of the responsible source, clinical actor and governed domain; they confer no additional real-world management authority.


{\footnotesize \href{../../../../continuum/strategy/capability/b825ebab42f444d8b738b4c1f489f785.md}{Canonical record: b825ebab42f444d8b738b4c1f489f785}; full qualifications and provenance.}\par
\subsubsection*{II.1.3.3.7 [CT2] Emergency Care}
\phantomsection
\label{navigation:II.1.3.3.7}
\addcontentsline{toc}{subsubsection}{II.1.3.3.7 [CT2] Emergency Care}

{\small Historical predecessor: R.2.78 \hyperref[element:e5edbb54258c490aaba0816d980dc624]{[CT2] Emergency Care}; Canonical ID: \texttt{\seqsplit{e5edbb54258c490aaba0816d980dc624}}. Superseded historical identity — direct correspondence · superseded.}\par
\label{element:aa088f0647ad4b78966fe28fa80af674}
\index{[CT2] Emergency Care (II.1.3.3.7)}
{\small Capability · Proposed · current}\par

Identify and track the operational and infrastructure dependencies necessary for emergency care services to function, retaining each dependency provider's responsibility.



\textbf{Overview}



Identify and track the operational and infrastructure dependencies necessary for emergency care services to function, retaining each dependency provider's responsibility. The reviewed hierarchy assigns this enduring ability to its stated business-enabling context.



\textbf{Intent}



Identify and track the operational and infrastructure dependencies necessary for emergency care services to function, retaining each dependency provider's responsibility.



\textbf{Boundary}



Information administration and coordination preserve the authority of the responsible source, clinical actor and governed domain; they confer no additional real-world management authority.


{\footnotesize \href{../../../../continuum/strategy/capability/aa088f0647ad4b78966fe28fa80af674.md}{Canonical record: aa088f0647ad4b78966fe28fa80af674}; full qualifications and provenance.}\par
\subsubsection*{II.1.3.3.8 [CT2] Inpatient Care}
\phantomsection
\label{navigation:II.1.3.3.8}
\addcontentsline{toc}{subsubsection}{II.1.3.3.8 [CT2] Inpatient Care}

{\small Historical predecessor: R.2.79 \hyperref[element:bcc40a14465e48aca42c9505a6853081]{[CT2] Inpatient Care}; Canonical ID: \texttt{\seqsplit{bcc40a14465e48aca42c9505a6853081}}. Superseded historical identity — direct correspondence · superseded.}\par
\label{element:b39fb49c01b545a5b51d3313531efb1b}
\index{[CT2] Inpatient Care (II.1.3.3.8)}
{\small Capability · Proposed · current}\par

Identify and track the operational and infrastructure dependencies necessary for inpatient care services to function, retaining each dependency provider's responsibility.



\textbf{Overview}



Identify and track the operational and infrastructure dependencies necessary for inpatient care services to function, retaining each dependency provider's responsibility. The reviewed hierarchy assigns this enduring ability to its stated business-enabling context.



\textbf{Intent}



Identify and track the operational and infrastructure dependencies necessary for inpatient care services to function, retaining each dependency provider's responsibility.



\textbf{Boundary}



Information administration and coordination preserve the authority of the responsible source, clinical actor and governed domain; they confer no additional real-world management authority.


{\footnotesize \href{../../../../continuum/strategy/capability/b39fb49c01b545a5b51d3313531efb1b.md}{Canonical record: b39fb49c01b545a5b51d3313531efb1b}; full qualifications and provenance.}\par
\subsubsection*{II.1.3.3.9 [CT2] Medication Therapy}
\phantomsection
\label{navigation:II.1.3.3.9}
\addcontentsline{toc}{subsubsection}{II.1.3.3.9 [CT2] Medication Therapy}

{\small Historical predecessor: R.2.80 \hyperref[element:db064e0cd4794479be1374e97a3d7e63]{[CT2] Medication Therapy}; Canonical ID: \texttt{\seqsplit{db064e0cd4794479be1374e97a3d7e63}}. Superseded historical identity — direct correspondence · superseded.}\par
\label{element:bb04996164c94e0c827a020432f4e498}
\index{[CT2] Medication Therapy (II.1.3.3.9)}
{\small Capability · Proposed · current}\par

Identify and track the operational and infrastructure dependencies necessary for medication therapy services to function, retaining each dependency provider's responsibility.



\textbf{Overview}



Identify and track the operational and infrastructure dependencies necessary for medication therapy services to function, retaining each dependency provider's responsibility. The reviewed hierarchy assigns this enduring ability to its stated business-enabling context.



\textbf{Intent}



Identify and track the operational and infrastructure dependencies necessary for medication therapy services to function, retaining each dependency provider's responsibility.



\textbf{Boundary}



Information administration and coordination preserve the authority of the responsible source, clinical actor and governed domain; they confer no additional real-world management authority.


{\footnotesize \href{../../../../continuum/strategy/capability/bb04996164c94e0c827a020432f4e498.md}{Canonical record: bb04996164c94e0c827a020432f4e498}; full qualifications and provenance.}\par
\subsubsection*{II.1.3.3.10 [CT2] Outreach Care}
\phantomsection
\label{navigation:II.1.3.3.10}
\addcontentsline{toc}{subsubsection}{II.1.3.3.10 [CT2] Outreach Care}

{\small Historical predecessor: R.2.81 \hyperref[element:75f8b32a3d5048ccbdfb30d61ae73dad]{[CT2] Outreach Care}; Canonical ID: \texttt{\seqsplit{75f8b32a3d5048ccbdfb30d61ae73dad}}. Superseded historical identity — direct correspondence · superseded.}\par
\label{element:b9ddbe2dd6a7485ba9cd3482215ea09e}
\index{[CT2] Outreach Care (II.1.3.3.10)}
{\small Capability · Proposed · current}\par

Identify and track the operational and infrastructure dependencies necessary for outreach care services to function, retaining each dependency provider's responsibility.



\textbf{Overview}



Identify and track the operational and infrastructure dependencies necessary for outreach care services to function, retaining each dependency provider's responsibility. The reviewed hierarchy assigns this enduring ability to its stated business-enabling context.



\textbf{Intent}



Identify and track the operational and infrastructure dependencies necessary for outreach care services to function, retaining each dependency provider's responsibility.



\textbf{Boundary}



Information administration and coordination preserve the authority of the responsible source, clinical actor and governed domain; they confer no additional real-world management authority.


{\footnotesize \href{../../../../continuum/strategy/capability/b9ddbe2dd6a7485ba9cd3482215ea09e.md}{Canonical record: b9ddbe2dd6a7485ba9cd3482215ea09e}; full qualifications and provenance.}\par
\subsubsection*{II.1.3.3.11 [CT2] Preventive Care}
\phantomsection
\label{navigation:II.1.3.3.11}
\addcontentsline{toc}{subsubsection}{II.1.3.3.11 [CT2] Preventive Care}

{\small Historical predecessor: R.2.82 \hyperref[element:97f0d1acc37b4ca994c6c1ca006f2daa]{[CT2] Preventive Care}; Canonical ID: \texttt{\seqsplit{97f0d1acc37b4ca994c6c1ca006f2daa}}. Superseded historical identity — direct correspondence · superseded.}\par
\label{element:dfdc4f0273ba4910848b3b65c9ef6ab2}
\index{[CT2] Preventive Care (II.1.3.3.11)}
{\small Capability · Proposed · current}\par

Identify and track the operational and infrastructure dependencies necessary for preventive care services to function, retaining each dependency provider's responsibility.



\textbf{Overview}



Identify and track the operational and infrastructure dependencies necessary for preventive care services to function, retaining each dependency provider's responsibility. The reviewed hierarchy assigns this enduring ability to its stated business-enabling context.



\textbf{Intent}



Identify and track the operational and infrastructure dependencies necessary for preventive care services to function, retaining each dependency provider's responsibility.



\textbf{Boundary}



Information administration and coordination preserve the authority of the responsible source, clinical actor and governed domain; they confer no additional real-world management authority.


{\footnotesize \href{../../../../continuum/strategy/capability/dfdc4f0273ba4910848b3b65c9ef6ab2.md}{Canonical record: dfdc4f0273ba4910848b3b65c9ef6ab2}; full qualifications and provenance.}\par
\subsubsection*{II.1.3.3.12 [CT2] Primary Care}
\phantomsection
\label{navigation:II.1.3.3.12}
\addcontentsline{toc}{subsubsection}{II.1.3.3.12 [CT2] Primary Care}

{\small Historical predecessor: R.2.83 \hyperref[element:fe9d58b5cfe14e8cb7a83a1ac6686d65]{[CT2] Primary Care}; Canonical ID: \texttt{\seqsplit{fe9d58b5cfe14e8cb7a83a1ac6686d65}}. Superseded historical identity — direct correspondence · superseded.}\par
\label{element:d118673d77d74c8cb8cec24f03adddd8}
\index{[CT2] Primary Care (II.1.3.3.12)}
{\small Capability · Proposed · current}\par

Identify and track the operational and infrastructure dependencies necessary for primary care services to function, retaining each dependency provider's responsibility.



\textbf{Overview}



Identify and track the operational and infrastructure dependencies necessary for primary care services to function, retaining each dependency provider's responsibility. The reviewed hierarchy assigns this enduring ability to its stated business-enabling context.



\textbf{Intent}



Identify and track the operational and infrastructure dependencies necessary for primary care services to function, retaining each dependency provider's responsibility.



\textbf{Boundary}



Information administration and coordination preserve the authority of the responsible source, clinical actor and governed domain; they confer no additional real-world management authority.


{\footnotesize \href{../../../../continuum/strategy/capability/d118673d77d74c8cb8cec24f03adddd8.md}{Canonical record: d118673d77d74c8cb8cec24f03adddd8}; full qualifications and provenance.}\par
\subsubsection*{II.1.3.3.13 [CT2] Rehabilitation}
\phantomsection
\label{navigation:II.1.3.3.13}
\addcontentsline{toc}{subsubsection}{II.1.3.3.13 [CT2] Rehabilitation}

{\small Historical predecessor: R.2.84 \hyperref[element:28ec21d5918a49aaac12740c69f42d77]{[CT2] Rehabilitation}; Canonical ID: \texttt{\seqsplit{28ec21d5918a49aaac12740c69f42d77}}. Superseded historical identity — direct correspondence · superseded.}\par
\label{element:99ed51abcc5e4e8ba3b3d0161b17544a}
\index{[CT2] Rehabilitation (II.1.3.3.13)}
{\small Capability · Canonical · current}\par

Identify and track the operational and infrastructure dependencies necessary for rehabilitation services to function, retaining each dependency provider's responsibility.



\textbf{Overview}



Identify and track the operational and infrastructure dependencies necessary for rehabilitation services to function, retaining each dependency provider's responsibility.



\textbf{Intent}



Identify and track the operational and infrastructure dependencies necessary for rehabilitation services to function, retaining each dependency provider's responsibility.



\textbf{Boundary}



Information administration does not confer authority over the represented person, organisation, service, place, device or clinical activity.



\medskip\hrule\medskip

{\footnotesize \href{../../../../continuum/strategy/capability/99ed51abcc5e4e8ba3b3d0161b17544a.md}{Canonical record: 99ed51abcc5e4e8ba3b3d0161b17544a}; full qualifications and provenance.}\par
\subsubsection*{II.1.3.3.14 [CT2] Remote \& Virtual Care}
\phantomsection
\label{navigation:II.1.3.3.14}
\addcontentsline{toc}{subsubsection}{II.1.3.3.14 [CT2] Remote \& Virtual Care}

{\small Historical predecessor: R.2.85 \hyperref[element:75e2b447659c4fdf94107d94758337ed]{[CT2] Remote \& Virtual Care}; Canonical ID: \texttt{\seqsplit{75e2b447659c4fdf94107d94758337ed}}. Superseded historical identity — direct correspondence · superseded.}\par
\label{element:30dff22d40b64da29ede7c41673db085}
\index{[CT2] Remote \& Virtual Care (II.1.3.3.14)}
{\small Capability · Proposed · current}\par

Identify and track the operational and infrastructure dependencies necessary for remote \& virtual care services to function, retaining each dependency provider's responsibility.



\textbf{Overview}



Identify and track the operational and infrastructure dependencies necessary for remote \& virtual care services to function, retaining each dependency provider's responsibility. The reviewed hierarchy assigns this enduring ability to its stated business-enabling context.



\textbf{Intent}



Identify and track the operational and infrastructure dependencies necessary for remote \& virtual care services to function, retaining each dependency provider's responsibility.



\textbf{Boundary}



Information administration and coordination preserve the authority of the responsible source, clinical actor and governed domain; they confer no additional real-world management authority.


{\footnotesize \href{../../../../continuum/strategy/capability/30dff22d40b64da29ede7c41673db085.md}{Canonical record: 30dff22d40b64da29ede7c41673db085}; full qualifications and provenance.}\par
\subsubsection*{II.1.3.3.15 [CT2] Surgical \& Procedural Care}
\phantomsection
\label{navigation:II.1.3.3.15}
\addcontentsline{toc}{subsubsection}{II.1.3.3.15 [CT2] Surgical \& Procedural Care}

{\small Historical predecessor: R.2.86 \hyperref[element:8bc2a106165a4702b5f72fe46a206b69]{[CT2] Surgical \& Procedural Care}; Canonical ID: \texttt{\seqsplit{8bc2a106165a4702b5f72fe46a206b69}}. Superseded historical identity — direct correspondence · superseded.}\par
\label{element:c9a26d1aa1014545b9ae759b663c807b}
\index{[CT2] Surgical \& Procedural Care (II.1.3.3.15)}
{\small Capability · Canonical · current}\par

Identify and track the operational and infrastructure dependencies necessary for surgical \& procedural care services to function, retaining each dependency provider's responsibility.



\textbf{Overview}



Identify and track the operational and infrastructure dependencies necessary for surgical \& procedural care services to function, retaining each dependency provider's responsibility.



\textbf{Intent}



Identify and track the operational and infrastructure dependencies necessary for surgical \& procedural care services to function, retaining each dependency provider's responsibility.



\textbf{Boundary}



Information administration does not confer authority over the represented person, organisation, service, place, device or clinical activity.



\medskip\hrule\medskip

{\footnotesize \href{../../../../continuum/strategy/capability/c9a26d1aa1014545b9ae759b663c807b.md}{Canonical record: c9a26d1aa1014545b9ae759b663c807b}; full qualifications and provenance.}\par
\subsubsection*{II.1.3.3.16 [CT2] Therapeutic Services}
\phantomsection
\label{navigation:II.1.3.3.16}
\addcontentsline{toc}{subsubsection}{II.1.3.3.16 [CT2] Therapeutic Services}

{\small Historical predecessor: R.2.87 \hyperref[element:54032b52ee444bfe8b611fb9e8e2ba10]{[CT2] Therapeutic Services}; Canonical ID: \texttt{\seqsplit{54032b52ee444bfe8b611fb9e8e2ba10}}. Superseded historical identity — direct correspondence · superseded.}\par
\label{element:b512b6d326004cb7b91a627fe6dbd8be}
\index{[CT2] Therapeutic Services (II.1.3.3.16)}
{\small Capability · Canonical · current}\par

Identify and track the operational and infrastructure dependencies necessary for therapeutic services services to function, retaining each dependency provider's responsibility.



\textbf{Overview}



Identify and track the operational and infrastructure dependencies necessary for therapeutic services services to function, retaining each dependency provider's responsibility.



\textbf{Intent}



Identify and track the operational and infrastructure dependencies necessary for therapeutic services services to function, retaining each dependency provider's responsibility.



\textbf{Boundary}



Information administration does not confer authority over the represented person, organisation, service, place, device or clinical activity.



\medskip\hrule\medskip

{\footnotesize \href{../../../../continuum/strategy/capability/b512b6d326004cb7b91a627fe6dbd8be.md}{Canonical record: b512b6d326004cb7b91a627fe6dbd8be}; full qualifications and provenance.}\par
\subsection*{II.1.3.4 [CT1] Entity Information Enablement}
\phantomsection
\label{navigation:II.1.3.4}
\addcontentsline{toc}{subsection}{II.1.3.4 [CT1] Entity Information Enablement}

{\small Historical predecessor: R.2.114 \hyperref[element:e4d959d6c71d498c9ce00d2be4c048e3]{[CT1] Entity Information Enablement}; Canonical ID: \texttt{\seqsplit{e4d959d6c71d498c9ce00d2be4c048e3}}. Superseded historical identity — direct correspondence · superseded.}\par
\label{element:b5a92874f8ac4509a30af5253590edc1}
\index{[CT1] Entity Information Enablement (II.1.3.4)}
{\small Capability · Proposed · current}\par

Administer identifying, descriptive and qualified relationship information concerning healthcare participants, organisations, locations, services, materiel and devices.



\textbf{Overview}



Administer identifying, descriptive and qualified relationship information concerning healthcare participants, organisations, locations, services, materiel and devices.



\textbf{Intent}



Administer identifying, descriptive and qualified relationship information concerning healthcare participants, organisations, locations, services, materiel and devices.



\textbf{Boundary}



Registry information does not establish clinical privilege, service readiness, occupancy or originating national identity authority.


{\footnotesize \href{../../../../continuum/strategy/capability/b5a92874f8ac4509a30af5253590edc1.md}{Canonical record: b5a92874f8ac4509a30af5253590edc1}; full qualifications and provenance.}\par
\subsubsection*{II.1.3.4.1 [CT2] Client Administration}
\phantomsection
\label{navigation:II.1.3.4.1}
\addcontentsline{toc}{subsubsection}{II.1.3.4.1 [CT2] Client Administration}

{\small Historical predecessor: R.2.93 \hyperref[element:70d2a44a0442441f8b306b82b297bb8a]{[CT2] Client Administration}; Canonical ID: \texttt{\seqsplit{70d2a44a0442441f8b306b82b297bb8a}}. Superseded historical identity — direct correspondence · superseded.}\par
\label{element:f0885472cced4bccbcce1568bf0360dc}
\index{[CT2] Client Administration (II.1.3.4.1)}
{\small Capability · Canonical · current}\par

Client Administration provides the representation and demographic context needed to recognise individuals receiving care. It brings identity, healthcare-subject context, relationships and privacy into the same administrative frame while preserving their different responsibilities. The established partial chain to Person Identity is retained; the catalogue's other contextual descriptions do not supply a complete capability hierarchy.



\textbf{Overview}



Client Administration provides the representation and demographic context needed to recognise individuals receiving care. It brings identity, healthcare-subject context, relationships and privacy into the same administrative frame while preserving their different responsibilities. The established partial chain to Person Identity is retained; the catalogue's other contextual descriptions do not supply a complete capability hierarchy.



\textbf{Intent}



Client Administration provides the representation and demographic context needed to recognise individuals receiving care. It brings identity, healthcare-subject context, relationships and privacy into the same administrative frame while preserving their different responsibilities. The established partial chain to Person Identity is retained; the catalogue's other contextual descriptions do not supply a complete capability hierarchy.



\textbf{Boundary}



Information administration does not confer authority over the represented person, organisation, service, place, device or clinical activity.



\medskip\hrule\medskip

{\small Recorded unresolved matters: source-qualified matters; exact scope remains in the canonical record.}\par
{\footnotesize \href{../../../../continuum/strategy/capability/f0885472cced4bccbcce1568bf0360dc.md}{Canonical record: f0885472cced4bccbcce1568bf0360dc}; full qualifications and provenance.}\par
\paragraph*{II.1.3.4.1.1 [CT3] Client Privacy Guardrail Management}
\phantomsection
\label{navigation:II.1.3.4.1.1}
\addcontentsline{toc}{paragraph}{II.1.3.4.1.1 [CT3] Client Privacy Guardrail Management}

{\small Historical predecessor: R.2.89 \hyperref[element:d2e6078b5db041cb904e0cadcebce6eb]{[CT3] Client Privacy Guardrail Management}; Canonical ID: \texttt{\seqsplit{d2e6078b5db041cb904e0cadcebce6eb}}. Superseded historical identity — direct correspondence · superseded.}\par
\label{element:7577a6e1815243648dd5356c21cca151}
\index{[CT3] Client Privacy Guardrail Management (II.1.3.4.1.1)}
{\small Capability · Canonical · current}\par

Enforcing client-level privacy preferences, consent directives, and confidentiality markers.



Definition — reviewed candidate



Manage client-associated privacy constraints, preferences, permissions, restrictions and consent directives, and evaluate their applicability to information use.



\textbf{Overview}



Manage client-associated privacy constraints, preferences, permissions, restrictions and consent directives, and evaluate their applicability to information use.



\textbf{Intent}



Manage client-associated privacy constraints, preferences, permissions, restrictions and consent directives, and evaluate their applicability to information use.



\textbf{Boundary}



Owns client-associated constraints and consent directives; general enterprise privacy policy remains outside this capability.



\medskip\hrule\medskip

{\footnotesize \href{../../../../continuum/strategy/capability/7577a6e1815243648dd5356c21cca151.md}{Canonical record: 7577a6e1815243648dd5356c21cca151}; full qualifications and provenance.}\par
\subparagraph*{II.1.3.4.1.1.1 [Feature] Confidentiality Flag Enforcement}
\phantomsection
\label{navigation:II.1.3.4.1.1.1}
\addcontentsline{toc}{subparagraph}{II.1.3.4.1.1.1 [Feature] Confidentiality Flag Enforcement}

\label{element:1c18e7c2ccf640c9afc2525503b505a6}
\index{[Feature] Confidentiality Flag Enforcement (II.1.3.4.1.1.1)}
{\small Capability / Feature · Canonical · current}\par

Enforce restricted access flags (e.g., VIP, sensitive diagnosis, restricted record) across query and presentation interfaces.



Confidentiality enforcement makes a restricted-record condition effective across discovery and presentation paths. Its purpose is consistent governed treatment of sensitive information, rather than a flag displayed in one interface while other paths disclose the same content without restriction.


{\footnotesize \href{../../../../continuum/strategy/capability/feature/1c18e7c2ccf640c9afc2525503b505a6.md}{Canonical record: 1c18e7c2ccf640c9afc2525503b505a6}; full qualifications and provenance.}\par
\subparagraph*{II.1.3.4.1.1.2 [Feature] Consent Directive Evaluation}
\phantomsection
\label{navigation:II.1.3.4.1.1.2}
\addcontentsline{toc}{subparagraph}{II.1.3.4.1.1.2 [Feature] Consent Directive Evaluation}

\label{element:baf6dc6ef38a4256a9a67a65fd3f31ba}
\index{[Feature] Consent Directive Evaluation (II.1.3.4.1.1.2)}
{\small Capability / Feature · Canonical · current}\par

Evaluate client consent directives against requested information disclosure and exchange actions.



Consent evaluation brings the client's applicable directives into consideration of a requested disclosure or exchange. It distinguishes the recorded directive from the particular action being evaluated, preserving the need for established policy and authority rather than treating consent data as a universal permission token.


{\footnotesize \href{../../../../continuum/strategy/capability/feature/baf6dc6ef38a4256a9a67a65fd3f31ba.md}{Canonical record: baf6dc6ef38a4256a9a67a65fd3f31ba}; full qualifications and provenance.}\par
\paragraph*{II.1.3.4.1.2 [CT3] Client Relationship Administration}
\phantomsection
\label{navigation:II.1.3.4.1.2}
\addcontentsline{toc}{paragraph}{II.1.3.4.1.2 [CT3] Client Relationship Administration}

{\small Historical predecessor: R.2.90 \hyperref[element:2c799ec0c4c74c97a75c12e949887b95]{[CT3] Client Relationship Administration}; Canonical ID: \texttt{\seqsplit{2c799ec0c4c74c97a75c12e949887b95}}. Superseded historical identity — direct correspondence · superseded.}\par
\label{element:d682c11d3eb04770b1f20be478fd475e}
\index{[CT3] Client Relationship Administration (II.1.3.4.1.2)}
{\small Capability · Canonical · current}\par

Representing relationships between the client and carers, guardians, or contacts.



Definition — reviewed candidate



Administer qualified client relationships with carers, guardians, contacts and other nominated care participants.



\textbf{Overview}



Administer qualified client relationships with carers, guardians, contacts and other nominated care participants.



\textbf{Intent}



Administer qualified client relationships with carers, guardians, contacts and other nominated care participants.



\textbf{Boundary}



Information administration does not confer authority over the represented person, organisation, service, place, device or clinical activity.



\medskip\hrule\medskip

{\footnotesize \href{../../../../continuum/strategy/capability/d682c11d3eb04770b1f20be478fd475e.md}{Canonical record: d682c11d3eb04770b1f20be478fd475e}; full qualifications and provenance.}\par
\subparagraph*{II.1.3.4.1.2.1 [Feature] Next-of-Kin \& Guardian Association}
\phantomsection
\label{navigation:II.1.3.4.1.2.1}
\addcontentsline{toc}{subparagraph}{II.1.3.4.1.2.1 [Feature] Next-of-Kin \& Guardian Association}

\label{element:05de74a8a7dc4334ba60543a228e73cb}
\index{[Feature] Next-of-Kin \& Guardian Association (II.1.3.4.1.2.1)}
{\small Capability / Feature · Canonical · current}\par

Associate nominated emergency contacts, legal guardians, and care agents with the client record.



This Feature records the nominated people and capacities through which a client can be contacted or represented. Association preserves who the contact or guardian is in relation to the client; it does not make every nominated person an unrestricted recipient of clinical information.


{\footnotesize \href{../../../../continuum/strategy/capability/feature/05de74a8a7dc4334ba60543a228e73cb.md}{Canonical record: 05de74a8a7dc4334ba60543a228e73cb}; full qualifications and provenance.}\par
\subparagraph*{II.1.3.4.1.2.2 [Feature] Relationship Validity Tracking}
\phantomsection
\label{navigation:II.1.3.4.1.2.2}
\addcontentsline{toc}{subparagraph}{II.1.3.4.1.2.2 [Feature] Relationship Validity Tracking}

\label{element:309603a8330545abb5db5f9889aa60d9}
\index{[Feature] Relationship Validity Tracking (II.1.3.4.1.2.2)}
{\small Capability / Feature · Canonical · current}\par

Maintain effective date intervals and authority scopes for nominated client relationships.



Relationship validity makes effective periods and authority scope part of the relationship's meaning. A contact association that was legitimate in one period or capacity need not be treated as currently valid for every purpose.


{\footnotesize \href{../../../../continuum/strategy/capability/feature/309603a8330545abb5db5f9889aa60d9.md}{Canonical record: 309603a8330545abb5db5f9889aa60d9}; full qualifications and provenance.}\par
\paragraph*{II.1.3.4.1.3 [CT3] Healthcare Subject Profile Administration}
\phantomsection
\label{navigation:II.1.3.4.1.3}
\addcontentsline{toc}{paragraph}{II.1.3.4.1.3 [CT3] Healthcare Subject Profile Administration}

{\small Historical predecessor: R.2.91 \hyperref[element:d653c3247ddd49479f985f6b435933a0]{[CT3] Healthcare Subject Profile Administration}; Canonical ID: \texttt{\seqsplit{d653c3247ddd49479f985f6b435933a0}}. Superseded historical identity — direct correspondence · superseded.}\par
\label{element:df214b5499694b768eb4a7a175a9d471}
\index{[CT3] Healthcare Subject Profile Administration (II.1.3.4.1.3)}
{\small Capability · Canonical · current}\par

Managing the core healthcare subject context bound to clinical records and operational activities.



Definition — reviewed candidate



Administer demographic and subject-of-care profile information bound to clinical records and operational activities.



\textbf{Overview}



Administer demographic and subject-of-care profile information bound to clinical records and operational activities.



\textbf{Intent}



Administer demographic and subject-of-care profile information bound to clinical records and operational activities.



\textbf{Boundary}



Information administration does not confer authority over the represented person, organisation, service, place, device or clinical activity.



\medskip\hrule\medskip

{\footnotesize \href{../../../../continuum/strategy/capability/df214b5499694b768eb4a7a175a9d471.md}{Canonical record: df214b5499694b768eb4a7a175a9d471}; full qualifications and provenance.}\par
\subparagraph*{II.1.3.4.1.3.1 [Feature] Demographic Context Governance}
\phantomsection
\label{navigation:II.1.3.4.1.3.1}
\addcontentsline{toc}{subparagraph}{II.1.3.4.1.3.1 [Feature] Demographic Context Governance}

\label{element:3d0377451a1b44d0bb41404207cdd867}
\index{[Feature] Demographic Context Governance (II.1.3.4.1.3.1)}
{\small Capability / Feature · Canonical · current}\par

Manage essential demographic attributes (date of birth, sex, vital status) with full provenance and timestamping.



Demographic governance retains the provenance and temporal context of attributes used to interpret the subject of care. A changed date, status or demographic assertion can therefore be understood as sourced information rather than an unexplained overwrite of previously known context.


{\footnotesize \href{../../../../continuum/strategy/capability/feature/3d0377451a1b44d0bb41404207cdd867.md}{Canonical record: 3d0377451a1b44d0bb41404207cdd867}; full qualifications and provenance.}\par
\subparagraph*{II.1.3.4.1.3.2 [Feature] Subject Context Binding}
\phantomsection
\label{navigation:II.1.3.4.1.3.2}
\addcontentsline{toc}{subparagraph}{II.1.3.4.1.3.2 [Feature] Subject Context Binding}

\label{element:5d956689bc714e2d8be4fff7766c0b2d}
\index{[Feature] Subject Context Binding (II.1.3.4.1.3.2)}
{\small Capability / Feature · Canonical · current}\par

Bind authenticated person identity to clinical encounters, orders, and observations as the enduring subject of care with verifiable association.



Subject binding connects resolved identity with the records and activities concerning that person. It makes the association verifiable across encounters, orders and observations rather than assuming that proximity in a message or screen establishes the enduring subject of care.


{\footnotesize \href{../../../../continuum/strategy/capability/feature/5d956689bc714e2d8be4fff7766c0b2d.md}{Canonical record: 5d956689bc714e2d8be4fff7766c0b2d}; full qualifications and provenance.}\par
\paragraph*{II.1.3.4.1.4 [CT3] Person Identity Administration}
\phantomsection
\label{navigation:II.1.3.4.1.4}
\addcontentsline{toc}{paragraph}{II.1.3.4.1.4 [CT3] Person Identity Administration}

{\small Historical predecessor: R.2.92 \hyperref[element:7550fd5b17e643118fa76d2616e1ba20]{[CT3] Person Identity Administration}; Canonical ID: \texttt{\seqsplit{7550fd5b17e643118fa76d2616e1ba20}}. Superseded historical identity — direct correspondence · superseded.}\par
\label{element:6c3de7750ef1436cad6a9533731782b6}
\index{[CT3] Person Identity Administration (II.1.3.4.1.4)}
{\small Capability · Canonical · current}\par

Resolving, maintaining, and correlating identities and identifiers assigned to an individual across different authorities.



Definition — reviewed candidate



Resolve, maintain and correlate identities and identifiers assigned to a person across issuing authorities, retaining source-qualified correction history.



\textbf{Overview}



Resolve, maintain and correlate identities and identifiers assigned to a person across issuing authorities, retaining source-qualified correction history.



\textbf{Intent}



Resolve, maintain and correlate identities and identifiers assigned to a person across issuing authorities, retaining source-qualified correction history.



\textbf{Boundary}



Information administration does not confer authority over the represented person, organisation, service, place, device or clinical activity.



\medskip\hrule\medskip

{\footnotesize \href{../../../../continuum/strategy/capability/6c3de7750ef1436cad6a9533731782b6.md}{Canonical record: 6c3de7750ef1436cad6a9533731782b6}; full qualifications and provenance.}\par
\subparagraph*{II.1.3.4.1.4.1 [Feature] Cross-Authority Identifier Correlation}
\phantomsection
\label{navigation:II.1.3.4.1.4.1}
\addcontentsline{toc}{subparagraph}{II.1.3.4.1.4.1 [Feature] Cross-Authority Identifier Correlation}

\label{element:6d74d143c7454fa0bf30743ddf7ba094}
\index{[Feature] Cross-Authority Identifier Correlation (II.1.3.4.1.4.1)}
{\small Capability / Feature · Canonical · current}\par

Correlate disparate local identifiers (e.g., hospital MRN, clinic client ID, national IHI) belonging to the same individual.



Cross-authority correlation preserves the association between identifiers used for one person by different participants. It enables information to be related across regional systems without erasing the separate identifier namespaces or claiming that Harmonia originated the identity decision.


{\footnotesize \href{../../../../continuum/strategy/capability/feature/6d74d143c7454fa0bf30743ddf7ba094.md}{Canonical record: 6d74d143c7454fa0bf30743ddf7ba094}; full qualifications and provenance.}\par
\subparagraph*{II.1.3.4.1.4.2 [Feature] Governed Identity Correction}
\phantomsection
\label{navigation:II.1.3.4.1.4.2}
\addcontentsline{toc}{subparagraph}{II.1.3.4.1.4.2 [Feature] Governed Identity Correction}

\label{element:f7ab7a24e8234129b2ab02e9abd4fba2}
\index{[Feature] Governed Identity Correction (II.1.3.4.1.4.2)}
{\small Capability / Feature · Canonical · current}\par

Apply administrative identity updates, linked-record corrections, and split/merge notifications propagated from authoritative sources.



Governed correction applies externally authoritative identity updates and their consequences to Harmonia-managed associations. Administrative review concerns the legitimate application of the source outcome; it does not grant Harmonia an EMPI merge engine or authority to select a golden person record.


{\footnotesize \href{../../../../continuum/strategy/capability/feature/f7ab7a24e8234129b2ab02e9abd4fba2.md}{Canonical record: f7ab7a24e8234129b2ab02e9abd4fba2}; full qualifications and provenance.}\par
\subparagraph*{II.1.3.4.1.4.3 [Feature] Identifier Capture}
\phantomsection
\label{navigation:II.1.3.4.1.4.3}
\addcontentsline{toc}{subparagraph}{II.1.3.4.1.4.3 [Feature] Identifier Capture}

\label{element:3cae3c4810974c5f8ee5a4c942f38b07}
\index{[Feature] Identifier Capture (II.1.3.4.1.4.3)}
{\small Capability / Feature · Accepted bounded reconciliation · current}\par

The governed source has no established definition. This absence is retained for Solution Architect review; no definition is inferred from the name.


{\small Recorded unresolved matters: definition; exact scope remains in the canonical record.}\par
{\footnotesize \href{../../../../continuum/strategy/capability/feature/3cae3c4810974c5f8ee5a4c942f38b07.md}{Canonical record: 3cae3c4810974c5f8ee5a4c942f38b07}; full qualifications and provenance.}\par
\subparagraph*{II.1.3.4.1.4.4 [Feature] Identifier Resolution}
\phantomsection
\label{navigation:II.1.3.4.1.4.4}
\addcontentsline{toc}{subparagraph}{II.1.3.4.1.4.4 [Feature] Identifier Resolution}

\label{element:ae034169e1394ccdb3a2ce9651eaf305}
\index{[Feature] Identifier Resolution (II.1.3.4.1.4.4)}
{\small Capability / Feature · Canonical · current}\par

Ingest and resolve external identifiers against registered issuing authorities and namespace systems.



Identifier Resolution interprets an identifier in its issuing-authority and namespace context, so that the same local value is not assumed to identify the same person everywhere. It provides an identifiable subject reference without deciding an authoritative master-person match.


{\footnotesize \href{../../../../continuum/strategy/capability/feature/ae034169e1394ccdb3a2ce9651eaf305.md}{Canonical record: ae034169e1394ccdb3a2ce9651eaf305}; full qualifications and provenance.}\par
\subparagraph*{II.1.3.4.1.4.5 [Feature] Identity Alias Association}
\phantomsection
\label{navigation:II.1.3.4.1.4.5}
\addcontentsline{toc}{subparagraph}{II.1.3.4.1.4.5 [Feature] Identity Alias Association}

\label{element:97721dbd017e4c2daa73e1e87a8499b9}
\index{[Feature] Identity Alias Association (II.1.3.4.1.4.5)}
{\small Capability / Feature · Canonical · current}\par

Maintain explicit associations between active identifiers and historical or demographic aliases without destructive overwriting.



Alias association keeps historical and demographic identity references navigable alongside active identifiers. The non-destructive relationship preserves how earlier information identified the person rather than rewriting that history whenever the current identity representation changes.


{\footnotesize \href{../../../../continuum/strategy/capability/feature/97721dbd017e4c2daa73e1e87a8499b9.md}{Canonical record: 97721dbd017e4c2daa73e1e87a8499b9}; full qualifications and provenance.}\par
\subparagraph*{II.1.3.4.1.4.6 [Feature] Identity Source \& Correction Lineage}
\phantomsection
\label{navigation:II.1.3.4.1.4.6}
\addcontentsline{toc}{subparagraph}{II.1.3.4.1.4.6 [Feature] Identity Source \& Correction Lineage}

\label{element:d72134e75001449a9e7ffae53f61458f}
\index{[Feature] Identity Source \& Correction Lineage (II.1.3.4.1.4.6)}
{\small Capability / Feature · Proposed · current}\par

Retain source-attributed identity assertions and externally decided correction history.



\textbf{Overview}



Identifier correlations can change when an upstream authority corrects a relationship. Reviewers must be able to distinguish the source decision, its effective standing and the managed links affected.



\textbf{Boundary}



No probabilistic matching, golden-record selection or autonomous merge authority.


{\footnotesize \href{../../../../continuum/strategy/capability/feature/d72134e75001449a9e7ffae53f61458f.md}{Canonical record: d72134e75001449a9e7ffae53f61458f}; full qualifications and provenance.}\par
\subsubsection*{II.1.3.4.2 [CT2] Clinical Device Administration}
\phantomsection
\label{navigation:II.1.3.4.2}
\addcontentsline{toc}{subsubsection}{II.1.3.4.2 [CT2] Clinical Device Administration}

{\small Historical predecessor: R.2.98 \hyperref[element:3792be10541e4d68b3f740b1ba964966]{[CT2] Clinical Device Administration}; Canonical ID: \texttt{\seqsplit{3792be10541e4d68b3f740b1ba964966}}. Superseded historical identity — direct correspondence · superseded.}\par
\label{element:31bcd6fda0d14cc8b992cdca39a3cbed}
\index{[CT2] Clinical Device Administration (II.1.3.4.2)}
{\small Capability · Canonical · current}\par

Clinical Device Administration represents device types, identifiable instances, patient or location associations and communication endpoints. This lets device-originated information be interpreted in relation to the actual device and care context, rather than merely a network address. Type, instance, assignment and connectivity each answer a different question about the source of device data.



\textbf{Overview}



Clinical Device Administration represents device types, identifiable instances, patient or location associations and communication endpoints. This lets device-originated information be interpreted in relation to the actual device and care context, rather than merely a network address. Type, instance, assignment and connectivity each answer a different question about the source of device data.



\textbf{Intent}



Clinical Device Administration represents device types, identifiable instances, patient or location associations and communication endpoints. This lets device-originated information be interpreted in relation to the actual device and care context, rather than merely a network address. Type, instance, assignment and connectivity each answer a different question about the source of device data.



\textbf{Boundary}



Information administration does not confer authority over the represented person, organisation, service, place, device or clinical activity.



\medskip\hrule\medskip

{\small Recorded unresolved matters: source-qualified matters; exact scope remains in the canonical record.}\par
{\footnotesize \href{../../../../continuum/strategy/capability/31bcd6fda0d14cc8b992cdca39a3cbed.md}{Canonical record: 31bcd6fda0d14cc8b992cdca39a3cbed}; full qualifications and provenance.}\par
\paragraph*{II.1.3.4.2.1 [CT3] Clinical Device Profile Administration}
\phantomsection
\label{navigation:II.1.3.4.2.1}
\addcontentsline{toc}{paragraph}{II.1.3.4.2.1 [CT3] Clinical Device Profile Administration}

{\small Historical predecessor: R.2.94 \hyperref[element:e8c72952da144d6883caba7fd2afb88d]{[CT3] Clinical Device Profile Administration}; Canonical ID: \texttt{\seqsplit{e8c72952da144d6883caba7fd2afb88d}}. Superseded historical identity — direct correspondence · superseded.}\par
\label{element:d0a1ba80c790476ea6042290eb7ea09e}
\index{[CT3] Clinical Device Profile Administration (II.1.3.4.2.1)}
{\small Capability · Proposed · current}\par

Administer device type, instance, identifying and descriptive information needed to recognise clinical equipment and its stated capabilities.



\textbf{Overview}



Administer device type, instance, identifying and descriptive information needed to recognise clinical equipment and its stated capabilities.



\textbf{Intent}



Administer device type, instance, identifying and descriptive information needed to recognise clinical equipment and its stated capabilities.



\textbf{Boundary}



Describes device types and capabilities without taking regulatory approval or clinical-use authority.


{\footnotesize \href{../../../../continuum/strategy/capability/d0a1ba80c790476ea6042290eb7ea09e.md}{Canonical record: d0a1ba80c790476ea6042290eb7ea09e}; full qualifications and provenance.}\par
\subparagraph*{II.1.3.4.2.1.1 [Feature] Device Serial Tracking}
\phantomsection
\label{navigation:II.1.3.4.2.1.1}
\addcontentsline{toc}{subparagraph}{II.1.3.4.2.1.1 [Feature] Device Serial Tracking}

\label{element:df3dae0cf3044059875dbb888ea6a9ba}
\index{[Feature] Device Serial Tracking (II.1.3.4.2.1.1)}
{\small Capability / Feature · Canonical · current}\par

Ingest and maintain unique device identifier (UDI) records and serial numbers.



Serial tracking makes individual hardware distinguishable within the device estate. A unique device reference supports subsequent assignment and source interpretation, while remaining separate from evidence that the device is maintained, connected or fit for a particular clinical use.


{\footnotesize \href{../../../../continuum/strategy/capability/feature/df3dae0cf3044059875dbb888ea6a9ba.md}{Canonical record: df3dae0cf3044059875dbb888ea6a9ba}; full qualifications and provenance.}\par
\subparagraph*{II.1.3.4.2.1.2 [Feature] Device Type Specification}
\phantomsection
\label{navigation:II.1.3.4.2.1.2}
\addcontentsline{toc}{subparagraph}{II.1.3.4.2.1.2 [Feature] Device Type Specification}

\label{element:ca6d91b03b56415790810846649a3dca}
\index{[Feature] Device Type Specification (II.1.3.4.2.1.2)}
{\small Capability / Feature · Canonical · current}\par

Catalog medical device models, regulatory approvals, and capabilities.



Device type specification provides model-level capabilities and regulatory context against which individual devices can be understood. It describes the kind of equipment without substituting for the release, condition or deployment of a particular instance.


{\footnotesize \href{../../../../continuum/strategy/capability/feature/ca6d91b03b56415790810846649a3dca.md}{Canonical record: ca6d91b03b56415790810846649a3dca}; full qualifications and provenance.}\par
\paragraph*{II.1.3.4.2.2 [CT3] Device Association [superseded]}
\phantomsection
\label{navigation:II.1.3.4.2.2}
\addcontentsline{toc}{paragraph}{II.1.3.4.2.2 [CT3] Device Association [superseded]}

\label{element:902daf2ce9714876bbfee52eef1a7ab0}
\index{[CT3] Device Association (II.1.3.4.2.2)}
{\small Capability · Proposed · superseded}\par
\href{../../../../continuum/strategy/capability/902daf2ce9714876bbfee52eef1a7ab0.md}{Historical canonical definition and provenance}.\par
\paragraph*{II.1.3.4.2.3 [CT3] Device Connectivity [superseded]}
\phantomsection
\label{navigation:II.1.3.4.2.3}
\addcontentsline{toc}{paragraph}{II.1.3.4.2.3 [CT3] Device Connectivity [superseded]}

\label{element:658d4b6770a943bcb5801fb7ca9b47be}
\index{[CT3] Device Connectivity (II.1.3.4.2.3)}
{\small Capability · Proposed · superseded}\par
\href{../../../../continuum/strategy/capability/658d4b6770a943bcb5801fb7ca9b47be.md}{Historical canonical definition and provenance}.\par
\paragraph*{II.1.3.4.2.4 [CT3] Device Connectivity Compliance Management}
\phantomsection
\label{navigation:II.1.3.4.2.4}
\addcontentsline{toc}{paragraph}{II.1.3.4.2.4 [CT3] Device Connectivity Compliance Management}

{\small Historical predecessor: R.2.95 \hyperref[element:9225f083d3474715a75af2a827a4016d]{[CT3] Device Connectivity Compliance Management}; Canonical ID: \texttt{\seqsplit{9225f083d3474715a75af2a827a4016d}}. Superseded historical identity — direct correspondence · superseded.}\par
\label{element:b16ab359030a4583868f3aac2221074e}
\index{[CT3] Device Connectivity Compliance Management (II.1.3.4.2.4)}
{\small Capability · Proposed · current}\par

Manage information about whether device connectivity satisfies applicable standards, policies and integration constraints.



\textbf{Overview}



Manage information about whether device connectivity satisfies applicable standards, policies and integration constraints. The reviewed hierarchy assigns this enduring ability to its stated business-enabling context.



\textbf{Intent}



Manage information about whether device connectivity satisfies applicable standards, policies and integration constraints.



\textbf{Boundary}



Endpoint registration alone does not establish compliance; conformity concerns applicable standards, policies and integration constraints.


{\footnotesize \href{../../../../continuum/strategy/capability/b16ab359030a4583868f3aac2221074e.md}{Canonical record: b16ab359030a4583868f3aac2221074e}; full qualifications and provenance.}\par
\paragraph*{II.1.3.4.2.5 [CT3] Device Instance [superseded]}
\phantomsection
\label{navigation:II.1.3.4.2.5}
\addcontentsline{toc}{paragraph}{II.1.3.4.2.5 [CT3] Device Instance [superseded]}

\label{element:dd860926038f41e082c56e38cbb5ab95}
\index{[CT3] Device Instance (II.1.3.4.2.5)}
{\small Capability · Proposed · superseded}\par
\href{../../../../continuum/strategy/capability/dd860926038f41e082c56e38cbb5ab95.md}{Historical canonical definition and provenance}.\par
\paragraph*{II.1.3.4.2.6 [CT3] Device OA\&M Metadata Management}
\phantomsection
\label{navigation:II.1.3.4.2.6}
\addcontentsline{toc}{paragraph}{II.1.3.4.2.6 [CT3] Device OA\&M Metadata Management}

{\small Historical predecessor: R.2.96 \hyperref[element:a398c72bef094dee96c96d97994449fb]{[CT3] Device OA\&M Metadata Management}; Canonical ID: \texttt{\seqsplit{a398c72bef094dee96c96d97994449fb}}. Superseded historical identity — direct correspondence · superseded.}\par
\label{element:659bb194a7ef42ab9ba8a938adb5c179}
\index{[CT3] Device OA\&M Metadata Management (II.1.3.4.2.6)}
{\small Capability · Proposed · current}\par

Manage operational and administrative metadata needed to understand how a clinical device is connected, monitored and administered.



\textbf{Overview}



Manage operational and administrative metadata needed to understand how a clinical device is connected, monitored and administered. The reviewed hierarchy assigns this enduring ability to its stated business-enabling context.



\textbf{Intent}



Manage operational and administrative metadata needed to understand how a clinical device is connected, monitored and administered.



\textbf{Boundary}



Connectivity, monitoring and administration metadata do not perform device maintenance, interpret clinical telemetry or determine conformance.


{\footnotesize \href{../../../../continuum/strategy/capability/659bb194a7ef42ab9ba8a938adb5c179.md}{Canonical record: 659bb194a7ef42ab9ba8a938adb5c179}; full qualifications and provenance.}\par
\subparagraph*{II.1.3.4.2.6.1 [Feature] Device Endpoint Registration}
\phantomsection
\label{navigation:II.1.3.4.2.6.1}
\addcontentsline{toc}{subparagraph}{II.1.3.4.2.6.1 [Feature] Device Endpoint Registration}

\label{element:d32c856559c54b41941fc51fbd019195}
\index{[Feature] Device Endpoint Registration (II.1.3.4.2.6.1)}
{\small Capability / Feature · Canonical · current}\par

Register communication interfaces and security credentials for diagnostic telemetry feeds.



Endpoint registration supplies the interfaces and credential context through which a diagnostic or telemetry feed can participate. It describes provisioned connectivity rather than proving that data has arrived, been clinically interpreted or acquired authoritative standing.


{\footnotesize \href{../../../../continuum/strategy/capability/feature/d32c856559c54b41941fc51fbd019195.md}{Canonical record: d32c856559c54b41941fc51fbd019195}; full qualifications and provenance.}\par
\paragraph*{II.1.3.4.2.7 [CT3] Device Utilisation Administration}
\phantomsection
\label{navigation:II.1.3.4.2.7}
\addcontentsline{toc}{paragraph}{II.1.3.4.2.7 [CT3] Device Utilisation Administration}

{\small Historical predecessor: R.2.97 \hyperref[element:a0af0a86667d48f4ab95ed2920e878cb]{[CT3] Device Utilisation Administration}; Canonical ID: \texttt{\seqsplit{a0af0a86667d48f4ab95ed2920e878cb}}. Superseded historical identity — direct correspondence · superseded.}\par
\label{element:1a3732b5911c45b3aece38b2db2d0f15}
\index{[CT3] Device Utilisation Administration (II.1.3.4.2.7)}
{\small Capability · Proposed · current}\par

Administer qualified device-to-subject/place allocation, location, utilisation and operational readiness information, including the time at which an association applies.



\textbf{Overview}



Administer qualified device-to-subject/place allocation, location, utilisation and operational readiness information, including the time at which an association applies. The reviewed hierarchy assigns this enduring ability to its stated business-enabling context.



\textbf{Intent}



Administer qualified device-to-subject/place allocation, location, utilisation and operational readiness information, including the time at which an association applies.



\textbf{Boundary}



Information administration and coordination preserve the authority of the responsible source, clinical actor and governed domain; they confer no additional real-world management authority.


{\footnotesize \href{../../../../continuum/strategy/capability/1a3732b5911c45b3aece38b2db2d0f15.md}{Canonical record: 1a3732b5911c45b3aece38b2db2d0f15}; full qualifications and provenance.}\par
\subparagraph*{II.1.3.4.2.7.1 [Feature] Device Association Lineage}
\phantomsection
\label{navigation:II.1.3.4.2.7.1}
\addcontentsline{toc}{subparagraph}{II.1.3.4.2.7.1 [Feature] Device Association Lineage}

\label{element:849889c1a9e4495c86a6b32e49c449e3}
\index{[Feature] Device Association Lineage (II.1.3.4.2.7.1)}
{\small Capability / Feature · Proposed · current}\par

Preserve the attributed history and effective basis of device-to-subject/place associations.



\textbf{Overview}



A telemetry record needs to be interpreted against the association applicable when it arose, rather than the device current location alone.



\textbf{Boundary}



Association history does not interpret physiological observations or prove clinical correctness.


{\footnotesize \href{../../../../continuum/strategy/capability/feature/849889c1a9e4495c86a6b32e49c449e3.md}{Canonical record: 849889c1a9e4495c86a6b32e49c449e3}; full qualifications and provenance.}\par
\subparagraph*{II.1.3.4.2.7.2 [Feature] Device-Patient/Location Assignment}
\phantomsection
\label{navigation:II.1.3.4.2.7.2}
\addcontentsline{toc}{subparagraph}{II.1.3.4.2.7.2 [Feature] Device-Patient/Location Assignment}

\label{element:891ed75bdd3c490296e5f0b7cacf6c87}
\index{[Feature] Device-Patient/Location Assignment (II.1.3.4.2.7.2)}
{\small Capability / Feature · Canonical · current}\par

Track dynamic association between a telemetry monitor, patient, and physical bed.



Dynamic device association preserves which patient or physical place a device is linked to as assignments change. This is essential context for interpreting telemetry: a stable device identity does not imply that every observation concerns the same patient over time.


{\footnotesize \href{../../../../continuum/strategy/capability/feature/891ed75bdd3c490296e5f0b7cacf6c87.md}{Canonical record: 891ed75bdd3c490296e5f0b7cacf6c87}; full qualifications and provenance.}\par
\subparagraph*{II.1.3.4.2.7.3 [Feature] Mobile Clinical Equipment Tracking}
\phantomsection
\label{navigation:II.1.3.4.2.7.3}
\addcontentsline{toc}{subparagraph}{II.1.3.4.2.7.3 [Feature] Mobile Clinical Equipment Tracking}

\label{element:a013a8f41ce245f48aed8d9fe33615c5}
\index{[Feature] Mobile Clinical Equipment Tracking (II.1.3.4.2.7.3)}
{\small Capability / Feature · Canonical · current}\par

Track physical location and allocation status of shared clinical equipment (infusion pumps, ventilators, mobile telemetry).



Equipment tracking provides the dynamic location and allocation context of shared devices. It distinguishes an identifiable asset from one available for the present operational need, while leaving maintenance, ownership and clinical authorisation with their respective responsibilities.


{\footnotesize \href{../../../../continuum/strategy/capability/feature/a013a8f41ce245f48aed8d9fe33615c5.md}{Canonical record: a013a8f41ce245f48aed8d9fe33615c5}; full qualifications and provenance.}\par
\subsubsection*{II.1.3.4.3 [CT2] Health Materiel Administration}
\phantomsection
\label{navigation:II.1.3.4.3}
\addcontentsline{toc}{subsubsection}{II.1.3.4.3 [CT2] Health Materiel Administration}

{\small Historical predecessor: R.2.102 \hyperref[element:4952f86078c54a8d898ac992e6f4f577]{[CT2] Health Materiel Administration}; Canonical ID: \texttt{\seqsplit{4952f86078c54a8d898ac992e6f4f577}}. Superseded historical identity — direct correspondence · superseded.}\par
\label{element:e7081070a9f8444f81c1d8fce37ef664}
\index{[CT2] Health Materiel Administration (II.1.3.4.3)}
{\small Capability · Canonical · current}\par

Health Product Administration retains the reference context for medications, implants, blood products and consumables. Product definitions, formularies and batch or serial information help explain what is being referenced without assigning a supply-chain implementation to Harmonia. No Feature decomposition is established for this responsibility in the present baseline.



\textbf{Overview}



Health Product Administration retains the reference context for medications, implants, blood products and consumables. Product definitions, formularies and batch or serial information help explain what is being referenced without assigning a supply-chain implementation to Harmonia. No Feature decomposition is established for this responsibility in the present baseline.



\textbf{Intent}



Health Product Administration retains the reference context for medications, implants, blood products and consumables. Product definitions, formularies and batch or serial information help explain what is being referenced without assigning a supply-chain implementation to Harmonia. No Feature decomposition is established for this responsibility in the present baseline.



\textbf{Boundary}



Information administration does not confer authority over the represented person, organisation, service, place, device or clinical activity.



\medskip\hrule\medskip

{\small Recorded unresolved matters: source-qualified matters; exact scope remains in the canonical record.}\par
{\footnotesize \href{../../../../continuum/strategy/capability/e7081070a9f8444f81c1d8fce37ef664.md}{Canonical record: e7081070a9f8444f81c1d8fce37ef664}; full qualifications and provenance.}\par
\paragraph*{II.1.3.4.3.1 [CT3] Health Materiel Catalogue \& Formulary Administration}
\phantomsection
\label{navigation:II.1.3.4.3.1}
\addcontentsline{toc}{paragraph}{II.1.3.4.3.1 [CT3] Health Materiel Catalogue \& Formulary Administration}

{\small Historical predecessor: R.2.99 \hyperref[element:5c691959e2cc4a289fc25b9f3497eed9]{[CT3] Health Materiel Catalogue \& Formulary Administration}; Canonical ID: \texttt{\seqsplit{5c691959e2cc4a289fc25b9f3497eed9}}. Superseded historical identity — direct correspondence · superseded.}\par
\label{element:28b327760af94b8b986357469524d793}
\index{[CT3] Health Materiel Catalogue \& Formulary Administration (II.1.3.4.3.1)}
{\small Capability · Proposed · current}\par

Administer health-materiel catalogue and externally approved formulary information, preserving approval and version standing.



\textbf{Overview}



Administer health-materiel catalogue and externally approved formulary information, preserving approval and version standing.



\textbf{Intent}



Administer health-materiel catalogue and externally approved formulary information, preserving approval and version standing.



\textbf{Boundary}



Represents externally approved formularies; it does not acquire formulary approval or dispensing authority.


{\footnotesize \href{../../../../continuum/strategy/capability/28b327760af94b8b986357469524d793.md}{Canonical record: 28b327760af94b8b986357469524d793}; full qualifications and provenance.}\par
\paragraph*{II.1.3.4.3.2 [CT3] Health Materiel Profile Administration}
\phantomsection
\label{navigation:II.1.3.4.3.2}
\addcontentsline{toc}{paragraph}{II.1.3.4.3.2 [CT3] Health Materiel Profile Administration}

{\small Historical predecessor: R.2.100 \hyperref[element:9216f9cc905e4526b1d98cd4ab382cbf]{[CT3] Health Materiel Profile Administration}; Canonical ID: \texttt{\seqsplit{9216f9cc905e4526b1d98cd4ab382cbf}}. Superseded historical identity — direct correspondence · superseded.}\par
\label{element:9626edae517d40d2bfae619860f18e20}
\index{[CT3] Health Materiel Profile Administration (II.1.3.4.3.2)}
{\small Capability · Proposed · current}\par

Administer identifying and descriptive information about health materiel referenced in healthcare information and activity.



\textbf{Overview}



Administer identifying and descriptive information about health materiel referenced in healthcare information and activity.



\textbf{Intent}



Administer identifying and descriptive information about health materiel referenced in healthcare information and activity.



\textbf{Boundary}



Reference representation is distinct from procurement, manufacturing and prescribing.


{\footnotesize \href{../../../../continuum/strategy/capability/9626edae517d40d2bfae619860f18e20.md}{Canonical record: 9626edae517d40d2bfae619860f18e20}; full qualifications and provenance.}\par
\paragraph*{II.1.3.4.3.3 [CT3] Health Materiel Traceability Administration}
\phantomsection
\label{navigation:II.1.3.4.3.3}
\addcontentsline{toc}{paragraph}{II.1.3.4.3.3 [CT3] Health Materiel Traceability Administration}

{\small Historical predecessor: R.2.101 \hyperref[element:f1b888e8d50c4019ac1afe162d6223f3]{[CT3] Health Materiel Traceability Administration}; Canonical ID: \texttt{\seqsplit{f1b888e8d50c4019ac1afe162d6223f3}}. Superseded historical identity — direct correspondence · superseded.}\par
\label{element:2d41ed0f00a94a4b90ffe87d6fd4e9fe}
\index{[CT3] Health Materiel Traceability Administration (II.1.3.4.3.3)}
{\small Capability · Proposed · current}\par

Administer batch/lot, serial, supplier, manufacturer, receipt/delivery and relevant movement/provenance information concerning health materiel.



\textbf{Overview}



Administer batch/lot, serial, supplier, manufacturer, receipt/delivery and relevant movement/provenance information concerning health materiel.



\textbf{Intent}



Administer batch/lot, serial, supplier, manufacturer, receipt/delivery and relevant movement/provenance information concerning health materiel.



\textbf{Boundary}



Materiel traceability does not make Harmonia a procurement, stock-control or supply-chain manager; specimen movement remains Clinical Logistics Coordination.


{\footnotesize \href{../../../../continuum/strategy/capability/2d41ed0f00a94a4b90ffe87d6fd4e9fe.md}{Canonical record: 2d41ed0f00a94a4b90ffe87d6fd4e9fe}; full qualifications and provenance.}\par
\subsubsection*{II.1.3.4.4 [CT2] Health Service Administration}
\phantomsection
\label{navigation:II.1.3.4.4}
\addcontentsline{toc}{subsubsection}{II.1.3.4.4 [CT2] Health Service Administration}

{\small Historical predecessor: R.2.104 \hyperref[element:f1ab1258cf2944efb50105e7885f8027]{[CT2] Health Service Administration}; Canonical ID: \texttt{\seqsplit{f1ab1258cf2944efb50105e7885f8027}}. Superseded historical identity — direct correspondence · superseded.}\par
\label{element:ba68a52913a0499d8377553bddfa74a4}
\index{[CT2] Health Service Administration (II.1.3.4.4)}
{\small Capability · Canonical · current}\par

Health Service Administration represents identifiable services and the conditions under which organisations, locations and practitioner roles provide them. Catalogue publishing, service-location-provider association, operating schedules and eligibility definitions allow a service to be understood as more than a location or specialty label. The service context remains material wherever activity meaning or accountability depends on it.



\textbf{Overview}



Health Service Administration represents identifiable services and the conditions under which organisations, locations and practitioner roles provide them. Catalogue publishing, service-location-provider association, operating schedules and eligibility definitions allow a service to be understood as more than a location or specialty label. The service context remains material wherever activity meaning or accountability depends on it.



\textbf{Intent}



Health Service Administration represents identifiable services and the conditions under which organisations, locations and practitioner roles provide them. Catalogue publishing, service-location-provider association, operating schedules and eligibility definitions allow a service to be understood as more than a location or specialty label. The service context remains material wherever activity meaning or accountability depends on it.



\textbf{Boundary}



Information administration does not confer authority over the represented person, organisation, service, place, device or clinical activity.



\medskip\hrule\medskip

{\small Recorded unresolved matters: source-qualified matters; exact scope remains in the canonical record.}\par
{\footnotesize \href{../../../../continuum/strategy/capability/ba68a52913a0499d8377553bddfa74a4.md}{Canonical record: ba68a52913a0499d8377553bddfa74a4}; full qualifications and provenance.}\par
\paragraph*{II.1.3.4.4.1 [CT3] Healthcare Service Profile Administration}
\phantomsection
\label{navigation:II.1.3.4.4.1}
\addcontentsline{toc}{paragraph}{II.1.3.4.4.1 [CT3] Healthcare Service Profile Administration}

{\small Historical predecessor: R.2.103 \hyperref[element:7780ca2fed2346439d5f559d00cc4ee4]{[CT3] Healthcare Service Profile Administration}; Canonical ID: \texttt{\seqsplit{7780ca2fed2346439d5f559d00cc4ee4}}. Superseded historical identity — direct correspondence · superseded.}\par
\label{element:e75d5eb30936476ea7a9cbadfa43d5ef}
\index{[CT3] Healthcare Service Profile Administration (II.1.3.4.4.1)}
{\small Capability · Proposed · current}\par

Administer service definitions, provision relationships, defined availability and eligibility information describing offered healthcare services.



\textbf{Overview}



Administer service definitions, provision relationships, defined availability and eligibility information describing offered healthcare services.



\textbf{Intent}



Administer service definitions, provision relationships, defined availability and eligibility information describing offered healthcare services.



\textbf{Boundary}



Defined availability is not dynamic operational capacity, collective readiness or authority to accept clinical care.


{\footnotesize \href{../../../../continuum/strategy/capability/e75d5eb30936476ea7a9cbadfa43d5ef.md}{Canonical record: e75d5eb30936476ea7a9cbadfa43d5ef}; full qualifications and provenance.}\par
\subparagraph*{II.1.3.4.4.1.1 [Feature] Eligibility \& Catchment Specification}
\phantomsection
\label{navigation:II.1.3.4.4.1.1}
\addcontentsline{toc}{subparagraph}{II.1.3.4.4.1.1 [Feature] Eligibility \& Catchment Specification}

\label{element:c65feabade3d47c39e7d7af9bc16ebdc}
\index{[Feature] Eligibility \& Catchment Specification (II.1.3.4.4.1.1)}
{\small Capability / Feature · Canonical · current}\par

Define clinical referral thresholds and geographic catchment boundaries.



Eligibility specification captures the clinical thresholds and geographical bounds applicable to service access. It provides the criteria context needed to consider a request without making publication of those criteria an individual eligibility decision or a completed allocation.


{\footnotesize \href{../../../../continuum/strategy/capability/feature/c65feabade3d47c39e7d7af9bc16ebdc.md}{Canonical record: c65feabade3d47c39e7d7af9bc16ebdc}; full qualifications and provenance.}\par
\subparagraph*{II.1.3.4.4.1.2 [Feature] Operating Schedule Specification}
\phantomsection
\label{navigation:II.1.3.4.4.1.2}
\addcontentsline{toc}{subparagraph}{II.1.3.4.4.1.2 [Feature] Operating Schedule Specification}

\label{element:4753953665d343c5a97bc4965be3f615}
\index{[Feature] Operating Schedule Specification (II.1.3.4.4.1.2)}
{\small Capability / Feature · Canonical · current}\par

Define temporal availability, clinic operating hours, and on-call service windows.



Operating schedules make temporal availability and exceptions explicit for a service. They distinguish an existing offering from one available at the relevant time, while leaving actual resource readiness and operational capacity as separate concerns.


{\footnotesize \href{../../../../continuum/strategy/capability/feature/4753953665d343c5a97bc4965be3f615.md}{Canonical record: 4753953665d343c5a97bc4965be3f615}; full qualifications and provenance.}\par
\subparagraph*{II.1.3.4.4.1.3 [Feature] Service Catalogue Publishing}
\phantomsection
\label{navigation:II.1.3.4.4.1.3}
\addcontentsline{toc}{subparagraph}{II.1.3.4.4.1.3 [Feature] Service Catalogue Publishing}

\label{element:d397e670ae524f399e91176c64390fa7}
\index{[Feature] Service Catalogue Publishing (II.1.3.4.4.1.3)}
{\small Capability / Feature · Canonical · current}\par

Publish structured healthcare service definitions, clinical specialties, and intake criteria to regional directories.



Service catalogue publishing makes identifiable service definitions and intake context available to regional participants. A published service can be discovered as an offering rather than inferred from an organisation's name; publication alone does not establish present availability or collective readiness.


{\footnotesize \href{../../../../continuum/strategy/capability/feature/d397e670ae524f399e91176c64390fa7.md}{Canonical record: d397e670ae524f399e91176c64390fa7}; full qualifications and provenance.}\par
\subparagraph*{II.1.3.4.4.1.4 [Feature] Service-Location-Provider Binding}
\phantomsection
\label{navigation:II.1.3.4.4.1.4}
\addcontentsline{toc}{subparagraph}{II.1.3.4.4.1.4 [Feature] Service-Location-Provider Binding}

\label{element:37188bb916244f41b1b7d988ad52a2eb}
\index{[Feature] Service-Location-Provider Binding (II.1.3.4.4.1.4)}
{\small Capability / Feature · Canonical · current}\par

Correlate which organization provides a named health service at which physical location with which practitioner roles.



Service-location-provider binding explains which participant provides an identifiable service in which place and acting capacities. The mapping associations belong to Health Service Administration, while the organisations, locations and practitioner entities retain their own information responsibility.


{\footnotesize \href{../../../../continuum/strategy/capability/feature/37188bb916244f41b1b7d988ad52a2eb.md}{Canonical record: 37188bb916244f41b1b7d988ad52a2eb}; full qualifications and provenance.}\par
\paragraph*{II.1.3.4.4.2 [CT3] Service Availability Definition [superseded]}
\phantomsection
\label{navigation:II.1.3.4.4.2}
\addcontentsline{toc}{paragraph}{II.1.3.4.4.2 [CT3] Service Availability Definition [superseded]}

\label{element:9e5329a42f964199b4b0a87ab1cacc2e}
\index{[CT3] Service Availability Definition (II.1.3.4.4.2)}
{\small Capability · Proposed · superseded}\par
\href{../../../../continuum/strategy/capability/9e5329a42f964199b4b0a87ab1cacc2e.md}{Historical canonical definition and provenance}.\par
\paragraph*{II.1.3.4.4.3 [CT3] Service Eligibility Definition [superseded]}
\phantomsection
\label{navigation:II.1.3.4.4.3}
\addcontentsline{toc}{paragraph}{II.1.3.4.4.3 [CT3] Service Eligibility Definition [superseded]}

\label{element:8a3f29e5948141e2a3d0c72539759daf}
\index{[CT3] Service Eligibility Definition (II.1.3.4.4.3)}
{\small Capability · Proposed · superseded}\par
\href{../../../../continuum/strategy/capability/8a3f29e5948141e2a3d0c72539759daf.md}{Historical canonical definition and provenance}.\par
\paragraph*{II.1.3.4.4.4 [CT3] Service Provision Relationship [superseded]}
\phantomsection
\label{navigation:II.1.3.4.4.4}
\addcontentsline{toc}{paragraph}{II.1.3.4.4.4 [CT3] Service Provision Relationship [superseded]}

\label{element:718ebe1fa8d44ad9a6632a5932277a8e}
\index{[CT3] Service Provision Relationship (II.1.3.4.4.4)}
{\small Capability · Proposed · superseded}\par
\href{../../../../continuum/strategy/capability/718ebe1fa8d44ad9a6632a5932277a8e.md}{Historical canonical definition and provenance}.\par
\subsubsection*{II.1.3.4.5 [CT2] Location Administration}
\phantomsection
\label{navigation:II.1.3.4.5}
\addcontentsline{toc}{subsubsection}{II.1.3.4.5 [CT2] Location Administration}

{\small Historical predecessor: R.2.107 \hyperref[element:d412aff3a5c64e6aaba910d4c2c9356d]{[CT2] Location Administration}; Canonical ID: \texttt{\seqsplit{d412aff3a5c64e6aaba910d4c2c9356d}}. Superseded historical identity — direct correspondence · superseded.}\par
\label{element:c410005667444f99a221483bf577fad2}
\index{[CT2] Location Administration (II.1.3.4.5)}
{\small Capability · Canonical · current}\par

Location Administration represents the places and physical relationships relevant to healthcare activity, from facilities to particular care places. Address and geospatial context answer where a place is; hierarchy explains its surroundings; care-place definition records the local characteristics relevant to care. These are stable descriptive responsibilities, distinct from the changing operational state of a bed or ward.



\textbf{Overview}



Location Administration represents the places and physical relationships relevant to healthcare activity, from facilities to particular care places. Address and geospatial context answer where a place is; hierarchy explains its surroundings; care-place definition records the local characteristics relevant to care. These are stable descriptive responsibilities, distinct from the changing operational state of a bed or ward.



\textbf{Intent}



Location Administration represents the places and physical relationships relevant to healthcare activity, from facilities to particular care places. Address and geospatial context answer where a place is; hierarchy explains its surroundings; care-place definition records the local characteristics relevant to care. These are stable descriptive responsibilities, distinct from the changing operational state of a bed or ward.



\textbf{Boundary}



Information administration does not confer authority over the represented person, organisation, service, place, device or clinical activity.



\medskip\hrule\medskip

{\small Recorded unresolved matters: source-qualified matters; exact scope remains in the canonical record.}\par
{\footnotesize \href{../../../../continuum/strategy/capability/c410005667444f99a221483bf577fad2.md}{Canonical record: c410005667444f99a221483bf577fad2}; full qualifications and provenance.}\par
\paragraph*{II.1.3.4.5.1 [CT3] Care-Place Definition Administration}
\phantomsection
\label{navigation:II.1.3.4.5.1}
\addcontentsline{toc}{paragraph}{II.1.3.4.5.1 [CT3] Care-Place Definition Administration}

{\small Historical predecessor: R.2.105 \hyperref[element:a3464bd9c0dc4a01b50803acc7705b4e]{[CT3] Care-Place Definition Administration}; Canonical ID: \texttt{\seqsplit{a3464bd9c0dc4a01b50803acc7705b4e}}. Superseded historical identity — direct correspondence · superseded.}\par
\label{element:629b493c7d9844baa25c029a87ed04f6}
\index{[CT3] Care-Place Definition Administration (II.1.3.4.5.1)}
{\small Capability · Proposed · current}\par

Administer the stable description, characteristics and physical location of a bed, bay, chair or other care-place.



\textbf{Overview}



Administer the stable description, characteristics and physical location of a bed, bay, chair or other care-place.



\textbf{Intent}



Administer the stable description, characteristics and physical location of a bed, bay, chair or other care-place.



\textbf{Boundary}



Describes a place and its capabilities; operational bed state and readiness remain separately managed.


{\footnotesize \href{../../../../continuum/strategy/capability/629b493c7d9844baa25c029a87ed04f6.md}{Canonical record: 629b493c7d9844baa25c029a87ed04f6}; full qualifications and provenance.}\par
\subparagraph*{II.1.3.4.5.1.1 [Feature] Atomic Care-Place Modeling}
\phantomsection
\label{navigation:II.1.3.4.5.1.1}
\addcontentsline{toc}{subparagraph}{II.1.3.4.5.1.1 [Feature] Atomic Care-Place Modeling}

\label{element:ddf34f0844114798a4bd84ee72bbe497}
\index{[Feature] Atomic Care-Place Modeling (II.1.3.4.5.1.1)}
{\small Capability / Feature · Canonical · current}\par

Define specific beds, care-bays, and examination rooms with their physical capabilities (e.g., telemetry-capable, negative-pressure).



Care-place modelling identifies the particular physical place and characteristics relevant to care, such as isolation or telemetry capability. It supplies the descriptive frame against which changing occupancy and readiness can be considered without making those states part of the place's identity.


{\footnotesize \href{../../../../continuum/strategy/capability/feature/ddf34f0844114798a4bd84ee72bbe497.md}{Canonical record: ddf34f0844114798a4bd84ee72bbe497}; full qualifications and provenance.}\par
\paragraph*{II.1.3.4.5.2 [CT3] Healthcare Location Profile Administration}
\phantomsection
\label{navigation:II.1.3.4.5.2}
\addcontentsline{toc}{paragraph}{II.1.3.4.5.2 [CT3] Healthcare Location Profile Administration}

{\small Historical predecessor: R.2.106 \hyperref[element:8a04d1f16d0849e786d2aaec33e0ba87]{[CT3] Healthcare Location Profile Administration}; Canonical ID: \texttt{\seqsplit{8a04d1f16d0849e786d2aaec33e0ba87}}. Superseded historical identity — direct correspondence · superseded.}\par
\label{element:e579e54a08bb4f899a504aa998ba246b}
\index{[CT3] Healthcare Location Profile Administration (II.1.3.4.5.2)}
{\small Capability · Proposed · current}\par

Administer the identity, address, geospatial description and physical containment relationships of healthcare locations.



\textbf{Overview}



Administer the identity, address, geospatial description and physical containment relationships of healthcare locations.



\textbf{Intent}



Administer the identity, address, geospatial description and physical containment relationships of healthcare locations.



\textbf{Boundary}



Location description does not establish occupancy, cleanliness, service capacity or service readiness.


{\footnotesize \href{../../../../continuum/strategy/capability/e579e54a08bb4f899a504aa998ba246b.md}{Canonical record: e579e54a08bb4f899a504aa998ba246b}; full qualifications and provenance.}\par
\subparagraph*{II.1.3.4.5.2.1 [Feature] Location Geospatial \& Address Resolution}
\phantomsection
\label{navigation:II.1.3.4.5.2.1}
\addcontentsline{toc}{subparagraph}{II.1.3.4.5.2.1 [Feature] Location Geospatial \& Address Resolution}

\label{element:a27d1f97313f4ba5907c7cf7e2ec8858}
\index{[Feature] Location Geospatial \& Address Resolution (II.1.3.4.5.2.1)}
{\small Capability / Feature · Canonical · current}\par

Validate physical addresses, geospatial coordinates, and location operational status.



Address and geospatial resolution gives a recognised healthcare location an intelligible physical reference. It supports place interpretation without treating a valid address or recorded operating status as proof that a particular service or care place is currently ready.


{\footnotesize \href{../../../../continuum/strategy/capability/feature/a27d1f97313f4ba5907c7cf7e2ec8858.md}{Canonical record: a27d1f97313f4ba5907c7cf7e2ec8858}; full qualifications and provenance.}\par
\subparagraph*{II.1.3.4.5.2.2 [Feature] Physical Hierarchy Modeling}
\phantomsection
\label{navigation:II.1.3.4.5.2.2}
\addcontentsline{toc}{subparagraph}{II.1.3.4.5.2.2 [Feature] Physical Hierarchy Modeling}

\label{element:6166149dbd7b431985063d92e78ecdb7}
\index{[Feature] Physical Hierarchy Modeling (II.1.3.4.5.2.2)}
{\small Capability / Feature · Canonical · current}\par

Represent campus, building, wing, floor, and ward architectural relationships.



Physical hierarchy modelling relates a care environment to its campus, building, floor or ward context. The relationships describe where places sit physically; they do not automatically reproduce organisational hierarchy, service responsibility or operational capacity.


{\footnotesize \href{../../../../continuum/strategy/capability/feature/6166149dbd7b431985063d92e78ecdb7.md}{Canonical record: 6166149dbd7b431985063d92e78ecdb7}; full qualifications and provenance.}\par
\paragraph*{II.1.3.4.5.3 [CT3] Location Structure [superseded]}
\phantomsection
\label{navigation:II.1.3.4.5.3}
\addcontentsline{toc}{paragraph}{II.1.3.4.5.3 [CT3] Location Structure [superseded]}

\label{element:2436c87195334daf85ed4e2de4ec4bd5}
\index{[CT3] Location Structure (II.1.3.4.5.3)}
{\small Capability · Proposed · superseded}\par
\href{../../../../continuum/strategy/capability/2436c87195334daf85ed4e2de4ec4bd5.md}{Historical canonical definition and provenance}.\par
\subsubsection*{II.1.3.4.6 [CT2] Organisation Administration}
\phantomsection
\label{navigation:II.1.3.4.6}
\addcontentsline{toc}{subsubsection}{II.1.3.4.6 [CT2] Organisation Administration}

{\small Historical predecessor: R.2.109 \hyperref[element:1b0eacfbb27d4631a24f9279dab6681d]{[CT2] Organisation Administration}; Canonical ID: \texttt{\seqsplit{1b0eacfbb27d4631a24f9279dab6681d}}. Superseded historical identity — direct correspondence · superseded.}\par
\label{element:14fbdc5ce3dd4c2092c5397dab3ed743}
\index{[CT2] Organisation Administration (II.1.3.4.6)}
{\small Capability · Canonical · current}\par

Organisation Administration maintains the identity, profiles, structural relationships and contact context of participating healthcare organisations. It enables a legal or operational entity to be distinguished from a department, network relationship or communication channel. This supports consistent interpretation of providers and destinations across information flows rather than allowing each interface to create an incompatible organisational account.



\textbf{Overview}



Organisation Administration maintains the identity, profiles, structural relationships and contact context of participating healthcare organisations. It enables a legal or operational entity to be distinguished from a department, network relationship or communication channel. This supports consistent interpretation of providers and destinations across information flows rather than allowing each interface to create an incompatible organisational account.



\textbf{Intent}



Organisation Administration maintains the identity, profiles, structural relationships and contact context of participating healthcare organisations. It enables a legal or operational entity to be distinguished from a department, network relationship or communication channel. This supports consistent interpretation of providers and destinations across information flows rather than allowing each interface to create an incompatible organisational account.



\textbf{Boundary}



Information administration does not confer authority over the represented person, organisation, service, place, device or clinical activity.



\medskip\hrule\medskip

{\small Recorded unresolved matters: source-qualified matters; exact scope remains in the canonical record.}\par
{\footnotesize \href{../../../../continuum/strategy/capability/14fbdc5ce3dd4c2092c5397dab3ed743.md}{Canonical record: 14fbdc5ce3dd4c2092c5397dab3ed743}; full qualifications and provenance.}\par
\paragraph*{II.1.3.4.6.1 [CT3] Healthcare Organisation Profile Administration}
\phantomsection
\label{navigation:II.1.3.4.6.1}
\addcontentsline{toc}{paragraph}{II.1.3.4.6.1 [CT3] Healthcare Organisation Profile Administration}

{\small Historical predecessor: R.2.108 \hyperref[element:e485424b799c42aca7399d341ca02d0b]{[CT3] Healthcare Organisation Profile Administration}; Canonical ID: \texttt{\seqsplit{e485424b799c42aca7399d341ca02d0b}}. Superseded historical identity — direct correspondence · superseded.}\par
\label{element:f2de9e9394b849f2943355c7a2882a15}
\index{[CT3] Healthcare Organisation Profile Administration (II.1.3.4.6.1)}
{\small Capability · Proposed · current}\par

Administer identifying, descriptive, structural, location, contact, registered/legal and recognised endpoint information needed to represent and interact with a healthcare organisation.



\textbf{Overview}



Administer identifying, descriptive, structural, location, contact, registered/legal and recognised endpoint information needed to represent and interact with a healthcare organisation.



\textbf{Intent}



Administer identifying, descriptive, structural, location, contact, registered/legal and recognised endpoint information needed to represent and interact with a healthcare organisation.



\textbf{Boundary}



Maintains verified legal and operational identity information without corporate management authority.


{\footnotesize \href{../../../../continuum/strategy/capability/f2de9e9394b849f2943355c7a2882a15.md}{Canonical record: f2de9e9394b849f2943355c7a2882a15}; full qualifications and provenance.}\par
\subparagraph*{II.1.3.4.6.1.1 [Feature] Directory Contact Management}
\phantomsection
\label{navigation:II.1.3.4.6.1.1}
\addcontentsline{toc}{subparagraph}{II.1.3.4.6.1.1 [Feature] Directory Contact Management}

\label{element:1cdac17b7547448daa20b74ed45ec366}
\index{[Feature] Directory Contact Management (II.1.3.4.6.1.1)}
{\small Capability / Feature · Canonical · current}\par

Manage verified phone, email, and administrative contact channels for healthcare facilities.



Contact management provides verified administrative and clinical communication channels for organisations. It supports reaching the appropriate participant while distinguishing a directory contact from the role-specific destination and authority relevant to a particular clinical exchange.


{\footnotesize \href{../../../../continuum/strategy/capability/feature/1cdac17b7547448daa20b74ed45ec366.md}{Canonical record: 1cdac17b7547448daa20b74ed45ec366}; full qualifications and provenance.}\par
\subparagraph*{II.1.3.4.6.1.2 [Feature] Hierarchical Structure Navigation}
\phantomsection
\label{navigation:II.1.3.4.6.1.2}
\addcontentsline{toc}{subparagraph}{II.1.3.4.6.1.2 [Feature] Hierarchical Structure Navigation}

\label{element:7ded6f7865294fb2a696a58160d577f7}
\index{[Feature] Hierarchical Structure Navigation (II.1.3.4.6.1.2)}
{\small Capability / Feature · Canonical · current}\par

Model, traverse, and resolve nested departmental, divisional, and regional healthcare networks.



Hierarchy navigation makes departmental, divisional and network relationships usable without flattening them into unrelated organisations. Traversal reflects established structure; it does not create corporate authority or infer responsibility from the way a directory happens to display a tree.


{\footnotesize \href{../../../../continuum/strategy/capability/feature/7ded6f7865294fb2a696a58160d577f7.md}{Canonical record: 7ded6f7865294fb2a696a58160d577f7}; full qualifications and provenance.}\par
\subparagraph*{II.1.3.4.6.1.3 [Feature] National Organisation Verification}
\phantomsection
\label{navigation:II.1.3.4.6.1.3}
\addcontentsline{toc}{subparagraph}{II.1.3.4.6.1.3 [Feature] National Organisation Verification}

\label{element:0f17fe5439fe45dc8f4edbff3cca0776}
\index{[Feature] National Organisation Verification (II.1.3.4.6.1.3)}
{\small Capability / Feature · Canonical · current}\par

Validate national organization identifiers (e.g., HPI-O) and legal entity classifications.



National organisation verification grounds an organisation identifier and entity classification in their recognised context. This supports reliable interpretation of the healthcare participant without assuming that its identity establishes every service, location or authority associated with it.


{\footnotesize \href{../../../../continuum/strategy/capability/feature/0f17fe5439fe45dc8f4edbff3cca0776.md}{Canonical record: 0f17fe5439fe45dc8f4edbff3cca0776}; full qualifications and provenance.}\par
\subparagraph*{II.1.3.4.6.1.4 [Feature] Organisation Profile Management}
\phantomsection
\label{navigation:II.1.3.4.6.1.4}
\addcontentsline{toc}{subparagraph}{II.1.3.4.6.1.4 [Feature] Organisation Profile Management}

\label{element:6a4c44cfb1414f3bae5ba0fce1e58970}
\index{[Feature] Organisation Profile Management (II.1.3.4.6.1.4)}
{\small Capability / Feature · Canonical · current}\par

Govern directory profiles, operational status, and published service descriptions.



Organisation profiles provide a governed account of the participant's status and published service descriptions. They support consistent directory use across systems while leaving the organisation's real-world governance and the definition of individual offerings distinguishable.


{\footnotesize \href{../../../../continuum/strategy/capability/feature/6a4c44cfb1414f3bae5ba0fce1e58970.md}{Canonical record: 6a4c44cfb1414f3bae5ba0fce1e58970}; full qualifications and provenance.}\par
\paragraph*{II.1.3.4.6.2 [CT3] Organisation Contact [superseded]}
\phantomsection
\label{navigation:II.1.3.4.6.2}
\addcontentsline{toc}{paragraph}{II.1.3.4.6.2 [CT3] Organisation Contact [superseded]}

\label{element:b7f5108880eb429bb407c87a37b0e23f}
\index{[CT3] Organisation Contact (II.1.3.4.6.2)}
{\small Capability · Proposed · superseded}\par
\href{../../../../continuum/strategy/capability/b7f5108880eb429bb407c87a37b0e23f.md}{Historical canonical definition and provenance}.\par
\paragraph*{II.1.3.4.6.3 [CT3] Organisational Structure [superseded]}
\phantomsection
\label{navigation:II.1.3.4.6.3}
\addcontentsline{toc}{paragraph}{II.1.3.4.6.3 [CT3] Organisational Structure [superseded]}

\label{element:19d211cf601c49e5a65665330d0c5c3d}
\index{[CT3] Organisational Structure (II.1.3.4.6.3)}
{\small Capability · Proposed · superseded}\par
\href{../../../../continuum/strategy/capability/19d211cf601c49e5a65665330d0c5c3d.md}{Historical canonical definition and provenance}.\par
\subsubsection*{II.1.3.4.7 [CT2] Provider Administration}
\phantomsection
\label{navigation:II.1.3.4.7}
\addcontentsline{toc}{subsubsection}{II.1.3.4.7 [CT2] Provider Administration}

{\small Historical predecessor: R.2.113 \hyperref[element:a352fd50419d451887c85fcb98566cd5]{[CT2] Provider Administration}; Canonical ID: \texttt{\seqsplit{a352fd50419d451887c85fcb98566cd5}}. Superseded historical identity — direct correspondence · superseded.}\par
\label{element:d6312f85968a4209bafff5ed009297a6}
\index{[CT2] Provider Administration (II.1.3.4.7)}
{\small Capability · Canonical · current}\par

Provider Administration maintains the practitioner, role, affiliation and engagement context needed to interpret healthcare participation. National identifiers, professional profiles, role-specific endpoints and privilege records support different questions: who the practitioner is, in what capacity they act and how they can be reached. Directory representation makes those distinctions usable without replacing the business governance that establishes professional standing.



\textbf{Overview}



Provider Administration maintains the practitioner, role, affiliation and engagement context needed to interpret healthcare participation. National identifiers, professional profiles, role-specific endpoints and privilege records support different questions: who the practitioner is, in what capacity they act and how they can be reached. Directory representation makes those distinctions usable without replacing the business governance that establishes professional standing.



\textbf{Intent}



Provider Administration maintains the practitioner, role, affiliation and engagement context needed to interpret healthcare participation. National identifiers, professional profiles, role-specific endpoints and privilege records support different questions: who the practitioner is, in what capacity they act and how they can be reached. Directory representation makes those distinctions usable without replacing the business governance that establishes professional standing.



\textbf{Boundary}



Information administration does not confer authority over the represented person, organisation, service, place, device or clinical activity.



\medskip\hrule\medskip

{\small Recorded unresolved matters: source-qualified matters; exact scope remains in the canonical record.}\par
{\footnotesize \href{../../../../continuum/strategy/capability/d6312f85968a4209bafff5ed009297a6.md}{Canonical record: d6312f85968a4209bafff5ed009297a6}; full qualifications and provenance.}\par
\paragraph*{II.1.3.4.7.1 [CT3] Practitioner Profile Administration}
\phantomsection
\label{navigation:II.1.3.4.7.1}
\addcontentsline{toc}{paragraph}{II.1.3.4.7.1 [CT3] Practitioner Profile Administration}

{\small Historical predecessor: R.2.110 \hyperref[element:2415ddcbc73c47e1a63be310bd4342ee]{[CT3] Practitioner Profile Administration}; Canonical ID: \texttt{\seqsplit{2415ddcbc73c47e1a63be310bd4342ee}}. Superseded historical identity — direct correspondence · superseded.}\par
\label{element:40759a835f7b456e9257c38b67b0a380}
\index{[CT3] Practitioner Profile Administration (II.1.3.4.7.1)}
{\small Capability · Proposed · current}\par

Administer practitioner identity, professional profile and supplied registration information with its issuing-authority attribution.



\textbf{Overview}



Administer practitioner identity, professional profile and supplied registration information with its issuing-authority attribution.



\textbf{Intent}



Administer practitioner identity, professional profile and supplied registration information with its issuing-authority attribution.



\textbf{Boundary}



Professional registration and profile representation do not grant clinical or operational privilege.


{\footnotesize \href{../../../../continuum/strategy/capability/40759a835f7b456e9257c38b67b0a380.md}{Canonical record: 40759a835f7b456e9257c38b67b0a380}; full qualifications and provenance.}\par
\subparagraph*{II.1.3.4.7.1.1 [Feature] National Practitioner Verification}
\phantomsection
\label{navigation:II.1.3.4.7.1.1}
\addcontentsline{toc}{subparagraph}{II.1.3.4.7.1.1 [Feature] National Practitioner Verification}

\label{element:12c5a17dcc564619a6a9222dd2e15528}
\index{[Feature] National Practitioner Verification (II.1.3.4.7.1.1)}
{\small Capability / Feature · Canonical · current}\par

Ingest, validate, and verify national practitioner identifiers (e.g., HPI-I) and professional registration status.



National verification connects practitioner identifiers and registration information to their recognised source context. It helps distinguish an identifiable professional from an unverified directory assertion without implying that verification alone grants local clinical privileges.


{\footnotesize \href{../../../../continuum/strategy/capability/feature/12c5a17dcc564619a6a9222dd2e15528.md}{Canonical record: 12c5a17dcc564619a6a9222dd2e15528}; full qualifications and provenance.}\par
\subparagraph*{II.1.3.4.7.1.2 [Feature] Practitioner Profile Governance}
\phantomsection
\label{navigation:II.1.3.4.7.1.2}
\addcontentsline{toc}{subparagraph}{II.1.3.4.7.1.2 [Feature] Practitioner Profile Governance}

\label{element:d23be319c6114c28bcb9e24e1258a0fc}
\index{[Feature] Practitioner Profile Governance (II.1.3.4.7.1.2)}
{\small Capability / Feature · Canonical · current}\par

Maintain core professional demographic details, qualifications, and active practice status.



Practitioner profile governance maintains the professional context used across care information and communication. Qualifications and practice status remain attributable facts, distinct from the organisation-specific role or privilege under which the practitioner may act.


{\footnotesize \href{../../../../continuum/strategy/capability/feature/d23be319c6114c28bcb9e24e1258a0fc.md}{Canonical record: d23be319c6114c28bcb9e24e1258a0fc}; full qualifications and provenance.}\par
\subparagraph*{II.1.3.4.7.1.3 [Feature] Provider Directory Discovery \& Resolution}
\phantomsection
\label{navigation:II.1.3.4.7.1.3}
\addcontentsline{toc}{subparagraph}{II.1.3.4.7.1.3 [Feature] Provider Directory Discovery \& Resolution}

\label{element:cd3bd35f7f9e4852bbbce299bb795d8a}
\index{[Feature] Provider Directory Discovery \& Resolution (II.1.3.4.7.1.3)}
{\small Capability / Feature · Proposed · current}\par

Expose qualified practitioner discovery using identifiers, profile criteria and effective standing.



\textbf{Overview}



Routing and care participation need to find a described provider matching stated criteria, beyond merely ingesting a national identifier or updating a profile.



\textbf{Boundary}



Discovery supplies information; it does not determine privilege or availability.


{\footnotesize \href{../../../../continuum/strategy/capability/feature/cd3bd35f7f9e4852bbbce299bb795d8a.md}{Canonical record: cd3bd35f7f9e4852bbbce299bb795d8a}; full qualifications and provenance.}\par
\paragraph*{II.1.3.4.7.2 [CT3] Practitioner Role Administration}
\phantomsection
\label{navigation:II.1.3.4.7.2}
\addcontentsline{toc}{paragraph}{II.1.3.4.7.2 [CT3] Practitioner Role Administration}

{\small Historical predecessor: R.2.111 \hyperref[element:6b777f1201db474eb22dd89f2b1f68c6]{[CT3] Practitioner Role Administration}; Canonical ID: \texttt{\seqsplit{6b777f1201db474eb22dd89f2b1f68c6}}. Superseded historical identity — direct correspondence · superseded.}\par
\label{element:6651471e67ae46fb8d1c1241c8590be9}
\index{[CT3] Practitioner Role Administration (II.1.3.4.7.2)}
{\small Capability · Proposed · current}\par

Administer practitioner roles, specialties, organisational affiliations, operational capacities and recognised role endpoints.



\textbf{Overview}



Administer practitioner roles, specialties, organisational affiliations, operational capacities and recognised role endpoints.



\textbf{Intent}



Administer practitioner roles, specialties, organisational affiliations, operational capacities and recognised role endpoints.



\textbf{Boundary}



Role, specialty, capacity and endpoint information supply context; they do not constitute a privilege determination.


{\footnotesize \href{../../../../continuum/strategy/capability/6651471e67ae46fb8d1c1241c8590be9.md}{Canonical record: 6651471e67ae46fb8d1c1241c8590be9}; full qualifications and provenance.}\par
\subparagraph*{II.1.3.4.7.2.1 [Feature] Electronic Endpoint Binding}
\phantomsection
\label{navigation:II.1.3.4.7.2.1}
\addcontentsline{toc}{subparagraph}{II.1.3.4.7.2.1 [Feature] Electronic Endpoint Binding}

\label{element:5dc3a424e68d47aea64ec573dc4b1225}
\index{[Feature] Electronic Endpoint Binding (II.1.3.4.7.2.1)}
{\small Capability / Feature · Canonical · current}\par

Associate validated communication endpoints (secure messaging addresses, direct routing queues) with specific practitioner roles.



Endpoint binding makes a verified communication destination specific to the practitioner's acting role. It supports routing to the intended professional capacity rather than assuming that every address attached to the person is interchangeable across organisations and services.


{\footnotesize \href{../../../../continuum/strategy/capability/feature/5dc3a424e68d47aea64ec573dc4b1225.md}{Canonical record: 5dc3a424e68d47aea64ec573dc4b1225}; full qualifications and provenance.}\par
\subparagraph*{II.1.3.4.7.2.2 [Feature] Practitioner-Organisation Affiliation}
\phantomsection
\label{navigation:II.1.3.4.7.2.2}
\addcontentsline{toc}{subparagraph}{II.1.3.4.7.2.2 [Feature] Practitioner-Organisation Affiliation}

\label{element:9501e88ea4b7474a8aa63191687db51e}
\index{[Feature] Practitioner-Organisation Affiliation (II.1.3.4.7.2.2)}
{\small Capability / Feature · Canonical · current}\par

Maintain temporal affiliations between practitioners and healthcare organizations.



Affiliation records keep practitioner participation in an organisation temporally intelligible. They support interpretation of an actor's organisational context while preserving that affiliation, employment, qualification and permission to undertake a particular activity answer different questions.


{\footnotesize \href{../../../../continuum/strategy/capability/feature/9501e88ea4b7474a8aa63191687db51e.md}{Canonical record: 9501e88ea4b7474a8aa63191687db51e}; full qualifications and provenance.}\par
\subparagraph*{II.1.3.4.7.2.3 [Feature] Role \& Specialty Modeling}
\phantomsection
\label{navigation:II.1.3.4.7.2.3}
\addcontentsline{toc}{subparagraph}{II.1.3.4.7.2.3 [Feature] Role \& Specialty Modeling}

\label{element:ef859d4934514f4eb601c6abf9708760}
\index{[Feature] Role \& Specialty Modeling (II.1.3.4.7.2.3)}
{\small Capability / Feature · Canonical · current}\par

Define and associate clinical specialties, practice roles, and service capacities for a practitioner.



Role and specialty modelling describes the capacity in which a practitioner participates at an organisation or location. This keeps the person separate from each acting role and its service context; role resolution itself establishes neither Clinical Privilege nor Operational Privilege.


{\footnotesize \href{../../../../continuum/strategy/capability/feature/ef859d4934514f4eb601c6abf9708760.md}{Canonical record: ef859d4934514f4eb601c6abf9708760}; full qualifications and provenance.}\par
\paragraph*{II.1.3.4.7.3 [CT3] Provider Engagement [superseded]}
\phantomsection
\label{navigation:II.1.3.4.7.3}
\addcontentsline{toc}{paragraph}{II.1.3.4.7.3 [CT3] Provider Engagement [superseded]}

\label{element:a0777dbd9b0c4528817e99a7f34aa444}
\index{[CT3] Provider Engagement (II.1.3.4.7.3)}
{\small Capability · Proposed · superseded}\par
\href{../../../../continuum/strategy/capability/a0777dbd9b0c4528817e99a7f34aa444.md}{Historical canonical definition and provenance}.\par
\paragraph*{II.1.3.4.7.4 [CT3] Provider Relationship [superseded]}
\phantomsection
\label{navigation:II.1.3.4.7.4}
\addcontentsline{toc}{paragraph}{II.1.3.4.7.4 [CT3] Provider Relationship [superseded]}

\label{element:3bb8a3d5b62a4b779442290fe8d5106c}
\index{[CT3] Provider Relationship (II.1.3.4.7.4)}
{\small Capability · Proposed · superseded}\par
\href{../../../../continuum/strategy/capability/3bb8a3d5b62a4b779442290fe8d5106c.md}{Historical canonical definition and provenance}.\par
\paragraph*{II.1.3.4.7.5 [CT3] Vendor Liaison}
\phantomsection
\label{navigation:II.1.3.4.7.5}
\addcontentsline{toc}{paragraph}{II.1.3.4.7.5 [CT3] Vendor Liaison}

{\small Historical predecessor: R.2.112 \hyperref[element:29076cbb9fc5442fb25faf43947574f7]{[CT3] Vendor Liaison}; Canonical ID: \texttt{\seqsplit{29076cbb9fc5442fb25faf43947574f7}}. Superseded historical identity — direct correspondence · superseded.}\par
\label{element:66358bd6698e4050b12342522fbbb923}
\index{[CT3] Vendor Liaison (II.1.3.4.7.5)}
{\small Capability · Proposed · current}\par

Maintain the information and liaison arrangements needed to interact with vendors supporting healthcare provision.



\textbf{Overview}



Maintain the information and liaison arrangements needed to interact with vendors supporting healthcare provision. The reviewed hierarchy assigns this enduring ability to its stated business-enabling context.



\textbf{Intent}



Maintain the information and liaison arrangements needed to interact with vendors supporting healthcare provision.



\textbf{Boundary}



Information administration and coordination preserve the authority of the responsible source, clinical actor and governed domain; they confer no additional real-world management authority.


{\footnotesize \href{../../../../continuum/strategy/capability/66358bd6698e4050b12342522fbbb923.md}{Canonical record: 66358bd6698e4050b12342522fbbb923}; full qualifications and provenance.}\par
\subsection*{II.1.3.5 [CT1] Health Informatics Management}
\phantomsection
\label{navigation:II.1.3.5}
\addcontentsline{toc}{subsection}{II.1.3.5 [CT1] Health Informatics Management}

{\small Historical predecessor: R.2.118 \hyperref[element:50c56bac0dc540a2b392fc2376efc129]{[CT1] Health Informatics Management}; Canonical ID: \texttt{\seqsplit{50c56bac0dc540a2b392fc2376efc129}}. Superseded historical identity — direct correspondence · superseded.}\par
\label{element:a790b5f2c5a1497fa6ec09d7824da32e}
\index{[CT1] Health Informatics Management (II.1.3.5)}
{\small Capability · Proposed · current}\par

Establish and govern semantically meaningful health-information assets suitable for downstream research, population-health reporting and service planning in an adjacent analytical environment.



\textbf{Overview}



Establish and govern semantically meaningful health-information assets suitable for downstream research, population-health reporting and service planning in an adjacent analytical environment.



\textbf{Intent}



Establish and govern semantically meaningful health-information assets suitable for downstream research, population-health reporting and service planning in an adjacent analytical environment.



\textbf{Boundary}



Harmonia is adjacent to a Big Data/analytical environment. It establishes governed assets and reporting inputs; it does not become that analytical platform.


{\footnotesize \href{../../../../continuum/strategy/capability/a790b5f2c5a1497fa6ec09d7824da32e.md}{Canonical record: a790b5f2c5a1497fa6ec09d7824da32e}; full qualifications and provenance.}\par
\subsubsection*{II.1.3.5.1 [CT2] Population Health Information Reporting}
\phantomsection
\label{navigation:II.1.3.5.1}
\addcontentsline{toc}{subsubsection}{II.1.3.5.1 [CT2] Population Health Information Reporting}

{\small Historical predecessor: R.2.115 \hyperref[element:ced472b2d1a54016878ee541de969d9b]{[CT2] Population Health Information Reporting}; Canonical ID: \texttt{\seqsplit{ced472b2d1a54016878ee541de969d9b}}. Superseded historical identity — direct correspondence · superseded.}\par
\label{element:2135b45eab034370b0063b20e821985d}
\index{[CT2] Population Health Information Reporting (II.1.3.5.1)}
{\small Capability · Proposed · current}\par

Establish and provide governed information assets and reporting inputs for population-health purposes, including applicable notification obligations.



\textbf{Overview}



Establish and provide governed information assets and reporting inputs for population-health purposes, including applicable notification obligations.



\textbf{Intent}



Establish and provide governed information assets and reporting inputs for population-health purposes, including applicable notification obligations.



\textbf{Boundary}



Does not determine population-health findings, create statutory obligations or equate transport delivery with discharge of notification duty.


{\footnotesize \href{../../../../continuum/strategy/capability/2135b45eab034370b0063b20e821985d.md}{Canonical record: 2135b45eab034370b0063b20e821985d}; full qualifications and provenance.}\par
\subsubsection*{II.1.3.5.2 [CT2] Research Information Curation}
\phantomsection
\label{navigation:II.1.3.5.2}
\addcontentsline{toc}{subsubsection}{II.1.3.5.2 [CT2] Research Information Curation}

{\small Historical predecessor: R.2.116 \hyperref[element:11203b9c80964a9b96883eeab07f41d5]{[CT2] Research Information Curation}; Canonical ID: \texttt{\seqsplit{11203b9c80964a9b96883eeab07f41d5}}. Superseded historical identity — direct correspondence · superseded.}\par
\label{element:4f5b8cec90324c2989cdc30e80bc8cb9}
\index{[CT2] Research Information Curation (II.1.3.5.2)}
{\small Capability · Proposed · current}\par

Establish and curate governed data assets suitable for explicitly authorised research use, retaining source attribution and applicable release constraints.



\textbf{Overview}



Establish and curate governed data assets suitable for explicitly authorised research use, retaining source attribution and applicable release constraints.



\textbf{Intent}



Establish and curate governed data assets suitable for explicitly authorised research use, retaining source attribution and applicable release constraints.



\textbf{Boundary}



Does not grant ethics approval, recruit trial participants, manage trials or become the authority for derived research conclusions.


{\footnotesize \href{../../../../continuum/strategy/capability/4f5b8cec90324c2989cdc30e80bc8cb9.md}{Canonical record: 4f5b8cec90324c2989cdc30e80bc8cb9}; full qualifications and provenance.}\par
\subsubsection*{II.1.3.5.3 [CT2] Service Planning Data Management}
\phantomsection
\label{navigation:II.1.3.5.3}
\addcontentsline{toc}{subsubsection}{II.1.3.5.3 [CT2] Service Planning Data Management}

{\small Historical predecessor: R.2.117 \hyperref[element:19263e0802cc4b88af7a345f1410cffd]{[CT2] Service Planning Data Management}; Canonical ID: \texttt{\seqsplit{19263e0802cc4b88af7a345f1410cffd}}. Superseded historical identity — direct correspondence · superseded.}\par
\label{element:4873cd4b960d45d8b93560d5c8b98500}
\index{[CT2] Service Planning Data Management (II.1.3.5.3)}
{\small Capability · Proposed · current}\par

Establish and maintain governed information assets suitable for health-service planning, modelling and analysis.



\textbf{Overview}



Establish and maintain governed information assets suitable for health-service planning, modelling and analysis.



\textbf{Intent}



Establish and maintain governed information assets suitable for health-service planning, modelling and analysis.



\textbf{Boundary}



Observed utilisation, operational capacity and planned capacity are distinct; information provision does not commission services or allocate resources.


{\footnotesize \href{../../../../continuum/strategy/capability/4873cd4b960d45d8b93560d5c8b98500.md}{Canonical record: 4873cd4b960d45d8b93560d5c8b98500}; full qualifications and provenance.}\par
\subsection*{II.1.3.6 [CT1] Health Service Operations Enablement}
\phantomsection
\label{navigation:II.1.3.6}
\addcontentsline{toc}{subsection}{II.1.3.6 [CT1] Health Service Operations Enablement}

{\small Historical predecessor: R.2.160 \hyperref[element:a98424298dac46d8b2ddd109d928270a]{[CT1] Health Service Operations Enablement}; Canonical ID: \texttt{\seqsplit{a98424298dac46d8b2ddd109d928270a}}. Superseded historical identity — direct correspondence · superseded.}\par
\label{element:b13ee2c5a4ef4fd6be8a44795696dfc2}
\index{[CT1] Health Service Operations Enablement (II.1.3.6)}
{\small Capability · Proposed · current}\par

Enable health-service operations through information about flow, care-places, capacity, workforce capability, operational work, logistics and discharge progression.



\textbf{Overview}



Enable health-service operations through information about flow, care-places, capacity, workforce capability, operational work, logistics and discharge progression.



\textbf{Intent}



Enable health-service operations through information about flow, care-places, capacity, workforce capability, operational work, logistics and discharge progression.



\textbf{Boundary}



Coordination owns only explicitly assigned operational activity. It does not originate clinical clearance, determine clinical work or take workforce and resource governance from their sources.


{\footnotesize \href{../../../../continuum/strategy/capability/b13ee2c5a4ef4fd6be8a44795696dfc2.md}{Canonical record: b13ee2c5a4ef4fd6be8a44795696dfc2}; full qualifications and provenance.}\par
\subsubsection*{II.1.3.6.1 [CT2] Bed \& Care-Place Management}
\phantomsection
\label{navigation:II.1.3.6.1}
\addcontentsline{toc}{subsubsection}{II.1.3.6.1 [CT2] Bed \& Care-Place Management}

{\small Historical predecessor: R.2.121 \hyperref[element:c94d29b7d1c045dabc7814ceea795484]{[CT2] Bed \& Care-Place Management}; Canonical ID: \texttt{\seqsplit{c94d29b7d1c045dabc7814ceea795484}}. Superseded historical identity — direct correspondence · superseded.}\par
\label{element:f117ccec0bfa4f738e6ddb646b3534a2}
\index{[CT2] Bed \& Care-Place Management (II.1.3.6.1)}
{\small Capability · Canonical · current}\par

Bed \& Care-Place Management maintains operational place state and the turnover progression that makes a place usable again. Vacancy, reservation, cleaning and readiness must remain distinguishable so that an empty bed is not automatically treated as available for allocation. Notifications expose meaningful readiness changes to the responsible operational participants.



\textbf{Overview}



Bed \& Care-Place Management maintains operational place state and the turnover progression that makes a place usable again. Vacancy, reservation, cleaning and readiness must remain distinguishable so that an empty bed is not automatically treated as available for allocation. Notifications expose meaningful readiness changes to the responsible operational participants.



\textbf{Intent}



Bed \& Care-Place Management maintains operational place state and the turnover progression that makes a place usable again. Vacancy, reservation, cleaning and readiness must remain distinguishable so that an empty bed is not automatically treated as available for allocation. Notifications expose meaningful readiness changes to the responsible operational participants.



\textbf{Boundary}



Information administration does not confer authority over the represented person, organisation, service, place, device or clinical activity.



\medskip\hrule\medskip

{\footnotesize \href{../../../../continuum/strategy/capability/f117ccec0bfa4f738e6ddb646b3534a2.md}{Canonical record: f117ccec0bfa4f738e6ddb646b3534a2}; full qualifications and provenance.}\par
\paragraph*{II.1.3.6.1.1 [CT3] Bed \& Care-Place State Administration}
\phantomsection
\label{navigation:II.1.3.6.1.1}
\addcontentsline{toc}{paragraph}{II.1.3.6.1.1 [CT3] Bed \& Care-Place State Administration}

{\small Historical predecessor: R.2.119 \hyperref[element:dd2e2ee1a68845c087ba93610752181b]{[CT3] Bed \& Care-Place State Administration}; Canonical ID: \texttt{\seqsplit{dd2e2ee1a68845c087ba93610752181b}}. Superseded historical identity — direct correspondence · superseded.}\par
\label{element:26a9383d7bb146a19141558d232843de}
\index{[CT3] Bed \& Care-Place State Administration (II.1.3.6.1.1)}
{\small Capability · Proposed · current}\par

Administer the current operational occupancy, reservation, blocking and readiness states of beds and other care-places.



\textbf{Overview}



Administer the current operational occupancy, reservation, blocking and readiness states of beds and other care-places. The reviewed hierarchy assigns this enduring ability to its stated business-enabling context.



\textbf{Intent}



Administer the current operational occupancy, reservation, blocking and readiness states of beds and other care-places.



\textbf{Boundary}



Information administration and coordination preserve the authority of the responsible source, clinical actor and governed domain; they confer no additional real-world management authority.


{\footnotesize \href{../../../../continuum/strategy/capability/26a9383d7bb146a19141558d232843de.md}{Canonical record: 26a9383d7bb146a19141558d232843de}; full qualifications and provenance.}\par
\subparagraph*{II.1.3.6.1.1.1 [Feature] Bed Availability Tracking}
\phantomsection
\label{navigation:II.1.3.6.1.1.1}
\addcontentsline{toc}{subparagraph}{II.1.3.6.1.1.1 [Feature] Bed Availability Tracking}

\label{element:7c61f21e558c44ab92c267a781116d1d}
\index{[Feature] Bed Availability Tracking (II.1.3.6.1.1.1)}
{\small Capability / Feature · Canonical · current}\par

Maintain operational status of all registered beds and care-places (VACANT, OCCUPIED, BLOCKED, RESERVED, DIRTY, CLEANING, CLEAN).



Availability tracking preserves the operational distinctions among vacancy, occupancy, reservation, blockage and cleaning conditions. These states support meaningful care-place coordination without equating an unoccupied bed with a ready place or usable service capacity.


{\footnotesize \href{../../../../continuum/strategy/capability/feature/7c61f21e558c44ab92c267a781116d1d.md}{Canonical record: 7c61f21e558c44ab92c267a781116d1d}; full qualifications and provenance.}\par
\paragraph*{II.1.3.6.1.2 [CT3] Bed \& Care-Place Turnover Coordination}
\phantomsection
\label{navigation:II.1.3.6.1.2}
\addcontentsline{toc}{paragraph}{II.1.3.6.1.2 [CT3] Bed \& Care-Place Turnover Coordination}

{\small Historical predecessor: R.2.120 \hyperref[element:a50662edda954bc58c44184d84591d5b]{[CT3] Bed \& Care-Place Turnover Coordination}; Canonical ID: \texttt{\seqsplit{a50662edda954bc58c44184d84591d5b}}. Superseded historical identity — direct correspondence · superseded.}\par
\label{element:3c9464f9eae1415d8fa041bd1f3da4f2}
\index{[CT3] Bed \& Care-Place Turnover Coordination (II.1.3.6.1.2)}
{\small Capability · Proposed · current}\par

Coordinate the operational transition of a released care-place through cleaning, preparation and readiness notifications.



\textbf{Overview}



Coordinate the operational transition of a released care-place through cleaning, preparation and readiness notifications. The reviewed hierarchy assigns this enduring ability to its stated business-enabling context.



\textbf{Intent}



Coordinate the operational transition of a released care-place through cleaning, preparation and readiness notifications.



\textbf{Boundary}



Information administration and coordination preserve the authority of the responsible source, clinical actor and governed domain; they confer no additional real-world management authority.


{\footnotesize \href{../../../../continuum/strategy/capability/3c9464f9eae1415d8fa041bd1f3da4f2.md}{Canonical record: 3c9464f9eae1415d8fa041bd1f3da4f2}; full qualifications and provenance.}\par
\subparagraph*{II.1.3.6.1.2.1 [Feature] Bed Readiness Notification}
\phantomsection
\label{navigation:II.1.3.6.1.2.1}
\addcontentsline{toc}{subparagraph}{II.1.3.6.1.2.1 [Feature] Bed Readiness Notification}

\label{element:6a9cf80d6f82425cb2aaae27f4c5c349}
\index{[Feature] Bed Readiness Notification (II.1.3.6.1.2.1)}
{\small Capability / Feature · Canonical · current}\par

Notify bed allocation managers and clinical units when a cleaned bed achieves READY status.



Readiness notification exposes the established readiness of a care place to the responsible allocation participants. It makes a relevant transition actionable without deciding the clinical suitability of the next allocation or establishing whole-service readiness.


{\footnotesize \href{../../../../continuum/strategy/capability/feature/6a9cf80d6f82425cb2aaae27f4c5c349.md}{Canonical record: 6a9cf80d6f82425cb2aaae27f4c5c349}; full qualifications and provenance.}\par
\subparagraph*{II.1.3.6.1.2.2 [Feature] Bed State Progression}
\phantomsection
\label{navigation:II.1.3.6.1.2.2}
\addcontentsline{toc}{subparagraph}{II.1.3.6.1.2.2 [Feature] Bed State Progression}

\label{element:c903b83181bb456a89dc9d450063eff6}
\index{[Feature] Bed State Progression (II.1.3.6.1.2.2)}
{\small Capability / Feature · Canonical · current}\par

Progress bed states in response to clinical and environmental events (e.g., patient discharge automatically transitions bed to DIRTY).



Bed-state progression applies the consequences of relevant clinical and environmental events to the managed place state. The resulting change is an operational fact, such as a need for cleaning after departure, rather than a transfer of clinical responsibility.


{\footnotesize \href{../../../../continuum/strategy/capability/feature/c903b83181bb456a89dc9d450063eff6.md}{Canonical record: c903b83181bb456a89dc9d450063eff6}; full qualifications and provenance.}\par
\subsubsection*{II.1.3.6.2 [CT2] Care Movement \& Transition Coordination [superseded]}
\phantomsection
\label{navigation:II.1.3.6.2}
\addcontentsline{toc}{subsubsection}{II.1.3.6.2 [CT2] Care Movement \& Transition Coordination [superseded]}

\label{element:59b9331fa3c04ef09bf963ad8a5d0d71}
\index{[CT2] Care Movement \& Transition Coordination (II.1.3.6.2)}
{\small Capability · Proposed · superseded}\par
\href{../../../../continuum/strategy/capability/59b9331fa3c04ef09bf963ad8a5d0d71.md}{Historical canonical definition and provenance}.\par
\subsubsection*{II.1.3.6.3 [CT2] Care Resource \& Capacity Coordination [superseded]}
\phantomsection
\label{navigation:II.1.3.6.3}
\addcontentsline{toc}{subsubsection}{II.1.3.6.3 [CT2] Care Resource \& Capacity Coordination [superseded]}

\label{element:96f71cb8fa6844d7aa7e9f16759a3a62}
\index{[CT2] Care Resource \& Capacity Coordination (II.1.3.6.3)}
{\small Capability · Proposed · superseded}\par
\href{../../../../continuum/strategy/capability/96f71cb8fa6844d7aa7e9f16759a3a62.md}{Historical canonical definition and provenance}.\par
\paragraph*{II.1.3.6.3.1 [CT3] Clinical Resource Management [superseded]}
\phantomsection
\label{navigation:II.1.3.6.3.1}
\addcontentsline{toc}{paragraph}{II.1.3.6.3.1 [CT3] Clinical Resource Management [superseded]}

\label{element:a5faedcd97064c28b9e7e551e9106ea1}
\index{[CT3] Clinical Resource Management (II.1.3.6.3.1)}
{\small Capability · Canonical · superseded}\par
\href{../../../../continuum/strategy/capability/a5faedcd97064c28b9e7e551e9106ea1.md}{Historical canonical definition and provenance}.\par
\paragraph*{II.1.3.6.3.2 [CT3] Service Offering Readiness Coordination}
\phantomsection
\label{navigation:II.1.3.6.3.2}
\addcontentsline{toc}{paragraph}{II.1.3.6.3.2 [CT3] Service Offering Readiness Coordination}

\label{element:0caecaba5d3e4781b8671b12c13ee283}
\index{[CT3] Service Offering Readiness Coordination (II.1.3.6.3.2)}
{\small Capability · Proposed · pending\_adjudication}\par

Assemble the dependency evidence needed to determine and communicate whether a defined service offering is operationally ready in a particular context.



\textbf{Overview}



Collective readiness depends on suitable workforce, spaces, equipment, materials and supporting services. This proposal evaluates their joint availability for an offering without treating one available bed, an operating-hours declaration or aggregate capacity as service readiness.



\textbf{Intent}



Distinguish collective service deliverability from nominal availability or readiness of any single dependency.



\textbf{Boundary}



Contributing capabilities retain source readiness and resource management. Clinical go/no-go decisions and service governance are not granted by this information-based coordination.


{\small Recorded unresolved matters: source-qualified matters; exact scope remains in the canonical record.}\par
{\footnotesize \href{../../../../continuum/strategy/capability/0caecaba5d3e4781b8671b12c13ee283.md}{Canonical record: 0caecaba5d3e4781b8671b12c13ee283}; full qualifications and provenance.}\par
\subparagraph*{II.1.3.6.3.2.1 [Feature] Service Readiness Dependency Assessment}
\phantomsection
\label{navigation:II.1.3.6.3.2.1}
\addcontentsline{toc}{subparagraph}{II.1.3.6.3.2.1 [Feature] Service Readiness Dependency Assessment}

\label{element:4d917929d5cf43218049a84cf46662b0}
\index{[Feature] Service Readiness Dependency Assessment (II.1.3.6.3.2.1)}
{\small Capability / Feature · Proposed · pending\_adjudication}\par

Associate time- and context-qualified dependency evidence to establish a bounded readiness view for an offering.



\textbf{Overview}



An offering may be in operating hours but not deliverable because staffing, equipment, materials or supporting services are missing. A readiness view retains insufficiency and the owners of those dependency assertions.



\textbf{Boundary}



Readiness evidence does not originate clinical clearance or transfer resource management.


{\small Recorded unresolved matters: source-qualified matters; exact scope remains in the canonical record.}\par
{\footnotesize \href{../../../../continuum/strategy/capability/feature/4d917929d5cf43218049a84cf46662b0.md}{Canonical record: 4d917929d5cf43218049a84cf46662b0}; full qualifications and provenance.}\par
\subsubsection*{II.1.3.6.4 [CT2] Care Setting Operations Enablement [superseded]}
\phantomsection
\label{navigation:II.1.3.6.4}
\addcontentsline{toc}{subsubsection}{II.1.3.6.4 [CT2] Care Setting Operations Enablement [superseded]}

\label{element:a03741e6e98642bea44680d23fdcb434}
\index{[CT2] Care Setting Operations Enablement (II.1.3.6.4)}
{\small Capability · Proposed · superseded}\par
\href{../../../../continuum/strategy/capability/a03741e6e98642bea44680d23fdcb434.md}{Historical canonical definition and provenance}.\par
\subsubsection*{II.1.3.6.5 [CT2] Clinic \& Practice Operations}
\phantomsection
\label{navigation:II.1.3.6.5}
\addcontentsline{toc}{subsubsection}{II.1.3.6.5 [CT2] Clinic \& Practice Operations}

{\small Historical predecessor: R.2.123 \hyperref[element:0856d2daf2ba4b92b5390b38b98efb54]{[CT2] Clinic \& Practice Operations}; Canonical ID: \texttt{\seqsplit{0856d2daf2ba4b92b5390b38b98efb54}}. Superseded historical identity — direct correspondence · superseded.}\par
\label{element:3dffe63e03124efab255204ee4e9a5d9}
\index{[CT2] Clinic \& Practice Operations (II.1.3.6.5)}
{\small Capability · Canonical · current}\par

Clinic \& Practice Operations tracks operating sessions and the operational progression of people through the clinic environment. It makes delays, waiting and completion context visible alongside service-specific capacity responsibilities. This is real-world clinic operation, distinct from the technical queueing machinery that might realise the information flow.



\textbf{Overview}



Clinic \& Practice Operations tracks operating sessions and the operational progression of people through the clinic environment. It makes delays, waiting and completion context visible alongside service-specific capacity responsibilities. This is real-world clinic operation, distinct from the technical queueing machinery that might realise the information flow.



\textbf{Intent}



Clinic \& Practice Operations tracks operating sessions and the operational progression of people through the clinic environment. It makes delays, waiting and completion context visible alongside service-specific capacity responsibilities. This is real-world clinic operation, distinct from the technical queueing machinery that might realise the information flow.



\textbf{Boundary}



Information administration does not confer authority over the represented person, organisation, service, place, device or clinical activity.



\medskip\hrule\medskip

{\footnotesize \href{../../../../continuum/strategy/capability/3dffe63e03124efab255204ee4e9a5d9.md}{Canonical record: 3dffe63e03124efab255204ee4e9a5d9}; full qualifications and provenance.}\par
\paragraph*{II.1.3.6.5.1 [CT3] Session \& Queue Coordination}
\phantomsection
\label{navigation:II.1.3.6.5.1}
\addcontentsline{toc}{paragraph}{II.1.3.6.5.1 [CT3] Session \& Queue Coordination}

{\small Historical predecessor: R.2.122 \hyperref[element:5298c6170cef4f05a6d956da716071bd]{[CT3] Session \& Queue Coordination}; Canonical ID: \texttt{\seqsplit{5298c6170cef4f05a6d956da716071bd}}. Superseded historical identity — direct correspondence · superseded.}\par
\label{element:22981a95ca6e466582363706fde8fceb}
\index{[CT3] Session \& Queue Coordination (II.1.3.6.5.1)}
{\small Capability · Proposed · current}\par

Coordinate clinic/practice session state, operational queue progression and context-specific session capacity information.



\textbf{Overview}



Coordinate clinic/practice session state, operational queue progression and context-specific session capacity information. The reviewed hierarchy assigns this enduring ability to its stated business-enabling context.



\textbf{Intent}



Coordinate clinic/practice session state, operational queue progression and context-specific session capacity information.



\textbf{Boundary}



Information administration and coordination preserve the authority of the responsible source, clinical actor and governed domain; they confer no additional real-world management authority.


{\footnotesize \href{../../../../continuum/strategy/capability/22981a95ca6e466582363706fde8fceb.md}{Canonical record: 22981a95ca6e466582363706fde8fceb}; full qualifications and provenance.}\par
\subparagraph*{II.1.3.6.5.1.1 [Feature] Clinic Session Capacity Management}
\phantomsection
\label{navigation:II.1.3.6.5.1.1}
\addcontentsline{toc}{subparagraph}{II.1.3.6.5.1.1 [Feature] Clinic Session Capacity Management}

\label{element:02bdd6a6f22e47f5a29d722ab57f9915}
\index{[Feature] Clinic Session Capacity Management (II.1.3.6.5.1.1)}
{\small Capability / Feature · Proposed · current}\par

Represent available, committed and utilised session capacity under clinic-specific constraints.



\textbf{Overview}



Clinic capacity depends on session staffing and service mix as well as places; queue progression alone is insufficient.



\textbf{Boundary}



Local responsibility retains its capacity meaning; Service Capacity Management consumes it without acquiring ownership or imposing one common capacity model.


{\footnotesize \href{../../../../continuum/strategy/capability/feature/02bdd6a6f22e47f5a29d722ab57f9915.md}{Canonical record: 02bdd6a6f22e47f5a29d722ab57f9915}; full qualifications and provenance.}\par
\subparagraph*{II.1.3.6.5.1.2 [Feature] Clinic Session Tracking}
\phantomsection
\label{navigation:II.1.3.6.5.1.2}
\addcontentsline{toc}{subparagraph}{II.1.3.6.5.1.2 [Feature] Clinic Session Tracking}

\label{element:03e185a6e0d145c3975b3f4715c94207}
\index{[Feature] Clinic Session Tracking (II.1.3.6.5.1.2)}
{\small Capability / Feature · Canonical · current}\par

Ingest and track clinic operating session start, delay, and completion states.



Clinic-session tracking records the operating session's start, delay and completion as service-level context. It is distinct from an individual patient's queue state: knowing where one patient is waiting does not establish whether the clinic session is running as planned.


{\footnotesize \href{../../../../continuum/strategy/capability/feature/03e185a6e0d145c3975b3f4715c94207.md}{Canonical record: 03e185a6e0d145c3975b3f4715c94207}; full qualifications and provenance.}\par
\subparagraph*{II.1.3.6.5.1.3 [Feature] Operational Queue Progression}
\phantomsection
\label{navigation:II.1.3.6.5.1.3}
\addcontentsline{toc}{subparagraph}{II.1.3.6.5.1.3 [Feature] Operational Queue Progression}

\label{element:5852c7d4c9ac4bbf99e7df008d0ed770}
\index{[Feature] Operational Queue Progression (II.1.3.6.5.1.3)}
{\small Capability / Feature · Canonical · current}\par

Track client queue progression through reception, waiting area, consultation, and billing.



Operational queue progression makes the patient's passage through clinic stages visible. The queue concerns real-world arrival, waiting and consultation context, rather than a broker queue, and does not give Harmonia authority to determine clinical priority or own the clinician's worklist.


{\footnotesize \href{../../../../continuum/strategy/capability/feature/5852c7d4c9ac4bbf99e7df008d0ed770.md}{Canonical record: 5852c7d4c9ac4bbf99e7df008d0ed770}; full qualifications and provenance.}\par
\subsubsection*{II.1.3.6.6 [CT2] Clinical Competency Administration}
\phantomsection
\label{navigation:II.1.3.6.6}
\addcontentsline{toc}{subsubsection}{II.1.3.6.6 [CT2] Clinical Competency Administration}

{\small Historical predecessor: R.2.124 \hyperref[element:8a4b918e1c9d47bb99daf2b42008ab39]{[CT2] Clinical Competency Administration}; Canonical ID: \texttt{\seqsplit{8a4b918e1c9d47bb99daf2b42008ab39}}. Superseded historical identity — direct correspondence · superseded.}\par
\label{element:8a99f4a6e5054d6c997bacec04805c01}
\index{[CT2] Clinical Competency Administration (II.1.3.6.6)}
{\small Capability · Canonical · current}\par

Administer information about what clinical capability an organisation can provide and what a practitioner is competent to perform.



\textbf{Overview}



Administer information about what clinical capability an organisation can provide and what a practitioner is competent to perform.



\textbf{Intent}



Administer information about what clinical capability an organisation can provide and what a practitioner is competent to perform.



\textbf{Boundary}



Information administration does not confer authority over the represented person, organisation, service, place, device or clinical activity.



\medskip\hrule\medskip

{\footnotesize \href{../../../../continuum/strategy/capability/8a99f4a6e5054d6c997bacec04805c01.md}{Canonical record: 8a99f4a6e5054d6c997bacec04805c01}; full qualifications and provenance.}\par
\subsubsection*{II.1.3.6.7 [CT2] Clinical Logistics Coordination}
\phantomsection
\label{navigation:II.1.3.6.7}
\addcontentsline{toc}{subsubsection}{II.1.3.6.7 [CT2] Clinical Logistics Coordination}

{\small Historical predecessor: R.2.127 \hyperref[element:90a782ff2d004fa984ae0465528b2253]{[CT2] Clinical Logistics Coordination}; Canonical ID: \texttt{\seqsplit{90a782ff2d004fa984ae0465528b2253}}. Superseded historical identity — direct correspondence · superseded.}\par
\label{element:0c27be11b4a44cb19e992309b8146a17}
\index{[CT2] Clinical Logistics Coordination (II.1.3.6.7)}
{\small Capability · Canonical · current}\par

Coordinate specimen metadata and specimen movement traceability across responsible collection, courier and diagnostic services.



\textbf{Overview}



Coordinate specimen metadata and specimen movement traceability across responsible collection, courier and diagnostic services.



\textbf{Intent}



Coordinate specimen metadata and specimen movement traceability across responsible collection, courier and diagnostic services.



\textbf{Boundary}



Includes specimen metadata and specimen movement traceability. Supply/materiel movement belongs to Health Materiel Traceability Administration.



\medskip\hrule\medskip

{\footnotesize \href{../../../../continuum/strategy/capability/0c27be11b4a44cb19e992309b8146a17.md}{Canonical record: 0c27be11b4a44cb19e992309b8146a17}; full qualifications and provenance.}\par
\paragraph*{II.1.3.6.7.1 [CT3] Specimen Metadata Management}
\phantomsection
\label{navigation:II.1.3.6.7.1}
\addcontentsline{toc}{paragraph}{II.1.3.6.7.1 [CT3] Specimen Metadata Management}

{\small Historical predecessor: R.2.125 \hyperref[element:3b000d3712534d16a48470f79df8e33a]{[CT3] Specimen Metadata Management}; Canonical ID: \texttt{\seqsplit{3b000d3712534d16a48470f79df8e33a}}. Superseded historical identity — direct correspondence · superseded.}\par
\label{element:738736425b7446fdae47b06c91c4f69c}
\index{[CT3] Specimen Metadata Management (II.1.3.6.7.1)}
{\small Capability · Proposed · current}\par

Manage the identifying, descriptive, custody and relevant handling information needed to interpret a transported specimen.



\textbf{Overview}



Manage the identifying, descriptive, custody and relevant handling information needed to interpret a transported specimen. The reviewed hierarchy assigns this enduring ability to its stated business-enabling context.



\textbf{Intent}



Manage the identifying, descriptive, custody and relevant handling information needed to interpret a transported specimen.



\textbf{Boundary}



Information administration and coordination preserve the authority of the responsible source, clinical actor and governed domain; they confer no additional real-world management authority.


{\footnotesize \href{../../../../continuum/strategy/capability/738736425b7446fdae47b06c91c4f69c.md}{Canonical record: 738736425b7446fdae47b06c91c4f69c}; full qualifications and provenance.}\par
\paragraph*{II.1.3.6.7.2 [CT3] Specimen Movement Traceability}
\phantomsection
\label{navigation:II.1.3.6.7.2}
\addcontentsline{toc}{paragraph}{II.1.3.6.7.2 [CT3] Specimen Movement Traceability}

{\small Historical predecessor: R.2.126 \hyperref[element:129119d7b2744ffd805c58f25557fa27]{[CT3] Specimen Movement Traceability}; Canonical ID: \texttt{\seqsplit{129119d7b2744ffd805c58f25557fa27}}. Superseded historical identity — direct correspondence · superseded.}\par
\label{element:bfa2e9a84ddd4287897573dc0d917a00}
\index{[CT3] Specimen Movement Traceability (II.1.3.6.7.2)}
{\small Capability · Proposed · current}\par

Maintain attributable pickup, custody, transport-condition and delivery traceability for clinical specimens.



\textbf{Overview}



Maintain attributable pickup, custody, transport-condition and delivery traceability for clinical specimens. The reviewed hierarchy assigns this enduring ability to its stated business-enabling context.



\textbf{Intent}



Maintain attributable pickup, custody, transport-condition and delivery traceability for clinical specimens.



\textbf{Boundary}



Information administration and coordination preserve the authority of the responsible source, clinical actor and governed domain; they confer no additional real-world management authority.


{\footnotesize \href{../../../../continuum/strategy/capability/bfa2e9a84ddd4287897573dc0d917a00.md}{Canonical record: bfa2e9a84ddd4287897573dc0d917a00}; full qualifications and provenance.}\par
\subparagraph*{II.1.3.6.7.2.1 [Feature] Pathology Specimen Transport Tracking}
\phantomsection
\label{navigation:II.1.3.6.7.2.1}
\addcontentsline{toc}{subparagraph}{II.1.3.6.7.2.1 [Feature] Pathology Specimen Transport Tracking}

\label{element:34f182cfc6b945d090a880943bbb4959}
\index{[Feature] Pathology Specimen Transport Tracking (II.1.3.6.7.2.1)}
{\small Capability / Feature · Canonical · current}\par

Track pickup, custody chain, transport conditions, and laboratory delivery of clinical specimens.



Specimen transport tracking retains pickup, custody and condition context through laboratory delivery. It supports interpretation of the material's journey without treating a delivered package as proof of specimen suitability, completed analysis or an authorised diagnostic report.


{\footnotesize \href{../../../../continuum/strategy/capability/feature/34f182cfc6b945d090a880943bbb4959.md}{Canonical record: 34f182cfc6b945d090a880943bbb4959}; full qualifications and provenance.}\par
\subsubsection*{II.1.3.6.8 [CT2] Clinical Qualification Administration}
\phantomsection
\label{navigation:II.1.3.6.8}
\addcontentsline{toc}{subsubsection}{II.1.3.6.8 [CT2] Clinical Qualification Administration}

{\small Historical predecessor: R.2.130 \hyperref[element:5fea7aaac4384245b1f08a3d8a076bd9]{[CT2] Clinical Qualification Administration}; Canonical ID: \texttt{\seqsplit{5fea7aaac4384245b1f08a3d8a076bd9}}. Superseded historical identity — direct correspondence · superseded.}\par
\label{element:0511a705d74a4f48a0ce92590ae764fa}
\index{[CT2] Clinical Qualification Administration (II.1.3.6.8)}
{\small Capability · Canonical · current}\par

Administer organisation clinical certification/compliance information and provider qualification evidence with its externally established standing.



\textbf{Overview}



Administer organisation clinical certification/compliance information and provider qualification evidence with its externally established standing.



\textbf{Intent}



Administer organisation clinical certification/compliance information and provider qualification evidence with its externally established standing.



\textbf{Boundary}



Information administration does not confer authority over the represented person, organisation, service, place, device or clinical activity.



\medskip\hrule\medskip

{\footnotesize \href{../../../../continuum/strategy/capability/0511a705d74a4f48a0ce92590ae764fa.md}{Canonical record: 0511a705d74a4f48a0ce92590ae764fa}; full qualifications and provenance.}\par
\paragraph*{II.1.3.6.8.1 [CT3] Organisation Clinical Certification \& Compliance Support}
\phantomsection
\label{navigation:II.1.3.6.8.1}
\addcontentsline{toc}{paragraph}{II.1.3.6.8.1 [CT3] Organisation Clinical Certification \& Compliance Support}

{\small Historical predecessor: R.2.128 \hyperref[element:cee76a930a4e4e7d9c49357ba290c8bf]{[CT3] Organisation Clinical Certification \& Compliance Support}; Canonical ID: \texttt{\seqsplit{cee76a930a4e4e7d9c49357ba290c8bf}}. Superseded historical identity — direct correspondence · superseded.}\par
\label{element:4de624e7bf49472f91e3e5c734ee8725}
\index{[CT3] Organisation Clinical Certification \& Compliance Support (II.1.3.6.8.1)}
{\small Capability · Proposed · current}\par

Administer supplied organisation clinical certification and conformance evidence for use in service qualification and participation.



\textbf{Overview}



Administer supplied organisation clinical certification and conformance evidence for use in service qualification and participation. The reviewed hierarchy assigns this enduring ability to its stated business-enabling context.



\textbf{Intent}



Administer supplied organisation clinical certification and conformance evidence for use in service qualification and participation.



\textbf{Boundary}



Information administration and coordination preserve the authority of the responsible source, clinical actor and governed domain; they confer no additional real-world management authority.


{\footnotesize \href{../../../../continuum/strategy/capability/4de624e7bf49472f91e3e5c734ee8725.md}{Canonical record: 4de624e7bf49472f91e3e5c734ee8725}; full qualifications and provenance.}\par
\paragraph*{II.1.3.6.8.2 [CT3] Provider Clinical Qualification Tracking}
\phantomsection
\label{navigation:II.1.3.6.8.2}
\addcontentsline{toc}{paragraph}{II.1.3.6.8.2 [CT3] Provider Clinical Qualification Tracking}

{\small Historical predecessor: R.2.129 \hyperref[element:ab0415e2ea0a4da4ad3dc46997d48bc3]{[CT3] Provider Clinical Qualification Tracking}; Canonical ID: \texttt{\seqsplit{ab0415e2ea0a4da4ad3dc46997d48bc3}}. Superseded historical identity — direct correspondence · superseded.}\par
\label{element:8fd5c78eba3046fa9153fb4efb44e9c8}
\index{[CT3] Provider Clinical Qualification Tracking (II.1.3.6.8.2)}
{\small Capability · Proposed · current}\par

Track supplied provider clinical qualifications, competency certificates and privilege evidence with effective dates and issuing-authority standing.



\textbf{Overview}



Track supplied provider clinical qualifications, competency certificates and privilege evidence with effective dates and issuing-authority standing. The reviewed hierarchy assigns this enduring ability to its stated business-enabling context.



\textbf{Intent}



Track supplied provider clinical qualifications, competency certificates and privilege evidence with effective dates and issuing-authority standing.



\textbf{Boundary}



Information administration and coordination preserve the authority of the responsible source, clinical actor and governed domain; they confer no additional real-world management authority.


{\footnotesize \href{../../../../continuum/strategy/capability/8fd5c78eba3046fa9153fb4efb44e9c8.md}{Canonical record: 8fd5c78eba3046fa9153fb4efb44e9c8}; full qualifications and provenance.}\par
\subparagraph*{II.1.3.6.8.2.1 [Feature] Clinical Privileging Record}
\phantomsection
\label{navigation:II.1.3.6.8.2.1}
\addcontentsline{toc}{subparagraph}{II.1.3.6.8.2.1 [Feature] Clinical Privileging Record}

\label{element:5073df3c6bb94c218ca19b74550b19fa}
\index{[Feature] Clinical Privileging Record (II.1.3.6.8.2.1)}
{\small Capability / Feature · Canonical · current}\par

Track active clinical admitting and practice privileges within healthcare facilities.



The privileging record makes applicable admitting and practice permissions available as contextual evidence. It represents the health service's established privilege standing rather than granting new authority merely because a record exists in the directory.


{\footnotesize \href{../../../../continuum/strategy/capability/feature/5073df3c6bb94c218ca19b74550b19fa.md}{Canonical record: 5073df3c6bb94c218ca19b74550b19fa}; full qualifications and provenance.}\par
\subparagraph*{II.1.3.6.8.2.2 [Feature] Operational Credential Check}
\phantomsection
\label{navigation:II.1.3.6.8.2.2}
\addcontentsline{toc}{subparagraph}{II.1.3.6.8.2.2 [Feature] Operational Credential Check}

\label{element:1a76831985fa461c891a3428e5dea86d}
\index{[Feature] Operational Credential Check (II.1.3.6.8.2.2)}
{\small Capability / Feature · Canonical · current}\par

Validate that a practitioner possesses active credentials and required competency certs prior to dispatching specialized work.



The qualification checkpoint considers active professional credentials and competency evidence before specialised dispatch. Its historical credential terminology does not make it the operational badge or physical-access responsibility of Credential Management, nor does a qualification check alone establish all required privilege.


{\footnotesize \href{../../../../continuum/strategy/capability/feature/1a76831985fa461c891a3428e5dea86d.md}{Canonical record: 1a76831985fa461c891a3428e5dea86d}; full qualifications and provenance.}\par
\subsubsection*{II.1.3.6.9 [CT2] Clinical Rostering}
\phantomsection
\label{navigation:II.1.3.6.9}
\addcontentsline{toc}{subsubsection}{II.1.3.6.9 [CT2] Clinical Rostering}

{\small Historical predecessor: R.2.132 \hyperref[element:b6527516a69c499d8cc9c8fe881169fe]{[CT2] Clinical Rostering}; Canonical ID: \texttt{\seqsplit{b6527516a69c499d8cc9c8fe881169fe}}. Superseded historical identity — direct correspondence · superseded.}\par
\label{element:192ca88a289544dcbf6adaf643f39969}
\index{[CT2] Clinical Rostering (II.1.3.6.9)}
{\small Capability · Canonical · current}\par

Clinical Rostering provides the Adjacent context of planned workforce availability. It is distinct from consuming roster projections for operational views or resolving who is currently on call. The present architecture retains the context without establishing a Harmonia Feature decomposition for authoring or managing the authoritative roster.



\textbf{Overview}



Clinical Rostering provides the Adjacent context of planned workforce availability. It is distinct from consuming roster projections for operational views or resolving who is currently on call. The present architecture retains the context without establishing a Harmonia Feature decomposition for authoring or managing the authoritative roster.



\textbf{Intent}



Clinical Rostering provides the Adjacent context of planned workforce availability. It is distinct from consuming roster projections for operational views or resolving who is currently on call. The present architecture retains the context without establishing a Harmonia Feature decomposition for authoring or managing the authoritative roster.



\textbf{Boundary}



Information administration does not confer authority over the represented person, organisation, service, place, device or clinical activity.



\medskip\hrule\medskip

{\footnotesize \href{../../../../continuum/strategy/capability/192ca88a289544dcbf6adaf643f39969.md}{Canonical record: 192ca88a289544dcbf6adaf643f39969}; full qualifications and provenance.}\par
\paragraph*{II.1.3.6.9.1 [CT3] Clinical Roster Capture \& Communication}
\phantomsection
\label{navigation:II.1.3.6.9.1}
\addcontentsline{toc}{paragraph}{II.1.3.6.9.1 [CT3] Clinical Roster Capture \& Communication}

{\small Historical predecessor: R.2.131 \hyperref[element:2054d57bb33a4361b3ef0824545a1784]{[CT3] Clinical Roster Capture \& Communication}; Canonical ID: \texttt{\seqsplit{2054d57bb33a4361b3ef0824545a1784}}. Superseded historical identity — direct correspondence · superseded.}\par
\label{element:8243cba81fb6449781d10823ae194237}
\index{[CT3] Clinical Roster Capture \& Communication (II.1.3.6.9.1)}
{\small Capability · Proposed · current}\par

Capture and communicate supplied clinical roster information, its effective periods and subsequent coverage changes.



\textbf{Overview}



Capture and communicate supplied clinical roster information, its effective periods and subsequent coverage changes. The reviewed hierarchy assigns this enduring ability to its stated business-enabling context.



\textbf{Intent}



Capture and communicate supplied clinical roster information, its effective periods and subsequent coverage changes.



\textbf{Boundary}



Information administration and coordination preserve the authority of the responsible source, clinical actor and governed domain; they confer no additional real-world management authority.


{\footnotesize \href{../../../../continuum/strategy/capability/8243cba81fb6449781d10823ae194237.md}{Canonical record: 8243cba81fb6449781d10823ae194237}; full qualifications and provenance.}\par
\subsubsection*{II.1.3.6.10 [CT2] Clinical Supervision}
\phantomsection
\label{navigation:II.1.3.6.10}
\addcontentsline{toc}{subsubsection}{II.1.3.6.10 [CT2] Clinical Supervision}

{\small Historical predecessor: R.2.134 \hyperref[element:50c838981f3946d78e8f78d18e31c70c]{[CT2] Clinical Supervision}; Canonical ID: \texttt{\seqsplit{50c838981f3946d78e8f78d18e31c70c}}. Superseded historical identity — direct correspondence · superseded.}\par
\label{element:f697b1ac79b544749315f31819b17969}
\index{[CT2] Clinical Supervision (II.1.3.6.10)}
{\small Capability · Canonical · current}\par

Enable responsible clinical participants to identify and use the applicable clinical escalation path.



\textbf{Overview}



Enable responsible clinical participants to identify and use the applicable clinical escalation path.



\textbf{Intent}



Enable responsible clinical participants to identify and use the applicable clinical escalation path.



\textbf{Boundary}



Information administration does not confer authority over the represented person, organisation, service, place, device or clinical activity.



\medskip\hrule\medskip

{\footnotesize \href{../../../../continuum/strategy/capability/f697b1ac79b544749315f31819b17969.md}{Canonical record: f697b1ac79b544749315f31819b17969}; full qualifications and provenance.}\par
\paragraph*{II.1.3.6.10.1 [CT3] Clinical Escalation Path Management}
\phantomsection
\label{navigation:II.1.3.6.10.1}
\addcontentsline{toc}{paragraph}{II.1.3.6.10.1 [CT3] Clinical Escalation Path Management}

{\small Historical predecessor: R.2.133 \hyperref[element:759a2e52172345e8ad1bc96d7e5ddad3]{[CT3] Clinical Escalation Path Management}; Canonical ID: \texttt{\seqsplit{759a2e52172345e8ad1bc96d7e5ddad3}}. Superseded historical identity — direct correspondence · superseded.}\par
\label{element:cca86a32ba8c4dda84caeaa248e9f444}
\index{[CT3] Clinical Escalation Path Management (II.1.3.6.10.1)}
{\small Capability · Proposed · current}\par

Manage clinical escalation paths and their applicability, contacts and progression conditions for use by responsible clinical participants.



\textbf{Overview}



Manage clinical escalation paths and their applicability, contacts and progression conditions for use by responsible clinical participants. The reviewed hierarchy assigns this enduring ability to its stated business-enabling context.



\textbf{Intent}



Manage clinical escalation paths and their applicability, contacts and progression conditions for use by responsible clinical participants.



\textbf{Boundary}



Harmonia manages the escalation path; responsible clinical participants perform supervision and retain clinical judgement.


{\footnotesize \href{../../../../continuum/strategy/capability/cca86a32ba8c4dda84caeaa248e9f444.md}{Canonical record: cca86a32ba8c4dda84caeaa248e9f444}; full qualifications and provenance.}\par
\subsubsection*{II.1.3.6.11 [CT2] Clinical Training \& Education}
\phantomsection
\label{navigation:II.1.3.6.11}
\addcontentsline{toc}{subsubsection}{II.1.3.6.11 [CT2] Clinical Training \& Education}

{\small Historical predecessor: R.2.135 \hyperref[element:df8bff7df9c1469da16cac6a5eb2bf76]{[CT2] Clinical Training \& Education}; Canonical ID: \texttt{\seqsplit{df8bff7df9c1469da16cac6a5eb2bf76}}. Superseded historical identity — direct correspondence · superseded.}\par
\label{element:d016f35ffbc345d99d7281e968c14ddb}
\index{[CT2] Clinical Training \& Education (II.1.3.6.11)}
{\small Capability · Canonical · current}\par

Clinical Training \& Education remains a Reference context for developing practitioners' knowledge and skills. The catalogue recognises its enterprise significance without assigning training delivery or a Feature decomposition to Harmonia. Its separation from qualification checks prevents a represented certificate or training record from being mistaken for the education activity itself.



\textbf{Overview}



Clinical Training \& Education remains a Reference context for developing practitioners' knowledge and skills. The catalogue recognises its enterprise significance without assigning training delivery or a Feature decomposition to Harmonia. Its separation from qualification checks prevents a represented certificate or training record from being mistaken for the education activity itself.



\textbf{Intent}



Clinical Training \& Education remains a Reference context for developing practitioners' knowledge and skills. The catalogue recognises its enterprise significance without assigning training delivery or a Feature decomposition to Harmonia. Its separation from qualification checks prevents a represented certificate or training record from being mistaken for the education activity itself.



\textbf{Boundary}



Information administration does not confer authority over the represented person, organisation, service, place, device or clinical activity.



\medskip\hrule\medskip

{\footnotesize \href{../../../../continuum/strategy/capability/d016f35ffbc345d99d7281e968c14ddb.md}{Canonical record: d016f35ffbc345d99d7281e968c14ddb}; full qualifications and provenance.}\par
\subsubsection*{II.1.3.6.12 [CT2] Clinical Workforce Capability Monitoring}
\phantomsection
\label{navigation:II.1.3.6.12}
\addcontentsline{toc}{subsubsection}{II.1.3.6.12 [CT2] Clinical Workforce Capability Monitoring}

{\small Historical predecessor: R.2.136 \hyperref[element:24c015676c254e2e82e8e4287c295179]{[CT2] Clinical Workforce Capability Monitoring}; Canonical ID: \texttt{\seqsplit{24c015676c254e2e82e8e4287c295179}}. Superseded historical identity — direct correspondence · superseded.}\par
\label{element:367b3e5643f44fa785f093db8bb9a6e8}
\index{[CT2] Clinical Workforce Capability Monitoring (II.1.3.6.12)}
{\small Capability · Canonical · current}\par

Monitor supplied information about the clinical capability available from the workforce without acquiring employment or clinical allocation authority.



\textbf{Overview}



Monitor supplied information about the clinical capability available from the workforce without acquiring employment or clinical allocation authority.



\textbf{Intent}



Monitor supplied information about the clinical capability available from the workforce without acquiring employment or clinical allocation authority.



\textbf{Boundary}



Information administration does not confer authority over the represented person, organisation, service, place, device or clinical activity.



\medskip\hrule\medskip

{\footnotesize \href{../../../../continuum/strategy/capability/367b3e5643f44fa785f093db8bb9a6e8.md}{Canonical record: 367b3e5643f44fa785f093db8bb9a6e8}; full qualifications and provenance.}\par
\subsubsection*{II.1.3.6.13 [CT2] Credential Management}
\phantomsection
\label{navigation:II.1.3.6.13}
\addcontentsline{toc}{subsubsection}{II.1.3.6.13 [CT2] Credential Management}

{\small Historical predecessor: R.2.138 \hyperref[element:bab3d0ab6eb748978b962ea20fb31640]{[CT2] Credential Management}; Canonical ID: \texttt{\seqsplit{bab3d0ab6eb748978b962ea20fb31640}}. Superseded historical identity — direct correspondence · superseded.}\par
\label{element:21eed4fc8c924496b16c6ddaa32f3d01}
\index{[CT2] Credential Management (II.1.3.6.13)}
{\small Capability · Proposed · current}\par

Administer definitions of clinical-service credentials and their required applicability without granting professional privilege or managing physical badges.



\textbf{Overview}



Administer definitions of clinical-service credentials and their required applicability without granting professional privilege or managing physical badges. The reviewed hierarchy assigns this enduring ability to its stated business-enabling context.



\textbf{Intent}



Administer definitions of clinical-service credentials and their required applicability without granting professional privilege or managing physical badges.



\textbf{Boundary}



Information administration and coordination preserve the authority of the responsible source, clinical actor and governed domain; they confer no additional real-world management authority.


{\footnotesize \href{../../../../continuum/strategy/capability/21eed4fc8c924496b16c6ddaa32f3d01.md}{Canonical record: 21eed4fc8c924496b16c6ddaa32f3d01}; full qualifications and provenance.}\par
\paragraph*{II.1.3.6.13.1 [CT3] Clinical Service Credential Definition Management}
\phantomsection
\label{navigation:II.1.3.6.13.1}
\addcontentsline{toc}{paragraph}{II.1.3.6.13.1 [CT3] Clinical Service Credential Definition Management}

{\small Historical predecessor: R.2.137 \hyperref[element:0115ed50867d468a95c0fafa1c5d79e2]{[CT3] Clinical Service Credential Definition Management}; Canonical ID: \texttt{\seqsplit{0115ed50867d468a95c0fafa1c5d79e2}}. Superseded historical identity — direct correspondence · superseded.}\par
\label{element:2e5c97f51e7e48d1b865ac416dd8099d}
\index{[CT3] Clinical Service Credential Definition Management (II.1.3.6.13.1)}
{\small Capability · Proposed · current}\par

Manage the definitions of credentials required for participation in a clinical service, distinct from physical badges and authentication credentials.



\textbf{Overview}



Manage the definitions of credentials required for participation in a clinical service, distinct from physical badges and authentication credentials. The reviewed hierarchy assigns this enduring ability to its stated business-enabling context.



\textbf{Intent}



Manage the definitions of credentials required for participation in a clinical service, distinct from physical badges and authentication credentials.



\textbf{Boundary}



Information administration and coordination preserve the authority of the responsible source, clinical actor and governed domain; they confer no additional real-world management authority.


{\footnotesize \href{../../../../continuum/strategy/capability/2e5c97f51e7e48d1b865ac416dd8099d.md}{Canonical record: 2e5c97f51e7e48d1b865ac416dd8099d}; full qualifications and provenance.}\par
\subsubsection*{II.1.3.6.14 [CT2] Discharge Management}
\phantomsection
\label{navigation:II.1.3.6.14}
\addcontentsline{toc}{subsubsection}{II.1.3.6.14 [CT2] Discharge Management}

{\small Historical predecessor: R.2.140 \hyperref[element:c57c0721ab0040ccb24e88e5e92e8cf8]{[CT2] Discharge Management}; Canonical ID: \texttt{\seqsplit{c57c0721ab0040ccb24e88e5e92e8cf8}}. Superseded historical identity — direct correspondence · superseded.}\par
\label{element:ec5e8d6732764aa489c2dd7630fc8fef}
\index{[CT2] Discharge Management (II.1.3.6.14)}
{\small Capability · Canonical · current}\par

Discharge Management coordinates the reported milestones needed around discharge, including medication supply, transport and home-care confirmations. It makes readiness and impediments visible across disciplines and communicates departure information when physical exit occurs. The responsibility preserves a distinction between gathering confirmations and making the clinical decision that discharge is appropriate.



\textbf{Overview}



Discharge Management coordinates the reported milestones needed around discharge, including medication supply, transport and home-care confirmations. It makes readiness and impediments visible across disciplines and communicates departure information when physical exit occurs. The responsibility preserves a distinction between gathering confirmations and making the clinical decision that discharge is appropriate.



\textbf{Intent}



Discharge Management coordinates the reported milestones needed around discharge, including medication supply, transport and home-care confirmations. It makes readiness and impediments visible across disciplines and communicates departure information when physical exit occurs. The responsibility preserves a distinction between gathering confirmations and making the clinical decision that discharge is appropriate.



\textbf{Boundary}



Information administration does not confer authority over the represented person, organisation, service, place, device or clinical activity.



\medskip\hrule\medskip

{\footnotesize \href{../../../../continuum/strategy/capability/ec5e8d6732764aa489c2dd7630fc8fef.md}{Canonical record: ec5e8d6732764aa489c2dd7630fc8fef}; full qualifications and provenance.}\par
\paragraph*{II.1.3.6.14.1 [CT3] Discharge Progression \& Coordination}
\phantomsection
\label{navigation:II.1.3.6.14.1}
\addcontentsline{toc}{paragraph}{II.1.3.6.14.1 [CT3] Discharge Progression \& Coordination}

{\small Historical predecessor: R.2.139 \hyperref[element:68344272ca6049648a3ddae20c150f37]{[CT3] Discharge Progression \& Coordination}; Canonical ID: \texttt{\seqsplit{68344272ca6049648a3ddae20c150f37}}. Superseded historical identity — direct correspondence · superseded.}\par
\label{element:d6048993146b45639bab27860f94f34f}
\index{[CT3] Discharge Progression \& Coordination (II.1.3.6.14.1)}
{\small Capability · Proposed · current}\par

Coordinate externally established discharge milestones, multidisciplinary readiness information and departure notifications.



\textbf{Overview}



Coordinate externally established discharge milestones, multidisciplinary readiness information and departure notifications. The reviewed hierarchy assigns this enduring ability to its stated business-enabling context.



\textbf{Intent}



Coordinate externally established discharge milestones, multidisciplinary readiness information and departure notifications.



\textbf{Boundary}



Information administration and coordination preserve the authority of the responsible source, clinical actor and governed domain; they confer no additional real-world management authority.


{\footnotesize \href{../../../../continuum/strategy/capability/d6048993146b45639bab27860f94f34f.md}{Canonical record: d6048993146b45639bab27860f94f34f}; full qualifications and provenance.}\par
\subparagraph*{II.1.3.6.14.1.1 [Feature] Discharge Milestone Tracking}
\phantomsection
\label{navigation:II.1.3.6.14.1.1}
\addcontentsline{toc}{subparagraph}{II.1.3.6.14.1.1 [Feature] Discharge Milestone Tracking}

\label{element:887eec1336934b13a3f21a8e3eb5f131}
\index{[Feature] Discharge Milestone Tracking (II.1.3.6.14.1.1)}
{\small Capability / Feature · Canonical · current}\par

Track clinical criteria-led discharge milestones (medications dispensed, transportation arranged, home care confirmed).



Discharge milestone tracking records confirmations contributed by the relevant clinical and operational participants. It makes outstanding requirements visible without allowing the tracking mechanism to declare clinical readiness or assume that every local milestone defines a universal discharge pathway.


{\footnotesize \href{../../../../continuum/strategy/capability/feature/887eec1336934b13a3f21a8e3eb5f131.md}{Canonical record: 887eec1336934b13a3f21a8e3eb5f131}; full qualifications and provenance.}\par
\subparagraph*{II.1.3.6.14.1.2 [Feature] Discharge Notification Dispatch}
\phantomsection
\label{navigation:II.1.3.6.14.1.2}
\addcontentsline{toc}{subparagraph}{II.1.3.6.14.1.2 [Feature] Discharge Notification Dispatch}

\label{element:5476307ca5a34590aa7025ab59ca0c10}
\index{[Feature] Discharge Notification Dispatch (II.1.3.6.14.1.2)}
{\small Capability / Feature · Canonical · current}\par

Dispatch departure notifications, updated bed states, and finalized discharge summaries to external care providers upon physical patient exit.



Departure notification communicates the physical exit and associated record and place-state context to the applicable participants. It distinguishes departure information from earlier readiness and from the clinical decision, while retaining the separate question of responsibility actually assumed elsewhere.


{\footnotesize \href{../../../../continuum/strategy/capability/feature/5476307ca5a34590aa7025ab59ca0c10.md}{Canonical record: 5476307ca5a34590aa7025ab59ca0c10}; full qualifications and provenance.}\par
\subparagraph*{II.1.3.6.14.1.3 [Feature] Discharge Readiness \& Coordination Oversight}
\phantomsection
\label{navigation:II.1.3.6.14.1.3}
\addcontentsline{toc}{subparagraph}{II.1.3.6.14.1.3 [Feature] Discharge Readiness \& Coordination Oversight}

\label{element:5499422ed3644e0bbf7645f9ae3f20cb}
\index{[Feature] Discharge Readiness \& Coordination Oversight (II.1.3.6.14.1.3)}
{\small Capability / Feature · Canonical · current}\par

Track and aggregate multidisciplinary discharge milestones, home care confirmations, and transit readiness across care teams.



Readiness oversight brings multidisciplinary, home-care and transport confirmations into one coordination view. Aggregation supports the responsible participants' decisions but does not itself confer authority to discharge or establish acceptance of continuing care by another provider.


{\footnotesize \href{../../../../continuum/strategy/capability/feature/5499422ed3644e0bbf7645f9ae3f20cb.md}{Canonical record: 5499422ed3644e0bbf7645f9ae3f20cb}; full qualifications and provenance.}\par
\subsubsection*{II.1.3.6.15 [CT2] Emergency Department Operations}
\phantomsection
\label{navigation:II.1.3.6.15}
\addcontentsline{toc}{subsubsection}{II.1.3.6.15 [CT2] Emergency Department Operations}

{\small Historical predecessor: R.2.142 \hyperref[element:8963b571e9cc40e49c5d738e31d2f97e]{[CT2] Emergency Department Operations}; Canonical ID: \texttt{\seqsplit{8963b571e9cc40e49c5d738e31d2f97e}}. Superseded historical identity — direct correspondence · superseded.}\par
\label{element:24a4b5815faa4d369bc756a6e97d3e9e}
\index{[CT2] Emergency Department Operations (II.1.3.6.15)}
{\small Capability · Canonical · current}\par

This capability maintains the operational context of emergency places, waits and length-of-stay milestones. It supports understanding of occupancy and impediments in the department while keeping capacity's emergency-service meaning intact. Its concern differs from Emergency Care's clinical presentation context, even though both can refer to the same patient encounter.



\textbf{Overview}



This capability maintains the operational context of emergency places, waits and length-of-stay milestones. It supports understanding of occupancy and impediments in the department while keeping capacity's emergency-service meaning intact. Its concern differs from Emergency Care's clinical presentation context, even though both can refer to the same patient encounter.



\textbf{Intent}



This capability maintains the operational context of emergency places, waits and length-of-stay milestones. It supports understanding of occupancy and impediments in the department while keeping capacity's emergency-service meaning intact. Its concern differs from Emergency Care's clinical presentation context, even though both can refer to the same patient encounter.



\textbf{Boundary}



Information administration does not confer authority over the represented person, organisation, service, place, device or clinical activity.



\medskip\hrule\medskip

{\footnotesize \href{../../../../continuum/strategy/capability/24a4b5815faa4d369bc756a6e97d3e9e.md}{Canonical record: 24a4b5815faa4d369bc756a6e97d3e9e}; full qualifications and provenance.}\par
\paragraph*{II.1.3.6.15.1 [CT3] Emergency Care Progression}
\phantomsection
\label{navigation:II.1.3.6.15.1}
\addcontentsline{toc}{paragraph}{II.1.3.6.15.1 [CT3] Emergency Care Progression}

\label{element:7125e459a48946d0ba075b5ecc3571b4}
\index{[CT3] Emergency Care Progression (II.1.3.6.15.1)}
{\small Capability · Canonical · current}\par

Definition unresolved in the baseline. No definition has been synthesised from the name.



\medskip\hrule\medskip

{\footnotesize \href{../../../../continuum/strategy/capability/7125e459a48946d0ba075b5ecc3571b4.md}{Canonical record: 7125e459a48946d0ba075b5ecc3571b4}; full qualifications and provenance.}\par
\subparagraph*{II.1.3.6.15.1.1 [Feature] Emergency Care Cross-Facility Correlation}
\phantomsection
\label{navigation:II.1.3.6.15.1.1}
\addcontentsline{toc}{subparagraph}{II.1.3.6.15.1.1 [Feature] Emergency Care Cross-Facility Correlation}

\label{element:6140586e2b18424d845a1f07018e6b49}
\index{[Feature] Emergency Care Cross-Facility Correlation (II.1.3.6.15.1.1)}
{\small Capability / Feature · Canonical · current}\par

Correlate incoming emergency presentations with recent regional admissions, prior discharge summaries, and known allergy profiles.



Cross-facility correlation connects an emergency presentation with relevant recent history and allergy information. It supports situational understanding across institutional boundaries without deciding the clinical significance of the combined evidence or assigning emergency treatment.


{\footnotesize \href{../../../../continuum/strategy/capability/feature/6140586e2b18424d845a1f07018e6b49.md}{Canonical record: 6140586e2b18424d845a1f07018e6b49}; full qualifications and provenance.}\par
\subparagraph*{II.1.3.6.15.1.2 [Feature] Emergency Department Arrival \& Triage Ingestion}
\phantomsection
\label{navigation:II.1.3.6.15.1.2}
\addcontentsline{toc}{subparagraph}{II.1.3.6.15.1.2 [Feature] Emergency Department Arrival \& Triage Ingestion}

\label{element:1eb4053512164f6f976a9767da9ba7a7}
\index{[Feature] Emergency Department Arrival \& Triage Ingestion (II.1.3.6.15.1.2)}
{\small Capability / Feature · Canonical · current}\par

Ingest emergency presentation, triage category, and chief complaint records.



Emergency presentation ingestion retains the arrival, triage category and presenting complaint reported by the emergency context. It makes that event available for regional correlation while preserving that the clinical triage judgement was established by the responsible source.


{\footnotesize \href{../../../../continuum/strategy/capability/feature/1eb4053512164f6f976a9767da9ba7a7.md}{Canonical record: 1eb4053512164f6f976a9767da9ba7a7}; full qualifications and provenance.}\par
\paragraph*{II.1.3.6.15.2 [CT3] Emergency Department Operational Coordination}
\phantomsection
\label{navigation:II.1.3.6.15.2}
\addcontentsline{toc}{paragraph}{II.1.3.6.15.2 [CT3] Emergency Department Operational Coordination}

{\small Historical predecessor: R.2.141 \hyperref[element:15b6d2518e60444098b7c546ca31c49b]{[CT3] Emergency Department Operational Coordination}; Canonical ID: \texttt{\seqsplit{15b6d2518e60444098b7c546ca31c49b}}. Superseded historical identity — direct correspondence · superseded.}\par
\label{element:8f0b0f3d9e8d44a7827b0835870b4a4f}
\index{[CT3] Emergency Department Operational Coordination (II.1.3.6.15.2)}
{\small Capability · Proposed · current}\par

Coordinate emergency-department bed/bay occupancy, operational waits, length-of-stay milestones and local capacity information.



\textbf{Overview}



Coordinate emergency-department bed/bay occupancy, operational waits, length-of-stay milestones and local capacity information. The reviewed hierarchy assigns this enduring ability to its stated business-enabling context.



\textbf{Intent}



Coordinate emergency-department bed/bay occupancy, operational waits, length-of-stay milestones and local capacity information.



\textbf{Boundary}



Information administration and coordination preserve the authority of the responsible source, clinical actor and governed domain; they confer no additional real-world management authority.


{\footnotesize \href{../../../../continuum/strategy/capability/8f0b0f3d9e8d44a7827b0835870b4a4f.md}{Canonical record: 8f0b0f3d9e8d44a7827b0835870b4a4f}; full qualifications and provenance.}\par
\subparagraph*{II.1.3.6.15.2.1 [Feature] ED Bed \& Bay Tracking}
\phantomsection
\label{navigation:II.1.3.6.15.2.1}
\addcontentsline{toc}{subparagraph}{II.1.3.6.15.2.1 [Feature] ED Bed \& Bay Tracking}

\label{element:7c0251128fc942d88bf25141421fd909}
\index{[Feature] ED Bed \& Bay Tracking (II.1.3.6.15.2.1)}
{\small Capability / Feature · Canonical · current}\par

Track occupancy and status of emergency cubicles, resuscitation bays, and short-stay beds.



Emergency-place tracking distinguishes occupancy and status across cubicles, resuscitation bays and short-stay beds. It provides place context for departmental operations while leaving the service's usable capacity and clinical triage decisions separate from an inventory of spaces.


{\footnotesize \href{../../../../continuum/strategy/capability/feature/7c0251128fc942d88bf25141421fd909.md}{Canonical record: 7c0251128fc942d88bf25141421fd909}; full qualifications and provenance.}\par
\subparagraph*{II.1.3.6.15.2.2 [Feature] ED Length of Stay Tracking}
\phantomsection
\label{navigation:II.1.3.6.15.2.2}
\addcontentsline{toc}{subparagraph}{II.1.3.6.15.2.2 [Feature] ED Length of Stay Tracking}

\label{element:e544f3d76db84dfaad57880d13a873c4}
\index{[Feature] ED Length of Stay Tracking (II.1.3.6.15.2.2)}
{\small Capability / Feature · Canonical · current}\par

Monitor operational wait times and length-of-stay milestones, generating operational escalation notifications.



Length-of-stay tracking makes waits and operational timing thresholds visible for escalation. The timing information supports management of impediments without determining clinical urgency or making an elapsed-time alert an independent assurance finding.


{\footnotesize \href{../../../../continuum/strategy/capability/feature/e544f3d76db84dfaad57880d13a873c4.md}{Canonical record: e544f3d76db84dfaad57880d13a873c4}; full qualifications and provenance.}\par
\subparagraph*{II.1.3.6.15.2.3 [Feature] Emergency Treatment Capacity Management}
\phantomsection
\label{navigation:II.1.3.6.15.2.3}
\addcontentsline{toc}{subparagraph}{II.1.3.6.15.2.3 [Feature] Emergency Treatment Capacity Management}

\label{element:502876430cf747f4b3bb102fb69557a5}
\index{[Feature] Emergency Treatment Capacity Management (II.1.3.6.15.2.3)}
{\small Capability / Feature · Proposed · current}\par

Represent current emergency treatment capacity under space, staffing and acuity constraints.



\textbf{Overview}



Bay state and length of stay contribute evidence but do not alone establish the ability to accept another care demand.



\textbf{Boundary}



Local responsibility retains its capacity meaning; Service Capacity Management consumes it without acquiring ownership or imposing one common capacity model.


{\footnotesize \href{../../../../continuum/strategy/capability/feature/502876430cf747f4b3bb102fb69557a5.md}{Canonical record: 502876430cf747f4b3bb102fb69557a5}; full qualifications and provenance.}\par
\subsubsection*{II.1.3.6.16 [CT2] Mobile Staff Management}
\phantomsection
\label{navigation:II.1.3.6.16}
\addcontentsline{toc}{subsubsection}{II.1.3.6.16 [CT2] Mobile Staff Management}

{\small Historical predecessor: R.2.144 \hyperref[element:32118856636f400a95fc28daacea61c8]{[CT2] Mobile Staff Management}; Canonical ID: \texttt{\seqsplit{32118856636f400a95fc28daacea61c8}}. Superseded historical identity — direct correspondence · superseded.}\par
\label{element:79a4a48fddf74c69acf7e16404d06ea5}
\index{[CT2] Mobile Staff Management (II.1.3.6.16)}
{\small Capability · Canonical · current}\par

Mobile Staff Management integrates with external business mobile-work mechanisms and receives visit progression and safety milestones. Its established dispatch Feature exchanges assignment, dispatch and status information with the relevant worker context. This is an information integration boundary around mobile activity, rather than a platform claim to decide the clinician's work.



\textbf{Overview}



Mobile Staff Management integrates with external business mobile-work mechanisms and receives visit progression and safety milestones. Its established dispatch Feature exchanges assignment, dispatch and status information with the relevant worker context. This is an information integration boundary around mobile activity, rather than a platform claim to decide the clinician's work.



\textbf{Intent}



Mobile Staff Management integrates with external business mobile-work mechanisms and receives visit progression and safety milestones. Its established dispatch Feature exchanges assignment, dispatch and status information with the relevant worker context. This is an information integration boundary around mobile activity, rather than a platform claim to decide the clinician's work.



\textbf{Boundary}



Information administration does not confer authority over the represented person, organisation, service, place, device or clinical activity.



\medskip\hrule\medskip

{\footnotesize \href{../../../../continuum/strategy/capability/79a4a48fddf74c69acf7e16404d06ea5.md}{Canonical record: 79a4a48fddf74c69acf7e16404d06ea5}; full qualifications and provenance.}\par
\paragraph*{II.1.3.6.16.1 [CT3] Mobile Worker Dispatch \& Safety Coordination}
\phantomsection
\label{navigation:II.1.3.6.16.1}
\addcontentsline{toc}{paragraph}{II.1.3.6.16.1 [CT3] Mobile Worker Dispatch \& Safety Coordination}

{\small Historical predecessor: R.2.143 \hyperref[element:5aafd1f4bff046b2b532c20dfb191473]{[CT3] Mobile Worker Dispatch \& Safety Coordination}; Canonical ID: \texttt{\seqsplit{5aafd1f4bff046b2b532c20dfb191473}}. Superseded historical identity — direct correspondence · superseded.}\par
\label{element:d72164668dfe448bbebe74e9e1a75233}
\index{[CT3] Mobile Worker Dispatch \& Safety Coordination (II.1.3.6.16.1)}
{\small Capability · Proposed · current}\par

Exchange externally established mobile work assignments, visit progression and safety check-in information with qualified worker context.



\textbf{Overview}



Exchange externally established mobile work assignments, visit progression and safety check-in information with qualified worker context. The reviewed hierarchy assigns this enduring ability to its stated business-enabling context.



\textbf{Intent}



Exchange externally established mobile work assignments, visit progression and safety check-in information with qualified worker context.



\textbf{Boundary}



Information administration and coordination preserve the authority of the responsible source, clinical actor and governed domain; they confer no additional real-world management authority.


{\footnotesize \href{../../../../continuum/strategy/capability/d72164668dfe448bbebe74e9e1a75233.md}{Canonical record: d72164668dfe448bbebe74e9e1a75233}; full qualifications and provenance.}\par
\subparagraph*{II.1.3.6.16.1.1 [Feature] Mobile Service Capacity Management}
\phantomsection
\label{navigation:II.1.3.6.16.1.1}
\addcontentsline{toc}{subparagraph}{II.1.3.6.16.1.1 [Feature] Mobile Service Capacity Management}

\label{element:abdcb95b3725405f878f199b07779d19}
\index{[Feature] Mobile Service Capacity Management (II.1.3.6.16.1.1)}
{\small Capability / Feature · Proposed · current}\par

Represent mobile service capacity with travel, coverage, work commitment and availability context.



\textbf{Overview}



Mobile availability cannot be interpreted as a simple on-duty count because distance, assignments and suitable skills affect deliverability.



\textbf{Boundary}



Local responsibility retains its capacity meaning; Service Capacity Management consumes it without acquiring ownership or imposing one common capacity model.


{\footnotesize \href{../../../../continuum/strategy/capability/feature/abdcb95b3725405f878f199b07779d19.md}{Canonical record: abdcb95b3725405f878f199b07779d19}; full qualifications and provenance.}\par
\subparagraph*{II.1.3.6.16.1.2 [Feature] Mobile Visit Status Tracking}
\phantomsection
\label{navigation:II.1.3.6.16.1.2}
\addcontentsline{toc}{subparagraph}{II.1.3.6.16.1.2 [Feature] Mobile Visit Status Tracking}

\label{element:3cd49cb00159437ea25048c1158172c1}
\index{[Feature] Mobile Visit Status Tracking (II.1.3.6.16.1.2)}
{\small Capability / Feature · Canonical · current}\par

Ingest arrival, visit progression, and safety check-in milestones from roving staff.



Mobile visit tracking receives arrival, progression and safety checkpoints from roving workers. It makes reported activity visible in the wider operational context without originating the clinical visit, deciding its clinical content or assuming that a checkpoint proves the intended clinical outcome.


{\footnotesize \href{../../../../continuum/strategy/capability/feature/3cd49cb00159437ea25048c1158172c1.md}{Canonical record: 3cd49cb00159437ea25048c1158172c1}; full qualifications and provenance.}\par
\subparagraph*{II.1.3.6.16.1.3 [Feature] Mobile Worker Task Dispatch}
\phantomsection
\label{navigation:II.1.3.6.16.1.3}
\addcontentsline{toc}{subparagraph}{II.1.3.6.16.1.3 [Feature] Mobile Worker Task Dispatch}

\label{element:426363f2f02a4625a4feab89a2adc755}
\index{[Feature] Mobile Worker Task Dispatch (II.1.3.6.16.1.3)}
{\small Capability / Feature · Canonical · current}\par

Dispatch community care tasks and urgent visits to roving clinicians' mobile devices.



Mobile dispatch exchanges externally established assignment, dispatch and progression information with business mobile-work mechanisms. The retained name does not grant Harmonia authority to allocate clinical work, manage handover or own EMR worklists, and external task terminology does not redefine Synthetic Task.



\textbf{Boundary}



The proposed replacement name remains unconfirmed. This entry preserves the existing identity and the human-adjudicated integration boundary.


{\footnotesize \href{../../../../continuum/strategy/capability/feature/426363f2f02a4625a4feab89a2adc755.md}{Canonical record: 426363f2f02a4625a4feab89a2adc755}; full qualifications and provenance.}\par
\subsubsection*{II.1.3.6.17 [CT2] On-Call Management}
\phantomsection
\label{navigation:II.1.3.6.17}
\addcontentsline{toc}{subsubsection}{II.1.3.6.17 [CT2] On-Call Management}

{\small Historical predecessor: R.2.146 \hyperref[element:c8e6e66e8718493c9093f37d129f9ff2]{[CT2] On-Call Management}; Canonical ID: \texttt{\seqsplit{c8e6e66e8718493c9093f37d129f9ff2}}. Superseded historical identity — direct correspondence · superseded.}\par
\label{element:b16cef8f2bfd4b92948ec75cbcfa7b9a}
\index{[CT2] On-Call Management (II.1.3.6.17)}
{\small Capability · Canonical · current}\par

On-Call Management resolves the active provider for a specialty, service or hospital at the relevant time. It turns changing coverage information into usable routing context so that operational communication can reach the intended capacity. This differs from authoring the roster or deciding who is professionally qualified for the requested activity.



\textbf{Overview}



On-Call Management resolves the active provider for a specialty, service or hospital at the relevant time. It turns changing coverage information into usable routing context so that operational communication can reach the intended capacity. This differs from authoring the roster or deciding who is professionally qualified for the requested activity.



\textbf{Intent}



On-Call Management resolves the active provider for a specialty, service or hospital at the relevant time. It turns changing coverage information into usable routing context so that operational communication can reach the intended capacity. This differs from authoring the roster or deciding who is professionally qualified for the requested activity.



\textbf{Boundary}



Information administration does not confer authority over the represented person, organisation, service, place, device or clinical activity.



\medskip\hrule\medskip

{\footnotesize \href{../../../../continuum/strategy/capability/b16cef8f2bfd4b92948ec75cbcfa7b9a.md}{Canonical record: b16cef8f2bfd4b92948ec75cbcfa7b9a}; full qualifications and provenance.}\par
\paragraph*{II.1.3.6.17.1 [CT3] On-Call Availability \& Contact Resolution}
\phantomsection
\label{navigation:II.1.3.6.17.1}
\addcontentsline{toc}{paragraph}{II.1.3.6.17.1 [CT3] On-Call Availability \& Contact Resolution}

{\small Historical predecessor: R.2.145 \hyperref[element:a03075a46f75497f8c6ba95d9c9e27c4]{[CT3] On-Call Availability \& Contact Resolution}; Canonical ID: \texttt{\seqsplit{a03075a46f75497f8c6ba95d9c9e27c4}}. Superseded historical identity — direct correspondence · superseded.}\par
\label{element:677a816c42414ea694791b4ee5a14285}
\index{[CT3] On-Call Availability \& Contact Resolution (II.1.3.6.17.1)}
{\small Capability · Proposed · current}\par

Resolve currently available on-call providers and their contact paths for the applicable specialty, service and time.



\textbf{Overview}



Resolve currently available on-call providers and their contact paths for the applicable specialty, service and time. The reviewed hierarchy assigns this enduring ability to its stated business-enabling context.



\textbf{Intent}



Resolve currently available on-call providers and their contact paths for the applicable specialty, service and time.



\textbf{Boundary}



Information administration and coordination preserve the authority of the responsible source, clinical actor and governed domain; they confer no additional real-world management authority.


{\footnotesize \href{../../../../continuum/strategy/capability/677a816c42414ea694791b4ee5a14285.md}{Canonical record: 677a816c42414ea694791b4ee5a14285}; full qualifications and provenance.}\par
\subparagraph*{II.1.3.6.17.1.1 [Feature] Active On-Call Provider Resolution}
\phantomsection
\label{navigation:II.1.3.6.17.1.1}
\addcontentsline{toc}{subparagraph}{II.1.3.6.17.1.1 [Feature] Active On-Call Provider Resolution}

\label{element:b588f3cc869d460a9561657352fb189a}
\index{[Feature] Active On-Call Provider Resolution (II.1.3.6.17.1.1)}
{\small Capability / Feature · Canonical · current}\par

Resolve currently active on-call specialists, registrars, and operational managers for a given specialty and hospital at any point in time.



Active on-call resolution identifies the current provider in the required specialty and hospital context. It supports correct routing at the relevant time without authoring the roster or treating current coverage as proof of qualification, privilege or responsibility acceptance.


{\footnotesize \href{../../../../continuum/strategy/capability/feature/b588f3cc869d460a9561657352fb189a.md}{Canonical record: b588f3cc869d460a9561657352fb189a}; full qualifications and provenance.}\par
\subsubsection*{II.1.3.6.18 [CT2] Operational Workforce Coordination [superseded]}
\phantomsection
\label{navigation:II.1.3.6.18}
\addcontentsline{toc}{subsubsection}{II.1.3.6.18 [CT2] Operational Workforce Coordination [superseded]}

\label{element:66e72222f2064a23adbdb0311092918b}
\index{[CT2] Operational Workforce Coordination (II.1.3.6.18)}
{\small Capability · Proposed · superseded}\par
\href{../../../../continuum/strategy/capability/66e72222f2064a23adbdb0311092918b.md}{Historical canonical definition and provenance}.\par
\subsubsection*{II.1.3.6.19 [CT2] Outpatient Operations}
\phantomsection
\label{navigation:II.1.3.6.19}
\addcontentsline{toc}{subsubsection}{II.1.3.6.19 [CT2] Outpatient Operations}

{\small Historical predecessor: R.2.148 \hyperref[element:3db41bc06ade4f9088eec7b8aedd7448]{[CT2] Outpatient Operations}; Canonical ID: \texttt{\seqsplit{3db41bc06ade4f9088eec7b8aedd7448}}. Superseded historical identity — direct correspondence · superseded.}\par
\label{element:21399d52b9624ddfbeb786f1d72e40c3}
\index{[CT2] Outpatient Operations (II.1.3.6.19)}
{\small Capability · Canonical · current}\par

Outpatient Operations integrates arrival events and follow-up requirements generated as clinics progress. It connects a person's operational presence with the relevant service context and makes post-clinic booking or diagnostic needs available for subsequent activity. The responsibility retains outpatient-specific capacity semantics without requiring a separate arrival Function where approved reusable composition suffices.



\textbf{Overview}



Outpatient Operations integrates arrival events and follow-up requirements generated as clinics progress. It connects a person's operational presence with the relevant service context and makes post-clinic booking or diagnostic needs available for subsequent activity. The responsibility retains outpatient-specific capacity semantics without requiring a separate arrival Function where approved reusable composition suffices.



\textbf{Intent}



Outpatient Operations integrates arrival events and follow-up requirements generated as clinics progress. It connects a person's operational presence with the relevant service context and makes post-clinic booking or diagnostic needs available for subsequent activity. The responsibility retains outpatient-specific capacity semantics without requiring a separate arrival Function where approved reusable composition suffices.



\textbf{Boundary}



Information administration does not confer authority over the represented person, organisation, service, place, device or clinical activity.



\medskip\hrule\medskip

{\footnotesize \href{../../../../continuum/strategy/capability/21399d52b9624ddfbeb786f1d72e40c3.md}{Canonical record: 21399d52b9624ddfbeb786f1d72e40c3}; full qualifications and provenance.}\par
\paragraph*{II.1.3.6.19.1 [CT3] Outpatient Flow Coordination}
\phantomsection
\label{navigation:II.1.3.6.19.1}
\addcontentsline{toc}{paragraph}{II.1.3.6.19.1 [CT3] Outpatient Flow Coordination}

{\small Historical predecessor: R.2.147 \hyperref[element:e6c59c7b72fe44169a5fdfd08a5e91c4]{[CT3] Outpatient Flow Coordination}; Canonical ID: \texttt{\seqsplit{e6c59c7b72fe44169a5fdfd08a5e91c4}}. Superseded historical identity — direct correspondence · superseded.}\par
\label{element:799ea278af014168b9b534e5fdadc6d1}
\index{[CT3] Outpatient Flow Coordination (II.1.3.6.19.1)}
{\small Capability · Proposed · current}\par

Coordinate outpatient arrival, clinic flow, follow-up requirements and context-specific outpatient capacity information.



\textbf{Overview}



Coordinate outpatient arrival, clinic flow, follow-up requirements and context-specific outpatient capacity information. The reviewed hierarchy assigns this enduring ability to its stated business-enabling context.



\textbf{Intent}



Coordinate outpatient arrival, clinic flow, follow-up requirements and context-specific outpatient capacity information.



\textbf{Boundary}



Information administration and coordination preserve the authority of the responsible source, clinical actor and governed domain; they confer no additional real-world management authority.


{\footnotesize \href{../../../../continuum/strategy/capability/799ea278af014168b9b534e5fdadc6d1.md}{Canonical record: 799ea278af014168b9b534e5fdadc6d1}; full qualifications and provenance.}\par
\subparagraph*{II.1.3.6.19.1.1 [Feature] Outpatient Appointment Capacity Management}
\phantomsection
\label{navigation:II.1.3.6.19.1.1}
\addcontentsline{toc}{subparagraph}{II.1.3.6.19.1.1 [Feature] Outpatient Appointment Capacity Management}

\label{element:3c4f8b0fc23044cbbac32602e1272680}
\index{[Feature] Outpatient Appointment Capacity Management (II.1.3.6.19.1.1)}
{\small Capability / Feature · Proposed · current}\par

Represent usable outpatient session capacity from booking, staffing and local flow information.



\textbf{Overview}



The projection distinguishes scheduled slots from currently usable care capacity without taking source booking authority.



\textbf{Boundary}



Local responsibility retains its capacity meaning; Service Capacity Management consumes it without acquiring ownership or imposing one common capacity model.


{\footnotesize \href{../../../../continuum/strategy/capability/feature/3c4f8b0fc23044cbbac32602e1272680.md}{Canonical record: 3c4f8b0fc23044cbbac32602e1272680}; full qualifications and provenance.}\par
\subparagraph*{II.1.3.6.19.1.2 [Feature] Outpatient Arrival Registration}
\phantomsection
\label{navigation:II.1.3.6.19.1.2}
\addcontentsline{toc}{subparagraph}{II.1.3.6.19.1.2 [Feature] Outpatient Arrival Registration}

\label{element:4ea71fb4950848cfb4f6d916190edad0}
\index{[Feature] Outpatient Arrival Registration (II.1.3.6.19.1.2)}
{\small Capability / Feature · Canonical · current}\par

Ingest self-service check-in and clinic arrival events, notifying treating clinicians.



Outpatient arrival registration connects a check-in event with the treating context and required notification. Approved reusable composition supplies this integration responsibility; it does not require a specialised arrival Function or prove that a consultation has begun or finished.


{\footnotesize \href{../../../../continuum/strategy/capability/feature/4ea71fb4950848cfb4f6d916190edad0.md}{Canonical record: 4ea71fb4950848cfb4f6d916190edad0}; full qualifications and provenance.}\par
\subparagraph*{II.1.3.6.19.1.3 [Feature] Post-Clinic Order Coordination}
\phantomsection
\label{navigation:II.1.3.6.19.1.3}
\addcontentsline{toc}{subparagraph}{II.1.3.6.19.1.3 [Feature] Post-Clinic Order Coordination}

\label{element:439b3530b7c9443ab3e48d2cc50a3b0c}
\index{[Feature] Post-Clinic Order Coordination (II.1.3.6.19.1.3)}
{\small Capability / Feature · Canonical · current}\par

Capture follow-up booking and diagnostic order requirements generated at clinic completion.



Post-clinic coordination captures the follow-up bookings and diagnostic requirements established at clinic completion. It connects those requirements with subsequent operational activity while preserving that capture is not the clinical decision, a completed booking or diagnostic fulfilment.


{\footnotesize \href{../../../../continuum/strategy/capability/feature/439b3530b7c9443ab3e48d2cc50a3b0c.md}{Canonical record: 439b3530b7c9443ab3e48d2cc50a3b0c}; full qualifications and provenance.}\par
\subsubsection*{II.1.3.6.20 [CT2] Patient Transport}
\phantomsection
\label{navigation:II.1.3.6.20}
\addcontentsline{toc}{subsubsection}{II.1.3.6.20 [CT2] Patient Transport}

{\small Historical predecessor: R.2.150 \hyperref[element:e9c1a897884f4d758f5654b824364030]{[CT2] Patient Transport}; Canonical ID: \texttt{\seqsplit{e9c1a897884f4d758f5654b824364030}}. Superseded historical identity — direct correspondence · superseded.}\par
\label{element:6041d7c0cc5d4519aaea584bb4e6d365}
\index{[CT2] Patient Transport (II.1.3.6.20)}
{\small Capability · Canonical · current}\par

Patient Transport coordinates the operational progression of movement requests, including the patient's applicable mobility constraints. Departure, transit and arrival checkpoints make the physical journey visible to participants. The responsibility supports internal and inter-facility movement without treating arrival as acceptance of clinical responsibility at the destination.



\textbf{Overview}



Patient Transport coordinates the operational progression of movement requests, including the patient's applicable mobility constraints. Departure, transit and arrival checkpoints make the physical journey visible to participants. The responsibility supports internal and inter-facility movement without treating arrival as acceptance of clinical responsibility at the destination.



\textbf{Intent}



Patient Transport coordinates the operational progression of movement requests, including the patient's applicable mobility constraints. Departure, transit and arrival checkpoints make the physical journey visible to participants. The responsibility supports internal and inter-facility movement without treating arrival as acceptance of clinical responsibility at the destination.



\textbf{Boundary}



Information administration does not confer authority over the represented person, organisation, service, place, device or clinical activity.



\medskip\hrule\medskip

{\footnotesize \href{../../../../continuum/strategy/capability/6041d7c0cc5d4519aaea584bb4e6d365.md}{Canonical record: 6041d7c0cc5d4519aaea584bb4e6d365}; full qualifications and provenance.}\par
\paragraph*{II.1.3.6.20.1 [CT3] Patient Transfer Coordination}
\phantomsection
\label{navigation:II.1.3.6.20.1}
\addcontentsline{toc}{paragraph}{II.1.3.6.20.1 [CT3] Patient Transfer Coordination}

{\small Historical predecessor: R.2.149 \hyperref[element:b5a76d4877424b2ea9b01fb5fa22772c]{[CT3] Patient Transfer Coordination}; Canonical ID: \texttt{\seqsplit{b5a76d4877424b2ea9b01fb5fa22772c}}. Superseded historical identity — direct correspondence · superseded.}\par
\label{element:c7c7aa8e34df4568ae4b91cdde380a11}
\index{[CT3] Patient Transfer Coordination (II.1.3.6.20.1)}
{\small Capability · Proposed · current}\par

Coordinate patient transport requests, dispatch and observed departure, transit and arrival milestones.



\textbf{Overview}



Coordinate patient transport requests, dispatch and observed departure, transit and arrival milestones. The reviewed hierarchy assigns this enduring ability to its stated business-enabling context.



\textbf{Intent}



Coordinate patient transport requests, dispatch and observed departure, transit and arrival milestones.



\textbf{Boundary}



Information administration and coordination preserve the authority of the responsible source, clinical actor and governed domain; they confer no additional real-world management authority.


{\footnotesize \href{../../../../continuum/strategy/capability/c7c7aa8e34df4568ae4b91cdde380a11.md}{Canonical record: c7c7aa8e34df4568ae4b91cdde380a11}; full qualifications and provenance.}\par
\subparagraph*{II.1.3.6.20.1.1 [Feature] Transport Dispatch \& Progress Tracking}
\phantomsection
\label{navigation:II.1.3.6.20.1.1}
\addcontentsline{toc}{subparagraph}{II.1.3.6.20.1.1 [Feature] Transport Dispatch \& Progress Tracking}

\label{element:d8ef9994298f466bafce919c5e6cb701}
\index{[Feature] Transport Dispatch \& Progress Tracking (II.1.3.6.20.1.1)}
{\small Capability / Feature · Canonical · current}\par

Dispatch transport orders, tracking departure, transit, and destination arrival checkpoints.



Transport progression connects dispatch, departure, transit and arrival evidence for a physical patient journey. The checkpoints make movement intelligible while preserving that reaching a destination does not establish clinical handover or assumption of responsibility there.


{\footnotesize \href{../../../../continuum/strategy/capability/feature/d8ef9994298f466bafce919c5e6cb701.md}{Canonical record: d8ef9994298f466bafce919c5e6cb701}; full qualifications and provenance.}\par
\subparagraph*{II.1.3.6.20.1.2 [Feature] Transport Request Ingestion}
\phantomsection
\label{navigation:II.1.3.6.20.1.2}
\addcontentsline{toc}{subparagraph}{II.1.3.6.20.1.2 [Feature] Transport Request Ingestion}

\label{element:7104472d065840eb91ae4cb7a3a0cb43}
\index{[Feature] Transport Request Ingestion (II.1.3.6.20.1.2)}
{\small Capability / Feature · Canonical · current}\par

Ingest internal (portering) and external (ambulance, inter-hospital) patient transfer requests with clinical mobility requirements.



Transport ingestion establishes the requested movement and applicable clinical mobility constraints. It gives operational coordination the context needed to arrange transit without making the request itself a departure, arrival or accepted care transition.


{\footnotesize \href{../../../../continuum/strategy/capability/feature/7104472d065840eb91ae4cb7a3a0cb43.md}{Canonical record: 7104472d065840eb91ae4cb7a3a0cb43}; full qualifications and provenance.}\par
\subsubsection*{II.1.3.6.21 [CT2] Service Capacity Management}
\phantomsection
\label{navigation:II.1.3.6.21}
\addcontentsline{toc}{subsubsection}{II.1.3.6.21 [CT2] Service Capacity Management}

{\small Historical predecessor: R.2.152 \hyperref[element:ac205f374c9d41c3b6d18a6b7c356152]{[CT2] Service Capacity Management}; Canonical ID: \texttt{\seqsplit{ac205f374c9d41c3b6d18a6b7c356152}}. Superseded historical identity — direct correspondence · superseded.}\par
\label{element:d86105bf9c86445394088f65c895d37f}
\index{[CT2] Service Capacity Management (II.1.3.6.21)}
{\small Capability · Canonical · current}\par

Service Capacity Management provides a higher-order view by aggregating capacity information across services and facilities. It can consider occupancy, waits and diversion while preserving the service-specific meaning of the underlying capacity. Its purpose is not to erase differences between clinic, ward, theatre, emergency, outpatient and mobile-service operation.



\textbf{Overview}



Service Capacity Management provides a higher-order view by aggregating capacity information across services and facilities. It can consider occupancy, waits and diversion while preserving the service-specific meaning of the underlying capacity. Its purpose is not to erase differences between clinic, ward, theatre, emergency, outpatient and mobile-service operation.



\textbf{Intent}



Service Capacity Management provides a higher-order view by aggregating capacity information across services and facilities. It can consider occupancy, waits and diversion while preserving the service-specific meaning of the underlying capacity. Its purpose is not to erase differences between clinic, ward, theatre, emergency, outpatient and mobile-service operation.



\textbf{Boundary}



Information administration does not confer authority over the represented person, organisation, service, place, device or clinical activity.



\medskip\hrule\medskip

{\footnotesize \href{../../../../continuum/strategy/capability/d86105bf9c86445394088f65c895d37f.md}{Canonical record: d86105bf9c86445394088f65c895d37f}; full qualifications and provenance.}\par
\paragraph*{II.1.3.6.21.1 [CT3] Service Demand \& Capacity Coordination}
\phantomsection
\label{navigation:II.1.3.6.21.1}
\addcontentsline{toc}{paragraph}{II.1.3.6.21.1 [CT3] Service Demand \& Capacity Coordination}

{\small Historical predecessor: R.2.151 \hyperref[element:b57081e650f74ad488549fdaa5305bcb]{[CT3] Service Demand \& Capacity Coordination}; Canonical ID: \texttt{\seqsplit{b57081e650f74ad488549fdaa5305bcb}}. Superseded historical identity — direct correspondence · superseded.}\par
\label{element:93b5548ce8d04fcd999b97536c1eebf9}
\index{[CT3] Service Demand \& Capacity Coordination (II.1.3.6.21.1)}
{\small Capability · Proposed · current}\par

Coordinate demand and capacity information across services while retaining the originating responsibilities of context-specific capacity providers.



\textbf{Overview}



Coordinate demand and capacity information across services while retaining the originating responsibilities of context-specific capacity providers. The reviewed hierarchy assigns this enduring ability to its stated business-enabling context.



\textbf{Intent}



Coordinate demand and capacity information across services while retaining the originating responsibilities of context-specific capacity providers.



\textbf{Boundary}



Information administration and coordination preserve the authority of the responsible source, clinical actor and governed domain; they confer no additional real-world management authority.


{\footnotesize \href{../../../../continuum/strategy/capability/93b5548ce8d04fcd999b97536c1eebf9.md}{Canonical record: 93b5548ce8d04fcd999b97536c1eebf9}; full qualifications and provenance.}\par
\subparagraph*{II.1.3.6.21.1.1 [Feature] Operational Capacity Metric Aggregation}
\phantomsection
\label{navigation:II.1.3.6.21.1.1}
\addcontentsline{toc}{subparagraph}{II.1.3.6.21.1.1 [Feature] Operational Capacity Metric Aggregation}

\label{element:1fe9199a49744eeaa5d063d7a84544ea}
\index{[Feature] Operational Capacity Metric Aggregation (II.1.3.6.21.1.1)}
{\small Capability / Feature · Canonical · current}\par

Aggregate operational capacity telemetry (occupancy rates, wait times, diversion status) across facilities.



Capacity aggregation combines occupancy, waiting and diversion context across facilities. The higher-order view must retain the originating service-specific capacity meaning; a shared metric does not transfer clinic, ward or theatre capacity responsibility to the aggregator.


{\footnotesize \href{../../../../continuum/strategy/capability/feature/1fe9199a49744eeaa5d063d7a84544ea.md}{Canonical record: 1fe9199a49744eeaa5d063d7a84544ea}; full qualifications and provenance.}\par
\subsubsection*{II.1.3.6.22 [CT2] Service Establishment Support}
\phantomsection
\label{navigation:II.1.3.6.22}
\addcontentsline{toc}{subsubsection}{II.1.3.6.22 [CT2] Service Establishment Support}

{\small Historical predecessor: R.2.153 \hyperref[element:8a63b08f34fa42c5bda2a3e12b9ecd55]{[CT2] Service Establishment Support}; Canonical ID: \texttt{\seqsplit{8a63b08f34fa42c5bda2a3e12b9ecd55}}. Superseded historical identity — direct correspondence · superseded.}\par
\label{element:20e67a41e7f44398b0f0cf3d379d8b9c}
\index{[CT2] Service Establishment Support (II.1.3.6.22)}
{\small Capability · Proposed · current}\par

Support the information needed to establish a health service in its intended operating context.



\textbf{Overview}



Support the information needed to establish a health service in its intended operating context. The reviewed hierarchy assigns this enduring ability to its stated business-enabling context.



\textbf{Intent}



Support the information needed to establish a health service in its intended operating context.



\textbf{Boundary}



Information administration and coordination preserve the authority of the responsible source, clinical actor and governed domain; they confer no additional real-world management authority.


{\footnotesize \href{../../../../continuum/strategy/capability/20e67a41e7f44398b0f0cf3d379d8b9c.md}{Canonical record: 20e67a41e7f44398b0f0cf3d379d8b9c}; full qualifications and provenance.}\par
\subsubsection*{II.1.3.6.23 [CT2] Theatre Operations}
\phantomsection
\label{navigation:II.1.3.6.23}
\addcontentsline{toc}{subsubsection}{II.1.3.6.23 [CT2] Theatre Operations}

{\small Historical predecessor: R.2.155 \hyperref[element:9ac72fea821c45e2b18632ae9a9867dd]{[CT2] Theatre Operations}; Canonical ID: \texttt{\seqsplit{9ac72fea821c45e2b18632ae9a9867dd}}. Superseded historical identity — direct correspondence · superseded.}\par
\label{element:4bf7a8260e1444fdac44e97c3557b533}
\index{[CT2] Theatre Operations (II.1.3.6.23)}
{\small Capability · Canonical · current}\par

Theatre Operations integrates surgical case progression and the operational consequences relevant to portering, sterilisation and recovery participants. It gives reported milestones an operational coordination context without originating the surgical acts they describe. Theatre-specific capacity remains a business responsibility even when shared information machinery supports several service contexts.



\textbf{Overview}



Theatre Operations integrates surgical case progression and the operational consequences relevant to portering, sterilisation and recovery participants. It gives reported milestones an operational coordination context without originating the surgical acts they describe. Theatre-specific capacity remains a business responsibility even when shared information machinery supports several service contexts.



\textbf{Intent}



Theatre Operations integrates surgical case progression and the operational consequences relevant to portering, sterilisation and recovery participants. It gives reported milestones an operational coordination context without originating the surgical acts they describe. Theatre-specific capacity remains a business responsibility even when shared information machinery supports several service contexts.



\textbf{Boundary}



Information administration does not confer authority over the represented person, organisation, service, place, device or clinical activity.



\medskip\hrule\medskip

{\footnotesize \href{../../../../continuum/strategy/capability/4bf7a8260e1444fdac44e97c3557b533.md}{Canonical record: 4bf7a8260e1444fdac44e97c3557b533}; full qualifications and provenance.}\par
\paragraph*{II.1.3.6.23.1 [CT3] Perioperative Logistics Coordination}
\phantomsection
\label{navigation:II.1.3.6.23.1}
\addcontentsline{toc}{paragraph}{II.1.3.6.23.1 [CT3] Perioperative Logistics Coordination}

{\small Historical predecessor: R.2.154 \hyperref[element:e7813a06b9df462fa135c299261f192e]{[CT3] Perioperative Logistics Coordination}; Canonical ID: \texttt{\seqsplit{e7813a06b9df462fa135c299261f192e}}. Superseded historical identity — direct correspondence · superseded.}\par
\label{element:6b7c08b60b5b48519bfe4b6f993dd6fb}
\index{[CT3] Perioperative Logistics Coordination (II.1.3.6.23.1)}
{\small Capability · Proposed · current}\par

Coordinate theatre case progression information and notifications to supporting perioperative operational teams.



\textbf{Overview}



Coordinate theatre case progression information and notifications to supporting perioperative operational teams. The reviewed hierarchy assigns this enduring ability to its stated business-enabling context.



\textbf{Intent}



Coordinate theatre case progression information and notifications to supporting perioperative operational teams.



\textbf{Boundary}



Information administration and coordination preserve the authority of the responsible source, clinical actor and governed domain; they confer no additional real-world management authority.


{\footnotesize \href{../../../../continuum/strategy/capability/6b7c08b60b5b48519bfe4b6f993dd6fb.md}{Canonical record: 6b7c08b60b5b48519bfe4b6f993dd6fb}; full qualifications and provenance.}\par
\subparagraph*{II.1.3.6.23.1.1 [Feature] Operating Theatre Case Progression}
\phantomsection
\label{navigation:II.1.3.6.23.1.1}
\addcontentsline{toc}{subparagraph}{II.1.3.6.23.1.1 [Feature] Operating Theatre Case Progression}

\label{element:2fd65f09f7384f0486280e56ec5a9e84}
\index{[Feature] Operating Theatre Case Progression (II.1.3.6.23.1.1)}
{\small Capability / Feature · Canonical · current}\par

Track surgical case state transitions (CALLED → IN\_THEATRE → ANAESTHETISED → INCISION → CLOSED → RECOVERY).



Case progression records externally established surgical milestones in an operationally meaningful sequence. It makes the theatre situation visible without authorising or performing the clinical acts denoted by anaesthesia, incision, closure or recovery.


{\footnotesize \href{../../../../continuum/strategy/capability/feature/2fd65f09f7384f0486280e56ec5a9e84.md}{Canonical record: 2fd65f09f7384f0486280e56ec5a9e84}; full qualifications and provenance.}\par
\subparagraph*{II.1.3.6.23.1.2 [Feature] Perioperative Resource Notification}
\phantomsection
\label{navigation:II.1.3.6.23.1.2}
\addcontentsline{toc}{subparagraph}{II.1.3.6.23.1.2 [Feature] Perioperative Resource Notification}

\label{element:1f2446ccc59348a6bc54d8268ef4e6e1}
\index{[Feature] Perioperative Resource Notification (II.1.3.6.23.1.2)}
{\small Capability / Feature · Canonical · current}\par

Dispatch alerts to portering, sterilization, and recovery teams based on surgical case progression milestones.



Perioperative notification connects relevant case milestones with portering, sterilisation and recovery participants. Observation alone does not realise the Feature's dispatch responsibility; the established case-progression behaviour includes the required alerts without adding another Function or Service.


{\footnotesize \href{../../../../continuum/strategy/capability/feature/1f2446ccc59348a6bc54d8268ef4e6e1.md}{Canonical record: 1f2446ccc59348a6bc54d8268ef4e6e1}; full qualifications and provenance.}\par
\subparagraph*{II.1.3.6.23.1.3 [Feature] Theatre Session Capacity Management}
\phantomsection
\label{navigation:II.1.3.6.23.1.3}
\addcontentsline{toc}{subparagraph}{II.1.3.6.23.1.3 [Feature] Theatre Session Capacity Management}

\label{element:258c27077caf4bf88e562e8dd607bdfd}
\index{[Feature] Theatre Session Capacity Management (II.1.3.6.23.1.3)}
{\small Capability / Feature · Proposed · current}\par

Represent available and committed procedural-session capacity with its resource dependencies.



\textbf{Overview}



Theatre capacity involves session time, team capability, equipment and recovery provision rather than a count of theatres.



\textbf{Boundary}



Local responsibility retains its capacity meaning; Service Capacity Management consumes it without acquiring ownership or imposing one common capacity model.


{\footnotesize \href{../../../../continuum/strategy/capability/feature/258c27077caf4bf88e562e8dd607bdfd.md}{Canonical record: 258c27077caf4bf88e562e8dd607bdfd}; full qualifications and provenance.}\par
\subsubsection*{II.1.3.6.24 [CT2] Ward Operations}
\phantomsection
\label{navigation:II.1.3.6.24}
\addcontentsline{toc}{subsubsection}{II.1.3.6.24 [CT2] Ward Operations}

{\small Historical predecessor: R.2.157 \hyperref[element:bf3d2e3581f24beda20a6a3a3348460e]{[CT2] Ward Operations}; Canonical ID: \texttt{\seqsplit{bf3d2e3581f24beda20a6a3a3348460e}}. Superseded historical identity — direct correspondence · superseded.}\par
\label{element:f996ecca728f4942a88b034aca818a99}
\index{[CT2] Ward Operations (II.1.3.6.24)}
{\small Capability · Canonical · current}\par

Ward Operations coordinates ward occupancy and care-place context and assembles the information needed around shift changes. Its responsibility includes the ward-specific meaning of capacity and patient flow rather than a generic resource counter. It can make pending concerns and tasks visible without becoming the authority managing clinical handover.



\textbf{Overview}



Ward Operations coordinates ward occupancy and care-place context and assembles the information needed around shift changes. Its responsibility includes the ward-specific meaning of capacity and patient flow rather than a generic resource counter. It can make pending concerns and tasks visible without becoming the authority managing clinical handover.



\textbf{Intent}



Ward Operations coordinates ward occupancy and care-place context and assembles the information needed around shift changes. Its responsibility includes the ward-specific meaning of capacity and patient flow rather than a generic resource counter. It can make pending concerns and tasks visible without becoming the authority managing clinical handover.



\textbf{Boundary}



Information administration does not confer authority over the represented person, organisation, service, place, device or clinical activity.



\medskip\hrule\medskip

{\footnotesize \href{../../../../continuum/strategy/capability/f996ecca728f4942a88b034aca818a99.md}{Canonical record: f996ecca728f4942a88b034aca818a99}; full qualifications and provenance.}\par
\paragraph*{II.1.3.6.24.1 [CT3] Multidisciplinary Inpatient Care Coordination}
\phantomsection
\label{navigation:II.1.3.6.24.1}
\addcontentsline{toc}{paragraph}{II.1.3.6.24.1 [CT3] Multidisciplinary Inpatient Care Coordination}

\label{element:c798e38e12e5428dbb8b3cabc977d430}
\index{[CT3] Multidisciplinary Inpatient Care Coordination (II.1.3.6.24.1)}
{\small Capability · Canonical · current}\par

Definition unresolved in the baseline. No definition has been synthesised from the name.



\medskip\hrule\medskip

{\footnotesize \href{../../../../continuum/strategy/capability/c798e38e12e5428dbb8b3cabc977d430.md}{Canonical record: c798e38e12e5428dbb8b3cabc977d430}; full qualifications and provenance.}\par
\subparagraph*{II.1.3.6.24.1.1 [Feature] Inpatient Clinical Progression Tracking}
\phantomsection
\label{navigation:II.1.3.6.24.1.1}
\addcontentsline{toc}{subparagraph}{II.1.3.6.24.1.1 [Feature] Inpatient Clinical Progression Tracking}

\label{element:863265027ecb432595b9c6d45adf1190}
\index{[Feature] Inpatient Clinical Progression Tracking (II.1.3.6.24.1.1)}
{\small Capability / Feature · Canonical · current}\par

Capture inpatient clinical milestone completions (e.g., medical stability achieved, allied health clearance).



Inpatient progression tracking captures clinical milestones reported by the responsible participants. General clinical-information capabilities and configured workflow composition suffice for the established Feature; capture does not require inventing a specialised Function or make Harmonia the authority declaring clinical stability.


{\footnotesize \href{../../../../continuum/strategy/capability/feature/863265027ecb432595b9c6d45adf1190.md}{Canonical record: 863265027ecb432595b9c6d45adf1190}; full qualifications and provenance.}\par
\subparagraph*{II.1.3.6.24.1.2 [Feature] Inpatient Rounding Context Provision}
\phantomsection
\label{navigation:II.1.3.6.24.1.2}
\addcontentsline{toc}{subparagraph}{II.1.3.6.24.1.2 [Feature] Inpatient Rounding Context Provision}

\label{element:c0dfa2ef22354340a4f1488790359bbd}
\index{[Feature] Inpatient Rounding Context Provision (II.1.3.6.24.1.2)}
{\small Capability / Feature · Canonical · current}\par

Aggregate diagnostic results, current medication administration records, and nursing notes for multidisciplinary inpatient rounds.



Rounding context brings results, medication records and nursing notes into an accessible multidisciplinary view. Assembly serves the clinical discussion but does not perform rounds, determine pending clinical work or establish the team's professional conclusions.


{\footnotesize \href{../../../../continuum/strategy/capability/feature/c0dfa2ef22354340a4f1488790359bbd.md}{Canonical record: c0dfa2ef22354340a4f1488790359bbd}; full qualifications and provenance.}\par
\paragraph*{II.1.3.6.24.2 [CT3] Ward Logistics Coordination}
\phantomsection
\label{navigation:II.1.3.6.24.2}
\addcontentsline{toc}{paragraph}{II.1.3.6.24.2 [CT3] Ward Logistics Coordination}

{\small Historical predecessor: R.2.156 \hyperref[element:8a9dcc3b5f0143458ec71bc188a56020]{[CT3] Ward Logistics Coordination}; Canonical ID: \texttt{\seqsplit{8a9dcc3b5f0143458ec71bc188a56020}}. Superseded historical identity — direct correspondence · superseded.}\par
\label{element:3249647f7b9846c899f982ab339c81f9}
\index{[CT3] Ward Logistics Coordination (II.1.3.6.24.2)}
{\small Capability · Proposed · current}\par

Coordinate ward patient flow, handover context and ward-specific operational logistics and capacity information.



\textbf{Overview}



Coordinate ward patient flow, handover context and ward-specific operational logistics and capacity information. The reviewed hierarchy assigns this enduring ability to its stated business-enabling context.



\textbf{Intent}



Coordinate ward patient flow, handover context and ward-specific operational logistics and capacity information.



\textbf{Boundary}



Information administration and coordination preserve the authority of the responsible source, clinical actor and governed domain; they confer no additional real-world management authority.


{\footnotesize \href{../../../../continuum/strategy/capability/3249647f7b9846c899f982ab339c81f9.md}{Canonical record: 3249647f7b9846c899f982ab339c81f9}; full qualifications and provenance.}\par
\subparagraph*{II.1.3.6.24.2.1 [Feature] Clinical Handover Context Assembly}
\phantomsection
\label{navigation:II.1.3.6.24.2.1}
\addcontentsline{toc}{subparagraph}{II.1.3.6.24.2.1 [Feature] Clinical Handover Context Assembly}

\label{element:f8c748aa08444723b784887174831d51}
\index{[Feature] Clinical Handover Context Assembly (II.1.3.6.24.2.1)}
{\small Capability / Feature · Canonical · current}\par

Collate current ward patient lists, active medical concerns, and pending tasks for nursing shift handovers.



Handover context assembly brings patient lists, current concerns and pending tasks together for the participants conducting a shift handover. Assembly is the established responsibility: it does not perform clinical review, decide pending clinical work or manage the handover itself.


{\footnotesize \href{../../../../continuum/strategy/capability/feature/f8c748aa08444723b784887174831d51.md}{Canonical record: f8c748aa08444723b784887174831d51}; full qualifications and provenance.}\par
\subparagraph*{II.1.3.6.24.2.2 [Feature] Ward Care Capacity Management}
\phantomsection
\label{navigation:II.1.3.6.24.2.2}
\addcontentsline{toc}{subparagraph}{II.1.3.6.24.2.2 [Feature] Ward Care Capacity Management}

\label{element:0d45b0afc0584afba4bd0414f5f10e37}
\index{[Feature] Ward Care Capacity Management (II.1.3.6.24.2.2)}
{\small Capability / Feature · Proposed · current}\par

Represent ward care capacity in relation to staffing, care-place suitability and patient dependency.



\textbf{Overview}



A clean bed can remain unusable when staffing or patient needs constrain the ward; source clinical facts remain attributed.



\textbf{Boundary}



Local responsibility retains its capacity meaning; Service Capacity Management consumes it without acquiring ownership or imposing one common capacity model.


{\footnotesize \href{../../../../continuum/strategy/capability/feature/0d45b0afc0584afba4bd0414f5f10e37.md}{Canonical record: 0d45b0afc0584afba4bd0414f5f10e37}; full qualifications and provenance.}\par
\subparagraph*{II.1.3.6.24.2.3 [Feature] Ward Patient Flow Coordination}
\phantomsection
\label{navigation:II.1.3.6.24.2.3}
\addcontentsline{toc}{subparagraph}{II.1.3.6.24.2.3 [Feature] Ward Patient Flow Coordination}

\label{element:011d8c6481264ff58b7866a2c74a235b}
\index{[Feature] Ward Patient Flow Coordination (II.1.3.6.24.2.3)}
{\small Capability / Feature · Canonical · current}\par

Coordinate inpatient ward occupancy, care-place allocation, and patient bed status across care facilities.



Ward flow coordination connects occupancy, place allocation and bed status within the ward's operational context. Those facts contribute to ward capacity but do not exhaust its acuity, staffing and cohorting responsibilities or prove that another patient can be accepted.


{\footnotesize \href{../../../../continuum/strategy/capability/feature/011d8c6481264ff58b7866a2c74a235b.md}{Canonical record: 011d8c6481264ff58b7866a2c74a235b}; full qualifications and provenance.}\par
\subsubsection*{II.1.3.6.25 [CT2] Work Allocation \& Dispatch}
\phantomsection
\label{navigation:II.1.3.6.25}
\addcontentsline{toc}{subsubsection}{II.1.3.6.25 [CT2] Work Allocation \& Dispatch}

{\small Historical predecessor: R.2.159 \hyperref[element:f8c13b6f74b840cc96fbeada5865455a]{[CT2] Work Allocation \& Dispatch}; Canonical ID: \texttt{\seqsplit{f8c13b6f74b840cc96fbeada5865455a}}. Superseded historical identity — direct correspondence · superseded.}\par
\label{element:152a5252952c405994e9b171234efcb0}
\index{[CT2] Work Allocation \& Dispatch (II.1.3.6.25)}
{\small Capability · Canonical · current}\par

Work Allocation \& Dispatch progresses authorised operational work from a recognised requirement through matching, dispatch, acknowledgement and oversight. It distinguishes an instantiated work item from its assignment and from evidence of performance, allowing stalled activity to be identified and escalated. Examples such as cleaning or transport concern real-world support work, not technical message queues.



\textbf{Overview}



Work Allocation \& Dispatch progresses authorised operational work from a recognised requirement through matching, dispatch, acknowledgement and oversight. It distinguishes an instantiated work item from its assignment and from evidence of performance, allowing stalled activity to be identified and escalated. Examples such as cleaning or transport concern real-world support work, not technical message queues.



\textbf{Intent}



Work Allocation \& Dispatch progresses authorised operational work from a recognised requirement through matching, dispatch, acknowledgement and oversight. It distinguishes an instantiated work item from its assignment and from evidence of performance, allowing stalled activity to be identified and escalated. Examples such as cleaning or transport concern real-world support work, not technical message queues.



\textbf{Boundary}



Information administration does not confer authority over the represented person, organisation, service, place, device or clinical activity.



\medskip\hrule\medskip

{\footnotesize \href{../../../../continuum/strategy/capability/152a5252952c405994e9b171234efcb0.md}{Canonical record: 152a5252952c405994e9b171234efcb0}; full qualifications and provenance.}\par
\paragraph*{II.1.3.6.25.1 [CT3] Operational Work Assignment \& Dispatch Coordination}
\phantomsection
\label{navigation:II.1.3.6.25.1}
\addcontentsline{toc}{paragraph}{II.1.3.6.25.1 [CT3] Operational Work Assignment \& Dispatch Coordination}

{\small Historical predecessor: R.2.158 \hyperref[element:b42afda9e876463ba56fa8f73fe7a32a]{[CT3] Operational Work Assignment \& Dispatch Coordination}; Canonical ID: \texttt{\seqsplit{b42afda9e876463ba56fa8f73fe7a32a}}. Superseded historical identity — direct correspondence · superseded.}\par
\label{element:3b58935cfa974176989c3c75d76a00fe}
\index{[CT3] Operational Work Assignment \& Dispatch Coordination (II.1.3.6.25.1)}
{\small Capability · Proposed · current}\par

Coordinate qualified operational work assignment, dispatch, acknowledgement, progression and escalation within assigned authority.



\textbf{Overview}



Coordinate qualified operational work assignment, dispatch, acknowledgement, progression and escalation within assigned authority. The reviewed hierarchy assigns this enduring ability to its stated business-enabling context.



\textbf{Intent}



Coordinate qualified operational work assignment, dispatch, acknowledgement, progression and escalation within assigned authority.



\textbf{Boundary}



Information administration and coordination preserve the authority of the responsible source, clinical actor and governed domain; they confer no additional real-world management authority.


{\footnotesize \href{../../../../continuum/strategy/capability/3b58935cfa974176989c3c75d76a00fe.md}{Canonical record: 3b58935cfa974176989c3c75d76a00fe}; full qualifications and provenance.}\par
\subparagraph*{II.1.3.6.25.1.1 [Feature] Work Dispatch Delivery}
\phantomsection
\label{navigation:II.1.3.6.25.1.1}
\addcontentsline{toc}{subparagraph}{II.1.3.6.25.1.1 [Feature] Work Dispatch Delivery}

\label{element:095826f788b64784af38af31bbe77105}
\index{[Feature] Work Dispatch Delivery (II.1.3.6.25.1.1)}
{\small Capability / Feature · Canonical · current}\par

Deliver work orders to target workers, devices, or team queues with technical and operational acknowledgement capture.



Dispatch delivery brings assigned operational work to the intended worker or team and retains acknowledgement evidence. Technical receipt and operational acknowledgement must remain distinguishable; neither alone proves that the work was performed or its intended outcome achieved.


{\footnotesize \href{../../../../continuum/strategy/capability/feature/095826f788b64784af38af31bbe77105.md}{Canonical record: 095826f788b64784af38af31bbe77105}; full qualifications and provenance.}\par
\subparagraph*{II.1.3.6.25.1.2 [Feature] Work Item Instantiation}
\phantomsection
\label{navigation:II.1.3.6.25.1.2}
\addcontentsline{toc}{subparagraph}{II.1.3.6.25.1.2 [Feature] Work Item Instantiation}

\label{element:321fa656c71f4019842b5b991dbc9ff4}
\index{[Feature] Work Item Instantiation (II.1.3.6.25.1.2)}
{\small Capability / Feature · Canonical · current}\par

Instantiate operational units of work from clinical triggers (e.g., bed cleaning required, specimen transport required).



Work-item instantiation gives an authorised operational requirement an identifiable unit of work. It distinguishes recognising that support activity is needed from assigning it, dispatching it and recording its completion; a clinical trigger does not make the resulting support item clinical work authority.


{\footnotesize \href{../../../../continuum/strategy/capability/feature/321fa656c71f4019842b5b991dbc9ff4.md}{Canonical record: 321fa656c71f4019842b5b991dbc9ff4}; full qualifications and provenance.}\par
\subparagraph*{II.1.3.6.25.1.3 [Feature] Work Progress Oversight \& Escalation}
\phantomsection
\label{navigation:II.1.3.6.25.1.3}
\addcontentsline{toc}{subparagraph}{II.1.3.6.25.1.3 [Feature] Work Progress Oversight \& Escalation}

\label{element:abc12faeffd24917aad6748b0f0ca8d7}
\index{[Feature] Work Progress Oversight \& Escalation (II.1.3.6.25.1.3)}
{\small Capability / Feature · Canonical · current}\par

Monitor work progression against time-to-completion thresholds, escalating stalled or unacknowledged tasks.



Progress oversight identifies stalled or unacknowledged operational work against applicable timing expectations. Escalation serves completion management and does not itself constitute independent assurance or transfer authority over clinical decisions related to the support activity.


{\footnotesize \href{../../../../continuum/strategy/capability/feature/abc12faeffd24917aad6748b0f0ca8d7.md}{Canonical record: abc12faeffd24917aad6748b0f0ca8d7}; full qualifications and provenance.}\par
\subparagraph*{II.1.3.6.25.1.4 [Feature] Worker Matching \& Allocation}
\phantomsection
\label{navigation:II.1.3.6.25.1.4}
\addcontentsline{toc}{subparagraph}{II.1.3.6.25.1.4 [Feature] Worker Matching \& Allocation}

\label{element:caac8a83dc974cdb9c9bc1b778b9cee7}
\index{[Feature] Worker Matching \& Allocation (II.1.3.6.25.1.4)}
{\small Capability / Feature · Canonical · current}\par

Match work items against available staff roles, credentials, and locations based on allocation policies.



Worker matching relates an operational requirement to eligible roles, credentials, locations and availability under policy. It consumes those constraints rather than changing their authoritative records, and does not turn operational allocation into permission to decide clinical work.


{\footnotesize \href{../../../../continuum/strategy/capability/feature/caac8a83dc974cdb9c9bc1b778b9cee7.md}{Canonical record: caac8a83dc974cdb9c9bc1b778b9cee7}; full qualifications and provenance.}\par
\subsubsection*{II.1.3.6.26 [CT2] Workforce Qualification \& Development Enablement [superseded]}
\phantomsection
\label{navigation:II.1.3.6.26}
\addcontentsline{toc}{subsubsection}{II.1.3.6.26 [CT2] Workforce Qualification \& Development Enablement [superseded]}

\label{element:d24a1f4c4f674c45a95c317998ff4186}
\index{[CT2] Workforce Qualification \& Development Enablement (II.1.3.6.26)}
{\small Capability · Proposed · superseded}\par
\href{../../../../continuum/strategy/capability/d24a1f4c4f674c45a95c317998ff4186.md}{Historical canonical definition and provenance}.\par
\paragraph*{II.1.3.6.26.1 [CT3] Credential Management}
\phantomsection
\label{navigation:II.1.3.6.26.1}
\addcontentsline{toc}{paragraph}{II.1.3.6.26.1 [CT3] Credential Management}

\label{element:3a7a4d2d0ce8446b8a6293284fb16b83}
\index{[CT3] Credential Management (II.1.3.6.26.1)}
{\small Capability · Canonical · pending\_adjudication}\par

Definition unresolved in the baseline. No definition has been synthesised from the name.



\medskip\hrule\medskip

{\small Recorded unresolved matters: source-qualified matters; exact scope remains in the canonical record.}\par
{\footnotesize \href{../../../../continuum/strategy/capability/3a7a4d2d0ce8446b8a6293284fb16b83.md}{Canonical record: 3a7a4d2d0ce8446b8a6293284fb16b83}; full qualifications and provenance.}\par
\subsection*{II.1.3.7 [CT1] Human Collaboration \& Work Enablement [superseded]}
\phantomsection
\label{navigation:II.1.3.7}
\addcontentsline{toc}{subsection}{II.1.3.7 [CT1] Human Collaboration \& Work Enablement [superseded]}

\label{element:c9347f73b21242c5b8cce41252b74da0}
\index{[CT1] Human Collaboration \& Work Enablement (II.1.3.7)}
{\small Capability · Proposed · superseded}\par
\href{../../../../continuum/strategy/capability/c9347f73b21242c5b8cce41252b74da0.md}{Historical canonical definition and provenance}.\par
\subsubsection*{II.1.3.7.1 [CT2] Collaborative Interaction Enablement [superseded]}
\phantomsection
\label{navigation:II.1.3.7.1}
\addcontentsline{toc}{subsubsection}{II.1.3.7.1 [CT2] Collaborative Interaction Enablement [superseded]}

\label{element:6140b21c874241d092708e2411caceed}
\index{[CT2] Collaborative Interaction Enablement (II.1.3.7.1)}
{\small Capability · Proposed · superseded}\par
\href{../../../../continuum/strategy/capability/6140b21c874241d092708e2411caceed.md}{Historical canonical definition and provenance}.\par
\subsubsection*{II.1.3.7.2 [CT2] Operational Activity \& Temporal Enablement [superseded]}
\phantomsection
\label{navigation:II.1.3.7.2}
\addcontentsline{toc}{subsubsection}{II.1.3.7.2 [CT2] Operational Activity \& Temporal Enablement [superseded]}

\label{element:ca96b2b2ff434854995711e403e95f6f}
\index{[CT2] Operational Activity \& Temporal Enablement (II.1.3.7.2)}
{\small Capability · Proposed · superseded}\par
\href{../../../../continuum/strategy/capability/ca96b2b2ff434854995711e403e95f6f.md}{Historical canonical definition and provenance}.\par
\paragraph*{II.1.3.7.2.1 [CT3] Entity-Specific Operational Coordination}
\phantomsection
\label{navigation:II.1.3.7.2.1}
\addcontentsline{toc}{paragraph}{II.1.3.7.2.1 [CT3] Entity-Specific Operational Coordination}

\label{element:9f0c1429b64f4f88982fbbda6eaec233}
\index{[CT3] Entity-Specific Operational Coordination (II.1.3.7.2.1)}
{\small Capability · Proposed · pending\_adjudication}\par

Coordinate governed operational progression relative to one represented entity, including interested behaviour participation and controlled recovery.



\textbf{Overview}



Activity concerning the same patient, practitioner or bed needs a coherent entity-relative frame. Interest fan-out, serial progression, fresh information context and uncertain effects remain distinguishable from generic multi-work-unit progression.



\textbf{Intent}



Keep the progression of one represented entity coherent across interested operational behaviours.



\textbf{Boundary}



Entity coordination does not own all entity information, determine clinical work or become the source of clinical decisions. No Twin component allocation follows from this proposal.


{\small Recorded unresolved matters: source-qualified matters; exact scope remains in the canonical record.}\par
{\footnotesize \href{../../../../continuum/strategy/capability/9f0c1429b64f4f88982fbbda6eaec233.md}{Canonical record: 9f0c1429b64f4f88982fbbda6eaec233}; full qualifications and provenance.}\par
\subparagraph*{II.1.3.7.2.1.1 [Feature] Entity Behaviour Failure \& Recovery Governance}
\phantomsection
\label{navigation:II.1.3.7.2.1.1}
\addcontentsline{toc}{subparagraph}{II.1.3.7.2.1.1 [Feature] Entity Behaviour Failure \& Recovery Governance}

\label{element:bdfaeebab74e4d4d96a5e3799a3aa07b}
\index{[Feature] Entity Behaviour Failure \& Recovery Governance (II.1.3.7.2.1.1)}
{\small Capability / Feature · Proposed · pending\_adjudication}\par

Apply explicit behaviour-failure dispositions and controlled, authorised resumption for entity-relative processing.



\textbf{Overview}



Failure may stop a behaviour, an input or the entity progression and may leave effects uncertain. Resumption must preserve that distinction rather than imply rollback or automatic safe retry.



\textbf{Boundary}



Recovery cannot silently declare no effects, bypass governance or resume a halted entity automatically.


{\small Recorded unresolved matters: source-qualified matters; exact scope remains in the canonical record.}\par
{\footnotesize \href{../../../../continuum/strategy/capability/feature/bdfaeebab74e4d4d96a5e3799a3aa07b.md}{Canonical record: bdfaeebab74e4d4d96a5e3799a3aa07b}; full qualifications and provenance.}\par
\subparagraph*{II.1.3.7.2.1.2 [Feature] Entity-Relative Activity Coordination}
\phantomsection
\label{navigation:II.1.3.7.2.1.2}
\addcontentsline{toc}{subparagraph}{II.1.3.7.2.1.2 [Feature] Entity-Relative Activity Coordination}

\label{element:b9960d7de097464ea6de94929ca48c89}
\index{[Feature] Entity-Relative Activity Coordination (II.1.3.7.2.1.2)}
{\small Capability / Feature · Proposed · pending\_adjudication}\par

Sequence applicable behaviours for an identifiable entity input while preserving its operational reference frame.



\textbf{Overview}



Interest matching can select multiple behaviours while the entity progresses one input at a time. Each behaviour uses a fresh managed information context.



\textbf{Boundary}



Fan-out is not concurrent execution; clinical decisions and originating information remain with their owners.


{\small Recorded unresolved matters: source-qualified matters; exact scope remains in the canonical record.}\par
{\footnotesize \href{../../../../continuum/strategy/capability/feature/b9960d7de097464ea6de94929ca48c89.md}{Canonical record: b9960d7de097464ea6de94929ca48c89}; full qualifications and provenance.}\par
\paragraph*{II.1.3.7.2.2 [CT3] Governed Operational Activity Design}
\phantomsection
\label{navigation:II.1.3.7.2.2}
\addcontentsline{toc}{paragraph}{II.1.3.7.2.2 [CT3] Governed Operational Activity Design}

\label{element:3a96040821c14fe38830411dab4f8453}
\index{[CT3] Governed Operational Activity Design (II.1.3.7.2.2)}
{\small Capability · Proposed · pending\_adjudication}\par

Define the information context, expected outcomes, continuation conditions and housekeeping of governed operational activity.



\textbf{Overview}



An activity needs a fulfilment definition even when only one bounded execution unit is involved. This capability specifies its semantic envelope while execution and event delivery remain separate responsibilities.



\textbf{Intent}



Make the expected information and fulfilment basis of an undertaking explicit before it executes.



\textbf{Boundary}



An activity definition is neither an orchestrator nor a choreographer. It does not choose clinical care or dynamically invent execution logic.


{\small Recorded unresolved matters: source-qualified matters; exact scope remains in the canonical record.}\par
{\footnotesize \href{../../../../continuum/strategy/capability/3a96040821c14fe38830411dab4f8453.md}{Canonical record: 3a96040821c14fe38830411dab4f8453}; full qualifications and provenance.}\par
\subparagraph*{II.1.3.7.2.2.1 [Feature] Activity Outcome \& Continuation Rules}
\phantomsection
\label{navigation:II.1.3.7.2.2.1}
\addcontentsline{toc}{subparagraph}{II.1.3.7.2.2.1 [Feature] Activity Outcome \& Continuation Rules}

\label{element:34c8309244c142ec8772dcbb11ecc4ef}
\index{[Feature] Activity Outcome \& Continuation Rules (II.1.3.7.2.2.1)}
{\small Capability / Feature · Proposed · pending\_adjudication}\par

Define how output roles, expected outcomes and continuation conditions are interpreted for an undertaking.



\textbf{Overview}



An output may contribute to completion, subsequent activity or retained evidence under different rules. Those meanings need to be specified rather than inferred from a returned payload.



\textbf{Boundary}



Completion policy is not business outcome truth; uncertain effects remain explicit.


{\small Recorded unresolved matters: source-qualified matters; exact scope remains in the canonical record.}\par
{\footnotesize \href{../../../../continuum/strategy/capability/feature/34c8309244c142ec8772dcbb11ecc4ef.md}{Canonical record: 34c8309244c142ec8772dcbb11ecc4ef}; full qualifications and provenance.}\par
\subparagraph*{II.1.3.7.2.2.2 [Feature] Operational Activity Context Definition}
\phantomsection
\label{navigation:II.1.3.7.2.2.2}
\addcontentsline{toc}{subparagraph}{II.1.3.7.2.2.2 [Feature] Operational Activity Context Definition}

\label{element:58dcd59389664b7d83d14b536d7d03d0}
\index{[Feature] Operational Activity Context Definition (II.1.3.7.2.2.2)}
{\small Capability / Feature · Proposed · pending\_adjudication}\par

Define the anticipated information context and expected fulfilment basis of an operational undertaking.



\textbf{Overview}



Even a one-unit activity needs a declared information and outcome frame. The definition is separate from loading values or running the business logic.



\textbf{Boundary}



No dynamically invented resource set or orchestration engine is implied.


{\small Recorded unresolved matters: source-qualified matters; exact scope remains in the canonical record.}\par
{\footnotesize \href{../../../../continuum/strategy/capability/feature/58dcd59389664b7d83d14b536d7d03d0.md}{Canonical record: 58dcd59389664b7d83d14b536d7d03d0}; full qualifications and provenance.}\par
\paragraph*{II.1.3.7.2.3 [CT3] Workflow \& Activity Coordination}
\phantomsection
\label{navigation:II.1.3.7.2.3}
\addcontentsline{toc}{paragraph}{II.1.3.7.2.3 [CT3] Workflow \& Activity Coordination}

\label{element:94c683c7588f475e973b58cb226345db}
\index{[CT3] Workflow \& Activity Coordination (II.1.3.7.2.3)}
{\small Capability · Canonical · pending\_adjudication}\par

Definition unresolved in the baseline. No definition has been synthesised from the name.



\medskip\hrule\medskip

{\small Recorded unresolved matters: source-qualified matters; exact scope remains in the canonical record.}\par
{\footnotesize \href{../../../../continuum/strategy/capability/94c683c7588f475e973b58cb226345db.md}{Canonical record: 94c683c7588f475e973b58cb226345db}; full qualifications and provenance.}\par
\subparagraph*{II.1.3.7.2.3.1 [Feature] Controlled Administrative Recovery}
\phantomsection
\label{navigation:II.1.3.7.2.3.1}
\addcontentsline{toc}{subparagraph}{II.1.3.7.2.3.1 [Feature] Controlled Administrative Recovery}

\label{element:21d0601a852043a4b0b1d4fe841108d2}
\index{[Feature] Controlled Administrative Recovery (II.1.3.7.2.3.1)}
{\small Capability / Feature · Proposed · pending\_adjudication}\par

Perform authorised recovery as identifiable governed activity through normal execution and evidence boundaries.



\textbf{Overview}



Operators may need to resume, reconcile or repair an operational concern. That is first-class work with permission and outcome standing, rather than direct mutation through a console.



\textbf{Boundary}



Administrative interaction grants no raw cache/database or subsystem-internal bypass.


{\small Recorded unresolved matters: source-qualified matters; exact scope remains in the canonical record.}\par
{\footnotesize \href{../../../../continuum/strategy/capability/feature/21d0601a852043a4b0b1d4fe841108d2.md}{Canonical record: 21d0601a852043a4b0b1d4fe841108d2}; full qualifications and provenance.}\par
\subparagraph*{II.1.3.7.2.3.2 [Feature] Synthetic Task Orchestration}
\phantomsection
\label{navigation:II.1.3.7.2.3.2}
\addcontentsline{toc}{subparagraph}{II.1.3.7.2.3.2 [Feature] Synthetic Task Orchestration}

\label{element:7cf9dd3aeeb84a83b2445f897db9bd4e}
\index{[Feature] Synthetic Task Orchestration (II.1.3.7.2.3.2)}
{\small Capability / Feature · Canonical · pending\_adjudication}\par

Execute non-human, automated platform activity units and sequence orchestration.



\medskip\hrule\medskip

{\small Recorded unresolved matters: source-qualified matters; exact scope remains in the canonical record.}\par
{\footnotesize \href{../../../../continuum/strategy/capability/feature/7cf9dd3aeeb84a83b2445f897db9bd4e.md}{Canonical record: 7cf9dd3aeeb84a83b2445f897db9bd4e}; full qualifications and provenance.}\par
\subparagraph*{II.1.3.7.2.3.3 [Feature] To Do Decision \& Approval Coordination}
\phantomsection
\label{navigation:II.1.3.7.2.3.3}
\addcontentsline{toc}{subparagraph}{II.1.3.7.2.3.3 [Feature] To Do Decision \& Approval Coordination}

\label{element:377a40e8f2d64840a10dd55f87a9d005}
\index{[Feature] To Do Decision \& Approval Coordination (II.1.3.7.2.3.3)}
{\small Capability / Feature · Canonical · pending\_adjudication}\par

Coordinate human clinical review items, sign-offs, and approval workflows.



\medskip\hrule\medskip

{\small Recorded unresolved matters: source-qualified matters; exact scope remains in the canonical record.}\par
{\footnotesize \href{../../../../continuum/strategy/capability/feature/377a40e8f2d64840a10dd55f87a9d005.md}{Canonical record: 377a40e8f2d64840a10dd55f87a9d005}; full qualifications and provenance.}\par
\subparagraph*{II.1.3.7.2.3.4 [Feature] Uncertain Activity Outcome Reconciliation}
\phantomsection
\label{navigation:II.1.3.7.2.3.4}
\addcontentsline{toc}{subparagraph}{II.1.3.7.2.3.4 [Feature] Uncertain Activity Outcome Reconciliation}

\label{element:eb06ecc956404d0397701ceafb5fbd46}
\index{[Feature] Uncertain Activity Outcome Reconciliation (II.1.3.7.2.3.4)}
{\small Capability / Feature · Proposed · pending\_adjudication}\par

Reconcile indeterminate activity effects through authoritative observation before establishing a safe next disposition.



\textbf{Overview}



A timeout or absent response leaves a materially significant outcome unknown. Reconciliation can confirm an effect, establish absence or retain uncertainty without automatically retrying the original work.



\textbf{Boundary}



Execution uncertainty differs from insufficient assurance evidence; a retry does not prove the earlier attempt had no effect.


{\small Recorded unresolved matters: source-qualified matters; exact scope remains in the canonical record.}\par
{\footnotesize \href{../../../../continuum/strategy/capability/feature/eb06ecc956404d0397701ceafb5fbd46.md}{Canonical record: eb06ecc956404d0397701ceafb5fbd46}; full qualifications and provenance.}\par
\subparagraph*{II.1.3.7.2.3.5 [Feature] Work Order Progression}
\phantomsection
\label{navigation:II.1.3.7.2.3.5}
\addcontentsline{toc}{subparagraph}{II.1.3.7.2.3.5 [Feature] Work Order Progression}

\label{element:5ca494e85869419c867618fd9648d03a}
\index{[Feature] Work Order Progression (II.1.3.7.2.3.5)}
{\small Capability / Feature · Canonical · pending\_adjudication}\par

Manage execution states and dispatch of human operational work orders.



\medskip\hrule\medskip

{\small Recorded unresolved matters: source-qualified matters; exact scope remains in the canonical record.}\par
{\footnotesize \href{../../../../continuum/strategy/capability/feature/5ca494e85869419c867618fd9648d03a.md}{Canonical record: 5ca494e85869419c867618fd9648d03a}; full qualifications and provenance.}\par
\subparagraph*{II.1.3.7.2.3.6 [Feature] Workflow Timeout \& Escalation Oversight}
\phantomsection
\label{navigation:II.1.3.7.2.3.6}
\addcontentsline{toc}{subparagraph}{II.1.3.7.2.3.6 [Feature] Workflow Timeout \& Escalation Oversight}

\label{element:ec55d06e393e4988a84c6157b719bd40}
\index{[Feature] Workflow Timeout \& Escalation Oversight (II.1.3.7.2.3.6)}
{\small Capability / Feature · Canonical · pending\_adjudication}\par

Supervise operational deadlines, executing automated timeout notifications and supervisory escalations.



\medskip\hrule\medskip

{\small Recorded unresolved matters: source-qualified matters; exact scope remains in the canonical record.}\par
{\footnotesize \href{../../../../continuum/strategy/capability/feature/ec55d06e393e4988a84c6157b719bd40.md}{Canonical record: ec55d06e393e4988a84c6157b719bd40}; full qualifications and provenance.}\par
\subsection*{II.1.3.8 [CT1] Independent Assurance Enablement}
\phantomsection
\label{navigation:II.1.3.8}
\addcontentsline{toc}{subsection}{II.1.3.8 [CT1] Independent Assurance Enablement}

{\small Historical predecessor: R.2.166 \hyperref[element:88bf57b7a92c4c1da1b5f0a67d3ebd82]{[CT1] Independent Assurance Enablement}; Canonical ID: \texttt{\seqsplit{88bf57b7a92c4c1da1b5f0a67d3ebd82}}. Superseded historical identity — direct correspondence · superseded.}\par
\label{element:616682b612d244429e656b040e002d56}
\index{[CT1] Independent Assurance Enablement (II.1.3.8)}
{\small Capability · Proposed · current}\par

Enable assurance specification and independent review of the assurance service while preserving the distinct Canonical responsibility for assurance against a governed subject.



\textbf{Overview}



Enable assurance specification and independent review of the assurance service while preserving the distinct Canonical responsibility for assurance against a governed subject.



\textbf{Intent}



Enable assurance specification and independent review of the assurance service while preserving the distinct Canonical responsibility for assurance against a governed subject.



\textbf{Boundary}



Assurance cannot manage, remediate or approve its subject merely because it observes it. Clinical adequacy, financial governance and source-information authority remain outside this mandate.


{\footnotesize \href{../../../../continuum/strategy/capability/616682b612d244429e656b040e002d56.md}{Canonical record: 616682b612d244429e656b040e002d56}; full qualifications and provenance.}\par
\subsubsection*{II.1.3.8.1 [CT2] Assurance Service Review}
\phantomsection
\label{navigation:II.1.3.8.1}
\addcontentsline{toc}{subsubsection}{II.1.3.8.1 [CT2] Assurance Service Review}

{\small Historical predecessor: R.2.162 \hyperref[element:b345cb12b45746409302899514126170]{[CT2] Assurance Service Review}; Canonical ID: \texttt{\seqsplit{b345cb12b45746409302899514126170}}. Superseded historical identity — direct correspondence · superseded.}\par
\label{element:b1a791458d644d4794895e4fa42ec83c}
\index{[CT2] Assurance Service Review (II.1.3.8.1)}
{\small Capability · Proposed · current}\par

Independently review the operation, effectiveness, appropriateness and conformance of the assurance service itself.



\textbf{Overview}



Independently review the operation, effectiveness, appropriateness and conformance of the assurance service itself. The reviewed hierarchy assigns this enduring ability to its stated business-enabling context.



\textbf{Intent}



Independently review the operation, effectiveness, appropriateness and conformance of the assurance service itself.



\textbf{Boundary}



Reviews the assurance service itself. Canonical evidence association, subject assessment, adjudication, findings and conclusions are retained separately pending adjudication.


{\footnotesize \href{../../../../continuum/strategy/capability/b1a791458d644d4794895e4fa42ec83c.md}{Canonical record: b1a791458d644d4794895e4fa42ec83c}; full qualifications and provenance.}\par
\paragraph*{II.1.3.8.1.1 [CT3] Assurance Service Evaluation}
\phantomsection
\label{navigation:II.1.3.8.1.1}
\addcontentsline{toc}{paragraph}{II.1.3.8.1.1 [CT3] Assurance Service Evaluation}

{\small Historical predecessor: R.2.161 \hyperref[element:be34b54314c6444f8418d645f6c83b6d]{[CT3] Assurance Service Evaluation}; Canonical ID: \texttt{\seqsplit{be34b54314c6444f8418d645f6c83b6d}}. Superseded historical identity — direct correspondence · superseded.}\par
\label{element:039343518db542d48cde6afab84a3e1e}
\index{[CT3] Assurance Service Evaluation (II.1.3.8.1.1)}
{\small Capability · Proposed · current}\par

Evaluate the operation, effectiveness, appropriateness and conformance of the assurance service itself through independent review.



\textbf{Overview}



Evaluate the operation, effectiveness, appropriateness and conformance of the assurance service itself through independent review. The reviewed hierarchy assigns this enduring ability to its stated business-enabling context.



\textbf{Intent}



Evaluate the operation, effectiveness, appropriateness and conformance of the assurance service itself through independent review.



\textbf{Boundary}



Evaluation concerns operation, effectiveness, appropriateness and conformance of the assurance service; it is not execution of assurance against a subject.


{\footnotesize \href{../../../../continuum/strategy/capability/039343518db542d48cde6afab84a3e1e.md}{Canonical record: 039343518db542d48cde6afab84a3e1e}; full qualifications and provenance.}\par
\subsubsection*{II.1.3.8.2 [CT2] Assurance Specification}
\phantomsection
\label{navigation:II.1.3.8.2}
\addcontentsline{toc}{subsubsection}{II.1.3.8.2 [CT2] Assurance Specification}

{\small Historical predecessor: R.2.165 \hyperref[element:8006ee0dd95f4d7fb1f408fa09b6ce19]{[CT2] Assurance Specification}; Canonical ID: \texttt{\seqsplit{8006ee0dd95f4d7fb1f408fa09b6ce19}}. Superseded historical identity — direct correspondence · superseded.}\par
\label{element:b81775d89d774cf8ba61731d13c67a54}
\index{[CT2] Assurance Specification (II.1.3.8.2)}
{\small Capability · Proposed · current}\par

Model how satisfaction will be assured and maintain applicable, temporally qualified criteria.



\textbf{Overview}



Model how satisfaction will be assured and maintain applicable, temporally qualified criteria.



\textbf{Intent}



Model how satisfaction will be assured and maintain applicable, temporally qualified criteria.



\textbf{Boundary}



Modelling and criteria stewardship do not approve governing requirements or perform independent conclusion-making.


{\footnotesize \href{../../../../continuum/strategy/capability/b81775d89d774cf8ba61731d13c67a54.md}{Canonical record: b81775d89d774cf8ba61731d13c67a54}; full qualifications and provenance.}\par
\paragraph*{II.1.3.8.2.1 [CT3] Assurance Criteria Administration}
\phantomsection
\label{navigation:II.1.3.8.2.1}
\addcontentsline{toc}{paragraph}{II.1.3.8.2.1 [CT3] Assurance Criteria Administration}

{\small Historical predecessor: R.2.163 \hyperref[element:6dab08ab1e3249d7b314843ca9ed0149]{[CT3] Assurance Criteria Administration}; Canonical ID: \texttt{\seqsplit{6dab08ab1e3249d7b314843ca9ed0149}}. Superseded historical identity — direct correspondence · superseded.}\par
\label{element:1bfd71f24321428c8871b587d413f47e}
\index{[CT3] Assurance Criteria Administration (II.1.3.8.2.1)}
{\small Capability · Canonical · current}\par

The capability to establish and manage reusable, governed and temporally identifiable assurance criteria through which subjects may be evaluated.



Definition — reviewed candidate



Administer reusable governed assurance criteria and their version-specific and temporal applicability without changing governing requirement authority.



\textbf{Overview}



Administer reusable governed assurance criteria and their version-specific and temporal applicability without changing governing requirement authority.



\textbf{Intent}



Administer reusable governed assurance criteria and their version-specific and temporal applicability without changing governing requirement authority.



\textbf{Boundary}



Information administration does not confer authority over the represented person, organisation, service, place, device or clinical activity.



\medskip\hrule\medskip

{\small Recorded unresolved matters: source-qualified matters; exact scope remains in the canonical record.}\par
{\footnotesize \href{../../../../continuum/strategy/capability/1bfd71f24321428c8871b587d413f47e.md}{Canonical record: 1bfd71f24321428c8871b587d413f47e}; full qualifications and provenance.}\par
\subparagraph*{II.1.3.8.2.1.1 [Feature] Temporal Assurance Criteria Stewardship}
\phantomsection
\label{navigation:II.1.3.8.2.1.1}
\addcontentsline{toc}{subparagraph}{II.1.3.8.2.1.1 [Feature] Temporal Assurance Criteria Stewardship}

\label{element:e345a4b1eef846a7b77568f519e2c75c}
\index{[Feature] Temporal Assurance Criteria Stewardship (II.1.3.8.2.1.1)}
{\small Capability / Feature · Proposed · current}\par

Manage identifiable assurance criteria and their relevant temporal and version applicability.



\textbf{Overview}



Evidence must be evaluated against the criteria applicable to its purpose and time rather than an unqualified current rule.



\textbf{Boundary}



Criteria stewardship does not approve governing requirements or select a universal result taxonomy.


{\footnotesize \href{../../../../continuum/strategy/capability/feature/e345a4b1eef846a7b77568f519e2c75c.md}{Canonical record: e345a4b1eef846a7b77568f519e2c75c}; full qualifications and provenance.}\par
\paragraph*{II.1.3.8.2.2 [CT3] Assurance Design}
\phantomsection
\label{navigation:II.1.3.8.2.2}
\addcontentsline{toc}{paragraph}{II.1.3.8.2.2 [CT3] Assurance Design}

{\small Historical predecessor: R.2.164 \hyperref[element:e1b8db7c95fa4e7ea296351a03532e24]{[CT3] Assurance Design}; Canonical ID: \texttt{\seqsplit{e1b8db7c95fa4e7ea296351a03532e24}}. Superseded historical identity — direct correspondence · superseded.}\par
\label{element:71c8a612b88f4ecba049e94d506bd0c2}
\index{[CT3] Assurance Design (II.1.3.8.2.2)}
{\small Capability · Canonical · current}\par

The capability to determine and define how satisfaction of a governed requirement, constraint or expected behaviour will be assured, including its assurance disposition, applicable criteria, required evidence and evaluation expectations.



Definition — reviewed candidate



Model how satisfaction of a governed subject or requirement will be assured, preserving applicable assurance dispositions and boundaries.



\textbf{Overview}



Model how satisfaction of a governed subject or requirement will be assured, preserving applicable assurance dispositions and boundaries.



\textbf{Intent}



Model how satisfaction of a governed subject or requirement will be assured, preserving applicable assurance dispositions and boundaries.



\textbf{Boundary}



Information administration does not confer authority over the represented person, organisation, service, place, device or clinical activity.



\medskip\hrule\medskip

{\small Recorded unresolved matters: source-qualified matters; exact scope remains in the canonical record.}\par
{\footnotesize \href{../../../../continuum/strategy/capability/71c8a612b88f4ecba049e94d506bd0c2.md}{Canonical record: 71c8a612b88f4ecba049e94d506bd0c2}; full qualifications and provenance.}\par
\subparagraph*{II.1.3.8.2.2.1 [Feature] Assurance Disposition \& Design Modelling}
\phantomsection
\label{navigation:II.1.3.8.2.2.1}
\addcontentsline{toc}{subparagraph}{II.1.3.8.2.2.1 [Feature] Assurance Disposition \& Design Modelling}

\label{element:c36180ebffe44b8db587a00e78b4b358}
\index{[Feature] Assurance Disposition \& Design Modelling (II.1.3.8.2.2.1)}
{\small Capability / Feature · Proposed · current}\par

Model how satisfaction of an applicable governed concern will be independently established.



\textbf{Overview}



This proposal provides a Feature frame for the established design responsibility, including disposition, criteria and evidence expectations, without creating an additional assurance Function.



\textbf{Boundary}



Modelling is distinct from requirement approval, subject management and evaluation.


{\footnotesize \href{../../../../continuum/strategy/capability/feature/c36180ebffe44b8db587a00e78b4b358.md}{Canonical record: c36180ebffe44b8db587a00e78b4b358}; full qualifications and provenance.}\par
\subsubsection*{II.1.3.8.3 [CT2] Independent Assurance Evaluation Enablement}
\phantomsection
\label{navigation:II.1.3.8.3}
\addcontentsline{toc}{subsubsection}{II.1.3.8.3 [CT2] Independent Assurance Evaluation Enablement}

\label{element:7fa5112efb724b90a17a03c7c19c516a}
\index{[CT2] Independent Assurance Evaluation Enablement (II.1.3.8.3)}
{\small Capability · Proposed · pending\_adjudication}\par

Associate trustworthy evidence, assess it and adjudicate independently reportable findings about governed service behaviour.



\textbf{Overview}



Associate trustworthy evidence, assess it and adjudicate independently reportable findings about governed service behaviour. Its constituent abilities distinguish Governed Assurance.



\textbf{Intent}



Produce findings supported by applicable evidence and criteria while retaining inability to conclude.



\textbf{Boundary}



The subject can supply evidence but cannot control its conclusion. Evidence insufficiency remains distinct from execution uncertainty and clinical judgement.


{\small Recorded unresolved matters: source-qualified matters; exact scope remains in the canonical record.}\par
{\footnotesize \href{../../../../continuum/strategy/capability/7fa5112efb724b90a17a03c7c19c516a.md}{Canonical record: 7fa5112efb724b90a17a03c7c19c516a}; full qualifications and provenance.}\par
\paragraph*{II.1.3.8.3.1 [CT3] Governed Assurance}
\phantomsection
\label{navigation:II.1.3.8.3.1}
\addcontentsline{toc}{paragraph}{II.1.3.8.3.1 [CT3] Governed Assurance}

\label{element:da1a983d4dcf4d5bb772b8f69bc08426}
\index{[CT3] Governed Assurance (II.1.3.8.3.1)}
{\small Capability · Canonical · pending\_adjudication}\par

The capability to independently evaluate a governed subject against applicable assurance criteria through the assembly and evaluation of sufficient trustworthy evidence, establishing explicit findings, exceptions and assurance conclusions.



\medskip\hrule\medskip

{\small Recorded unresolved matters: source-qualified matters; exact scope remains in the canonical record.}\par
{\footnotesize \href{../../../../continuum/strategy/capability/da1a983d4dcf4d5bb772b8f69bc08426.md}{Canonical record: da1a983d4dcf4d5bb772b8f69bc08426}; full qualifications and provenance.}\par
\subparagraph*{II.1.3.8.3.1.1 [Feature] Independent Evidence Assessment \& Adjudication}
\phantomsection
\label{navigation:II.1.3.8.3.1.1}
\addcontentsline{toc}{subparagraph}{II.1.3.8.3.1.1 [Feature] Independent Evidence Assessment \& Adjudication}

\label{element:3aedee8f78484b81955866bcf16b8a6d}
\index{[Feature] Independent Evidence Assessment \& Adjudication (II.1.3.8.3.1.1)}
{\small Capability / Feature · Proposed · pending\_adjudication}\par

Independently establish assurance context, assess evidence and adjudicate findings under applicable criteria.



\textbf{Overview}



The Feature unifies the review frame while retaining the three distinct canonical Functions and the boundary between criterion-level assessment and reportable conclusion.



\textbf{Boundary}



No new assurance Function or clinical-adequacy mandate is introduced; insufficient evidence establishes neither satisfaction nor non-satisfaction.


{\small Recorded unresolved matters: source-qualified matters; exact scope remains in the canonical record.}\par
{\footnotesize \href{../../../../continuum/strategy/capability/feature/3aedee8f78484b81955866bcf16b8a6d.md}{Canonical record: 3aedee8f78484b81955866bcf16b8a6d}; full qualifications and provenance.}\par
\subsection*{II.1.3.9 [CT1] Shared Information \& Coordination Enablement}
\phantomsection
\label{navigation:II.1.3.9}
\addcontentsline{toc}{subsection}{II.1.3.9 [CT1] Shared Information \& Coordination Enablement}

{\small Historical predecessor: R.2.214 \hyperref[element:eeabb9c3c526404cbe4302906560fa7e]{[CT1] Shared Information \& Coordination Enablement}; Canonical ID: \texttt{\seqsplit{eeabb9c3c526404cbe4302906560fa7e}}. Superseded historical identity — direct correspondence · superseded.}\par
\label{element:0062fcacae1740988b6be2bda2aa1674}
\index{[CT1] Shared Information \& Coordination Enablement (II.1.3.9)}
{\small Capability · Proposed · current}\par

Enable governed shared health-information use, exchange, stewardship, knowledge, collaboration, external work coordination, calendars, presentation and information design.



\textbf{Overview}



Enable governed shared health-information use, exchange, stewardship, knowledge, collaboration, external work coordination, calendars, presentation and information design.



\textbf{Intent}



Enable governed shared health-information use, exchange, stewardship, knowledge, collaboration, external work coordination, calendars, presentation and information design.



\textbf{Boundary}



Preservation, custody, discoverability and presentation do not confer originating fact authority. Clinical documentation administration and clinical-context use retain their own responsibilities.


{\footnotesize \href{../../../../continuum/strategy/capability/0062fcacae1740988b6be2bda2aa1674.md}{Canonical record: 0062fcacae1740988b6be2bda2aa1674}; full qualifications and provenance.}\par
\subsubsection*{II.1.3.9.1 [CT2] Calendar Management}
\phantomsection
\label{navigation:II.1.3.9.1}
\addcontentsline{toc}{subsubsection}{II.1.3.9.1 [CT2] Calendar Management}

{\small Historical predecessor: R.2.168 \hyperref[element:8b42aac6083441caa78b59f5e62eebd9]{[CT2] Calendar Management}; Canonical ID: \texttt{\seqsplit{8b42aac6083441caa78b59f5e62eebd9}}. Superseded historical identity — direct correspondence · superseded.}\par
\label{element:7669036f7e6749219412fc64df5f7064}
\index{[CT2] Calendar Management (II.1.3.9.1)}
{\small Capability · Canonical · current}\par

Calendar Management projects schedules and correlates events with temporal commitments. It brings rosters, on-call windows, facility calendars and appointment milestones into operational views without replacing their authoritative sources. This distinguishes the temporal context used for coordination from the business decisions that establish the schedule.



\textbf{Overview}



Calendar Management projects schedules and correlates events with temporal commitments. It brings rosters, on-call windows, facility calendars and appointment milestones into operational views without replacing their authoritative sources. This distinguishes the temporal context used for coordination from the business decisions that establish the schedule.



\textbf{Intent}



Calendar Management projects schedules and correlates events with temporal commitments. It brings rosters, on-call windows, facility calendars and appointment milestones into operational views without replacing their authoritative sources. This distinguishes the temporal context used for coordination from the business decisions that establish the schedule.



\textbf{Boundary}



Information administration does not confer authority over the represented person, organisation, service, place, device or clinical activity.



\medskip\hrule\medskip

{\footnotesize \href{../../../../continuum/strategy/capability/7669036f7e6749219412fc64df5f7064.md}{Canonical record: 7669036f7e6749219412fc64df5f7064}; full qualifications and provenance.}\par
\paragraph*{II.1.3.9.1.1 [CT3] Operational \& Clinical Calendar Coordination}
\phantomsection
\label{navigation:II.1.3.9.1.1}
\addcontentsline{toc}{paragraph}{II.1.3.9.1.1 [CT3] Operational \& Clinical Calendar Coordination}

{\small Historical predecessor: R.2.167 \hyperref[element:c0f6e5f7f5a94d1580fa6e0e8b82c93a]{[CT3] Operational \& Clinical Calendar Coordination}; Canonical ID: \texttt{\seqsplit{c0f6e5f7f5a94d1580fa6e0e8b82c93a}}. Superseded historical identity — direct correspondence · superseded.}\par
\label{element:f6d8f1310ba543f69e8b57bf05603db3}
\index{[CT3] Operational \& Clinical Calendar Coordination (II.1.3.9.1.1)}
{\small Capability · Proposed · current}\par

Coordinate projections of operational and clinical calendars and correlate temporal events without taking ownership of source bookings or rosters.



\textbf{Overview}



Coordinate projections of operational and clinical calendars and correlate temporal events without taking ownership of source bookings or rosters. The reviewed hierarchy assigns this enduring ability to its stated business-enabling context.



\textbf{Intent}



Coordinate projections of operational and clinical calendars and correlate temporal events without taking ownership of source bookings or rosters.



\textbf{Boundary}



Information administration and coordination preserve the authority of the responsible source, clinical actor and governed domain; they confer no additional real-world management authority.


{\footnotesize \href{../../../../continuum/strategy/capability/f6d8f1310ba543f69e8b57bf05603db3.md}{Canonical record: f6d8f1310ba543f69e8b57bf05603db3}; full qualifications and provenance.}\par
\subparagraph*{II.1.3.9.1.1.1 [Feature] Schedule Projection Integration}
\phantomsection
\label{navigation:II.1.3.9.1.1.1}
\addcontentsline{toc}{subparagraph}{II.1.3.9.1.1.1 [Feature] Schedule Projection Integration}

\label{element:883a68ab282b462988d5ae7a6ea2e980}
\index{[Feature] Schedule Projection Integration (II.1.3.9.1.1.1)}
{\small Capability / Feature · Canonical · current}\par

Project clinician rosters, on-call schedules, and facility operating calendars into operational views.



Schedule projection makes external roster, on-call and facility calendars available in operational views. It preserves their temporal usefulness without making the projection the authoritative schedule or granting authority to alter staffing and coverage decisions.


{\footnotesize \href{../../../../continuum/strategy/capability/feature/883a68ab282b462988d5ae7a6ea2e980.md}{Canonical record: 883a68ab282b462988d5ae7a6ea2e980}; full qualifications and provenance.}\par
\subparagraph*{II.1.3.9.1.1.2 [Feature] Temporal Event Correlation}
\phantomsection
\label{navigation:II.1.3.9.1.1.2}
\addcontentsline{toc}{subparagraph}{II.1.3.9.1.1.2 [Feature] Temporal Event Correlation}

\label{element:05e1c23baf1c47fbac8457f991efa2be}
\index{[Feature] Temporal Event Correlation (II.1.3.9.1.1.2)}
{\small Capability / Feature · Canonical · current}\par

Correlate incoming clinical events with scheduled appointments and operational calendar milestones.



Temporal correlation relates incoming events to the relevant appointment or calendar milestone. Association provides context for interpretation and coordination without assuming that temporal proximity proves attendance, fulfilment or the clinical outcome expected at that time.


{\footnotesize \href{../../../../continuum/strategy/capability/feature/05e1c23baf1c47fbac8457f991efa2be.md}{Canonical record: 05e1c23baf1c47fbac8457f991efa2be}; full qualifications and provenance.}\par
\subsubsection*{II.1.3.9.2 [CT2] Clinical Collaboration}
\phantomsection
\label{navigation:II.1.3.9.2}
\addcontentsline{toc}{subsubsection}{II.1.3.9.2 [CT2] Clinical Collaboration}

{\small Historical predecessor: R.2.170 \hyperref[element:c5adc1955b6046d8b920b0176cebc9de]{[CT2] Clinical Collaboration}; Canonical ID: \texttt{\seqsplit{c5adc1955b6046d8b920b0176cebc9de}}. Superseded historical identity — direct correspondence · superseded.}\par
\label{element:e6429225463a44d39f82e490e254e5d4}
\index{[CT2] Clinical Collaboration (II.1.3.9.2)}
{\small Capability · Canonical · current}\par

Clinical Collaboration provides secure multidisciplinary spaces bound to clinical context, including governed access to longitudinal information during discussion. Its purpose is contextual conversation between care participants rather than a separate source of authoritative clinical facts. A statement in a conversation remains discourse until explicitly authored through the governed record boundary.



\textbf{Overview}



Clinical Collaboration provides secure multidisciplinary spaces bound to clinical context, including governed access to longitudinal information during discussion. Its purpose is contextual conversation between care participants rather than a separate source of authoritative clinical facts. A statement in a conversation remains discourse until explicitly authored through the governed record boundary.



\textbf{Intent}



Clinical Collaboration provides secure multidisciplinary spaces bound to clinical context, including governed access to longitudinal information during discussion. Its purpose is contextual conversation between care participants rather than a separate source of authoritative clinical facts. A statement in a conversation remains discourse until explicitly authored through the governed record boundary.



\textbf{Boundary}



Information administration does not confer authority over the represented person, organisation, service, place, device or clinical activity.



\medskip\hrule\medskip

{\footnotesize \href{../../../../continuum/strategy/capability/e6429225463a44d39f82e490e254e5d4.md}{Canonical record: e6429225463a44d39f82e490e254e5d4}; full qualifications and provenance.}\par
\paragraph*{II.1.3.9.2.1 [CT3] Clinical Context Collaboration}
\phantomsection
\label{navigation:II.1.3.9.2.1}
\addcontentsline{toc}{paragraph}{II.1.3.9.2.1 [CT3] Clinical Context Collaboration}

{\small Historical predecessor: R.2.169 \hyperref[element:750ef0d8e1504a8e9a1fde9bf13a6d0c]{[CT3] Clinical Context Collaboration}; Canonical ID: \texttt{\seqsplit{750ef0d8e1504a8e9a1fde9bf13a6d0c}}. Superseded historical identity — direct correspondence · superseded.}\par
\label{element:f11617eed6ff4d6fa237644ec95b28b2}
\index{[CT3] Clinical Context Collaboration (II.1.3.9.2.1)}
{\small Capability · Proposed · current}\par

Enable governed multidisciplinary collaboration with clinical context while preserving the distinction between discourse and authored clinical records.



\textbf{Overview}



Enable governed multidisciplinary collaboration with clinical context while preserving the distinction between discourse and authored clinical records. The reviewed hierarchy assigns this enduring ability to its stated business-enabling context.



\textbf{Intent}



Enable governed multidisciplinary collaboration with clinical context while preserving the distinction between discourse and authored clinical records.



\textbf{Boundary}



Information administration and coordination preserve the authority of the responsible source, clinical actor and governed domain; they confer no additional real-world management authority.


{\footnotesize \href{../../../../continuum/strategy/capability/f11617eed6ff4d6fa237644ec95b28b2.md}{Canonical record: f11617eed6ff4d6fa237644ec95b28b2}; full qualifications and provenance.}\par
\subparagraph*{II.1.3.9.2.1.1 [Feature] Care-Team Clinical Collaboration Coordination}
\phantomsection
\label{navigation:II.1.3.9.2.1.1}
\addcontentsline{toc}{subparagraph}{II.1.3.9.2.1.1 [Feature] Care-Team Clinical Collaboration Coordination}

\label{element:b0860d6b68f34b8cb2f93b25cdf2aa97}
\index{[Feature] Care-Team Clinical Collaboration Coordination (II.1.3.9.2.1.1)}
{\small Capability / Feature · Canonical · current}\par

Coordinate secure multidisciplinary clinical communication and consultation exchange bound to patient encounters or clinical topics.



Clinical collaboration coordination binds secure multidisciplinary exchange to the relevant encounter or clinical topic. It supports conversation among authorised participants without making discourse an authoritative observation or acquiring the clinical handover authority excluded by the metamodel.


{\footnotesize \href{../../../../continuum/strategy/capability/feature/b0860d6b68f34b8cb2f93b25cdf2aa97.md}{Canonical record: b0860d6b68f34b8cb2f93b25cdf2aa97}; full qualifications and provenance.}\par
\subparagraph*{II.1.3.9.2.1.2 [Feature] Collaboration Interaction Capture Qualification}
\phantomsection
\label{navigation:II.1.3.9.2.1.2}
\addcontentsline{toc}{subparagraph}{II.1.3.9.2.1.2 [Feature] Collaboration Interaction Capture Qualification}

\label{element:8e2a760f15bd454c8aba1526ee69740c}
\index{[Feature] Collaboration Interaction Capture Qualification (II.1.3.9.2.1.2)}
{\small Capability / Feature · Proposed · current}\par

Qualify when a defined collaboration interaction creates a new governed information or activity consequence.



\textbf{Overview}



A message may convey existing information or an explicit interaction may derive a new managed consequence. The distinction preserves collaboration-native records while routing governed consequences to their accountable responsibility.



\textbf{Boundary}



Conversation presence, display or read receipt does not itself make a clinical decision, approval or final record.


{\footnotesize \href{../../../../continuum/strategy/capability/feature/8e2a760f15bd454c8aba1526ee69740c.md}{Canonical record: 8e2a760f15bd454c8aba1526ee69740c}; full qualifications and provenance.}\par
\subparagraph*{II.1.3.9.2.1.3 [Feature] In-Conversation LHR Query Resolution}
\phantomsection
\label{navigation:II.1.3.9.2.1.3}
\addcontentsline{toc}{subparagraph}{II.1.3.9.2.1.3 [Feature] In-Conversation LHR Query Resolution}

\label{element:d4437b477aab42c79b5f45150c5d7f3f}
\index{[Feature] In-Conversation LHR Query Resolution (II.1.3.9.2.1.3)}
{\small Capability / Feature · Canonical · current}\par

Resolve clinical data queries within an authorized collaboration context, presenting verified clinical summaries directly to care team members.



In-conversation query supplies governed clinical summaries within an authorised collaboration context. The convenience of retrieving information during discussion does not weaken consent or access controls, and displaying a summary does not promote the surrounding conversation into the clinical record.


{\footnotesize \href{../../../../continuum/strategy/capability/feature/d4437b477aab42c79b5f45150c5d7f3f.md}{Canonical record: d4437b477aab42c79b5f45150c5d7f3f}; full qualifications and provenance.}\par
\subparagraph*{II.1.3.9.2.1.4 [Feature] Zero-PHI Collaboration Metadata Governance}
\phantomsection
\label{navigation:II.1.3.9.2.1.4}
\addcontentsline{toc}{subparagraph}{II.1.3.9.2.1.4 [Feature] Zero-PHI Collaboration Metadata Governance}

\label{element:87e9246c944a40d9af5c4d0a3a7d4e9a}
\index{[Feature] Zero-PHI Collaboration Metadata Governance (II.1.3.9.2.1.4)}
{\small Capability / Feature · Canonical · current}\par

Ensure collaboration channel titles, topics, and aliases use opaque identifiers free of patient names, MRNs, or sensitive details.



Metadata governance keeps patient names, identifiers and sensitive detail out of collaboration titles, topics and aliases. Opaque references preserve useful contextual association without placing PHI in broadly visible channel metadata or diagnostic surfaces.


{\footnotesize \href{../../../../continuum/strategy/capability/feature/87e9246c944a40d9af5c4d0a3a7d4e9a.md}{Canonical record: 87e9246c944a40d9af5c4d0a3a7d4e9a}; full qualifications and provenance.}\par
\subsubsection*{II.1.3.9.3 [CT2] Clinical Information Stewardship Enablement [superseded]}
\phantomsection
\label{navigation:II.1.3.9.3}
\addcontentsline{toc}{subsubsection}{II.1.3.9.3 [CT2] Clinical Information Stewardship Enablement [superseded]}

\label{element:3811c758ff2140c39a801bf89f75098a}
\index{[CT2] Clinical Information Stewardship Enablement (II.1.3.9.3)}
{\small Capability · Proposed · superseded}\par
\href{../../../../continuum/strategy/capability/3811c758ff2140c39a801bf89f75098a.md}{Historical canonical definition and provenance}.\par
\paragraph*{II.1.3.9.3.1 [CT3] Information Authority Representation}
\phantomsection
\label{navigation:II.1.3.9.3.1}
\addcontentsline{toc}{paragraph}{II.1.3.9.3.1 [CT3] Information Authority Representation}

\label{element:8c850faafb334f97a81a0b4cd8fdb65f}
\index{[CT3] Information Authority Representation (II.1.3.9.3.1)}
{\small Capability · Proposed · pending\_adjudication}\par

Represent the source, credibility and qualified authority standing of competing health-information assertions.



\textbf{Overview}



A persisted assertion may be authoritative, informational or anecdotal. This ability carries that standing into interpretation without treating transport reliability, freshness or successful commit as proof of credibility.



\textbf{Intent}



Keep an assertion technical acceptance separate from the standing of the source that asserted it.



\textbf{Boundary}



Classification does not originate clinical facts or universally adjudicate conflicts. Policy determines how qualified assertions are used.


{\small Recorded unresolved matters: source-qualified matters; exact scope remains in the canonical record.}\par
{\footnotesize \href{../../../../continuum/strategy/capability/8c850faafb334f97a81a0b4cd8fdb65f.md}{Canonical record: 8c850faafb334f97a81a0b4cd8fdb65f}; full qualifications and provenance.}\par
\subparagraph*{II.1.3.9.3.1.1 [Feature] Information Authority \& Source Standing Representation}
\phantomsection
\label{navigation:II.1.3.9.3.1.1}
\addcontentsline{toc}{subparagraph}{II.1.3.9.3.1.1 [Feature] Information Authority \& Source Standing Representation}

\label{element:70a444156cd246cb97608a782e4bdea8}
\index{[Feature] Information Authority \& Source Standing Representation (II.1.3.9.3.1.1)}
{\small Capability / Feature · Proposed · pending\_adjudication}\par

Attach and preserve qualified source-authority and credibility standing for managed assertions.



\textbf{Overview}



The user of an assertion needs to distinguish authoritative, informational and anecdotal standing even when each was technically valid and durably accepted.



\textbf{Boundary}



Authority classification is not a clinical truth judgement or a universal rule for choosing competing assertions.


{\small Recorded unresolved matters: source-qualified matters; exact scope remains in the canonical record.}\par
{\footnotesize \href{../../../../continuum/strategy/capability/feature/70a444156cd246cb97608a782e4bdea8.md}{Canonical record: 70a444156cd246cb97608a782e4bdea8}; full qualifications and provenance.}\par
\subsubsection*{II.1.3.9.4 [CT2] Clinical Knowledge Services}
\phantomsection
\label{navigation:II.1.3.9.4}
\addcontentsline{toc}{subsubsection}{II.1.3.9.4 [CT2] Clinical Knowledge Services}

{\small Historical predecessor: R.2.174 \hyperref[element:5b8478b08c3f42c0a3480edb8f0b9ea7]{[CT2] Clinical Knowledge Services}; Canonical ID: \texttt{\seqsplit{5b8478b08c3f42c0a3480edb8f0b9ea7}}. Superseded historical identity — direct correspondence · superseded.}\par
\label{element:bbd2bd21859f40e39963bdbe7a851686}
\index{[CT2] Clinical Knowledge Services (II.1.3.9.4)}
{\small Capability · Canonical · current}\par

Administer clinical knowledge requests, distribute governed knowledge artefacts and administer approved reference datasets.



\textbf{Overview}



Administer clinical knowledge requests, distribute governed knowledge artefacts and administer approved reference datasets.



\textbf{Intent}



Administer clinical knowledge requests, distribute governed knowledge artefacts and administer approved reference datasets.



\textbf{Boundary}



Information administration does not confer authority over the represented person, organisation, service, place, device or clinical activity.



\medskip\hrule\medskip

{\footnotesize \href{../../../../continuum/strategy/capability/bbd2bd21859f40e39963bdbe7a851686.md}{Canonical record: bbd2bd21859f40e39963bdbe7a851686}; full qualifications and provenance.}\par
\paragraph*{II.1.3.9.4.1 [CT3] Clinical Knowledge Artefact Distribution Management}
\phantomsection
\label{navigation:II.1.3.9.4.1}
\addcontentsline{toc}{paragraph}{II.1.3.9.4.1 [CT3] Clinical Knowledge Artefact Distribution Management}

{\small Historical predecessor: R.2.171 \hyperref[element:d81e0d6d76014066a682e3d9e03e5f8a]{[CT3] Clinical Knowledge Artefact Distribution Management}; Canonical ID: \texttt{\seqsplit{d81e0d6d76014066a682e3d9e03e5f8a}}. Superseded historical identity — direct correspondence · superseded.}\par
\label{element:cd29a9eb798949bc8bc81bb5d434cabf}
\index{[CT3] Clinical Knowledge Artefact Distribution Management (II.1.3.9.4.1)}
{\small Capability · Proposed · current}\par

Manage distribution of governed clinical knowledge artefacts to intended consumers under their applicable use contracts.



\textbf{Overview}



Manage distribution of governed clinical knowledge artefacts to intended consumers under their applicable use contracts. The reviewed hierarchy assigns this enduring ability to its stated business-enabling context.



\textbf{Intent}



Manage distribution of governed clinical knowledge artefacts to intended consumers under their applicable use contracts.



\textbf{Boundary}



Information administration and coordination preserve the authority of the responsible source, clinical actor and governed domain; they confer no additional real-world management authority.


{\footnotesize \href{../../../../continuum/strategy/capability/cd29a9eb798949bc8bc81bb5d434cabf.md}{Canonical record: cd29a9eb798949bc8bc81bb5d434cabf}; full qualifications and provenance.}\par
\paragraph*{II.1.3.9.4.2 [CT3] Clinical Knowledge Request Administration}
\phantomsection
\label{navigation:II.1.3.9.4.2}
\addcontentsline{toc}{paragraph}{II.1.3.9.4.2 [CT3] Clinical Knowledge Request Administration}

{\small Historical predecessor: R.2.172 \hyperref[element:eb370c01fbdb4d22956d97e9ecd74b57]{[CT3] Clinical Knowledge Request Administration}; Canonical ID: \texttt{\seqsplit{eb370c01fbdb4d22956d97e9ecd74b57}}. Superseded historical identity — direct correspondence · superseded.}\par
\label{element:543c53483a2249d78044c6635b7b7c48}
\index{[CT3] Clinical Knowledge Request Administration (II.1.3.9.4.2)}
{\small Capability · Proposed · current}\par

Administer clinical knowledge requests and their qualified resolution to applicable governed knowledge or reference content.



\textbf{Overview}



Administer clinical knowledge requests and their qualified resolution to applicable governed knowledge or reference content. The reviewed hierarchy assigns this enduring ability to its stated business-enabling context.



\textbf{Intent}



Administer clinical knowledge requests and their qualified resolution to applicable governed knowledge or reference content.



\textbf{Boundary}



Information administration and coordination preserve the authority of the responsible source, clinical actor and governed domain; they confer no additional real-world management authority.


{\footnotesize \href{../../../../continuum/strategy/capability/543c53483a2249d78044c6635b7b7c48.md}{Canonical record: 543c53483a2249d78044c6635b7b7c48}; full qualifications and provenance.}\par
\subparagraph*{II.1.3.9.4.2.1 [Feature] Bidirectional Terminology Mapping}
\phantomsection
\label{navigation:II.1.3.9.4.2.1}
\addcontentsline{toc}{subparagraph}{II.1.3.9.4.2.1 [Feature] Bidirectional Terminology Mapping}

\label{element:6134dc91af8d45e8b81cc28d18a5251f}
\index{[Feature] Bidirectional Terminology Mapping (II.1.3.9.4.2.1)}
{\small Capability / Feature · Canonical · current}\par

Translate between legacy local code systems and canonical national healthcare profiles.



Terminology mapping applies governed correspondence between local and canonical semantic systems. Translation preserves the limitations and authority of the map; it does not justify silently treating related source concepts as identical or clinically adjudicating the underlying information.


{\footnotesize \href{../../../../continuum/strategy/capability/feature/6134dc91af8d45e8b81cc28d18a5251f.md}{Canonical record: 6134dc91af8d45e8b81cc28d18a5251f}; full qualifications and provenance.}\par
\subparagraph*{II.1.3.9.4.2.2 [Feature] Clinical Concept Resolution}
\phantomsection
\label{navigation:II.1.3.9.4.2.2}
\addcontentsline{toc}{subparagraph}{II.1.3.9.4.2.2 [Feature] Clinical Concept Resolution}

\label{element:30cdb9b2a810467c9b56e9870de4367e}
\index{[Feature] Clinical Concept Resolution (II.1.3.9.4.2.2)}
{\small Capability / Feature · Canonical · current}\par

Resolve local clinical codes against standard terminologies (SNOMED CT, LOINC, ICD-10).



Concept resolution makes a coded clinical reference interpretable against its governed terminology. It validates the reference context without making a recognised code proof of clinical truth or giving the consuming system authority to redefine the concept.


{\footnotesize \href{../../../../continuum/strategy/capability/feature/30cdb9b2a810467c9b56e9870de4367e.md}{Canonical record: 30cdb9b2a810467c9b56e9870de4367e}; full qualifications and provenance.}\par
\paragraph*{II.1.3.9.4.3 [CT3] Governed Reference Dataset Administration}
\phantomsection
\label{navigation:II.1.3.9.4.3}
\addcontentsline{toc}{paragraph}{II.1.3.9.4.3 [CT3] Governed Reference Dataset Administration}

{\small Historical predecessor: R.2.173 \hyperref[element:15d972931ed6447797c30a5ddb71737e]{[CT3] Governed Reference Dataset Administration}; Canonical ID: \texttt{\seqsplit{15d972931ed6447797c30a5ddb71737e}}. Superseded historical identity — direct correspondence · superseded.}\par
\label{element:682b8684ad9e49d5a5c348a148081607}
\index{[CT3] Governed Reference Dataset Administration (II.1.3.9.4.3)}
{\small Capability · Proposed · current}\par

Administer approved runtime reference-dataset content and revisions separately from the reusable models defining its meaning.



\textbf{Overview}



Administer approved runtime reference-dataset content and revisions separately from the reusable models defining its meaning.



\textbf{Intent}



Administer approved runtime reference-dataset content and revisions separately from the reusable models defining its meaning.



\textbf{Boundary}



Does not own external ontology authority, execute authorisation policy or establish durable truth independently of the normal persistence boundary.


{\footnotesize \href{../../../../continuum/strategy/capability/682b8684ad9e49d5a5c348a148081607.md}{Canonical record: 682b8684ad9e49d5a5c348a148081607}; full qualifications and provenance.}\par
\subparagraph*{II.1.3.9.4.3.1 [Feature] Governed Dataset Content Revision}
\phantomsection
\label{navigation:II.1.3.9.4.3.1}
\addcontentsline{toc}{subparagraph}{II.1.3.9.4.3.1 [Feature] Governed Dataset Content Revision}

\label{element:89123c868b8a46c38800916110cd918d}
\index{[Feature] Governed Dataset Content Revision (II.1.3.9.4.3.1)}
{\small Capability / Feature · Proposed · current}\par

Review, accept and expose authorised changes to mutable reference content with effective revision standing.



\textbf{Overview}



A model definition can remain stable while approved terminology membership or mapping content changes. This Feature represents that operational governance without making every definition runtime mutable.



\textbf{Boundary}



Does not invent an approval authority or bypass normal durable establishment.


{\footnotesize \href{../../../../continuum/strategy/capability/feature/89123c868b8a46c38800916110cd918d.md}{Canonical record: 89123c868b8a46c38800916110cd918d}; full qualifications and provenance.}\par
\subsubsection*{II.1.3.9.5 [CT2] General Collaboration}
\phantomsection
\label{navigation:II.1.3.9.5}
\addcontentsline{toc}{subsubsection}{II.1.3.9.5 [CT2] General Collaboration}

{\small Historical predecessor: R.2.176 \hyperref[element:4f58c82aca7a4b54aeb6e2920ce5c12a]{[CT2] General Collaboration}; Canonical ID: \texttt{\seqsplit{4f58c82aca7a4b54aeb6e2920ce5c12a}}. Superseded historical identity — direct correspondence · superseded.}\par
\label{element:3bb740e6db284112b054905c0c6c80be}
\index{[CT2] General Collaboration (II.1.3.9.5)}
{\small Capability · Canonical · current}\par

General Collaboration provides secure operational messaging for departmental, administrative and facility-support participants. It gives non-clinical coordination an interaction context without assuming every discussion requires a patient-linked clinical workspace. Its role is communication enablement rather than ownership of the support activity being discussed.



\textbf{Overview}



General Collaboration provides secure operational messaging for departmental, administrative and facility-support participants. It gives non-clinical coordination an interaction context without assuming every discussion requires a patient-linked clinical workspace. Its role is communication enablement rather than ownership of the support activity being discussed.



\textbf{Intent}



General Collaboration provides secure operational messaging for departmental, administrative and facility-support participants. It gives non-clinical coordination an interaction context without assuming every discussion requires a patient-linked clinical workspace. Its role is communication enablement rather than ownership of the support activity being discussed.



\textbf{Boundary}



Information administration does not confer authority over the represented person, organisation, service, place, device or clinical activity.



\medskip\hrule\medskip

{\footnotesize \href{../../../../continuum/strategy/capability/3bb740e6db284112b054905c0c6c80be.md}{Canonical record: 3bb740e6db284112b054905c0c6c80be}; full qualifications and provenance.}\par
\paragraph*{II.1.3.9.5.1 [CT3] Operational Collaboration}
\phantomsection
\label{navigation:II.1.3.9.5.1}
\addcontentsline{toc}{paragraph}{II.1.3.9.5.1 [CT3] Operational Collaboration}

{\small Historical predecessor: R.2.175 \hyperref[element:99e9d9763ec94140b7d7642f4b5966f6]{[CT3] Operational Collaboration}; Canonical ID: \texttt{\seqsplit{99e9d9763ec94140b7d7642f4b5966f6}}. Superseded historical identity — direct correspondence · superseded.}\par
\label{element:631fdb8b47084008ac9b5de28b26972b}
\index{[CT3] Operational Collaboration (II.1.3.9.5.1)}
{\small Capability · Proposed · current}\par

Enable operational team discussion and collaboration without promoting conversational content to clinical record authority.



\textbf{Overview}



Enable operational team discussion and collaboration without promoting conversational content to clinical record authority. The reviewed hierarchy assigns this enduring ability to its stated business-enabling context.



\textbf{Intent}



Enable operational team discussion and collaboration without promoting conversational content to clinical record authority.



\textbf{Boundary}



Information administration and coordination preserve the authority of the responsible source, clinical actor and governed domain; they confer no additional real-world management authority.


{\footnotesize \href{../../../../continuum/strategy/capability/631fdb8b47084008ac9b5de28b26972b.md}{Canonical record: 631fdb8b47084008ac9b5de28b26972b}; full qualifications and provenance.}\par
\subparagraph*{II.1.3.9.5.1.1 [Feature] Operational Team Chat Enablement}
\phantomsection
\label{navigation:II.1.3.9.5.1.1}
\addcontentsline{toc}{subparagraph}{II.1.3.9.5.1.1 [Feature] Operational Team Chat Enablement}

\label{element:367a390e7a464003beba189428c5faca}
\index{[Feature] Operational Team Chat Enablement (II.1.3.9.5.1.1)}
{\small Capability / Feature · Canonical · current}\par

Provide secure instant messaging channels for departmental, administrative, and facility support staff.



Operational chat provides secure communication for administrative and support teams. It enables participants to discuss work while leaving assignment, completion and the authority of operational decisions with the capabilities responsible for that activity.


{\footnotesize \href{../../../../continuum/strategy/capability/feature/367a390e7a464003beba189428c5faca.md}{Canonical record: 367a390e7a464003beba189428c5faca}; full qualifications and provenance.}\par
\subsubsection*{II.1.3.9.6 [CT2] Health Information Access}
\phantomsection
\label{navigation:II.1.3.9.6}
\addcontentsline{toc}{subsubsection}{II.1.3.9.6 [CT2] Health Information Access}

{\small Historical predecessor: R.2.178 \hyperref[element:5019ae9f2cdc45a1b2de97bc0f5a313d]{[CT2] Health Information Access}; Canonical ID: \texttt{\seqsplit{5019ae9f2cdc45a1b2de97bc0f5a313d}}. Superseded historical identity — direct correspondence · superseded.}\par
\label{element:cc253aea3c5f489da0502df74dde1669}
\index{[CT2] Health Information Access (II.1.3.9.6)}
{\small Capability · Canonical · current}\par

Health Information Access enables authorised discovery and retrieval across clinical information. It combines semantic search and discrete retrieval with access-qualified filtering so that a discoverable item is not automatically disclosable. It consumes established policy and consent determinations rather than making the search interface an alternative security authority.



\textbf{Overview}



Health Information Access enables authorised discovery and retrieval across clinical information. It combines semantic search and discrete retrieval with access-qualified filtering so that a discoverable item is not automatically disclosable. It consumes established policy and consent determinations rather than making the search interface an alternative security authority.



\textbf{Intent}



Health Information Access enables authorised discovery and retrieval across clinical information. It combines semantic search and discrete retrieval with access-qualified filtering so that a discoverable item is not automatically disclosable. It consumes established policy and consent determinations rather than making the search interface an alternative security authority.



\textbf{Boundary}



Information administration does not confer authority over the represented person, organisation, service, place, device or clinical activity.



\medskip\hrule\medskip

{\footnotesize \href{../../../../continuum/strategy/capability/cc253aea3c5f489da0502df74dde1669.md}{Canonical record: cc253aea3c5f489da0502df74dde1669}; full qualifications and provenance.}\par
\paragraph*{II.1.3.9.6.1 [CT3] Governed Health Information Retrieval}
\phantomsection
\label{navigation:II.1.3.9.6.1}
\addcontentsline{toc}{paragraph}{II.1.3.9.6.1 [CT3] Governed Health Information Retrieval}

{\small Historical predecessor: R.2.177 \hyperref[element:90b838d703dd472b8b6d74961f3a748c]{[CT3] Governed Health Information Retrieval}; Canonical ID: \texttt{\seqsplit{90b838d703dd472b8b6d74961f3a748c}}. Superseded historical identity — direct correspondence · superseded.}\par
\label{element:76e74e46ca4d44a1a7bb116b358f231c}
\index{[CT3] Governed Health Information Retrieval (II.1.3.9.6.1)}
{\small Capability · Proposed · current}\par

Retrieve and search health information and apply access-qualified result filtering under established access and consent determinations.



\textbf{Overview}



Retrieve and search health information and apply access-qualified result filtering under established access and consent determinations. The reviewed hierarchy assigns this enduring ability to its stated business-enabling context.



\textbf{Intent}



Retrieve and search health information and apply access-qualified result filtering under established access and consent determinations.



\textbf{Boundary}



Information administration and coordination preserve the authority of the responsible source, clinical actor and governed domain; they confer no additional real-world management authority.


{\footnotesize \href{../../../../continuum/strategy/capability/76e74e46ca4d44a1a7bb116b358f231c.md}{Canonical record: 76e74e46ca4d44a1a7bb116b358f231c}; full qualifications and provenance.}\par
\subparagraph*{II.1.3.9.6.1.1 [Feature] Access-Controlled Information Filtering}
\phantomsection
\label{navigation:II.1.3.9.6.1.1}
\addcontentsline{toc}{subparagraph}{II.1.3.9.6.1.1 [Feature] Access-Controlled Information Filtering}

\label{element:a9f1d2100faa44beb69add62bb4f4b3f}
\index{[Feature] Access-Controlled Information Filtering (II.1.3.9.6.1.1)}
{\small Capability / Feature · Canonical · current}\par

Filter search results dynamically according to requesting user authorizations and patient consent directives.



Access-qualified filtering restricts the disclosed result to the verified caller's authority and applicable consent. The responsibility is more specific than retrieving a resource; the latter's shared owner and subject matter do not establish a Feature realisation relationship.


{\footnotesize \href{../../../../continuum/strategy/capability/feature/a9f1d2100faa44beb69add62bb4f4b3f.md}{Canonical record: a9f1d2100faa44beb69add62bb4f4b3f}; full qualifications and provenance.}\par
\subparagraph*{II.1.3.9.6.1.2 [Feature] Federated Clinical Query}
\phantomsection
\label{navigation:II.1.3.9.6.1.2}
\addcontentsline{toc}{subparagraph}{II.1.3.9.6.1.2 [Feature] Federated Clinical Query}

\label{element:cfcc4d64af434418a567ccf46d8052dd}
\index{[Feature] Federated Clinical Query (II.1.3.9.6.1.2)}
{\small Capability / Feature · Canonical · current}\par

Execute bounded, indexed queries across distributed clinical resources and directory graphs.



Federated query makes relevant clinical and directory information discoverable across distributed contexts using bounded criteria. Discovery is a governed information-access responsibility rather than authority to disclose every match or reinterpret the clinical meaning of the results.


{\footnotesize \href{../../../../continuum/strategy/capability/feature/cfcc4d64af434418a567ccf46d8052dd.md}{Canonical record: cfcc4d64af434418a567ccf46d8052dd}; full qualifications and provenance.}\par
\subsubsection*{II.1.3.9.7 [CT2] Health Information Communication}
\phantomsection
\label{navigation:II.1.3.9.7}
\addcontentsline{toc}{subsubsection}{II.1.3.9.7 [CT2] Health Information Communication}

{\small Historical predecessor: R.2.183 \hyperref[element:bcefb378424842fcb6d703208c3f8773]{[CT2] Health Information Communication}; Canonical ID: \texttt{\seqsplit{bcefb378424842fcb6d703208c3f8773}}. Superseded historical identity — direct correspondence · superseded.}\par
\label{element:7925605bf60f4b288e826020409a09d9}
\index{[CT2] Health Information Communication (II.1.3.9.7)}
{\small Capability · Proposed · current}\par

Enable patient notifications, patient task communication, patient self-service and access to a patient health library under applicable information-use controls.



\textbf{Overview}



Enable patient notifications, patient task communication, patient self-service and access to a patient health library under applicable information-use controls. The reviewed hierarchy assigns this enduring ability to its stated business-enabling context.



\textbf{Intent}



Enable patient notifications, patient task communication, patient self-service and access to a patient health library under applicable information-use controls.



\textbf{Boundary}



Information administration and coordination preserve the authority of the responsible source, clinical actor and governed domain; they confer no additional real-world management authority.


{\footnotesize \href{../../../../continuum/strategy/capability/7925605bf60f4b288e826020409a09d9.md}{Canonical record: 7925605bf60f4b288e826020409a09d9}; full qualifications and provenance.}\par
\paragraph*{II.1.3.9.7.1 [CT3] Patient Health Library Access}
\phantomsection
\label{navigation:II.1.3.9.7.1}
\addcontentsline{toc}{paragraph}{II.1.3.9.7.1 [CT3] Patient Health Library Access}

{\small Historical predecessor: R.2.179 \hyperref[element:b132224cfa6f435f8f2a66a3878d6df6]{[CT3] Patient Health Library Access}; Canonical ID: \texttt{\seqsplit{b132224cfa6f435f8f2a66a3878d6df6}}. Superseded historical identity — direct correspondence · superseded.}\par
\label{element:808e5a3c3f6b496fa329962088f0bee0}
\index{[CT3] Patient Health Library Access (II.1.3.9.7.1)}
{\small Capability · Proposed · current}\par

Enable authorised patient access to appropriately governed health knowledge and library information.



\textbf{Overview}



Enable authorised patient access to appropriately governed health knowledge and library information. The reviewed hierarchy assigns this enduring ability to its stated business-enabling context.



\textbf{Intent}



Enable authorised patient access to appropriately governed health knowledge and library information.



\textbf{Boundary}



Information administration and coordination preserve the authority of the responsible source, clinical actor and governed domain; they confer no additional real-world management authority.


{\footnotesize \href{../../../../continuum/strategy/capability/808e5a3c3f6b496fa329962088f0bee0.md}{Canonical record: 808e5a3c3f6b496fa329962088f0bee0}; full qualifications and provenance.}\par
\paragraph*{II.1.3.9.7.2 [CT3] Patient Notification Services}
\phantomsection
\label{navigation:II.1.3.9.7.2}
\addcontentsline{toc}{paragraph}{II.1.3.9.7.2 [CT3] Patient Notification Services}

{\small Historical predecessor: R.2.180 \hyperref[element:35c5be472bf2433fbbde9646d87aefb3]{[CT3] Patient Notification Services}; Canonical ID: \texttt{\seqsplit{35c5be472bf2433fbbde9646d87aefb3}}. Superseded historical identity — direct correspondence · superseded.}\par
\label{element:9d172cd511b742499c4e7c23f6630a35}
\index{[CT3] Patient Notification Services (II.1.3.9.7.2)}
{\small Capability · Proposed · current}\par

Enable governed patient notifications concerning health information and applicable care or operational events.



\textbf{Overview}



Enable governed patient notifications concerning health information and applicable care or operational events. The reviewed hierarchy assigns this enduring ability to its stated business-enabling context.



\textbf{Intent}



Enable governed patient notifications concerning health information and applicable care or operational events.



\textbf{Boundary}



Information administration and coordination preserve the authority of the responsible source, clinical actor and governed domain; they confer no additional real-world management authority.


{\footnotesize \href{../../../../continuum/strategy/capability/9d172cd511b742499c4e7c23f6630a35.md}{Canonical record: 9d172cd511b742499c4e7c23f6630a35}; full qualifications and provenance.}\par
\paragraph*{II.1.3.9.7.3 [CT3] Patient Self-Service}
\phantomsection
\label{navigation:II.1.3.9.7.3}
\addcontentsline{toc}{paragraph}{II.1.3.9.7.3 [CT3] Patient Self-Service}

{\small Historical predecessor: R.2.181 \hyperref[element:60f427e6b5764757b909624085a2af89]{[CT3] Patient Self-Service}; Canonical ID: \texttt{\seqsplit{60f427e6b5764757b909624085a2af89}}. Superseded historical identity — direct correspondence · superseded.}\par
\label{element:d4956cab8b2b46c396899949c9486a1c}
\index{[CT3] Patient Self-Service (II.1.3.9.7.3)}
{\small Capability · Proposed · current}\par

Enable patients to undertake permitted information-access and administrative self-service interactions.



\textbf{Overview}



Enable patients to undertake permitted information-access and administrative self-service interactions. The reviewed hierarchy assigns this enduring ability to its stated business-enabling context.



\textbf{Intent}



Enable patients to undertake permitted information-access and administrative self-service interactions.



\textbf{Boundary}



Information administration and coordination preserve the authority of the responsible source, clinical actor and governed domain; they confer no additional real-world management authority.


{\footnotesize \href{../../../../continuum/strategy/capability/d4956cab8b2b46c396899949c9486a1c.md}{Canonical record: d4956cab8b2b46c396899949c9486a1c}; full qualifications and provenance.}\par
\paragraph*{II.1.3.9.7.4 [CT3] Patient Task Communication Services}
\phantomsection
\label{navigation:II.1.3.9.7.4}
\addcontentsline{toc}{paragraph}{II.1.3.9.7.4 [CT3] Patient Task Communication Services}

{\small Historical predecessor: R.2.182 \hyperref[element:0914f2bb267449e282e9ae0f7407d9fe]{[CT3] Patient Task Communication Services}; Canonical ID: \texttt{\seqsplit{0914f2bb267449e282e9ae0f7407d9fe}}. Superseded historical identity — direct correspondence · superseded.}\par
\label{element:ba8ce1d54ace4bc9aa3521dd954f36cb}
\index{[CT3] Patient Task Communication Services (II.1.3.9.7.4)}
{\small Capability · Proposed · current}\par

Communicate patient-facing work requests and their reported progression under the applicable work and disclosure contracts.



\textbf{Overview}



Communicate patient-facing work requests and their reported progression under the applicable work and disclosure contracts. The reviewed hierarchy assigns this enduring ability to its stated business-enabling context.



\textbf{Intent}



Communicate patient-facing work requests and their reported progression under the applicable work and disclosure contracts.



\textbf{Boundary}



Information administration and coordination preserve the authority of the responsible source, clinical actor and governed domain; they confer no additional real-world management authority.


{\footnotesize \href{../../../../continuum/strategy/capability/ba8ce1d54ace4bc9aa3521dd954f36cb.md}{Canonical record: ba8ce1d54ace4bc9aa3521dd954f36cb}; full qualifications and provenance.}\par
\subsubsection*{II.1.3.9.8 [CT2] Health Information Connectivity Enablement [superseded]}
\phantomsection
\label{navigation:II.1.3.9.8}
\addcontentsline{toc}{subsubsection}{II.1.3.9.8 [CT2] Health Information Connectivity Enablement [superseded]}

\label{element:6a1030483dfc4ae6af6998b587129d1c}
\index{[CT2] Health Information Connectivity Enablement (II.1.3.9.8)}
{\small Capability · Proposed · superseded}\par
\href{../../../../continuum/strategy/capability/6a1030483dfc4ae6af6998b587129d1c.md}{Historical canonical definition and provenance}.\par
\paragraph*{II.1.3.9.8.1 [CT3] Health Information Communication}
\phantomsection
\label{navigation:II.1.3.9.8.1}
\addcontentsline{toc}{paragraph}{II.1.3.9.8.1 [CT3] Health Information Communication}

\label{element:9794b2aff43443b495a19aefe0efa67b}
\index{[CT3] Health Information Communication (II.1.3.9.8.1)}
{\small Capability · Canonical · pending\_adjudication}\par

Definition unresolved in the baseline. No definition has been synthesised from the name.



\medskip\hrule\medskip

{\small Recorded unresolved matters: source-qualified matters; exact scope remains in the canonical record.}\par
{\footnotesize \href{../../../../continuum/strategy/capability/9794b2aff43443b495a19aefe0efa67b.md}{Canonical record: 9794b2aff43443b495a19aefe0efa67b}; full qualifications and provenance.}\par
\subparagraph*{II.1.3.9.8.1.1 [Feature] Standards-Based Boundary Exchange}
\phantomsection
\label{navigation:II.1.3.9.8.1.1}
\addcontentsline{toc}{subparagraph}{II.1.3.9.8.1.1 [Feature] Standards-Based Boundary Exchange}

\label{element:3192ba1541ff41458d176184119c30bc}
\index{[Feature] Standards-Based Boundary Exchange (II.1.3.9.8.1.1)}
{\small Capability / Feature · Canonical · pending\_adjudication}\par

Provide standards-conformant boundary exchange for healthcare messaging and resource representations.



\medskip\hrule\medskip

{\small Recorded unresolved matters: source-qualified matters; exact scope remains in the canonical record.}\par
{\footnotesize \href{../../../../continuum/strategy/capability/feature/3192ba1541ff41458d176184119c30bc.md}{Canonical record: 3192ba1541ff41458d176184119c30bc}; full qualifications and provenance.}\par
\subsubsection*{II.1.3.9.9 [CT2] Health Information Control}
\phantomsection
\label{navigation:II.1.3.9.9}
\addcontentsline{toc}{subsubsection}{II.1.3.9.9 [CT2] Health Information Control}

{\small Historical predecessor: R.2.187 \hyperref[element:a3196aef21d443519de874c24f64a185]{[CT2] Health Information Control}; Canonical ID: \texttt{\seqsplit{a3196aef21d443519de874c24f64a185}}. Superseded historical identity — direct correspondence · superseded.}\par
\label{element:44cc0e17ad534675bf26dce6e46a68ff}
\index{[CT2] Health Information Control (II.1.3.9.9)}
{\small Capability · Canonical · current}\par

Govern patient health-information restrictions, filtering and disclosure using established information-access authority, client privacy/consent determinations and applicable security controls.



\textbf{Overview}



Govern patient health-information restrictions, filtering and disclosure using established information-access authority, client privacy/consent determinations and applicable security controls.



\textbf{Intent}



Govern patient health-information restrictions, filtering and disclosure using established information-access authority, client privacy/consent determinations and applicable security controls.



\textbf{Boundary}



Information administration does not confer authority over the represented person, organisation, service, place, device or clinical activity. Client consent directives remain owned by Client Privacy Guardrail Management (FEAT-EM-09); this capability references and respects those determinations without defining a second consent owner.



\medskip\hrule\medskip

{\footnotesize \href{../../../../continuum/strategy/capability/44cc0e17ad534675bf26dce6e46a68ff.md}{Canonical record: 44cc0e17ad534675bf26dce6e46a68ff}; full qualifications and provenance.}\par
\paragraph*{II.1.3.9.9.1 [CT3] Patient Health Information Disclosure Management}
\phantomsection
\label{navigation:II.1.3.9.9.1}
\addcontentsline{toc}{paragraph}{II.1.3.9.9.1 [CT3] Patient Health Information Disclosure Management}

{\small Historical predecessor: R.2.184 \hyperref[element:07cf507f183d4905a7ad2adbb8555df5]{[CT3] Patient Health Information Disclosure Management}; Canonical ID: \texttt{\seqsplit{07cf507f183d4905a7ad2adbb8555df5}}. Superseded historical identity — direct correspondence · superseded.}\par
\label{element:7f918639e23842e78fa92f1bf47c5a7b}
\index{[CT3] Patient Health Information Disclosure Management (II.1.3.9.9.1)}
{\small Capability · Proposed · current}\par

Manage permitted patient health-information disclosure and its applicable recipient, purpose and restriction conditions.



\textbf{Overview}



Manage permitted patient health-information disclosure and its applicable recipient, purpose and restriction conditions. The reviewed hierarchy assigns this enduring ability to its stated business-enabling context.



\textbf{Intent}



Manage permitted patient health-information disclosure and its applicable recipient, purpose and restriction conditions.



\textbf{Boundary}



Information administration and coordination preserve the authority of the responsible source, clinical actor and governed domain; they confer no additional real-world management authority. Client consent directives remain owned by Client Privacy Guardrail Management (FEAT-EM-09); this capability references and respects those determinations without defining a second consent owner.


{\footnotesize \href{../../../../continuum/strategy/capability/7f918639e23842e78fa92f1bf47c5a7b.md}{Canonical record: 7f918639e23842e78fa92f1bf47c5a7b}; full qualifications and provenance.}\par
\paragraph*{II.1.3.9.9.2 [CT3] Patient Health Information Filter Management}
\phantomsection
\label{navigation:II.1.3.9.9.2}
\addcontentsline{toc}{paragraph}{II.1.3.9.9.2 [CT3] Patient Health Information Filter Management}

{\small Historical predecessor: R.2.185 \hyperref[element:589f192797454d9e8960fb63774e5e9c]{[CT3] Patient Health Information Filter Management}; Canonical ID: \texttt{\seqsplit{589f192797454d9e8960fb63774e5e9c}}. Superseded historical identity — direct correspondence · superseded.}\par
\label{element:62f673613c76494eba972cad3e4e8d0b}
\index{[CT3] Patient Health Information Filter Management (II.1.3.9.9.2)}
{\small Capability · Proposed · current}\par

Manage patient health-information filtering rules and qualified projections while respecting established client privacy and consent determinations.



\textbf{Overview}



Manage patient health-information filtering rules and qualified projections while respecting established client privacy and consent determinations. The reviewed hierarchy assigns this enduring ability to its stated business-enabling context.



\textbf{Intent}



Manage patient health-information filtering rules and qualified projections while respecting established client privacy and consent determinations.



\textbf{Boundary}



Information administration and coordination preserve the authority of the responsible source, clinical actor and governed domain; they confer no additional real-world management authority. Client consent directives remain owned by Client Privacy Guardrail Management (FEAT-EM-09); this capability references and respects those determinations without defining a second consent owner.


{\footnotesize \href{../../../../continuum/strategy/capability/62f673613c76494eba972cad3e4e8d0b.md}{Canonical record: 62f673613c76494eba972cad3e4e8d0b}; full qualifications and provenance.}\par
\paragraph*{II.1.3.9.9.3 [CT3] Patient Health Information Restriction Management}
\phantomsection
\label{navigation:II.1.3.9.9.3}
\addcontentsline{toc}{paragraph}{II.1.3.9.9.3 [CT3] Patient Health Information Restriction Management}

{\small Historical predecessor: R.2.186 \hyperref[element:5ad91d045474435a96b560e66ff7de05]{[CT3] Patient Health Information Restriction Management}; Canonical ID: \texttt{\seqsplit{5ad91d045474435a96b560e66ff7de05}}. Superseded historical identity — direct correspondence · superseded.}\par
\label{element:262366990bda4ac4ace5429269cf6a98}
\index{[CT3] Patient Health Information Restriction Management (II.1.3.9.9.3)}
{\small Capability · Proposed · current}\par

Manage application of patient health-information restrictions using the client-associated privacy and consent responsibilities retained by Client Privacy Guardrail Management.



\textbf{Overview}



Manage application of patient health-information restrictions using the client-associated privacy and consent responsibilities retained by Client Privacy Guardrail Management. The reviewed hierarchy assigns this enduring ability to its stated business-enabling context.



\textbf{Intent}



Manage application of patient health-information restrictions using the client-associated privacy and consent responsibilities retained by Client Privacy Guardrail Management.



\textbf{Boundary}



Information administration and coordination preserve the authority of the responsible source, clinical actor and governed domain; they confer no additional real-world management authority. Client consent directives remain owned by Client Privacy Guardrail Management (FEAT-EM-09); this capability references and respects those determinations without defining a second consent owner.


{\footnotesize \href{../../../../continuum/strategy/capability/262366990bda4ac4ace5429269cf6a98.md}{Canonical record: 262366990bda4ac4ace5429269cf6a98}; full qualifications and provenance.}\par
\subparagraph*{II.1.3.9.9.3.1 [Feature] Health Information Access \& Consent Control}
\phantomsection
\label{navigation:II.1.3.9.9.3.1}
\addcontentsline{toc}{subparagraph}{II.1.3.9.9.3.1 [Feature] Health Information Access \& Consent Control}

\label{element:e2643fa59303401abb1b778f3295e536}
\index{[Feature] Health Information Access \& Consent Control (II.1.3.9.9.3.1)}
{\small Capability / Feature · Canonical · current}\par

Evaluate patient consent preferences, practitioner authority boundaries, and access control policies for managed clinical records.



Access and consent control evaluates whether an information operation is legitimate in its practitioner, patient and purpose context. It supplies policy determinations consumed by access and presentation paths rather than allowing those consumers to establish authority from caller-supplied content.


{\footnotesize \href{../../../../continuum/strategy/capability/feature/e2643fa59303401abb1b778f3295e536.md}{Canonical record: e2643fa59303401abb1b778f3295e536}; full qualifications and provenance.}\par
\paragraph*{II.1.3.9.9.4 [Feature] Governed Policy Assignment \& Revocation}
\phantomsection
\label{navigation:II.1.3.9.9.4}
\addcontentsline{toc}{paragraph}{II.1.3.9.9.4 [Feature] Governed Policy Assignment \& Revocation}

\label{element:d6c271ae95ea44449c5e22665d1d9bc2}
\index{[Feature] Governed Policy Assignment \& Revocation (II.1.3.9.9.4)}
{\small Capability / Feature · Proposed · pending\_adjudication}\par

Manage approved permission criteria and identity/role assignments, including temporal applicability and revocation.



\textbf{Overview}



An assignment becoming inapplicable by expiry differs from an explicit revocation; both must affect the effective authority used for operations.



\textbf{Boundary}



Management does not grant unconstrained administrative authority or derive permission from resource tags.


{\small Recorded unresolved matters: source-qualified matters; exact scope remains in the canonical record.}\par
{\footnotesize \href{../../../../continuum/strategy/capability/feature/d6c271ae95ea44449c5e22665d1d9bc2.md}{Canonical record: d6c271ae95ea44449c5e22665d1d9bc2}; full qualifications and provenance.}\par
\paragraph*{II.1.3.9.9.5 [Feature] Non-PHI Compliance Auditing}
\phantomsection
\label{navigation:II.1.3.9.9.5}
\addcontentsline{toc}{paragraph}{II.1.3.9.9.5 [Feature] Non-PHI Compliance Auditing}

\label{element:08e37326908a49988dfb787f68df840a}
\index{[Feature] Non-PHI Compliance Auditing (II.1.3.9.9.5)}
{\small Capability / Feature · Canonical · pending\_adjudication}\par

Generate tamper-evident, verifiable audit records of security-significant events without emitting unmasked PHI.



\medskip\hrule\medskip

{\small Recorded unresolved matters: source-qualified matters; exact scope remains in the canonical record.}\par
{\footnotesize \href{../../../../continuum/strategy/capability/feature/08e37326908a49988dfb787f68df840a.md}{Canonical record: 08e37326908a49988dfb787f68df840a}; full qualifications and provenance.}\par
\paragraph*{II.1.3.9.9.6 [Feature] Operation-Bound Authorisation Enforcement}
\phantomsection
\label{navigation:II.1.3.9.9.6}
\addcontentsline{toc}{paragraph}{II.1.3.9.9.6 [Feature] Operation-Bound Authorisation Enforcement}

\label{element:3b7dcee7d58d4d1c960dc92abd58b4c3}
\index{[Feature] Operation-Bound Authorisation Enforcement (II.1.3.9.9.6)}
{\small Capability / Feature · Proposed · pending\_adjudication}\par

Apply effective policy at the governed action boundary and reject actions lacking sufficient authority.



\textbf{Overview}



Consistent enforcement is required across mutation, access, execution and administration rather than relying on each caller to voluntarily invoke a check.



\textbf{Boundary}



Consumes trusted identity and criteria; it does not authenticate or establish clinical privilege from directory roles.


{\small Recorded unresolved matters: source-qualified matters; exact scope remains in the canonical record.}\par
{\footnotesize \href{../../../../continuum/strategy/capability/feature/3b7dcee7d58d4d1c960dc92abd58b4c3.md}{Canonical record: 3b7dcee7d58d4d1c960dc92abd58b4c3}; full qualifications and provenance.}\par
\paragraph*{II.1.3.9.9.7 [Feature] Security Context Propagation}
\phantomsection
\label{navigation:II.1.3.9.9.7}
\addcontentsline{toc}{paragraph}{II.1.3.9.9.7 [Feature] Security Context Propagation}

\label{element:2b071c4fb10c4f61afecd18d09387214}
\index{[Feature] Security Context Propagation (II.1.3.9.9.7)}
{\small Capability / Feature · Canonical · pending\_adjudication}\par

Propagate tamper-evident, attributable security context across asynchronous processing boundaries.



\medskip\hrule\medskip

{\small Recorded unresolved matters: source-qualified matters; exact scope remains in the canonical record.}\par
{\footnotesize \href{../../../../continuum/strategy/capability/feature/2b071c4fb10c4f61afecd18d09387214.md}{Canonical record: 2b071c4fb10c4f61afecd18d09387214}; full qualifications and provenance.}\par
\paragraph*{II.1.3.9.9.8 [Feature] Trusted Identity Context Intake}
\phantomsection
\label{navigation:II.1.3.9.9.8}
\addcontentsline{toc}{paragraph}{II.1.3.9.9.8 [Feature] Trusted Identity Context Intake}

\label{element:2ef8f67907c34cb5a6a37fe7181434e2}
\index{[Feature] Trusted Identity Context Intake (II.1.3.9.9.8)}
{\small Capability / Feature · Proposed · pending\_adjudication}\par

Accept and qualify externally authenticated identity context for Harmonia-governed operations.



\textbf{Overview}



The boundary between WHO and WHAT remains explicit: an external identity assertion is verified as trusted input before internal authorisation is resolved.



\textbf{Boundary}



Harmonia does not become an authentication or credential-issuance provider.


{\small Recorded unresolved matters: source-qualified matters; exact scope remains in the canonical record.}\par
{\footnotesize \href{../../../../continuum/strategy/capability/feature/2ef8f67907c34cb5a6a37fe7181434e2.md}{Canonical record: 2ef8f67907c34cb5a6a37fe7181434e2}; full qualifications and provenance.}\par
\subsubsection*{II.1.3.9.10 [CT2] Health Information Exchange Enablement}
\phantomsection
\label{navigation:II.1.3.9.10}
\addcontentsline{toc}{subsubsection}{II.1.3.9.10 [CT2] Health Information Exchange Enablement}

{\small Historical predecessor: R.2.191 \hyperref[element:330d1c43856a4322a2e9f331b1a5ed20]{[CT2] Health Information Exchange}; Canonical ID: \texttt{\seqsplit{330d1c43856a4322a2e9f331b1a5ed20}}. Superseded historical identity — evolved correspondence · superseded.}\par
\label{element:a61a15908ed544e6acf3e9cab63a0267}
\index{[CT2] Health Information Exchange Enablement (II.1.3.9.10)}
{\small Capability · Canonical · current}\par

As a BEC, Health Information Exchange provides reusable information-movement enablement through submission, retrieval mediation, addressed distribution and syndication. The models distinguish sending to a participant, requesting information, deliberately distributing to named recipients and publishing changes to matching subscribers. Their common responsibility is the exchange context, including routing and delivery state, rather than the business meaning of every payload.



\textbf{Overview}



As a BEC, Health Information Exchange provides reusable information-movement enablement through submission, retrieval mediation, addressed distribution and syndication. The models distinguish sending to a participant, requesting information, deliberately distributing to named recipients and publishing changes to matching subscribers. Their common responsibility is the exchange context, including routing and delivery state, rather than the business meaning of every payload.



\textbf{Intent}



As a BEC, Health Information Exchange provides reusable information-movement enablement through submission, retrieval mediation, addressed distribution and syndication. The models distinguish sending to a participant, requesting information, deliberately distributing to named recipients and publishing changes to matching subscribers. Their common responsibility is the exchange context, including routing and delivery state, rather than the business meaning of every payload.



\textbf{Boundary}



Information administration does not confer authority over the represented person, organisation, service, place, device or clinical activity.



\medskip\hrule\medskip

{\small Recorded unresolved matters: source-qualified matters; exact scope remains in the canonical record.}\par
{\footnotesize \href{../../../../continuum/strategy/capability/a61a15908ed544e6acf3e9cab63a0267.md}{Canonical record: a61a15908ed544e6acf3e9cab63a0267}; full qualifications and provenance.}\par
\paragraph*{II.1.3.9.10.1 [CT3] HIE Endpoint State Management}
\phantomsection
\label{navigation:II.1.3.9.10.1}
\addcontentsline{toc}{paragraph}{II.1.3.9.10.1 [CT3] HIE Endpoint State Management}

{\small Historical predecessor: R.2.188 \hyperref[element:697de08ac49a4b7f8df7a1ed6f04210e]{[CT3] HIE Endpoint State Management}; Canonical ID: \texttt{\seqsplit{697de08ac49a4b7f8df7a1ed6f04210e}}. Superseded historical identity — evolved correspondence · superseded.}\par
\label{element:1bec951a1c98450b8d4672c5c424f9bd}
\index{[CT3] HIE Endpoint State Management (II.1.3.9.10.1)}
{\small Capability · Proposed · current}\par

Manage recognised exchange endpoint registrations, operational standing and contract-qualified availability information.



\textbf{Overview}



Manage recognised exchange endpoint registrations, operational standing and contract-qualified availability information. The reviewed hierarchy assigns this enduring ability to its stated business-enabling context.



\textbf{Intent}



Manage recognised exchange endpoint registrations, operational standing and contract-qualified availability information.



\textbf{Boundary}



Information administration and coordination preserve the authority of the responsible source, clinical actor and governed domain; they confer no additional real-world management authority.


{\footnotesize \href{../../../../continuum/strategy/capability/1bec951a1c98450b8d4672c5c424f9bd.md}{Canonical record: 1bec951a1c98450b8d4672c5c424f9bd}; full qualifications and provenance.}\par
\paragraph*{II.1.3.9.10.2 [CT3] HIE Implementation Management}
\phantomsection
\label{navigation:II.1.3.9.10.2}
\addcontentsline{toc}{paragraph}{II.1.3.9.10.2 [CT3] HIE Implementation Management}

{\small Historical predecessor: R.2.189 \hyperref[element:a99b4119339545a89dc54ebb216ddb85]{[CT3] HIE Implementation Management}; Canonical ID: \texttt{\seqsplit{a99b4119339545a89dc54ebb216ddb85}}. Superseded historical identity — evolved correspondence · superseded.}\par
\label{element:8ce29be373e04252b7aa881b6a72f375}
\index{[CT3] HIE Implementation Management (II.1.3.9.10.2)}
{\small Capability · Proposed · current}\par

Manage implementation information and progression for established exchange requirements, including submission, retrieval, distribution and syndication interactions.



\textbf{Overview}



Manage implementation information and progression for established exchange requirements, including submission, retrieval, distribution and syndication interactions. The reviewed hierarchy assigns this enduring ability to its stated business-enabling context.



\textbf{Intent}



Manage implementation information and progression for established exchange requirements, including submission, retrieval, distribution and syndication interactions.



\textbf{Boundary}



Information administration and coordination preserve the authority of the responsible source, clinical actor and governed domain; they confer no additional real-world management authority.


{\footnotesize \href{../../../../continuum/strategy/capability/8ce29be373e04252b7aa881b6a72f375.md}{Canonical record: 8ce29be373e04252b7aa881b6a72f375}; full qualifications and provenance.}\par
\subparagraph*{II.1.3.9.10.2.1 [Feature] Addressed Clinical Distribution}
\phantomsection
\label{navigation:II.1.3.9.10.2.1}
\addcontentsline{toc}{subparagraph}{II.1.3.9.10.2.1 [Feature] Addressed Clinical Distribution}

\label{element:b476c86edab644c394da33d3de7bd0bd}
\index{[Feature] Addressed Clinical Distribution (II.1.3.9.10.2.1)}
{\small Capability / Feature · Canonical · current}\par

Fan out clinical events and diagnostic reports to designated external and internal recipients with per-destination delivery tracking.



Addressed distribution sends information to explicitly designated recipients and tracks each destination's delivery separately. This preserves partial outcomes in a fan-out rather than collapsing several deliveries into a single success that hides a missing recipient.


{\footnotesize \href{../../../../continuum/strategy/capability/feature/b476c86edab644c394da33d3de7bd0bd.md}{Canonical record: b476c86edab644c394da33d3de7bd0bd}; full qualifications and provenance.}\par
\subparagraph*{II.1.3.9.10.2.2 [Feature] Clinical Record Retrieval Mediation}
\phantomsection
\label{navigation:II.1.3.9.10.2.2}
\addcontentsline{toc}{subparagraph}{II.1.3.9.10.2.2 [Feature] Clinical Record Retrieval Mediation}

\label{element:6f4a9c8f7916488bb0dd838ccccb1df4}
\index{[Feature] Clinical Record Retrieval Mediation (II.1.3.9.10.2.2)}
{\small Capability / Feature · Canonical · current}\par

Mediate query-based retrieval requests across external and internal repositories.



Retrieval mediation connects a request with information held in participating repositories. It makes distributed retrieval possible without transferring repository authority, acquiring the clinical payload or treating a resolvable request as automatic permission to disclose.


{\footnotesize \href{../../../../continuum/strategy/capability/feature/6f4a9c8f7916488bb0dd838ccccb1df4.md}{Canonical record: 6f4a9c8f7916488bb0dd838ccccb1df4}; full qualifications and provenance.}\par
\subparagraph*{II.1.3.9.10.2.3 [Feature] Direct Submission Ingestion \& Delivery}
\phantomsection
\label{navigation:II.1.3.9.10.2.3}
\addcontentsline{toc}{subparagraph}{II.1.3.9.10.2.3 [Feature] Direct Submission Ingestion \& Delivery}

\label{element:e8bfad88315c4833b6772ef580ca55a5}
\index{[Feature] Direct Submission Ingestion \& Delivery (II.1.3.9.10.2.3)}
{\small Capability / Feature · Canonical · current}\par

Support push submissions of clinical documents and messages across boundaries.



Direct submission supports a sender's deliberate push of information across an exchange boundary. It preserves the distinction between ingress or delivery acceptance and the recipient's business handling, under the actual contract rather than an assumed meaning of transport success.


{\footnotesize \href{../../../../continuum/strategy/capability/feature/e8bfad88315c4833b6772ef580ca55a5.md}{Canonical record: e8bfad88315c4833b6772ef580ca55a5}; full qualifications and provenance.}\par
\subparagraph*{II.1.3.9.10.2.4 [Feature] Event-Driven Syndication}
\phantomsection
\label{navigation:II.1.3.9.10.2.4}
\addcontentsline{toc}{subparagraph}{II.1.3.9.10.2.4 [Feature] Event-Driven Syndication}

\label{element:30bb720ed4e64a47a1fcf3be811c1931}
\index{[Feature] Event-Driven Syndication (II.1.3.9.10.2.4)}
{\small Capability / Feature · Canonical · current}\par

Publish clinical and directory change events to subscribed systems matching interest filters.



Syndication connects published changes with subscribers whose registered interests match them. It differs from named-recipient distribution in how consumers are selected, while retaining governed delivery and the distinction between notification and the subscriber's eventual business action.


{\footnotesize \href{../../../../continuum/strategy/capability/feature/30bb720ed4e64a47a1fcf3be811c1931.md}{Canonical record: 30bb720ed4e64a47a1fcf3be811c1931}; full qualifications and provenance.}\par
\paragraph*{II.1.3.9.10.3 [CT3] HIE Requirement Management}
\phantomsection
\label{navigation:II.1.3.9.10.3}
\addcontentsline{toc}{paragraph}{II.1.3.9.10.3 [CT3] HIE Requirement Management}

{\small Historical predecessor: R.2.190 \hyperref[element:a32c8bff0c284f5c9d6ab1f46991c232]{[CT3] HIE Requirement Management}; Canonical ID: \texttt{\seqsplit{a32c8bff0c284f5c9d6ab1f46991c232}}. Superseded historical identity — evolved correspondence · superseded.}\par
\label{element:1c852a1799984feba50011172983c865}
\index{[CT3] HIE Requirement Management (II.1.3.9.10.3)}
{\small Capability · Proposed · current}\par

Manage the identified information-exchange requirements, interaction contracts and participant obligations applicable to a health-information exchange.



\textbf{Overview}



Manage the identified information-exchange requirements, interaction contracts and participant obligations applicable to a health-information exchange. The reviewed hierarchy assigns this enduring ability to its stated business-enabling context.



\textbf{Intent}



Manage the identified information-exchange requirements, interaction contracts and participant obligations applicable to a health-information exchange.



\textbf{Boundary}



Information administration and coordination preserve the authority of the responsible source, clinical actor and governed domain; they confer no additional real-world management authority.


{\footnotesize \href{../../../../continuum/strategy/capability/1c852a1799984feba50011172983c865.md}{Canonical record: 1c852a1799984feba50011172983c865}; full qualifications and provenance.}\par
\subsubsection*{II.1.3.9.11 [CT2] Information Asset Stewardship Management}
\phantomsection
\label{navigation:II.1.3.9.11}
\addcontentsline{toc}{subsubsection}{II.1.3.9.11 [CT2] Information Asset Stewardship Management}

{\small Historical predecessor: R.2.194 \hyperref[element:cee085c690314802851de71f4741d8db]{[CT2] Information Asset Stewardship Management}; Canonical ID: \texttt{\seqsplit{cee085c690314802851de71f4741d8db}}. Superseded historical identity — direct correspondence · superseded.}\par
\label{element:5bac15316d8d4f4db6e7a253f1acfbe3}
\index{[CT2] Information Asset Stewardship Management (II.1.3.9.11)}
{\small Capability · Proposed · current}\par

Manage definitions and tags that make governed information assets recognisable and interpretable for their permitted uses.



\textbf{Overview}



Manage definitions and tags that make governed information assets recognisable and interpretable for their permitted uses. The reviewed hierarchy assigns this enduring ability to its stated business-enabling context.



\textbf{Intent}



Manage definitions and tags that make governed information assets recognisable and interpretable for their permitted uses.



\textbf{Boundary}



Information administration and coordination preserve the authority of the responsible source, clinical actor and governed domain; they confer no additional real-world management authority.


{\footnotesize \href{../../../../continuum/strategy/capability/5bac15316d8d4f4db6e7a253f1acfbe3.md}{Canonical record: 5bac15316d8d4f4db6e7a253f1acfbe3}; full qualifications and provenance.}\par
\paragraph*{II.1.3.9.11.1 [CT3] Information Asset Definition}
\phantomsection
\label{navigation:II.1.3.9.11.1}
\addcontentsline{toc}{paragraph}{II.1.3.9.11.1 [CT3] Information Asset Definition}

{\small Historical predecessor: R.2.192 \hyperref[element:166a22463d2c4534bda7ae31273398b6]{[CT3] Information Asset Definition}; Canonical ID: \texttt{\seqsplit{166a22463d2c4534bda7ae31273398b6}}. Superseded historical identity — direct correspondence · superseded.}\par
\label{element:05c58b2ccfcc45e9a6c70d9a569c4f61}
\index{[CT3] Information Asset Definition (II.1.3.9.11.1)}
{\small Capability · Proposed · current}\par

Define the identifying, semantic and governed-use characteristics of an information asset.



\textbf{Overview}



Define the identifying, semantic and governed-use characteristics of an information asset. The reviewed hierarchy assigns this enduring ability to its stated business-enabling context.



\textbf{Intent}



Define the identifying, semantic and governed-use characteristics of an information asset.



\textbf{Boundary}



Information administration and coordination preserve the authority of the responsible source, clinical actor and governed domain; they confer no additional real-world management authority.


{\footnotesize \href{../../../../continuum/strategy/capability/05c58b2ccfcc45e9a6c70d9a569c4f61.md}{Canonical record: 05c58b2ccfcc45e9a6c70d9a569c4f61}; full qualifications and provenance.}\par
\paragraph*{II.1.3.9.11.2 [CT3] Information Asset Tag Management}
\phantomsection
\label{navigation:II.1.3.9.11.2}
\addcontentsline{toc}{paragraph}{II.1.3.9.11.2 [CT3] Information Asset Tag Management}

{\small Historical predecessor: R.2.193 \hyperref[element:ae44a4ee78414d1098814315793bd751]{[CT3] Information Asset Tag Management}; Canonical ID: \texttt{\seqsplit{ae44a4ee78414d1098814315793bd751}}. Superseded historical identity — direct correspondence · superseded.}\par
\label{element:df27c3ba6e164e56af7bff4e1eb1954a}
\index{[CT3] Information Asset Tag Management (II.1.3.9.11.2)}
{\small Capability · Proposed · current}\par

Manage attributable tags describing information assets without treating caller-controlled tags as security authority.



\textbf{Overview}



Manage attributable tags describing information assets without treating caller-controlled tags as security authority. The reviewed hierarchy assigns this enduring ability to its stated business-enabling context.



\textbf{Intent}



Manage attributable tags describing information assets without treating caller-controlled tags as security authority.



\textbf{Boundary}



Information administration and coordination preserve the authority of the responsible source, clinical actor and governed domain; they confer no additional real-world management authority.


{\footnotesize \href{../../../../continuum/strategy/capability/df27c3ba6e164e56af7bff4e1eb1954a.md}{Canonical record: df27c3ba6e164e56af7bff4e1eb1954a}; full qualifications and provenance.}\par
\subsubsection*{II.1.3.9.12 [CT2] Information Content Maintainer Management}
\phantomsection
\label{navigation:II.1.3.9.12}
\addcontentsline{toc}{subsubsection}{II.1.3.9.12 [CT2] Information Content Maintainer Management}

{\small Historical predecessor: R.2.198 \hyperref[element:acb977e879204f1491293e5452e0ddf0]{[CT2] Information Content Maintainer Management}; Canonical ID: \texttt{\seqsplit{acb977e879204f1491293e5452e0ddf0}}. Superseded historical identity — direct correspondence · superseded.}\par
\label{element:b81a1ab42af24c2e9896b8f3f78d291a}
\index{[CT2] Information Content Maintainer Management (II.1.3.9.12)}
{\small Capability · Proposed · current}\par

Administer information-content maintainer participation, source provenance and governed information products.



\textbf{Overview}



Administer information-content maintainer participation, source provenance and governed information products. The reviewed hierarchy assigns this enduring ability to its stated business-enabling context.



\textbf{Intent}



Administer information-content maintainer participation, source provenance and governed information products.



\textbf{Boundary}



Information administration and coordination preserve the authority of the responsible source, clinical actor and governed domain; they confer no additional real-world management authority.


{\footnotesize \href{../../../../continuum/strategy/capability/b81a1ab42af24c2e9896b8f3f78d291a.md}{Canonical record: b81a1ab42af24c2e9896b8f3f78d291a}; full qualifications and provenance.}\par
\paragraph*{II.1.3.9.12.1 [CT3] Data Provenance Capture}
\phantomsection
\label{navigation:II.1.3.9.12.1}
\addcontentsline{toc}{paragraph}{II.1.3.9.12.1 [CT3] Data Provenance Capture}

{\small Historical predecessor: R.2.195 \hyperref[element:61258c7fe41f4541871a67ec9d7185cf]{[CT3] Data Provenance Capture}; Canonical ID: \texttt{\seqsplit{61258c7fe41f4541871a67ec9d7185cf}}. Superseded historical identity — direct correspondence · superseded.}\par
\label{element:44a7937520a54945b06b5fd4a2bbd870}
\index{[CT3] Data Provenance Capture (II.1.3.9.12.1)}
{\small Capability · Proposed · current}\par

Capture meaningful source attribution and transformation lineage for governed information content without transferring originating authority.



\textbf{Overview}



Capture meaningful source attribution and transformation lineage for governed information content without transferring originating authority.



\textbf{Intent}



Capture meaningful source attribution and transformation lineage for governed information content without transferring originating authority.



\textbf{Boundary}



Lineage does not grant originating authority, decide clinical correctness or replace activity audit and independent assurance.


{\footnotesize \href{../../../../continuum/strategy/capability/44a7937520a54945b06b5fd4a2bbd870.md}{Canonical record: 44a7937520a54945b06b5fd4a2bbd870}; full qualifications and provenance.}\par
\subparagraph*{II.1.3.9.12.1.1 [Feature] Information Transformation Lineage}
\phantomsection
\label{navigation:II.1.3.9.12.1.1}
\addcontentsline{toc}{subparagraph}{II.1.3.9.12.1.1 [Feature] Information Transformation Lineage}

\label{element:cc47f49b85a34e87bafc49de7f0426e2}
\index{[Feature] Information Transformation Lineage (II.1.3.9.12.1.1)}
{\small Capability / Feature · Proposed · current}\par

Record meaningful source-to-derived information relationships and transformation attribution.



\textbf{Overview}



A consumer should be able to explain which source assertions contributed to a managed representation and what significant transformation occurred.



\textbf{Boundary}



Routine cache reconstruction, retries and diagnostics do not automatically become permanent lineage.


{\footnotesize \href{../../../../continuum/strategy/capability/feature/cc47f49b85a34e87bafc49de7f0426e2.md}{Canonical record: cc47f49b85a34e87bafc49de7f0426e2}; full qualifications and provenance.}\par
\paragraph*{II.1.3.9.12.2 [CT3] Data Scientist Administration}
\phantomsection
\label{navigation:II.1.3.9.12.2}
\addcontentsline{toc}{paragraph}{II.1.3.9.12.2 [CT3] Data Scientist Administration}

{\small Historical predecessor: R.2.196 \hyperref[element:1cb497aec93c4b82ba7f5b622813e54e]{[CT3] Data Scientist Administration}; Canonical ID: \texttt{\seqsplit{1cb497aec93c4b82ba7f5b622813e54e}}. Superseded historical identity — direct correspondence · superseded.}\par
\label{element:3bb1a11fba234717bf40c9e6452f401a}
\index{[CT3] Data Scientist Administration (II.1.3.9.12.2)}
{\small Capability · Proposed · current}\par

Administer profiles and qualified participation information for data scientists maintaining governed information content.



\textbf{Overview}



Administer profiles and qualified participation information for data scientists maintaining governed information content. The reviewed hierarchy assigns this enduring ability to its stated business-enabling context.



\textbf{Intent}



Administer profiles and qualified participation information for data scientists maintaining governed information content.



\textbf{Boundary}



Information administration and coordination preserve the authority of the responsible source, clinical actor and governed domain; they confer no additional real-world management authority.


{\footnotesize \href{../../../../continuum/strategy/capability/3bb1a11fba234717bf40c9e6452f401a.md}{Canonical record: 3bb1a11fba234717bf40c9e6452f401a}; full qualifications and provenance.}\par
\paragraph*{II.1.3.9.12.3 [CT3] Information Product Administration}
\phantomsection
\label{navigation:II.1.3.9.12.3}
\addcontentsline{toc}{paragraph}{II.1.3.9.12.3 [CT3] Information Product Administration}

{\small Historical predecessor: R.2.197 \hyperref[element:5d9b7f8d2d7b41759b6d4931c1d2aaa2]{[CT3] Information Product Administration}; Canonical ID: \texttt{\seqsplit{5d9b7f8d2d7b41759b6d4931c1d2aaa2}}. Superseded historical identity — direct correspondence · superseded.}\par
\label{element:b5a6a31855834ea58870ebb6c51ed4b0}
\index{[CT3] Information Product Administration (II.1.3.9.12.3)}
{\small Capability · Proposed · current}\par

Administer identifiable governed information products, their purpose, content standing and qualified release/use context.



\textbf{Overview}



Administer identifiable governed information products, their purpose, content standing and qualified release/use context. The reviewed hierarchy assigns this enduring ability to its stated business-enabling context.



\textbf{Intent}



Administer identifiable governed information products, their purpose, content standing and qualified release/use context.



\textbf{Boundary}



Information administration and coordination preserve the authority of the responsible source, clinical actor and governed domain; they confer no additional real-world management authority.


{\footnotesize \href{../../../../continuum/strategy/capability/b5a6a31855834ea58870ebb6c51ed4b0.md}{Canonical record: b5a6a31855834ea58870ebb6c51ed4b0}; full qualifications and provenance.}\par
\subparagraph*{II.1.3.9.12.3.1 [Feature] Authorised Research Release Assembly}
\phantomsection
\label{navigation:II.1.3.9.12.3.1}
\addcontentsline{toc}{subparagraph}{II.1.3.9.12.3.1 [Feature] Authorised Research Release Assembly}

\label{element:dcebc2f2b5b7494a8e80a02402664078}
\index{[Feature] Authorised Research Release Assembly (II.1.3.9.12.3.1)}
{\small Capability / Feature · Proposed · current}\par

Assemble and qualify the existing information allowed by a particular authorised research purpose.



\textbf{Overview}



The release carries purpose and applicable privacy/consent/use constraints; the requester does not gain unrestricted clinical-record access by having a research role.



\textbf{Boundary}



No ethics approval, de-identification guarantee or clinical-trial management is inferred.


{\footnotesize \href{../../../../continuum/strategy/capability/feature/dcebc2f2b5b7494a8e80a02402664078.md}{Canonical record: dcebc2f2b5b7494a8e80a02402664078}; full qualifications and provenance.}\par
\subparagraph*{II.1.3.9.12.3.2 [Feature] Demand \& Capacity Planning Projection}
\phantomsection
\label{navigation:II.1.3.9.12.3.2}
\addcontentsline{toc}{subparagraph}{II.1.3.9.12.3.2 [Feature] Demand \& Capacity Planning Projection}

\label{element:4f97baa3aed24b8385bf3ecda4eab622}
\index{[Feature] Demand \& Capacity Planning Projection (II.1.3.9.12.3.2)}
{\small Capability / Feature · Proposed · current}\par

Project qualified observed demand, capacity and utilisation for authorised planning consumers.



\textbf{Overview}



Comparisons need service, temporal and context qualifiers so that planned future capacity is not confused with present operational capacity.



\textbf{Boundary}



Projection does not establish a planning conclusion or allocate services and staff.


{\footnotesize \href{../../../../continuum/strategy/capability/feature/4f97baa3aed24b8385bf3ecda4eab622.md}{Canonical record: 4f97baa3aed24b8385bf3ecda4eab622}; full qualifications and provenance.}\par
\subparagraph*{II.1.3.9.12.3.3 [Feature] Population Notification Obligation Tracking}
\phantomsection
\label{navigation:II.1.3.9.12.3.3}
\addcontentsline{toc}{subparagraph}{II.1.3.9.12.3.3 [Feature] Population Notification Obligation Tracking}

\label{element:62f6a91a10684876b7d08b536700f12f}
\index{[Feature] Population Notification Obligation Tracking (II.1.3.9.12.3.3)}
{\small Capability / Feature · Proposed · current}\par

Track the relation between an applicable notification obligation, the supplied source evidence and observed notification disposition.



\textbf{Overview}



Sending a required notification and discharging the governing obligation are different assertions. This Feature keeps the explicit contract and indeterminate responses visible.



\textbf{Boundary}



Does not define notifiability, author statutory rules or own public-health response.


{\footnotesize \href{../../../../continuum/strategy/capability/feature/62f6a91a10684876b7d08b536700f12f.md}{Canonical record: 62f6a91a10684876b7d08b536700f12f}; full qualifications and provenance.}\par
\subsubsection*{II.1.3.9.13 [CT2] Information Design Governance}
\phantomsection
\label{navigation:II.1.3.9.13}
\addcontentsline{toc}{subsubsection}{II.1.3.9.13 [CT2] Information Design Governance}

{\small Historical predecessor: R.2.203 \hyperref[element:0bdd7f8c0cd74d9fabc49801af037fc2]{[CT2] Information Design Governance}; Canonical ID: \texttt{\seqsplit{0bdd7f8c0cd74d9fabc49801af037fc2}}. Superseded historical identity — direct correspondence · superseded.}\par
\label{element:f39ed9ebcc724857a4bf7a2b8e045c7a}
\index{[CT2] Information Design Governance (II.1.3.9.13)}
{\small Capability · Canonical · current}\par

Information Design Governance maintains canonical definitions, semantic bindings, exchange profiles and the conformance rules applied to representations. It makes meaning and constraints available across regional information flows so that an accepted technical format is not mistaken for semantic correctness. Governance and verification remain distinguishable from the clinical facts represented.



\textbf{Overview}



Information Design Governance maintains canonical definitions, semantic bindings, exchange profiles and the conformance rules applied to representations. It makes meaning and constraints available across regional information flows so that an accepted technical format is not mistaken for semantic correctness. Governance and verification remain distinguishable from the clinical facts represented.



\textbf{Intent}



Information Design Governance maintains canonical definitions, semantic bindings, exchange profiles and the conformance rules applied to representations. It makes meaning and constraints available across regional information flows so that an accepted technical format is not mistaken for semantic correctness. Governance and verification remain distinguishable from the clinical facts represented.



\textbf{Boundary}



Information administration does not confer authority over the represented person, organisation, service, place, device or clinical activity.



\medskip\hrule\medskip

{\footnotesize \href{../../../../continuum/strategy/capability/f39ed9ebcc724857a4bf7a2b8e045c7a.md}{Canonical record: f39ed9ebcc724857a4bf7a2b8e045c7a}; full qualifications and provenance.}\par
\paragraph*{II.1.3.9.13.1 [CT3] Enterprise Canonical/Reusable Model Design}
\phantomsection
\label{navigation:II.1.3.9.13.1}
\addcontentsline{toc}{paragraph}{II.1.3.9.13.1 [CT3] Enterprise Canonical/Reusable Model Design}

{\small Historical predecessor: R.2.199 \hyperref[element:dbc37445ad944a85a7c4f26f410967b2]{[CT3] Enterprise Canonical/Reusable Model Design}; Canonical ID: \texttt{\seqsplit{dbc37445ad944a85a7c4f26f410967b2}}. Superseded historical identity — direct correspondence · superseded.}\par
\label{element:90e9cd401ddf4c7690c978d8dd131413}
\index{[CT3] Enterprise Canonical/Reusable Model Design (II.1.3.9.13.1)}
{\small Capability · Proposed · current}\par

Design canonical and reusable enterprise information models and semantic bindings under applicable information-design governance.



\textbf{Overview}



Design canonical and reusable enterprise information models and semantic bindings under applicable information-design governance. The reviewed hierarchy assigns this enduring ability to its stated business-enabling context.



\textbf{Intent}



Design canonical and reusable enterprise information models and semantic bindings under applicable information-design governance.



\textbf{Boundary}



Information administration and coordination preserve the authority of the responsible source, clinical actor and governed domain; they confer no additional real-world management authority.


{\footnotesize \href{../../../../continuum/strategy/capability/90e9cd401ddf4c7690c978d8dd131413.md}{Canonical record: 90e9cd401ddf4c7690c978d8dd131413}; full qualifications and provenance.}\par
\subparagraph*{II.1.3.9.13.1.1 [Feature] Exchange Profile Governance}
\phantomsection
\label{navigation:II.1.3.9.13.1.1}
\addcontentsline{toc}{subparagraph}{II.1.3.9.13.1.1 [Feature] Exchange Profile Governance}

\label{element:5e59d8f2e23e4fb5b59a5bd1dd3f31fa}
\index{[Feature] Exchange Profile Governance (II.1.3.9.13.1.1)}
{\small Capability / Feature · Proposed · current}\par

Govern the semantic and conformance constraints of an interoperability profile.



\textbf{Overview}



An external contract must state which representations and extensions are permitted; validation applies that contract but does not define it.



\textbf{Boundary}



Profile governance does not transmit messages or publish internal operational metadata.


{\footnotesize \href{../../../../continuum/strategy/capability/feature/5e59d8f2e23e4fb5b59a5bd1dd3f31fa.md}{Canonical record: 5e59d8f2e23e4fb5b59a5bd1dd3f31fa}; full qualifications and provenance.}\par
\subparagraph*{II.1.3.9.13.1.2 [Feature] Information Definition \& Schema Governance}
\phantomsection
\label{navigation:II.1.3.9.13.1.2}
\addcontentsline{toc}{subparagraph}{II.1.3.9.13.1.2 [Feature] Information Definition \& Schema Governance}

\label{element:3317bbb7e2db4384abf44a50bf67422c}
\index{[Feature] Information Definition \& Schema Governance (II.1.3.9.13.1.2)}
{\small Capability / Feature · Proposed · current}\par

Govern structural definitions, data dictionaries and constraint meaning.



\textbf{Overview}



Existing definition-governance behaviour needs a distinct Feature frame from general semantic-model authorship and runtime validation.



\textbf{Boundary}



Governed definitions do not own runtime instances or their clinical facts.


{\footnotesize \href{../../../../continuum/strategy/capability/feature/3317bbb7e2db4384abf44a50bf67422c.md}{Canonical record: 3317bbb7e2db4384abf44a50bf67422c}; full qualifications and provenance.}\par
\subparagraph*{II.1.3.9.13.1.3 [Feature] Information Standards \& Semantic Governance}
\phantomsection
\label{navigation:II.1.3.9.13.1.3}
\addcontentsline{toc}{subparagraph}{II.1.3.9.13.1.3 [Feature] Information Standards \& Semantic Governance}

\label{element:af5d5249302f47ebb393c91afcc60392}
\index{[Feature] Information Standards \& Semantic Governance (II.1.3.9.13.1.3)}
{\small Capability / Feature · Canonical · current}\par

Govern canonical clinical data definitions, terminology bindings, and exchange profiles across regional systems.



Semantic governance provides canonical definitions, bindings and profiles through which regional representations can be understood consistently. Its purpose concerns meaning and constraints rather than a particular schema technology, and consuming a definition does not transfer its governing authority.


{\footnotesize \href{../../../../continuum/strategy/capability/feature/af5d5249302f47ebb393c91afcc60392.md}{Canonical record: af5d5249302f47ebb393c91afcc60392}; full qualifications and provenance.}\par
\subparagraph*{II.1.3.9.13.1.4 [Feature] Terminology Binding Governance}
\phantomsection
\label{navigation:II.1.3.9.13.1.4}
\addcontentsline{toc}{subparagraph}{II.1.3.9.13.1.4 [Feature] Terminology Binding Governance}

\label{element:b999df692ca64010b4adeedbcb0d3a5a}
\index{[Feature] Terminology Binding Governance (II.1.3.9.13.1.4)}
{\small Capability / Feature · Proposed · current}\par

Govern how information definitions are bound to applicable terminology and value sets.



\textbf{Overview}



Binding a data element to a vocabulary is distinct from resolving a patient code or changing the runtime membership of a dataset.



\textbf{Boundary}



Definition of a binding does not take external ontology authority or execute policy.


{\footnotesize \href{../../../../continuum/strategy/capability/feature/b999df692ca64010b4adeedbcb0d3a5a.md}{Canonical record: b999df692ca64010b4adeedbcb0d3a5a}; full qualifications and provenance.}\par
\paragraph*{II.1.3.9.13.2 [CT3] Enterprise Canonical/Reusable Model Library Management}
\phantomsection
\label{navigation:II.1.3.9.13.2}
\addcontentsline{toc}{paragraph}{II.1.3.9.13.2 [CT3] Enterprise Canonical/Reusable Model Library Management}

{\small Historical predecessor: R.2.200 \hyperref[element:f5eccd4f65df49bba98104e5abfa5e19]{[CT3] Enterprise Canonical/Reusable Model Library Management}; Canonical ID: \texttt{\seqsplit{f5eccd4f65df49bba98104e5abfa5e19}}. Superseded historical identity — direct correspondence · superseded.}\par
\label{element:76943b4b68644c4dac1368fce8bb20d9}
\index{[CT3] Enterprise Canonical/Reusable Model Library Management (II.1.3.9.13.2)}
{\small Capability · Proposed · current}\par

Manage the identifiable library of canonical and reusable enterprise model definitions and their governed revisions.



\textbf{Overview}



Manage the identifiable library of canonical and reusable enterprise model definitions and their governed revisions. The reviewed hierarchy assigns this enduring ability to its stated business-enabling context.



\textbf{Intent}



Manage the identifiable library of canonical and reusable enterprise model definitions and their governed revisions.



\textbf{Boundary}



Information administration and coordination preserve the authority of the responsible source, clinical actor and governed domain; they confer no additional real-world management authority.


{\footnotesize \href{../../../../continuum/strategy/capability/76943b4b68644c4dac1368fce8bb20d9.md}{Canonical record: 76943b4b68644c4dac1368fce8bb20d9}; full qualifications and provenance.}\par
\paragraph*{II.1.3.9.13.3 [CT3] Enterprise Canonical/Reusable Model Utilisation Support}
\phantomsection
\label{navigation:II.1.3.9.13.3}
\addcontentsline{toc}{paragraph}{II.1.3.9.13.3 [CT3] Enterprise Canonical/Reusable Model Utilisation Support}

{\small Historical predecessor: R.2.201 \hyperref[element:7ae95677a00346f18a0ba01522a20153]{[CT3] Enterprise Canonical/Reusable Model Utilisation Support}; Canonical ID: \texttt{\seqsplit{7ae95677a00346f18a0ba01522a20153}}. Superseded historical identity — direct correspondence · superseded.}\par
\label{element:5744cb9d86494064af87afb3cbf2cf97}
\index{[CT3] Enterprise Canonical/Reusable Model Utilisation Support (II.1.3.9.13.3)}
{\small Capability · Proposed · current}\par

Support correct utilisation and conformance of canonical and reusable models across their intended consumers.



\textbf{Overview}



Support correct utilisation and conformance of canonical and reusable models across their intended consumers. The reviewed hierarchy assigns this enduring ability to its stated business-enabling context.



\textbf{Intent}



Support correct utilisation and conformance of canonical and reusable models across their intended consumers.



\textbf{Boundary}



Information administration and coordination preserve the authority of the responsible source, clinical actor and governed domain; they confer no additional real-world management authority.


{\footnotesize \href{../../../../continuum/strategy/capability/5744cb9d86494064af87afb3cbf2cf97.md}{Canonical record: 5744cb9d86494064af87afb3cbf2cf97}; full qualifications and provenance.}\par
\subparagraph*{II.1.3.9.13.3.1 [Feature] Semantic Conformance Verification}
\phantomsection
\label{navigation:II.1.3.9.13.3.1}
\addcontentsline{toc}{subparagraph}{II.1.3.9.13.3.1 [Feature] Semantic Conformance Verification}

\label{element:d170a68efcae4aa59350f565ca69443a}
\index{[Feature] Semantic Conformance Verification (II.1.3.9.13.3.1)}
{\small Capability / Feature · Canonical · current}\par

Verify inbound and outbound representations against published semantic rules and conformance constraints.



Conformance verification evaluates a representation against published semantic rules and constraints. It distinguishes meaningful conformity from simple parseability while retaining that a conformant message is not necessarily clinically true, authoritative or permitted for disclosure.


{\footnotesize \href{../../../../continuum/strategy/capability/feature/d170a68efcae4aa59350f565ca69443a.md}{Canonical record: d170a68efcae4aa59350f565ca69443a}; full qualifications and provenance.}\par
\paragraph*{II.1.3.9.13.4 [CT3] Enterprise Reference Data Stewardship}
\phantomsection
\label{navigation:II.1.3.9.13.4}
\addcontentsline{toc}{paragraph}{II.1.3.9.13.4 [CT3] Enterprise Reference Data Stewardship}

{\small Historical predecessor: R.2.202 \hyperref[element:fd9f3239fdb64ea687825b66bb4d1f17]{[CT3] Enterprise Reference Data Stewardship}; Canonical ID: \texttt{\seqsplit{fd9f3239fdb64ea687825b66bb4d1f17}}. Superseded historical identity — direct correspondence · superseded.}\par
\label{element:971d4b079ad24e89b70a8ed8dec0c8bf}
\index{[CT3] Enterprise Reference Data Stewardship (II.1.3.9.13.4)}
{\small Capability · Proposed · current}\par

Steward enterprise reference-data definitions, bindings and applicable governance independently of mutable runtime dataset administration.



\textbf{Overview}



Steward enterprise reference-data definitions, bindings and applicable governance independently of mutable runtime dataset administration. The reviewed hierarchy assigns this enduring ability to its stated business-enabling context.



\textbf{Intent}



Steward enterprise reference-data definitions, bindings and applicable governance independently of mutable runtime dataset administration.



\textbf{Boundary}



Information administration and coordination preserve the authority of the responsible source, clinical actor and governed domain; they confer no additional real-world management authority.


{\footnotesize \href{../../../../continuum/strategy/capability/971d4b079ad24e89b70a8ed8dec0c8bf.md}{Canonical record: 971d4b079ad24e89b70a8ed8dec0c8bf}; full qualifications and provenance.}\par
\subsubsection*{II.1.3.9.14 [CT2] Information Protection \& Accountability Enablement [superseded]}
\phantomsection
\label{navigation:II.1.3.9.14}
\addcontentsline{toc}{subsubsection}{II.1.3.9.14 [CT2] Information Protection \& Accountability Enablement [superseded]}

\label{element:e922cba7cf1d4ebeb9abf6472a3803ce}
\index{[CT2] Information Protection \& Accountability Enablement (II.1.3.9.14)}
{\small Capability · Proposed · superseded}\par
\href{../../../../continuum/strategy/capability/e922cba7cf1d4ebeb9abf6472a3803ce.md}{Historical canonical definition and provenance}.\par
\paragraph*{II.1.3.9.14.1 [CT3] Governance Requirement Publication Support}
\phantomsection
\label{navigation:II.1.3.9.14.1}
\addcontentsline{toc}{paragraph}{II.1.3.9.14.1 [CT3] Governance Requirement Publication Support}

\label{element:7e780751d50e4ffe942f8443b7cb6567}
\index{[CT3] Governance Requirement Publication Support (II.1.3.9.14.1)}
{\small Capability · Proposed · pending\_adjudication}\par

Represent and publish versioned governing requirements for use by authorised operational and assurance consumers.



\textbf{Overview}



Consumers need an identifiable statement of which rule applies and who established it. This proposal fills BC-17 information support without giving the platform authority to make legislation, approve governance requirements or reinterpret their force.



\textbf{Intent}



Let operational and assurance consumers identify the governing statement and applicability basis they actually used.



\textbf{Boundary}



Governance approval remains externally mandated; operational policy, assurance criteria and requirement authority are not interchangeable.


{\small Recorded unresolved matters: source-qualified matters; exact scope remains in the canonical record.}\par
{\footnotesize \href{../../../../continuum/strategy/capability/7e780751d50e4ffe942f8443b7cb6567.md}{Canonical record: 7e780751d50e4ffe942f8443b7cb6567}; full qualifications and provenance.}\par
\subparagraph*{II.1.3.9.14.1.1 [Feature] Governance Requirement Version Publication}
\phantomsection
\label{navigation:II.1.3.9.14.1.1}
\addcontentsline{toc}{subparagraph}{II.1.3.9.14.1.1 [Feature] Governance Requirement Version Publication}

\label{element:001a067a4b3b46bebfe945f1bef541ce}
\index{[Feature] Governance Requirement Version Publication (II.1.3.9.14.1.1)}
{\small Capability / Feature · Proposed · pending\_adjudication}\par

Publish attributable versions and applicability of externally approved governing requirements.



\textbf{Overview}



Operational and assurance consumers need to identify the governing basis they used rather than silently substitute a current rule for a historical one.



\textbf{Boundary}



Publication does not confer requirement approval, authorisation policy or assurance-criteria authority.


{\small Recorded unresolved matters: source-qualified matters; exact scope remains in the canonical record.}\par
{\footnotesize \href{../../../../continuum/strategy/capability/feature/001a067a4b3b46bebfe945f1bef541ce.md}{Canonical record: 001a067a4b3b46bebfe945f1bef541ce}; full qualifications and provenance.}\par
\paragraph*{II.1.3.9.14.2 [CT3] Significant Activity Audit Capture}
\phantomsection
\label{navigation:II.1.3.9.14.2}
\addcontentsline{toc}{paragraph}{II.1.3.9.14.2 [CT3] Significant Activity Audit Capture}

\label{element:d903948169eb4785b0f03d916c842a8a}
\index{[CT3] Significant Activity Audit Capture (II.1.3.9.14.2)}
{\small Capability · Proposed · pending\_adjudication}\par

Capture required security- and information-significant activity as governed, attributable audit evidence.



\textbf{Overview}



Mandatory evidence capture has its own policy-governed platform authority and must remain possible when the audited identity is denied an action. The captured fact is distinguished from debug output and the identity whose activity is recorded.



\textbf{Intent}



Ensure required significant activity evidence can be captured under the authority specifically assigned to that responsibility.



\textbf{Boundary}



The specific intrinsic authority for mandatory audit capture is not generalised to other activities. Capture is not an assurance finding and cannot expose unmasked PHI in diagnostic logs.


{\small Recorded unresolved matters: source-qualified matters; exact scope remains in the canonical record.}\par
{\footnotesize \href{../../../../continuum/strategy/capability/d903948169eb4785b0f03d916c842a8a.md}{Canonical record: d903948169eb4785b0f03d916c842a8a}; full qualifications and provenance.}\par
\subparagraph*{II.1.3.9.14.2.1 [Feature] Mandatory Significant Activity Audit Capture}
\phantomsection
\label{navigation:II.1.3.9.14.2.1}
\addcontentsline{toc}{subparagraph}{II.1.3.9.14.2.1 [Feature] Mandatory Significant Activity Audit Capture}

\label{element:b69936757c4446a7a1b57a455aae7a5f}
\index{[Feature] Mandatory Significant Activity Audit Capture (II.1.3.9.14.2.1)}
{\small Capability / Feature · Proposed · pending\_adjudication}\par

Capture required significant activity evidence using the specifically governed authoritative audit-capture boundary.



\textbf{Overview}



The actor whose action is recorded need not be entitled to persist the audit. Required capture uses its own established authority and failure can block the governed operation.



\textbf{Boundary}



Specific audit authority is not borrowed from the audited identity or generalised to other operations.


{\small Recorded unresolved matters: source-qualified matters; exact scope remains in the canonical record.}\par
{\footnotesize \href{../../../../continuum/strategy/capability/feature/b69936757c4446a7a1b57a455aae7a5f.md}{Canonical record: b69936757c4446a7a1b57a455aae7a5f}; full qualifications and provenance.}\par
\subsubsection*{II.1.3.9.15 [CT2] Patient Clinical Record}
\phantomsection
\label{navigation:II.1.3.9.15}
\addcontentsline{toc}{subsubsection}{II.1.3.9.15 [CT2] Patient Clinical Record}

{\small Historical predecessor: R.2.206 \hyperref[element:2e4de521bda24ad7b0ea64ae535b7f84]{[CT2] Patient Clinical Record}; Canonical ID: \texttt{\seqsplit{2e4de521bda24ad7b0ea64ae535b7f84}}. Superseded historical identity — direct correspondence · superseded.}\par
\label{element:5478878e0e6e4d34a8d5023ac6ac2d54}
\index{[CT2] Patient Clinical Record (II.1.3.9.15)}
{\small Capability · Canonical · current}\par

Patient Clinical Record provides a governed longitudinal representation across encounters, providers and source systems. Assembly supplies the timeline, active access supports current care context, and durable preservation retains the history of what Harmonia received, knew or managed. These contributions are complementary: a responsive working view must not become the authority for durable history.



\textbf{Overview}



Patient Clinical Record provides a governed longitudinal representation across encounters, providers and source systems. Assembly supplies the timeline, active access supports current care context, and durable preservation retains the history of what Harmonia received, knew or managed. These contributions are complementary: a responsive working view must not become the authority for durable history.



\textbf{Intent}



Patient Clinical Record provides a governed longitudinal representation across encounters, providers and source systems. Assembly supplies the timeline, active access supports current care context, and durable preservation retains the history of what Harmonia received, knew or managed. These contributions are complementary: a responsive working view must not become the authority for durable history.



\textbf{Boundary}



Information administration does not confer authority over the represented person, organisation, service, place, device or clinical activity.



\medskip\hrule\medskip

{\footnotesize \href{../../../../continuum/strategy/capability/5478878e0e6e4d34a8d5023ac6ac2d54.md}{Canonical record: 5478878e0e6e4d34a8d5023ac6ac2d54}; full qualifications and provenance.}\par
\paragraph*{II.1.3.9.15.1 [CT3] Linked Information Curation}
\phantomsection
\label{navigation:II.1.3.9.15.1}
\addcontentsline{toc}{paragraph}{II.1.3.9.15.1 [CT3] Linked Information Curation}

{\small Historical predecessor: R.2.204 \hyperref[element:9d0925bb39664a5985cacf3cea7442d1]{[CT3] Linked Information Curation}; Canonical ID: \texttt{\seqsplit{9d0925bb39664a5985cacf3cea7442d1}}. Superseded historical identity — direct correspondence · superseded.}\par
\label{element:5edb2a8b5ee54ddda8104d1fab323a94}
\index{[CT3] Linked Information Curation (II.1.3.9.15.1)}
{\small Capability · Proposed · current}\par

Curate associations to information and binary content needed to interpret and use a clinical record, retaining custody and source-authority distinctions.



\textbf{Overview}



Curate associations to information and binary content needed to interpret and use a clinical record, retaining custody and source-authority distinctions.



\textbf{Intent}



Curate associations to information and binary content needed to interpret and use a clinical record, retaining custody and source-authority distinctions.



\textbf{Boundary}



Association is not physical custody. Persistence authority, source clinical meaning and external repository authority remain separate.


{\footnotesize \href{../../../../continuum/strategy/capability/5edb2a8b5ee54ddda8104d1fab323a94.md}{Canonical record: 5edb2a8b5ee54ddda8104d1fab323a94}; full qualifications and provenance.}\par
\subparagraph*{II.1.3.9.15.1.1 [Feature] Binary Content Resolution \& Assembly}
\phantomsection
\label{navigation:II.1.3.9.15.1.1}
\addcontentsline{toc}{subparagraph}{II.1.3.9.15.1.1 [Feature] Binary Content Resolution \& Assembly}

\label{element:8fb4036f7a9c47c3883b2507528fab5e}
\index{[Feature] Binary Content Resolution \& Assembly (II.1.3.9.15.1.1)}
{\small Capability / Feature · Proposed · current}\par

Resolve governed payload associations and retrieve or assemble the content required by a requesting information context.



\textbf{Overview}



A logical clinical item may reference externally held or multipart payloads. The Feature supplies a usable content result while retaining custody and missing-content standing.



\textbf{Boundary}



No assumption of local copying, shared storage, universal content retention or originating clinical authority.


{\footnotesize \href{../../../../continuum/strategy/capability/feature/8fb4036f7a9c47c3883b2507528fab5e.md}{Canonical record: 8fb4036f7a9c47c3883b2507528fab5e}; full qualifications and provenance.}\par
\paragraph*{II.1.3.9.15.2 [CT3] Longitudinal Clinical Record Administration}
\phantomsection
\label{navigation:II.1.3.9.15.2}
\addcontentsline{toc}{paragraph}{II.1.3.9.15.2 [CT3] Longitudinal Clinical Record Administration}

{\small Historical predecessor: R.2.205 \hyperref[element:bd91d848cf424c42bec7c80cd32edd28]{[CT3] Longitudinal Clinical Record Administration}; Canonical ID: \texttt{\seqsplit{bd91d848cf424c42bec7c80cd32edd28}}. Superseded historical identity — direct correspondence · superseded.}\par
\label{element:14748dc6117f4b6ba7c0fc554b76bd89}
\index{[CT3] Longitudinal Clinical Record Administration (II.1.3.9.15.2)}
{\small Capability · Proposed · current}\par

Administer longitudinal record assembly, active clinical summaries and durable received/managed history with source attribution and supersession.



\textbf{Overview}



Administer longitudinal record assembly, active clinical summaries and durable received/managed history with source attribution and supersession. The reviewed hierarchy assigns this enduring ability to its stated business-enabling context.



\textbf{Intent}



Administer longitudinal record assembly, active clinical summaries and durable received/managed history with source attribution and supersession.



\textbf{Boundary}



Information administration and coordination preserve the authority of the responsible source, clinical actor and governed domain; they confer no additional real-world management authority.


{\footnotesize \href{../../../../continuum/strategy/capability/14748dc6117f4b6ba7c0fc554b76bd89.md}{Canonical record: 14748dc6117f4b6ba7c0fc554b76bd89}; full qualifications and provenance.}\par
\subparagraph*{II.1.3.9.15.2.1 [Feature] Active Clinical Record Access}
\phantomsection
\label{navigation:II.1.3.9.15.2.1}
\addcontentsline{toc}{subparagraph}{II.1.3.9.15.2.1 [Feature] Active Clinical Record Access}

\label{element:76b35ecac42b4781bfd7911c9931f5f1}
\index{[Feature] Active Clinical Record Access (II.1.3.9.15.2.1)}
{\small Capability / Feature · Canonical · current}\par

Provide low-latency, active distributed access to consolidated patient clinical summaries.



Active access supports responsive use of consolidated clinical summaries in ongoing care. The working representation remains reconstructable active state; speed and availability do not promote it to authoritative durable truth when persistence or convergence is unavailable.


{\footnotesize \href{../../../../continuum/strategy/capability/feature/76b35ecac42b4781bfd7911c9931f5f1.md}{Canonical record: 76b35ecac42b4781bfd7911c9931f5f1}; full qualifications and provenance.}\par
\subparagraph*{II.1.3.9.15.2.2 [Feature] Clinical Information Historical Reconstruction}
\phantomsection
\label{navigation:II.1.3.9.15.2.2}
\addcontentsline{toc}{subparagraph}{II.1.3.9.15.2.2 [Feature] Clinical Information Historical Reconstruction}

\label{element:838728ef4fd04b4da0c28dc4fce0cfe9}
\index{[Feature] Clinical Information Historical Reconstruction (II.1.3.9.15.2.2)}
{\small Capability / Feature · Proposed · current}\par

Reconstruct what information Harmonia received, knew or managed at a relevant historical basis.



\textbf{Overview}



Current state and historical operational truth answer different questions. Reconstruction retains superseded and corrected knowledge without rewriting the earlier record.



\textbf{Boundary}



Does not infer that historical information was clinically true or that a present source state applied in the past.


{\footnotesize \href{../../../../continuum/strategy/capability/feature/838728ef4fd04b4da0c28dc4fce0cfe9.md}{Canonical record: 838728ef4fd04b4da0c28dc4fce0cfe9}; full qualifications and provenance.}\par
\subparagraph*{II.1.3.9.15.2.3 [Feature] Durable Clinical Record Preservation}
\phantomsection
\label{navigation:II.1.3.9.15.2.3}
\addcontentsline{toc}{subparagraph}{II.1.3.9.15.2.3 [Feature] Durable Clinical Record Preservation}

\label{element:aba48dbf729d477aa87e112725dc1511}
\index{[Feature] Durable Clinical Record Preservation (II.1.3.9.15.2.3)}
{\small Capability / Feature · Canonical · current}\par

Durably maintain governed, vendor-neutral longitudinal clinical representation and preservation with verifiable integrity controls.



Durable preservation retains the governed longitudinal representation and the history of what Harmonia received or managed. Corrections establish changed knowledge without retrospectively rewriting that history, and preservation neither originates clinical facts nor automatically establishes legal qualification of the record.


{\footnotesize \href{../../../../continuum/strategy/capability/feature/aba48dbf729d477aa87e112725dc1511.md}{Canonical record: aba48dbf729d477aa87e112725dc1511}; full qualifications and provenance.}\par
\subparagraph*{II.1.3.9.15.2.4 [Feature] Longitudinal Health Record (LHR) Assembly}
\phantomsection
\label{navigation:II.1.3.9.15.2.4}
\addcontentsline{toc}{subparagraph}{II.1.3.9.15.2.4 [Feature] Longitudinal Health Record (LHR) Assembly}

\label{element:c9c1c7a58cb5470db0141329ecc349f4}
\index{[Feature] Longitudinal Health Record (LHR) Assembly (II.1.3.9.15.2.4)}
{\small Capability / Feature · Canonical · current}\par

Assemble disparate encounters, diagnostic reports, medications, and clinical documents into a unified, vendor-neutral longitudinal timeline.



Longitudinal assembly relates contributions from disparate systems into a usable chronology. The unified view preserves source and subject context rather than turning consolidation into originating clinical authority or erasing differences in the standing of individual assertions.


{\footnotesize \href{../../../../continuum/strategy/capability/feature/c9c1c7a58cb5470db0141329ecc349f4.md}{Canonical record: c9c1c7a58cb5470db0141329ecc349f4}; full qualifications and provenance.}\par
\paragraph*{II.1.3.9.15.3 [CT3] Primary Care Coordination}
\phantomsection
\label{navigation:II.1.3.9.15.3}
\addcontentsline{toc}{paragraph}{II.1.3.9.15.3 [CT3] Primary Care Coordination}

\label{element:f80e73ee91624ce5baa11ae15480fd35}
\index{[CT3] Primary Care Coordination (II.1.3.9.15.3)}
{\small Capability · Canonical · current}\par

Definition unresolved in the baseline. No definition has been synthesised from the name.



\medskip\hrule\medskip

{\footnotesize \href{../../../../continuum/strategy/capability/f80e73ee91624ce5baa11ae15480fd35.md}{Canonical record: f80e73ee91624ce5baa11ae15480fd35}; full qualifications and provenance.}\par
\subparagraph*{II.1.3.9.15.3.1 [Feature] General Practice Shared Record Ingestion}
\phantomsection
\label{navigation:II.1.3.9.15.3.1}
\addcontentsline{toc}{subparagraph}{II.1.3.9.15.3.1 [Feature] General Practice Shared Record Ingestion}

\label{element:869b77462ece488788c5a20598e7d587}
\index{[Feature] General Practice Shared Record Ingestion (II.1.3.9.15.3.1)}
{\small Capability / Feature · Canonical · current}\par

Ingest shared health summaries, encounter notes, and immunization records from primary care practice management systems.



Shared-record ingestion brings general-practice contributions into cross-provider continuity. The summary or encounter note retains its clinical source while becoming available alongside regional information; consolidation does not transfer the GP's professional authority to the platform.


{\footnotesize \href{../../../../continuum/strategy/capability/feature/869b77462ece488788c5a20598e7d587.md}{Canonical record: 869b77462ece488788c5a20598e7d587}; full qualifications and provenance.}\par
\subparagraph*{II.1.3.9.15.3.2 [Feature] Primary Care Notification Dispatch}
\phantomsection
\label{navigation:II.1.3.9.15.3.2}
\addcontentsline{toc}{subparagraph}{II.1.3.9.15.3.2 [Feature] Primary Care Notification Dispatch}

\label{element:56739765fcfd4084b38d4550593ee74d}
\index{[Feature] Primary Care Notification Dispatch (II.1.3.9.15.3.2)}
{\small Capability / Feature · Canonical · current}\par

Dispatch emergency department arrival, acute admission, and discharge notifications to the patient's registered general practitioner.



Primary-care notifications inform the registered GP of significant acute-care events. This outward notification responsibility is distinct from querying a chronic-care timeline: the latter shares context but has no established equivalence to dispatching the required arrival, admission or discharge notice.


{\footnotesize \href{../../../../continuum/strategy/capability/feature/56739765fcfd4084b38d4550593ee74d.md}{Canonical record: 56739765fcfd4084b38d4550593ee74d}; full qualifications and provenance.}\par
\paragraph*{II.1.3.9.15.4 [CT3] Remote Patient Monitoring}
\phantomsection
\label{navigation:II.1.3.9.15.4}
\addcontentsline{toc}{paragraph}{II.1.3.9.15.4 [CT3] Remote Patient Monitoring}

\label{element:2d1addc7dd1d4821b71a11c692d35b27}
\index{[CT3] Remote Patient Monitoring (II.1.3.9.15.4)}
{\small Capability · Canonical · current}\par

Definition unresolved in the baseline. No definition has been synthesised from the name.



\medskip\hrule\medskip

{\footnotesize \href{../../../../continuum/strategy/capability/2d1addc7dd1d4821b71a11c692d35b27.md}{Canonical record: 2d1addc7dd1d4821b71a11c692d35b27}; full qualifications and provenance.}\par
\subparagraph*{II.1.3.9.15.4.1 [Feature] Remote Telemetry Ingestion}
\phantomsection
\label{navigation:II.1.3.9.15.4.1}
\addcontentsline{toc}{subparagraph}{II.1.3.9.15.4.1 [Feature] Remote Telemetry Ingestion}

\label{element:969dde193fab4c5d82d7dc210c7f5a71}
\index{[Feature] Remote Telemetry Ingestion (II.1.3.9.15.4.1)}
{\small Capability / Feature · Canonical · current}\par

Ingest home patient monitoring data, biometric telemetry, and patient-reported outcome measures (PROMs).



Remote telemetry ingestion accepts measurements and patient-reported contributions from outside the usual care site. The information can inform subsequent monitoring and interpretation, but acquisition alone does not establish clinical significance or authorise an intervention.


{\footnotesize \href{../../../../continuum/strategy/capability/feature/969dde193fab4c5d82d7dc210c7f5a71.md}{Canonical record: 969dde193fab4c5d82d7dc210c7f5a71}; full qualifications and provenance.}\par
\subparagraph*{II.1.3.9.15.4.2 [Feature] Virtual Care Session Context Binding}
\phantomsection
\label{navigation:II.1.3.9.15.4.2}
\addcontentsline{toc}{subparagraph}{II.1.3.9.15.4.2 [Feature] Virtual Care Session Context Binding}

\label{element:c25d66fd6aff4299857eadc7c0243544}
\index{[Feature] Virtual Care Session Context Binding (II.1.3.9.15.4.2)}
{\small Capability / Feature · Canonical · current}\par

Associate synchronous virtual/telehealth consultations with the longitudinal health record and active care episode.



Virtual-session binding associates a consultation with the enduring record and the relevant episode. It preserves the clinical context across delivery modes without turning the communication session itself into proof of clinical conclusions or adequacy of care.


{\footnotesize \href{../../../../continuum/strategy/capability/feature/c25d66fd6aff4299857eadc7c0243544.md}{Canonical record: c25d66fd6aff4299857eadc7c0243544}; full qualifications and provenance.}\par
\subsubsection*{II.1.3.9.16 [CT2] Presentation Services}
\phantomsection
\label{navigation:II.1.3.9.16}
\addcontentsline{toc}{subsubsection}{II.1.3.9.16 [CT2] Presentation Services}

{\small Historical predecessor: R.2.208 \hyperref[element:c89e7c8b277247d9965fdb534aad2d1f]{[CT2] Presentation Services}; Canonical ID: \texttt{\seqsplit{c89e7c8b277247d9965fdb534aad2d1f}}. Superseded historical identity — direct correspondence · superseded.}\par
\label{element:e9b6ba1917ba44cca7466444d0836c98}
\index{[CT2] Presentation Services (II.1.3.9.16)}
{\small Capability · Canonical · current}\par

Presentation Services supplies contextual human interaction through clinical timelines, directory views and operational dashboards. It makes governed information usable by people while consuming the identity, security and clinical context established elsewhere. The responsibility is the interaction and presentation experience rather than a new authority for the information shown.



\textbf{Overview}



Presentation Services supplies contextual human interaction through clinical timelines, directory views and operational dashboards. It makes governed information usable by people while consuming the identity, security and clinical context established elsewhere. The responsibility is the interaction and presentation experience rather than a new authority for the information shown.



\textbf{Intent}



Presentation Services supplies contextual human interaction through clinical timelines, directory views and operational dashboards. It makes governed information usable by people while consuming the identity, security and clinical context established elsewhere. The responsibility is the interaction and presentation experience rather than a new authority for the information shown.



\textbf{Boundary}



Information administration does not confer authority over the represented person, organisation, service, place, device or clinical activity.



\medskip\hrule\medskip

{\footnotesize \href{../../../../continuum/strategy/capability/e9b6ba1917ba44cca7466444d0836c98.md}{Canonical record: e9b6ba1917ba44cca7466444d0836c98}; full qualifications and provenance.}\par
\paragraph*{II.1.3.9.16.1 [CT3] Contextual Information Presentation}
\phantomsection
\label{navigation:II.1.3.9.16.1}
\addcontentsline{toc}{paragraph}{II.1.3.9.16.1 [CT3] Contextual Information Presentation}

{\small Historical predecessor: R.2.207 \hyperref[element:e349d5f6a8774bbe8d150f9c51af770c]{[CT3] Contextual Information Presentation}; Canonical ID: \texttt{\seqsplit{e349d5f6a8774bbe8d150f9c51af770c}}. Superseded historical identity — direct correspondence · superseded.}\par
\label{element:9f900c53ae56488cad69e53c2205e898}
\index{[CT3] Contextual Information Presentation (II.1.3.9.16.1)}
{\small Capability · Proposed · current}\par

Present governed longitudinal clinical and operational information in its applicable human interaction context.



\textbf{Overview}



Present governed longitudinal clinical and operational information in its applicable human interaction context. The reviewed hierarchy assigns this enduring ability to its stated business-enabling context.



\textbf{Intent}



Present governed longitudinal clinical and operational information in its applicable human interaction context.



\textbf{Boundary}



Information administration and coordination preserve the authority of the responsible source, clinical actor and governed domain; they confer no additional real-world management authority.


{\footnotesize \href{../../../../continuum/strategy/capability/9f900c53ae56488cad69e53c2205e898.md}{Canonical record: 9f900c53ae56488cad69e53c2205e898}; full qualifications and provenance.}\par
\subparagraph*{II.1.3.9.16.1.1 [Feature] Longitudinal Clinical \& Operational Presentation}
\phantomsection
\label{navigation:II.1.3.9.16.1.1}
\addcontentsline{toc}{subparagraph}{II.1.3.9.16.1.1 [Feature] Longitudinal Clinical \& Operational Presentation}

\label{element:e82eaafa2e4046f48d4c86e15f282a0b}
\index{[Feature] Longitudinal Clinical \& Operational Presentation (II.1.3.9.16.1.1)}
{\small Capability / Feature · Canonical · current}\par

Deliver contextual clinical timelines, directory administration views, and operational management dashboards.



Contextual presentation makes clinical timelines, directories and operational information usable for people. It consumes established identity and security context while preserving the originating authority of what is shown; a user interface does not become an alternative persistence or policy boundary.


{\footnotesize \href{../../../../continuum/strategy/capability/feature/e82eaafa2e4046f48d4c86e15f282a0b.md}{Canonical record: e82eaafa2e4046f48d4c86e15f282a0b}; full qualifications and provenance.}\par
\subsubsection*{II.1.3.9.17 [CT2] Semantic \& Reference Knowledge Enablement [superseded]}
\phantomsection
\label{navigation:II.1.3.9.17}
\addcontentsline{toc}{subsubsection}{II.1.3.9.17 [CT2] Semantic \& Reference Knowledge Enablement [superseded]}

\label{element:92a71292710c41b29775d19b98c197d6}
\index{[CT2] Semantic \& Reference Knowledge Enablement (II.1.3.9.17)}
{\small Capability · Proposed · superseded}\par
\href{../../../../continuum/strategy/capability/92a71292710c41b29775d19b98c197d6.md}{Historical canonical definition and provenance}.\par
\subsubsection*{II.1.3.9.18 [CT2] Workflow \& Activity Coordination}
\phantomsection
\label{navigation:II.1.3.9.18}
\addcontentsline{toc}{subsubsection}{II.1.3.9.18 [CT2] Workflow \& Activity Coordination}

{\small Historical predecessor: R.2.213 \hyperref[element:4422d865978c4353ad575759c6026d27]{[CT2] Workflow \& Activity Coordination}; Canonical ID: \texttt{\seqsplit{4422d865978c4353ad575759c6026d27}}. Superseded historical identity — direct correspondence · superseded.}\par
\label{element:c724f3d38bfa4567983a11ad22095d61}
\index{[CT2] Workflow \& Activity Coordination (II.1.3.9.18)}
{\small Capability · Proposed · current}\par

Coordinate tasks entering or leaving the Harmonia execution boundary and work undertaken by other systems, with responsibility/destination routing, external status/traceability and governance of the completion bridge.



\textbf{Overview}



Coordinate tasks entering or leaving the Harmonia execution boundary and work undertaken by other systems, with responsibility/destination routing, external status/traceability and governance of the completion bridge. The reviewed hierarchy assigns this enduring ability to its stated business-enabling context.



\textbf{Intent}



Coordinate tasks entering or leaving the Harmonia execution boundary and work undertaken by other systems, with responsibility/destination routing, external status/traceability and governance of the completion bridge.



\textbf{Boundary}



External means outside Harmonia execution, potentially inside the same health organisation. Work coordination includes completion responsibilities and traceability; it is not generic interoperability or message routing.


{\footnotesize \href{../../../../continuum/strategy/capability/c724f3d38bfa4567983a11ad22095d61.md}{Canonical record: c724f3d38bfa4567983a11ad22095d61}; full qualifications and provenance.}\par
\paragraph*{II.1.3.9.18.1 [CT3] Activity Routing Management}
\phantomsection
\label{navigation:II.1.3.9.18.1}
\addcontentsline{toc}{paragraph}{II.1.3.9.18.1 [CT3] Activity Routing Management}

{\small Historical predecessor: R.2.209 \hyperref[element:243ea93e7b7448fe97dc74f90bfc8946]{[CT3] Activity Routing Management}; Canonical ID: \texttt{\seqsplit{243ea93e7b7448fe97dc74f90bfc8946}}. Superseded historical identity — direct correspondence · superseded.}\par
\label{element:27411c143e6e41ffba1af139fbbcd7e2}
\index{[CT3] Activity Routing Management (II.1.3.9.18.1)}
{\small Capability · Proposed · current}\par

Manage responsibility and destination routing for activity crossing the Harmonia execution boundary, respecting the intended work and completion responsibility.



\textbf{Overview}



Manage responsibility and destination routing for activity crossing the Harmonia execution boundary, respecting the intended work and completion responsibility. The reviewed hierarchy assigns this enduring ability to its stated business-enabling context.



\textbf{Intent}



Manage responsibility and destination routing for activity crossing the Harmonia execution boundary, respecting the intended work and completion responsibility.



\textbf{Boundary}



Information administration and coordination preserve the authority of the responsible source, clinical actor and governed domain; they confer no additional real-world management authority.


{\footnotesize \href{../../../../continuum/strategy/capability/27411c143e6e41ffba1af139fbbcd7e2.md}{Canonical record: 27411c143e6e41ffba1af139fbbcd7e2}; full qualifications and provenance.}\par
\paragraph*{II.1.3.9.18.2 [CT3] External Task Coordination}
\phantomsection
\label{navigation:II.1.3.9.18.2}
\addcontentsline{toc}{paragraph}{II.1.3.9.18.2 [CT3] External Task Coordination}

{\small Historical predecessor: R.2.210 \hyperref[element:32be309d65594592bac4d7b068def0d2]{[CT3] External Task Coordination}; Canonical ID: \texttt{\seqsplit{32be309d65594592bac4d7b068def0d2}}. Superseded historical identity — direct correspondence · superseded.}\par
\label{element:ff3447c929724245a13ee864874b5561}
\index{[CT3] External Task Coordination (II.1.3.9.18.2)}
{\small Capability · Proposed · current}\par

Coordinate work entering or leaving the Harmonia execution boundary with the system or actor responsible for undertaking and completing that work.



\textbf{Overview}



Coordinate work entering or leaving the Harmonia execution boundary with the system or actor responsible for undertaking and completing that work. The reviewed hierarchy assigns this enduring ability to its stated business-enabling context.



\textbf{Intent}



Coordinate work entering or leaving the Harmonia execution boundary with the system or actor responsible for undertaking and completing that work.



\textbf{Boundary}



Information administration and coordination preserve the authority of the responsible source, clinical actor and governed domain; they confer no additional real-world management authority.


{\footnotesize \href{../../../../continuum/strategy/capability/ff3447c929724245a13ee864874b5561.md}{Canonical record: ff3447c929724245a13ee864874b5561}; full qualifications and provenance.}\par
\paragraph*{II.1.3.9.18.3 [CT3] External Task Status \& Traceability Management}
\phantomsection
\label{navigation:II.1.3.9.18.3}
\addcontentsline{toc}{paragraph}{II.1.3.9.18.3 [CT3] External Task Status \& Traceability Management}

{\small Historical predecessor: R.2.211 \hyperref[element:95985d244a234571892a9e9b4024f594]{[CT3] External Task Status \& Traceability Management}; Canonical ID: \texttt{\seqsplit{95985d244a234571892a9e9b4024f594}}. Superseded historical identity — direct correspondence · superseded.}\par
\label{element:b25bb04a8aa94dcfac105a646878ccfc}
\index{[CT3] External Task Status \& Traceability Management (II.1.3.9.18.3)}
{\small Capability · Proposed · current}\par

Manage the observable status and attributable traceability of work as it progresses through systems outside the Harmonia execution boundary.



\textbf{Overview}



Manage the observable status and attributable traceability of work as it progresses through systems outside the Harmonia execution boundary. The reviewed hierarchy assigns this enduring ability to its stated business-enabling context.



\textbf{Intent}



Manage the observable status and attributable traceability of work as it progresses through systems outside the Harmonia execution boundary.



\textbf{Boundary}



Information administration and coordination preserve the authority of the responsible source, clinical actor and governed domain; they confer no additional real-world management authority.


{\footnotesize \href{../../../../continuum/strategy/capability/b25bb04a8aa94dcfac105a646878ccfc.md}{Canonical record: b25bb04a8aa94dcfac105a646878ccfc}; full qualifications and provenance.}\par
\paragraph*{II.1.3.9.18.4 [CT3] Work Boundary Governance}
\phantomsection
\label{navigation:II.1.3.9.18.4}
\addcontentsline{toc}{paragraph}{II.1.3.9.18.4 [CT3] Work Boundary Governance}

{\small Historical predecessor: R.2.212 \hyperref[element:b4d14abf930a4230a7f5b93e9ce3df91]{[CT3] Work Boundary Governance}; Canonical ID: \texttt{\seqsplit{b4d14abf930a4230a7f5b93e9ce3df91}}. Superseded historical identity — direct correspondence · superseded.}\par
\label{element:35b66a0b38db4568849ecb842ffdcbf3}
\index{[CT3] Work Boundary Governance (II.1.3.9.18.4)}
{\small Capability · Proposed · current}\par

Govern boundary conditions and completion responsibilities of the bridge between Harmonia and another system undertaking work.



\textbf{Overview}



Govern boundary conditions and completion responsibilities of the bridge between Harmonia and another system undertaking work. The reviewed hierarchy assigns this enduring ability to its stated business-enabling context.



\textbf{Intent}



Govern boundary conditions and completion responsibilities of the bridge between Harmonia and another system undertaking work.



\textbf{Boundary}



Information administration and coordination preserve the authority of the responsible source, clinical actor and governed domain; they confer no additional real-world management authority.


{\footnotesize \href{../../../../continuum/strategy/capability/35b66a0b38db4568849ecb842ffdcbf3.md}{Canonical record: 35b66a0b38db4568849ecb842ffdcbf3}; full qualifications and provenance.}\par
\chapter*{II.2 Courses of Action}
\phantomsection
\label{navigation:II.2}
\addcontentsline{toc}{chapter}{II.2 Courses of Action}

\section*{II.2.1 Boundary Membrane Sovereignty}
\phantomsection
\label{navigation:II.2.1}
\addcontentsline{toc}{section}{II.2.1 Boundary Membrane Sovereignty}

\label{element:25d8509146fd4f2794a17f4a5253b7a3}
\index{Boundary Membrane Sovereignty (II.2.1)}
{\small CourseOfAction · Canonical · current}\par

\par\medskip\textbf{Strategic Approach}\par


\begin{quote}


\textbf{Preserve strict separation between external standards-compliant exchange representations and internal Harmonia-governed operational semantics.}



\end{quote}


Harmonia acts as an integration environment bridging heterogeneous external healthcare systems and internal integration services. This course of action establishes that:
- External exchange contracts (such as FHIR REST interactions, MLLP message streams, or directory protocols) govern only the interaction semantics at the enterprise boundary.
- External or standards-based representations MAY remain native Harmonia information representations where appropriate. Their use SHALL NOT cause an external exchange contract or representation to define Harmonia-private management, execution, governance or coordination semantics.
- Egress across the boundary terminates Harmonia's governance of the emitted representation.



\par\medskip\textbf{Motivational Grounding \& Traceability}\par


\begin{itemize}


\item \textbf{Architectural Axioms}:

\item \texttt{\seqsplit{AX-02}} (Boundary Membrane): External representations are boundary projections; internal semantics are sovereign.

\item \texttt{\seqsplit{AX-03}} (Native Standards): Standards govern interoperability at the boundary; standards are not distorted to fit private operational models.

\item \texttt{\seqsplit{AX-13}} (Explicit Management Boundary): Egress terminates platform management; internal operational metadata does not leak across the boundary.

\item \texttt{\seqsplit{AX-14}} (Preserve Distinctions): Distinctions between external exchange syntax and internal operational meaning are preserved.

\item \textbf{Strategic Drivers \& Goals}: Interoperability Mandates, Vendor Independence, Unified Ingress/Egress Governance.

\end{itemize}


\par\medskip\textbf{Capability Realisation \& Configuration}\par


\begin{itemize}


\item \textbf{Primary Capabilities}: Shapes \texttt{\seqsplit{EC-08}} (Interoperability \& Exchange) and \texttt{\seqsplit{EC-13}} (Semantic Governance \& Conformance).

\item \textbf{Component Boundary Impact}: Informs the boundary responsibility of \textbf{Pylai}, establishing it as the standards-facing membrane rather than an internal processing engine or generic transport layer.

\end{itemize}


\par\medskip\textbf{Downstream Direction (Domains 03–13)}\par


\begin{itemize}


\item Application Architecture (Domain 05) and Integration Architecture (Domain 06) must implement fail-closed boundary projections, ensuring private coordination metadata (e.g., distributed transaction IDs, internal routing tags) is never emitted externally.

\item Information Architecture (Domain 04) preserves conceptual meaning and responsibility independently of external contracts. Native standards representations remain permitted under AX-03; no parallel information model is required merely to accommodate Harmonia-private management concerns.

\end{itemize}

{\small Recorded unresolved matters: canonicalName; exact scope remains in the canonical record.}\par
{\footnotesize \href{../../../../continuum/strategy/course-of-action/25d8509146fd4f2794a17f4a5253b7a3.md}{Canonical record: 25d8509146fd4f2794a17f4a5253b7a3}; full qualifications and provenance.}\par
\section*{II.2.2 Collaborative Reusable Capability Realisation}
\phantomsection
\label{navigation:II.2.2}
\addcontentsline{toc}{section}{II.2.2 Collaborative Reusable Capability Realisation}

\label{element:503939e1fa464b78b1d980ddedfeb3cf}
\index{Collaborative Reusable Capability Realisation (II.2.2)}
{\small CourseOfAction · Canonical · current}\par

\par\medskip\textbf{Strategic Approach}\par


\begin{quote}


\textbf{Realize bounded reusable platform capabilities (context, policy, provenance, EC-12 Operational Assurance) collaboratively across participating components rather than through centralized, bottlenecked runtime services.}



\end{quote}


Harmonia requires pervasive platform capabilities—such as security policy evaluation, contextual propagation, operational assurance, and provenance tracking—across all transactions. This course of action establishes that:
- \textbf{Reusable Capability \$\textbackslash{}neq\$ Centralised Service (Guardrail G1)}: The existence of a reusable enterprise capability does not imply that a single monolithic service must execute it centrally.
- Reusable capabilities define bounded responsibilities and common semantic rules. Participating components consume or enforce them through governed contracts without acquiring shared semantic ownership. \textbf{Distributed participation does not imply distributed responsibility. Cross-cutting concern does not imply cross-cutting responsibility.}
- Centralized choke points and synchronous runtime bottlenecks are strictly avoided, ensuring horizontal scalability and fault isolation.



\par\medskip\textbf{Motivational Grounding \& Traceability}\par


\begin{itemize}


\item \textbf{Architectural Axioms}:

\item \texttt{\seqsplit{AX-04}} (Own Semantics / Use Machinery): Maintain sovereign capability rules while delegating execution across components.

\item \texttt{\seqsplit{AX-07}} (Intrinsic Security): Security and policy evaluation must be pervasive and intrinsic to all operations, not an external proxy.

\item \texttt{\seqsplit{AX-11}} (High Availability \& Responsiveness): Eliminate single points of congestion and failure across the integration fabric.

\item \textbf{Strategic Drivers \& Goals}: High Scalability, Fault Tolerance, Zero Bottleneck Architecture, Decentralised Resilience.

\end{itemize}


\par\medskip\textbf{Capability Realisation \& Configuration}\par


\begin{itemize}


\item \textbf{Primary Capabilities}: Shapes \texttt{\seqsplit{EC-02}} (Context Management), \texttt{\seqsplit{EC-06}} (Policy \& Control), \texttt{\seqsplit{EC-07}} (Provenance \& Traceability), and \texttt{\seqsplit{EC-12}} (Operational Assurance). Assurance here means the established EC-12 operational contribution; no independent Governed Assurance or EC-14 relationship is established.

\item \textbf{Component Boundary Impact}: Prevents components like \textbf{Themis} (Policy \& Control) from becoming synchronous network choke points for every micro-operation; policy contracts are distributed and evaluated collaboratively.

\end{itemize}


\par\medskip\textbf{Downstream Direction (Domains 03–13)}\par


\begin{itemize}


\item Application Architecture (Domain 05) and Technology Architecture (Domain 07) realise reusable responsibilities through governed contracts while preserving bounded ownership, scalability and fault isolation. Existing \texttt{\seqsplit{themis-api}} and \texttt{\seqsplit{petasos-api}} packages are illustrative downstream traceability only; Strategy does not mandate API packaging, local embedding, frameworks or deployment choices.

\end{itemize}

{\small Recorded unresolved matters: canonicalName; exact scope remains in the canonical record.}\par
{\footnotesize \href{../../../../continuum/strategy/course-of-action/503939e1fa464b78b1d980ddedfeb3cf.md}{Canonical record: 503939e1fa464b78b1d980ddedfeb3cf}; full qualifications and provenance.}\par
\section*{II.2.3 Distinct Management and Durable Preservation of Information and State}
\phantomsection
\label{navigation:II.2.3}
\addcontentsline{toc}{section}{II.2.3 Distinct Management and Durable Preservation of Information and State}

\label{element:4c71b89477b34733b31e478485e89b1a}
\index{Distinct Management and Durable Preservation of Information and State (II.2.3)}
{\small CourseOfAction · Canonical · current}\par

\par\medskip\textbf{Strategic Approach}\par


\begin{quote}


\textbf{Maintain a clear separation between the governed management of Harmonia information and state and its durable preservation and recovery, allowing runtime management and persistence concerns to evolve independently without compromising information integrity.}



\end{quote}


This course of action establishes that:
- Runtime management of information, active relationships, context, and state is an operational governance concern requiring high availability, concurrent access, and responsive query navigation.
- Authoritative durable state establishment, preservation and recovery atomically establish authoritative state and authoritative version progression, preserve the durable management metadata required to interpret that state, and ensure historical retention, immutability, and state reconstruction.
- The distinction is strictly one of \textbf{architectural responsibility}, not a crude division between "ephemeral data" and "durable data". Information managed at runtime frequently represents durable healthcare concepts, but runtime management and durable preservation fulfill distinct architectural purposes.



\par\medskip\textbf{Motivational Grounding \& Traceability}\par


\begin{itemize}


\item \textbf{Architectural Axioms}:

\item \texttt{\seqsplit{AX-05}} (Active vs Durable State): Application-facing active state access/coordination and authoritative durable state/version establishment remain separate architectural concerns.

\item \texttt{\seqsplit{AX-09}} (Ephemeral Operational State): Transient execution states are reconstructable and decoupled from durable health records.

\item \texttt{\seqsplit{AX-11}} (High Availability \& Responsiveness): Information access must not be blocked by backend persistence latencies or recovery locks.

\item \texttt{\seqsplit{AX-14}} (Preserve Distinctions): Preserve the boundary between runtime information governance and storage mechanics.

\item \textbf{Strategic Drivers \& Goals}: 24/7 Clinical Operational Continuity, High-Throughput Access, Authoritative Preservation.

\end{itemize}


\par\medskip\textbf{Capability Realisation \& Configuration}\par


\begin{itemize}


\item \textbf{Primary Capabilities}: Shapes \texttt{\seqsplit{EC-03}} (Managed State \& Lifecycle) and \texttt{\seqsplit{EC-04}} (Information Management).

\item \textbf{Component Boundary Impact}: Explicitly establishes the seam between \textbf{Mneme} (application-facing access and active information/state management) and \textbf{Mnemosyne} (authoritative durable state/version establishment, preservation and recovery). Neither component subsumes the other; Mnemosyne does not own operational activity progression, workflow execution, active distributed state, Digital Twin coordination, or application-facing query interfaces. Mneme may reject or coordinate a proposed state progression before persistence; only Mnemosyne establishes the new authoritative durable state. Following authoritative commit, Mneme converges its active representation toward that state. Mneme active-state generation and Mnemosyne authoritative version remain distinct concurrency domains (\texttt{\seqsplit{AX-05}}).

\end{itemize}


\par\medskip\textbf{Downstream Direction (Domains 03–13)}\par


\begin{itemize}


\item Application Architecture (Domain 05) must forbid direct presentation/application access to database/JPA layers (enforcing Invariant 8).

\item Technology Architecture (Domain 07) must allow distributed in-memory data fabrics (for runtime state) and durable relational/document stores (for preservation) to be configured and scaled independently.

\end{itemize}

{\small Recorded unresolved matters: canonicalName; exact scope remains in the canonical record.}\par
{\footnotesize \href{../../../../continuum/strategy/course-of-action/4c71b89477b34733b31e478485e89b1a.md}{Canonical record: 4c71b89477b34733b31e478485e89b1a}; full qualifications and provenance.}\par
\section*{II.2.4 Entity-Centred Operational Coordination}
\phantomsection
\label{navigation:II.2.4}
\addcontentsline{toc}{section}{II.2.4 Entity-Centred Operational Coordination}

\label{element:f601fe1f39a3447f8e9fb2cc4ffe9dee}
\index{Entity-Centred Operational Coordination (II.2.4)}
{\small CourseOfAction · Canonical · current}\par

\par\medskip\textbf{Strategic Approach}\par


\begin{quote}


\textbf{Coordinate governed information, state and operational activity around real-world healthcare entities where those entities are operationally significant in their own right and require active management, without requiring every managed entity to maintain a permanently active execution construct.}



\end{quote}


Healthcare operations naturally center around authentic real-world entities: patients, practitioners, care teams, wards, and beds. This course of action establishes that:
- Harmonia organizes the coordination of state and ongoing activities around the operational identity of the real-world entity itself.
- \textbf{The Digital Twin is the architectural construct} through which this entity-centred coordination strategy is realised. A Twin coordinates information, state, and activity across the Mneme/Ponos seam for a specific real-world entity.
- \textbf{Demand-Driven Lifecycle}: A managed entity does not require a permanently running execution thread or process. Twins are activated when entity-centred operational activity or active monitoring is required, and become quiescent when work concludes.
- \textbf{Representation \$\textbackslash{}neq\$ Twin}: Simply possessing a database record or FHIR resource does not justify or constitute a Digital Twin. Coordination is justified only when the entity is operationally active and evolving in state.



\par\medskip\textbf{Motivational Grounding \& Traceability}\par


\begin{itemize}


\item \textbf{Architectural Axioms}:

\item \texttt{\seqsplit{AX-01}} (Health-Information Centric): Architecture aligns with genuine healthcare entities and clinical relationships.

\item \texttt{\seqsplit{AX-11}} (High Availability \& Responsiveness): High concurrency is achieved by isolating entity coordination without central serialisation bottlenecks.

\item \texttt{\seqsplit{AX-16}} (Activity \& State Progress Together): Entity state and active operational workflows are coordinated dynamically.

\item \textbf{Strategic Drivers \& Goals}: Holistic Patient-Centric Care, High Concurrency without Serialization, Real-Time Healthcare Coordination.

\end{itemize}


\par\medskip\textbf{Capability Realisation \& Configuration}\par


\begin{itemize}


\item \textbf{Primary Capabilities}: Shapes \texttt{\seqsplit{EC-01}} (Managed Entity \& Relationship), \texttt{\seqsplit{EC-02}} (Context Management), and \texttt{\seqsplit{EC-10}} (Activity \& Execution).

\item \textbf{Component Boundary Impact}: Validates the \textbf{Digital Twin} as an active management construct bridging the \textbf{Mneme ↔ Ponos} boundary, while asserting that the Twin is an architectural construct rather than an independent deployable platform component.

\end{itemize}


\par\medskip\textbf{Downstream Direction (Domains 03–13)}\par


\begin{itemize}


\item Application Architecture (Domain 05) must model Twin archetypes using composite context (combining multiple resources such as Practitioner, PractitionerRole, and Location) rather than equating twins to single FHIR resources or permanent background threads.

\end{itemize}

{\small Recorded unresolved matters: canonicalName; exact scope remains in the canonical record.}\par
{\footnotesize \href{../../../../continuum/strategy/course-of-action/f601fe1f39a3447f8e9fb2cc4ffe9dee.md}{Canonical record: f601fe1f39a3447f8e9fb2cc4ffe9dee}; full qualifications and provenance.}\par
\section*{II.2.5 Governed Asynchronous Activity Progression}
\phantomsection
\label{navigation:II.2.5}
\addcontentsline{toc}{section}{II.2.5 Governed Asynchronous Activity Progression}

\label{element:a8ab09ee29454376bd45ba90b320d6ec}
\index{Governed Asynchronous Activity Progression (II.2.5)}
{\small CourseOfAction · Canonical · current}\par

\par\medskip\textbf{Strategic Approach}\par


\begin{quote}


\textbf{Progress multi-stage operational activities through explicit, observable unit-of-work transitions coordinated with entity state, separating execution from boundary exchange.}



\end{quote}


Complex healthcare integrations (e.g., closed-loop diagnostic ordering, multi-destination clinical document fan-out) span multiple distributed systems and cannot rely on brittle synchronous distributed transactions. This course of action establishes that:
- Work progresses through discrete, observable units of work whose lifecycle state transitions are explicitly tracked and auditable.
- Operational progression progresses in lockstep with governed entity state (\texttt{\seqsplit{AX-16}}), preserving explicit uncertainty (\texttt{\seqsplit{AX-15}}) when external responses are pending or failed.
- Execution progression is strictly decoupled from boundary transport and exchange mechanics: activity execution determines \emph{that} an external interaction is needed, but does not own transport protocols or wire-level retries.



\par\medskip\textbf{Motivational Grounding \& Traceability}\par


\begin{itemize}


\item \textbf{Architectural Axioms}:

\item \texttt{\seqsplit{AX-10}} (Distribution \& Normal Failure): Systems fail normally; activity progression must be resilient, idempotent, and restartable.

\item \texttt{\seqsplit{AX-15}} (Preserve Uncertainty): Never assume success or failure; pending, timed-out, or ambiguous outcomes must be explicitly modeled.

\item \texttt{\seqsplit{AX-16}} (Activity \& State Progress Together): Operational activity and managed information state advance in coordinated lockstep.

\item \textbf{Strategic Drivers \& Goals}: Resilient Regional Integration, Asynchronous Clinical Decoupling, Bounded Workflow Recovery.

\end{itemize}


\par\medskip\textbf{Capability Realisation \& Configuration}\par


\begin{itemize}


\item \textbf{Primary Capabilities}: Shapes \texttt{\seqsplit{EC-03}} (Managed State \& Lifecycle), \texttt{\seqsplit{EC-09}} (Event \& Subscription), and \texttt{\seqsplit{EC-10}} (Activity \& Execution).

\item \textbf{Component Boundary Impact}: Grounds the responsibility of \textbf{Ponos} as the execution engine for governed units of work, while strictly enforcing Guardrail G4 (Ponos progresses execution, but does not acquire transport adapters or wire protocols).

\end{itemize}


\par\medskip\textbf{Downstream Direction (Domains 03–13)}\par


\begin{itemize}


\item Application Architecture (Domain 05) and Integration Architecture (Domain 06) realise governed work context, recoverable activity handoffs and explicit per-destination progression according to the applicable acceptance and delivery contracts (preserving Invariants 4 and 5). Established downstream constructs such as Pragma envelopes and Petasos messaging illustrate traceability; their APIs, queues, packaging and runtime mechanisms are not selected by this strategic responsibility.

\end{itemize}

{\small Recorded unresolved matters: canonicalName; exact scope remains in the canonical record.}\par
{\footnotesize \href{../../../../continuum/strategy/course-of-action/a8ab09ee29454376bd45ba90b320d6ec.md}{Canonical record: a8ab09ee29454376bd45ba90b320d6ec}; full qualifications and provenance.}\par
\section*{II.2.6 Meaning-Centric Provenance and Traceability}
\phantomsection
\label{navigation:II.2.6}
\addcontentsline{toc}{section}{II.2.6 Meaning-Centric Provenance and Traceability}

\label{element:fd00bdbdc2234a208c5045aa69279db8}
\index{Meaning-Centric Provenance and Traceability (II.2.6)}
{\small CourseOfAction · Canonical · current}\par

\par\medskip\textbf{Strategic Approach}\par


\begin{quote}


\textbf{Preserve provenance, authority, attribution and traceability for information-significant and business-significant actions and state changes, while avoiding unnecessary elevation of transient operational mechanics into enduring business evidence.}



\end{quote}


Healthcare integration requires undeniable accountability without overwhelming storage or logging systems with low-level operational noise. This course of action establishes that:
- Harmonia SHALL preserve provenance, audit and operational evidence where required by governing obligations, architectural significance, accountability, reconstruction, security or assurance needs. Routine technical activity SHALL NOT automatically require durable evidentiary recording.
- Required semantic attribution records who, what, when, why and under whose authority. Telemetry, operational facts, provenance, audit and assurance evidence remain distinct; being evidence is contextual, not intrinsic to every recorded fact.
- Transient operational mechanics (e.g., thread switches, socket retries, queue polling, network pinging) represent operational telemetry, not business or clinical evidence.
- Strategy sets the architectural requirement for semantic attribution and non-repudiation without dictating specific implementation mechanisms (such as cryptographic signatures, write-once ledgers, or specific storage engines) at this layer.



\par\medskip\textbf{Motivational Grounding \& Traceability}\par


\begin{itemize}


\item \textbf{Architectural Axioms}:

\item \texttt{\seqsplit{AX-06}} (Explicit Authority): Information authority and credibility remain explicit and independent of technical transport or persistence. Required attribution preserves the applicable source or actor authority without making every access evaluation durable evidence.

\item \texttt{\seqsplit{AX-07}} (Intrinsic Security): Security, authorization, and provenance are embedded within operations, not bolted on.

\item \texttt{\seqsplit{AX-08}} (Meaning over Machinery): Record semantic and business meaning, not transient technical machinery.

\item \texttt{\seqsplit{AX-14}} (Preserve Distinctions): Distinguish enduring clinical/business evidence from transient technical logs.

\item \textbf{Strategic Drivers \& Goals}: Regulatory Accountability, Clinical Governance, Legal Auditability, Evidentiary Integrity.

\end{itemize}


\par\medskip\textbf{Capability Realisation \& Configuration}\par


\begin{itemize}


\item \textbf{Primary Capabilities}: Shapes \texttt{\seqsplit{EC-06}} (Policy \& Control) and \texttt{\seqsplit{EC-07}} (Provenance \& Traceability).

\item \textbf{Component Boundary Impact}: Governs how all logical components emit evidence, ensuring that audit trails capture governed business assertions rather than unmasked PHI or low-level machine noise (supporting Invariants 6 and 7).

\end{itemize}


\par\medskip\textbf{Downstream Direction (Domains 03–13)}\par


\begin{itemize}


\item Security Architecture (Domain 08) and Information Architecture (Domain 04) define the governed audit and provenance representations (with Kleio as established downstream evidence traceability), preserving PHI-safe capture and verifiable integrity according to applicable obligations. Strategy does not select a verification technology.

\end{itemize}

{\small Recorded unresolved matters: canonicalName; exact scope remains in the canonical record.}\par
{\footnotesize \href{../../../../continuum/strategy/course-of-action/fd00bdbdc2234a208c5045aa69279db8.md}{Canonical record: fd00bdbdc2234a208c5045aa69279db8}; full qualifications and provenance.}\par
\section*{II.2.7 Implementation Guardrails}
\phantomsection
\label{navigation:II.2.7}
\addcontentsline{toc}{section}{II.2.7 Implementation Guardrails}

\subsection*{II.2.7.1 [ImplementationGuideline] Default Development Test and POC Identity Provider}
\phantomsection
\label{navigation:II.2.7.1}
\addcontentsline{toc}{subsection}{II.2.7.1 [ImplementationGuideline] Default Development Test and POC Identity Provider}

\label{element:f8874dd3bc5c4b68b1ece4eb2859d1d9}
\index{[ImplementationGuideline] Default Development Test and POC Identity Provider (II.2.7.1)}
{\small CourseOfAction / ImplementationGuideline · Adjudicated Domain05 direction — U12/U15/U16; holistic result acceptance pending · current}\par

\begin{quote}


\textbf{Keycloak + PostgreSQL SHALL provide the default FOSS Identity Provider for Harmonia development, test and POC deployments. Keycloak remains an external, replaceable implementation component and SHALL NOT introduce Keycloak-specific dependencies into Harmonia's Identity.Authentication architecture.}



\end{quote}


This assertion is normalised from the governed DesignNote under U12/U15/U16 and the ARB release Praxis. Original candidate/release qualifications and separately deferred implementation details remain in force.



Architectural subject: \hyperref[element:e4998b0a132c41afbe610b11869c3297]{SEC-001 — Trusted Authentication}.


{\small Recorded unresolved matters: canonicalName; exact scope remains in the canonical record.}\par
{\footnotesize \href{../../../../continuum/strategy/course-of-action/implementation-guideline/f8874dd3bc5c4b68b1ece4eb2859d1d9.md}{Canonical record: f8874dd3bc5c4b68b1ece4eb2859d1d9}; full qualifications and provenance.}\par
\chapter*{II.3 Resources}
\phantomsection
\label{navigation:II.3}
\addcontentsline{toc}{chapter}{II.3 Resources}

\section*{II.3.1 Australian National Healthcare Directory \& Identifier Specifications}
\phantomsection
\label{navigation:II.3.1}
\addcontentsline{toc}{section}{II.3.1 Australian National Healthcare Directory \& Identifier Specifications}

\label{element:eac71c56624d43bda2b9d3648f14603a}
\index{Australian National Healthcare Directory \& Identifier Specifications (II.3.1)}
{\small Resource · Canonical · current}\par

\par\medskip\textbf{Strategic Definition}\par


\textbf{Australian National Healthcare Directory \& Identifier Specifications} comprise the jurisdictional normative specifications, identifier schemes, and federation protocols defining how healthcare providers, organisations, and service delivery endpoints are identified, addressed, and verified across Australian digital health.



\par\medskip\textbf{Enduring Strategic Role \& Distinct Facets}\par


Harmonia depends upon external national specifications to participate in federated healthcare discovery and directed clinical message routing. To maintain architectural clarity, Harmonia distinguishes the complementary facets of this resource:
1. \textbf{National Healthcare Directory \& Locator Specifications}:
   - Normative specifications governing federated directory lookup, provider service discovery, and electronic service delivery endpoint resolution (e.g., Australian National Directory specifications, Endpoint Directory Service, and Location Directory Service).
   - Enables Harmonia to resolve destination endpoints for clinical communications without relying on brittle point-to-point configurations.
2. \textbf{Healthcare Identifier (HI) Ecosystem Specifications}:
   - Normative identifier schemes, metadata structures, and verification service specifications for Australian healthcare entities:
     - \textbf{HPI-I}: Healthcare Provider Identifier - Individual (practitioners).
     - \textbf{HPI-O}: Healthcare Provider Identifier - Organisation (health service organisations).
     - \textbf{IHI}: Individual Healthcare Identifier (patients / healthcare recipients).
   - Establishes jurisdictional identity anchors that Harmonia must correlate across diverse regional identifiers (e.g., hospital MRNs, practice local IDs).



\par\medskip\textbf{Traceability}\par


\begin{itemize}


\item \textbf{Architectural Axioms}: \texttt{\seqsplit{AX-01}} (Health-Information Centric), \texttt{\seqsplit{AX-06}} (Explicit Authority), \texttt{\seqsplit{AX-14}} (Preserve Distinctions).

\item \textbf{Enterprise Capabilities}: \texttt{\seqsplit{EC-01}} (Managed Entity \& Relationship), \texttt{\seqsplit{EC-08}} (Interoperability \& Exchange).

\item \textbf{Courses of Action}: \texttt{\seqsplit{COA-01}} (Boundary Membrane Sovereignty), \texttt{\seqsplit{COA-05}} (Entity-Centred Operational Coordination).

\end{itemize}

{\small Recorded unresolved matters: canonicalName; exact scope remains in the canonical record.}\par
{\footnotesize \href{../../../../continuum/strategy/resource/eac71c56624d43bda2b9d3648f14603a.md}{Canonical record: eac71c56624d43bda2b9d3648f14603a}; full qualifications and provenance.}\par
\section*{II.3.2 Healthcare Interoperability Standards and Specifications}
\phantomsection
\label{navigation:II.3.2}
\addcontentsline{toc}{section}{II.3.2 Healthcare Interoperability Standards and Specifications}

\label{element:4d725c804fb34dfa90486b21118bab8e}
\index{Healthcare Interoperability Standards and Specifications (II.3.2)}
{\small Resource · Canonical · current}\par

\par\medskip\textbf{Strategic Definition}\par


\textbf{Healthcare Interoperability Standards and Specifications} comprise the external, open, normative specifications governing healthcare data structures, clinical resource models, exchange semantics, and conformance criteria across enterprise and regional boundaries.



\par\medskip\textbf{Enduring Strategic Role}\par


Harmonia is architected to facilitate vendor-neutral health data exchange across acute, primary, community, and diagnostic healthcare systems. This interoperability cannot be accomplished using proprietary vendor protocols. Harmonia depends directly on consensus-based, open standards to establish:
- Standardised structural and semantic representations for clinical entities (encounters, observations, conditions, medications, diagnostic reports).
- Conformance validation frameworks to assert that emitted and ingested payloads comply with normative specifications.
- Standards-conformant external interaction contracts across the Harmonia boundary.



\par\medskip\textbf{Current Strategic Instance}\par


\begin{itemize}


\item \textbf{HL7® FHIR® Release 5 Core \& Australian Profile Specifications}: The principal contemporary instance of this resource is HL7 FHIR Release 5, alongside Australian jurisdictional profile specifications (AU Base and AU Core).

\item \textbf{Technology Independence}: The strategic resource is the \textbf{specification and standard itself}, not any specific implementation framework or runtime library (such as HAPI FHIR). The strategic dependency endures regardless of which runtime software serialises or deserialises the specifications.

\end{itemize}


\par\medskip\textbf{Traceability}\par


\begin{itemize}


\item \textbf{Architectural Axioms}: \texttt{\seqsplit{AX-02}} (Boundary Membrane), \texttt{\seqsplit{AX-03}} (Native Standards), \texttt{\seqsplit{AX-13}} (Explicit Management Boundary).

\item \textbf{Enterprise Capabilities}: \texttt{\seqsplit{EC-08}} (Interoperability \& Exchange), \texttt{\seqsplit{EC-13}} (Semantic Governance \& Conformance).

\item \textbf{Courses of Action}: \texttt{\seqsplit{COA-01}} (Boundary Membrane Sovereignty).

\end{itemize}

{\small Recorded unresolved matters: canonicalName; exact scope remains in the canonical record.}\par
{\footnotesize \href{../../../../continuum/strategy/resource/4d725c804fb34dfa90486b21118bab8e.md}{Canonical record: 4d725c804fb34dfa90486b21118bab8e}; full qualifications and provenance.}\par
\section*{II.3.3 National Clinical Terminology Assets (SNOMED CT-AU / AMT)}
\phantomsection
\label{navigation:II.3.3}
\addcontentsline{toc}{section}{II.3.3 National Clinical Terminology Assets (SNOMED CT-AU / AMT)}

\label{element:430f641e40aa4bd4a8e85a5cbc1f93be}
\index{National Clinical Terminology Assets (SNOMED CT-AU / AMT) (II.3.3)}
{\small Resource · Canonical · current}\par

\par\medskip\textbf{Strategic Definition}\par


\textbf{National Clinical Terminology Assets} comprise the authoritative, governed clinical vocabularies, concept identifiers, descriptions, and semantic relationships used to capture clinical meaning unambiguously across healthcare systems.



\par\medskip\textbf{Enduring Strategic Role}\par


Harmonia requires authoritative clinical vocabularies to govern semantic meaning during data ingestion, correlation, and egress:
- \textbf{SNOMED CT-AU}: The Australian extension to SNOMED CT, providing the clinical concept reference hierarchy for clinical findings, diagnoses, procedures, and body structures.
- \textbf{Australian Medicines Terminology (AMT)}: The national standard terminology for medicinal products, supporting unambiguous medication reconciliation, prescribing concepts, and dispensing data.



These terminology assets enable Harmonia to evaluate clinical semantic equivalence, enforce conformance rules, and execute meaning-preserving transformations without corrupting clinical intent.



\emph{(Note on LOINC: While Logical Observation Identifiers Names and Codes (LOINC) is a recognised international standard for laboratory and diagnostic observations, its formal admission as an enduring foundational resource is governed by jurisdictional adoption and is not pre-empted ahead of explicit canonical scope).}



\par\medskip\textbf{Traceability}\par


\begin{itemize}


\item \textbf{Architectural Axioms}: \texttt{\seqsplit{AX-04}} (Own Semantics / Use Machinery), \texttt{\seqsplit{AX-08}} (Meaning over Machinery), \texttt{\seqsplit{AX-14}} (Preserve Distinctions).

\item \textbf{Enterprise Capabilities}: \texttt{\seqsplit{EC-13}} (Semantic Governance \& Conformance).

\item \textbf{Courses of Action}: \texttt{\seqsplit{COA-03}} (Meaning-Centric Provenance and Traceability).

\end{itemize}

{\small Recorded unresolved matters: canonicalName; exact scope remains in the canonical record.}\par
{\footnotesize \href{../../../../continuum/strategy/resource/430f641e40aa4bd4a8e85a5cbc1f93be.md}{Canonical record: 430f641e40aa4bd4a8e85a5cbc1f93be}; full qualifications and provenance.}\par
